% proofs.tex -- written proofs of Lean-checked results, Stage 1 (symmetric point sets).
% Build: pdflatex -interaction=nonstopmode -halt-on-error proofs  (three passes).
\documentclass[10pt]{amsart}
% preamble.tex -- shared preamble of the written proofs of the Lean-checked results.
% Every section file is written for this preamble; do not define new packages or macros inside a section file.
\usepackage[margin=0.9in]{geometry}
\usepackage[T1]{fontenc}
\usepackage{lmodern}
\usepackage{amsmath,amssymb,amsthm,mathtools}
\usepackage{booktabs,array,longtable}
\usepackage{enumitem}
\usepackage{microtype}
\usepackage{url}
\usepackage[hidelinks]{hyperref}

\setlist{itemsep=0pt,topsep=2pt,parsep=0pt}
\emergencystretch=2em
\def\UrlBreaks{\do\.\do\/\do\_}

% Statement heads carry long Lean names. amsthm sets the head in a box and unboxes it into the first
% paragraph; a high tolerance inside every statement lets TeX use the break points of the head.
\makeatletter
\def\es@sloppy{\tolerance 9999 \hfuzz .5\p@ \ifdim\emergencystretch<3em \emergencystretch 3em\fi}
\g@addto@macro\th@plain{\es@sloppy}
\g@addto@macro\th@definition{\es@sloppy}
\g@addto@macro\th@remark{\es@sloppy}

% A run-in heading (\subsection, \paragraph) followed directly by a statement, a list or a proof is set as a
% one-line paragraph, after which the list adds \@beginparpenalty (negative), an invitation to break the page
% below the heading. Forbid that break; and until the next paragraph starts, let a \label (or another command
% that ends with \@esphack) in vertical mode add \nobreak after its \write, which would otherwise be a break point.
\newif\ifes@keephead
\let\es@trivlist@orig\@trivlist
\def\@trivlist{\if@noskipsec
    \@beginparpenalty\@M\es@trivlist@orig\global\es@keepheadtrue
  \else\es@trivlist@orig\fi}
\let\es@esphack@orig\@esphack
\def\@esphack{\es@esphack@orig\ifvmode\ifes@keephead\nobreak\fi\fi}
\AddToHook{para/begin}{\global\es@keepheadfalse}

% Paragraph headings: bold and run in. A \paragraph that directly follows a run-in heading keeps that heading
% on its page (no \@secpenalty between them).
\def\paragraph{\if@noskipsec\leavevmode\par\nobreak\@nobreaktrue\fi
  \@startsection{paragraph}{4}\z@\z@{-\fontdimen2\font}{\normalfont\bfseries}}
\makeatother
% Hyphenation exceptions (not macros): the default patterns would break ``namespace'' as ``names-pace''.
\hyphenation{name-space name-spaces}

\newtheorem{theorem}{Theorem}[section]
\newtheorem{lemma}[theorem]{Lemma}
\newtheorem{proposition}[theorem]{Proposition}
\newtheorem{corollary}[theorem]{Corollary}
\newtheorem{fact}[theorem]{Fact}
\theoremstyle{definition}
\newtheorem{definition}[theorem]{Definition}
\theoremstyle{remark}
\newtheorem{remark}[theorem]{Remark}

% Lean declaration names, directory and file names: typewriter text that may break after . / _ and, less readily,
% after : (never after the last character, so that a following punctuation mark stays on the line). The break points
% are explicit penalties, so they also work in statement heads, which amsthm builds in a box (the math-mode break
% points of \nolinkurl are lost there). \kern0pt after each character prevents ligatures and hyphenation. The
% arguments are plain names (letters, digits and . / _ : - ' only).
\makeatletter
\def\es@ttpen{0}
\def\es@ttcolpen{100}
\def\es@tt#1{{\upshape\ttfamily\let\es@ttpending\relax\expandafter\es@ttloop\detokenize{#1}\es@ttstop}}
\def\es@ttloop#1{\ifx\es@ttstop#1\else\es@ttpending\es@ttchar{#1}\expandafter\es@ttloop\fi}
\def\es@ttchar#1{#1\kern\z@
  \if.#1\def\es@ttpending{\penalty\es@ttpen\relax}\else
  \if/#1\def\es@ttpending{\penalty\es@ttpen\relax}\else
  \if_#1\def\es@ttpending{\penalty\es@ttpen\relax}\else
  \if:#1\def\es@ttpending{\penalty\es@ttcolpen\relax}\else
  \let\es@ttpending\relax\fi\fi\fi\fi}
\let\es@ttstop\relax
\DeclareRobustCommand{\lean}[1]{\ifmmode\text{\es@tt{#1}}\else\es@tt{#1}\fi}
\DeclareRobustCommand{\dirn}[1]{\ifmmode\text{\es@tt{#1}}\else\es@tt{#1}\fi}
% For the statement map (appendix): an underscore in roman text becomes a permitted line break, as in identifiers.
\newcommand{\esmapbreaks}{\protected\def\_{\ifmmode\nfss@text{\textunderscore}\else\textunderscore\allowbreak\fi}}
\makeatother
% The source of a statement: Lean name and file:line, printed after the statement head.
\newcommand{\leanref}[2]{Lean: \lean{#1}, \dirn{#2}}

% Geometry and combinatorics.
\newcommand{\orient}{\operatorname{orient}}
\newcommand{\conv}{\operatorname{conv}}
\newcommand{\hv}{\operatorname{hv}}
\newcommand{\peel}{\operatorname{peel}}
\newcommand{\sgn}{\operatorname{sgn}}
\newcommand{\R}{\mathbb R}
\newcommand{\Z}{\mathbb Z}
\newcommand{\N}{\mathbb N}
\newcommand{\C}{\mathbb C}
\newcommand{\abs}[1]{\lvert#1\rvert}
\newcommand{\card}[1]{\lvert#1\rvert}

\hypersetup{pdftitle={Written proofs of Lean-checked results, Stage 1: symmetric point sets}}

\title[Written proofs of Lean-checked results, Stage~1: symmetric point sets]{Written proofs of Lean-checked
  results,\\ Stage~1: symmetric point sets}
\date{8 November 2026}

\begin{document}

\begin{abstract}
The Erd\H{o}s--Szekeres project checks some of its results in Lean~4. This document is Stage~1 of a series of written
proofs of these results. It treats the Lean files listed under ``Scope of Stage~1'' below, which concern mainly
symmetric point sets, and the foundations they use; other Lean-checked results of the project, among them further
results on symmetric point sets, are left to later stages. Each section restates the Lean-checked theorems of one
part, faithfully to the Lean declarations, and gives a written proof that a mathematician can check without Lean,
with two exceptions. Two published theorems that the Lean code uses, Baek's theorem and the cups--caps theorem of
Erd\H{o}s and Szekeres, are cited and not reproved. Statements marked computer-assisted are reduced to a finite
computation that is stated precisely and carried out by a machine; it is not carried out by hand, except where a
section gives a proof by hand as well. The Lean checks remain the primary evidence for these theorems; the proofs
here are written proofs and have not been refereed. Where a Lean theorem takes unproved propositions or
satisfiability verdicts as hypotheses, it is stated with those hypotheses and is established here only as that
conditional statement.
\end{abstract}

\maketitle
\tableofcontents

\phantomsection
\section*{Scope of Stage~1}

Paths follow the convention of Section~\ref{fd:sec}: a path that begins with \dirn{ErdosSzekeres/}, \dirn{docs/},
\dirn{experiments/} or \dirn{scripts/} is relative to the root of the project repository, and every other path is
relative to its research-records directory \dirn{cloud/es7-2026-10-19/results}.

\paragraph{Treated here.}
The sections restate the Lean-checked results of the following files and give written proofs of them. They also
restate declarations of the main library \dirn{ErdosSzekeres/} that these files use: the foundations
(Section~\ref{fd:sec}), and the formalizations of Baek's theorem and of the cups--caps theorem with their bridge files
(Section~\ref{mb:sec}); a few such declarations are also restated in Sections~\ref{sr:sec} and~\ref{mm:sec}. These
four sections list the main-library files concerned, with their hashes.
\begin{itemize}
\item \dirn{LAYERLEAN/LayerLemma.lean}, mainly in Sections~\ref{ll:sec} and~\ref{lg:sec};
\item the subdirectories \dirn{orbit}, \dirn{merged}, \dirn{keys}, \dirn{phase}, \dirn{linked}, \dirn{capstone} and
  \dirn{es7corr} of \dirn{LEANEXT}, in Sections~\ref{oc:sec}, \ref{pc:sec} and~\ref{ec:sec};
\item \dirn{ES8SYM/lean} and \dirn{PRIMEFOLD/compat/lean}, in Section~\ref{ec:sec};
\item \dirn{SYMLEAN/SymmetricReduction.lean} and \dirn{SYMLEAN/AffineReduction.lean}, in Sections~\ref{sr:sec}
  and~\ref{ar:sec}, and \dirn{SYMLEAN/mirror/ReductionThreeClasses.lean}, in Section~\ref{mb:sec};
\item \dirn{MIRRORETV/lean}, in Sections~\ref{mm:sec} and~\ref{mb:sec};
\item \dirn{C3LEAN} and \dirn{C2FINAL} (the declarations of the latter are Lean-checked according to the record), in
  Section~\ref{cb:sec};
\item \dirn{AUDIT3/A1_reduction} and \dirn{ES7SPLIT/review/equivalence}, in Section~\ref{sa:sec}.
\end{itemize}
Not restated, although they belong to these records: verbatim copies, tests and controls, \texttt{\#print} and
\texttt{\#check} commands, the review files in \dirn{LEANEXT/review}, \dirn{SYMLEAN/review} and
\dirn{MIRRORETV/review} (re-checks, independent formalizations and tests; a few of their theorems are mentioned in
remarks), \dirn{MIRRORETV/lean/FarFlat.lean} (estimates for a lower-bound construction) and a few further
declarations. The Coverage paragraph at the end of each section gives the details and the reasons. The Lean files of
\dirn{C3L} (a generator of CNF formulas) and of \dirn{AUDIT2} (audit scripts) state no theorems and are not treated.

\paragraph{Left to later stages.}
The following Lean-checked work of the project is not treated in this document. Planned for Stage~2, the SAT
encodings:
\begin{itemize}
\item \dirn{ENCBLEAN}: encoder~B of the sixfold case of the $k=8$ problem as a Lean program, with its soundness, and
  the geometric assembly in \dirn{ENCBLEAN/geo}, which derives its copy of \lean{Sixfold66}
  (Definition~\ref{ll:def:66-70}) from the unsatisfiability of $64$ cube formulas
  (\lean{EncB.Geo.sixfold66_of_cubes64}; see Remark~\ref{ll:rem:hyp});
\item \dirn{H} and \dirn{H2}: Lean versions of clause generators of the project's SAT encodings, among them
  mirror-symmetric and $D_3$-symmetric ones, with their soundness;
\item \dirn{H3}, \dirn{V1-L}, \dirn{V2-L} and \dirn{V2-L2}: the derivation of the mirror statement
  \lean{Mir.MirrorExclusionLine} from hypotheses asserting that certain CNF formulas are unsatisfiable
  (Remark~\ref{sr:rem:mirror-line}).
\end{itemize}
Planned for Stage~3:
\begin{itemize}
\item \dirn{experiments/omission-compact-checker-2026-11-05}: the compact checker of omission certificates, with its
  Lean check of $\mathrm{ES}(8)\le311$;
\item the main library beyond the declarations treated here.
\end{itemize}
The Lean files of \dirn{ES8SLACK} are not treated here; written proofs of its lemmas are in
\dirn{ES8SLACK/proofs.md}.

\paragraph{Evidence outside Lean.}
Some hypotheses of conditional theorems are supported in the project by SAT certificates. For these the project
counts only the verdicts of the proof checker cake\_lpr~\cite{CakeLPR}; the checker lrat-check of
drat-trim~\cite{DT} is a secondary check, whose verdicts are not counted (\dirn{WRITEUP/SYMMETRIC.md}, \S3, and
\dirn{LRATCHECK20261105}). None of this evidence is checked in this document.

\paragraph{Works cited.}
The sections cite published works inline, by author together with title or arXiv number, by author alone or by
name; the References list them. They are: Lean~4
\cite{Lean4} and Mathlib \cite{Mathlib}; textbooks for standard facts \cite{Roc70,Rud76}; the theorems of
Baek~\cite{Bae25} and of Erd\H{o}s and Szekeres~\cite{ES35}, which are cited and not reproved
(Theorems~\ref{mb:thm:baek} and~\ref{mb:thm:capcup}); earlier work with priority for the fivefold case
(Dumitru~\cite{Dum25}) and for the fourfold case (Subercaseaux, Mackey, Qian and Heule~\cite{SMQH25}; see
Remark~\ref{sr:rem:records}); and the archive of M.~Schroeder~\cite{Sch26}, from which six of the bridge files of
Section~\ref{mb:sec} are adapted.

\section{Conventions and foundations}\label{fd:sec}

\paragraph{Purpose of this document.}
This document collects human-readable proofs of results of a research project on the
Erd\H{o}s--Szekeres problem whose formal versions have been checked with the interactive theorem prover
Lean~4 and its mathematical library Mathlib. The Lean checks remain the primary evidence for these results. The
written proofs are meant to make the arguments accessible to a mathematician who does not read Lean:
every statement that corresponds to a Lean declaration is stated exactly as in Lean, with the name of
the declaration and its file and line, and it receives a complete proof that can be checked without
Lean. There are two exceptions. A statement labelled computer-assisted is reduced to a finite
computation, which is stated precisely and need not be carried out by hand; where this document also
gives a complete written proof by hand, the label is qualified as ``computer-assisted in Lean'' (see the
rigor labels below). A few
published theorems that the Lean code uses (for example Baek's theorem, Theorem~\ref{mb:thm:baek}) are
cited in the exact form in which they are used and are not reproved. A convenient variant of a Lean
statement, for example one in a different generality, is labelled as a variant. These written proofs
have not been refereed. Where a Lean theorem assumes propositions that are not proved in Lean (for
example verdicts of SAT solvers), the statement says so and the theorem is presented as a conditional
statement.

\paragraph{Files, commit and versions.}
All files belong to the project repository at commit \dirn{7e32ee259} (branch
\dirn{es7-night-2026-11-01}; Lean package \dirn{erdos_szekeres}). A path that begins with
\dirn{ErdosSzekeres/}, \dirn{docs/}, \dirn{experiments/} or \dirn{scripts/} is relative to the
repository root, and so are the files \dirn{lean-toolchain}, \dirn{lake-manifest.json} and
\dirn{lakefile.toml} at the top of the repository. The directory \dirn{ErdosSzekeres/} is the \emph{main
library}, the default build target of the package. Every other path, for example
\dirn{LAYERLEAN/LayerLemma.lean:298}, is relative to the research-records directory, whose path from the
repository root is \dirn{cloud/es7-2026-10-19/results}. The Lean files of the research records are not
part of the build target; they are checked separately, for example with \texttt{lake env lean} run from
the repository root (as \dirn{LAYERLEAN/logs/axioms_recheck_20261106/README.txt} records for a copy of
\dirn{LAYERLEAN/LayerLemma.lean} followed by \texttt{\#print axioms} commands), and each section names the
build and axioms logs it relies on. After a section has cited a file with its full path, it may refer to
it by its bare file name. A suffix \dirn{:N} is a line number; for a declaration it is the line that
carries the declaration keyword (\lean{theorem}, \lean{lemma}, \lean{def}, \lean{abbrev},
\lean{structure} or \lean{inductive}, possibly preceded by an attribute such as \texttt{@[simp]} or a
modifier such as \lean{noncomputable} or \lean{private}). A bare \dirn{:N} written after a file has been
cited refers to line~$N$ of that file. Lean declarations
are given by their full names in statement heads, for
example \lean{ErdosSzekeres.Suk.convexPosition_of_four_subsets}; in running text we sometimes omit the
common prefix \lean{ErdosSzekeres.} (so \lean{Suk.orient_reflectY} is
\lean{ErdosSzekeres.Suk.orient_reflectY}). Names in a namespace outside \lean{ErdosSzekeres}, such as
\lean{MirrorOdd.hasConvexNGon_of_le} in the top-level namespace \lean{MirrorOdd}, are always written in
full. Mathlib declarations are named as in Mathlib. Lean expressions are quoted in typewriter type, for
example \texttt{Finset Point}.
The Lean files were checked with Lean 4.24.0 (file \dirn{lean-toolchain}) and Mathlib at revision
\dirn{f897ebcf} (full hash \dirn{f897ebcf72cd16f89ab4577d0c826cd14afaafc7}, file
\dirn{lake-manifest.json}).

\paragraph{Rigor labels.}
We use exactly three labels; the first and the third also occur in a qualified form, defined below.
\begin{itemize}
\item \emph{Lean-checked}: the statement is a Lean declaration that Lean~4.24.0 with Mathlib
  \dirn{f897ebcf} accepted, using no axioms other than the three standard axioms \lean{propext},
  \lean{Classical.choice} and \lean{Quot.sound}; for a declaration whose proof uses Lean's
  \texttt{native\_decide}, directly or through a declaration on which it depends, the two further axioms
  \lean{Lean.ofReduceBool} and \lean{Lean.trustCompiler} that \texttt{native\_decide} adds are allowed as
  well, and the section that cites such a declaration names them (no declaration cited in this section is
  of this kind). The evidence is the saved output of
  Lean's command \texttt{\#print axioms}, for the declaration itself or for a declaration whose statement
  or proof uses it (this command reports every axiom in the dependency closure, which contains the
  declarations used in statements as well as in proofs). Each section names the build and axioms logs
  on which this label rests. Where a declaration is not covered by such an output, the section says so
  and names the evidence used instead, which is of one of two kinds:
  \begin{enumerate}[label=(\alph*)]
  \item evidence that the current version of the file that contains the declaration compiled without
    errors (exit status~$0$), namely a saved compile or build log of that file, or a project record of
    such a build when saved Lean output of a build that includes the file (for example an axiom audit)
    exists as well; together with a text search of the file and of the project files it imports that
    finds no \texttt{sorry} (a declaration that uses \texttt{sorry} compiles with a warning only, so the
    search, not the exit status, excludes it), no \texttt{admit}, no \texttt{axiom} declaration, no
    \texttt{native\_decide}, no \texttt{+native} (as in \texttt{decide +native}, of which
    \texttt{native\_decide} is a synonym), no \lean{Lean.ofReduceBool} or \lean{Lean.ofReduceNat} (the
    axioms through which values computed by Lean's compiler or interpreter enter a proof) and no
    \lean{debug.skipKernelTC} (an option that skips the kernel's type check);
  \item a \texttt{\#guard\_msgs} command, compiled in a build whose log is saved, that pins the output of
    \texttt{\#print axioms} for the declaration or for a declaration whose statement or proof uses it.
  \end{enumerate}
  A text search of every Lean file of the repository that this document cites, and of the project files
  these import, finds none of the strings \texttt{+native}, \texttt{ofReduceBool}, \texttt{ofReduceNat}
  and \texttt{skipKernelTC}, so every application of~(a) in this document remains valid, also where it
  reports only a search for \texttt{sorry}, \texttt{admit}, \texttt{axiom} declarations and
  \texttt{native\_decide}.
  If no Lean output of any kind (compile or build log, axiom output, pinned message) is saved for a
  build that includes a file, and the label of its declarations rests only on the statement of a project
  record such as a README, we write \emph{Lean-checked according to the record}.
\item \emph{Written proof (not refereed)}: every proof in this document. It is meant to be checkable by a
  mathematician without Lean; it has not been refereed.
\item \emph{Computer-assisted}: a statement that rests on a finite computation carried out by a machine.
  This is the case, first, when the Lean proof of the statement, or of one of its parts, is itself a
  finite computation: that statement or part is decided by kernel evaluation of a decision procedure (Lean's
  \texttt{decide} or \texttt{decide +kernel}, which add no axiom) or by compiled evaluation (Lean's
  \texttt{native\_decide}, which adds the axioms \lean{Lean.ofReduceBool} and \lean{Lean.trustCompiler}
  and trusts the compiler), or it is obtained by evaluating generated data or by replaying a certificate
  produced by an external program (for example an LRAT refutation found by a SAT solver, replayed in Lean
  and checked by the kernel). It is the case, second, when the written proof of the statement uses a
  machine computation that it does not carry out by hand, for example by applying a statement of the
  first kind that has no written proof by hand. For a computer-assisted statement we give the
  mathematical reduction to the computation and describe the computation precisely; we do not claim to
  have carried out by hand a computation for which no written proof by hand is given. If the Lean proof
  of a statement is a finite computation and this document also gives a complete written proof of the
  statement that does not use the computation, the statement is labelled \emph{computer-assisted in
  Lean}, usually with a pointer to the written proof (as in ``computer-assisted in Lean; written proof by
  hand''); that written proof can be checked without a machine. A Lean proof that uses \texttt{decide}
  only for side facts inside the proof of a statement that is not itself a finite computation (such as
  $0<1$ in $\mathrm{Fin}\,4$, or closed instances that the written proof checks by hand) does not make the
  statement computer-assisted. Certificates may also be checked by programs outside Lean. For the SAT
  certificates of the project only the checker \texttt{cake\_lpr} counts, and a verdict of it counts only
  if it prints the line \texttt{s VERIFIED UNSAT}; verdicts of \texttt{lrat-check} alone are not counted
  (\dirn{LRATCHECK20261105/README.md}; \dirn{WRITEUP/SYMMETRIC.md}, section~3). A section may report that
  a result of the project's records, such as an unsatisfiability verdict, is computer-assisted; this is a
  claim of the records, which this document does not prove, and a Lean theorem that takes such a result
  as a hypothesis is presented as conditional on it.
\end{itemize}
A statement, or a part of a statement, whose head names a declaration of the project in the form
``Lean: name, file:line'' is Lean-checked in this sense. For a definition this records only that Lean
accepted the definition: a definition asserts nothing, and what this document claims for it is that the
written definition restates the Lean definition. A part marked ``written step'' does not carry this
label, even where the text names a Lean counterpart of it elsewhere, and a Mathlib declaration named in a
head only identifies the Mathlib counterpart of a statement that is proved here. A Lean theorem that
takes unproved propositions as hypotheses is Lean-checked only as the conditional statement it is, and
the hypotheses are named. We use no other words of this kind (such as ``verified'' or ``certified'') for
the results of this document.

\paragraph{Lean's arithmetic, written explicitly.}
Lean statements often use natural numbers inside real expressions and subtract natural numbers. We make
both explicit. A natural number $m$ used as a real number is written $(m:\R)$ whenever the cast matters;
for example the Lean angle \texttt{2 * Real.pi / m} with \lean{m} a natural number is $2\pi/(m:\R)$. Real
division in Lean is total, with $x/0=0$; hence $2\pi/(m:\R)=0$ when $m=0$. Subtraction of natural
numbers in Lean is truncated ($a-b=0$ when $a<b$); whenever a Lean statement subtracts natural numbers we
write the difference as $\max(a-b,0)$, or as $a\mathbin{\dot-}b=\max(a-b,0)$ where a section defines this
symbol, or we state the inequality $b\le a$ that makes it an ordinary difference.
$\mathrm{Fin}\,n=\{0,1,\dots,n-1\}$ carries its natural order. For a map $f$ we write $f^{[n]}$ for its
$n$-th iterate ($f^{[0]}=\mathrm{id}$, $f^{[n+1]}=f^{[n]}\circ f$), as Lean does; the recursion
$f^{[n+1]}=f\circ f^{[n]}$ defines the same maps (see the proof of Lemma~\ref{fd:lem:inv-image}(c)).
Mathlib's \lean{Function.minimalPeriod} $f$ $x$ is the least $n>0$ with $f^{[n]}(x)=x$, or $0$ if there
is no such $n$.

\paragraph{Default notation.}
Points, finite sets and $\card S$ are as in Definition~\ref{fd:def:point}; $\orient$ is
Definition~\ref{fd:def:orient}; $\conv$ is the convex hull of Definition~\ref{fd:def:conv}. For a finite
set $X$ of points, $\hv(X)=\{x\in X: x\notin\conv(X\setminus\{x\})\}$ is the set of hull vertices
(Lean \lean{hullVerts}) and $\peel(X,0)=X$, $\peel(X,n+1)=\peel(X,n)\setminus\hv(\peel(X,n))$
(Lean \lean{peel}); these two notions are defined formally, with their Lean names, in
Definitions~\ref{ll:def:hv} and~\ref{ll:def:peel}. ``$S$ is $K$-free'' means
$\neg\,\mathrm{HasConvexNGon}\,K\,S$ (Definition~\ref{fd:def:ngon}). We write $\rho_{c,\theta}$ for the
rotation $\mathrm{rotAbout}\,c\,\theta$ of Definition~\ref{fd:def:rot}; rotations act counterclockwise
(Lemma~\ref{fd:lem:rot}(e)). In this section $R_\theta$ always denotes the rotation matrix of
Lemma~\ref{fd:lem:rot}, not a map. The notation of this paragraph and of the previous one is the
default for the whole document. A later section may write the same objects differently, for example
$\mathrm{rotAbout}(c,\theta)$ or $\mathrm{rotAbout}_{c,\theta}$ for $\rho_{c,\theta}$, a letter such as
$R$ with subscripts for the rotation map itself and another letter for its matrix, $f^n$ for $f^{[n]}$,
or $\mathrm{HasConvexNGon}$ in another typeface; such a section defines its symbols where it introduces
them.

\paragraph{This section.}
We define points, the orientation function, general position, convex position, convex $n$-gons,
$K$-freeness, forcing, invariance and rotations exactly as the Lean code does, and we prove the basic
facts that all later sections use: the symmetries of $\orient$ and its behaviour under affine maps; the
equivalence of $\orient(a,b,c)=0$ with collinearity; the hull form of convex position; monotonicity;
invariance of convex and general position under injective affine maps; Carath\'eodory's theorem in the
plane; the four-point criteria for convex position; the criterion that points whose increasing triples
are all counterclockwise are in convex position; and the separation of a point from a finite set. This
section relies on no other section. Its sources, with the first 12 hexadecimal digits of their SHA-256
hashes, are listed in the following table, the table of sources.
\begin{longtable}{ll}
\toprule
file & SHA-256 (prefix)\\
\midrule
\endfirsthead
\toprule
file & SHA-256 (prefix)\\
\midrule
\endhead
\bottomrule
\endfoot
\dirn{ErdosSzekeres/Definitions.lean} & \texttt{2425500eefd1}\\
\dirn{ErdosSzekeres/Orientation.lean} & \texttt{f67955ac7091}\\
\dirn{ErdosSzekeres/OrientationGeometry.lean} & \texttt{c7db9bbbb65e}\\
\dirn{ErdosSzekeres/Basic.lean} & \texttt{13667aa2d3dc}\\
\dirn{ErdosSzekeres/Suk/FourPoint.lean} & \texttt{f3783fe9766f}\\
\dirn{ErdosSzekeres/Suk/SupportingLines.lean} & \texttt{6a8e12677685}\\
\dirn{ErdosSzekeres/Suk/OrientationClique.lean} & \texttt{cf63d5219b47}\\
\dirn{ErdosSzekeres/ConvexCycleCuts.lean} & \texttt{b4ad0ba9e583}\\
\dirn{ErdosSzekeres/Suk/Reflection.lean} & \texttt{8b99cbf5fa2b}\\
\dirn{LAYERLEAN/LayerLemma.lean} & \texttt{afb8542dceb0}\\
\dirn{SYMLEAN/SymmetricReduction.lean} & \texttt{6d5ec90a4215}\\
\dirn{SYMLEAN/AffineReduction.lean} & \texttt{41cfc0b60aec}\\
\dirn{LEANEXT/orbit/OrbitConsistency.lean} & \texttt{9b843180ac38}\\
\dirn{MIRRORETV/lean/MirrorMain.lean} & \texttt{ca623e55a1b9}\\
\dirn{SYMLEAN/mirror/ReductionThreeClasses.lean} & \texttt{0cc83994fe8c}\\
\end{longtable}

\paragraph{Axioms evidence.}
{\raggedright
In this paragraph, ``line~$N$'' without a file name refers to the log, audit or report named last
before it.
A declaration said to be used ``at'' a file and line is used on that line, inside the proof named; the
line is not the line of its own declaration.
\begin{itemize}
\item \emph{Main library, printed theorems.} The saved axiom audit
  \dirn{docs/validation/axiom-audit-lower-bound-2026-10-11.txt} reports exactly the axioms \lean{propext},
  \lean{Classical.choice}, \lean{Quot.sound} for \lean{ErdosSzekeres.hasConvexNGon_of_positive_cycle}
  (line~87), \lean{ErdosSzekeres.Suk.convexPosition_of_four_subsets} (line~103),
  \lean{ErdosSzekeres.Suk.convexPosition_of_strict_supports} (line~113),
  \lean{ErdosSzekeres.Suk.convexPosition_four_of_positive} (line~114),
  \lean{ErdosSzekeres.Suk.convexPosition_of_positive_triples} (line~119),
  \lean{ErdosSzekeres.HasConvexNGon.linearEquiv} (line~175) and
  \lean{ErdosSzekeres.Suk.generalPosition_reflectY} (line~176).
\item \emph{Main library, declarations used by these theorems.} Since \texttt{\#print axioms} reports every
  axiom used anywhere in the dependency closure of a declaration, these lines also cover the declarations
  used in those proofs: \lean{Suk.convexPosition_of_subsets_card_le_four}, \lean{convexPosition_iff},
  \lean{ConvexPosition.mono}, \lean{Suk.convex_orient_positive}, \lean{Suk.orient_affine_combination},
  \lean{Suk.orient_support_sum}, \lean{Suk.reflectY}, \lean{Suk.orient_reflectY},
  \lean{generalPosition_iff_orient_ne_zero} and \lean{orient_eq_zero_iff_collinear}
  (\lean{Suk.generalPosition_reflectY} uses \lean{generalPosition_iff_orient_ne_zero} at
  \dirn{ErdosSzekeres/Suk/Reflection.lean:55}, and the proof of the latter uses
  \lean{orient_eq_zero_iff_collinear}).
\item \emph{Main library, further theorems of the audit.} Other theorems printed in the same audit, with
  the same three axioms, cover further declarations through their statements or proofs:
  \begin{itemize}
  \item line~725, \lean{conjecturedBound_le_of_forces}, uses \lean{ForcesConvexNGon.mono} at
    \dirn{ErdosSzekeres/LowerBound/Construction.lean:118};
  \item line~732, \lean{esNumber_three} (declared at \dirn{ErdosSzekeres/LowerBound/Number.lean:40}), uses
    \lean{forces_triangle}, whose proof uses \lean{forcesConvexNGon_iff_exact} and
    \lean{convexPosition_of_card_three} at \dirn{ErdosSzekeres/Basic.lean:73} and \dirn{:75}; the proof
    of \lean{forcesConvexNGon_iff_exact} uses \lean{GeneralPosition.mono} and \lean{HasConvexNGon.mono} at
    \dirn{ErdosSzekeres/Basic.lean:45};
  \item line~720, \lean{LowerBound.exists_cupCapFree}, uses \lean{LowerBound.generalPosition_translate}
    at \dirn{ErdosSzekeres/LowerBound/Construction.lean:54}, whose proof uses \lean{orient_translate} at
    \dirn{ErdosSzekeres/LowerBound/Placement.lean:20};
  \item line~624, \lean{Omission.Configuration.orient_ne_of_distinct}, uses \lean{orient_swap} and
    \lean{orient_cycle} at \dirn{ErdosSzekeres/OrientationObstructions.lean:64} and \dirn{:68};
  \item line~132, \lean{Suk.AdjacentCapFrames.glue_chain_caps}, uses \lean{ConvexPosition.hasConvexNGon}
    at \dirn{ErdosSzekeres/Suk/SpikeChain.lean:140};
  \item line~49, \lean{VisibleEdgeInsertion.convexPosition_insert_of_one_visible_edge}, uses
    \lean{Suk.convex_orient_nonnegative} at \dirn{ErdosSzekeres/VisibleEdgeInsertion.lean:77};
  \item line~178, \lean{Suk.horizontal_finite_density_bound}, uses \lean{Suk.hasConvexNGon_of_reflectY} at
    \dirn{ErdosSzekeres/Suk/FiniteDensityBound.lean:66};
  \item line~88 prints \lean{convexCycleCut_of_avoids}, whose statement contains \lean{ConvexCycleCut};
  \item line~529 prints \lean{Omission.BoundaryGeometry.exposed_not_in_hull} (declared at
    \dirn{ErdosSzekeres/Omission/BoundaryGeometry.lean:30}), which the Coverage paragraph names.
  \end{itemize}
  The definitions \lean{Point}, \lean{orient}, \lean{GeneralPosition}, \lean{ConvexPosition},
  \lean{HasConvexNGon} and \lean{ForcesConvexNGon} occur in the statements of the printed theorems above
  (for example in those of lines 87, 103, 176 and 725), so they are covered as well. Every main-library
  file named in this paragraph is unchanged since commit \dirn{7f0868fdb} (11~October 2026), which added
  the audit.
\item \emph{Main library, remaining declarations.} The remaining main-library declarations named in this
  section, \lean{HasConvexNGon.card_le}, \lean{Suk.convexPosition_of_negative_triples} and the three
  simplification lemmas \lean{Suk.reflectY_x}, \lean{Suk.reflectY_y}, \lean{Suk.reflectY_reflectY} of
  Remark~\ref{fd:rem:affine-cp-lean}, are not covered by a printed line that we identified. For them the
  evidence is of kind~(a) of the rigor labels: the build record \dirn{docs/VALIDATION.md:73} (on
  11~October 2026 the main library, which contains their files, builds with \texttt{lake build
  -{}-wfail}, which fails on any warning, in particular on the warning that Lean issues for a declaration
  that uses \texttt{sorry}); the audit above, saved Lean output of the same library, whose printed
  theorems include theorems of \dirn{ErdosSzekeres/Suk/OrientationClique.lean} and
  \dirn{ErdosSzekeres/Suk/Reflection.lean} and depend on declarations of \dirn{ErdosSzekeres/Basic.lean};
  and a text search of their files and of the main-library files these import
  (\dirn{ErdosSzekeres/Suk/BaseOrder.lean}, \dirn{ErdosSzekeres/Suk/BaseChart.lean},
  \dirn{ErdosSzekeres/Suk/TriangleOrder.lean}, \dirn{ErdosSzekeres/Suk/CupsCaps.lean} and the four files
  of \dirn{ErdosSzekeres/Vendor/ErdosSzekeresConvex/}, besides the files in the table of sources), which
  finds no \texttt{sorry}, \texttt{admit}, \texttt{axiom} declaration or \texttt{native\_decide}, and no
  \texttt{implemented\_by} or \texttt{extern} attribute.
\item \emph{The layer-lemma file.} The review report \dirn{LEANEXT/review/REPORT.md} (item~1) records that
  \dirn{LAYERLEAN/LayerLemma.lean}, in its current version (1148 lines, SHA-256 prefix \texttt{afb8542d}),
  compiles with exit code~0 and no messages. The re-check
  \dirn{LAYERLEAN/logs/axioms_recheck_20261106/out.txt} records exit status~0 for a file whose first 1148
  lines are byte-identical to \dirn{LAYERLEAN/LayerLemma.lean}, followed by 19 \texttt{\#print axioms}
  commands; the only Lean output in it is the 19 printed axiom lines. (The log
  \dirn{LAYERLEAN/logs/build.log} names no file hash; it was committed with the first, 454-line version of
  the file and has not been updated since, and we do not use it.) The review log
  \dirn{LEANEXT/review/axioms.log}, which prints the axioms of the theorems of
  \dirn{LAYERLEAN/LayerLemma.lean}, lists \lean{ErdosSzekeres.LayerLemma.iterate_mem} with \lean{propext}
  and \lean{Quot.sound} only (line~1) and \lean{ErdosSzekeres.LayerLemma.rotAbout_eq_crot} (line~22),
  \lean{ErdosSzekeres.LayerLemma.rotAbout_injective} (line~23),
  \lean{ErdosSzekeres.LayerLemma.hullVerts_eq_orbitFS_rotAbout} (line~26; its proof uses
  \lean{ErdosSzekeres.LayerLemma.rotAff} at \dirn{LAYERLEAN/LayerLemma.lean:368}) and
  \lean{ErdosSzekeres.LayerLemma.rotAbout_iterate} (line~33) with the three standard axioms.
\item \emph{The symmetric and affine reduction files.} \dirn{SYMLEAN/logs/runs.tsv} records that
  \dirn{SYMLEAN/SymmetricReduction.lean} (SHA-256 prefix \texttt{6d5ec90a}, the current file) compiles
  with exit code~0 and no output (line~2), and so does \dirn{SYMLEAN/AffineReduction.lean} (prefix
  \texttt{41cfc0b6}, the current file; line~4). The file \dirn{SYMLEAN/AffineReductionAxioms.lean} is
  \dirn{SYMLEAN/AffineReduction.lean} followed by nine \texttt{\#print axioms} commands and one
  \texttt{\#check} command; it compiles with exit code~0 (\dirn{SYMLEAN/logs/runs.tsv}, line~5), and its
  output \dirn{SYMLEAN/logs/AffineReductionAxioms.out} lists with the three standard axioms
  \begin{itemize}
  \item \lean{ErdosSzekeres.SymReduction.heptagon_of_nontrivial_affine_coords} (lines~9--11), whose proof
    uses \lean{ErdosSzekeres.SymReduction.hasConvexNGon_of_affine_coords_symmetry} at
    \dirn{SYMLEAN/AffineReduction.lean:1528}, whose proof in turn uses
    \lean{ErdosSzekeres.SymReduction.injective_of_affine} at \dirn{SYMLEAN/AffineReduction.lean:1462};
  \item \lean{ErdosSzekeres.SymReduction.generalPosition_image_affineEquiv} (lines~13--15), whose proof
    uses \lean{ErdosSzekeres.SymReduction.generalPosition_image_affine} at
    \dirn{SYMLEAN/AffineReduction.lean:1295}, whose proof in turn uses
    \lean{ErdosSzekeres.SymReduction.orient_affine} at \dirn{SYMLEAN/AffineReduction.lean:1281};
  \item \lean{ErdosSzekeres.SymReduction.hasConvexNGon_image_affine} (line~16), whose proof uses
    \lean{ErdosSzekeres.SymReduction.convexPosition_image_affine} at
    \dirn{SYMLEAN/AffineReduction.lean:1303}.
  \end{itemize}
\item \emph{The orbit-consistency file.} \dirn{LEANEXT/orbit/logs/build.log} records exit code~0 for the
  build of \dirn{LEANEXT/orbit/OrbitConsistency.lean}; the log contains only the exit code, and the report
  \dirn{LEANEXT/orbit/REPORT.md} that accompanies these logs identifies the checked file (line~3: 429
  lines, SHA-256 prefix \texttt{9b843180}, the current file). \dirn{LEANEXT/orbit/logs/axioms.log} lists
  the two theorems of that file named in this section,
  \lean{ErdosSzekeres.OrbitConsistency.orient_rotAbout_inv} (line~1) and
  \lean{ErdosSzekeres.OrbitConsistency.orient3_perm} (line~5), with the three standard axioms. From
  \dirn{SYMLEAN/SymmetricReduction.lean} and \dirn{LEANEXT/orbit/OrbitConsistency.lean} we otherwise cite
  only definitions.
\item \emph{The mirror files.} \lean{MirrorOdd.hasConvexNGon_of_le} is used at
  \dirn{MIRRORETV/lean/MirrorMain.lean:2520}, in the proof of \lean{MirrorOdd.mirror_odd_general}
  (declared at \dirn{MIRRORETV/lean/MirrorMain.lean:2498}), which the compile log
  \dirn{MIRRORETV/lean/MirrorMain.out} of that file (warnings, no errors) lists with the three standard
  axioms (line~858). Its verbatim copy in \dirn{SYMLEAN/mirror/ReductionThreeClasses.lean} is used in the
  same way, at \dirn{SYMLEAN/mirror/ReductionThreeClasses.lean:4079}, in the proof of the copy of
  \lean{MirrorOdd.mirror_odd_general} (declared at \dirn{SYMLEAN/mirror/ReductionThreeClasses.lean:4057}),
  which line~11 of \dirn{SYMLEAN/mirror/axioms.out} lists with the three standard axioms; and
  \dirn{SYMLEAN/mirror/README.md} records that this file compiles with exit code~0 and no errors.
\item \emph{Text search.} A text search of the six research files of the table of sources finds no
  \texttt{sorry}, \texttt{admit}, \texttt{axiom} declaration or \texttt{native\_decide}, and no
  \texttt{implemented\_by} or \texttt{extern} attribute.
\end{itemize}
\par}

%======================================================================
\subsection{Definitions}\label{fd:sec-defs}

\begin{definition}[Points and finite point sets; \leanref{ErdosSzekeres.Point}{ErdosSzekeres/Definitions.lean:9}]\label{fd:def:point}\sloppy
A \emph{point} is an element of $\R^2$. In Lean, \lean{Point} abbreviates the type of functions from
$\mathrm{Fin}\,2=\{0,1\}$ to $\R$; we write $p=(p_0,p_1)$, where $p_0=p(0)$ is the first
($x$-)coordinate and $p_1=p(1)$ the second ($y$-)coordinate. Addition, subtraction and multiplication by
real numbers are coordinatewise.

A \emph{finite set of points} is a Lean \texttt{Finset Point}: a finite set, without multiplicities. We use
ordinary set notation for finite sets: $S\subseteq T$, $S\cup T$, $S\setminus T$, $S\setminus\{c\}$
(Lean \texttt{S.erase c}), $\card S$ (Lean \lean{S.card}), and $f(S)=\{f(p):p\in S\}$ for a map $f$ (Lean
\texttt{S.image f}, again a finite set, with $\card{f(S)}=\card S$ when $f$ is injective). Where Lean needs
a set rather than a finite set, for example as the argument of a convex hull, it uses the coercion of $S$
to a \texttt{Set Point}, which has the same elements; we do not distinguish the two. Brace notation
$\{a,b,c\}$ denotes the set with these elements; it has fewer than three elements when two of $a,b,c$
coincide.
\end{definition}

\begin{remark}[Norm and topology]\label{fd:rem:norm}
Mathlib gives the Lean type of points its default sup norm $\max(\abs{p_0},\abs{p_1})$. No statement of
this section uses a norm or a metric; Euclidean quantities are always written out in coordinates (as in
Lemma~\ref{fd:lem:separation}). The only topological argument of this section is the compactness
argument in the proof of Lemma~\ref{fd:lem:separation}, which takes place in $\R^m$ with its usual
topology. (On $\R^2$ the sup norm and the Euclidean norm define the same, usual, topology.)
\end{remark}

\begin{definition}[Orientation; \leanref{ErdosSzekeres.orient}{ErdosSzekeres/Orientation.lean:9}]\label{fd:def:orient}\sloppy
Let $R$ be a commutative ring and let $a,b,c\in R^2$ (in Lean: functions from $\mathrm{Fin}\,2$ to
$R$). Then
\[
  \orient(a,b,c)=(b_0-a_0)(c_1-a_1)-(b_1-a_1)(c_0-a_0)\in R .
\]
Unless stated otherwise $R=\R$; Lean also evaluates $\orient$ over $\mathbb{Q}$ and $\Z$ in exact
certificate checks. Equivalently, $\orient(a,b,c)=\det\bigl(b-a\;\big|\;c-a\bigr)$ is the determinant of
the $2\times2$ matrix with columns $b-a$ and $c-a$, that is, twice the signed area of the triangle $abc$.
Positive orientation means counterclockwise; for example $\orient((0,0),(1,0),(0,1))=1>0$. For $a\ne b$
we say that $c$ lies (strictly) to the \emph{left} of the line through $a$ and $b$ directed from $a$ to
$b$ if $\orient(a,b,c)>0$, and to the \emph{right} if $\orient(a,b,c)<0$.
\end{definition}

\begin{definition}[Convex sets, convex hulls, affine maps; Mathlib \lean{Convex}, \lean{convexHull}, \lean{AffineIndependent}, \lean{AffineMap}, \lean{AffineEquiv}]\label{fd:def:conv}
A set $K\subseteq\R^2$ is \emph{convex} if $\lambda x+\mu y\in K$ for all $x,y\in K$ and all reals
$\lambda,\mu\ge0$ with $\lambda+\mu=1$ (Mathlib \lean{Convex}). The \emph{convex hull} $\conv Y$ of a set
$Y\subseteq\R^2$ is the intersection of all convex sets that contain $Y$ (Mathlib \lean{convexHull},
defined as the smallest convex superset). A \emph{convex combination} of points $y_1,\dots,y_k$
($k\ge1$) is a point $\sum_{i=1}^k\lambda_iy_i$ with all $\lambda_i\ge0$ and $\sum_i\lambda_i=1$. A finite
set $T\subseteq\R^2$ is \emph{affinely independent} if the only family of reals $(\mu_y)_{y\in T}$ with
$\sum_{y\in T}\mu_y=0$ and $\sum_{y\in T}\mu_y\,y=0$ is the zero family (equivalent to Mathlib's
\lean{AffineIndependent} for the inclusion map of $T$). An \emph{affine map} of $\R^2$ is a map
$g(x)=Ax+v$, where $A$ is a real $2\times2$ matrix (the \emph{linear part} of $g$) and $v\in\R^2$.
In Mathlib, an affine map $f$ from the plane to itself (\lean{AffineMap}) is a map $f\colon\R^2\to\R^2$
together with a linear map $f_{\mathrm{lin}}\colon\R^2\to\R^2$ (Lean \lean{f.linear}) such that
$f(w+p)=f_{\mathrm{lin}}(w)+f(p)$ for all $p,w\in\R^2$. These are exactly the affine maps just defined:
taking $p=0$ gives $f(x)=A_fx+f(0)$, where $A_f$ is the matrix of $f_{\mathrm{lin}}$; conversely,
$g(x)=Ax+v$ satisfies $g(w+p)=Aw+g(p)$, so it is a Mathlib affine map whose linear map is $x\mapsto Ax$.
A Mathlib \emph{affine equivalence} of the plane (\lean{AffineEquiv}) is a bijection $f$ of $\R^2$
together with a linear bijection $f_{\mathrm{lin}}$ of $\R^2$ with the same property; in particular it
is an injective affine map.
\end{definition}

\begin{definition}[General position; \leanref{ErdosSzekeres.GeneralPosition}{ErdosSzekeres/Definitions.lean:12}]\label{fd:def:gp}\sloppy
A set $X\subseteq\R^2$ is \emph{collinear} (Mathlib \lean{Collinear}) if the linear span of the set of
differences $\{p-q:p,q\in X\}$ (Mathlib \lean{vectorSpan}) has dimension at most~$1$ (in Lean:
\lean{Module.rank} of this span is at most $1$). A finite set $S$ of points is \emph{in general
position} if for all $a,b,c\in S$ with $a\ne b$, $a\ne c$ and $b\ne c$ the set $\{a,b,c\}$ is not
collinear.
\end{definition}

\begin{definition}[Convex position; \leanref{ErdosSzekeres.ConvexPosition}{ErdosSzekeres/Definitions.lean:17}]\label{fd:def:cp}\sloppy
A finite set $S$ of points is \emph{in convex position} if the inclusion map of $S$ into $\R^2$ is
convex independent over $\R$ (Mathlib \lean{ConvexIndependent}). Mathlib calls a family
$p\colon\iota\to\R^2$ convex independent if for every set $t\subseteq\iota$ and every $x\in\iota$,
$p(x)\in\conv p(t)$ implies $x\in t$. For the inclusion map of $S$ the index set is $S$ itself (in Lean,
the type of elements of $S$) and $p$ is the identity on it, so the definition says: \emph{for every
$T\subseteq S$ and every $x\in S$, if $x\in\conv T$ then $x\in T$.} Lemma~\ref{fd:lem:cp-hull} gives the
familiar form: no point of $S$ lies in the convex hull of the other points of $S$.
\end{definition}

\begin{definition}[Convex $n$-gons and $K$-free sets; \leanref{ErdosSzekeres.HasConvexNGon}{ErdosSzekeres/Definitions.lean:21}]\label{fd:def:ngon}\sloppy
For $n\in\N$ and a finite set $S$ of points, $\mathrm{HasConvexNGon}\,n\,S$ holds if there is a finite
set $T\subseteq S$ with $\card T=n$ that is in convex position. There is no emptiness condition: points
of $S\setminus T$ may lie in $\conv T$. For $K\in\N$ we say that $S$ is \emph{$K$-free} if
$\neg\,\mathrm{HasConvexNGon}\,K\,S$.
\end{definition}

\begin{remark}\label{fd:rem:ngon}
$\mathrm{HasConvexNGon}\,0\,S$ holds for every $S$ (Lemma~\ref{fd:lem:small}(a)), so no set is
$0$-free, and a hypothesis ``$S$ is $0$-free'' is contradictory. If $S$ is $K$-free and $K\le K'$, then
$S$ is $K'$-free (contrapositive of Lemma~\ref{fd:lem:mono}(d)). If $\mathrm{HasConvexNGon}\,n\,S$, then
$n\le\card S$ (Lemma~\ref{fd:lem:ngon-basic}(a)).
\end{remark}

\begin{definition}[Forcing a convex $n$-gon; \leanref{ErdosSzekeres.ForcesConvexNGon}{ErdosSzekeres/Definitions.lean:25}]\label{fd:def:forces}\sloppy
For $N,n\in\N$, $\mathrm{ForcesConvexNGon}\,N\,n$ holds if every finite set $S$ of points that is in
general position and satisfies $N\le\card S$ satisfies $\mathrm{HasConvexNGon}\,n\,S$.
\end{definition}

\begin{definition}[Invariant sets; \leanref{ErdosSzekeres.LayerLemma.Invariant}{LAYERLEAN/LayerLemma.lean:105}, \leanref{ErdosSzekeres.SymReduction.Invariant}{SYMLEAN/SymmetricReduction.lean:271}, \leanref{ErdosSzekeres.OrbitConsistency.Invariant}{LEANEXT/orbit/OrbitConsistency.lean:47}]\label{fd:def:inv}\sloppy
For a map $f\colon\R^2\to\R^2$ and a finite set $S$ of points, $\mathrm{Invariant}\,f\,S$ holds if
$f(p)\in S$ for every $p\in S$; that is, $f$ maps $S$ \emph{into} $S$. For injective $f$ this is
equivalent to $f(S)=S$, and then also $f^{-1}(S)=\{p\in\R^2:f(p)\in S\}=S$
(Lemma~\ref{fd:lem:inv-image}). The three cited Lean definitions are textually identical.
\end{definition}

\begin{lemma}[Invariance and images; (a), (b) written step; (c) \leanref{ErdosSzekeres.LayerLemma.iterate_mem}{LAYERLEAN/LayerLemma.lean:52}]\label{fd:lem:inv-image}\sloppy
Let $X$ be a set, let $f\colon X\to X$ be a map and let $S$ be a finite subset of $X$.
\begin{enumerate}[label=(\alph*)]
\item If $X=\R^2$ and $f$ is injective, then $\mathrm{Invariant}\,f\,S$ holds if and only if $f(S)=S$.
\item If $X=\R^2$, $f$ is injective and $\mathrm{Invariant}\,f\,S$ holds, then $\{p\in\R^2:f(p)\in S\}=S$.
\item If $f(y)\in S$ for every $y\in S$, then $f^{[i]}(x)\in S$ for every $x\in S$ and every $i\in\N$.
\end{enumerate}
Part (c) is stated for an arbitrary set $X$, as in Lean, where $X$ is an arbitrary type; parts (a) and
(b) take $X=\R^2$ because $\mathrm{Invariant}$ is defined for maps of the plane.
\end{lemma}

\begin{proof}
(a) If $f(S)=S$, then $f(p)\in f(S)=S$ for every $p\in S$. Conversely, if $\mathrm{Invariant}\,f\,S$
holds, then $f(S)\subseteq S$; since $f$ is injective, $\card{f(S)}=\card S$, and a subset of a finite
set with the same number of elements is the whole set, so $f(S)=S$. (Lean uses this argument in the step
at line~165 of the proof of \lean{ErdosSzekeres.LayerLemma.map_mem_hullVerts},
\dirn{LAYERLEAN/LayerLemma.lean:162}.)

(b) If $p\in S$ then $f(p)\in S$ by invariance. If $f(p)\in S$, then by (a) $f(p)\in f(S)$, so
$f(p)=f(q)$ for some $q\in S$, and $p=q\in S$ by injectivity.

(c) Induction on $i$. For $i=0$, $f^{[0]}(x)=x\in S$. If $f^{[i]}(x)\in S$, then
$f^{[i+1]}(x)=f\bigl(f^{[i]}(x)\bigr)\in S$ by the hypothesis applied to $y=f^{[i]}(x)$. Here we used
$f^{[i+1]}=f\circ f^{[i]}$, which follows from the definition $f^{[i+1]}=f^{[i]}\circ f$ by induction on
$i$ ($f^{[1]}=f^{[0]}\circ f=f=f\circ f^{[0]}$, and $f^{[i+2]}=f^{[i+1]}\circ f=f\circ f^{[i]}\circ
f=f\circ f^{[i+1]}$; Mathlib's \lean{Function.iterate_succ_apply'}). This is the Lean proof.
\end{proof}

\begin{definition}[Rotation about a point; \leanref{ErdosSzekeres.LayerLemma.rotAbout}{LAYERLEAN/LayerLemma.lean:298}, \leanref{ErdosSzekeres.SymReduction.rotAbout}{SYMLEAN/SymmetricReduction.lean:400}, \leanref{ErdosSzekeres.OrbitConsistency.rotAbout}{LEANEXT/orbit/OrbitConsistency.lean:93}]\label{fd:def:rot}\sloppy
For $c\in\R^2$ and $\theta\in\R$, the map $\rho_{c,\theta}=\mathrm{rotAbout}\,c\,\theta\colon\R^2\to\R^2$
is
\[
  \rho_{c,\theta}(p)=\bigl(c_0+\cos\theta\,(p_0-c_0)-\sin\theta\,(p_1-c_1),\;
  c_1+\sin\theta\,(p_0-c_0)+\cos\theta\,(p_1-c_1)\bigr).
\]
In complex notation, identifying $p$ with $p_0+ip_1$ and $c$ with $c_0+ic_1$, this is
$z\mapsto c+e^{i\theta}(z-c)$ (Lemma~\ref{fd:lem:rot}(a)): the counterclockwise rotation about $c$ by
the angle $\theta$ (Lemma~\ref{fd:lem:rot}(e)). The three cited Lean definitions are textually
identical; a text search of the research records finds 19 definitions of \lean{rotAbout}, all with this
same body.
\end{definition}

\begin{remark}[Angles $2\pi/m$]\label{fd:rem:angles}
In later sections the angle is usually written \texttt{2 * Real.pi / m} in Lean with \lean{m} a natural
number, that is, $\theta=2\pi/(m:\R)$. For $m=0$ Lean's division gives $\theta=0$, and $\rho_{c,0}$ is
the identity (Lemma~\ref{fd:lem:rot}(c)). The propositions \lean{Sixfold66} and \lean{Sevenfold70}
(Definition~\ref{ll:def:66-70}) write their angles as \texttt{2 * Real.pi / 6} and
\texttt{2 * Real.pi / 7} with real numerals; since the real numbers $6$ and $7$ are the casts of the
natural numbers $6$ and $7$, these are the angles $2\pi/(m:\R)$ for $m=6$ and $m=7$.
\end{remark}

\begin{lemma}[Basic properties of rotations; (a) \leanref{ErdosSzekeres.LayerLemma.rotAbout_eq_crot}{LAYERLEAN/LayerLemma.lean:302}; (b) \leanref{ErdosSzekeres.LayerLemma.rotAff}{LAYERLEAN/LayerLemma.lean:311}; (c) \leanref{ErdosSzekeres.LayerLemma.rotAbout_iterate}{LAYERLEAN/LayerLemma.lean:465}; (d) \leanref{ErdosSzekeres.LayerLemma.rotAbout_injective}{LAYERLEAN/LayerLemma.lean:319}; (e), (f) written step]\label{fd:lem:rot}\sloppy
Let $c\in\R^2$ and $\theta,\alpha,\beta\in\R$. Write $\rho_\theta=\rho_{c,\theta}$ and
$R_\theta=\begin{pmatrix}\cos\theta&-\sin\theta\\ \sin\theta&\cos\theta\end{pmatrix}$.
\begin{enumerate}[label=(\alph*)]
\item For every $p\in\R^2$, the complex number $\rho_\theta(p)_0+i\,\rho_\theta(p)_1$ equals
  $\gamma+e^{i\theta}(z-\gamma)$, where $z=p_0+ip_1$ and $\gamma=c_0+ic_1$. (In Lean:
  $\mathrm{rotAbout}\,c\,\theta=\mathrm{crot}\,(\mathrm{toC}\,c)\,(\exp(\theta i))$, where
  $\mathrm{toC}\,p=p_0+ip_1$ and $\mathrm{crot}\,\gamma\,w\,p$ is the point with complex coordinate
  $\gamma+w(\mathrm{toC}\,p-\gamma)$.)
\item $\rho_\theta(p)=c+R_\theta(p-c)=R_\theta\,p+(c-R_\theta\,c)$ for every $p$; thus $\rho_\theta$ is an
  affine map with linear part $R_\theta$, that is, $\rho_\theta(p+v)=\rho_\theta(p)+R_\theta v$ for all
  $p,v\in\R^2$. In particular $\rho_\theta(c)=c$. (The Lean definition \lean{LayerLemma.rotAff} cited for (b)
  packages $\rho_\theta$ as an affine map with linear part $R_\theta$, its defining property being the
  identity $\rho_\theta(p+v)=\rho_\theta(p)+R_\theta v$. The matrix form
  $\rho_\theta(p)=c+R_\theta(p-c)=R_\theta\,p+(c-R_\theta\,c)$ of Definition~\ref{fd:def:rot} and the
  identity $\rho_\theta(c)=c$ are written steps here; Lean proves the latter in the step at line~330 of the
  proof of \lean{ErdosSzekeres.LayerLemma.invariant_erase_centre}, \dirn{LAYERLEAN/LayerLemma.lean:325}.)
\item $\rho_\alpha\circ\rho_\beta=\rho_{\alpha+\beta}$ and $\rho_0=\mathrm{id}$. Hence
  $\rho_\alpha^{[j]}=\rho_{(j:\R)\alpha}$ for every $j\in\N$. (The Lean declaration cited for (c) is this
  iterate formula; the composition law and $\rho_0=\mathrm{id}$ are written steps here.)
\item $\rho_\theta$ is a bijection of $\R^2$ with inverse $\rho_{-\theta}$; in particular it is
  injective. (The Lean declaration cited for (d) states injectivity; bijectivity and the inverse are a
  written step here.)
\item $\orient\bigl(c,p,\rho_\theta(p)\bigr)=\sin\theta\,\bigl((p_0-c_0)^2+(p_1-c_1)^2\bigr)$ for every
  $p$. Hence, if $0<\theta<\pi$ and $p\ne c$, then $\orient(c,p,\rho_\theta(p))>0$: the rotation turns
  counterclockwise.
\item $(\rho_\theta(p)_0-c_0)^2+(\rho_\theta(p)_1-c_1)^2=(p_0-c_0)^2+(p_1-c_1)^2$ for every $p$.
\end{enumerate}
\end{lemma}

\begin{proof}
We give direct proofs in coordinates; the Lean proofs of (c) and (d) pass through the complex form (a).
Put $u=p-c$. By definition,
\[
  \rho_\theta(p)-c=\bigl(\cos\theta\,u_0-\sin\theta\,u_1,\;\sin\theta\,u_0+\cos\theta\,u_1\bigr)
  =R_\theta\,u .
\]
(b) This is the displayed identity; $R_\theta(p-c)=R_\theta p-R_\theta c$ by linearity. Hence
$\rho_\theta(p+v)=c+R_\theta(p-c)+R_\theta v=\rho_\theta(p)+R_\theta v$, and $p=c$ gives $u=0$ and
$\rho_\theta(c)=c$.

(a) By Euler's formula $e^{i\theta}=\cos\theta+i\sin\theta$, and
$(\cos\theta+i\sin\theta)(u_0+iu_1)=(\cos\theta\,u_0-\sin\theta\,u_1)+i(\sin\theta\,u_0+\cos\theta\,u_1)$.
Adding $\gamma=c_0+ic_1$ gives the two coordinates of $\rho_\theta(p)$.

(c) By (b), $\rho_\alpha(\rho_\beta(p))=c+R_\alpha R_\beta(p-c)$. By the addition formulas
$\cos(\alpha+\beta)=\cos\alpha\cos\beta-\sin\alpha\sin\beta$ and
$\sin(\alpha+\beta)=\sin\alpha\cos\beta+\cos\alpha\sin\beta$,
\[
  R_\alpha R_\beta=
  \begin{pmatrix}\cos\alpha\cos\beta-\sin\alpha\sin\beta&-(\sin\alpha\cos\beta+\cos\alpha\sin\beta)\\
  \sin\alpha\cos\beta+\cos\alpha\sin\beta&\cos\alpha\cos\beta-\sin\alpha\sin\beta\end{pmatrix}
  =R_{\alpha+\beta},
\]
so $\rho_\alpha\circ\rho_\beta=\rho_{\alpha+\beta}$. Since $R_0$ is the identity matrix, $\rho_0=\mathrm{id}$.
Now induct on $j$: $\rho_\alpha^{[0]}=\mathrm{id}=\rho_0=\rho_{(0:\R)\alpha}$, and
$\rho_\alpha^{[j+1]}=\rho_\alpha^{[j]}\circ\rho_\alpha=\rho_{(j:\R)\alpha}\circ\rho_\alpha
=\rho_{(j:\R)\alpha+\alpha}=\rho_{(j+1:\R)\alpha}$.

(d) By (c), $\rho_{-\theta}\circ\rho_\theta=\rho_0=\mathrm{id}$ and $\rho_\theta\circ\rho_{-\theta}=\mathrm{id}$.

(e) With $u=p-c$ as above,
\[
  \orient(c,p,\rho_\theta(p))=u_0(\sin\theta\,u_0+\cos\theta\,u_1)-u_1(\cos\theta\,u_0-\sin\theta\,u_1)
  =\sin\theta\,(u_0^2+u_1^2).
\]
If $0<\theta<\pi$ then $\sin\theta>0$, and if $p\ne c$ then $u_0^2+u_1^2>0$.

(f) $(\cos\theta\,u_0-\sin\theta\,u_1)^2+(\sin\theta\,u_0+\cos\theta\,u_1)^2
=(\cos^2\theta+\sin^2\theta)(u_0^2+u_1^2)=u_0^2+u_1^2$, the cross terms cancelling.
\end{proof}

%======================================================================
\subsection{Facts}\label{fd:sec-facts}

\begin{lemma}[Symmetries of the orientation; (a) \leanref{ErdosSzekeres.orient_swap}{ErdosSzekeres/Orientation.lean:12}; (b) \leanref{ErdosSzekeres.orient_cycle}{ErdosSzekeres/Orientation.lean:15}; (c) \leanref{ErdosSzekeres.orient_translate}{ErdosSzekeres/Orientation.lean:18}; (d), (e) written step]\label{fd:lem:orient}\sloppy
Let $R$ be a commutative ring and let $a,b,c,v\in R^2$.
\begin{enumerate}[label=(\alph*)]
\item $\orient(a,c,b)=-\orient(a,b,c)$.
\item $\orient(b,c,a)=\orient(a,b,c)$. Consequently $\orient(a,b,c)=\orient(b,c,a)=\orient(c,a,b)$.
\item $\orient(a+v,b+v,c+v)=\orient(a,b,c)$.
\item For all $p_1,p_2,p_3\in R^2$ and every permutation $\pi$ of $\{1,2,3\}$,
  $\orient(p_{\pi(1)},p_{\pi(2)},p_{\pi(3)})=\sgn(\pi)\,\orient(p_1,p_2,p_3)$.
\item $\orient(a,b,c)=0$ whenever two of $a,b,c$ are equal. Consequently, if $\orient(a,b,c)\ne0$, then
  $a,b,c$ are pairwise distinct.
\end{enumerate}
\end{lemma}

\begin{proof}
The Lean proofs of (a)--(c) expand both sides and compare polynomials (\texttt{unfold orient; ring} for
(a) and (b), \texttt{simp only [orient, Pi.add\_apply]; ring} for (c)); we organise the same
computation. For $x,y\in R^2$ put $[x,y]=x_0y_1-x_1y_0$. Then $[\cdot,\cdot]$ is
bilinear, $[y,x]=-[x,y]$ and $[x,x]=0$. Expanding the definition (the terms $a_0a_1$ cancel),
\[
  \orient(a,b,c)=a_0b_1-a_1b_0+b_0c_1-b_1c_0+c_0a_1-c_1a_0=[a,b]+[b,c]+[c,a].
\]
(b) $\orient(b,c,a)=[b,c]+[c,a]+[a,b]=\orient(a,b,c)$. Applying this twice gives the second statement.

(a) $\orient(a,c,b)=[a,c]+[c,b]+[b,a]=-[c,a]-[b,c]-[a,b]=-\orient(a,b,c)$.

(c) The defining formula involves only the differences $b-a$ and $c-a$, and
$(b+v)-(a+v)=b-a$, $(c+v)-(a+v)=c-a$.

(d) The permutation group of $\{1,2,3\}$ has six elements. The identity and the two $3$-cycles have sign
$+1$; they give $\orient(p_1,p_2,p_3)$, $\orient(p_2,p_3,p_1)$ and $\orient(p_3,p_1,p_2)$, which are equal
by (b). The three transpositions have sign $-1$; they give $\orient(p_1,p_3,p_2)$, $\orient(p_2,p_1,p_3)$
and $\orient(p_3,p_2,p_1)$. The first equals $-\orient(p_1,p_2,p_3)$ by (a). By (b) applied to the triple
$(p_2,p_1,p_3)$ we get $\orient(p_1,p_3,p_2)=\orient(p_2,p_1,p_3)$, and by (b) applied to
$(p_3,p_2,p_1)$ we get $\orient(p_2,p_1,p_3)=\orient(p_3,p_2,p_1)$; so all three equal
$-\orient(p_1,p_2,p_3)$.

(e) If $a=b$ the first factors $b_0-a_0$ and $b_1-a_1$ vanish; if $a=c$ the factors $c_1-a_1$ and
$c_0-a_0$ vanish; if $b=c$ the two products $(b_0-a_0)(b_1-a_1)$ and $(b_1-a_1)(b_0-a_0)$ are equal. In
each case $\orient(a,b,c)=0$.
\end{proof}

\begin{remark}[The real case of (d) in Lean]\label{fd:rem:orient-lean}\sloppy
For $R=\R$, part (d) of Lemma~\ref{fd:lem:orient} is the Lean theorem
\lean{ErdosSzekeres.OrbitConsistency.orient3_perm} (\dirn{LEANEXT/orbit/OrbitConsistency.lean:178};
Lemma~\ref{oc:lem:perm}). It is stated for ordered triples $t\colon\mathrm{Fin}\,3\to\R^2$, with
$\orient_3(t)=\orient(t(0),t(1),t(2))$ (Lean \lean{orient3},
\dirn{LEANEXT/orbit/OrbitConsistency.lean:168}), and says that
$\orient_3(t\circ\sigma)=\sgn(\sigma)\,\orient_3(t)$ for every permutation $\sigma$ of
$\mathrm{Fin}\,3=\{0,1,2\}$ and every $t$. Put $t(i)=p_{i+1}$ and $\sigma(i)=\pi(i+1)-1$; then
$\sigma=s^{-1}\circ\pi\circ s$ for the bijection $s(i)=i+1$ from $\{0,1,2\}$ to $\{1,2,3\}$, so $\sigma$
has the same cycle type and hence the same sign as $\pi$, and $(t\circ\sigma)(i)=p_{\pi(i+1)}$. So
$\orient_3(t\circ\sigma)=\orient(p_{\pi(1)},p_{\pi(2)},p_{\pi(3)})$ and $\orient_3(t)=\orient(p_1,p_2,p_3)$,
and the Lean statement is (d) for $R=\R$. Parts (d) for a general commutative ring and (e) are written
steps.
\end{remark}

\begin{lemma}[Orientation under affine maps; \leanref{ErdosSzekeres.SymReduction.orient_affine}{SYMLEAN/AffineReduction.lean:1241}; (i) \leanref{ErdosSzekeres.OrbitConsistency.orient_rotAbout_inv}{LEANEXT/orbit/OrbitConsistency.lean:117}; (ii) written step]\label{fd:lem:affine-orient}\sloppy
Let $g\colon\R^2\to\R^2$ be a map, let $A$ be a real $2\times2$ matrix and $v\in\R^2$, and suppose that
$g(x)=Ax+v$ for every $x\in\R^2$. Then for all $a,b,c\in\R^2$
\[
  \orient\bigl(g(a),g(b),g(c)\bigr)=\det(A)\,\orient(a,b,c).
\]
(In Lean, $g$ is an arbitrary map with the hypothesis that $g(p)=L\,p+w$ for every $p$, where $L$ is a
real $2\times2$ matrix and $w\in\R^2$ (Lean's \lean{b}), and the conclusion is
$\orient(g(x),g(y),g(z))=\det(L)\,\orient(x,y,z)$ for all $x,y,z$; this is the statement above with
$L=A$ and $w=v$.)
In particular:
\begin{enumerate}[label=(\roman*)]
\item (rotations have determinant $1$) for all $o\in\R^2$, $\theta\in\R$ and $a,b,d\in\R^2$,
  $\orient(\rho_{o,\theta}(a),\rho_{o,\theta}(b),\rho_{o,\theta}(d))=\orient(a,b,d)$, because the linear part
  $R_\theta$ of $\rho_{o,\theta}$ (Lemma~\ref{fd:lem:rot}(b)) has $\det R_\theta=\cos^2\theta+\sin^2\theta=1$;
\item (reflections have determinant $-1$) for $o\in\R^2$ and $\varphi\in\R$ let
  $F_\varphi=\begin{pmatrix}\cos2\varphi&\sin2\varphi\\ \sin2\varphi&-\cos2\varphi\end{pmatrix}$ and
  $\sigma(x)=o+F_\varphi(x-o)$, the reflection in the line through $o$ with direction
  $u=(\cos\varphi,\sin\varphi)$ (indeed $F_\varphi u=u$ and $F_\varphi u^\perp=-u^\perp$ for
  $u^\perp=(-\sin\varphi,\cos\varphi)$). Then
  $\orient(\sigma(a),\sigma(b),\sigma(d))=-\orient(a,b,d)$ for all $a,b,d$, because
  $\det F_\varphi=-\cos^22\varphi-\sin^22\varphi=-1$. For $o=0$ and $\varphi=0$, $\sigma$ is the reflection
  $(p_0,p_1)\mapsto(p_0,-p_1)$ in the first coordinate axis.
\end{enumerate}
\end{lemma}

\begin{proof}
By Definition~\ref{fd:def:orient}, $\orient(a,b,c)=\det(b-a\mid c-a)$. Since $g(b)-g(a)=A(b-a)$ and
$g(c)-g(a)=A(c-a)$, the matrix with columns $g(b)-g(a)$ and $g(c)-g(a)$ is the product $A\,(b-a\mid c-a)$,
and the determinant of a product of $2\times2$ matrices is the product of the determinants (a polynomial
identity in the eight entries, checked by expanding both sides). This proves the formula. (The Lean proof
substitutes $g(p)=Lp+w$, writes $\det L=L_{00}L_{11}-L_{01}L_{10}$, expands both sides in coordinates
and compares polynomials, which is the same computation.)

(i) $\rho_{o,\theta}(x)=R_\theta x+(o-R_\theta o)$ by Lemma~\ref{fd:lem:rot}(b), and
$\det R_\theta=\cos\theta\cdot\cos\theta-(-\sin\theta)\sin\theta=1$.

(ii) $\sigma(x)=F_\varphi x+(o-F_\varphi o)$ has linear part $F_\varphi$, and
$\det F_\varphi=-\cos^22\varphi-\sin^22\varphi=-1$. Using $\cos2\varphi\cos\varphi+\sin2\varphi\sin\varphi
=\cos\varphi$ and $\sin2\varphi\cos\varphi-\cos2\varphi\sin\varphi=\sin\varphi$ (the subtraction formulas for
$2\varphi-\varphi$) one gets $F_\varphi u=u$ and $F_\varphi u^\perp=(\sin\varphi,-\cos\varphi)=-u^\perp$, so
$\sigma$ fixes the line $o+\R u$ pointwise and reverses the perpendicular direction. For $\varphi=0$,
$F_0=\mathrm{diag}(1,-1)$.
\end{proof}

\begin{remark}[The formula and its special cases in Lean]\label{fd:rem:affine-orient-lean}\sloppy
The general formula is the Lean theorem \lean{ErdosSzekeres.SymReduction.orient_affine}, cited in the
head; a verbatim copy, with the same name, is at \dirn{SYMLEAN/mirror/ReductionThreeClasses.lean:1267}.
Part (i) is a separate Lean theorem, \lean{ErdosSzekeres.OrbitConsistency.orient_rotAbout_inv}, also
cited in the head (Lemma~\ref{oc:lem:rot-inv}). It is stated for the copy
\lean{ErdosSzekeres.OrbitConsistency.rotAbout} of the rotation of Definition~\ref{fd:def:rot}, and it
calls the centre $c$ instead of $o$. Part (ii) is a written step. Its case $o=0$, $\varphi=0$, the sign
change under $(p_0,p_1)\mapsto(p_0,-p_1)$, is the Lean theorem \lean{ErdosSzekeres.Suk.orient_reflectY}
(\dirn{ErdosSzekeres/Suk/Reflection.lean:47}).
\end{remark}

\begin{lemma}[Collinearity through a point; Mathlib \lean{collinear_iff_of_mem}; written step]\label{fd:lem:collinear-mem}
Let $X\subseteq\R^2$ and $p^*\in X$. Then $X$ is collinear if and only if there is $v\in\R^2$ such that
every $p\in X$ can be written $p=p^*+r\,v$ with $r\in\R$.
\end{lemma}

\begin{proof}
Let $W$ be the linear span of $\{p-q:p,q\in X\}$. Since $p^*\in X$, $W$ is also the span of
$\{p-p^*:p\in X\}$: these vectors are among the differences, and conversely
$p-q=(p-p^*)-(q-p^*)$. If $X$ is collinear, then $\dim W\le1$, so $W\subseteq\R v$ for some $v$ (take
$v=0$ if $W=0$, and a basis vector of $W$ if $\dim W=1$); then $p-p^*\in W$ gives $p=p^*+rv$ for some
$r$. Conversely, if every $p-p^*$ lies in $\R v$, then $W\subseteq\R v$ and $\dim W\le\dim\R v\le1$.
\end{proof}

\begin{lemma}[Orientation and collinearity; (a) \leanref{ErdosSzekeres.orient_eq_zero_iff_collinear}{ErdosSzekeres/OrientationGeometry.lean:36}; (b) \leanref{ErdosSzekeres.generalPosition_iff_orient_ne_zero}{ErdosSzekeres/OrientationGeometry.lean:66}]\label{fd:lem:collinear}\sloppy\tolerance=10000
The orientation detects collinearity, and hence general position.
\begin{enumerate}[label=(\alph*)]
\item For all $a,b,c\in\R^2$: $\orient(a,b,c)=0$ if and only if the set $\{a,b,c\}$ is collinear. This
  includes degenerate triples, in which two or three of the points coincide.
\item A finite set $S$ of points is in general position if and only if $\orient(a,b,c)\ne0$ for all
  $a,b,c\in S$ with $a\ne b$, $a\ne c$ and $b\ne c$.
\end{enumerate}
\end{lemma}

\begin{proof}
We follow the Lean proof. (a) We use Lemma~\ref{fd:lem:collinear-mem} with $X=\{a,b,c\}$ and $p^*=a$.

Suppose $\{a,b,c\}$ is collinear. Then there are $v$ and reals $r,s$ with $b=a+rv$ and $c=a+sv$, so
$\orient(a,b,c)=(rv_0)(sv_1)-(rv_1)(sv_0)=0$.

Suppose $\orient(a,b,c)=0$. If $b=a$, put $v=c-a$; then $a=a+0\cdot v$, $b=a+0\cdot v$ and $c=a+1\cdot v$, so
$\{a,b,c\}$ is collinear. If $b\ne a$, put $u=b-a\ne0$ and $w=c-a$; then
$u_0w_1-u_1w_0=\orient(a,b,c)=0$. We find $r$ with $w=ru$. If $u_0\ne0$, let $r=w_0/u_0$; then $ru_0=w_0$
and $ru_1=w_0u_1/u_0=u_0w_1/u_0=w_1$. If $u_0=0$, then $u_1\ne0$ and $0=u_0w_1-u_1w_0=-u_1w_0$ gives
$w_0=0$; let $r=w_1/u_1$, so that $ru_0=0=w_0$ and $ru_1=w_1$. Now $a=a+0\cdot u$, $b=a+1\cdot u$ and
$c=a+r\,u$, so $\{a,b,c\}$ is collinear.

(b) By Definition~\ref{fd:def:gp}, $S$ is in general position iff $\{a,b,c\}$ is not collinear for all
pairwise distinct $a,b,c\in S$; by (a) this says $\orient(a,b,c)\ne0$ for all such triples.
\end{proof}

\begin{lemma}[Convex hulls; written step]\label{fd:lem:hull}
Let $Y,Y'\subseteq\R^2$.
\begin{enumerate}[label=(\alph*)]
\item $Y\subseteq\conv Y$; $\conv Y$ is convex; and $\conv Y\subseteq K$ for every convex set $K$ with
  $Y\subseteq K$. If $Y\subseteq Y'$ then $\conv Y\subseteq\conv Y'$. Moreover $\conv\emptyset=\emptyset$
  and $\conv\{y\}=\{y\}$.
\item $\conv Y$ is the set of all convex combinations of finitely many points of $Y$.
\item If $F=\{y_1,\dots,y_m\}$ is a finite set of $m\ge1$ distinct points, then
  $\conv F=\bigl\{\sum_{i=1}^m\lambda_iy_i:\lambda_i\ge0,\ \sum_i\lambda_i=1\bigr\}$. In particular
  $\conv\{y,z\}=\{\lambda y+(1-\lambda)z:0\le\lambda\le1\}$.
\item For $w\in\R^2$ and $\alpha\in\R$, the open half-plane $\{z\in\R^2:w_0z_0+w_1z_1<\alpha\}$ is
  convex.
\end{enumerate}
\end{lemma}

\begin{proof}
(a) An intersection of convex sets is convex, so $\conv Y$ is convex; it contains $Y$; and it is contained
in each convex $K\supseteq Y$, being the intersection of all of them. If $Y\subseteq Y'$, every convex
set containing $Y'$ contains $Y$, so $\conv Y\subseteq\conv Y'$. The empty set and $\{y\}$ are convex
(for $\{y\}$: $\lambda y+\mu y=y$ when $\lambda+\mu=1$), which gives the last two claims.

(b) Let $C$ be the set of convex combinations of finitely many points of $Y$. Every $y\in Y$ is the convex
combination $1\cdot y$, so $Y\subseteq C$. $C$ is convex: if $x=\sum_i\lambda_iy_i$ and
$x'=\sum_j\lambda'_jy'_j$ are in $C$ and $\lambda,\mu\ge0$ with $\lambda+\mu=1$, then
$\lambda x+\mu x'=\sum_i(\lambda\lambda_i)y_i+\sum_j(\mu\lambda'_j)y'_j$ has nonnegative coefficients
with sum $\lambda+\mu=1$. Hence $\conv Y\subseteq C$ by (a). Conversely, let $K\supseteq Y$ be convex; we
show by induction on $k$ that every convex combination $x=\sum_{i=1}^k\lambda_iy_i$ of points of $Y$
lies in $K$. For $k=1$, $x=y_1\in K$. For $k+1$ terms: if $\lambda_{k+1}=1$, all other coefficients
vanish and $x=y_{k+1}\in K$; otherwise $s=\sum_{i\le k}\lambda_i=1-\lambda_{k+1}>0$, the point
$x'=\sum_{i\le k}(\lambda_i/s)y_i$ is a convex combination of $k$ points of $Y$, so $x'\in K$, and
$x=s\,x'+\lambda_{k+1}y_{k+1}$ with $s+\lambda_{k+1}=1$ lies in $K$ by convexity. Taking $K=\conv Y$
gives $C\subseteq\conv Y$.

(c) By (b), $\conv F$ is the set of convex combinations of points of $F$. In such a combination we may
merge the terms that use the same point of $F$ (adding their coefficients) and add the missing points of
$F$ with coefficient $0$; this gives the stated form, and conversely every point of the stated form is a
convex combination. For $F=\{y,z\}$ with $y\ne z$ this is the segment; for $y=z$ both sides are $\{y\}$.

(d) If $w_0z_0+w_1z_1<\alpha$ and $w_0z'_0+w_1z'_1<\alpha$, and $\lambda,\mu\ge0$ with $\lambda+\mu=1$,
then $w_0(\lambda z_0+\mu z'_0)+w_1(\lambda z_1+\mu z'_1)<\lambda\alpha+\mu\alpha=\alpha$, because at least
one of $\lambda,\mu$ is positive and the other term is at most the corresponding multiple of $\alpha$.
\end{proof}

\begin{lemma}[Convex position via hulls; \leanref{ErdosSzekeres.convexPosition_iff}{ErdosSzekeres/Basic.lean:15}]\label{fd:lem:cp-hull}\sloppy
A finite set $S$ of points is in convex position if and only if $p\notin\conv(S\setminus\{p\})$ for
every $p\in S$.
\end{lemma}

\begin{proof}
In Lean this is Mathlib's \lean{convexIndependent_set_iff_notMem_convexHull_diff} applied to $S$; we give
a direct proof from the form of Definition~\ref{fd:def:cp}.

Suppose $S$ is in convex position and $p\in S$. If $p\in\conv(S\setminus\{p\})$, then the definition with
$T=S\setminus\{p\}$ and $x=p$ gives $p\in S\setminus\{p\}$, which is absurd.

Conversely, suppose $p\notin\conv(S\setminus\{p\})$ for every $p\in S$. Let $T\subseteq S$ and $x\in S$
with $x\in\conv T$, and suppose $x\notin T$. Then $T\subseteq S\setminus\{x\}$, so
$x\in\conv T\subseteq\conv(S\setminus\{x\})$ by Lemma~\ref{fd:lem:hull}(a), a contradiction. Hence
$x\in T$, and $S$ is in convex position.
\end{proof}

\begin{lemma}[Monotonicity; (a) \leanref{ErdosSzekeres.ConvexPosition.mono}{ErdosSzekeres/Basic.lean:11}; (b) \leanref{ErdosSzekeres.GeneralPosition.mono}{ErdosSzekeres/Basic.lean:6}; (c) \leanref{ErdosSzekeres.HasConvexNGon.mono}{ErdosSzekeres/Basic.lean:20}; (d) \leanref{MirrorOdd.hasConvexNGon_of_le}{MIRRORETV/lean/MirrorMain.lean:2182}]\label{fd:lem:mono}\sloppy
Let $s,t,S$ be finite sets of points and $n,n'\in\N$.
\begin{enumerate}[label=(\alph*)]
\item If $s$ is in convex position and $t\subseteq s$, then $t$ is in convex position.
\item If $s$ is in general position and $t\subseteq s$, then $t$ is in general position.
\item If $\mathrm{HasConvexNGon}\,n\,s$ and $s\subseteq t$, then $\mathrm{HasConvexNGon}\,n\,t$.
\item If $\mathrm{HasConvexNGon}\,n\,S$ and $n'\le n$, then $\mathrm{HasConvexNGon}\,n'\,S$.
\end{enumerate}
\end{lemma}

\begin{proof}
(a) Let $p\in t$. Then $t\setminus\{p\}\subseteq s\setminus\{p\}$, so
$\conv(t\setminus\{p\})\subseteq\conv(s\setminus\{p\})$ (Lemma~\ref{fd:lem:hull}(a)), and the latter does
not contain $p$ by Lemma~\ref{fd:lem:cp-hull}, since $p\in s$. By Lemma~\ref{fd:lem:cp-hull} again, $t$ is
in convex position. (Lean uses Mathlib's \lean{ConvexIndependent.mono}.)

(b) Pairwise distinct $a,b,c\in t$ are pairwise distinct points of $s$, so $\{a,b,c\}$ is not collinear.
This is the Lean proof.

(c) A witness $T\subseteq s$ with $\card T=n$ in convex position satisfies $T\subseteq t$. This is the
Lean proof.

(d) Let $T\subseteq S$ with $\card T=n$ be in convex position. Since $n'\le n=\card T$, there is
$T'\subseteq T$ with $\card{T'}=n'$; it is in convex position by (a), and $T'\subseteq S$. This is the
Lean proof (it applies Lemma~\ref{fd:lem:ngon-basic}(b) to $T$ and then part (c) above). The same declaration, with the same name and
proof, also appears at \dirn{SYMLEAN/mirror/ReductionThreeClasses.lean:3741}.
\end{proof}

\begin{lemma}[Sizes of convex subsets; (a) \leanref{ErdosSzekeres.HasConvexNGon.card_le}{ErdosSzekeres/Basic.lean:25}; (b) \leanref{ErdosSzekeres.ConvexPosition.hasConvexNGon}{ErdosSzekeres/Basic.lean:31}]\label{fd:lem:ngon-basic}\sloppy
Let $s$ be a finite set of points and $n\in\N$.
\begin{enumerate}[label=(\alph*)]
\item If $\mathrm{HasConvexNGon}\,n\,s$, then $n\le\card s$.
\item If $s$ is in convex position and $n\le\card s$, then $\mathrm{HasConvexNGon}\,n\,s$.
\end{enumerate}
\end{lemma}

\begin{proof}
(a) A witness $T\subseteq s$ has $n=\card T\le\card s$. (b) Choose $T\subseteq s$ with $\card T=n$; it is
in convex position by Lemma~\ref{fd:lem:mono}(a). Both are the Lean proofs.
\end{proof}

\begin{lemma}[Forcing; (a) \leanref{ErdosSzekeres.ForcesConvexNGon.mono}{ErdosSzekeres/Basic.lean:47}; (b) \leanref{ErdosSzekeres.forcesConvexNGon_iff_exact}{ErdosSzekeres/Basic.lean:37}]\label{fd:lem:forces-basic}\sloppy
Let $N,M,n\in\N$.
\begin{enumerate}[label=(\alph*)]
\item If $\mathrm{ForcesConvexNGon}\,N\,n$ and $N\le M$, then $\mathrm{ForcesConvexNGon}\,M\,n$.
\item $\mathrm{ForcesConvexNGon}\,N\,n$ holds if and only if every finite set $S$ of points in general
  position with $\card S=N$ satisfies $\mathrm{HasConvexNGon}\,n\,S$.
\end{enumerate}
\end{lemma}

\begin{proof}
(a) If $S$ is in general position and $M\le\card S$, then $N\le\card S$. (b) The forward direction is the
case $\card S=N$. Conversely, let $S$ be in general position with $N\le\card S$. Choose $T\subseteq S$ with
$\card T=N$. By Lemma~\ref{fd:lem:mono}(b) $T$ is in general position, so $\mathrm{HasConvexNGon}\,n\,T$
by hypothesis, and $\mathrm{HasConvexNGon}\,n\,S$ by Lemma~\ref{fd:lem:mono}(c). Both are the Lean
proofs.
\end{proof}

\begin{lemma}[At most three points; written step, the last assertion \leanref{ErdosSzekeres.convexPosition_of_card_three}{ErdosSzekeres/Basic.lean:58}]\label{fd:lem:small}\sloppy
Sets of at most three points behave as follows.
\begin{enumerate}[label=(\alph*)]
\item For every finite set $S$ of points, $\mathrm{HasConvexNGon}\,0\,S$ holds.
\item Every finite set of at most two points is in convex position.
\item Let $a,b,c$ be pairwise distinct points. Then $\{a,b,c\}$ is in convex position if and only if
  $\{a,b,c\}$ is not collinear. In particular, if a finite set $s$ is in general position and
  $\card s=3$, then $s$ is in convex position.
\end{enumerate}
\end{lemma}

\begin{proof}
(a) $T=\emptyset$ is a subset of $S$ with $0$ elements, and it is in convex position by
Lemma~\ref{fd:lem:cp-hull} (there is nothing to check).

(b) We use Lemma~\ref{fd:lem:cp-hull} and Lemma~\ref{fd:lem:hull}(a). For $\emptyset$ there is nothing to
check; for $\{x\}$, $x\notin\conv\emptyset=\emptyset$; for $\{x,y\}$ with $x\ne y$,
$x\notin\conv\{y\}=\{y\}$ and $y\notin\conv\{x\}=\{x\}$.

(c) Since $a,b,c$ are pairwise distinct, $\{a,b,c\}\setminus\{a\}=\{b,c\}$, and similarly for $b$ and
$c$.

Suppose $\{a,b,c\}$ is not collinear. If $a\in\conv\{b,c\}$, then by Lemma~\ref{fd:lem:hull}(c)
$a=\lambda b+(1-\lambda)c$ for some $\lambda\in[0,1]$; with $v=b-c$ we get $a=c+\lambda v$, $b=c+1\cdot v$
and $c=c+0\cdot v$, so $\{a,b,c\}$ is collinear by Lemma~\ref{fd:lem:collinear-mem} (with $p^*=c$), a
contradiction. The same argument with the roles of $a,b,c$ exchanged shows $b\notin\conv\{a,c\}$ and
$c\notin\conv\{a,b\}$. By Lemma~\ref{fd:lem:cp-hull}, $\{a,b,c\}$ is in convex position.

Suppose $\{a,b,c\}$ is collinear. By Lemma~\ref{fd:lem:collinear-mem} with $p^*=a$ there are $v$ and
reals $r_b,r_c$ with $b=a+r_bv$ and $c=a+r_cv$; put $\tau(a)=0$, $\tau(b)=r_b$, $\tau(c)=r_c$, so that
every $x\in\{a,b,c\}$ equals $a+\tau(x)v$. Since the points are pairwise distinct, the three reals
$0,r_b,r_c$ are pairwise distinct (equal parameters would give equal points). Let $x$ be the point whose
parameter $\tau(x)$ is the middle one of the three, and let
$y,z$ be the other two points, with $\tau(y)<\tau(x)<\tau(z)$. With
$\lambda=(\tau(z)-\tau(x))/(\tau(z)-\tau(y))\in(0,1)$ we have
$\tau(x)=\lambda\tau(y)+(1-\lambda)\tau(z)$, hence $x=\lambda y+(1-\lambda)z\in\conv\{y,z\}
=\conv(\{a,b,c\}\setminus\{x\})$, and $\{a,b,c\}$ is not in convex position by
Lemma~\ref{fd:lem:cp-hull}.

Finally, if $s$ is in general position and $\card s=3$, then $s=\{a,b,c\}$ with $a,b,c$ pairwise distinct
and $\{a,b,c\}$ not collinear, so $s$ is in convex position. (The Lean proof of this last assertion
derives affine independence of $a,b,c$ from non-collinearity and uses that affinely independent families
are convex independent; we argued directly.)
\end{proof}

\begin{lemma}[Injective affine maps; (a) Mathlib \lean{AffineMap.image_convexHull}; (b$'$) \leanref{ErdosSzekeres.SymReduction.convexPosition_image_affine}{SYMLEAN/AffineReduction.lean:1263}; (c$'$) \leanref{ErdosSzekeres.SymReduction.generalPosition_image_affine}{SYMLEAN/AffineReduction.lean:1278}, \leanref{ErdosSzekeres.SymReduction.generalPosition_image_affineEquiv}{SYMLEAN/AffineReduction.lean:1290}; (d$'$) \leanref{ErdosSzekeres.SymReduction.hasConvexNGon_image_affine}{SYMLEAN/AffineReduction.lean:1298}; (e$'$) \leanref{ErdosSzekeres.SymReduction.injective_of_affine}{SYMLEAN/AffineReduction.lean:1305}; the rest written step]\label{fd:lem:affine-cp}\raggedright
Let $g\colon\R^2\to\R^2$ be an affine map, $g(x)=Ax+v$.
\begin{enumerate}[label=(\alph*)]
\item For every $Y\subseteq\R^2$, $g(\conv Y)=\conv\bigl(g(Y)\bigr)$.
\end{enumerate}
Suppose moreover that $g$ is injective. Then $\det A\ne0$, and $g$ is a bijection of $\R^2$ whose inverse
$y\mapsto A^{-1}(y-v)$ is affine. For every finite set $S$ of points and every $n\in\N$ (with $g(S)$ the
image set, $\card{g(S)}=\card S$):
\begin{enumerate}[label=(\alph*),resume]
\item $S$ is in convex position if and only if $g(S)$ is in convex position;
\item $S$ is in general position if and only if $g(S)$ is in general position;
\item $\mathrm{HasConvexNGon}\,n\,S$ holds if and only if $\mathrm{HasConvexNGon}\,n\,g(S)$ holds.
\end{enumerate}
The forward directions of (b)--(d), and the converse of the implication ``$g$ injective $\Rightarrow$
$\det A\ne0$'', are Lean theorems in the following forms. Here $S$ is a finite set of points, $n\in\N$,
$f$ is an affine map of $\R^2$ in the sense of Mathlib (\lean{AffineMap}; Definition~\ref{fd:def:conv}),
and $h\colon\R^2\to\R^2$ is a map with $h(p)=Lp+w$ for every $p\in\R^2$, where $L$ is a real $2\times2$
matrix and $w\in\R^2$ (Lean's \lean{b}).
\begin{enumerate}[label=(\alph*$'$),start=2]
\item If $S$ is in convex position and $f$ is injective, then $f(S)$ is in convex position.
\item If $\det L\ne0$ and $S$ is in general position, then $h(S)$ is in general position. If $S$ is in
  general position and $\phi$ is an affine equivalence of $\R^2$ in the sense of Mathlib
  (\lean{AffineEquiv}; Definition~\ref{fd:def:conv}), then $\phi(S)$ is in general position.
\item If $\mathrm{HasConvexNGon}\,n\,S$ and $f$ is injective, then $\mathrm{HasConvexNGon}\,n\,f(S)$.
\item If $\det L\ne0$, then $h$ is injective.
\end{enumerate}
\end{lemma}

\begin{proof}
For $x,x'\in\R^2$ and reals $\lambda,\mu$ with $\lambda+\mu=1$,
\begin{equation*}
  g(\lambda x+\mu x')=\lambda Ax+\mu Ax'+(\lambda+\mu)v=\lambda g(x)+\mu g(x'). \tag{$*$}
\end{equation*}

(a) By $(*)$, $g(\conv Y)$ is convex: any two of its points have the form $g(x),g(x')$ with
$x,x'\in\conv Y$, and for
$\lambda,\mu\ge0$ with $\lambda+\mu=1$, $\lambda g(x)+\mu g(x')=g(\lambda x+\mu x')$ with
$\lambda x+\mu x'\in\conv Y$. It contains $g(Y)$, so $\conv(g(Y))\subseteq g(\conv Y)$ by
Lemma~\ref{fd:lem:hull}(a). Conversely, the set $K=\{x:g(x)\in\conv(g(Y))\}$ is convex by $(*)$ and
contains $Y$, so $\conv Y\subseteq K$, that is, $g(\conv Y)\subseteq\conv(g(Y))$. (This is the proof of
Mathlib's \lean{AffineMap.image_convexHull}.)

Now let $g$ be injective, and write $A=\begin{pmatrix}\alpha&\beta\\ \gamma&\delta\end{pmatrix}$. If
$\det A=\alpha\delta-\beta\gamma=0$, there is $x\ne0$ with $Ax=0$: if $(\alpha,\beta)\ne(0,0)$ take
$x=(\beta,-\alpha)$, for which $Ax=(\alpha\beta-\beta\alpha,\ \gamma\beta-\delta\alpha)=(0,0)$; if
$\alpha=\beta=0$ take $x=(\delta,-\gamma)$ when $(\gamma,\delta)\ne(0,0)$, for which
$Ax=(0,\gamma\delta-\delta\gamma)=(0,0)$, and any $x\ne0$ when $A=0$. Then $g(x)=v=g(0)$ with $x\ne0$,
contradicting injectivity. So $\det A\ne0$, $A$ has the inverse
$A^{-1}=(\det A)^{-1}\begin{pmatrix}\delta&-\beta\\ -\gamma&\alpha\end{pmatrix}$, and $g$ is bijective with
inverse $y\mapsto A^{-1}(y-v)$, which is affine. Injectivity gives $\card{g(S)}=\card S$.

(b) The points of $g(S)$ are the $g(p)$ with $p\in S$, and $g(S)\setminus\{g(p)\}=g(S\setminus\{p\})$ by
injectivity. By (a), $\conv\bigl(g(S)\setminus\{g(p)\}\bigr)=g\bigl(\conv(S\setminus\{p\})\bigr)$, and by
injectivity $g(p)$ belongs to this set if and only if $p\in\conv(S\setminus\{p\})$. Lemma~\ref{fd:lem:cp-hull}
applied to $S$ and to $g(S)$ gives (b).

(c) Every triple of points of $g(S)$ has the form $(g(a),g(b),g(c))$ with $a,b,c\in S$; by injectivity
$g(a),g(b),g(c)$ are pairwise distinct if and only if $a,b,c$ are; and
$\orient(g(a),g(b),g(c))=\det(A)\,\orient(a,b,c)$ (Lemma~\ref{fd:lem:affine-orient}) with $\det A\ne0$, so
one orientation vanishes if and only if the other does. Lemma~\ref{fd:lem:collinear}(b) gives (c).

(d) If $T\subseteq S$, $\card T=n$, is in convex position, then $g(T)\subseteq g(S)$, $\card{g(T)}=n$, and
$g(T)$ is in convex position by (b). Conversely, let $T'\subseteq g(S)$ with $\card{T'}=n$ be in convex
position, and let $T=\{p\in S:g(p)\in T'\}$. Since every point of $T'$ is $g(p)$ for some $p\in S$, we have
$g(T)=T'$; hence $\card T=\card{T'}=n$ by injectivity, $T$ is in convex position by (b), and $T\subseteq S$.

(b$'$)--(e$'$). By Definition~\ref{fd:def:conv}, a Mathlib affine map $f$ of $\R^2$ has the form
$f(x)=A_fx+f(0)$, where $A_f$ is the matrix of its linear part (in Lean:
\lean{ErdosSzekeres.SymReduction.affineMap_eq_mulVec}, \dirn{SYMLEAN/AffineReduction.lean:1284}). So
(b$'$) and (d$'$) are the forward directions of (b) and (d) for $g=f$. (e$'$): if $h(p)=h(q)$, then
$L(p-q)=(Lp+w)-(Lq+w)=h(p)-h(q)=0$. Write $L=\begin{pmatrix}\alpha&\beta\\ \gamma&\delta\end{pmatrix}$; the
matrix $L^\sharp=\begin{pmatrix}\delta&-\beta\\ -\gamma&\alpha\end{pmatrix}$ satisfies
$L^\sharp L=(\det L)\,I$ (multiply out), so $(\det L)(p-q)=L^\sharp L(p-q)=0$, and $p=q$ because
$\det L\ne0$. (c$'$): if $\det L\ne0$, then $h$ is an affine map with linear part $L$, and it is injective
by (e$'$), so the forward direction of (c) applies to $g=h$. A Mathlib affine equivalence $\phi$ of
$\R^2$ is an injective affine map (Definition~\ref{fd:def:conv}), so the forward direction of (c) applies
to $g=\phi$ as well. (The Lean proofs follow the same routes:
\lean{SymReduction.convexPosition_image_affine} uses \lean{AffineMap.image_convexHull} and injectivity as
in (b); \lean{SymReduction.generalPosition_image_affine} uses \lean{SymReduction.orient_affine}, that is,
Lemma~\ref{fd:lem:affine-orient}; \lean{SymReduction.generalPosition_image_affineEquiv} derives
$\det\ne0$ for the matrix of the linear part of $\phi$ from the invertibility of that linear part.)
\end{proof}

\begin{remark}[Special cases in the main library]\label{fd:rem:affine-cp-lean}\sloppy
Besides the general forms (b$'$)--(e$'$), the main library proves special cases of
Lemma~\ref{fd:lem:affine-cp} in \dirn{ErdosSzekeres/Suk/Reflection.lean}.
\lean{ErdosSzekeres.HasConvexNGon.linearEquiv} (\dirn{ErdosSzekeres/Suk/Reflection.lean:6}) transports
$\mathrm{HasConvexNGon}\,n\,S$ to $\mathrm{HasConvexNGon}\,n\,e(S)$ for a linear bijection $e$ of the
plane (Mathlib \lean{LinearEquiv}): the forward direction of (d) for $g=e$. The reflection
$\sigma(p_0,p_1)=(p_0,-p_1)$ is the linear bijection \lean{ErdosSzekeres.Suk.reflectY}
(\dirn{ErdosSzekeres/Suk/Reflection.lean:34}); the simplification lemmas \lean{Suk.reflectY_x} and
\lean{Suk.reflectY_y} (\dirn{ErdosSzekeres/Suk/Reflection.lean:42} and \dirn{:43}) give its two
coordinates, and \lean{Suk.reflectY_reflectY} (\dirn{ErdosSzekeres/Suk/Reflection.lean:44}) says that
$\sigma(\sigma(p))=p$. \lean{ErdosSzekeres.Suk.generalPosition_reflectY}
(\dirn{ErdosSzekeres/Suk/Reflection.lean:52}) transports general position from $S$ to $\sigma(S)$ (the
forward direction of (c) for $g=\sigma$), and \lean{ErdosSzekeres.Suk.hasConvexNGon_of_reflectY}
(\dirn{ErdosSzekeres/Suk/Reflection.lean:64}) says that $\mathrm{HasConvexNGon}\,n\,\sigma(S)$ implies
$\mathrm{HasConvexNGon}\,n\,S$ (the reverse direction of (d) for $g=\sigma$).
\end{remark}

\begin{lemma}[Carath\'eodory's theorem in the plane; written step]\label{fd:lem:caratheodory}
Let $Y\subseteq\R^2$ and $x\in\conv Y$. Then there is a finite, affinely independent set $Y'\subseteq Y$
with $\card{Y'}\le3$ and $x\in\conv Y'$.
\end{lemma}

\begin{proof}
By Lemma~\ref{fd:lem:hull}(b), $x$ is a convex combination of finitely many points of $Y$; if $T_0$ is the
finite set of these points, then $x\in\conv T_0$ by Lemma~\ref{fd:lem:hull}(c). Among all finite sets
$T\subseteq Y$ with $x\in\conv T$ choose one, $T=\{y_1,\dots,y_k\}$ with distinct $y_i$, of least
cardinality $k$. Then $k\ge1$ because $\conv\emptyset=\emptyset$. By Lemma~\ref{fd:lem:hull}(c),
$x=\sum_i\lambda_iy_i$ with $\lambda_i\ge0$ and $\sum_i\lambda_i=1$. Every $\lambda_j$ is positive: if
$k=1$ then $\lambda_1=1$, and if $k\ge2$ and $\lambda_j=0$, then $x$ is a convex combination of the points
of $T\setminus\{y_j\}$, so $x\in\conv(T\setminus\{y_j\})$, contradicting minimality.

$T$ is affinely independent. Otherwise there are reals $\mu_1,\dots,\mu_k$, not all zero, with
$\sum_i\mu_i=0$ and $\sum_i\mu_iy_i=0$. Some $\mu_i$ is positive (they sum to $0$ and are not all zero).
Let $t=\min\{\lambda_i/\mu_i:\mu_i>0\}>0$, attained at $i=j$, and put $\lambda'_i=\lambda_i-t\mu_i$. Then
$\lambda'_i\ge0$ for all $i$ (if $\mu_i\le0$ then $\lambda'_i\ge\lambda_i$; if $\mu_i>0$ then
$t\le\lambda_i/\mu_i$), $\lambda'_j=0$, $\sum_i\lambda'_i=1-t\cdot0=1$ and
$\sum_i\lambda'_iy_i=x-t\cdot0=x$. So $x$ is a convex combination of the points of $T\setminus\{y_j\}$,
contradicting minimality again (note $k\ge2$ here, since $\sum\lambda'_i=1$ forces some
$\lambda'_i\ne0$, $i\ne j$).

Finally $k\le3$: if $k\ge4$, the $k-1\ge3$ vectors $y_2-y_1,\dots,y_k-y_1$ in the two-dimensional space
$\R^2$ are linearly dependent, $\sum_{i\ge2}\nu_i(y_i-y_1)=0$ with $\nu\ne0$; then $\mu_1=-\sum_{i\ge2}\nu_i$
and $\mu_i=\nu_i$ ($i\ge2$) satisfy $\sum_i\mu_i=0$ and $\sum_i\mu_iy_i=\sum_{i\ge2}\nu_i(y_i-y_1)=0$, with
$\mu\ne0$, contradicting affine independence. Take $Y'=T$.

The Lean code uses Mathlib's version of this theorem (module \lean{Mathlib.Analysis.Convex.Caratheodory},
with \lean{Caratheodory.minCardFinsetOfMemConvexHull}), whose proof is the same minimal-cardinality
argument; the bound $\card{Y'}\le3$ is derived in Lean from affine independence as
$\card{Y'}\le\dim+1=3$.
\end{proof}

\begin{lemma}[Four-point criteria; (a) \leanref{ErdosSzekeres.Suk.convexPosition_of_subsets_card_le_four}{ErdosSzekeres/Suk/FourPoint.lean:10}; (b) \leanref{ErdosSzekeres.Suk.convexPosition_of_four_subsets}{ErdosSzekeres/Suk/FourPoint.lean:40}]\label{fd:lem:four-subsets}\raggedright
Let $s$ be a finite set of points.
\begin{enumerate}[label=(\alph*)]
\item If every $t\subseteq s$ with $\card t\le4$ is in convex position, then $s$ is in convex position.
\item If $4\le\card s$ and every $t\subseteq s$ with $\card t=4$ is in convex position, then $s$ is in
  convex position.
\end{enumerate}
No general-position assumption is made.
\end{lemma}

\begin{proof}
We follow the Lean proofs. (a) By Lemma~\ref{fd:lem:cp-hull} we must show $p\notin\conv(s\setminus\{p\})$
for every $p\in s$. Suppose $p\in\conv(s\setminus\{p\})$. By Lemma~\ref{fd:lem:caratheodory} there is a
finite $Y'\subseteq s\setminus\{p\}$ with $\card{Y'}\le3$ and $p\in\conv Y'$. Let $u=Y'\cup\{p\}\subseteq s$.
Since $p\notin Y'$, $\card u=\card{Y'}+1\le4$ and $u\setminus\{p\}=Y'$. By hypothesis $u$ is in convex
position, so $p\notin\conv(u\setminus\{p\})=\conv Y'$ by Lemma~\ref{fd:lem:cp-hull}, a contradiction.

(b) Let $t\subseteq s$ with $\card t\le4$. Since $\card s\ge4$, we can add $4-\card t$ points of
$s\setminus t$ to $t$ (there are $\card s-\card t\ge4-\card t$ of them) and obtain $u$ with
$t\subseteq u\subseteq s$ and $\card u=4$. By hypothesis $u$ is in convex position, hence so is $t$
(Lemma~\ref{fd:lem:mono}(a)). Now (a) applies.
\end{proof}

\begin{lemma}[Half-planes of the orientation; (a) \leanref{ErdosSzekeres.Suk.orient_affine_combination}{ErdosSzekeres/Suk/SupportingLines.lean:6}; (b) \leanref{ErdosSzekeres.Suk.convex_orient_positive}{ErdosSzekeres/Suk/SupportingLines.lean:11}; (c) \leanref{ErdosSzekeres.Suk.convex_orient_nonnegative}{ErdosSzekeres/Suk/SupportingLines.lean:21}]\label{fd:lem:halfplane}\raggedright
The orientation is affine in its third argument, and its half-planes are convex.
\begin{enumerate}[label=(\alph*)]
\item For $p,q,a,b\in\R^2$ and $u,v\in\R$ with $u+v=1$:
  $\orient(p,q,u\,a+v\,b)=u\,\orient(p,q,a)+v\,\orient(p,q,b)$.
\item For $p,q\in\R^2$ the set $\{r\in\R^2:0<\orient(p,q,r)\}$ is convex.
\item For $p,q\in\R^2$ the set $\{r\in\R^2:0\le\orient(p,q,r)\}$ is convex.
\end{enumerate}
\end{lemma}

\begin{proof}
(a) $\orient(p,q,r)=(q_0-p_0)(r_1-p_1)-(q_1-p_1)(r_0-p_0)$ is a linear function of $r-p$. For
$r=ua+vb$ with $u+v=1$ we have $r-p=u(a-p)+v(b-p)$, and the claim follows. (Lean closes the same
polynomial identity with the certificate $(p_1q_0-p_0q_1)(u+v-1)$.)

(b) Let $\orient(p,q,a)>0$, $\orient(p,q,b)>0$ and $u,v\ge0$ with $u+v=1$. By (a) the value at $ua+vb$ is
$u\,\orient(p,q,a)+v\,\orient(p,q,b)$. If $u=0$ then $v=1$ and this is $\orient(p,q,b)>0$; otherwise
$u>0$, the first term is positive and the second is nonnegative.

(c) By (a) the value at $ua+vb$ is a sum of two nonnegative terms.
\end{proof}

\begin{lemma}[Strict supporting lines; \leanref{ErdosSzekeres.Suk.convexPosition_of_strict_supports}{ErdosSzekeres/Suk/SupportingLines.lean:28}]\label{fd:lem:strict-supports}\sloppy
Let $s$ be a finite set of points such that for every $p\in s$ there is $q\in\R^2$ with
$0<\orient(p,q,r)$ for every $r\in s$ with $r\ne p$. Then $s$ is in convex position.
\end{lemma}

\begin{proof}
We follow the Lean proof. Let $p\in s$ and choose $q$ as in the hypothesis. The set
$H=\{r:0<\orient(p,q,r)\}$ is convex (Lemma~\ref{fd:lem:halfplane}(b)) and contains $s\setminus\{p\}$, so
$\conv(s\setminus\{p\})\subseteq H$ by Lemma~\ref{fd:lem:hull}(a). But $\orient(p,q,p)=0$
(Lemma~\ref{fd:lem:orient}(e)), so $p\notin H$ and hence $p\notin\conv(s\setminus\{p\})$. By
Lemma~\ref{fd:lem:cp-hull}, $s$ is in convex position.
\end{proof}

\begin{lemma}[A support identity; \leanref{ErdosSzekeres.Suk.orient_support_sum}{ErdosSzekeres/Suk/SupportingLines.lean:40}]\label{fd:lem:support-sum}\sloppy
For all $a,p,b,r\in\R^2$: $\orient(p,\,p+b-a,\,r)=\orient(a,p,r)+\orient(p,b,r)$.
\end{lemma}

\begin{proof}
For $w\in\R^2$, $\orient(p,p+w,r)=\det(w\mid r-p)$ is linear in $w$. With $w=b-a=(p-a)+(b-p)$ we get
$\orient(p,p+b-a,r)=\det(p-a\mid r-p)+\det(b-p\mid r-p)$. The second term is $\orient(p,b,r)$. For the
first, $\orient(a,p,r)=\det(p-a\mid r-a)=\det(p-a\mid (r-p)+(p-a))=\det(p-a\mid r-p)$, since a determinant
with two equal columns vanishes. (The Lean proof expands both sides.)
\end{proof}

\begin{lemma}[Four counterclockwise points; \leanref{ErdosSzekeres.Suk.convexPosition_four_of_positive}{ErdosSzekeres/Suk/SupportingLines.lean:47}]\label{fd:lem:four-positive}\sloppy
Let $a,b,c,d\in\R^2$ with
\[
  0<\orient(a,b,c),\qquad 0<\orient(a,b,d),\qquad 0<\orient(a,c,d),\qquad 0<\orient(b,c,d).
\]
Then the finite set $\{a,b,c,d\}$ is in convex position.
\end{lemma}

\begin{proof}
We follow the Lean proof, which exhibits a strict supporting line at each of the four points and applies
Lemma~\ref{fd:lem:strict-supports}. (The hypotheses imply, by Lemma~\ref{fd:lem:orient}(e), that the four
points are pairwise distinct, but the argument does not need this.) Write $o_{abc}=\orient(a,b,c)$,
$o_{abd}=\orient(a,b,d)$, $o_{acd}=\orient(a,c,d)$ and $o_{bcd}=\orient(b,c,d)$, all positive. By
Lemma~\ref{fd:lem:orient}(b) every cyclic rotation of the arguments leaves $\orient$ unchanged, and by
Lemma~\ref{fd:lem:orient}(e) $\orient$ vanishes when two arguments are equal. For $p\in\{a,b,c,d\}$ and
$r\in\{a,b,c,d\}$ with $r\ne p$, the point $r$ is one of the three letters other than $p$ (as a point, $r$
differs from $p$), so it suffices to check the three letters other than $p$. Each line below applies
Lemma~\ref{fd:lem:support-sum}.
\begin{itemize}
\item $p=a$, $q=a+b-d$: $\orient(a,q,r)=\orient(d,a,r)+\orient(a,b,r)$. For $r=b$:
  $\orient(d,a,b)+0=o_{abd}$. For $r=c$: $\orient(d,a,c)+o_{abc}=o_{acd}+o_{abc}$. For $r=d$:
  $0+o_{abd}$.
\item $p=b$, $q=b+c-a$: $\orient(b,q,r)=\orient(a,b,r)+\orient(b,c,r)$. For $r=a$:
  $0+\orient(b,c,a)=o_{abc}$. For $r=c$: $o_{abc}+0$. For $r=d$: $o_{abd}+o_{bcd}$.
\item $p=c$, $q=c+d-b$: $\orient(c,q,r)=\orient(b,c,r)+\orient(c,d,r)$. For $r=a$:
  $\orient(b,c,a)+\orient(c,d,a)=o_{abc}+o_{acd}$. For $r=b$: $0+\orient(c,d,b)=o_{bcd}$. For $r=d$:
  $o_{bcd}+0$.
\item $p=d$, $q=d+a-c$: $\orient(d,q,r)=\orient(c,d,r)+\orient(d,a,r)$. For $r=a$:
  $\orient(c,d,a)+0=o_{acd}$. For $r=b$: $\orient(c,d,b)+\orient(d,a,b)=o_{bcd}+o_{abd}$. For $r=c$:
  $0+\orient(d,a,c)=o_{acd}$.
\end{itemize}
All twelve values are positive, so each of $a,b,c,d$ has a strict supporting line, and
Lemma~\ref{fd:lem:strict-supports} shows that $\{a,b,c,d\}$ is in convex position.
\end{proof}

\begin{lemma}[Positive triples; \leanref{ErdosSzekeres.Suk.convexPosition_of_positive_triples}{ErdosSzekeres/Suk/OrientationClique.lean:9}]\label{fd:lem:pos-triples}\sloppy
Let $\iota$ be a type with a linear order $<$, let $p\colon\iota\to\R^2$ be injective, and let $s$ be a
finite subset of $\iota$ with $4\le\card s$. Suppose that $0<\orient(p(i),p(j),p(k))$ for all
$i,j,k\in s$ with $i<j$ and $j<k$. Then the finite set $p(s)=\{p(i):i\in s\}$ is in convex position.
\end{lemma}

\begin{proof}
We follow the Lean proof. Since $p$ is injective, $\card{p(s)}=\card s\ge4$, so by
Lemma~\ref{fd:lem:four-subsets}(b) it suffices to show that every $t\subseteq p(s)$ with $\card t=4$ is in
convex position. Let $u=\{i\in s:p(i)\in t\}$. Since $t\subseteq p(s)$, $p(u)=t$, and $\card u=\card t=4$ by
injectivity. List the elements of the finite subset $u$ of the linearly ordered set $\iota$ in increasing
order, $e_0<e_1<e_2<e_3$. Then $t=\{p(e_0),p(e_1),p(e_2),p(e_3)\}$, and the hypothesis applied to the
increasing triples $(e_0,e_1,e_2)$, $(e_0,e_1,e_3)$, $(e_0,e_2,e_3)$ and $(e_1,e_2,e_3)$ of elements of
$s$ gives exactly the four hypotheses of Lemma~\ref{fd:lem:four-positive} for
$(a,b,c,d)=(p(e_0),p(e_1),p(e_2),p(e_3))$. Hence $t$ is in convex position.
\end{proof}

\begin{remark}\label{fd:rem:pos-triples}
(i) Only $s$ is assumed finite; $\iota$ may be infinite, and injectivity is assumed for $p$ on all of
$\iota$. (ii) The Lean proof discharges the order facts $0<1<2<3$ in $\mathrm{Fin}\,4$ and the equality
of the set of all elements of $\mathrm{Fin}\,4$ with $\{0,1,2,3\}$ by \texttt{decide}, which the kernel
evaluates without any additional axiom. These are side facts inside the proof of a statement that is not
itself a finite computation, and the written proof above needs no machine computation; so the statement
is not computer-assisted in the sense of the rigor labels. (iii) The
hypothesis $4\le\card s$ is needed by the Lean proof (through Lemma~\ref{fd:lem:four-subsets}(b)). As a
variant, not part of the Lean statement: the conclusion also holds when $\card s\le3$, by
Lemma~\ref{fd:lem:small}(b) when $\card{p(s)}\le2$, and when $\card s=3$ because the three points then
have nonzero orientation, hence are pairwise distinct and not collinear
(Lemmas~\ref{fd:lem:orient}(e) and~\ref{fd:lem:collinear}(a)), so Lemma~\ref{fd:lem:small}(c) applies.
\end{remark}

\begin{lemma}[Negative triples; \leanref{ErdosSzekeres.Suk.convexPosition_of_negative_triples}{ErdosSzekeres/Suk/OrientationClique.lean:33}]\label{fd:lem:neg-triples}\sloppy
Let $\iota$ be a type with a linear order $<$, let $p\colon\iota\to\R^2$ be injective, and let $s$ be a
finite subset of $\iota$ with $4\le\card s$. Suppose that $\orient(p(i),p(j),p(k))<0$ for all
$i,j,k\in s$ with $i<j$ and $j<k$. Then $p(s)$ is in convex position.
\end{lemma}

\begin{proof}
We follow the Lean proof, which applies Lemma~\ref{fd:lem:pos-triples} to the reversed order. Let
$i\prec j$ mean $j<i$; this is a linear order on $\iota$. If $i,j,k\in s$ with $i\prec j\prec k$, then
$k<j<i$, so $\orient(p(k),p(j),p(i))<0$ by hypothesis, and by Lemma~\ref{fd:lem:orient}(d) (the
transposition of the first and third arguments has sign $-1$)
$\orient(p(i),p(j),p(k))=-\orient(p(k),p(j),p(i))>0$. Lemma~\ref{fd:lem:pos-triples} for $(\iota,\prec)$
gives the claim.
\end{proof}

\begin{lemma}[Positive cycles; \leanref{ErdosSzekeres.hasConvexNGon_of_positive_cycle}{ErdosSzekeres/ConvexCycleCuts.lean:11}]\label{fd:lem:pos-cycle}\sloppy
Let $N,k\in\N$ with $4\le k$. Let $p\colon\mathrm{Fin}\,N\to\R^2$ be injective, let
$e\colon\mathrm{Fin}\,k\to\mathrm{Fin}\,N$ be an embedding (an injective map; Lean type
\texttt{Fin k}~$\hookrightarrow$~\texttt{Fin N}), and let $P$ be a finite set of points with $p(e(i))\in P$
for every $i\in\mathrm{Fin}\,k$. Suppose that
\[
  0<\orient\bigl(p(e(i)),\,p(e(j)),\,p(e(l))\bigr)\qquad\text{for all } i,j,l\in\mathrm{Fin}\,k
  \text{ with } i<j \text{ and } j<l .
\]
Then $\mathrm{HasConvexNGon}\,k\,P$.
\end{lemma}

\begin{proof}
We follow the Lean proof. The map $q=p\circ e\colon\mathrm{Fin}\,k\to\R^2$ is injective, as a composition of
injective maps. Let $T=q(\mathrm{Fin}\,k)=\{q(0),\dots,q(k-1)\}$. Then $T\subseteq P$ by hypothesis and
$\card T=k$ by injectivity. Lemma~\ref{fd:lem:pos-triples}, applied to $\iota=\mathrm{Fin}\,k$ with its
natural order, the injective map $q$ and $s=\mathrm{Fin}\,k$ (so $\card s=k\ge4$), shows that $T=q(s)$ is
in convex position; its hypothesis is exactly the assumed positivity. So $T$ witnesses
$\mathrm{HasConvexNGon}\,k\,P$.
\end{proof}

\begin{remark}\label{fd:rem:pos-cycle}\sloppy
The hypothesis concerns all increasing triples of the ordering $e$, not only consecutive ones. The name
refers to the clause family \lean{ConvexCycleCut} (\dirn{ErdosSzekeres/ConvexCycleCuts.lean:8}), which is
used as cuts in the SMT coordinate-repair searches described in \dirn{docs/SAT.md:168}--\dirn{233}:
\lean{convexCycleCut_of_avoids}
(\dirn{ErdosSzekeres/ConvexCycleCuts.lean:31}) is the
contrapositive of Lemma~\ref{fd:lem:pos-cycle}: under the same hypotheses on $k$, $p$, $e$ and $P$
(without the positivity hypothesis), if $P$ is $k$-free then some $i<j<l$ in $\mathrm{Fin}\,k$ has
$\orient(p(e(i)),p(e(j)),p(e(l)))\le0$.
\end{remark}

\begin{lemma}[Separation of a point from a finite set; written step]\label{fd:lem:separation}
Let $F\subseteq\R^2$ be finite and $x\in\R^2$. Then $x\notin\conv F$ if and only if there is a linear
functional $\ell\colon\R^2\to\R$ (that is, $\ell(z)=w_0z_0+w_1z_1$ for some $w\in\R^2$) with
$\ell(x)>\ell(y)$ for every $y\in F$.
\end{lemma}

\begin{proof}
Write $\langle w,z\rangle=w_0z_0+w_1z_1$ and $\lVert z\rVert^2=\langle z,z\rangle$ (the Euclidean norm,
written out).

($\Leftarrow$, elementary.) Let $\ell(z)=\langle w,z\rangle$ satisfy $\ell(x)>\ell(y)$ for all $y\in F$.
The set $H=\{z:\ell(z)<\ell(x)\}$ is convex by Lemma~\ref{fd:lem:hull}(d) and contains $F$, so
$\conv F\subseteq H$ by Lemma~\ref{fd:lem:hull}(a); since $x\notin H$, $x\notin\conv F$. (If $F=\emptyset$
this is also clear from $\conv\emptyset=\emptyset$.)

($\Rightarrow$.) Suppose $x\notin\conv F$. If $F=\emptyset$, the functional $\ell=0$ satisfies the
condition vacuously. Otherwise write $F=\{y_1,\dots,y_m\}$ with $m\ge1$ distinct points and let
$\Delta=\{\lambda\in\R^m:\lambda_i\ge0,\ \sum_i\lambda_i=1\}$. The set $\Delta$ is closed and bounded, hence
compact by the Heine--Borel theorem (a closed and bounded subset of $\R^m$ is compact; W.~Rudin,
\emph{Principles of Mathematical Analysis}, 3rd ed., Theorem~2.41), and the function
$h(\lambda)=\lVert x-\sum_i\lambda_iy_i\rVert^2$ is a polynomial in $\lambda$, hence continuous; so $h$
attains its minimum on the nonempty compact set $\Delta$ at some $\lambda^*$ (a continuous real function on
a nonempty compact metric space attains its minimum; Rudin, op.\ cit., Theorem~4.16). Put
$q=\sum_i\lambda^*_iy_i$. By Lemma~\ref{fd:lem:hull}(c), $\conv F=\{\sum_i\lambda_iy_i:\lambda\in\Delta\}$,
so $q\in\conv F$ and
\begin{equation*}
  \lVert x-q\rVert^2\le\lVert x-z\rVert^2\qquad\text{for every }z\in\conv F. \tag{$**$}
\end{equation*}
Since $x\notin\conv F$, $x\ne q$; put $w=x-q\ne0$ and $\ell(z)=\langle w,z\rangle$. Let $y\in F$ and
$0<t\le1$. The point $z_t=(1-t)q+ty$ lies in $\conv F$ (convexity, Lemma~\ref{fd:lem:hull}(a)), and
$x-z_t=w-t(y-q)$, so by $(**)$
\[
  \lVert w\rVert^2\le\lVert w-t(y-q)\rVert^2=\lVert w\rVert^2-2t\langle w,y-q\rangle+t^2\lVert y-q\rVert^2 .
\]
Hence $2\langle w,y-q\rangle\le t\lVert y-q\rVert^2$ for every $t\in(0,1]$, and letting $t\to0$ gives
$\langle w,y-q\rangle\le0$, that is, $\ell(y)\le\ell(q)$. Finally
$\ell(x)-\ell(q)=\langle w,x-q\rangle=\lVert w\rVert^2>0$, so $\ell(y)\le\ell(q)<\ell(x)$.
\end{proof}

\paragraph{Coverage.}
{\raggedright
\begin{itemize}
\item \emph{Covered definitions, with the label of the statement.} \lean{ErdosSzekeres.Point}
  (\ref{fd:def:point}), \lean{ErdosSzekeres.orient} (\ref{fd:def:orient}),
  \lean{ErdosSzekeres.GeneralPosition} (\ref{fd:def:gp}), \lean{ErdosSzekeres.ConvexPosition}
  (\ref{fd:def:cp}), \lean{ErdosSzekeres.HasConvexNGon} (\ref{fd:def:ngon}),
  \lean{ErdosSzekeres.ForcesConvexNGon} (\ref{fd:def:forces}), the three copies of \lean{Invariant}
  (\ref{fd:def:inv}) and the three cited copies of \lean{rotAbout} (\ref{fd:def:rot}), together with the
  Mathlib notions \lean{Convex}, \lean{convexHull}, \lean{AffineIndependent}, \lean{AffineMap},
  \lean{AffineEquiv} (\ref{fd:def:conv}), \lean{Collinear} (\ref{fd:def:gp}) and \lean{ConvexIndependent}
  (\ref{fd:def:cp}).
\item \emph{Covered facts, with the label of the statement.} \lean{orient_swap}, \lean{orient_cycle},
  \lean{orient_translate} (\ref{fd:lem:orient}); \lean{SymReduction.orient_affine} and, for part (i),
  \lean{OrbitConsistency.orient_rotAbout_inv} (\ref{fd:lem:affine-orient});
  \lean{orient_eq_zero_iff_collinear}, \lean{generalPosition_iff_orient_ne_zero} (\ref{fd:lem:collinear});
  \lean{convexPosition_iff} (\ref{fd:lem:cp-hull}); \lean{ConvexPosition.mono}, \lean{GeneralPosition.mono},
  \lean{HasConvexNGon.mono}, \lean{MirrorOdd.hasConvexNGon_of_le} (\ref{fd:lem:mono});
  \lean{HasConvexNGon.card_le}, \lean{ConvexPosition.hasConvexNGon} (\ref{fd:lem:ngon-basic});
  \lean{ForcesConvexNGon.mono}, \lean{forcesConvexNGon_iff_exact} (\ref{fd:lem:forces-basic});
  \lean{convexPosition_of_card_three} (\ref{fd:lem:small}); \lean{SymReduction.convexPosition_image_affine},
  \lean{SymReduction.generalPosition_image_affine}, \lean{SymReduction.generalPosition_image_affineEquiv},
  \lean{SymReduction.hasConvexNGon_image_affine}, \lean{SymReduction.injective_of_affine}
  (\ref{fd:lem:affine-cp}, parts (b$'$)--(e$'$)); \lean{Suk.convexPosition_of_subsets_card_le_four},
  \lean{Suk.convexPosition_of_four_subsets} (\ref{fd:lem:four-subsets}); \lean{Suk.orient_affine_combination},
  \lean{Suk.convex_orient_positive}, \lean{Suk.convex_orient_nonnegative} (\ref{fd:lem:halfplane});
  \lean{Suk.convexPosition_of_strict_supports} (\ref{fd:lem:strict-supports}); \lean{Suk.orient_support_sum}
  (\ref{fd:lem:support-sum}); \lean{Suk.convexPosition_four_of_positive} (\ref{fd:lem:four-positive});
  \lean{Suk.convexPosition_of_positive_triples} (\ref{fd:lem:pos-triples});
  \lean{Suk.convexPosition_of_negative_triples} (\ref{fd:lem:neg-triples});
  \lean{hasConvexNGon_of_positive_cycle} (\ref{fd:lem:pos-cycle}); \lean{LayerLemma.iterate_mem}
  (\ref{fd:lem:inv-image}); \lean{LayerLemma.rotAbout_eq_crot}, \lean{LayerLemma.rotAff},
  \lean{LayerLemma.rotAbout_iterate}, \lean{LayerLemma.rotAbout_injective} (\ref{fd:lem:rot}).
\item \emph{Written steps.} The following statements and parts carry no Lean tag in this section, and
  their proofs above are written steps. We found no declaration in the files of the table of sources that
  states one of them in the form given here. Several of them have Lean counterparts in another form,
  special cases or variants, in the project or in Mathlib; we name some of these, here or in the remark
  indicated, but the research records contain many further copies and variants of such elementary facts,
  and we do not list them all. Naming a counterpart does not make a written step Lean-checked in the sense
  of the rigor labels; we give axioms evidence for such a counterpart only where we say so. The written
  steps are:
  \begin{itemize}
  \item \ref{fd:lem:orient}(d) for a general commutative ring (its real case is
    \lean{OrbitConsistency.orient3_perm}, Remark~\ref{fd:rem:orient-lean}) and (e);
  \item part (ii) of \ref{fd:lem:affine-orient} (its case $o=0$, $\varphi=0$ is \lean{Suk.orient_reflectY},
    Remark~\ref{fd:rem:affine-orient-lean}; with the line given by a unit direction vector instead of the
    angle $\varphi$, (ii) is stated in the research file \dirn{H2/D3Encoding/Geometry.lean} as
    \lean{D3.orient_reflC} (\dirn{:87}));
  \item \ref{fd:lem:collinear-mem} (Mathlib's \lean{collinear_iff_of_mem}) and \ref{fd:lem:hull};
  \item \ref{fd:lem:inv-image}(a),(b) (the argument of (a) is a step of the proof of
    \lean{LayerLemma.map_mem_hullVerts});
  \item in \ref{fd:lem:rot}, the matrix form of Definition~\ref{fd:def:rot} and the identity
    $\rho_\theta(c)=c$ in (b) (the latter a step of the proof of \lean{LayerLemma.invariant_erase_centre}),
    the composition law and $\rho_0=\mathrm{id}$ in (c), the bijectivity and the inverse in (d), and (e),
    (f) (the identity $\rho_\theta(c)=c$ and the identity of (e) are stated, for the copy of
    \lean{rotAbout} in the review file \dirn{SYMLEAN/review/ReviewSanity.lean}, as
    \lean{SymReduction.Review.rotAbout_center} (\dirn{:1086}) and \lean{SymReduction.Review.orient_rotAbout}
    (\dirn{:1091}, which writes $(p_0-c_0)^2+(p_1-c_1)^2$ as \texttt{eDist2 p c}); we do not cite them as
    Lean-checked, because none of the theorems whose axioms the saved output
    \dirn{SYMLEAN/logs/ReviewSanity.out} of that file prints uses them);
  \item \ref{fd:lem:small}(a)--(c) apart from its last assertion;
  \item in \ref{fd:lem:affine-cp}, part (a) (Mathlib's \lean{AffineMap.image_convexHull}), the implication
    ``injective $\Rightarrow$ $\det A\ne0$'' with the inverse, and the reverse directions of (b)--(d) (for
    the reflection $(p_0,p_1)\mapsto(p_0,-p_1)$, the reverse direction of (d) is
    \lean{Suk.hasConvexNGon_of_reflectY}; Remark~\ref{fd:rem:affine-cp-lean} also names the
    forward-direction special cases \lean{HasConvexNGon.linearEquiv} and \lean{Suk.generalPosition_reflectY});
  \item \ref{fd:lem:caratheodory} (Mathlib's Carath\'eodory module);
  \item \ref{fd:lem:separation} (a variant of its easy direction is the main-library theorem
    \lean{Omission.BoundaryGeometry.exposed_not_in_hull},
    \dirn{ErdosSzekeres/Omission/BoundaryGeometry.lean:30}: if $S$ is a set of points, $p$ a point and
    $A,B,C$ reals with $Ap_0+Bp_1+C=0$ and $Aq_0+Bq_1+C>0$ for every $q\in S\setminus\{p\}$, then
    $p\notin\conv(S\setminus\{p\})$).
  \end{itemize}
\item \emph{Other declarations named.} Other declarations of the project that are named in statements,
  remarks, proofs or this paragraph as Lean counterparts, as definitions and lemmas that these rely on, or
  as the location of a proof step: \lean{OrbitConsistency.orient3_perm} and \lean{OrbitConsistency.orient3},
  \lean{Suk.orient_reflectY}, \lean{HasConvexNGon.linearEquiv}, \lean{Suk.reflectY}, \lean{Suk.reflectY_x},
  \lean{Suk.reflectY_y}, \lean{Suk.reflectY_reflectY}, \lean{Suk.generalPosition_reflectY},
  \lean{Suk.hasConvexNGon_of_reflectY}, \lean{ConvexCycleCut}, \lean{convexCycleCut_of_avoids},
  \lean{SymReduction.affineMap_eq_mulVec}, \lean{LayerLemma.map_mem_hullVerts},
  \lean{LayerLemma.invariant_erase_centre}, \lean{D3.orient_reflC}, \lean{SymReduction.Review.rotAbout_center},
  \lean{SymReduction.Review.orient_rotAbout} and \lean{Omission.BoundaryGeometry.exposed_not_in_hull}.
\item \emph{Deliberately not covered.} Of the declarations of the main-library files in the table of
  sources, all others are listed here: in \dirn{ErdosSzekeres/Definitions.lean}, the bound function
  \lean{conjecturedBound} ($n\mapsto2^{\max(n-2,0)}+1$) and the open conjectures
  \lean{ErdosSzekeresConjecture} and \lean{HeptagonConjecture} (propositions that are defined but not
  asserted), none of which is used by the facts of this section; in \dirn{ErdosSzekeres/Orientation.lean}
  and \dirn{ErdosSzekeres/OrientationGeometry.lean}, \lean{rationalToReal}, \lean{orient_rationalToReal},
  \lean{orient_rationalToReal_ne_zero}, \lean{RationalGeneralPosition} and
  \lean{generalPosition_of_rational} (transfer of exact rational certificates to real points, outside the
  foundations list), and the private helper \lean{exists_smul_of_det_zero} (absorbed into the proof of
  Lemma~\ref{fd:lem:collinear}); in \dirn{ErdosSzekeres/Basic.lean},
  \lean{convexIndependent_of_affineIndependent} (a bridge to Mathlib's affine independence, used in the
  main library only for \lean{convexPosition_of_card_three}, which we prove directly) and
  \lean{forces_triangle}, \lean{conjecture_case_three}, \lean{conjecture_implies_seven},
  \lean{counterexample_refutes_conjecture} (small corollaries about the conjecture, not used by the facts
  of this section; \lean{forces_triangle} is named above only in the axioms evidence); and in
  \dirn{ErdosSzekeres/Suk/OrientationClique.lean}, \lean{convexPosition_insert_of_positive_triples} and
  \lean{convexPosition_insert_of_negative_triples} (outside the foundations list). Of the research files
  in the table, only the declarations named in the statements, remarks or proofs of this section belong to
  its scope; the theorems named in the axioms-evidence paragraph only as evidence, and the other
  declarations of these files, are treated in other sections or lie outside Stage~1. In particular,
  \lean{toC}, \lean{ofC}, \lean{crot} and their lemmas in \dirn{LAYERLEAN/LayerLemma.lean} and its copies
  are implementation devices for complex coordinates, replaced here by the direct complex form in
  Lemma~\ref{fd:lem:rot}(a).
\end{itemize}
\par}

\section{The layer lemma and the reduction to exactly \texorpdfstring{$mt$}{mt} points}\label{ll:sec}

This section gives written proofs of the declarations of the Lean file \dirn{LAYERLEAN/LayerLemma.lean} (namespace
\lean{ErdosSzekeres.LayerLemma}) that concern hull vertices, orbits of a map of finite period, the layer lemma, the
peeling of convex layers, and the reduction of the sixfold and sevenfold octagon bounds to statements about sets of
exactly $mt$ points. The main results are:
\begin{itemize}
\item the layer lemma in abstract form (Theorem~\ref{ll:thm:core}): if an injective affine map $f$ sends a finite set
  $S$ into itself, every point of $S$ has minimal period $m>0$ under $f$, and no $K$ points of $S$ are in convex
  position for some $K\le 2m$, then for every hull vertex $x$ of $S$ the hull vertices of $S$ are exactly the points
  $x,f(x),\dots,f^{m-1}(x)$, so that, if $S$ is nonempty, they form a single $f$-orbit; no general position is
  assumed;
\item its rotation form (Theorem~\ref{ll:thm:core-rot}), the description of all convex layers
  (Theorem~\ref{ll:thm:layers}, Lemma~\ref{ll:lem:exists-layer}) and of their sizes (Theorem~\ref{ll:thm:card-peel});
\item the reduction to exactly $mt$ points (Theorem~\ref{ll:thm:reduction}) and the sixfold and sevenfold octagon
  theorems (Theorems~\ref{ll:thm:sixfold} and~\ref{ll:thm:sevenfold}). The last two are \emph{conditional}: they take
  the propositions \lean{Sixfold66} and \lean{Sevenfold70} (Definition~\ref{ll:def:66-70}) as hypotheses, and these
  propositions are not proved in Lean (Remark~\ref{ll:rem:hyp}).
\end{itemize}

\paragraph{Sources.} The Lean source is \dirn{LAYERLEAN/LayerLemma.lean} (1148 lines, SHA-256
\texttt{afb8542dceb0}\dots). This section treats its lines 1--469 and 1034--1148; the geometric lemmas of lines
471--1032 (strict nesting, the centre for even $m$, chains, polar rules, radius order) belong to other sections (for
example Theorem~\ref{lg:thm:nesting} and Lemma~\ref{lg:lem:centre}). The notions of point, general position, convex
position and \lean{HasConvexNGon} come from \dirn{ErdosSzekeres/Definitions.lean} (SHA-256 \texttt{2425500eefd1}\dots)
and are recalled in Definitions~\ref{fd:def:point}, \ref{fd:def:gp}, \ref{fd:def:cp} and~\ref{fd:def:ngon}.

\paragraph{Lean evidence.} The file was checked with Lean 4.24.0 and Mathlib \texttt{f897ebcf} (the repository's pins)
by \texttt{lake env lean}, run from the repository root. It is not a module of the default build target, so a plain
\texttt{lake build} does not re-check it. Two records show that the current file compiles. The review report
\dirn{LEANEXT/review/REPORT.md} of 4 November 2026 (item~1) records exit code~0 and no messages for the reviewed file,
which it identifies by its SHA-256 \texttt{afb8542d}\dots{} and its 1148 lines. The re-check
\dirn{LAYERLEAN/logs/axioms_recheck_20261106/out.txt} of 6 November 2026 records exit status~0 for a file whose first
1148 lines are byte-identical to \dirn{LAYERLEAN/LayerLemma.lean}, followed by 19 \texttt{\#print axioms} commands;
its only Lean output is the 19 printed axiom lines. The log \dirn{LAYERLEAN/logs/build.log} (exit code~0, no
messages) names no file hash; it was committed with the first version of the file (454 lines, SHA-256
\texttt{804b1d25f79a}\dots, commit \texttt{6429e9a70}) and was not updated by the five later commits that changed the
file, so we do not use it. Axioms: \dirn{LAYERLEAN/logs/axioms.log} (28 theorems) and the re-check output just named
(19 theorems) report exactly the axioms \lean{propext}, \lean{Classical.choice} and \lean{Quot.sound} for every
printed theorem. The review log
\dirn{LEANEXT/review/axioms.log} prints 59 theorems, namely all theorems of the file that are not \texttt{@[simp]}
lemmas, among them every theorem of this section except the seven \texttt{@[simp]} lemmas of
Lemma~\ref{ll:lem:complex}(a)--(b); it reports the same three axioms for each printed theorem, except that
\lean{iterate_mem} depends only on \lean{propext} and \lean{Quot.sound}. This log comes from the review whose report is
named above; the check file from which the axioms were printed is not kept in the record, so the identification of
the printed theorems with those of this file rests on the report. The seven \texttt{@[simp]} lemmas are printed in
none of the three logs; their proofs (lines 236--255) are
\texttt{rfl}, extensionality (\lean{Complex.ext}, \texttt{funext}) and, for \lean{toC_crot}, the lemma
\lean{toC_ofC}. The file contains no \texttt{sorry}, \texttt{admit}, \texttt{axiom} or \texttt{native\_decide}, and its
only \texttt{decide} (line 1125) proves that $6$ is even and is evaluated by the kernel. So no statement of this
section depends on \lean{sorryAx}, on \lean{Lean.ofReduceBool} (\texttt{native\_decide}) or on an axiom of the
project. Every lemma and theorem below whose head names a Lean declaration is \emph{Lean-checked} in this sense
(Theorems~\ref{ll:thm:sixfold} and~\ref{ll:thm:sevenfold} as implications from their hypotheses); for
Lemma~\ref{ll:lem:rot}(b) the Lean declaration is the definition \lean{rotAff}, whose affine-map property is a proof
obligation discharged inside the definition and checked with it. The definitions whose heads name a Lean declaration restate Lean
definitions, which the same build type-checks; being definitions, they are not Lean-checked claims. In particular,
Definition~\ref{ll:def:66-70} only introduces the propositions \lean{Sixfold66} and \lean{Sevenfold70}, which are not
proved (Remark~\ref{ll:rem:hyp}). The proofs below are \emph{written proofs (not refereed)}; for each statement that
corresponds to a Lean declaration we say whether its proof follows the Lean proof.

\paragraph{Dependencies.} From the foundations section we use Definitions~\ref{fd:def:point}, \ref{fd:def:orient},
\ref{fd:def:gp}, \ref{fd:def:cp}, \ref{fd:def:ngon}, \ref{fd:def:inv}, \ref{fd:def:rot}, Lemmas~\ref{fd:lem:cp-hull},
\ref{fd:lem:mono}, \ref{fd:lem:collinear} and~\ref{fd:lem:affine-cp}, and the lemma of that section titled ``Convex
hulls'' (in (C3) below). Theorem~\ref{ll:thm:sixfold} uses Lemma~\ref{lg:lem:centre} (exclusion of the centre for
even $m$); for the case needed we also include the short argument. Nothing else from other sections is used.

\subsection{Conventions}

Points are elements of $\R^2$ (Lean \lean{Point}, defined as \texttt{Fin 2} $\to\R$), written $p=(p_0,p_1)$. We use
the following conventions and standard facts.
\begin{enumerate}[label=(C\arabic*)]
\item \emph{Finite sets.} A finite set of points is a Lean \texttt{Finset Point}; $\card S$ is its cardinality;
  $S\setminus\{c\}$ is Lean \texttt{S.erase c}; ``$S$ is nonempty'' is Lean \lean{S.Nonempty}.
\item \emph{Iterates and minimal period.} For a set $\alpha$ and $f:\alpha\to\alpha$ put $f^0=\mathrm{id}$ and
  $f^{i+1}=f\circ f^i$. Lean's \texttt{f\textasciicircum[i]} (\lean{Nat.iterate}) is defined by the other recursion,
  $f^{i+1}=f^i\circ f$; an easy induction shows that both recursions define the same maps (Mathlib's
  \lean{Function.iterate_succ'}), and $f^{i+j}=f^i\circ f^j$ (Mathlib's \lean{Function.iterate_add}). The
  \emph{minimal period}
  $\operatorname{per}_f(x)$ of $x\in\alpha$ is Mathlib's \texttt{Function.minimalPeriod f x}: the least $n\ge1$ with
  $f^n(x)=x$ if such an $n$ exists, and $0$ otherwise. Directly from this definition:
  \begin{itemize}
  \item[(P1)] $f^{\operatorname{per}_f(x)}(x)=x$ for every $x$ (if no period exists, $f^0(x)=x$);
  \item[(P2)] if $n\ge1$ and $f^n(x)=x$, then $1\le\operatorname{per}_f(x)\le n$.
  \end{itemize}
\item \emph{Convex hulls.} For $A\subseteq\R^2$, $\conv(A)$ is the smallest convex set containing $A$ (Mathlib
  \lean{convexHull}). It is the set of all convex combinations $\sum_{i=1}^r\lambda_ia_i$ with $r\ge1$, $a_i\in A$,
  $\lambda_i\ge0$ and $\sum_i\lambda_i=1$. This is part~(b) of the lemma titled ``Convex hulls'' in the foundations
  section, where it is proved for the whole document (see also R.~T.~Rockafellar, \emph{Convex Analysis}, Princeton
  1970, Theorem~2.3). Hence: $A\subseteq\conv(A)$; $A\subseteq B$ implies $\conv(A)\subseteq\conv(B)$; if
  $A\subseteq C$ with $C$ convex
  then $\conv(A)\subseteq C$, so $\conv(\conv(A))=\conv(A)$; $\conv(\emptyset)=\emptyset$; a convex combination of
  points of a convex set $C$ lies in $C$. Merging equal points and deleting zero coefficients, every point of
  $\conv(A)$ is a convex combination of pairwise distinct points of $A$ with positive coefficients and at least one
  term.
\item \emph{Extreme points.} A point $x$ is an \emph{extreme point} of $A\subseteq\R^2$ (Mathlib
  \lean{Set.extremePoints}) if $x\in A$ and, whenever $x=ax_1+bx_2$ with $x_1,x_2\in A$, $a,b>0$ and $a+b=1$,
  we have $x_1=x$. This is Mathlib's definition at the pinned revision (with $x$ in the open segment between $x_1$ and
  $x_2$); exchanging the roles of $x_1$ and $x_2$ shows that then also $x_2=x$.
\item \emph{Affine maps.} A Lean affine map $f$, an element of the type \texttt{Point}~$\to^{\mathrm a}[\R]$~\texttt{Point},
  consists of a map $f:\R^2\to\R^2$ and a linear map $L$ with $f(v+p)=Lv+f(p)$ for all $v,p$; taking $p=0$,
$f(v)=Lv+b$ with $b=f(0)$, and conversely every map
  $p\mapsto Lp+b$ with $L$ linear is of this form. ``$f$ is injective'' refers to the underlying map on all of $\R^2$.
  If $\sum_i\lambda_i=1$ then $f(\sum_i\lambda_ia_i)=\sum_i\lambda_i(La_i+b)=\sum_i\lambda_if(a_i)$; hence, by (C3),
  $\conv(f(A))\subseteq f(\conv(A))$ for every $A\subseteq\R^2$ (this inclusion is part of
  Lemma~\ref{fd:lem:affine-cp}).
\item \emph{Convex position.} Convex position, general position and \lean{HasConvexNGon} are as in
  Definitions~\ref{fd:def:cp}, \ref{fd:def:gp} and~\ref{fd:def:ngon}: \lean{HasConvexNGon} $K$ $S$ means that some
  $T\subseteq S$ with $\card T=K$ is in convex position, and ``$S$ is $K$-free'' means $\neg\,$\lean{HasConvexNGon}
  $K$ $S$. We use Lemma~\ref{fd:lem:cp-hull} ($S$ is in convex position iff $p\notin\conv(S\setminus\{p\})$ for every
  $p\in S$) and Lemma~\ref{fd:lem:mono} (monotonicity). In particular the empty set is in convex position, so
  \lean{HasConvexNGon} $0$ $S$ holds for every $S$.
\item \emph{Angles.} In Lean the angle $2\pi/m$ is \texttt{2 * Real.pi / m} with $m\in\N$ cast to $\R$; division by
  $0$ gives $0$. For $t\in\N$ the subtraction $t\mathbin{\dot{-}}1$ is truncated ($0\mathbin{\dot{-}}1=0$).
\end{enumerate}

\subsection{Orbit segments}

\begin{definition}[orbit segment; \leanref{ErdosSzekeres.LayerLemma.orbitFS}{LAYERLEAN/LayerLemma.lean:48}]
\label{ll:def:orbit}
Let $\alpha$ be a set (Lean: a type with decidable equality), $f:\alpha\to\alpha$, $p\in\N$ and $x\in\alpha$. The
\emph{orbit segment} is the finite set
\[\operatorname{orb}_p(f,x)=\{f^i(x):\ i\in\N,\ 0\le i<p\}\]
(Lean \texttt{orbitFS f p x}, the image of $\{0,\dots,p-1\}$ under $i\mapsto f^i(x)$). It is empty for $p=0$.
\end{definition}

\begin{lemma}[orbit segments; \leanref{ErdosSzekeres.LayerLemma.iterate_mem}{LAYERLEAN/LayerLemma.lean:52},
\leanref{ErdosSzekeres.LayerLemma.mem_orbitFS_self}{LAYERLEAN/LayerLemma.lean:59},
\leanref{ErdosSzekeres.LayerLemma.orbitFS_subset}{LAYERLEAN/LayerLemma.lean:62},
\leanref{ErdosSzekeres.LayerLemma.orbitFS_invariant}{LAYERLEAN/LayerLemma.lean:68},
\leanref{ErdosSzekeres.LayerLemma.card_orbitFS_of_minimalPeriod}{LAYERLEAN/LayerLemma.lean:80},
\leanref{ErdosSzekeres.LayerLemma.mem_orbitFS_of_not_disjoint}{LAYERLEAN/LayerLemma.lean:91}]\label{ll:lem:orbit}
Let $\alpha$ be a set and $f:\alpha\to\alpha$.
\begin{enumerate}[label=(\alph*)]
\item (\lean{iterate_mem}) If $S\subseteq\alpha$ is finite, $f(y)\in S$ for every $y\in S$, and $x\in S$, then
  $f^i(x)\in S$ for every $i\in\N$.
\item (\lean{mem_orbitFS_self}) If $0<p$, then $x\in\operatorname{orb}_p(f,x)$ for every $x\in\alpha$.
\item (\lean{orbitFS_subset}) If $S\subseteq\alpha$ is finite, $f(y)\in S$ for every $y\in S$, and $x\in S$, then
  $\operatorname{orb}_p(f,x)\subseteq S$ for every $p\in\N$.
\item (\lean{orbitFS_invariant}) If $0<p$ and $f^p(x)=x$, then $f(y)\in\operatorname{orb}_p(f,x)$ for every
  $y\in\operatorname{orb}_p(f,x)$.
\item (\lean{card_orbitFS_of_minimalPeriod}) If $\operatorname{per}_f(x)=p$, then
  $\card{\operatorname{orb}_p(f,x)}=p$.
\item (\lean{mem_orbitFS_of_not_disjoint}) If $0<p$, $f^p(x)=x$, $f^p(y)=y$ and the sets $\operatorname{orb}_p(f,x)$
  and $\operatorname{orb}_p(f,y)$ are not disjoint, then $y\in\operatorname{orb}_p(f,x)$.
\end{enumerate}
\end{lemma}

\begin{proof}
We follow the Lean proofs.
(a) Induction on $i$: $f^0(x)=x\in S$, and if $f^i(x)\in S$ then $f^{i+1}(x)=f(f^i(x))\in S$.
(b) Take $i=0<p$: $f^0(x)=x$.
(c) Every element of $\operatorname{orb}_p(f,x)$ is some $f^i(x)$, which lies in $S$ by (a).
(d) Let $y=f^i(x)$ with $0\le i<p$; then $f(y)=f^{i+1}(x)$. If $i+1<p$ this lies in $\operatorname{orb}_p(f,x)$.
Otherwise $i+1=p$ and $f(y)=f^p(x)=x\in\operatorname{orb}_p(f,x)$ by (b).
(e) If $p=0$, then $\operatorname{orb}_0(f,x)=\emptyset$ has $0$ elements. Let $p\ge1$. It suffices to show that
$f^0(x),\dots,f^{p-1}(x)$ are pairwise distinct. By (P1), $f^p(x)=x$. Suppose $f^i(x)=f^j(x)$ with $0\le i<j\le p-1$.
Applying $f^{p-j}$ gives
\[f^{p-j+i}(x)=f^{p-j}(f^i(x))=f^{p-j}(f^j(x))=f^p(x)=x .\]
Since $j\le p-1$ and $i\le j-1$, we have $1\le p-j+i\le p-1$, so (P2) gives $\operatorname{per}_f(x)\le p-j+i<p$, a
contradiction.
(f) Let $z$ lie in both sets, say $z=f^j(y)$ with $0\le j<p$. Then $f^{p-j}(z)=f^{p}(y)=y$. By (d), $f$ maps the finite
set $\operatorname{orb}_p(f,x)$ into itself, and $z\in\operatorname{orb}_p(f,x)$; so by (a), applied to this set,
$y=f^{p-j}(z)\in\operatorname{orb}_p(f,x)$.
\end{proof}

\begin{definition}[invariance; \leanref{ErdosSzekeres.LayerLemma.Invariant}{LAYERLEAN/LayerLemma.lean:105}]
\label{ll:def:inv}
For a map $f:\R^2\to\R^2$ and a finite set $S\subseteq\R^2$, $S$ is \emph{$f$-invariant} (Lean
\texttt{Invariant f S}) if $f(p)\in S$ for every $p\in S$, that is, $f$ maps $S$ \emph{into} itself. This is the notion of
Definition~\ref{fd:def:inv}. (For injective $f$ it implies $f(S)=S$; see the proof of Lemma~\ref{ll:lem:map-hv}.)
\end{definition}

\subsection{Hull vertices}

\begin{definition}[hull vertices; \leanref{ErdosSzekeres.LayerLemma.hullVerts}{LAYERLEAN/LayerLemma.lean:111}]
\label{ll:def:hv}
For a finite set $S\subseteq\R^2$ let
\[\hv(S)=\{x\in S:\ x\notin\conv(S\setminus\{x\})\}\]
(Lean \texttt{hullVerts S}, the subset of $S$ cut out by this condition). Its elements are the \emph{hull vertices} of
$S$.
\end{definition}

\begin{lemma}[membership; \leanref{ErdosSzekeres.LayerLemma.mem_hullVerts}{LAYERLEAN/LayerLemma.lean:114},
\leanref{ErdosSzekeres.LayerLemma.hullVerts_subset}{LAYERLEAN/LayerLemma.lean:118}]\label{ll:lem:hv-basic}
Let $S$ be a finite set of points. (a) For every point $x$: $x\in\hv(S)$ if and only if $x\in S$ and
$x\notin\conv(S\setminus\{x\})$. (b) $\hv(S)\subseteq S$.
\end{lemma}

\begin{proof}
As in Lean: (a) is Definition~\ref{ll:def:hv}, and (b) follows from (a).
\end{proof}

\begin{lemma}[hull vertices; \leanref{ErdosSzekeres.LayerLemma.convexPosition_hullVerts}{LAYERLEAN/LayerLemma.lean:121},
\leanref{ErdosSzekeres.LayerLemma.extreme_mem_hullVerts}{LAYERLEAN/LayerLemma.lean:130},
\leanref{ErdosSzekeres.LayerLemma.subset_convexHull_hullVerts}{LAYERLEAN/LayerLemma.lean:146},
\leanref{ErdosSzekeres.LayerLemma.hullVerts_nonempty}{LAYERLEAN/LayerLemma.lean:139}]\label{ll:lem:hv}
Let $S$ be a finite set of points.
\begin{enumerate}[label=(\alph*)]
\item (\lean{convexPosition_hullVerts}) $\hv(S)$ is in convex position.
\item (\lean{extreme_mem_hullVerts}) Every extreme point of $\conv(S)$ belongs to $\hv(S)$.
\item (\lean{subset_convexHull_hullVerts}) $S\subseteq\conv(\hv(S))$.
\item (\lean{hullVerts_nonempty}) If $S$ is nonempty, then $\hv(S)$ is nonempty.
\end{enumerate}
\end{lemma}

\begin{proof}
The proof of (a) follows Lean. For (b)--(d) we give elementary proofs, different from the Lean proofs. Lean proves (b)
with two Mathlib facts: the extreme points of $\conv(S)$ lie in $S$, and ``$x$ is an extreme point of a convex set $C$
iff $x\in C$ and $C\setminus\{x\}$ is convex''. It deduces (d) from (b) and the existence of an extreme point of the
nonempty compact set $\conv(S)$ (Mathlib's Krein--Milman lemma), and (c) from (b) and the Krein--Milman theorem
($\conv(S)$ is the closure of the convex hull of its extreme points).

(a) By Lemma~\ref{fd:lem:cp-hull} we must show $x\notin\conv(\hv(S)\setminus\{x\})$ for every $x\in\hv(S)$. Since
$\hv(S)\setminus\{x\}\subseteq S\setminus\{x\}$, (C3) gives $\conv(\hv(S)\setminus\{x\})\subseteq\conv(S\setminus\{x\})$,
and $x\notin\conv(S\setminus\{x\})$ because $x\in\hv(S)$.

(b) Let $x$ be an extreme point of $C=\conv(S)$. We first show:
\begin{itemize}
\item[$(*)$] if $x=\sum_{i=1}^r\lambda_is_i$ with $r\ge1$, $s_i\in C$, $\lambda_i>0$ and $\sum_i\lambda_i=1$, then
  $s_i=x$ for every $i$.
\end{itemize}
If $r=1$ then $\lambda_1=1$ and $s_1=x$. Let $r\ge2$ and fix $j$. Then $0<\lambda_j<1$, since some other coefficient is
positive. The point $y_j=\sum_{i\ne j}\frac{\lambda_i}{1-\lambda_j}s_i$ is a convex combination of points of the convex
set $C$, so $y_j\in C$ by (C3), and $x=\lambda_js_j+(1-\lambda_j)y_j$. By (C4), applied with $x_1=s_j$, $x_2=y_j$,
$a=\lambda_j$, $b=1-\lambda_j$, we get $s_j=x$. This proves $(*)$.

Since $x\in C$, by (C3) $x=\sum_{i=1}^r\lambda_is_i$ with $r\ge1$, $s_i\in S\subseteq C$, $\lambda_i>0$,
$\sum_i\lambda_i=1$; by $(*)$, $x=s_1\in S$. Suppose $x\in\conv(S\setminus\{x\})$. Then $x=\sum_{i=1}^r\mu_it_i$ with
$r\ge1$, $t_i\in S\setminus\{x\}\subseteq C$, $\mu_i>0$, $\sum_i\mu_i=1$, and $(*)$ gives $t_1=x$, contradicting
$t_1\ne x$. Hence $x\notin\conv(S\setminus\{x\})$, and $x\in\hv(S)$.

(c) Strong induction on $\card S$. If $\hv(S)=S$, then $S\subseteq\conv(S)=\conv(\hv(S))$; this covers $S=\emptyset$.
Otherwise choose $y\in S\setminus\hv(S)$; then $y\in\conv(S')$ where $S'=S\setminus\{y\}$. As $\card{S'}<\card S$, the
induction hypothesis gives $S'\subseteq\conv(\hv(S'))$.

We claim $\hv(S')\subseteq\hv(S)$. Let $z\in\hv(S')$, so $z\in S'$ (in particular $z\in S$ and $z\ne y$) and
$z\notin\conv(A)$, where $A=S'\setminus\{z\}=S\setminus\{y,z\}$. Suppose $z\in\conv(S\setminus\{z\})$. Since
$S\setminus\{z\}=A\cup\{y\}$ and $S'=A\cup\{z\}$, (C3) and merging of equal points give
\[z=\sum_{a\in A}\alpha_aa+\lambda y,\qquad y=\sum_{a\in A}\beta_aa+\mu z,\]
with all $\alpha_a,\beta_a,\lambda,\mu\ge0$, $\sum_a\alpha_a+\lambda=1$ and $\sum_a\beta_a+\mu=1$ (the second
representation exists because $y\in\conv(S')$). Substituting the second into the first,
\[(1-\lambda\mu)\,z=\sum_{a\in A}(\alpha_a+\lambda\beta_a)\,a .\]
If $\lambda\mu<1$, the coefficients $(\alpha_a+\lambda\beta_a)/(1-\lambda\mu)$ are nonnegative with sum
$\bigl((1-\lambda)+\lambda(1-\mu)\bigr)/(1-\lambda\mu)=1$, so $z\in\conv(A)$, a contradiction. If $\lambda\mu=1$, then
$\lambda=\mu=1$ because $\lambda,\mu\in[0,1]$; so all $\alpha_a=0$ and $z=y$, contradicting $z\ne y$. This proves the
claim.

By the claim and (C3), $S'\subseteq\conv(\hv(S'))\subseteq\conv(\hv(S))$, and then
$y\in\conv(S')\subseteq\conv(\conv(\hv(S)))=\conv(\hv(S))$. Hence $S=S'\cup\{y\}\subseteq\conv(\hv(S))$.

(d) If $\hv(S)=\emptyset$, then by (c) $S\subseteq\conv(\emptyset)=\emptyset$, so $S$ is empty.
\end{proof}

\begin{lemma}[affine invariance of hull vertices;
\leanref{ErdosSzekeres.LayerLemma.map_mem_hullVerts}{LAYERLEAN/LayerLemma.lean:162}]\label{ll:lem:map-hv}
Let $f$ be an affine map of $\R^2$ (Lean: $f$ of type \texttt{Point}~$\to^{\mathrm a}[\R]$~\texttt{Point}) that is
injective, and let $S$ be a finite set
with $f(y)\in S$ for every $y\in S$. If $x\in\hv(S)$, then $f(x)\in\hv(S)$.
\end{lemma}

\begin{proof}
We follow the Lean proof. (1) $f(S)=S$: as $f$ is injective, $\card{f(S)}=\card S$, and $f(S)\subseteq S$ with $S$
finite. (2) $f(x)\in S$, since $x\in S$. (3) Suppose $f(x)\in\conv(S\setminus\{f(x)\})$. Every $y\in S$ with
$y\ne f(x)$ is, by (1), of the form $y=f(z)$ with $z\in S$, and $z\ne x$ since $y\ne f(x)$; thus
$S\setminus\{f(x)\}\subseteq f(S\setminus\{x\})$. By (C3) and (C5),
\[f(x)\in\conv\bigl(f(S\setminus\{x\})\bigr)\subseteq f\bigl(\conv(S\setminus\{x\})\bigr),\]
so $f(x)=f(u)$ for some $u\in\conv(S\setminus\{x\})$. Injectivity of $f$ on all of $\R^2$ (note that $u$ need not lie
in $S$) gives $x=u\in\conv(S\setminus\{x\})$, contradicting $x\in\hv(S)$. By Lemma~\ref{ll:lem:hv-basic}(a), (2) and
(3) give $f(x)\in\hv(S)$.
\end{proof}

\subsection{The core lemma}

\begin{theorem}[layer lemma, abstract form;
\leanref{ErdosSzekeres.LayerLemma.hullVerts_eq_orbitFS}{LAYERLEAN/LayerLemma.lean:189}]\label{ll:thm:core}
Let $f$ be an affine map of $\R^2$ (Lean: $f$ of type \texttt{Point}~$\to^{\mathrm a}[\R]$~\texttt{Point}) that is
injective, let $S$ be a finite set of
points, and let $m,K\in\N$ with $0<m$ and $K\le 2m$. Assume:
\begin{enumerate}[label=(\roman*)]
\item $f(y)\in S$ for every $y\in S$;
\item $\operatorname{per}_f(y)=m$ for every $y\in S$;
\item $S$ is $K$-free, i.e.\ $\neg\,$\lean{HasConvexNGon} $K$ $S$.
\end{enumerate}
Then for every $x\in\hv(S)$,
\[\hv(S)=\operatorname{orb}_m(f,x)=\{f^i(x):\ 0\le i<m\}.\]
\end{theorem}

\begin{proof}
We follow the Lean proof. Write $H=\hv(S)$, a finite set with $H\subseteq S$ (Lemma~\ref{ll:lem:hv-basic}).

Step 1. $f(y)\in H$ for every $y\in H$, by Lemma~\ref{ll:lem:map-hv} and (i).

Step 2. $f^m(y)=y$ for every $y\in S$, by (ii) and (P1).

Step 3. $\operatorname{orb}_m(f,x)\subseteq H$, by Lemma~\ref{ll:lem:orbit}(c) applied to the finite set $H$
(Step~1) and $x\in H$.

Step 4. $H\subseteq\operatorname{orb}_m(f,x)$. Let $y\in H$ and suppose $y\notin\operatorname{orb}_m(f,x)$. Both $x$ and
$y$ lie in $S$, so $f^m(x)=x$ and $f^m(y)=y$ by Step~2. If $\operatorname{orb}_m(f,x)$ and $\operatorname{orb}_m(f,y)$
were not disjoint, Lemma~\ref{ll:lem:orbit}(f) (with $p=m>0$) would give $y\in\operatorname{orb}_m(f,x)$; so they are
disjoint. Both are contained in $H$, by Step~1 and Lemma~\ref{ll:lem:orbit}(c), and each has exactly $m$ elements by
(ii) and Lemma~\ref{ll:lem:orbit}(e). Hence
\[\card H\ \ge\ \card{\operatorname{orb}_m(f,x)\cup\operatorname{orb}_m(f,y)}=2m\ \ge\ K .\]
Choose $T\subseteq H$ with $\card T=K$. By Lemma~\ref{ll:lem:hv}(a) and Lemma~\ref{fd:lem:mono}(a), $T$ is in convex
position, and $T\subseteq H\subseteq S$. So \lean{HasConvexNGon} $K$ $S$ holds, contradicting (iii). Hence
$y\in\operatorname{orb}_m(f,x)$.

Steps 3 and 4 give $H=\operatorname{orb}_m(f,x)$.
\end{proof}

\begin{lemma}[invariance of $S\setminus\hv(S)$;
\leanref{ErdosSzekeres.LayerLemma.invariant_sdiff_hullVerts}{LAYERLEAN/LayerLemma.lean:217}]\label{ll:lem:inv-sdiff}
Let $f$ be an affine map of $\R^2$ that is injective, $S$ a finite set and $m\in\N$ with $0<m$. Assume $f(y)\in S$ for
every $y\in S$ and $f^m(y)=y$ for every $y\in S$. Then $f(y)\in S\setminus\hv(S)$ for every $y\in S\setminus\hv(S)$.
\end{lemma}

\begin{proof}
We follow the Lean proof. Let $y\in S\setminus\hv(S)$. Then $f(y)\in S$. Suppose $f(y)\in\hv(S)$. By
Lemma~\ref{ll:lem:map-hv}, $f$ maps the finite set $\hv(S)$ into itself, so Lemma~\ref{ll:lem:orbit}(a) gives
$f^{m-1}(f(y))\in\hv(S)$. As $m\ge1$, $f^{m-1}(f(y))=f^m(y)=y$, so $y\in\hv(S)$, a contradiction. Hence
$f(y)\in S\setminus\hv(S)$.
\end{proof}

\begin{remark}[on Theorem~\ref{ll:thm:core}]\label{ll:rem:core}
(1) The proof uses only two geometric facts: an injective affine map that sends $S$ into itself sends $\hv(S)$ into
itself, hence permutes the finite set $\hv(S)$ (Lemma~\ref{ll:lem:map-hv}), and the hull vertices are in convex
position (Lemma~\ref{ll:lem:hv}(a)). No general
position is assumed, and $f$ need not be a rotation.
(2) If $K=0$, hypothesis (iii) is false by (C6), and the theorem is vacuous; if $\hv(S)=\emptyset$, which by
Lemma~\ref{ll:lem:hv}(d) happens only for $S=\emptyset$, there is no $x$.
(3) Under the hypotheses of the theorem, if $S\ne\emptyset$ then $\hv(S)$ contains some $x$ (Lemma~\ref{ll:lem:hv}(d)),
and the conclusion with Lemma~\ref{ll:lem:orbit}(e) and (ii) gives $\card{\hv(S)}=m$.
\end{remark}

\subsection{Rotations}

\begin{definition}[complex coordinates; \leanref{ErdosSzekeres.LayerLemma.toC}{LAYERLEAN/LayerLemma.lean:231},
\leanref{ErdosSzekeres.LayerLemma.ofC}{LAYERLEAN/LayerLemma.lean:234},
\leanref{ErdosSzekeres.LayerLemma.crot}{LAYERLEAN/LayerLemma.lean:252}]\label{ll:def:complex}
For $p\in\R^2$ let $\operatorname{toC}(p)=p_0+ip_1\in\C$; for $z\in\C$ let $\operatorname{ofC}(z)=(\operatorname{Re}z,
\operatorname{Im}z)\in\R^2$; for $\gamma,w\in\C$ and $p\in\R^2$ let
\[\operatorname{crot}_{\gamma,w}(p)=\operatorname{ofC}\bigl(\gamma+w\,(\operatorname{toC}(p)-\gamma)\bigr)\]
(Lean \texttt{crot c w p}, where $c\in\C$ plays the role of $\gamma$).
\end{definition}

\begin{lemma}[complex coordinates; \leanref{ErdosSzekeres.LayerLemma.toC_re}{LAYERLEAN/LayerLemma.lean:236},
\leanref{ErdosSzekeres.LayerLemma.toC_im}{LAYERLEAN/LayerLemma.lean:237},
\leanref{ErdosSzekeres.LayerLemma.ofC_zero}{LAYERLEAN/LayerLemma.lean:238},
\leanref{ErdosSzekeres.LayerLemma.ofC_one}{LAYERLEAN/LayerLemma.lean:239},
\leanref{ErdosSzekeres.LayerLemma.toC_ofC}{LAYERLEAN/LayerLemma.lean:241},
\leanref{ErdosSzekeres.LayerLemma.ofC_toC}{LAYERLEAN/LayerLemma.lean:244},
\leanref{ErdosSzekeres.LayerLemma.toC_injective}{LAYERLEAN/LayerLemma.lean:248},
\leanref{ErdosSzekeres.LayerLemma.toC_crot}{LAYERLEAN/LayerLemma.lean:254},
\leanref{ErdosSzekeres.LayerLemma.crot_one}{LAYERLEAN/LayerLemma.lean:257},
\leanref{ErdosSzekeres.LayerLemma.crot_comp}{LAYERLEAN/LayerLemma.lean:262},
\leanref{ErdosSzekeres.LayerLemma.crot_iterate}{LAYERLEAN/LayerLemma.lean:267},
\leanref{ErdosSzekeres.LayerLemma.crot_eq_self_iff}{LAYERLEAN/LayerLemma.lean:273},
\leanref{ErdosSzekeres.LayerLemma.crot_injective}{LAYERLEAN/LayerLemma.lean:287}]\label{ll:lem:complex}
For all $p\in\R^2$, $z,\gamma,w,v\in\C$ and $k\in\N$:
\begin{enumerate}[label=(\alph*)]
\item (\lean{toC_re}, \lean{toC_im}, \lean{ofC_zero}, \lean{ofC_one}, \lean{toC_ofC}, \lean{ofC_toC},
  \lean{toC_injective}) $\operatorname{Re}\operatorname{toC}(p)=p_0$, $\operatorname{Im}\operatorname{toC}(p)=p_1$,
  $\operatorname{ofC}(z)_0=\operatorname{Re}z$, $\operatorname{ofC}(z)_1=\operatorname{Im}z$,
  $\operatorname{toC}(\operatorname{ofC}(z))=z$ and $\operatorname{ofC}(\operatorname{toC}(p))=p$; in particular
  $\operatorname{toC}$ is injective;
\item (\lean{toC_crot}) $\operatorname{toC}(\operatorname{crot}_{\gamma,w}(p))=\gamma+w(\operatorname{toC}(p)-\gamma)$;
\item (\lean{crot_one}) $\operatorname{crot}_{\gamma,1}=\mathrm{id}$;
\item (\lean{crot_comp}) $\operatorname{crot}_{\gamma,w}(\operatorname{crot}_{\gamma,v}(p))
  =\operatorname{crot}_{\gamma,wv}(p)$;
\item (\lean{crot_iterate}) $(\operatorname{crot}_{\gamma,w})^k=\operatorname{crot}_{\gamma,w^k}$;
\item (\lean{crot_eq_self_iff}) if $w\ne1$, then $\operatorname{crot}_{\gamma,w}(p)=p$ if and only if
  $\operatorname{toC}(p)=\gamma$;
\item (\lean{crot_injective}) if $w\ne0$, then $\operatorname{crot}_{\gamma,w}$ is injective.
\end{enumerate}
\end{lemma}

\begin{proof}
We follow the Lean proofs. (a) The four coordinate identities are the definitions; they give
$\operatorname{toC}(\operatorname{ofC}(z))=\operatorname{Re}z+i\operatorname{Im}z=z$ and
$\operatorname{ofC}(\operatorname{toC}(p))=(p_0,p_1)=p$, and a map with a left inverse is injective.
(b) follows from $\operatorname{toC}\circ\operatorname{ofC}=\mathrm{id}$.
(c) By (b), $\operatorname{toC}(\operatorname{crot}_{\gamma,1}(p))=\gamma+(\operatorname{toC}(p)-\gamma)
=\operatorname{toC}(p)$; apply injectivity of $\operatorname{toC}$.
(d) By (b) twice, $\operatorname{toC}$ of the left side is $\gamma+w\bigl((\gamma+v(\operatorname{toC}(p)-\gamma))
-\gamma\bigr)=\gamma+wv(\operatorname{toC}(p)-\gamma)$, which is $\operatorname{toC}$ of the right side; apply
injectivity.
(e) Induction on $k$: $(\operatorname{crot}_{\gamma,w})^0=\mathrm{id}=\operatorname{crot}_{\gamma,1}$ by (c), and
$(\operatorname{crot}_{\gamma,w})^{k+1}(p)=\operatorname{crot}_{\gamma,w}(\operatorname{crot}_{\gamma,w^k}(p))
=\operatorname{crot}_{\gamma,w\,w^k}(p)$ by (d).
(f) Write $u=\operatorname{toC}(p)$. By (a) and (b), $\operatorname{crot}_{\gamma,w}(p)=p$ iff $\gamma+w(u-\gamma)=u$,
i.e.\ iff $(w-1)(u-\gamma)=0$, i.e.\ (as $w\ne1$ and $\C$ is a field) iff $u=\gamma$.
(g) If $\operatorname{crot}_{\gamma,w}(p)=\operatorname{crot}_{\gamma,w}(q)$, then by (b)
$w(\operatorname{toC}(p)-\operatorname{toC}(q))=0$, so $\operatorname{toC}(p)=\operatorname{toC}(q)$ as $w\ne0$, and
$p=q$ by (a).
\end{proof}

\begin{definition}[rotation about a point; \leanref{ErdosSzekeres.LayerLemma.rotAbout}{LAYERLEAN/LayerLemma.lean:298}]
\label{ll:def:rot}
For $c\in\R^2$ and $\theta\in\R$ let $\rho_{c,\theta}:\R^2\to\R^2$ (Lean \texttt{rotAbout c}~$\theta$) be
\[\rho_{c,\theta}(p)=\bigl(c_0+\cos\theta\,(p_0-c_0)-\sin\theta\,(p_1-c_1),\ \
c_1+\sin\theta\,(p_0-c_0)+\cos\theta\,(p_1-c_1)\bigr).\]
This is the map of Definition~\ref{fd:def:rot}, the counterclockwise rotation by $\theta$ about $c$. For $m\in\N$ the
map $\rho_{c,2\pi/m}$ is Lean's \texttt{rotAbout c (2 * Real.pi / m)}, with the convention (C7) on the angle.
\end{definition}

\begin{lemma}[rotations; \leanref{ErdosSzekeres.LayerLemma.rotAbout_eq_crot}{LAYERLEAN/LayerLemma.lean:302},
\leanref{ErdosSzekeres.LayerLemma.rotAff}{LAYERLEAN/LayerLemma.lean:311},
\leanref{ErdosSzekeres.LayerLemma.rotAbout_injective}{LAYERLEAN/LayerLemma.lean:319},
\leanref{ErdosSzekeres.LayerLemma.rotAbout_iterate}{LAYERLEAN/LayerLemma.lean:465}]\label{ll:lem:rot}
Let $c\in\R^2$, $\theta,\alpha\in\R$ and $j\in\N$.
\begin{enumerate}[label=(\alph*)]
\item (\lean{rotAbout_eq_crot}) $\rho_{c,\theta}=\operatorname{crot}_{\operatorname{toC}(c),\,e^{i\theta}}$.
\item (\lean{rotAff}) For all $p,v\in\R^2$, $\rho_{c,\theta}(v+p)=A_\theta v+\rho_{c,\theta}(p)$, where $A_\theta$ is the
  linear map with matrix $\begin{pmatrix}\cos\theta&-\sin\theta\\ \sin\theta&\cos\theta\end{pmatrix}$ acting on column
  vectors. So $\rho_{c,\theta}$ is an affine map with linear part $A_\theta$; Lean's \texttt{rotAff c}~$\theta$ is this
  affine map, and its underlying map is $\rho_{c,\theta}$.
\item (\lean{rotAbout_injective}) $\rho_{c,\theta}$ is injective.
\item (\lean{rotAbout_iterate}) $(\rho_{c,\alpha})^j=\rho_{c,j\alpha}$ (with $j$ cast to $\R$).
\end{enumerate}
\end{lemma}

\begin{proof}
We follow the Lean proofs.
(a) Let $u=\operatorname{toC}(p)=p_0+ip_1$, $\gamma=\operatorname{toC}(c)=c_0+ic_1$ and
$e^{i\theta}=\cos\theta+i\sin\theta$. Then
\[\gamma+e^{i\theta}(u-\gamma)=\bigl(c_0+\cos\theta(p_0-c_0)-\sin\theta(p_1-c_1)\bigr)
+i\bigl(c_1+\sin\theta(p_0-c_0)+\cos\theta(p_1-c_1)\bigr),\]
whose real and imaginary parts are the coordinates of $\rho_{c,\theta}(p)$. Applying $\operatorname{ofC}$
(Lemma~\ref{ll:lem:complex}(a)) gives $\operatorname{crot}_{\gamma,e^{i\theta}}(p)=\rho_{c,\theta}(p)$.
(b) The formula of Definition~\ref{ll:def:rot} reads $\rho_{c,\theta}(p)=A_\theta(p-c)+c$. Hence
$\rho_{c,\theta}(v+p)=A_\theta(v+p-c)+c=A_\theta v+\rho_{c,\theta}(p)$ by linearity of $A_\theta$.
(c) By (a), $\rho_{c,\theta}=\operatorname{crot}_{\gamma,e^{i\theta}}$ with $e^{i\theta}\ne0$; apply
Lemma~\ref{ll:lem:complex}(g).
(d) By (a) and Lemma~\ref{ll:lem:complex}(e), $(\rho_{c,\alpha})^j=\operatorname{crot}_{\gamma,(e^{i\alpha})^j}$, and
$(e^{i\alpha})^j=e^{ij\alpha}$; by (a) again this is $\rho_{c,j\alpha}$.
\end{proof}

\begin{lemma}[invariance of $S\setminus\{c\}$;
\leanref{ErdosSzekeres.LayerLemma.invariant_erase_centre}{LAYERLEAN/LayerLemma.lean:325}]\label{ll:lem:erase}
Let $S$ be a finite set, $c$ a point and $\theta\in\R$. If $S$ is $\rho_{c,\theta}$-invariant, then $S\setminus\{c\}$
is $\rho_{c,\theta}$-invariant.
\end{lemma}

\begin{proof}
We follow the Lean proof. Write $\rho=\rho_{c,\theta}$. Putting $p=c$ in Definition~\ref{ll:def:rot} gives $\rho(c)=c$.
Let $p\in S\setminus\{c\}$. Then $\rho(p)\in S$. If $\rho(p)=c=\rho(c)$, then $p=c$ by Lemma~\ref{ll:lem:rot}(c), which
is false. Hence $\rho(p)\in S\setminus\{c\}$.
\end{proof}

\begin{lemma}[exact period of the rotation;
\leanref{ErdosSzekeres.LayerLemma.minimalPeriod_rotAbout}{LAYERLEAN/LayerLemma.lean:336}]\label{ll:lem:minperiod}
Let $c$ be a point, $m\in\N$ with $0<m$, and $y$ a point with $y\ne c$. Then
$\operatorname{per}_{\rho}(y)=m$ for $\rho=\rho_{c,2\pi/m}$.
\end{lemma}

\begin{proof}
We follow the Lean proof. Let $\gamma=\operatorname{toC}(c)$ and $\omega=e^{2\pi i/m}$, where $2\pi/m$ is the real
number of (C7). By Lemma~\ref{ll:lem:rot}(a), $\rho=\operatorname{crot}_{\gamma,\omega}$, and by
Lemma~\ref{ll:lem:complex}(e), $\rho^k=\operatorname{crot}_{\gamma,\omega^k}$ for $k\in\N$. We have
$\omega^k=e^{2\pi ik/m}=\cos(2\pi k/m)+i\sin(2\pi k/m)$. Hence $\omega^m=\cos2\pi+i\sin2\pi=1$, and for $0<k<m$ we have
$0<2\pi k/m<2\pi$, so $\cos(2\pi k/m)<1$ and $\omega^k\ne1$. (That is, $\omega$ is a primitive $m$-th root of unity;
Lean uses Mathlib's \lean{Complex.isPrimitiveRoot_exp}.)

By Lemma~\ref{ll:lem:complex}(e) and~(c),
$\rho^m=\operatorname{crot}_{\gamma,\omega^m}=\operatorname{crot}_{\gamma,1}=\mathrm{id}$, so $\rho^m(y)=y$ with $m\ge1$, and
(P2) gives $1\le q\le m$ for $q=\operatorname{per}_\rho(y)$. Suppose $q<m$. By (P1), $\operatorname{crot}_{\gamma,
\omega^q}(y)=\rho^q(y)=y$, where $\omega^q\ne1$ because $0<q<m$. Lemma~\ref{ll:lem:complex}(f) gives
$\operatorname{toC}(y)=\gamma=\operatorname{toC}(c)$, hence $y=c$ by Lemma~\ref{ll:lem:complex}(a), a contradiction.
So $q=m$.
\end{proof}

\begin{theorem}[layer lemma, rotations;
\leanref{ErdosSzekeres.LayerLemma.hullVerts_eq_orbitFS_rotAbout}{LAYERLEAN/LayerLemma.lean:364}]\label{ll:thm:core-rot}
Let $m,K\in\N$ with $0<m$ and $K\le 2m$, let $S$ be a finite set and $c$ a point with $c\notin S$. Assume that $S$ is
$\rho$-invariant for $\rho=\rho_{c,2\pi/m}$ and that $S$ is $K$-free. Then for every $x\in\hv(S)$,
\[\hv(S)=\operatorname{orb}_m(\rho,x)=\{\rho^i(x):\ 0\le i<m\}.\]
\end{theorem}

\begin{proof}
As in Lean, apply Theorem~\ref{ll:thm:core} to the affine map $f=\rho$ (Lemma~\ref{ll:lem:rot}(b); Lean's
\texttt{rotAff c (2 * Real.pi / m)}, whose underlying map is $\rho$), which is injective by Lemma~\ref{ll:lem:rot}(c).
Its hypotheses hold: $0<m$ and $K\le2m$; (i) is the invariance of $S$; for (ii), every $y\in S$ satisfies $y\ne c$
because $c\notin S$, so $\operatorname{per}_\rho(y)=m$ by Lemma~\ref{ll:lem:minperiod}; (iii) is the $K$-freeness of
$S$.
\end{proof}

\subsection{Peeling the layers}

\begin{definition}[peeling; \leanref{ErdosSzekeres.LayerLemma.peel}{LAYERLEAN/LayerLemma.lean:374}]\label{ll:def:peel}
For a finite set $S$ define finite sets $\peel(S,n)$, $n\in\N$, by
\[\peel(S,0)=S,\qquad \peel(S,n+1)=\peel(S,n)\setminus\hv(\peel(S,n)).\]
So $\peel(S,n)$ is $S$ with its $n$ outermost convex layers removed, and $\hv(\peel(S,n))$ is the $n$-th convex layer
of $S$ (counting from $0$).
\end{definition}

\begin{lemma}[peeling, basic facts; \leanref{ErdosSzekeres.LayerLemma.peel_succ}{LAYERLEAN/LayerLemma.lean:378},
\leanref{ErdosSzekeres.LayerLemma.peel_subset}{LAYERLEAN/LayerLemma.lean:381}]\label{ll:lem:peel-basic}
Let $S$ be a finite set. (a) $\peel(S,n+1)=\peel(S,n)\setminus\hv(\peel(S,n))$ for every $n\in\N$. (b)
$\peel(S,n)\subseteq S$ for every $n\in\N$.
\end{lemma}

\begin{proof}
As in Lean: (a) is the definition. (b) By induction: $\peel(S,0)=S$, and
$\peel(S,n+1)\subseteq\peel(S,n)\subseteq S$.
\end{proof}

\begin{lemma}[peeling is antitone; \leanref{ErdosSzekeres.LayerLemma.peel_antitone}{LAYERLEAN/LayerLemma.lean:460}]
\label{ll:lem:antitone}
Let $S$ be a finite set and $n,n'\in\N$ with $n\le n'$. Then $\peel(S,n')\subseteq\peel(S,n)$.
\end{lemma}

\begin{proof}
Induction on $n'\ge n$, as in Lean: for $n'=n$ there is nothing to prove, and
$\peel(S,n'+1)\subseteq\peel(S,n')\subseteq\peel(S,n)$ by Lemma~\ref{ll:lem:peel-basic}(a).
\end{proof}

\begin{lemma}[the peels are invariant; \leanref{ErdosSzekeres.LayerLemma.invariant_peel}{LAYERLEAN/LayerLemma.lean:386}]
\label{ll:lem:inv-peel}
Let $m\in\N$ with $0<m$, let $S$ be a finite set and $c$ a point with $c\notin S$, and assume that $S$ is
$\rho$-invariant for $\rho=\rho_{c,2\pi/m}$. Then $\peel(S,n)$ is $\rho$-invariant for every $n\in\N$.
\end{lemma}

\begin{proof}
Induction on $n$, as in Lean. For $n=0$ this is the hypothesis. Assume $\peel(S,n)$ is $\rho$-invariant. Every
$y\in\peel(S,n)$ lies in $S$ (Lemma~\ref{ll:lem:peel-basic}(b)), so $y\ne c$, so $\operatorname{per}_\rho(y)=m$ by
Lemma~\ref{ll:lem:minperiod} and $\rho^m(y)=y$ by (P1). Lemma~\ref{ll:lem:inv-sdiff}, applied to the finite set
$\peel(S,n)$ and the injective affine map $\rho$ (Lemma~\ref{ll:lem:rot}(b),(c)), shows that $\rho$ maps
$\peel(S,n)\setminus\hv(\peel(S,n))=\peel(S,n+1)$ into itself.
\end{proof}

\begin{theorem}[every layer is one orbit;
\leanref{ErdosSzekeres.LayerLemma.layers_are_orbits}{LAYERLEAN/LayerLemma.lean:402}]\label{ll:thm:layers}
Let $m,K\in\N$ with $0<m$ and $K\le2m$, let $S$ be a finite set and $c$ a point with $c\notin S$, and write
$\rho=\rho_{c,2\pi/m}$. Assume that $S$ is $\rho$-invariant and $K$-free, and let $n\in\N$ be such that $\peel(S,n)$
is nonempty. Then there is $x\in\peel(S,n)$ such that
\begin{enumerate}[label=(\alph*)]
\item $\hv(\peel(S,n))=\operatorname{orb}_m(\rho,x)$;
\item $\card{\operatorname{orb}_m(\rho,x)}=m$;
\item $\peel(S,n+1)\subseteq\conv\bigl(\hv(\peel(S,n))\bigr)$.
\end{enumerate}
\end{theorem}

\begin{proof}
We follow the Lean proof. By Lemma~\ref{ll:lem:hv}(d) there is $x\in\hv(\peel(S,n))$, and $x\in\peel(S,n)$ by
Lemma~\ref{ll:lem:hv-basic}(b). The finite set $U=\peel(S,n)$ satisfies the hypotheses of
Theorem~\ref{ll:thm:core-rot}: $c\notin U$ since $U\subseteq S$ (Lemma~\ref{ll:lem:peel-basic}(b)); $U$ is
$\rho$-invariant (Lemma~\ref{ll:lem:inv-peel}); and $U$ is $K$-free, since a $K$-element subset of $U$ in convex position
would be one of $S$ (Lemma~\ref{fd:lem:mono}(c)). Theorem~\ref{ll:thm:core-rot} gives (a). Since $x\in U$, $x\ne c$, so
$\operatorname{per}_\rho(x)=m$ (Lemma~\ref{ll:lem:minperiod}) and (b) follows from Lemma~\ref{ll:lem:orbit}(e). For (c),
$\peel(S,n+1)=U\setminus\hv(U)\subseteq U\subseteq\conv(\hv(U))$ by Lemma~\ref{ll:lem:hv}(c).
\end{proof}

\begin{lemma}[the layers exhaust $S$; \leanref{ErdosSzekeres.LayerLemma.exists_layer}{LAYERLEAN/LayerLemma.lean:419}]
\label{ll:lem:exists-layer}
Let $S$ be a finite set of points and $y\in S$. Then there is $n\in\N$ with $y\in\hv(\peel(S,n))$.
\end{lemma}

\begin{proof}
We follow the Lean proof. There are no other hypotheses: no symmetry, no general position, no $K$-freeness.

Claim: for every $n$, either $\card{\peel(S,n)}+n\le\card S$ or $\peel(S,n)=\emptyset$. For $n=0$ the first alternative
holds. Assume the claim for $n$. If $\peel(S,n)=\emptyset$, then $\peel(S,n+1)=\emptyset\setminus\hv(\emptyset)
=\emptyset$. Otherwise the first alternative holds for $n$; moreover $\hv(\peel(S,n))$ is a nonempty subset of
$\peel(S,n)$ (Lemmas~\ref{ll:lem:hv}(d) and~\ref{ll:lem:hv-basic}(b)), so
$\card{\peel(S,n+1)}\le\card{\peel(S,n)}-1$ and $\card{\peel(S,n+1)}+(n+1)\le\card{\peel(S,n)}+n\le\card S$.

Taking $n=\card S+1$: if $y\in\peel(S,\card S+1)$, this set is nonempty, so
$\card{\peel(S,\card S+1)}+\card S+1\le\card S$, which is impossible. Hence $y\notin\peel(S,k)$ for some $k$; let $k$ be
the least such index. Since $\peel(S,0)=S\ni y$, $k\ne0$; write $k=n+1$. By minimality $y\in\peel(S,n)$, and
$y\notin\peel(S,n+1)=\peel(S,n)\setminus\hv(\peel(S,n))$, so $y\in\hv(\peel(S,n))$.
\end{proof}

\begin{theorem}[each peel removes exactly $m$ points;
\leanref{ErdosSzekeres.LayerLemma.card_peel}{LAYERLEAN/LayerLemma.lean:1037}]\label{ll:thm:card-peel}
Let $m,K\in\N$ with $0<m$ and $K\le2m$, let $S$ be a finite set and $c$ a point with $c\notin S$, and assume that $S$ is
$\rho_{c,2\pi/m}$-invariant and $K$-free. Let $t\in\N$ with $m\cdot(t\mathbin{\dot{-}}1)<\card S$. Then for every
$n\le t$,
\[\card{\peel(S,n)}+m\cdot n=\card S .\]
\end{theorem}

\begin{proof}
Induction on $n$, as in Lean. For $n=0$: $\card{\peel(S,0)}+0=\card S$. Let $n+1\le t$ and assume
$\card{\peel(S,n)}+mn=\card S$ (the case $n\le t$ of the induction hypothesis). Then $t\ge1$ and
$n\le t-1=t\mathbin{\dot{-}}1$, so $mn\le m(t\mathbin{\dot{-}}1)<\card S$ and $\card{\peel(S,n)}=\card S-mn>0$: the set
$\peel(S,n)$ is nonempty. Theorem~\ref{ll:thm:layers} gives $x$ with $\hv(\peel(S,n))=\operatorname{orb}_m(\rho,x)$ and
$\card{\operatorname{orb}_m(\rho,x)}=m$, so $\card{\hv(\peel(S,n))}=m$. As $\hv(\peel(S,n))\subseteq\peel(S,n)$, we have
$\card{\peel(S,n)}\ge m$ and
\[\card{\peel(S,n+1)}=\card{\peel(S,n)\setminus\hv(\peel(S,n))}=\card{\peel(S,n)}-m .\]
Hence $\card{\peel(S,n+1)}+m(n+1)=\card{\peel(S,n)}-m+mn+m=\card S$.
\end{proof}

The next corollary is not a Lean declaration; it assembles the preceding results into the layer decomposition of a set
of exactly $mt$ points.

\begin{corollary}[layer decomposition; written step]\label{ll:cor:decomp}
Let $m,K,t\in\N$ with $0<m$ and $K\le2m$, let $S$ be a finite set and $c$ a point with $c\notin S$, and write
$\rho=\rho_{c,2\pi/m}$. Assume that $S$ is $\rho$-invariant and $K$-free and that $\card S=mt$. For $n<t$ put
$L_n=\hv(\peel(S,n))$. Then:
\begin{enumerate}[label=(\alph*)]
\item for every $n<t$ there is $x_n\in\peel(S,n)$ with $L_n=\operatorname{orb}_m(\rho,x_n)$ and $\card{L_n}=m$;
\item for every $n\le t$, $\peel(S,n)$ is the disjoint union of $L_n,L_{n+1},\dots,L_{t-1}$; in particular
  $\peel(S,t)=\emptyset$ and $S=L_0\sqcup\dots\sqcup L_{t-1}$;
\item $L_j\subseteq\peel(S,i+1)\subseteq\conv(L_i)$ for all $i<j<t$.
\end{enumerate}
\end{corollary}

\begin{proof}
If $t=0$, then $S=\emptyset$; (a) and (c) are vacuous and (b) says $\peel(S,0)=\emptyset$. Let $t\ge1$. Then
$m(t\mathbin{\dot{-}}1)=mt-m<mt=\card S$, so Theorem~\ref{ll:thm:card-peel} gives $\card{\peel(S,n)}=m(t-n)$ for $n\le t$.
Thus $\peel(S,n)\ne\emptyset$ for $n<t$, and $\peel(S,t)=\emptyset$. (a) is Theorem~\ref{ll:thm:layers}(a),(b). For
(b): for $n<t$, $L_n\subseteq\peel(S,n)$ and $\peel(S,n+1)=\peel(S,n)\setminus L_n$, so $\peel(S,n)$ is the disjoint
union of $L_n$ and $\peel(S,n+1)$; downward induction on $n$, starting from $\peel(S,t)=\emptyset$, gives (b). For (c):
$L_j\subseteq\peel(S,j)\subseteq\peel(S,i+1)$ by Lemma~\ref{ll:lem:antitone} ($i+1\le j$), and
$\peel(S,i+1)\subseteq\conv(L_i)$ by Theorem~\ref{ll:thm:layers}(c), applicable since $\peel(S,i)\ne\emptyset$.
\end{proof}

\subsection{Reduction to exactly \texorpdfstring{$mt$}{mt} points and the octagon theorems}

\begin{theorem}[reduction to exactly $mt$ points;
\leanref{ErdosSzekeres.LayerLemma.convexNGon_of_layers}{LAYERLEAN/LayerLemma.lean:1066}]\label{ll:thm:reduction}
Let $m,K,t\in\N$ with $0<m$ and $K\le 2m$, and assume the premise
\begin{itemize}
\item[(H)] for every finite set $T$ and every point $c'$: if $T$ is in general position, $c'\notin T$, $T$ is
  $\rho_{c',2\pi/m}$-invariant and $\card T=m\cdot t$, then \lean{HasConvexNGon} $K$ $T$.
\end{itemize}
Then for every finite set $S$ and every point $c$: if $S$ is in general position, $S$ is $\rho_{c,2\pi/m}$-invariant,
and $m\cdot(t\mathbin{\dot{-}}1)<\card{S\setminus\{c\}}$, then \lean{HasConvexNGon} $K$ $S$.
\end{theorem}

In Lean the bound centre in (H) is also called $c$; we rename it $c'$. The size condition is on $S\setminus\{c\}$
(Lean \texttt{S.erase c}), not on $S$, and $c\in S$ is allowed. The premise (H) is an explicit hypothesis of the
theorem, and no theorem of \dirn{LAYERLEAN/LayerLemma.lean} proves it. In the two applications below it is supplied by
the propositions \lean{Sixfold66} and \lean{Sevenfold70}, which are themselves hypotheses (Remark~\ref{ll:rem:hyp}).

\begin{proof}
We follow the Lean proof. Write $\rho=\rho_{c,2\pi/m}$ and suppose that $S$ is $K$-free. Put $S'=S\setminus\{c\}$. Then
$c\notin S'$; $S'$ is $\rho$-invariant by Lemma~\ref{ll:lem:erase}; and $S'$ is $K$-free by
Lemma~\ref{fd:lem:mono}(c), since $S'\subseteq S$.

Theorem~\ref{ll:thm:card-peel} applies to $S'$ with the given $t$ (its size hypothesis $m(t\mathbin{\dot{-}}1)<\card{S'}$
is our hypothesis on $\card{S\setminus\{c\}}$), and with $n=t$ it gives $\card{\peel(S',t)}+mt=\card{S'}$.
Let $T=S'\setminus\peel(S',t)$. Since $\peel(S',t)\subseteq S'$ (Lemma~\ref{ll:lem:peel-basic}(b)),
$\card T=\card{S'}-\card{\peel(S',t)}=mt$.

$T$ is $\rho$-invariant. Let $y\in T$, so $y\in S'$ and $y\notin\peel(S',t)$. Then $\rho(y)\in S'$. Suppose
$\rho(y)\in\peel(S',t)$. By Lemma~\ref{ll:lem:inv-peel} (applied to $S'$; $c\notin S'$) the set $\peel(S',t)$ is
$\rho$-invariant, so $\rho^{m-1}(\rho(y))\in\peel(S',t)$ by Lemma~\ref{ll:lem:orbit}(a). As $y\in S'$, $y\ne c$, so
$\operatorname{per}_\rho(y)=m$ (Lemma~\ref{ll:lem:minperiod}) and, by (P1) and $m\ge1$,
$\rho^{m-1}(\rho(y))=\rho^m(y)=y$. Hence $y\in\peel(S',t)$, a contradiction. So $\rho(y)\in T$.

Moreover $T\subseteq S'\subseteq S$, so $T$ is in general position (Lemma~\ref{fd:lem:mono}(b)), and $c\notin T$. By (H)
with $c'=c$, \lean{HasConvexNGon} $K$ $T$, hence \lean{HasConvexNGon} $K$ $S$ (Lemma~\ref{fd:lem:mono}(c)),
contradicting the assumption. Hence \lean{HasConvexNGon} $K$ $S$.
\end{proof}

\begin{definition}[the propositions \lean{Sixfold66} and \lean{Sevenfold70};
\leanref{ErdosSzekeres.LayerLemma.Sixfold66}{LAYERLEAN/LayerLemma.lean:1102},
\leanref{ErdosSzekeres.LayerLemma.Sevenfold70}{LAYERLEAN/LayerLemma.lean:1107}]\label{ll:def:66-70}
\lean{Sixfold66} is the proposition: for every finite set $T$ and every point $c$, if $T$ is in general position,
$c\notin T$, $T$ is $\rho_{c,2\pi/6}$-invariant and $\card T=66$, then \lean{HasConvexNGon} $8$ $T$.

\lean{Sevenfold70} is the proposition: for every finite set $T$ and every point $c$, if $T$ is in general position,
$c\notin T$, $T$ is $\rho_{c,2\pi/7}$-invariant and $\card T=70$, then \lean{HasConvexNGon} $8$ $T$.

In both, the angle is written with a real numeral, \texttt{2 * Real.pi / 6} and \texttt{2 * Real.pi / 7}
(Remark~\ref{ll:rem:numerals}). Both are definitions of propositions; neither is proved in Lean
(Remark~\ref{ll:rem:hyp}).
\end{definition}

\begin{theorem}[sixfold octagon theorem, conditional on \lean{Sixfold66};
\leanref{ErdosSzekeres.LayerLemma.sixfold_octagon}{LAYERLEAN/LayerLemma.lean:1113}]\label{ll:thm:sixfold}
Assume the proposition \lean{Sixfold66} (a hypothesis, not proved in Lean; Remark~\ref{ll:rem:hyp}). Let $S$ be a finite
set of points and $c$ a point such that $S$ is in general position, $S$ is $\rho_{c,2\pi/6}$-invariant, and
$61\le\card S$. Then \lean{HasConvexNGon} $8$ $S$: some $8$ points of $S$ are in convex position.
\end{theorem}

\begin{proof}
We follow the Lean proof. Write $\rho=\rho_{c,2\pi/6}$; by Remark~\ref{ll:rem:numerals} this is $\rho_{c,2\pi/m}$ for the
natural number $m=6$.

Step 1. There is $p\in S$ with $p\ne c$: otherwise $S\subseteq\{c\}$ and $\card S\le1<61$.

Step 2. $c\notin S$. This is Lemma~\ref{lg:lem:centre} (Lean \lean{centre_not_mem_of_even}) with $m=6$, which is
positive and even, applied to the $\rho$-invariant set $S$ in general position and the point $p$ of Step~1. For
completeness, the argument in this case: suppose $c\in S$. By Lemma~\ref{ll:lem:rot}(d),
$\rho^3=\rho_{c,3\cdot2\pi/6}=\rho_{c,\pi}$, and since $\cos\pi=-1$, $\sin\pi=0$, Definition~\ref{ll:def:rot} gives
$q:=\rho^3(p)=2c-p$, which lies in $S$ by Lemma~\ref{ll:lem:orbit}(a). The points $p,c,q$ of $S$ are pairwise distinct:
$p\ne c$; $q=c$ and $q=p$ would each give $p=c$. But, by the formula of Definition~\ref{fd:def:orient},
\[\orient(p,c,2c-p)=(c_0-p_0)\cdot2(c_1-p_1)-(c_1-p_1)\cdot2(c_0-p_0)=0,\]
so $\{p,c,q\}$ is collinear by Lemma~\ref{fd:lem:collinear}, contradicting general position
(Definition~\ref{fd:def:gp}).

Step 3. Apply Theorem~\ref{ll:thm:reduction} with $m=6$, $K=8$, $t=11$: $0<6$ and $8\le12$. Premise (H) holds: for $T$
and $c'$ as there (the rotation $\rho_{c',2\pi/m}$ with $m=6$ is $\rho_{c',2\pi/6}$, Remark~\ref{ll:rem:numerals}),
$\card T=6\cdot11=66$, and \lean{Sixfold66} gives \lean{HasConvexNGon} $8$ $T$. $S$ is in general
position and $\rho$-invariant. By Step~2, $S\setminus\{c\}=S$, so $6\cdot(11\mathbin{\dot{-}}1)=60<61\le\card S
=\card{S\setminus\{c\}}$. Theorem~\ref{ll:thm:reduction} gives \lean{HasConvexNGon} $8$ $S$.
\end{proof}

\begin{theorem}[sevenfold octagon theorem, conditional on \lean{Sevenfold70};
\leanref{ErdosSzekeres.LayerLemma.sevenfold_octagon}{LAYERLEAN/LayerLemma.lean:1136}]\label{ll:thm:sevenfold}
Assume the proposition \lean{Sevenfold70} (a hypothesis, not proved in Lean; Remark~\ref{ll:rem:hyp}). Let $S$ be a
finite set of points and $c$ a point such that $S$ is in general position, $S$ is $\rho_{c,2\pi/7}$-invariant, and
$65\le\card S$. Then \lean{HasConvexNGon} $8$ $S$. (The centre $c$ may belong to $S$.)
\end{theorem}

\begin{proof}
We follow the Lean proof. Apply Theorem~\ref{ll:thm:reduction} with $m=7$, $K=8$, $t=10$ (Remark~\ref{ll:rem:numerals}
for the angle): $0<7$ and $8\le14$. Premise (H) holds: for $T$ and $c'$ as there (the rotation is
$\rho_{c',2\pi/7}$), $\card T=7\cdot10=70$, and
\lean{Sevenfold70} gives \lean{HasConvexNGon} $8$ $T$. $S$ is in general position and $\rho_{c,2\pi/7}$-invariant, and
$\card{S\setminus\{c\}}\ge\card S-1\ge64>63=7\cdot(10\mathbin{\dot{-}}1)$. Theorem~\ref{ll:thm:reduction} gives
\lean{HasConvexNGon} $8$ $S$. No exclusion of the centre is needed.
\end{proof}

\begin{remark}[status of \lean{Sixfold66} and \lean{Sevenfold70}]\label{ll:rem:hyp}
\lean{Sixfold66} and \lean{Sevenfold70} are Lean definitions of propositions (\texttt{def ... : Prop}), neither theorems
nor axioms. Nothing in \dirn{LAYERLEAN/LayerLemma.lean} proves them; Theorems~\ref{ll:thm:sixfold}
and~\ref{ll:thm:sevenfold} take them as explicit hypotheses, so they do not enter the axiom lists of these theorems.
Among the Lean files of the repository, the only theorems whose conclusion is \lean{Sixfold66} are in
\dirn{ENCBLEAN/geo/ES8Geo/Main.lean} and take further hypotheses (see below), and no theorem has the conclusion
\lean{Sevenfold70}. What is Lean-checked is the implication: \lean{Sixfold66} implies that every
$\rho_{c,2\pi/6}$-invariant set of at least $61$ points in general position has $8$ points in convex position, and
\lean{Sevenfold70} implies the same for $\rho_{c,2\pi/7}$ and at least $65$ points. These octagon statements are
therefore established here only conditionally.

In the project, the two propositions are supported as follows; none of this is part of the present section. The
record \dirn{C6UPPER} gives the unsatisfiability of SAT encodings of octagon-free configurations of eleven nested
regular hexagons in general position, by cake\_lpr verdicts on cube formulas whose cover is checked there: the formula
of encoder~B (\dirn{C6UPPER/encb/gen_b2.py}) alone, through $64$ cubes, and that of encoder~A
(\dirn{C7UPPER/gen.py}) with an added symmetry-breaking unit clause, through $32$ cubes. The record \dirn{C7UPPER}
gives the unsatisfiability of two separately implemented encodings of the same kind for ten nested regular heptagons
(encoder~A's with an added symmetry-breaking unit clause, encoder~B's alone), each by a cake\_lpr verdict. In the
terms of this document these unsatisfiability results are \emph{computer-assisted}: they rest on LRAT proofs checked
by cake\_lpr, outside Lean. The project's records call such a result a ``certificate'' (an unsatisfiability claim whose
LRAT proof cake\_lpr accepted; \dirn{WRITEUP/note/symmetric_extremal.tex}, lines 109--114); that is a label of the
project, not one of the labels of this document. The step from these unsatisfiability results to \lean{Sixfold66} and
\lean{Sevenfold70} is the written soundness proof of the encodings in \dirn{WRITEUP/note/k8_soundness.tex} (a written
proof, not refereed), which also covers the symmetry breaking and the
cube cover. That note lists the gaps that remain, (G1)--(G4): the correspondence between the generating programs and
their clauses is established by reading the source code; the passage from the formulas to \lean{Sixfold66} and
\lean{Sevenfold70} is a written proof; the angular part of encoder~A rests on binary64 floating-point arithmetic;
and several of its written proofs have had no internal review. For $m=6$ the unit \dirn{ENCBLEAN} also contains a Lean
derivation of \lean{Sixfold66}, \lean{EncB.Geo.sixfold66_of_cubes64}; its \lean{Sixfold66} is that of
\dirn{ENCBLEAN/geo/ES8Geo/Capstone2.lean}, which contains a verbatim copy of the declarations of this file (over a
byte-identical copy of \dirn{ErdosSzekeres/Definitions.lean}). Its hypothesis \lean{EncB.CubesUnsat64} is the
unsatisfiability of the $64$ cube formulas built from a Lean definition of encoder~B's formula (\lean{EncB.encB}).
Neither the identity of the printed texts of these formulas with the cube files that carry the cake\_lpr verdicts nor
the reading of those verdicts as unsatisfiability is Lean-checked; both rest on the trust items (E1)--(E6) listed
under gap~(G2) of that note.

The link with nested polygons is Corollary~\ref{ll:cor:decomp}: a set $T$ that violates \lean{Sixfold66} is in general
position, centre-free, $\rho_{c,2\pi/6}$-invariant, $8$-free and has $66=6\cdot11$ points, so (with $m=6$, $K=8\le12$,
$t=11$) it is the disjoint union of eleven $\rho_{c,2\pi/6}$-orbits $L_0,\dots,L_{10}$ of $6$ points each, with
$L_j\subseteq\conv(L_i)$ for $i<j$; likewise a set violating \lean{Sevenfold70} is the disjoint union of ten orbits of
$7$ points.
\end{remark}

\begin{remark}[real-cast numerals]\label{ll:rem:numerals}
In the general statements of this section the angle is \texttt{2 * Real.pi / m} with $m\in\N$ cast to $\R$ (C7). In
\lean{Sixfold66}, \lean{Sevenfold70}, \lean{sixfold_octagon} and \lean{sevenfold_octagon} the angle is
\texttt{2 * Real.pi / 6}, resp.\ \texttt{2 * Real.pi / 7}, where $6$ and $7$ are real numerals. The cast of the
natural number
$6$ is the real number $6$ (likewise for $7$), so $\rho_{c,2\pi/6}$ is the same map in both readings; the Lean proofs
bridge the two forms with \texttt{exact\_mod\_cast}, and the cardinality conditions $\card T=6\cdot 11$ and $\card T=66$
(resp.\ $7\cdot10$ and $70$) with \texttt{simpa}. Mathematically nothing changes.
\end{remark}

\paragraph{Coverage.}
All declarations of this section's scope in \dirn{LAYERLEAN/LayerLemma.lean} are covered:
\lean{orbitFS} (Definition~\ref{ll:def:orbit});
\lean{iterate_mem}, \lean{mem_orbitFS_self}, \lean{orbitFS_subset}, \lean{orbitFS_invariant},
\lean{card_orbitFS_of_minimalPeriod}, \lean{mem_orbitFS_of_not_disjoint} (Lemma~\ref{ll:lem:orbit});
\lean{Invariant} (Definition~\ref{ll:def:inv});
\lean{hullVerts} (Definition~\ref{ll:def:hv}); \lean{mem_hullVerts}, \lean{hullVerts_subset}
(Lemma~\ref{ll:lem:hv-basic}); \lean{convexPosition_hullVerts}, \lean{extreme_mem_hullVerts},
\lean{subset_convexHull_hullVerts}, \lean{hullVerts_nonempty} (Lemma~\ref{ll:lem:hv}; the proofs of (b)--(d) are
elementary and differ from the Lean proofs, which use two Mathlib facts on extreme points for (b), the Krein--Milman
theorem for (c) and the Krein--Milman lemma for (d));
\lean{map_mem_hullVerts} (Lemma~\ref{ll:lem:map-hv});
\lean{hullVerts_eq_orbitFS} (Theorem~\ref{ll:thm:core});
\lean{invariant_sdiff_hullVerts} (Lemma~\ref{ll:lem:inv-sdiff});
\lean{toC}, \lean{ofC}, \lean{crot} (Definition~\ref{ll:def:complex});
\lean{toC_re}, \lean{toC_im}, \lean{ofC_zero}, \lean{ofC_one}, \lean{toC_ofC}, \lean{ofC_toC}, \lean{toC_injective},
\lean{toC_crot}, \lean{crot_one}, \lean{crot_comp}, \lean{crot_iterate}, \lean{crot_eq_self_iff}, \lean{crot_injective}
(Lemma~\ref{ll:lem:complex});
\lean{rotAbout} (Definition~\ref{ll:def:rot});
\lean{rotAbout_eq_crot}, \lean{rotAff}, \lean{rotAbout_injective}, \lean{rotAbout_iterate} (Lemma~\ref{ll:lem:rot});
\lean{invariant_erase_centre} (Lemma~\ref{ll:lem:erase});
\lean{minimalPeriod_rotAbout} (Lemma~\ref{ll:lem:minperiod});
\lean{hullVerts_eq_orbitFS_rotAbout} (Theorem~\ref{ll:thm:core-rot});
\lean{peel} (Definition~\ref{ll:def:peel}); \lean{peel_succ}, \lean{peel_subset} (Lemma~\ref{ll:lem:peel-basic});
\lean{peel_antitone} (Lemma~\ref{ll:lem:antitone});
\lean{invariant_peel} (Lemma~\ref{ll:lem:inv-peel});
\lean{layers_are_orbits} (Theorem~\ref{ll:thm:layers});
\lean{exists_layer} (Lemma~\ref{ll:lem:exists-layer});
\lean{card_peel} (Theorem~\ref{ll:thm:card-peel});
\lean{convexNGon_of_layers} (Theorem~\ref{ll:thm:reduction});
\lean{Sixfold66}, \lean{Sevenfold70} (Definition~\ref{ll:def:66-70});
\lean{sixfold_octagon} (Theorem~\ref{ll:thm:sixfold}, conditional on \lean{Sixfold66});
\lean{sevenfold_octagon} (Theorem~\ref{ll:thm:sevenfold}, conditional on \lean{Sevenfold70}).
Corollary~\ref{ll:cor:decomp} is a written step, not a Lean declaration. Nothing in the scope is deliberately left out.
The declarations of lines 471--1032 of the same file (from \lean{orient_rotAbout} to \lean{dist2_le_of_later}) are
outside this section's scope; \lean{rotAbout_iterate} (line 465) is included here because it lies in the range of this
section, and Step~2 of our proof of Theorem~\ref{ll:thm:sixfold} uses it.

\section{Geometry of the layers: nesting, chains and polar rules}\label{lg:sec}

This section proves the geometric facts about the convex layers of a rotation-invariant finite point set that the
layered encoders rely on, together with the elementary identities behind them. In the statements about the layers
(Lemmas~\ref{lg:lem:layer} and~\ref{lg:lem:halfplane}, Theorems~\ref{lg:thm:nesting}, \ref{lg:thm:radius}
and~\ref{lg:thm:chain}) and in Lemma~\ref{lg:lem:centre}, $S$ is a finite set of points that is mapped into itself by
the counterclockwise rotation $\rho$ by $2\pi/m$ ($m\ge1$) about a point $c$; in all of them except
Lemma~\ref{lg:lem:centre}, also $c\notin S$ and $S$ has no $K$ points in convex position for some $K\le 2m$.
Theorem~\ref{lg:thm:fan}, the only other statement with a set $S$, involves no centre and no rotation. The section
contains:
\begin{itemize}
\item \emph{rotations and regular polygons} (\S\ref{lg:ssec:rot}): iterates of a rotation, the orientation of three
  vertices of a regular polygon (Lemmas~\ref{lg:lem:orient-rot}, \ref{lg:lem:trig}, \ref{lg:lem:regular},
  \ref{lg:lem:next-prev}) and the position of the vertices relative to the tangent at a vertex
  (Lemmas~\ref{lg:lem:lin-rot}, \ref{lg:lem:lin-orbit});
\item \emph{convexity} of the sets $\{q:0\le\orient(a,b,q)\}$ (closed half-planes when $a\ne b$) and of closed disks
  (\S\ref{lg:ssec:convex});
\item \emph{nesting, centre and radii} (\S\ref{lg:ssec:nest}): if $m\ge2$ and $S$ is in general position, every point
  of a later layer lies strictly to the left of every edge $x\to\rho(x)$ of an earlier layer
  (Theorem~\ref{lg:thm:nesting}); if $m\ge2$ is even and $S$ is in general position and has a point other than $c$,
  then $c\notin S$ (Lemma~\ref{lg:lem:centre}); the points of a layer and of all later layers are no farther from $c$
  than the vertices of that layer (Theorem~\ref{lg:thm:radius});
\item \emph{fans and chains} (\S\ref{lg:ssec:chain}): a pivot followed by $n\ge3$ chain points with positive fan and
  turn orientations, all in an open half-plane whose boundary line passes through the pivot, gives $n+1$ points in
  convex position (Theorem~\ref{lg:thm:fan}; by its proof, the pivot and the chain points themselves); if $m\ge3$ and
  $S$ is in general position, all other points of the layers $\ge i$ lie strictly on the inner side of the tangent at
  a vertex of layer $i$ (Lemma~\ref{lg:lem:halfplane}); hence, if moreover $K\ge4$, no vertex of layer $i$ is the pivot
  of a sequence of $K-1$ points of the layers $\ge i$, all different from it, with positive fan and turn orientations
  (Theorem~\ref{lg:thm:chain}), which is the geometric content of the chain clauses (see the paragraph
  \emph{Encoders} below);
\item \emph{polar rules} (\S\ref{lg:ssec:polar}): the orientation of three points given in polar coordinates about $c$
  (Lemma~\ref{lg:lem:polar}) and the angular rules A and B (Lemmas~\ref{lg:lem:ruleA}, \ref{lg:lem:ruleB}).
\end{itemize}

\paragraph{Sources.} Every Lean declaration that heads a statement of this section is in the file
\dirn{LAYERLEAN/LayerLemma.lean} (namespace \lean{ErdosSzekeres.LayerLemma}, 1148 lines, SHA-256
\texttt{afb8542dceb0}\dots), lines 465--1032; its README is \dirn{LAYERLEAN/README.md} (SHA-256
\texttt{9eaec241a1d2}\dots). The declarations of the same file that \S\ref{lg:ssec:notation} recalls lie outside
these lines and are restated in the section on the layer lemma. The Lean proofs of the declarations stated here use the following
declarations of the main library; the foundations section restates them and proves the lemmas among them.
\begin{itemize}
\item The definitions \lean{Point}, \lean{GeneralPosition}, \lean{ConvexPosition} and \lean{HasConvexNGon}
  of\newline \dirn{ErdosSzekeres/Definitions.lean} (SHA-256 \texttt{2425500eefd1}\dots).
\item The orientation \lean{ErdosSzekeres.orient} and the lemma \lean{ErdosSzekeres.orient_cycle},\newline
  in \dirn{ErdosSzekeres/Orientation.lean}, lines 9 and 15 (SHA-256 \texttt{f67955ac7091}\dots).
\item The lemma \lean{ErdosSzekeres.generalPosition_iff_orient_ne_zero},\newline
  in \dirn{ErdosSzekeres/OrientationGeometry.lean:66} (SHA-256 \texttt{c7db9bbbb65e}\dots).
\item The lemma \lean{ErdosSzekeres.Suk.convexPosition_of_positive_triples},\newline
  in \dirn{ErdosSzekeres/Suk/OrientationClique.lean:9} (SHA-256 \texttt{cf63d5219b47}\dots).
\end{itemize}
The written proofs use the results listed under \emph{Dependencies} below, among them Lemma~\ref{fd:lem:mono}(c)
(Lean \lean{HasConvexNGon.mono}, \dirn{ErdosSzekeres/Basic.lean:20}), which the Lean proofs do not use.

The file was checked with Lean~4.24.0 and Mathlib \texttt{f897ebcf} by \texttt{lake env lean} from the repository root;
it is not part of the default \texttt{lake build} target. The log \dirn{LAYERLEAN/logs/build.log} (exit code~0, no
messages) names no file hash; it was committed with the first version of the file (454 lines, commit
\texttt{6429e9a70}) and was not updated by the five later commits that changed the file, so it is not a record of the
present file, and we do not use it. The records for the present file are the review report
\dirn{LEANEXT/review/REPORT.md} of 4 November 2026, which reports exit~0 and no messages for the reviewed file, whose
SHA-256 it gives as \texttt{afb8542d}\dots{} (1148 lines), and the axiom re-check of 6 November 2026 (whose output is
named below), which compiled
\dirn{LAYERLEAN/logs/axioms_recheck_20261106/LayerLemmaAxioms.lean}, whose first 1148 lines are this file, with exit
status~0; the files \dirn{README.txt} and \dirn{run.sh} of that directory record the toolchain and the command given
above. The file contains no \lean{sorry}, \lean{admit}, \lean{axiom} or \lean{native_decide}.

\paragraph{Axioms.} The log \dirn{LAYERLEAN/logs/axioms.log} (SHA-256 \texttt{3c4b7aac502e}\dots) prints the axioms of
16 of the 23 theorems of this section: all except \lean{rotAbout_iterate}, \lean{orient_rotAbout},
\lean{convex_orient_nonneg}, \lean{orient_eq_lin}, \lean{lin_rotAbout}, \lean{lin_orbit_pos} and
\lean{orient_next_prev_pos}. The review log \dirn{LEANEXT/review/axioms.log} of 4 November 2026 (SHA-256
\texttt{386c4f578725}\dots) prints all 23. It covers 59 declarations of the reviewed file (SHA-256
\texttt{afb8542d}\dots, the file above), namely all its \texttt{theorem} declarations except the 7 marked
\texttt{@[simp]}; the file has 66 in all, and \dirn{LEANEXT/review/REPORT.md} speaks of ``all 59 theorem
declarations''. In each of the two logs, every theorem of this section that it prints depends only on the standard
axioms \lean{propext}, \lean{Classical.choice} and \lean{Quot.sound}. The re-check
\dirn{LAYERLEAN/logs/axioms_recheck_20261106/out.txt} (SHA-256 \texttt{a59c3816cbe7}\dots) gives the same result for
the 11 of them that it lists.

\paragraph{Rigor.} Each lemma or theorem below whose head carries a Lean name is \emph{Lean-checked}, with the axioms
just listed; none of them takes an unproved proposition as a hypothesis, and none is a finite computation. Each definition
whose head carries a Lean name restates that Lean definition. The proofs written
here are \emph{written proofs (not refereed)}. They follow the structure of the Lean proofs. Where Lean checks an
algebraic identity or inequality with an automatic tactic (such as \lean{ring}, \lean{linear_combination},
\lean{nlinarith} or \lean{positivity}), the written proof derives it by hand; every other difference from the Lean
proof is stated in a remark or in a parenthetical note in the proof. Statements tagged ``written step'' are not Lean
declarations, and they are proved here. Lemma~\ref{lg:lem:layer} collects the opening steps of the Lean proofs of
Theorems~\ref{lg:thm:nesting} and~\ref{lg:thm:radius} and Lemma~\ref{lg:lem:halfplane}. Lemma~\ref{lg:lem:rotfacts} collects
elementary facts about rotations that the Lean proofs obtain differently (through the complex form of the rotation,
by direct computation, or through minimal periods). Lemma~\ref{lg:lem:cramer} is a general identity of which Lean
checks only the three instances used, each by \lean{ring}. Corollary~\ref{lg:cor:period} restates
Lemma~\ref{lg:lem:rot-iter} for $\rho$ and adds $\rho^m=\mathrm{id}$, which replaces a minimal-period argument of the
Lean proof of Lemma~\ref{lg:lem:halfplane}.

\paragraph{Encoders.} Several results below are the geometric facts behind parts of the project's layered encoders:
Theorem~\ref{lg:thm:nesting} behind the nesting clauses, Theorem~\ref{lg:thm:chain} behind the chain clauses, and
Lemmas~\ref{lg:lem:ruleA} and~\ref{lg:lem:ruleB}, with the weak radius order of Theorem~\ref{lg:thm:radius} that
rule~B needs, behind the angular rules A and B, which the encoders apply to triples of points. Only the geometric
facts are stated and proved here. That the
encoders' programs emit exactly these clause families, and which rule they apply to which triple, are not part of
\dirn{LAYERLEAN/LayerLemma.lean}; \dirn{LAYERLEAN/README.md} lists both under ``What is not in Lean'' and gives their
status as ``written proof, review and tests''. For the encoders of the $k=8$ bounds the code-to-clause step is
recorded as gap~(G1) of \dirn{WRITEUP/note/k8_soundness.tex}: it rests on a written proof (not refereed) by reading
the code, with the Python interpreter trusted to run the programs as written. For encoder~B at $m=6$ that note also
records a second route on which the step is Lean-checked for a Lean definition of the formula (\dirn{ENCBLEAN}). On
that route the note lists, under gap~(G2), the trust items (E1)--(E6), which are not Lean-checked: among them, that the
texts printed from this definition are the $64$ cube files that carry the cake\_lpr verdicts is (E1), evidence from
running compiled code and not a Lean theorem, and those $64$ verdicts and their provenance are covered by (E2).

\paragraph{Dependencies.} The section uses the conventions and facts of the foundations section
(Definitions~\ref{fd:def:point}, \ref{fd:def:orient}, \ref{fd:def:gp}, \ref{fd:def:cp}, \ref{fd:def:ngon},
\ref{fd:def:inv}, \ref{fd:def:rot} and Lemmas~\ref{fd:lem:orient}, \ref{fd:lem:collinear}, \ref{fd:lem:mono},
\ref{fd:lem:pos-triples}) and, from the section on the layer lemma, Definitions~\ref{ll:def:hv} and~\ref{ll:def:peel},
Lemmas~\ref{ll:lem:hv}, \ref{ll:lem:antitone}, \ref{ll:lem:inv-peel} and Theorem~\ref{ll:thm:core-rot}. These,
except Lemma~\ref{fd:lem:pos-triples}, which is used only in Step~3 of the proof of Theorem~\ref{lg:thm:fan}, are
recalled in \S\ref{lg:ssec:notation}.

\subsection{Notation and the results used}\label{lg:ssec:notation}

\paragraph{Vectors.} Points are elements of $\R^2$, written $p=(p_0,p_1)$ (Definition~\ref{fd:def:point}; Lean
\lean{Point}). For $X,Y\in\R^2$ put
\[
  \det(X,Y)=X_0Y_1-X_1Y_0,\qquad \langle X,Y\rangle=X_0Y_0+X_1Y_1,\qquad \abs{X}^2=\langle X,X\rangle,\qquad
  JX=(-X_1,X_0).
\]
So $J$ is the counterclockwise quarter turn, $\det$ is bilinear and antisymmetric, and $\det(X,Y)=\langle JX,Y\rangle$.
In particular
\begin{equation}\label{lg:eq:J}
\begin{gathered}
  \det(X,X)=0,\qquad \det(X,JX)=\abs{X}^2,\qquad \det(JX,X)=-\abs{X}^2,\\
  \det(JX,JX)=0,\qquad \langle X,JX\rangle=0 .
\end{gathered}
\end{equation}

\paragraph{Orientation.} We use the following facts about $\orient$ (Definition~\ref{fd:def:orient}).
\begin{itemize}
\item[(O1)] $\orient(a,b,d)=(b_0-a_0)(d_1-a_1)-(b_1-a_1)(d_0-a_0)=\det(b-a,\,d-a)$.
\item[(O2)] $\orient$ is invariant under cyclic permutations of its three arguments (Lemma~\ref{fd:lem:orient}(b)).
  Hence $\orient(a,b,o)=\orient(o,a,b)=\det(a-o,\,b-o)$ for all points $a,b,o$.
\item[(O3)] $\orient(a,b,d)=0$ whenever two of $a,b,d$ coincide. Indeed, by (O1), $\orient(a,a,d)=\det(0,d-a)=0$,
  $\orient(a,b,a)=\det(b-a,0)=0$ and $\orient(a,b,b)=\det(b-a,b-a)=0$.
\item[(O4)] For fixed $a,b$ the map $q\mapsto\orient(a,b,q)=\det(b-a,\,q-a)$ is affine: if $u+v=1$ then
  $\orient(a,b,uq_1+vq_2)=u\orient(a,b,q_1)+v\orient(a,b,q_2)$, because $(uq_1+vq_2)-a=u(q_1-a)+v(q_2-a)$ and $\det$ is
  linear in its second argument.
\end{itemize}

\paragraph{Rotations.} For $\theta\in\R$ let $R_\theta$ be the matrix with rows $(\cos\theta,-\sin\theta)$ and
$(\sin\theta,\cos\theta)$, so that $R_\theta X=\cos\theta\,X+\sin\theta\,JX$. Definition~\ref{fd:def:rot} reads
\[
  \operatorname{rotAbout}(c,\theta)(p)=c+R_\theta(p-c)
  =\bigl(c_0+\cos\theta\,(p_0-c_0)-\sin\theta\,(p_1-c_1),\;c_1+\sin\theta\,(p_0-c_0)+\cos\theta\,(p_1-c_1)\bigr).
\]
For a map $f$ and $j\in\N$, $f^{j}$ is the $j$-fold iterate: $f^{0}=\mathrm{id}$ and $f^{j+1}=f\circ f^{j}$
($=f^{j}\circ f$; Lean \texttt{f\^{}[j]}). For $m\in\N$ the angle $2\pi/m$ is the real number $2\pi/(m:\R)$ (Lean
\texttt{2 * Real.pi / m}, with $m$ cast to $\R$; Lean's convention would give $0$ for $m=0$, but every statement of
this section in which $2\pi/m$ occurs assumes $m\ge1$).
Whenever $c$ and $m$ are given we write
\[
  \rho=\operatorname{rotAbout}(c,2\pi/m).
\]
A natural-number subtraction such as $m-1$ or $K-1$ is truncated at $0$ in Lean; in every use below the minuend is at
least the subtrahend, so it is the ordinary difference.

\paragraph{Sine and cosine.} Besides the values at $0$, $\pi$ and $2\pi$, the addition, subtraction and double-angle
formulas, $\sin^2+\cos^2=1$, the oddness of the sine and the $2\pi$-periodicity, we use two standard facts about the real sine
and cosine (see, e.g., W.~Rudin, \emph{Principles of Mathematical Analysis}, 3rd ed., Chapter~8; in Mathlib
\lean{Real.sin_pos_of_pos_of_lt_pi}, \lean{Real.cos_le_one} and \lean{Real.cos_eq_one_iff}):
\begin{itemize}
\item[(S1)] $\sin\theta>0$ for $0<\theta<\pi$; hence $\sin\theta\ge0$ for $0\le\theta\le\pi$, as $\sin0=\sin\pi=0$;
\item[(S2)] $\cos\theta\le1$ for every $\theta\in\R$, and $\cos\theta=1$ if and only if $\theta\in2\pi\Z$.
\end{itemize}

\paragraph{Hypotheses on $S$.} All point sets are finite (Lean \texttt{Finset Point}). ``$\rho$ maps $S$ into $S$'' is the
Lean hypothesis \texttt{Invariant (rotAbout c (2 * Real.pi / m)) S}, i.e.\ $\rho(p)\in S$ for all $p\in S$
(Definition~\ref{fd:def:inv}); by induction on $i$ it implies $\rho^{i}(p)\in S$ for all $p\in S$ and $i\in\N$.
``$S$ is $K$-free'' means $\neg$\,\texttt{HasConvexNGon K S}: no $K$-element subset of $S$ is in convex position
(Definitions~\ref{fd:def:cp}, \ref{fd:def:ngon}). By Lemma~\ref{fd:lem:mono}(c), every subset of a $K$-free set is
$K$-free. ``General position'' is Definition~\ref{fd:def:gp}; by Lemma~\ref{fd:lem:collinear}, $S$ is in general
position if and only if $\orient(a,b,d)\ne0$ for all pairwise distinct $a,b,d\in S$.

\paragraph{Convexity.} A set $C\subseteq\R^2$ is convex (Mathlib \lean{Convex}) if $uy_1+vy_2\in C$ whenever
$y_1,y_2\in C$, $u,v\ge0$ and $u+v=1$. The convex hull $\conv(A)$ is the smallest convex set containing $A$; so
$A\subseteq C$ with $C$ convex implies $\conv(A)\subseteq C$ (Mathlib \lean{convexHull_min}).

\paragraph{Layers.} $\hv(X)=\{x\in X: x\notin\conv(X\setminus\{x\})\}\subseteq X$ (Definition~\ref{ll:def:hv}), and
$\peel(S,0)=S$, $\peel(S,n+1)=\peel(S,n)\setminus\hv(\peel(S,n))$ (Definition~\ref{ll:def:peel}). For a map $f$, a
number $k\in\N$ and a point $x$, Lean's \texttt{orbitFS f k x} (\dirn{LAYERLEAN/LayerLemma.lean:48}) is the finite set
$\{f^{j}(x):0\le j<k\}$. We use the following results of the section on the layer lemma.
\begin{itemize}
\item[(L1)] For every finite $X$: $X\subseteq\conv(\hv(X))$ (Lemma~\ref{ll:lem:hv}(c), \lean{subset_convexHull_hullVerts}).
\item[(L2)] If $n\le n'$ then $\peel(S,n')\subseteq\peel(S,n)$ (Lemma~\ref{ll:lem:antitone}); in particular
  $\peel(S,n)\subseteq\peel(S,0)=S$.
\item[(L3)] If $0<m$, $c\notin S$ and $\rho$ maps $S$ into $S$, then $\rho$ maps $\peel(S,n)$ into $\peel(S,n)$ for every
  $n$ (Lemma~\ref{ll:lem:inv-peel}, \lean{invariant_peel}).
\item[(L4)] If $0<m$, $K\le 2m$, $c\notin S$, $\rho$ maps $S$ into $S$, $S$ is $K$-free and $x\in\hv(S)$, then
  $\hv(S)=\{\rho^{j}(x):0\le j<m\}$ (Theorem~\ref{ll:thm:core-rot}, \lean{hullVerts_eq_orbitFS_rotAbout}).
\end{itemize}

The Lean proofs of Theorems~\ref{lg:thm:nesting} and~\ref{lg:thm:radius} and of Lemma~\ref{lg:lem:halfplane} each
begin with the same few steps; we collect them once.

\begin{lemma}[layers of a layered configuration; written step]\label{lg:lem:layer}
Let $m,K\in\N$ with $0<m$ and $K\le 2m$, let $S$ be a finite set and $c$ a point with $c\notin S$, let
$\rho=\operatorname{rotAbout}(c,2\pi/m)$ map $S$ into $S$, and let $S$ be $K$-free. Let $n\in\N$, put
$L=\hv(\peel(S,n))$, and let $x\in L$. Then:
\begin{enumerate}[label=(\alph*)]
\item $L=\{\rho^{j}(x):0\le j<m\}$; in particular $\rho^{j}(x)\in L$ for $0\le j<m$;
\item $L\subseteq\peel(S,n)\subseteq S$, and $x\neq c$;
\item $\peel(S,n')\subseteq\conv(L)$ for every $n'\ge n$;
\item $\peel(S,n')\cap L=\emptyset$ for every $n'>n$.
\end{enumerate}
\end{lemma}

\begin{proof}
By (L2), $\peel(S,n)\subseteq S$. Hence $c\notin\peel(S,n)$, and $\peel(S,n)$ is $K$-free (a $K$-element subset of
$\peel(S,n)$ in convex position would be one of $S$). By (L3), $\rho$ maps $\peel(S,n)$ into itself. So (L4) applies to
the finite set $\peel(S,n)$ and its hull vertex $x$, and gives (a).
(b) $L\subseteq\peel(S,n)$ by the definition of $\hv$, and $\peel(S,n)\subseteq S$; since $x\in L\subseteq S$ and
$c\notin S$, $x\ne c$.
(c) If $n'\ge n$ then $\peel(S,n')\subseteq\peel(S,n)\subseteq\conv(\hv(\peel(S,n)))=\conv(L)$ by (L2) and (L1).
(d) If $n'>n$ then $n+1\le n'$, so $\peel(S,n')\subseteq\peel(S,n+1)=\peel(S,n)\setminus L$ by (L2) and the definition
of $\peel$.
\end{proof}

\subsection{Rotations and regular polygons}\label{lg:ssec:rot}

\begin{lemma}[rotation facts; written step]\label{lg:lem:rotfacts}
For all points $c,p$ and all $\alpha,\beta,\theta\in\R$:
\begin{enumerate}[label=(\alph*)]
\item $\operatorname{rotAbout}(c,0)(p)=p$;
\item $\operatorname{rotAbout}(c,\alpha)\bigl(\operatorname{rotAbout}(c,\beta)(p)\bigr)=\operatorname{rotAbout}(c,\alpha+\beta)(p)$;
\item $\operatorname{rotAbout}(c,2\pi)(p)=p$;
\item $\abs{\operatorname{rotAbout}(c,\theta)(p)-c}^2=\abs{p-c}^2$.
\end{enumerate}
\end{lemma}

\begin{proof}
Put $u=p-c$, so that $\operatorname{rotAbout}(c,\theta)(p)=c+R_\theta u$.
(a), (c): $\cos0=\cos2\pi=1$ and $\sin0=\sin2\pi=0$, so $R_0u=R_{2\pi}u=u$.
(b) $\operatorname{rotAbout}(c,\beta)(p)-c=R_\beta u$, hence
$\operatorname{rotAbout}(c,\alpha)(\operatorname{rotAbout}(c,\beta)(p))=c+R_\alpha R_\beta u$. The product
$R_\alpha R_\beta$ has rows $(\cos\alpha\cos\beta-\sin\alpha\sin\beta,\;-(\sin\alpha\cos\beta+\cos\alpha\sin\beta))$ and
$(\sin\alpha\cos\beta+\cos\alpha\sin\beta,\;\cos\alpha\cos\beta-\sin\alpha\sin\beta)$, which is $R_{\alpha+\beta}$ by
the addition formulas $\cos(\alpha+\beta)=\cos\alpha\cos\beta-\sin\alpha\sin\beta$ and
$\sin(\alpha+\beta)=\sin\alpha\cos\beta+\cos\alpha\sin\beta$.
(d) $\abs{R_\theta u}^2=(\cos\theta\,u_0-\sin\theta\,u_1)^2+(\sin\theta\,u_0+\cos\theta\,u_1)^2
=(\cos^2\theta+\sin^2\theta)(u_0^2+u_1^2)=\abs{u}^2$; the mixed terms cancel.
\end{proof}

\begin{lemma}[iterates of the rotation about a point; \leanref{ErdosSzekeres.LayerLemma.rotAbout_iterate}{LAYERLEAN/LayerLemma.lean:465}]\label{lg:lem:rot-iter}
For every point $c$, every $\alpha\in\R$ and every $j\in\N$,
\[
  \bigl(\operatorname{rotAbout}(c,\alpha)\bigr)^{j}=\operatorname{rotAbout}(c,\,j\alpha)
\]
as maps $\R^2\to\R^2$, where $j$ is cast to $\R$.
\end{lemma}

\begin{proof}
Induction on $j$. For $j=0$ the left side is the identity by definition of the iterate, and the right side is
$\operatorname{rotAbout}(c,0)=\mathrm{id}$ by Lemma~\ref{lg:lem:rotfacts}(a). If the claim holds for $j$, then
\[
  \bigl(\operatorname{rotAbout}(c,\alpha)\bigr)^{j+1}=\operatorname{rotAbout}(c,\alpha)\circ\operatorname{rotAbout}(c,j\alpha)
  =\operatorname{rotAbout}(c,\alpha+j\alpha)=\operatorname{rotAbout}(c,(j+1)\alpha)
\]
by Lemma~\ref{lg:lem:rotfacts}(b). (Lean argues instead with the complex form: in the complex coordinate $z$, with
$\gamma=c_0+ic_1$, the map $\operatorname{rotAbout}(c,\alpha)$ is $z\mapsto\gamma+e^{i\alpha}(z-\gamma)$
(\lean{rotAbout_eq_crot}, line 302), whose $j$-th iterate is $z\mapsto\gamma+(e^{i\alpha})^{j}(z-\gamma)$
(\lean{crot_iterate}, line 267), and $(e^{i\alpha})^{j}=e^{ij\alpha}$.)
\end{proof}

\begin{corollary}[$\rho^m$ is the identity; written step]\label{lg:cor:period}
Let $c$ be a point, $m\in\N$ with $0<m$, and $\rho=\operatorname{rotAbout}(c,2\pi/m)$. Then $\rho^{j}=
\operatorname{rotAbout}(c,2\pi j/m)$ for every $j\in\N$, and $\rho^{m}(p)=p$ for every point $p$.
\end{corollary}

\begin{proof}
By Lemma~\ref{lg:lem:rot-iter}, $\rho^{j}=\operatorname{rotAbout}(c,j\cdot(2\pi/m))$. For $j=m$ the angle is
$m\cdot(2\pi/m)=2\pi$ because $(m:\R)\ne0$, and $\operatorname{rotAbout}(c,2\pi)=\mathrm{id}$ by
Lemma~\ref{lg:lem:rotfacts}(c).
\end{proof}

\begin{lemma}[orientation of a regular-polygon triple; \leanref{ErdosSzekeres.LayerLemma.orient_rotAbout}{LAYERLEAN/LayerLemma.lean:472}]\label{lg:lem:orient-rot}
For all points $c,x$ and all $\alpha,\beta\in\R$,
\[
  \orient\bigl(x,\operatorname{rotAbout}(c,\alpha)(x),\operatorname{rotAbout}(c,\beta)(x)\bigr)
  =\bigl((x_0-c_0)^2+(x_1-c_1)^2\bigr)\bigl((\cos\alpha-1)\sin\beta-\sin\alpha\,(\cos\beta-1)\bigr).
\]
\end{lemma}

\begin{proof}
Put $u=x-c$, $b=\operatorname{rotAbout}(c,\alpha)(x)$ and $d=\operatorname{rotAbout}(c,\beta)(x)$. Then
$b-x=R_\alpha u-u=(\cos\alpha-1)u+\sin\alpha\,Ju$ and $d-x=(\cos\beta-1)u+\sin\beta\,Ju$. By (O1), bilinearity of
$\det$ and~\eqref{lg:eq:J},
\begin{align*}
  \orient(x,b,d)=\det(b-x,d-x)
  &=(\cos\alpha-1)(\cos\beta-1)\det(u,u)+(\cos\alpha-1)\sin\beta\,\det(u,Ju)\\
  &\quad+\sin\alpha\,(\cos\beta-1)\det(Ju,u)+\sin\alpha\sin\beta\,\det(Ju,Ju)\\
  &=\bigl((\cos\alpha-1)\sin\beta-\sin\alpha\,(\cos\beta-1)\bigr)\abs{u}^2,
\end{align*}
and $\abs{u}^2=(x_0-c_0)^2+(x_1-c_1)^2$. No relation between $\cos$ and $\sin$ is used: this is a polynomial identity in
$\cos\alpha,\sin\alpha,\cos\beta,\sin\beta$ and the coordinates, which is how Lean proves it (\lean{ring}).
\end{proof}

\begin{lemma}[a product formula for the orientation factor; \leanref{ErdosSzekeres.LayerLemma.trig_layer}{LAYERLEAN/LayerLemma.lean:480}]\label{lg:lem:trig}
For all $s,t\in\R$,
\[
  (\cos 2t-1)\sin 2s-\sin 2t\,(\cos 2s-1)=4\sin t\,\sin s\,\sin(s-t).
\]
\end{lemma}

\begin{proof}
By the double-angle formulas, $\cos2t-1=-2\sin^2t$, $\cos2s-1=-2\sin^2s$, $\sin2s=2\sin s\cos s$ and
$\sin2t=2\sin t\cos t$. Hence the left side equals
\[
  -4\sin^2t\,\sin s\cos s+4\sin t\cos t\,\sin^2s=4\sin t\,\sin s\,(\sin s\cos t-\cos s\sin t)=4\sin t\,\sin s\,\sin(s-t).
\]
(Lean rewrites with $\cos2x=2\cos^2x-1$ and $\sin2x=2\sin x\cos x$ for $x=s,t$ and with the subtraction formula for
$\sin(s-t)$, and closes the remaining polynomial identity by \lean{linear_combination} with $\sin^2+\cos^2=1$ for $s$
and for $t$.)
\end{proof}

\begin{lemma}[vertices lie weakly left of each edge; \leanref{ErdosSzekeres.LayerLemma.orient_regular_nonneg}{LAYERLEAN/LayerLemma.lean:488}]\label{lg:lem:regular}
Let $c,x$ be points, $m\in\N$ with $0<m$, $j\in\N$ with $j\le m$, and $\rho=\operatorname{rotAbout}(c,2\pi/m)$. Then
\[
  0\le\orient\bigl(x,\rho(x),\rho^{j}(x)\bigr).
\]
\end{lemma}

\begin{proof}
By Lemma~\ref{lg:lem:rot-iter}, $\rho^{j}(x)=\operatorname{rotAbout}(c,\beta)(x)$ with $\beta=j\cdot(2\pi/m)$. By
Lemma~\ref{lg:lem:orient-rot} with $\alpha=2\pi/m$,
\[
  \orient(x,\rho(x),\rho^{j}(x))=\abs{x-c}^2\,F,\qquad F=(\cos\alpha-1)\sin\beta-\sin\alpha\,(\cos\beta-1).
\]
The first factor is $\ge0$; we show $F\ge0$. If $j=0$ then $\beta=0$, so $\sin\beta=0=\cos\beta-1$ and $F=0$. Let
$1\le j\le m$ and put $t=\pi/m$. Then $\alpha=2t$ and $\beta=2(jt)$, so Lemma~\ref{lg:lem:trig} with $s=jt$ gives
\[
  F=4\sin t\,\sin(jt)\,\sin(jt-t).
\]
Since $m\ge1$, $0<t\le\pi$; since $j\le m$, $0\le jt=j\pi/m\le\pi$; and since $j\ge1$, $0\le jt-t=(j-1)t\le jt\le\pi$.
By (S1) the sine is $\ge0$ on $[0,\pi]$, so all three factors are $\ge0$, and $F\ge0$.
\end{proof}

\begin{remark}
The proof shows $\orient(x,\rho(x),\rho^{j}(x))=4\abs{x-c}^2\sin(\pi/m)\sin(j\pi/m)\sin((j-1)\pi/m)$ for
$1\le j\le m$: every vertex of the regular $m$-gon $\{\rho^{j}(x)\}$ lies weakly to the left of the directed edge
$x\to\rho(x)$. (If moreover $x\ne c$ and $m\ge2$, then by (S1) $\sin(\pi/m)>0$, $\sin(j\pi/m)=0$ only for $j=m$ and
$\sin((j-1)\pi/m)=0$ only for $j=1$, so the value is $0$ exactly for $j\in\{1,m\}$; this refinement is not used.)
\end{remark}

\begin{definition}[lin; \leanref{ErdosSzekeres.LayerLemma.lin}{LAYERLEAN/LayerLemma.lean:625}]\label{lg:def:lin}
For points $w,a,y$,
\[
  \operatorname{lin}(w,a,y)=w_0(y_0-a_0)+w_1(y_1-a_1)=\langle w,\,y-a\rangle .
\]
\end{definition}

For $w\ne0$, $\{y:\operatorname{lin}(w,a,y)>0\}$ is the open half-plane bounded by the line through $a$ orthogonal to
$w$, on the side into which $w$ points. Note $\operatorname{lin}(w,a,a)=0$.

\begin{lemma}[the tangent functional on the orbit; \leanref{ErdosSzekeres.LayerLemma.lin_rotAbout}{LAYERLEAN/LayerLemma.lean:794}]\label{lg:lem:lin-rot}
For all points $c,p$ and every $\beta\in\R$,
\[
  \operatorname{lin}\bigl(c-p,\;p,\;\operatorname{rotAbout}(c,\beta)(p)\bigr)=\bigl((p_0-c_0)^2+(p_1-c_1)^2\bigr)(1-\cos\beta),
\]
where $c-p$ is the coordinatewise difference.
\end{lemma}

\begin{proof}
Put $u=p-c$; then $c-p=-u$ and $\operatorname{rotAbout}(c,\beta)(p)-p=R_\beta u-u=(\cos\beta-1)u+\sin\beta\,Ju$. Hence,
by~\eqref{lg:eq:J},
\[
  \operatorname{lin}(c-p,p,\operatorname{rotAbout}(c,\beta)(p))=\langle-u,(\cos\beta-1)u+\sin\beta\,Ju\rangle
  =(1-\cos\beta)\abs{u}^2-\sin\beta\,\langle u,Ju\rangle=(1-\cos\beta)\abs{u}^2.\qedhere
\]
\end{proof}

\begin{lemma}[other vertices lie strictly inside the tangent; \leanref{ErdosSzekeres.LayerLemma.lin_orbit_pos}{LAYERLEAN/LayerLemma.lean:801}]\label{lg:lem:lin-orbit}
Let $c,p$ be points with $p\ne c$, let $m\in\N$ with $0<m$, let $j\in\N$ with $0<j$ and $j<m$, and let
$\rho=\operatorname{rotAbout}(c,2\pi/m)$. Then
\[
  0<\operatorname{lin}\bigl(c-p,\;p,\;\rho^{j}(p)\bigr).
\]
\end{lemma}

\begin{proof}
By Lemma~\ref{lg:lem:rot-iter}, $\rho^{j}(p)=\operatorname{rotAbout}(c,\theta)(p)$ with $\theta=j\cdot(2\pi/m)$, so by
Lemma~\ref{lg:lem:lin-rot} the value is $\abs{p-c}^2(1-\cos\theta)$. Since $p\ne c$, one of $p_0-c_0$, $p_1-c_1$ is
non-zero, so $\abs{p-c}^2>0$. Since $0<j<m$, $0<\theta<2\pi$. By (S2), $\cos\theta\le1$, and $\cos\theta=1$ only for
$\theta\in2\pi\Z$, which has no point in $(0,2\pi)$; hence $1-\cos\theta>0$, and the product is positive.
\end{proof}

\begin{lemma}[next and previous vertex; \leanref{ErdosSzekeres.LayerLemma.orient_next_prev_pos}{LAYERLEAN/LayerLemma.lean:830}]\label{lg:lem:next-prev}
Let $c,p$ be points with $p\ne c$, let $m\in\N$ with $3\le m$, and let $\rho=\operatorname{rotAbout}(c,2\pi/m)$. Then
\[
  0<\orient\bigl(p,\rho(p),\rho^{m-1}(p)\bigr).
\]
\end{lemma}

\begin{proof}
As $m\ge1$, the natural number $m-1$ has real cast $(m:\R)-1$. By Lemma~\ref{lg:lem:rot-iter},
$\rho^{m-1}(p)=\operatorname{rotAbout}(c,\beta)(p)$ with $\beta=(m-1)\cdot(2\pi/m)$. By Lemma~\ref{lg:lem:orient-rot}
with $\alpha=2\pi/m$, the orientation equals $\abs{p-c}^2F$ with $F$ as in the proof of Lemma~\ref{lg:lem:regular}.
Put $t=\pi/m$, so $mt=\pi$, $\alpha=2t$ and $\beta=2((m-1)t)$. Lemma~\ref{lg:lem:trig} with $s=(m-1)t$ gives
\[
  F=4\sin t\,\sin((m-1)t)\,\sin((m-1)t-t)=4\sin t\,\sin(\pi-t)\,\sin(\pi-2t).
\]
Since $m\ge3$, $0<t\le\pi/3$; hence $t$, $\pi-t$ and $\pi-2t$ (which is $\ge\pi/3$) all lie in the open interval
$(0,\pi)$, where the sine is positive by (S1). So $F>0$; also $\abs{p-c}^2>0$ as $p\ne c$. The product is positive.
\end{proof}

\subsection{Convexity of half-planes and disks}\label{lg:ssec:convex}

\begin{lemma}[$\{\orient(a,b,\cdot)\ge0\}$ is convex; \leanref{ErdosSzekeres.LayerLemma.convex_orient_nonneg}{LAYERLEAN/LayerLemma.lean:516}]\label{lg:lem:halfplane-convex}
For all points $a,b$ the set $\{q\in\R^2:0\le\orient(a,b,q)\}$ is convex.
\end{lemma}

\begin{proof}
Let $q_1,q_2$ be in the set and $u,v\ge0$ with $u+v=1$. By (O4),
$\orient(a,b,uq_1+vq_2)=u\orient(a,b,q_1)+v\orient(a,b,q_2)$, a sum of two products of non-negative numbers.
\end{proof}

For $a\ne b$ this set is the closed half-plane to the left of the directed line $a\to b$, the line included; for
$a=b$ it is the whole plane, since $\orient(a,a,q)=0$ by (O3).

\begin{definition}[dist2; \leanref{ErdosSzekeres.LayerLemma.dist2}{LAYERLEAN/LayerLemma.lean:999}]\label{lg:def:dist2}
For points $c,y$, $\operatorname{dist2}(c,y)=(y_0-c_0)^2+(y_1-c_1)^2=\abs{y-c}^2$, the squared Euclidean distance
(defined by this formula, not through a norm).
\end{definition}

\begin{lemma}[disks; \leanref{ErdosSzekeres.LayerLemma.convex_disk}{LAYERLEAN/LayerLemma.lean:1002}]\label{lg:lem:disk}
For every point $c$ and every $R\in\R$ the set $\{y\in\R^2:\operatorname{dist2}(c,y)\le R\}$ is convex.
\end{lemma}

\begin{proof}
Let $y_1,y_2$ be in the set and $u,v\ge0$ with $u+v=1$. Put $A=y_1-c$ and $B=y_2-c$. Since $u+v=1$,
$(uy_1+vy_2)-c=uA+vB$, and, using $1-u=v$ and $1-v=u$,
\[
  u\abs{A}^2+v\abs{B}^2-\abs{uA+vB}^2=u(1-u)\abs{A}^2+v(1-v)\abs{B}^2-2uv\langle A,B\rangle=uv\,\abs{A-B}^2\ge0.
\]
Hence $\operatorname{dist2}(c,uy_1+vy_2)=\abs{uA+vB}^2\le u\abs{A}^2+v\abs{B}^2\le uR+vR=R$. (Lean proves the same
inequality with \lean{nlinarith}, from the non-negativity of the squares of the coordinate differences of $y_1,y_2$ and
of $uv$.)
\end{proof}

For $R\ge0$ the set is the closed disk of radius $\sqrt R$ about $c$; for $R<0$ it is empty.

\subsection{Strict nesting, the centre, and the radius order}\label{lg:ssec:nest}

\begin{theorem}[strict nesting of the convex layers; \leanref{ErdosSzekeres.LayerLemma.strict_nesting}{LAYERLEAN/LayerLemma.lean:529}]\label{lg:thm:nesting}
Let $S$ be a finite set of points, $c$ a point, and $m,K\in\N$ with $2\le m$ and $K\le 2m$; let
$\rho=\operatorname{rotAbout}(c,2\pi/m)$. Assume that $c\notin S$, that $\rho$ maps $S$ into $S$, that $S$ is $K$-free,
and that $S$ is in general position. Let $n,n'\in\N$ with $n<n'$, let $x\in\hv(\peel(S,n))$ and $y\in\peel(S,n')$. Then
\[
  0<\orient\bigl(x,\rho(x),y\bigr).
\]
\end{theorem}

\begin{proof}
Since $m\ge2$, also $0<m$, so Lemma~\ref{lg:lem:layer} applies to $n$ and $x$. Put $L=\hv(\peel(S,n))$. By
Lemma~\ref{lg:lem:layer}(a), $L=\{\rho^{j}(x):0\le j<m\}$; by (c), applied with $n'\ge n$, $y\in\conv(L)$; by (d),
applied with $n'>n$, $y\notin L$.

\emph{Weak inequality.} Let $H=\{q:0\le\orient(x,\rho(x),q)\}$. Every element of $L$ is $\rho^{j}(x)$ with $j<m$, in
particular $j\le m$, so it lies in $H$ by Lemma~\ref{lg:lem:regular}. $H$ is convex by
Lemma~\ref{lg:lem:halfplane-convex} (with $a=x$, $b=\rho(x)$). Hence $\conv(L)\subseteq H$, and $y\in H$, that is,
$0\le\orient(x,\rho(x),y)$.

\emph{Strict inequality.} As $1<m$, $\rho(x)=\rho^{1}(x)\in L$ by Lemma~\ref{lg:lem:layer}(a). The three points
$x,\rho(x),y$ lie in $S$: $x,\rho(x)\in L\subseteq S$ by Lemma~\ref{lg:lem:layer}(b), and $y\in\peel(S,n')\subseteq S$
by (L2). They are pairwise distinct:
\begin{itemize}
\item $x\ne\rho(x)$: $x\ne c$ by Lemma~\ref{lg:lem:layer}(b), so Lemma~\ref{lg:lem:lin-orbit} with $p=x$ and $j=1$
  ($0<1<m$) gives $0<\operatorname{lin}(c-x,x,\rho(x))$, whereas $\operatorname{lin}(c-x,x,x)=0$;
\item $y\ne x$ and $y\ne\rho(x)$, because $x,\rho(x)\in L$ and $y\notin L$.
\end{itemize}
By general position and Lemma~\ref{fd:lem:collinear}, $\orient(x,\rho(x),y)\ne0$. Together with the weak inequality,
$0<\orient(x,\rho(x),y)$.
\end{proof}

\begin{remark}
(1) Lean proves $x\ne\rho(x)$ differently: $x\ne c$ has minimal period $m\ge2$ under $\rho$
(Lemma~\ref{ll:lem:minperiod}), while a fixed point of $\rho$ has minimal period $1$.
(2) For $m=2$ the hypotheses of Theorem~\ref{lg:thm:nesting} are never all satisfied. Indeed, then
$L=\{x,\rho(x)\}$ and $y\in\conv\{x,\rho(x)\}$. The segment $\{(1-\lambda)x+\lambda\rho(x):0\le\lambda\le1\}$ contains
$L$ and is convex (a convex combination of the points with parameters $\lambda_1,\lambda_2$ is the point with the
parameter given by the same combination of $\lambda_1,\lambda_2$), so it contains $\conv\{x,\rho(x)\}$, and
$y=(1-\lambda)x+\lambda\rho(x)$ with $0\le\lambda\le1$. By (O4)
and (O3) $\orient(x,\rho(x),y)=(1-\lambda)\orient(x,\rho(x),x)+\lambda\orient(x,\rho(x),\rho(x))=0$, contradicting the
conclusion. So for $m=2$ the theorem says that, under its hypotheses on $S$, $\peel(S,n')=\emptyset$ for all $n'>n$ as
soon as $\hv(\peel(S,n))\ne\emptyset$.
(3) The theorem is the orientation statement behind the nesting clauses of the layered encoders. That the encoders
emit exactly these clauses is not part of \dirn{LAYERLEAN/LayerLemma.lean}; see the paragraph \emph{Encoders} at the
beginning of this section.
\end{remark}

\begin{lemma}[even orders exclude the centre; \leanref{ErdosSzekeres.LayerLemma.centre_not_mem_of_even}{LAYERLEAN/LayerLemma.lean:582}]\label{lg:lem:centre}
Let $S$ be a finite set, $c$ a point, and $m\in\N$ with $0<m$ and $m$ even; let $\rho=\operatorname{rotAbout}(c,2\pi/m)$.
Assume that $\rho$ maps $S$ into $S$ and that $S$ is in general position. Let $p$ be a point with $p\in S$ and
$p\ne c$. Then $c\notin S$.
\end{lemma}

\begin{proof}
Suppose $c\in S$. Since $m$ is even, $m=k+k$ for some $k\in\N$ (Mathlib's \lean{Even}); $k\ge1$ because $m>0$, and the
natural-number quotient $m/2$ equals $k$. Let $q=\rho^{k}(p)$ (Lean: $\rho^{m/2}(p)$). Since $\rho$ maps $S$ into $S$,
$q\in S$. By Corollary~\ref{lg:cor:period}, $\rho^{k}=\operatorname{rotAbout}(c,2\pi k/m)$, and
$2\pi k/m=2\pi k/(2k)=\pi$. As $\cos\pi=-1$ and $\sin\pi=0$,
\[
  q=\operatorname{rotAbout}(c,\pi)(p)=\bigl(c_0-(p_0-c_0),\;c_1-(p_1-c_1)\bigr)=2c-p .
\]
By (O1), $\orient(p,c,q)=\det(c-p,\,2c-2p)=2\det(c-p,\,c-p)=0$. The points $p,c,q$ are pairwise distinct: $p\ne c$ by
hypothesis; $q=c$ would mean $2c-p=c$, i.e.\ $p=c$; and $p=q$ would mean $p=2c-p$, i.e.\ $p=c$. Since $p,c,q\in S$,
general position and Lemma~\ref{fd:lem:collinear} give $\orient(p,c,q)\ne0$, a contradiction. Hence $c\notin S$.
\end{proof}

\begin{remark}
There is no $K$-freeness hypothesis and no hypothesis $c\notin S$. The hypothesis that $S$ has a point $p\neq c$ cannot be
dropped: $S=\{c\}$ is mapped into itself by $\rho$, is in general position (it has no three distinct points), and
contains $c$.
\end{remark}

\begin{theorem}[the radius order of the layers; \leanref{ErdosSzekeres.LayerLemma.dist2_le_of_later}{LAYERLEAN/LayerLemma.lean:1012}]\label{lg:thm:radius}
Let $S$ be a finite set, $c$ a point, and $m,K\in\N$ with $0<m$ and $K\le 2m$; let $\rho=\operatorname{rotAbout}(c,2\pi/m)$.
Assume that $c\notin S$, that $\rho$ maps $S$ into $S$, and that $S$ is $K$-free. Let $n,n'\in\N$ with $n\le n'$, let
$x\in\hv(\peel(S,n))$ and $y\in\peel(S,n')$. Then
\[
  \operatorname{dist2}(c,y)\le\operatorname{dist2}(c,x).
\]
\end{theorem}

\begin{proof}
Put $L=\hv(\peel(S,n))$. By Lemma~\ref{lg:lem:layer}(a), $L=\{\rho^{j}(x):0\le j<m\}$. For each $j$,
$\rho^{j}(x)=\operatorname{rotAbout}(c,2\pi j/m)(x)$ by Corollary~\ref{lg:cor:period}, so
$\operatorname{dist2}(c,\rho^{j}(x))=\operatorname{dist2}(c,x)$ by Lemma~\ref{lg:lem:rotfacts}(d). Hence $L$ is
contained in $D=\{z:\operatorname{dist2}(c,z)\le\operatorname{dist2}(c,x)\}$, which is convex by Lemma~\ref{lg:lem:disk}
(with $R=\operatorname{dist2}(c,x)$). So $\conv(L)\subseteq D$. By Lemma~\ref{lg:lem:layer}(c) (as $n'\ge n$),
$y\in\conv(L)\subseteq D$.
\end{proof}

\begin{remark}
Only the weak inequality is stated, for $n\le n'$ (including $n=n'$), and general position is not assumed. Strictly
decreasing radii are not claimed here. Rule~B (Lemma~\ref{lg:lem:ruleB}) uses only weak inequalities between radii.
\end{remark}

\subsection{Fans and the chain clauses}\label{lg:ssec:chain}

The two orientation identities behind the fan lemma, and the identity behind Lemma~\ref{lg:lem:halfplane}, are all
instances of one vector identity in the plane.

\begin{lemma}[a Cramer identity; written step]\label{lg:lem:cramer}
For all $X,Y,Z\in\R^2$,
\[
  \det(X,Y)\,Z+\det(Y,Z)\,X+\det(Z,X)\,Y=0 .
\]
Consequently, for all $w,V\in\R^2$,
\begin{align*}
  \det(X,Y)\langle w,Z\rangle+\det(Y,Z)\langle w,X\rangle+\det(Z,X)\langle w,Y\rangle&=0,\\
  \det(X,Y)\det(V,Z)+\det(Y,Z)\det(V,X)+\det(Z,X)\det(V,Y)&=0.
\end{align*}
\end{lemma}

\begin{proof}
Fix $k\in\{0,1\}$ and consider the $3\times3$ matrix with rows $(X_k,Y_k,Z_k)$, $(X_0,Y_0,Z_0)$, $(X_1,Y_1,Z_1)$. Two of
its rows are equal, so its determinant vanishes. Expanding along the first row gives
$X_k\det(Y,Z)-Y_k\det(X,Z)+Z_k\det(X,Y)=0$, which is the $k$-th coordinate of the first identity because
$-\det(X,Z)=\det(Z,X)$. The other two identities follow by applying the linear maps $v\mapsto\langle w,v\rangle$ and
$v\mapsto\det(V,v)$.
\end{proof}

\begin{lemma}[transitivity in a half-plane; \leanref{ErdosSzekeres.LayerLemma.orient_trans_halfplane}{LAYERLEAN/LayerLemma.lean:629}]\label{lg:lem:trans}
Let $a,b,c,e,w$ be points (here $c$ is an arbitrary point, not a centre). If
\[
  0<\operatorname{lin}(w,a,b),\quad 0<\operatorname{lin}(w,a,c),\quad 0<\operatorname{lin}(w,a,e),\quad
  0<\orient(a,b,c),\quad 0<\orient(a,c,e),
\]
then $0<\orient(a,b,e)$.
\end{lemma}

\begin{proof}
Put $B=b-a$, $C=c-a$, $E=e-a$. By (O1) and Definition~\ref{lg:def:lin}, $\orient(a,b,c)=\det(B,C)$,
$\orient(a,c,e)=\det(C,E)$, $\orient(a,b,e)=\det(B,E)$, and $\operatorname{lin}(w,a,z)=\langle w,z-a\rangle$. The
second identity of Lemma~\ref{lg:lem:cramer} with $(X,Y,Z)=(B,C,E)$ reads
$\det(B,C)\langle w,E\rangle+\det(C,E)\langle w,B\rangle+\det(E,B)\langle w,C\rangle=0$. Since $\det(E,B)=-\det(B,E)$,
this is
\begin{equation}\label{lg:eq:trans}
  \orient(a,b,e)\operatorname{lin}(w,a,c)=\orient(a,b,c)\operatorname{lin}(w,a,e)+\orient(a,c,e)\operatorname{lin}(w,a,b).
\end{equation}
The right side is a sum of two products of positive numbers, hence positive. As $\operatorname{lin}(w,a,c)>0$, it
follows that $\orient(a,b,e)>0$. (Lean checks identity~\eqref{lg:eq:trans} directly by \lean{ring}, without
Lemma~\ref{lg:lem:cramer}.)
\end{proof}

In words: inside an open half-plane bounded by a line through $a$, the counterclockwise angular order around $a$ is
transitive.

\begin{lemma}[fan step; \leanref{ErdosSzekeres.LayerLemma.fan_step}{LAYERLEAN/LayerLemma.lean:641}]\label{lg:lem:fan-step}
Let $p,a,b,c,e$ be points (here $c$ is an arbitrary point, not a centre). If
\[
  0<\orient(b,c,e),\quad 0<\orient(p,c,e),\quad 0<\orient(p,b,c),\quad 0<\orient(a,b,c),\quad 0<\orient(p,a,c),
\]
then $0<\orient(a,c,e)$.
\end{lemma}

\begin{proof}
Put $P=p-c$, $A=a-c$, $B=b-c$, $E=e-c$. By (O2) (cyclic invariance, and $\orient(x,y,c)=\det(x-c,y-c)$ for all
$x,y$),
\begin{gather*}
  \orient(p,b,c)=\det(P,B),\qquad \orient(a,b,c)=\det(A,B),\qquad \orient(p,a,c)=\det(P,A),\\
  \orient(a,c,e)=\orient(e,a,c)=\det(E,A),\qquad \orient(p,c,e)=\orient(e,p,c)=\det(E,P),\\
  \orient(b,c,e)=\orient(e,b,c)=\det(E,B).
\end{gather*}
The third identity of Lemma~\ref{lg:lem:cramer} with $(X,Y,Z,V)=(P,B,A,E)$ reads
\[
  \det(P,B)\det(E,A)+\det(B,A)\det(E,P)+\det(A,P)\det(E,B)=0.
\]
Since $\det(B,A)=-\det(A,B)$ and $\det(A,P)=-\det(P,A)$, this is
\begin{equation}\label{lg:eq:fanstep}
  \orient(p,b,c)\orient(a,c,e)=\orient(p,c,e)\orient(a,b,c)+\orient(b,c,e)\orient(p,a,c).
\end{equation}
By the hypotheses the right side is positive, and $\orient(p,b,c)>0$; hence $\orient(a,c,e)>0$. (Lean checks
identity~\eqref{lg:eq:fanstep} directly by \lean{ring}, without Lemma~\ref{lg:lem:cramer}.)
\end{proof}

\begin{lemma}[orientation as a linear functional; \leanref{ErdosSzekeres.LayerLemma.orient_eq_lin}{LAYERLEAN/LayerLemma.lean:653}]\label{lg:lem:orient-lin}
For all points $p,x,y$,
\[
  \orient(p,x,y)=\operatorname{lin}\bigl((-(x_1-p_1),\;x_0-p_0),\;x,\;y\bigr).
\]
\end{lemma}

\begin{proof}
By (O1), $\orient(p,x,y)=\det(x-p,\,y-p)=\det(x-p,\,y-x)+\det(x-p,\,x-p)=\det(x-p,\,y-x)=\langle J(x-p),\,y-x\rangle$,
and $J(x-p)=(-(x_1-p_1),\,x_0-p_0)$.
\end{proof}

So, for $x\ne p$, $\{y:\orient(p,x,y)>0\}$ is the open half-plane to the left of the directed line $p\to x$; its
bounding line passes through $x$.

\begin{lemma}[fan chains have all triples positive; \leanref{ErdosSzekeres.LayerLemma.fan_positive}{LAYERLEAN/LayerLemma.lean:658}]\label{lg:lem:fan-pos}
Let $p,w$ be points, $q:\N\to\R^2$ a sequence of points (write $q_s=q(s)$), and $n\in\N$. Assume
\begin{itemize}
\item[(H)] $0<\operatorname{lin}(w,p,q_s)$ for all $s<n$;
\item[(F)] $0<\orient(p,q_s,q_{s+1})$ for all $s$ with $s+1<n$;
\item[(T)] $0<\orient(q_s,q_{s+1},q_{s+2})$ for all $s$ with $s+2<n$.
\end{itemize}
Then (a) $0<\orient(p,q_j,q_k)$ for all $j,k$ with $j<k$ and $k<n$; and (b) $0<\orient(q_i,q_j,q_k)$ for all $i,j,k$
with $i<j$, $j<k$ and $k<n$.
\end{lemma}

\begin{proof}
(a) Fix $j$. We show by induction on $k\ge j+1$ that $k<n$ implies $\orient(p,q_j,q_k)>0$. For $k=j+1$ this is (F)
with $s=j$. Let $k\ge j+1$ satisfy the claim and let $k+1<n$. Then $k<n$, so $\orient(p,q_j,q_k)>0$; (F) with $s=k$
gives $\orient(p,q_k,q_{k+1})>0$; and (H) gives $\operatorname{lin}(w,p,q_t)>0$ for $t\in\{j,k,k+1\}$, all of which
are $<n$. Lemma~\ref{lg:lem:trans} with $(a,b,c,e)=(p,q_j,q_k,q_{k+1})$ gives $\orient(p,q_j,q_{k+1})>0$.

(b$_0$) We first show: for all $i<s$ with $s+1<n$, $\orient(q_i,q_s,q_{s+1})>0$. Fix $i$ and induct on $s\ge i+1$. For
$s=i+1$ this is (T) with $s=i$ (note $i+2<n$). Let $s\ge i+1$ satisfy the claim and let $s+2<n$. Apply
Lemma~\ref{lg:lem:fan-step} with $(p,a,b,c,e)=(p,q_i,q_s,q_{s+1},q_{s+2})$; its five hypotheses hold:
$\orient(q_s,q_{s+1},q_{s+2})>0$ by (T); $\orient(p,q_{s+1},q_{s+2})>0$ by (F) with $s+1$; $\orient(p,q_s,q_{s+1})>0$ by
(F) with $s$; $\orient(q_i,q_s,q_{s+1})>0$ by the induction hypothesis (as $s+1<n$); and $\orient(p,q_i,q_{s+1})>0$ by
(a), since $i<s+1<n$. The conclusion is $\orient(q_i,q_{s+1},q_{s+2})>0$.

(b) Fix $i<j$ and induct on $k\ge j+1$, showing that $k<n$ implies $\orient(q_i,q_j,q_k)>0$. For $k=j+1$ this is
(b$_0$) with $s=j$. Let $k\ge j+1$ satisfy the claim and let $k+1<n$. Put $w'=(-(q_{i,1}-p_1),\,q_{i,0}-p_0)$, where
$q_i=(q_{i,0},q_{i,1})$. By Lemma~\ref{lg:lem:orient-lin}, $\operatorname{lin}(w',q_i,y)=\orient(p,q_i,y)$ for every $y$,
so by (a) $\operatorname{lin}(w',q_i,q_t)>0$ for every $t$ with $i<t<n$, in particular for $t\in\{j,k,k+1\}$. Moreover
$\orient(q_i,q_j,q_k)>0$ by the induction hypothesis (as $k<n$), and $\orient(q_i,q_k,q_{k+1})>0$ by (b$_0$) with
$s=k$ (as $i<k$ and $k+1<n$). Lemma~\ref{lg:lem:trans} with $(a,b,c,e,w)=(q_i,q_j,q_k,q_{k+1},w')$ gives
$\orient(q_i,q_j,q_{k+1})>0$.
\end{proof}

\begin{definition}[fan points; \leanref{ErdosSzekeres.LayerLemma.fanPt}{LAYERLEAN/LayerLemma.lean:694}]\label{lg:def:fanpt}
For $n\in\N$, a point $p$, $q:\N\to\R^2$ and $i\in\{0,1,\dots,n\}$ (Lean \texttt{Fin (n + 1)}; $n$ is an implicit
argument),
\[
  \operatorname{fanPt}(p,q)(i)=\begin{cases} p & \text{if } i=0,\\ q(i-1) & \text{otherwise.}\end{cases}
\]
\end{definition}

\begin{theorem}[fan chains give convex polygons; \leanref{ErdosSzekeres.LayerLemma.hasConvexNGon_of_fan}{LAYERLEAN/LayerLemma.lean:700}]\label{lg:thm:fan}
Let $S$ be a finite set, $p,w$ points, $q:\N\to\R^2$ (write $q_s=q(s)$) and $n\in\N$ with $3\le n$. Assume that $p\in S$,
that $q_s\in S$ for all $s<n$, and that (H), (F) and (T) of Lemma~\ref{lg:lem:fan-pos} hold:
$0<\operatorname{lin}(w,p,q_s)$ for all $s<n$; $0<\orient(p,q_s,q_{s+1})$ for all $s$ with $s+1<n$; and
$0<\orient(q_s,q_{s+1},q_{s+2})$ for all $s$ with $s+2<n$. Then \texttt{HasConvexNGon (n + 1) S}: some $(n+1)$-element
subset of $S$ is in convex position.
\end{theorem}

\begin{proof}
Let $x_i=\operatorname{fanPt}(p,q)(i)$ for $i\in\{0,\dots,n\}$; so $x_0=p$ and $x_i=q_{i-1}$ for $1\le i\le n$.

\emph{Step 1: every increasing triple is positive.} Let $0\le i<j<k\le n$. If $i=0$, then $j,k\ge1$ and
$\orient(x_0,x_j,x_k)=\orient(p,q_{j-1},q_{k-1})>0$ by Lemma~\ref{lg:lem:fan-pos}(a), since $j-1<k-1<n$. If $i\ge1$,
then $\orient(x_i,x_j,x_k)=\orient(q_{i-1},q_{j-1},q_{k-1})>0$ by Lemma~\ref{lg:lem:fan-pos}(b), since
$i-1<j-1<k-1<n$.

\emph{Step 2: the points $x_0,\dots,x_n$ are pairwise distinct.} Suppose $x_i=x_j$ with $i\ne j$. As $n+1\ge4>2$, there
is $k\in\{0,\dots,n\}$ different from $i$ and $j$. Arrange $i,j,k$ increasingly as $i'<j'<k'$. By Step~1,
$\orient(x_{i'},x_{j'},x_{k'})>0$; but two of these three points are equal (namely $x_i$ and $x_j$), so the orientation
is $0$ by (O3), a contradiction. (Lean chooses $k$ explicitly and treats the possible orders of $i,j,k$ one by one,
using \lean{orient_cycle} and $\orient(a,a,d)=0$; the argument is the same.)

\emph{Step 3: convex position.} Apply Lemma~\ref{fd:lem:pos-triples} (\lean{Suk.convexPosition_of_positive_triples})
with the linearly ordered index set $\{0,\dots,n\}$ (Lean \texttt{Fin (n + 1)} with its usual order), the injective map
$i\mapsto x_i$ (Step~2), and the full index set, which has $n+1\ge4$ elements; its hypothesis is Step~1. Hence
$T=\{x_0,\dots,x_n\}$ is in convex position.

\emph{Step 4.} $T\subseteq S$, since $x_0=p\in S$ and $x_i=q_{i-1}\in S$ for $1\le i\le n$ (as $i-1<n$); and
$\card T=n+1$ by Step~2. Hence \texttt{HasConvexNGon (n + 1) S}.
\end{proof}

\begin{remark}
General position is not assumed, and the chain points are not assumed distinct: distinctness of $p,q_0,\dots,q_{n-1}$
follows from the positivity of the orientations (Step~2), and the proof shows that $\{p,q_0,\dots,q_{n-1}\}$ itself is
the convex $(n+1)$-gon.
\end{remark}

\begin{lemma}[inner side of the tangent in the layers; \leanref{ErdosSzekeres.LayerLemma.lin_pos_of_peel}{LAYERLEAN/LayerLemma.lean:864}]\label{lg:lem:halfplane}
Let $S$ be a finite set, $c$ a point, and $m,K\in\N$ with $3\le m$ and $K\le 2m$; let $\rho=\operatorname{rotAbout}(c,2\pi/m)$.
Assume that $c\notin S$, that $\rho$ maps $S$ into $S$, that $S$ is $K$-free, and that $S$ is in general position. Let
$i\in\N$, $p\in\hv(\peel(S,i))$, and $y\in\peel(S,i)$ with $y\ne p$. Then
\[
  0<\operatorname{lin}(c-p,\;p,\;y).
\]
\end{lemma}

\begin{proof}
Since $m\ge3$, $0<m$ and Lemma~\ref{lg:lem:layer} applies to $i$ and $p$. Put $L=\hv(\peel(S,i))$; then
$L=\{\rho^{j}(p):0\le j<m\}$ and $p\ne c$.

\emph{Case $y\in L$.} Then $y=\rho^{j}(p)$ for some $j<m$, and $j\ne0$ because $\rho^{0}(p)=p\ne y$. So $0<j<m$, and
Lemma~\ref{lg:lem:lin-orbit} gives $0<\operatorname{lin}(c-p,p,y)$.

\emph{Case $y\notin L$.} Then $y\in\peel(S,i)\setminus L=\peel(S,i+1)$. Let $p'=\rho^{m-1}(p)$. Since $m-1<m$,
$p'\in L$ by Lemma~\ref{lg:lem:layer}(a); and $\rho(p')=\rho^{m}(p)=p$ by Corollary~\ref{lg:cor:period}. We collect
five positive quantities.
\begin{itemize}
\item[(e1)] $\orient(p,\rho(p),y)>0$: Theorem~\ref{lg:thm:nesting} with $(n,n',x)=(i,i+1,p)$ (its hypotheses hold:
  $2\le m$, $K\le2m$, $c\notin S$, invariance, $K$-freeness, general position, $p\in\hv(\peel(S,i))$,
  $y\in\peel(S,i+1)$).
\item[(e2)] $\orient(p',p,y)>0$: Theorem~\ref{lg:thm:nesting} with $(n,n',x)=(i,i+1,p')$ gives
  $\orient(p',\rho(p'),y)>0$, and $\rho(p')=p$.
\item[(w1)] $\operatorname{lin}(c-p,p,\rho(p))>0$: Lemma~\ref{lg:lem:lin-orbit} with $j=1$ (as $0<1<m$).
\item[(w2)] $\operatorname{lin}(c-p,p,p')>0$: Lemma~\ref{lg:lem:lin-orbit} with $j=m-1$ (as $0<m-1<m$).
\item[(o)] $\orient(p,\rho(p),p')>0$: Lemma~\ref{lg:lem:next-prev} (as $p\ne c$ and $m\ge3$).
\end{itemize}
Now put $B=\rho(p)-p$, $C=p'-p$, $E=y-p$ and $w=c-p$. By (O1) and (O2),
$\orient(p,\rho(p),p')=\det(B,C)$, $\orient(p,\rho(p),y)=\det(B,E)$ and $\orient(p',p,y)=\orient(p,y,p')=\det(E,C)$;
and $\operatorname{lin}(w,p,z)=\langle w,z-p\rangle$. The second identity of Lemma~\ref{lg:lem:cramer} with
$(X,Y,Z)=(B,C,E)$ gives $\det(B,C)\langle w,E\rangle=-\det(C,E)\langle w,B\rangle-\det(E,B)\langle w,C\rangle
=\det(E,C)\langle w,B\rangle+\det(B,E)\langle w,C\rangle$, that is,
\begin{equation}\label{lg:eq:key}
  \orient(p,\rho(p),p')\operatorname{lin}(c-p,p,y)
  =\orient(p',p,y)\operatorname{lin}(c-p,p,\rho(p))+\orient(p,\rho(p),y)\operatorname{lin}(c-p,p,p').
\end{equation}
By (e2), (w1), (e1) and (w2) the right side is positive, and by (o) $\orient(p,\rho(p),p')>0$. Hence
$\operatorname{lin}(c-p,p,y)>0$.
\end{proof}

\begin{remark}
(1) Geometrically, $\{y:\operatorname{lin}(c-p,p,y)=0\}$ is the tangent at $p$ to the circle about $c$ through $p$, and
the lemma says that every other point of the layers $\ge i$ lies strictly on the side of the centre. Identity
\eqref{lg:eq:key} holds for arbitrary points $p,\rho(p),p',y,c$; Lean proves it by \lean{ring}.
(2) Lean obtains $\rho(p')=\rho^{m}(p)=p$ from the minimal period of $p$ (Lemma~\ref{ll:lem:minperiod}); we used
Corollary~\ref{lg:cor:period} instead.
\end{remark}

\begin{theorem}[geometric chain soundness; \leanref{ErdosSzekeres.LayerLemma.chain_clauses_sound}{LAYERLEAN/LayerLemma.lean:914}]\label{lg:thm:chain}
Let $S$ be a finite set, $c$ a point, and $m,K\in\N$ with $3\le m$, $4\le K$ and $K\le 2m$; let
$\rho=\operatorname{rotAbout}(c,2\pi/m)$. Assume that $c\notin S$, that $\rho$ maps $S$ into $S$, that $S$ is $K$-free,
and that $S$ is in general position. Let $i\in\N$, $p\in\hv(\peel(S,i))$, and $q:\N\to\R^2$ (write $q_s=q(s)$) be such
that
\begin{itemize}
\item[(C1)] $q_s\in\peel(S,i)$ for all $s<K-1$;
\item[(C2)] $q_s\ne p$ for all $s<K-1$;
\item[(C3)] $0<\orient(p,q_s,q_{s+1})$ for all $s$ with $s+1<K-1$;
\item[(C4)] $0<\orient(q_s,q_{s+1},q_{s+2})$ for all $s$ with $s+2<K-1$.
\end{itemize}
Then a contradiction follows (the Lean conclusion is \lean{False}). In other words, no such $i$, $p$ and $q$ exist.
\end{theorem}

\begin{proof}
Since $K\ge4$, $K-1\ge3$ is the ordinary difference. We have $p\in\peel(S,i)$, as $\hv(\peel(S,i))\subseteq\peel(S,i)$.
For every $s<K-1$, Lemma~\ref{lg:lem:halfplane} with $y=q_s$ (which lies in $\peel(S,i)$ by (C1) and differs from $p$
by (C2)) gives $0<\operatorname{lin}(c-p,p,q_s)$. Apply Theorem~\ref{lg:thm:fan} to the finite set $\peel(S,i)$, the
pivot $p$, the vector $w=c-p$, the sequence $q$ and $n=K-1\ge3$: its hypotheses are $p\in\peel(S,i)$, (C1), the
half-plane inequalities just proved, (C3) and (C4). Hence $\peel(S,i)$ contains $(K-1)+1=K$ points in convex position.
Since $\peel(S,i)\subseteq S$ by (L2), so does $S$ (Lemma~\ref{fd:lem:mono}(c)), contradicting the $K$-freeness of
$S$. (Lean unfolds the definition of \lean{HasConvexNGon} here instead of citing Lemma~\ref{fd:lem:mono}(c).)
\end{proof}

\begin{remark}
The chain points $q_s$ are not assumed distinct, and they may lie in any layer $\ge i$. The theorem is the geometric
statement behind the chain clauses of the layered encoders (``no pivot of layer $i$ with $K-1$ chain points whose
consecutive fan and turn orientations are positive''). That the encoders' clauses express exactly this family
(including, where auxiliary chain variables are used, that their least-fixpoint values satisfy every chain clause;
\dirn{LEANEXT/review/REPORT.md}, finding F5) is a code-level step outside \dirn{LAYERLEAN/LayerLemma.lean}; see the
paragraph \emph{Encoders} at the beginning of this section.
\end{remark}

\subsection{Polar coordinates and the angular rules}\label{lg:ssec:polar}

\begin{definition}[polar; \leanref{ErdosSzekeres.LayerLemma.polar}{LAYERLEAN/LayerLemma.lean:932}]\label{lg:def:polar}
For a point $c$ and $r,\theta\in\R$,
\[
  \operatorname{polar}(c,r,\theta)=\bigl(c_0+r\cos\theta,\;c_1+r\sin\theta\bigr).
\]
There is no sign condition on $r$.
\end{definition}

\begin{lemma}[orientation of three points in polar form; \leanref{ErdosSzekeres.LayerLemma.orient_polar}{LAYERLEAN/LayerLemma.lean:935}]\label{lg:lem:polar}
For every point $c$ and all $r_a,r_b,r_d,\alpha,\beta,\delta\in\R$,
\[
  \orient\bigl(\operatorname{polar}(c,r_a,\alpha),\operatorname{polar}(c,r_b,\beta),\operatorname{polar}(c,r_d,\delta)\bigr)
  =r_ar_b\sin(\beta-\alpha)+r_br_d\sin(\delta-\beta)+r_dr_a\sin(\alpha-\delta).
\]
\end{lemma}

\begin{proof}
Write $e(\varphi)=(\cos\varphi,\sin\varphi)$ and let $a,b,d$ be the three points; then $a-c=A:=r_ae(\alpha)$,
$b-c=B:=r_be(\beta)$, $d-c=D:=r_de(\delta)$, so $b-a=B-A$ and $d-a=D-A$. By (O1) and bilinearity and antisymmetry of
$\det$,
\begin{align*}
  \orient(a,b,d)=\det(B-A,D-A)&=\det(B,D)-\det(B,A)-\det(A,D)+\det(A,A)\\
  &=\det(A,B)+\det(B,D)+\det(D,A).
\end{align*}
Finally, for all $r,s,\varphi,\psi\in\R$,
$\det(re(\varphi),se(\psi))=rs(\cos\varphi\sin\psi-\sin\varphi\cos\psi)=rs\sin(\psi-\varphi)$; applied to the three
terms this gives the formula. (Lean expands the three sines on the right by the subtraction formula and concludes with
\lean{ring}.)
\end{proof}

\begin{lemma}[rule A, centre inside; \leanref{ErdosSzekeres.LayerLemma.orient_polar_centre_inside}{LAYERLEAN/LayerLemma.lean:943}]\label{lg:lem:ruleA}
(Angular rule A: centre inside the triangle.)
Let $c$ be a point and $r_a,r_b,r_d,\alpha,\beta,\delta\in\R$ with $0<r_a$, $0<r_b$, $0<r_d$, and
\[
  0<\beta-\alpha<\pi,\qquad 0<\delta-\beta<\pi,\qquad 0<\alpha-\delta+2\pi<\pi .
\]
Then $0<\orient\bigl(\operatorname{polar}(c,r_a,\alpha),\operatorname{polar}(c,r_b,\beta),\operatorname{polar}(c,r_d,\delta)\bigr)$.
\end{lemma}

\begin{proof}
By Lemma~\ref{lg:lem:polar} the orientation is $r_ar_b\sin(\beta-\alpha)+r_br_d\sin(\delta-\beta)+r_dr_a\sin(\alpha-\delta)$.
By (S1) the sine is positive on $(0,\pi)$, so $\sin(\beta-\alpha)>0$, $\sin(\delta-\beta)>0$, and, by
$2\pi$-periodicity, $\sin(\alpha-\delta)=\sin(\alpha-\delta+2\pi)>0$. With $r_a,r_b,r_d>0$ all three terms are
positive.
\end{proof}

The hypotheses say $\alpha<\beta<\delta<\alpha+2\pi$ with all three counterclockwise gaps $\beta-\alpha$,
$\delta-\beta$, $\alpha+2\pi-\delta$ (which sum to $2\pi$) in $(0,\pi)$. The name of the rule refers to the position of
$c$ inside the triangle; that geometric fact is neither stated nor used here, and only the orientation statement is
proved. Lean states only the positive orientation.

\begin{lemma}[sine of a sum below the sum of the sines; \leanref{ErdosSzekeres.LayerLemma.sin_add_lt}{LAYERLEAN/LayerLemma.lean:956}]\label{lg:lem:sinadd}
For all $x,y\in\R$ with $0<x<\pi$ and $0<y<\pi$: $\sin(x+y)<\sin x+\sin y$.
\end{lemma}

\begin{proof}
By (S1), $\sin x>0$ and $\sin y>0$. Moreover $\cos x<1$: by (S2), $\cos x\le1$, and $\cos x=1$ only for
$x\in2\pi\Z$, which has no point in $(0,\pi)$; likewise $\cos y<1$. Hence $\sin x\cos y<\sin x$ and $\cos x\sin y<\sin y$, and
$\sin(x+y)=\sin x\cos y+\cos x\sin y<\sin x+\sin y$.
\end{proof}

\begin{lemma}[rule B, middle point outermost; \leanref{ErdosSzekeres.LayerLemma.orient_polar_halfplane}{LAYERLEAN/LayerLemma.lean:974}]\label{lg:lem:ruleB}
(Angular rule B: open half-plane, angularly middle point outermost.)
Let $c$ be a point and $r_a,r_b,r_d,\alpha,\beta,\delta\in\R$ with
\[
  0<r_a,\quad 0<r_d,\quad r_a\le r_b,\quad r_d\le r_b,\quad 0<\beta-\alpha,\quad 0<\delta-\beta,\quad \delta-\alpha<\pi .
\]
Then $0<\orient\bigl(\operatorname{polar}(c,r_a,\alpha),\operatorname{polar}(c,r_b,\beta),\operatorname{polar}(c,r_d,\delta)\bigr)$.
\end{lemma}

\begin{proof}
Let $a,b,d$ denote the three points $\operatorname{polar}(c,r_a,\alpha)$, $\operatorname{polar}(c,r_b,\beta)$,
$\operatorname{polar}(c,r_d,\delta)$. Put $x=\beta-\alpha>0$ and $y=\delta-\beta>0$. Then $x+y=\delta-\alpha<\pi$, so
$x=(x+y)-y<\pi$ and $y=(x+y)-x<\pi$. Since $\alpha-\delta=-(x+y)$ and the sine is odd, Lemma~\ref{lg:lem:polar} gives
\begin{align*}
  \orient(a,b,d)&=r_ar_b\sin x+r_br_d\sin y-r_dr_a\sin(x+y)\\
  &=r_a(r_b-r_d)\sin x+r_d(r_b-r_a)\sin y+r_ar_d\bigl(\sin x+\sin y-\sin(x+y)\bigr).
\end{align*}
Here $\sin x,\sin y>0$ by (S1) (arguments in $(0,\pi)$), $r_a,r_d>0$ and $r_b-r_d,\,r_b-r_a\ge0$, so the first two terms are
$\ge0$; the third is $>0$ by Lemma~\ref{lg:lem:sinadd} and $r_ar_d>0$. Hence the orientation is positive. (Lean uses
$r_ar_b\sin x\ge r_ar_d\sin x$, $r_br_d\sin y\ge r_ar_d\sin y$ and Lemma~\ref{lg:lem:sinadd}, and concludes with
\lean{nlinarith}; this is the same estimate.)
\end{proof}

\begin{remark}
The radius hypotheses are weak inequalities ($r_b>0$ follows from $0<r_a\le r_b$). This is the form of the radius
order of Theorem~\ref{lg:thm:radius}, which gives $\operatorname{dist2}(c,y)\le\operatorname{dist2}(c,x)$, that is,
$r_y\le r_x$ for the radii $r=\sqrt{\operatorname{dist2}(c,\cdot)}$. The rules apply to points given in polar form
with explicit angle representatives. No Lean statement of this file writes the points of a layered configuration in
that form (a point $y$ as $\operatorname{polar}(c,r,\theta)$ with $r^2=\operatorname{dist2}(c,y)$, and a layer as
$\{\operatorname{polar}(c,r_i,\theta_i+2\pi k/m)\}$; \dirn{LEANEXT/review/REPORT.md}, finding F2); that step, and the
choice of which rule applies to which triple, are outside these statements (see the paragraph \emph{Encoders} at the
beginning of this section).
\end{remark}

\paragraph{Coverage.}
All 27 declarations of the scope of this section, which are all the declarations of lines 465--1032 of
\dirn{LAYERLEAN/LayerLemma.lean} (23 theorems and 4 definitions), are covered:
\lean{rotAbout_iterate} (Lemma~\ref{lg:lem:rot-iter}; the section on the layer lemma restates it as well),
\lean{orient_rotAbout} (Lemma~\ref{lg:lem:orient-rot}),
\lean{trig_layer} (Lemma~\ref{lg:lem:trig}, stated just before Lemma~\ref{lg:lem:regular}, whose proof uses it),
\lean{orient_regular_nonneg} (Lemma~\ref{lg:lem:regular}),
\lean{lin} (Definition~\ref{lg:def:lin}),
\lean{lin_rotAbout} (Lemma~\ref{lg:lem:lin-rot}),
\lean{lin_orbit_pos} (Lemma~\ref{lg:lem:lin-orbit}),
\lean{orient_next_prev_pos} (Lemma~\ref{lg:lem:next-prev}),
\lean{convex_orient_nonneg} (Lemma~\ref{lg:lem:halfplane-convex}),
\lean{dist2} (Definition~\ref{lg:def:dist2}),
\lean{convex_disk} (Lemma~\ref{lg:lem:disk}),
\lean{strict_nesting} (Theorem~\ref{lg:thm:nesting}),
\lean{centre_not_mem_of_even} (Lemma~\ref{lg:lem:centre}),
\lean{dist2_le_of_later} (Theorem~\ref{lg:thm:radius}),
\lean{orient_trans_halfplane} (Lemma~\ref{lg:lem:trans}),
\lean{fan_step} (Lemma~\ref{lg:lem:fan-step}),
\lean{orient_eq_lin} (Lemma~\ref{lg:lem:orient-lin}),
\lean{fan_positive} (Lemma~\ref{lg:lem:fan-pos}),
\lean{fanPt} (Definition~\ref{lg:def:fanpt}),
\lean{hasConvexNGon_of_fan} (Theorem~\ref{lg:thm:fan}),
\lean{lin_pos_of_peel} (Lemma~\ref{lg:lem:halfplane}),
\lean{chain_clauses_sound} (Theorem~\ref{lg:thm:chain}),
\lean{polar} (Definition~\ref{lg:def:polar}),
\lean{orient_polar} (Lemma~\ref{lg:lem:polar}),
\lean{orient_polar_centre_inside} (Lemma~\ref{lg:lem:ruleA}),
\lean{sin_add_lt} (Lemma~\ref{lg:lem:sinadd}) and
\lean{orient_polar_halfplane} (Lemma~\ref{lg:lem:ruleB}); all in \dirn{LAYERLEAN/LayerLemma.lean}, namespace
\lean{ErdosSzekeres.LayerLemma}.
The written steps (Lemmas~\ref{lg:lem:layer}, \ref{lg:lem:rotfacts}, \ref{lg:lem:cramer} and
Corollary~\ref{lg:cor:period}) are not Lean declarations; the paragraph \emph{Rigor} at the beginning of the section
says how each relates to the Lean proofs.
Deliberately not covered here:
\begin{itemize}
\item \lean{peel_antitone} (line 460) and \lean{card_peel} (line 1037, docstring at line 1036), which adjoin these
  lines but belong to the section on the layer lemma (Lemma~\ref{ll:lem:antitone}, Theorem~\ref{ll:thm:card-peel});
\item the verbatim copies of these declarations in other units, which are duplicates. The following eight files
  contain lines 11--1147 of \dirn{LAYERLEAN/LayerLemma.lean} verbatim (checked with \texttt{diff}); a search of the
  uncompressed files of the repository finds no further copy:\newline
  \dirn{LEANEXT/capstone/Capstone.lean};\newline
  \dirn{LEANEXT/capstone/Capstone2.lean};\newline
  \dirn{LEANEXT/merged/Merged.lean};\newline
  \dirn{LEANEXT/linked/Linked.lean};\newline
  \dirn{LEANEXT/es7corr/ES7Corollaries.lean};\newline
  \dirn{LEANEXT/review/VacuityTest.lean} (all 1148 lines, followed by the tests of the next item);\newline
  \dirn{ENCBLEAN/geo/ES8Geo/Capstone2.lean};\newline
  \dirn{LAYERLEAN/logs/axioms_recheck_20261106/LayerLemmaAxioms.lean};
\item the 19 declarations of \dirn{LEANEXT/review/tests_tail.lean} (namespace
  \lean{ErdosSzekeres.LayerLemma.Review}), which no section covers: the definitions \lean{L4}, \lean{S4}, \lean{H6} and
  the theorems \lean{mem_S4}, \lean{card_S4}, \lean{angle4}, \lean{rot4}, \lean{hc4}, \lean{inv4}, \lean{gp4},
  \lean{inner4}, \lean{hno4}, \lean{hv4}, \lean{peel1_4}, \lean{nest4}, \lean{orbit4}, \lean{layers4}, \lean{chain4}
  and \lean{inv6}. They are non-vacuity tests, not results (\dirn{LEANEXT/review/REPORT.md}, item~6). For the eight
  points $S_4=\{(\pm2,0),(0,\pm2),(\pm1/2,\pm1/2)\}$, the centre $c=0$, $m=4$ and $K=8$, the theorems \lean{hc4},
  \lean{inv4}, \lean{gp4} and \lean{hno4} state that $c\notin S_4$, that $\operatorname{rotAbout}(0,2\pi/4)$ maps $S_4$
  into $S_4$, that $S_4$ is in general position and that $S_4$ is $8$-free; \lean{hv4} and \lean{peel1_4} state that
  $(2,0)\in\hv(S_4)$ and $(1/2,1/2)\in\peel(S_4,1)$; and \lean{nest4}, \lean{orbit4}, \lean{layers4} and \lean{chain4}
  apply \lean{strict_nesting} (Theorem~\ref{lg:thm:nesting}), \lean{hullVerts_eq_orbitFS_rotAbout}
  (Theorem~\ref{ll:thm:core-rot}), \lean{layers_are_orbits} (Theorem~\ref{ll:thm:layers}) and
  \lean{chain_clauses_sound} (Theorem~\ref{lg:thm:chain}) to this instance. \lean{inv6} states that
  $\operatorname{rotAbout}(0,2\pi/6)$ maps the set \lean{H6} of the six points $(\cos(k\pi/3),\sin(k\pi/3))$,
  $0\le k<6$ (a regular hexagon), into itself. \dirn{LEANEXT/review/VacuityTest.lean} consists of all 1148 lines of
  \dirn{LAYERLEAN/LayerLemma.lean} followed by the 149 lines of \dirn{LEANEXT/review/tests_tail.lean}, as lines
  1149--1297 (both parts checked with \texttt{cmp}), so the same 19 declarations stand there at lines 1160--1276. Both
  files end with \texttt{\#print axioms}, \texttt{\#check} and \texttt{\#print} commands, which are not declarations;
\item declarations in other namespaces that share a short name with one of the 27 (in
  \lean{ErdosSzekeres.OrbitConsistency}, \lean{PhaseClassification}, \lean{C7} and
  \lean{ErdosSzekeres.SymReduction.Review}); they are separate Lean declarations. Those of
  \dirn{LEANEXT/phase/PhaseClassification.lean} (\lean{polar}, \lean{orient_polar}, \lean{orient_polar_centre_inside},
  \lean{sin_add_lt} and \lean{orient_polar_halfplane}) are treated in the section on phase patterns.
\end{itemize}

\section{Orbit literals, rays and phases, key order}\label{oc:sec}
% Long Lean names and file paths: allow loose lines in this section (group closed at the end).
\begingroup\sloppy\emergencystretch=15em

This section gives written proofs of the Lean-checked results of three files of the \dirn{LEANEXT}
unit. They are the geometric and arithmetic facts behind three features of the layered SAT encoders
for point sets with $m$-fold rotational symmetry about a centre $c$.
\begin{enumerate}[label=(\Alph*)]
\item \emph{Orbit literals} (\S\ref{oc:sub:rot}, \S\ref{oc:sub:perm}). Rotations about $c$ preserve
  $\orient$, and permuting an ordered triple multiplies its orientation
  $\orient_3(t)=\orient(t(0),t(1),t(2))$ by the sign of the permutation. Hence, if an ordered triple
  $C$ of points arises from an ordered triple $t$ by a permutation $\sigma$ of the positions followed by
  a composition of rotations by integer multiples of $2\pi/m$, and $\orient_3(t)\ne0$, then
  ``$\orient_3(t)>0$'' is equivalent to ``$\orient_3(C)>0$, negated when $\sgn\sigma=-1$''
  (Theorems~\ref{oc:thm:literal} and~\ref{oc:thm:encoder}). This is the rule ``one SAT variable per
  orbit of ordered triples''.
\item \emph{Rays and phases for even $m>0$} (\S\ref{oc:sub:even}, \S\ref{oc:sub:layers}). For even
  $m>0$, a finite set in general position that is mapped into itself by the rotation $\rho$ by $2\pi/m$
  about $c$ contains the antipode $2c-x$ of each of its points, and no two of its points lie on one open
  ray from $c$ (Lemma~\ref{oc:lem:same-ray}). Consequently two of its points whose directions agree
  modulo $2\pi/m$ lie in one $\rho$-orbit (Theorem~\ref{oc:thm:phase-orbit}). Combined with the layer
  lemma (under its hypotheses: $c$ is not in the set, and the set has no $K$ points in convex position
  for some $K\le2m$), hull vertices of distinct layers never lie on one open ray from $c$ through a
  rotated point and have directions that are incongruent modulo $2\pi/m$
  (Theorems~\ref{oc:thm:layers-ray} and~\ref{oc:thm:layers-phase}). Thus, for even $m>0$, this
  distinctness of phases is obtained without any perturbation argument.
\item \emph{Key order} (\S\ref{oc:sub:keys}). For $u>0$, real numbers $Ju+\psi$ with $J\in\Z$ and
  $0\le\psi<u$ are ordered lexicographically by $(J,\psi)$, or by $(J,\operatorname{rank}\psi)$ for a
  strictly increasing rank; reduction modulo $2\pi$ into $[0,2\pi)$ replaces $J$ by $J\bmod N$ when
  $Nu=2\pi$; and the antipode $\theta+\pi$ has an explicit cell (Theorem~\ref{oc:thm:key},
  Lemma~\ref{oc:lem:antipode-cell}).
\end{enumerate}

\paragraph{Sources.} File paths are relative to the results directory of the project, except that paths
beginning with \dirn{ErdosSzekeres/} are relative to the repository root; a bare file name such as
\dirn{OrbitConsistency.lean}, \dirn{Merged.lean} or \dirn{KeyOrder.lean} refers to the file of that
name listed below. SHA-256 digests are abbreviated to their first 12 hexadecimal digits.
\begin{itemize}
\item \dirn{LEANEXT/orbit/OrbitConsistency.lean} (429 lines, SHA-256 \texttt{9b843180ac38}), namespace
  \lean{ErdosSzekeres.OrbitConsistency}. The new content starts at line 114. Lines 39--112 repeat
  definitions and basic lemmas of \dirn{LAYERLEAN/LayerLemma.lean} as separate constants
  (Remark~\ref{oc:rem:copies}).
\item \dirn{LEANEXT/merged/Merged.lean} (1385 lines, SHA-256 \texttt{245b7d65407a}), namespace
  \lean{ErdosSzekeres.LayerLemma}, section~13 (lines 1149--1383). Lines 1--1147 of this file coincide
  line by line with lines 1--1147 of \dirn{LAYERLEAN/LayerLemma.lean} (SHA-256 \texttt{afb8542dceb0});
  hull vertices, peels and the layer lemma used in \S\ref{oc:sub:layers} come from there. The seven
  declarations at lines 1160--1311 of \dirn{Merged.lean} restate those at lines 273--427 of
  \dirn{OrbitConsistency.lean} over the definitions of the namespace \lean{ErdosSzekeres.LayerLemma};
  the proofs are identical and only the doc comments differ.
\item \dirn{LEANEXT/keys/KeyOrder.lean} (158 lines, SHA-256 \texttt{6e34096cff52}), namespace
  \lean{KeyOrder}; it imports Mathlib only.
\end{itemize}
Byte-identical copies of \dirn{OrbitConsistency.lean} and \dirn{KeyOrder.lean} (same digests) are
\dirn{ENCBLEAN/geo/ES8Geo/OrbitConsistency.lean} and \dirn{ENCBLEAN/geo/ES8Geo/KeyOrder.lean}. Lines
1148--1383 of \dirn{Merged.lean} reappear verbatim as lines 1148--1383 of
\dirn{LEANEXT/linked/Linked.lean} and as lines 1156--1391 of \dirn{LEANEXT/capstone/Capstone.lean} and
of \dirn{LEANEXT/capstone/Capstone2.lean}; \dirn{ENCBLEAN/geo/ES8Geo/Capstone2.lean} is a byte-identical
copy of the last (SHA-256 \texttt{f8c3930c2427}). The files are compiled from the repository root
(\dirn{LEANEXT/README.md:6}), whose pinned toolchain is Lean 4.24.0 with Mathlib \texttt{f897ebcf}.

\paragraph{Axioms.} \dirn{LEANEXT/orbit/logs/axioms.log} records \texttt{\#print axioms} for the 15 main
theorems of \dirn{OrbitConsistency.lean} (\lean{orient_rotAbout_inv}, \lean{RotComp.preservesOrient},
\lean{rotComp_iterate}, \lean{orient3_swap}, \lean{orient3_perm}, \lean{orient3_rot_perm},
\lean{sign_eq_of_two_reps}, \lean{orient3_of_common_rep}, \lean{literal_sound},
\lean{encoder_literal_sound}, \lean{antipode_mem_of_even}, \lean{same_ray_false},
\lean{eq_iterate_of_phase_eq}, \lean{eq_iterate_of_phase_eq_int}, \lean{phase_ne_of_not_mem_orbit}),
each depending on \texttt{[propext, Classical.choice, Quot.sound]}, with exit status 0.
\dirn{LEANEXT/merged/axioms.log} records the same three axioms for
\lean{ErdosSzekeres.LayerLemma.phase_ne_of_distinct_layers},
\lean{ErdosSzekeres.LayerLemma.not_same_ray_of_distinct_layers} and
\lean{ErdosSzekeres.LayerLemma.phase_ne_of_not_mem_orbit} (and for the copies in \dirn{Merged.lean} of
three theorems of \dirn{LAYERLEAN/LayerLemma.lean}: \lean{layers_are_orbits},
\lean{convexNGon_of_layers} and \lean{sevenfold_octagon}). Lines 151--158 of \dirn{KeyOrder.lean} print
the axioms of its eight theorems; that output, again the same three axioms, is recorded for the
byte-identical copy in \dirn{ENCBLEAN/logs/m4_build.log} (lines 3--10), not inside \dirn{LEANEXT}. The
remaining theorems of the scope that carry a Lean reference below (the copied lemmas
\lean{iterate_mem} and \lean{rotAbout_iterate} of \dirn{OrbitConsistency.lean}, both used by
\lean{antipode_mem_of_even}; the helpers \lean{preservesOrient_id}, \lean{PreservesOrient.comp},
\lean{preservesOrient_rotAbout}, \lean{rotAbout_polar}, \lean{polar_add_int}; the other
\dirn{Merged.lean} restatements; and \lean{not_mem_hullVerts_peel_of_lt},
\lean{not_mem_hullVerts_peel_of_ne}, \lean{iterate_mem_hullVerts_peel}) are used by printed theorems,
and \texttt{\#print axioms} reports the axioms of all transitive dependencies; so their axioms are among
the three listed. There are two exceptions. (i) \lean{PreservesOrient.iterate} is used by no printed
theorem, and no axiom output is recorded for it; that it compiles is recorded by the build log of the
whole file, \dirn{LEANEXT/orbit/logs/build.log} (exit status 0). (ii) Lemmas~\ref{oc:lem:iterate}
and~\ref{oc:lem:iter-mem} also cite the \dirn{LAYERLEAN} declarations
\lean{ErdosSzekeres.LayerLemma.rotAbout_iterate} and \lean{ErdosSzekeres.LayerLemma.iterate_mem}. Their
axioms are recorded in \dirn{LEANEXT/review/axioms.log} (lines 33 and 1), from the review of 4 November
2026 of \dirn{LAYERLEAN/LayerLemma.lean} (\dirn{LEANEXT/review/REPORT.md}, lines 3--4): the three axioms
above for \lean{rotAbout_iterate}, and only \lean{propext} and \lean{Quot.sound} for
\lean{iterate_mem}. None of the three files contains \lean{sorry}, \lean{admit}, an \lean{axiom}
declaration or \lean{native_decide}. Within the scope of this section the only \lean{decide} (kernel
evaluation, no extra axiom) is the side goal $-1\ne1$ in $\Z^\times$ at
\dirn{OrbitConsistency.lean:236}; the copied part of \dirn{Merged.lean} has one more, for
$\mathrm{Even}(6)$ at line 1125.

\paragraph{Rigor.} Every statement below that carries a Lean reference is \emph{Lean-checked} in the
sense of the records above (for a definition: it is part of the compiled file). The proofs in this
section are \emph{written proofs (not refereed)}. They
follow the Lean proofs; where they take a different route, this is said in the proof
(Lemma~\ref{oc:lem:iterate}: addition formulas instead of complex exponentials;
Lemma~\ref{oc:lem:swap}: the symmetries of $\orient$ instead of a six-case expansion;
Lemma~\ref{oc:lem:closed}: $\rho^{[m]}=\mathrm{id}$ directly instead of the minimal period). No
statement of this section takes an unproved proposition or a SAT verdict as a hypothesis.

\paragraph{Dependencies.} From the foundations section we use the definitions of points, $\orient$,
general position, \lean{HasConvexNGon}, invariance and rotations (Definitions~\ref{fd:def:point},
\ref{fd:def:orient}, \ref{fd:def:gp}, \ref{fd:def:ngon}, \ref{fd:def:inv}, \ref{fd:def:rot}), the
symmetries of $\orient$ (Lemma~\ref{fd:lem:orient}), the orientation test for general position
(Lemma~\ref{fd:lem:collinear}) and the monotonicity of \lean{HasConvexNGon} (Lemma~\ref{fd:lem:mono});
Lemma~\ref{fd:lem:affine-orient} is cited only in a note on notation in
Definition~\ref{oc:def:basic}(b) and in a parenthetical comparison in the proof of
Lemma~\ref{oc:lem:rot-inv}, and the example in Remark~\ref{oc:rem:layers} uses
Lemmas~\ref{fd:lem:cp-hull} and~\ref{fd:lem:separation}. From the layer-lemma section we use the hull
vertices and peels (Definitions~\ref{ll:def:hv} and~\ref{ll:def:peel}), Lemma~\ref{ll:lem:antitone},
Lemma~\ref{ll:lem:inv-peel} and Theorem~\ref{ll:thm:core-rot}; Lemma~\ref{ll:lem:minperiod} is cited
only in a comparison with the Lean proof (proof of Lemma~\ref{oc:lem:closed}). From the layer-geometry
section we use the definition of $\operatorname{polar}$ (Definition~\ref{lg:def:polar}) and, only in
comments on hypotheses (Lemma~\ref{oc:lem:same-ray} and Remark~\ref{oc:rem:layers}),
Lemma~\ref{lg:lem:centre}. \S\ref{oc:sub:keys} uses no other section.

\subsection{Conventions and copied definitions}\label{oc:sub:conv}

Points are elements $p=(p_0,p_1)$ of $\R^2$ (Definition~\ref{fd:def:point}); sums, differences and scalar multiples
of points are taken coordinatewise, and the Lean term \texttt{![u, v]} denotes the point $(u,v)$. All sets
of points in this section are finite (Lean \texttt{Finset Point}).
$\orient(a,b,d)=(b_0-a_0)(d_1-a_1)-(b_1-a_1)(d_0-a_0)$ (Definition~\ref{fd:def:orient}); equivalently
$\orient(a,b,d)=\det\bigl(b-a\mid d-a\bigr)$, the determinant of the $2\times2$ matrix with columns $b-a$
and $d-a$. For a map $f:X\to X$ the iterates are $f^{[0]}=\mathrm{id}$ and
$f^{[s+1]}=f\circ f^{[s]}$ ($=f^{[s]}\circ f$). For $m\in\N$, the real number $2\pi/m$ is $2\pi$ divided
by $m$ cast to $\R$, with Lean's convention $x/0=0$, so that $2\pi/0=0$. ``$m$ is even'' means $m=k+k$
for some $k\in\N$ (Mathlib's \lean{Even}); then the natural-number quotient $m/2$ equals $k$.

\begin{definition}[invariance, rotations, polar points; \leanref{ErdosSzekeres.OrbitConsistency.Invariant}{LEANEXT/orbit/OrbitConsistency.lean:47}, \leanref{ErdosSzekeres.OrbitConsistency.rotAbout}{LEANEXT/orbit/OrbitConsistency.lean:93}, \leanref{ErdosSzekeres.OrbitConsistency.polar}{LEANEXT/orbit/OrbitConsistency.lean:112}]\label{oc:def:basic}
Let $f:\R^2\to\R^2$, $S$ a finite set of points, $c,p\in\R^2$ and $r,\theta\in\R$.
\begin{enumerate}[label=(\alph*)]
\item $f$ \emph{maps $S$ into $S$} (Lean \texttt{Invariant f S}) if $f(p)\in S$ for every $p\in S$
  (Definition~\ref{fd:def:inv}).
\item The rotation by $\theta$ about $c$ (Lean \lean{rotAbout}; Definition~\ref{fd:def:rot}) is
  \[R_{c,\theta}(p)=\bigl(c_0+\cos\theta\,(p_0-c_0)-\sin\theta\,(p_1-c_1),\;
     c_1+\sin\theta\,(p_0-c_0)+\cos\theta\,(p_1-c_1)\bigr),\]
  that is, $R_{c,\theta}(p)=c+M_\theta(p-c)$ with $M_\theta=\begin{pmatrix}\cos\theta&-\sin\theta\\
  \sin\theta&\cos\theta\end{pmatrix}$; it turns counterclockwise. For $m\in\N$ we write
  $\rho=R_{c,2\pi/m}$ when $c$ and $m$ are clear. (Notation: Definition~\ref{fd:def:rot} and the
  layer-lemma section write $\rho_{c,\theta}$ for the map $R_{c,\theta}$, so that our $\rho$ is their
  $\rho_{c,2\pi/m}$; the foundations and layer-geometry sections write $R_\theta$ for the matrix
  $M_\theta$, for instance in Lemma~\ref{fd:lem:affine-orient}. In this section $R_{c,\theta}$ always
  denotes the map, never a matrix.)
\item $\operatorname{polar}(c,r,\theta)=(c_0+r\cos\theta,\;c_1+r\sin\theta)$, for every real $r$
  (Definition~\ref{lg:def:polar}).
\end{enumerate}
\end{definition}

\begin{remark}[two copies of the definitions]\label{oc:rem:copies}
Lines 39--112 of \dirn{OrbitConsistency.lean} define their own \lean{iterate_mem}, \lean{Invariant},
\lean{toC}, \lean{ofC}, \lean{crot}, \lean{rotAbout}, \lean{rotAbout_iterate}, \lean{polar} and some
helper lemmas, with the same declaration text as the corresponding declarations of
\dirn{LAYERLEAN/LayerLemma.lean} (\lean{iterate_mem} at line 52, \lean{Invariant} at 105, the
complex-coordinate helpers among lines 231--308, \lean{rotAbout} at 298, \lean{rotAbout_iterate} at
465, \lean{polar} at 932), except that \lean{iterate_mem} declares its type variable in its own binder
(in \dirn{LAYERLEAN/LayerLemma.lean} it is a section variable) and that the doc comment of \lean{polar} is
worded differently; neither difference affects the meaning. These are distinct Lean constants with identical
defining formulas; for instance, the Lean theorem \lean{EncB.Geo.rotAbout_eq_orbitConsistency}
(\dirn{ENCBLEAN/geo/ES8Geo/Orbit.lean:13}) states the equality of the two \lean{rotAbout}, with proof
\lean{rfl} (in that build \lean{ErdosSzekeres.LayerLemma.rotAbout} is the copy in
\dirn{ENCBLEAN/geo/ES8Geo/Capstone2.lean}, whose text equals that of
\dirn{LAYERLEAN/LayerLemma.lean:298}). The theorems of \dirn{OrbitConsistency.lean} are stated with the
\lean{OrbitConsistency} copies, those of \dirn{Merged.lean} with the \lean{LayerLemma} copies
(\lean{ErdosSzekeres.LayerLemma.Invariant}, \dirn{LAYERLEAN/LayerLemma.lean:105};
\lean{ErdosSzekeres.LayerLemma.rotAbout}, \dirn{LAYERLEAN/LayerLemma.lean:298};
\lean{ErdosSzekeres.LayerLemma.polar}, \dirn{LAYERLEAN/LayerLemma.lean:932}). Mathematically each pair
is one and the same function, so we state each result once and list both Lean names where both exist;
the statement then reads verbatim for either copy.
\end{remark}

\begin{lemma}[composition of rotations about one centre; written step]\label{oc:lem:rot-add}
For all $c,p\in\R^2$ and $\alpha,\beta\in\R$:
\begin{enumerate}[label=(\alph*)]
\item $R_{c,\alpha}(R_{c,\beta}(p))=R_{c,\alpha+\beta}(p)$;
\item $R_{c,0}=\mathrm{id}$, $R_{c,2\pi}=\mathrm{id}$ and $R_{c,\pi}(p)=(2c_0-p_0,\;2c_1-p_1)$.
\end{enumerate}
\end{lemma}
\begin{proof}
(a) Put $v=p-c$. Then $R_{c,\beta}(p)-c=(\cos\beta\,v_0-\sin\beta\,v_1,\;\sin\beta\,v_0+\cos\beta\,v_1)$, and
applying $R_{c,\alpha}$ gives $c$ plus the vector with coordinates
\begin{align*}
&\cos\alpha(\cos\beta\,v_0-\sin\beta\,v_1)-\sin\alpha(\sin\beta\,v_0+\cos\beta\,v_1)
 =\cos(\alpha+\beta)\,v_0-\sin(\alpha+\beta)\,v_1,\\
&\sin\alpha(\cos\beta\,v_0-\sin\beta\,v_1)+\cos\alpha(\sin\beta\,v_0+\cos\beta\,v_1)
 =\sin(\alpha+\beta)\,v_0+\cos(\alpha+\beta)\,v_1,
\end{align*}
by the addition formulas $\cos(\alpha+\beta)=\cos\alpha\cos\beta-\sin\alpha\sin\beta$ and
$\sin(\alpha+\beta)=\sin\alpha\cos\beta+\cos\alpha\sin\beta$. This is $R_{c,\alpha+\beta}(p)-c$.
(b) Substitute $\cos0=\cos2\pi=1$, $\sin0=\sin2\pi=0$, and $\cos\pi=-1$, $\sin\pi=0$: then
$R_{c,0}(p)=R_{c,2\pi}(p)=c+(p-c)=p$ and $R_{c,\pi}(p)=c-(p-c)=2c-p$.
\end{proof}

\begin{lemma}[iterates of a rotation; \leanref{ErdosSzekeres.OrbitConsistency.rotAbout_iterate}{LEANEXT/orbit/OrbitConsistency.lean:105}, \leanref{ErdosSzekeres.LayerLemma.rotAbout_iterate}{LAYERLEAN/LayerLemma.lean:465}]\label{oc:lem:iterate}
For all $c\in\R^2$, $\alpha\in\R$ and $j\in\N$: $(R_{c,\alpha})^{[j]}=R_{c,j\alpha}$.
\end{lemma}
\begin{proof}
Induction on $j$. For $j=0$ both sides are the identity (Lemma~\ref{oc:lem:rot-add}(b)). If the claim
holds for $j$, then $(R_{c,\alpha})^{[j+1]}=R_{c,\alpha}\circ(R_{c,\alpha})^{[j]}=R_{c,\alpha}\circ
R_{c,j\alpha}=R_{c,\alpha+j\alpha}=R_{c,(j+1)\alpha}$ by Lemma~\ref{oc:lem:rot-add}(a). (The Lean proof
writes $R_{c,\alpha}$ as $z\mapsto c+e^{i\alpha}(z-c)$ and uses $(e^{i\alpha})^j=e^{ij\alpha}$; the
argument above is an equivalent elementary route.)
\end{proof}

\begin{corollary}[period and half-turn; written step]\label{oc:cor:period}
Let $c\in\R^2$, $m\in\N$ with $0<m$, and $\rho=R_{c,2\pi/m}$. Then $\rho^{[m]}=\mathrm{id}$. If moreover
$m$ is even, then $\rho^{[m/2]}(p)=(2c_0-p_0,\;2c_1-p_1)$ for every $p\in\R^2$.
\end{corollary}
\begin{proof}
Since $m\ne0$, $m\cdot(2\pi/m)=2\pi$, so $\rho^{[m]}=R_{c,2\pi}=\mathrm{id}$ by
Lemmas~\ref{oc:lem:iterate} and~\ref{oc:lem:rot-add}(b). If $m=k+k$, then $k\ge1$ (as $m>0$), $m/2=k$
and $k\cdot\frac{2\pi}{2k}=\pi$, so $\rho^{[m/2]}=R_{c,\pi}$, and Lemma~\ref{oc:lem:rot-add}(b) applies.
\end{proof}

\begin{lemma}[iterates stay in a set mapped into itself; \leanref{ErdosSzekeres.OrbitConsistency.iterate_mem}{LEANEXT/orbit/OrbitConsistency.lean:39}, \leanref{ErdosSzekeres.LayerLemma.iterate_mem}{LAYERLEAN/LayerLemma.lean:52}]\label{oc:lem:iter-mem}
Let $X$ be a type, $f:X\to X$, $S$ a finite subset of $X$ with $f(y)\in S$ for every $y\in S$, and $x\in
S$. Then $f^{[i]}(x)\in S$ for every $i\in\N$.
\end{lemma}
\begin{proof}
Induction on $i$: $f^{[0]}(x)=x\in S$, and if $f^{[i]}(x)\in S$ then $f^{[i+1]}(x)=f(f^{[i]}(x))\in S$.
\end{proof}

\subsection{Rotations preserve orientation}\label{oc:sub:rot}

\begin{lemma}[rotation invariance of the orientation; \leanref{ErdosSzekeres.OrbitConsistency.orient_rotAbout_inv}{LEANEXT/orbit/OrbitConsistency.lean:117}]\label{oc:lem:rot-inv}
For all $c\in\R^2$, $\theta\in\R$ and $a,b,d\in\R^2$,
\[\orient\bigl(R_{c,\theta}(a),\,R_{c,\theta}(b),\,R_{c,\theta}(d)\bigr)=\orient(a,b,d).\]
There are no further hypotheses.
\end{lemma}
\begin{proof}
Write $M=M_\theta$. Since $R_{c,\theta}(p)=c+M(p-c)$, we have $R_{c,\theta}(b)-R_{c,\theta}(a)=M(b-a)$
and $R_{c,\theta}(d)-R_{c,\theta}(a)=M(d-a)$. Hence
\begin{align*}
\orient\bigl(R_{c,\theta}(a),R_{c,\theta}(b),R_{c,\theta}(d)\bigr)
 &=\det\bigl(M(b-a)\mid M(d-a)\bigr)=\det M\cdot\det\bigl(b-a\mid d-a\bigr)\\
 &=(\cos^2\theta+\sin^2\theta)\,\orient(a,b,d),
\end{align*}
and $\cos^2\theta+\sin^2\theta=1$. (This is the case $\det=1$ of Lemma~\ref{fd:lem:affine-orient},
item~(i), which writes the map $R_{o,\theta}$ as $\rho_{o,\theta}$ and the matrix $M_\theta$ as
$R_\theta$.) The Lean proof is the same computation in expanded form: the difference of the two sides equals
$\orient(a,b,d)\,(\sin^2\theta+\cos^2\theta-1)$ as a polynomial identity in the coordinates and in
$\cos\theta,\sin\theta$ (\lean{linear_combination} with \lean{Real.sin_sq_add_cos_sq}).
\end{proof}

\begin{definition}[orientation-preserving maps; \leanref{ErdosSzekeres.OrbitConsistency.PreservesOrient}{LEANEXT/orbit/OrbitConsistency.lean:124}]\label{oc:def:preserves}
A map $f:\R^2\to\R^2$ \emph{preserves orientation} if $\orient(f(a),f(b),f(d))=\orient(a,b,d)$ for all
$a,b,d\in\R^2$.
\end{definition}

\begin{definition}[compositions of rotations by multiples of $2\pi/m$; \leanref{ErdosSzekeres.OrbitConsistency.RotComp}{LEANEXT/orbit/OrbitConsistency.lean:145}]\label{oc:def:rotcomp}
Let $c\in\R^2$ and $m\in\N$, and put $u=2\pi/m$ (so $u=0$ if $m=0$). $\mathrm{RotComp}(c,m)$ is the
inductively defined set of maps $\R^2\to\R^2$ given by two rules:
\begin{itemize}
\item[(id)] $\mathrm{id}\in\mathrm{RotComp}(c,m)$;
\item[(step)] if $g\in\mathrm{RotComp}(c,m)$ and $k\in\Z$, then $R_{c,ku}\circ g\in\mathrm{RotComp}(c,m)$.
\end{itemize}
It is the smallest set of maps that contains $\mathrm{id}$ and is closed under (step). Its elements are
exactly the maps $R_{c,k_nu}\circ\dots\circ R_{c,k_1u}$ with $n\ge0$ and $k_1,\dots,k_n\in\Z$ (the
identity for $n=0$): this set of compositions contains $\mathrm{id}$ and is closed under (step), and each
of its elements arises from $\mathrm{id}$ by $n$ applications of (step). Proofs about
$\mathrm{RotComp}(c,m)$ use \emph{rule induction}: a property of maps that holds for $\mathrm{id}$ and
passes from $g$ to $R_{c,ku}\circ g$ holds for every element of $\mathrm{RotComp}(c,m)$.
\end{definition}

\begin{lemma}[orientation-preserving maps and rotation compositions; \leanref{ErdosSzekeres.OrbitConsistency.preservesOrient_id}{LEANEXT/orbit/OrbitConsistency.lean:127}, \leanref{ErdosSzekeres.OrbitConsistency.PreservesOrient.comp}{LEANEXT/orbit/OrbitConsistency.lean:129}, \leanref{ErdosSzekeres.OrbitConsistency.PreservesOrient.iterate}{LEANEXT/orbit/OrbitConsistency.lean:134}, \leanref{ErdosSzekeres.OrbitConsistency.preservesOrient_rotAbout}{LEANEXT/orbit/OrbitConsistency.lean:141}, \leanref{ErdosSzekeres.OrbitConsistency.RotComp.preservesOrient}{LEANEXT/orbit/OrbitConsistency.lean:150}, \leanref{ErdosSzekeres.OrbitConsistency.rotComp_iterate}{LEANEXT/orbit/OrbitConsistency.lean:157}]\label{oc:lem:rotcomp-pres}
The following properties of orientation-preserving maps (Definition~\ref{oc:def:preserves}) and of
the sets $\mathrm{RotComp}(c,m)$ (Definition~\ref{oc:def:rotcomp}) hold.
\begin{enumerate}[label=(\alph*)]
\item The identity preserves orientation. If $f$ and $g$ preserve orientation, so does $f\circ g$. If $f$
  preserves orientation, so does $f^{[n]}$ for every $n\in\N$. For all $c\in\R^2$ and $\theta\in\R$ the
  rotation $R_{c,\theta}$ preserves orientation.
\item Let $c\in\R^2$ and $m\in\N$ (any $m$, including $m=0$). Every $g\in\mathrm{RotComp}(c,m)$ preserves
  orientation.
\item Let $c\in\R^2$, $m\in\N$ (any $m$) and $\rho=R_{c,2\pi/m}$. Then $\rho^{[s]}\in\mathrm{RotComp}(c,m)$
  for every $s\in\N$.
\end{enumerate}
\end{lemma}
\begin{proof}
(a) For the identity there is nothing to show. If $f$ and $g$ preserve orientation, then for all $a,b,d$,
$\orient(f(g(a)),f(g(b)),f(g(d)))=\orient(g(a),g(b),g(d))=\orient(a,b,d)$, using first $f$ (at the
points $g(a),g(b),g(d)$) and then $g$. For the iterates use induction on $n$: $f^{[0]}=\mathrm{id}$, and
$f^{[n+1]}=f^{[n]}\circ f$ is a composition of two orientation-preserving maps. Rotations preserve
orientation by Lemma~\ref{oc:lem:rot-inv}.

(b) Rule induction (Definition~\ref{oc:def:rotcomp}): the identity preserves orientation, and if $g$
does, then so does $R_{c,ku}\circ g$, as a composition of two orientation-preserving maps by (a).

(c) Induction on $s$. $\rho^{[0]}=\mathrm{id}\in\mathrm{RotComp}(c,m)$ by (id). If
$\rho^{[s]}\in\mathrm{RotComp}(c,m)$, then $\rho^{[s+1]}=\rho\circ\rho^{[s]}=R_{c,1\cdot u}\circ\rho^{[s]}$
with $u=2\pi/m$, which lies in $\mathrm{RotComp}(c,m)$ by (step) with $k=1$.
\end{proof}

For $m=0$ we have $u=0$ and $R_{c,0}=\mathrm{id}$, so every element of $\mathrm{RotComp}(c,0)$ is the
identity; the statements of this subsection and the next hold for every $m\in\N$.

\subsection{Permutations and orbit literals}\label{oc:sub:perm}

\begin{definition}[orientation of an ordered triple; \leanref{ErdosSzekeres.OrbitConsistency.orient3}{LEANEXT/orbit/OrbitConsistency.lean:168}]\label{oc:def:orient3}
An \emph{ordered triple} is a map $t:\{0,1,2\}\to\R^2$ (in Lean a function \texttt{Fin 3} $\to$
\lean{Point}), and $\orient_3(t)=\orient(t(0),t(1),t(2))$. For a permutation $\sigma$ of $\{0,1,2\}$ and a
map $g:\R^2\to\R^2$, the triples $t\circ\sigma$ and $g\circ t\circ\sigma$ are $i\mapsto t(\sigma(i))$ and
$i\mapsto g(t(\sigma(i)))$.
\end{definition}

\begin{fact}[the sign of a permutation; cited, Mathlib \lean{Equiv.Perm.sign}]\label{oc:fact:sign}
Let $F$ be a finite set and $\mathrm{Perm}(F)$ the group of its permutations, with product
$(\sigma\tau)(i)=\sigma(\tau(i))$. Mathlib's sign is a group homomorphism
$\sgn:\mathrm{Perm}(F)\to\Z^\times$ with the following properties, which are all we use.
\begin{itemize}
\item[(S1)] $\Z^\times=\{1,-1\}$ (\lean{Int.units_eq_one_or}).
\item[(S2)] $\sgn(\sigma\tau)=\sgn(\sigma)\sgn(\tau)$ (\lean{Equiv.Perm.sign_mul}); hence, with (S1),
  $\sgn(\mathrm{id})=1$ and $\sgn(\sigma^{-1})=\sgn(\sigma)^{-1}=\sgn(\sigma)$.
\item[(S3)] $\sgn(\mathrm{swap}(x,y))=-1$ for $x\ne y$, where $\mathrm{swap}(x,y)$ exchanges $x$ and $y$
  (\lean{Equiv.Perm.sign_swap}).
\item[(S4)] Every permutation of $F$ is a product of transpositions; in the form used: if a property $P$
  of permutations holds for $\mathrm{id}$, and $P(f)$ implies $P(\mathrm{swap}(x,y)\,f)$ for every $f$ and
  all $x\ne y$, then $P$ holds for every permutation (\lean{Equiv.Perm.swap_induction_on}).
\end{itemize}
In real formulas $\sgn\sigma$ stands for the image of the unit $\sgn\sigma$ under $\Z^\times\to\Z\to\R$,
the real number $1$ or $-1$. This map is multiplicative, as a composite of monoid homomorphisms, so
(S2) also holds for these real numbers; and it is injective because $1\ne-1$ in $\R$.
\end{fact}

\begin{lemma}[a transposition negates the orientation; \leanref{ErdosSzekeres.OrbitConsistency.orient3_swap}{LEANEXT/orbit/OrbitConsistency.lean:171}]\label{oc:lem:swap}
For every ordered triple $t$ and all $x,y\in\{0,1,2\}$ with $x\ne y$:
$\orient_3(t\circ\mathrm{swap}(x,y))=-\orient_3(t)$.
\end{lemma}
\begin{proof}
Since $\mathrm{swap}(x,y)=\mathrm{swap}(y,x)$, it suffices to treat $\{x,y\}=\{1,2\},\{0,1\},\{0,2\}$.
Write $t_i=t(i)$ and $o=\orient(t_0,t_1,t_2)$; we use $\orient(a,c,b)=-\orient(a,b,c)$ and
$\orient(b,c,a)=\orient(a,b,c)$ (Lemma~\ref{fd:lem:orient}(a),(b)).
For $\{1,2\}$: $t\circ\mathrm{swap}(1,2)=(t_0,t_2,t_1)$ and $\orient(t_0,t_2,t_1)=-o$.
For $\{0,1\}$: $t\circ\mathrm{swap}(0,1)=(t_1,t_0,t_2)$, and the cyclic rule applied to $(t_1,t_0,t_2)$
gives $\orient(t_1,t_0,t_2)=\orient(t_0,t_2,t_1)=-o$.
For $\{0,2\}$: $t\circ\mathrm{swap}(0,2)=(t_2,t_1,t_0)$, and the cyclic rule applied to $(t_2,t_1,t_0)$
gives $\orient(t_2,t_1,t_0)=\orient(t_1,t_0,t_2)=-o$ by the previous case.
(The Lean proof runs through the six ordered pairs $(x,y)$ and expands $\orient$.)
\end{proof}

\begin{lemma}[permutations multiply the orientation by the sign; \leanref{ErdosSzekeres.OrbitConsistency.orient3_perm}{LEANEXT/orbit/OrbitConsistency.lean:178}]\label{oc:lem:perm}
For every permutation $\sigma$ of $\{0,1,2\}$ and every ordered triple $t$,
\[\orient_3(t\circ\sigma)=\sgn(\sigma)\cdot\orient_3(t).\]
\end{lemma}
\begin{proof}
We follow the Lean proof, an induction over transpositions (Fact~\ref{oc:fact:sign}(S4)) for the
property $P(\sigma)$: ``$\orient_3(t\circ\sigma)=\sgn(\sigma)\orient_3(t)$ for \emph{every} triple
$t$''. $P(\mathrm{id})$ holds since $t\circ\mathrm{id}=t$ and $\sgn(\mathrm{id})=1$. Assume $P(f)$ and
$x\ne y$, and let $t$ be any triple. Since $(\mathrm{swap}(x,y)\,f)(i)=\mathrm{swap}(x,y)(f(i))$, we have
$t\circ(\mathrm{swap}(x,y)\,f)=(t\circ\mathrm{swap}(x,y))\circ f$. By $P(f)$ for the triple
$t\circ\mathrm{swap}(x,y)$ and by Lemma~\ref{oc:lem:swap},
\[\orient_3\bigl(t\circ(\mathrm{swap}(x,y)\,f)\bigr)=\sgn(f)\,\orient_3(t\circ\mathrm{swap}(x,y))
 =-\sgn(f)\,\orient_3(t),\]
and $\sgn(\mathrm{swap}(x,y)\,f)=\sgn(\mathrm{swap}(x,y))\sgn(f)=-\sgn(f)$ by (S2) and (S3). So
$P(\mathrm{swap}(x,y)\,f)$ holds, and (S4) gives $P(\sigma)$ for all $\sigma$.
\end{proof}

\begin{lemma}[rotations and permutations together; \leanref{ErdosSzekeres.OrbitConsistency.orient3_rot_perm}{LEANEXT/orbit/OrbitConsistency.lean:194}]\label{oc:lem:rot-perm}
Let $c\in\R^2$, $m\in\N$, $g\in\mathrm{RotComp}(c,m)$, $\sigma$ a permutation of $\{0,1,2\}$, and $t,t'$
ordered triples with $t'=g\circ t\circ\sigma$. Then $\orient_3(t')=\sgn(\sigma)\cdot\orient_3(t)$.
\end{lemma}
\begin{proof}
By Lemma~\ref{oc:lem:rotcomp-pres}(b), $g$ preserves orientation, so
$\orient_3(t')=\orient\bigl(g(t(\sigma(0))),g(t(\sigma(1))),g(t(\sigma(2)))\bigr)
=\orient\bigl(t(\sigma(0)),t(\sigma(1)),t(\sigma(2))\bigr)=\orient_3(t\circ\sigma)$, and
$\orient_3(t\circ\sigma)=\sgn(\sigma)\orient_3(t)$ by Lemma~\ref{oc:lem:perm}.
\end{proof}

\begin{lemma}[two representations of one triple; \leanref{ErdosSzekeres.OrbitConsistency.sign_eq_of_two_reps}{LEANEXT/orbit/OrbitConsistency.lean:204}, \leanref{ErdosSzekeres.OrbitConsistency.orient3_of_common_rep}{LEANEXT/orbit/OrbitConsistency.lean:216}]\label{oc:lem:two-reps}
Let $c\in\R^2$, $m\in\N$ and $g_1,g_2\in\mathrm{RotComp}(c,m)$, and let $\sigma_1,\sigma_2$ be permutations
of $\{0,1,2\}$.
\begin{enumerate}[label=(\alph*)]
\item (\lean{sign_eq_of_two_reps}) If $t,t'$ are ordered triples with $t'=g_1\circ t\circ\sigma_1$ and
  $t'=g_2\circ t\circ\sigma_2$, and $\orient_3(t)\ne0$, then $\sgn\sigma_1=\sgn\sigma_2$ in $\Z^\times$.
\item (\lean{orient3_of_common_rep}) If $t_1,t_2,C$ are ordered triples with $C=g_1\circ t_1\circ\sigma_1$
  and $C=g_2\circ t_2\circ\sigma_2$, then $\sgn(\sigma_1)\orient_3(t_1)=\sgn(\sigma_2)\orient_3(t_2)$ in
  $\R$.
\end{enumerate}
\end{lemma}
\begin{proof}
(a) Lemma~\ref{oc:lem:rot-perm}, applied to both representations, gives
$\sgn(\sigma_1)\orient_3(t)=\orient_3(t')=\sgn(\sigma_2)\orient_3(t)$. Since $\orient_3(t)\ne0$, the real
numbers $\sgn\sigma_1$ and $\sgn\sigma_2$ are equal, and so are the units (Fact~\ref{oc:fact:sign}: the
map $\Z^\times\to\R$ is injective). (b) By Lemma~\ref{oc:lem:rot-perm} both sides equal $\orient_3(C)$.
\end{proof}

\begin{definition}[the literal of the encoder; \leanref{ErdosSzekeres.OrbitConsistency.litVal}{LEANEXT/orbit/OrbitConsistency.lean:223}]\label{oc:def:litval}
For $\varepsilon\in\Z^\times$ and an ordered triple $C$, $\operatorname{litVal}(\varepsilon,C)$ is the
proposition ``$0<\orient_3(C)$'' if $\varepsilon=1$, and the proposition ``not $0<\orient_3(C)$''
otherwise (that is, if $\varepsilon=-1$). It models the SAT literal ``variable of the triple $C$, negated
when the parity $\varepsilon$ is odd'', where the variable of $C$ stands for $0<\orient_3(C)$.
\end{definition}

\begin{theorem}[literal soundness; \leanref{ErdosSzekeres.OrbitConsistency.literal_sound}{LEANEXT/orbit/OrbitConsistency.lean:228}]\label{oc:thm:literal}
Let $c\in\R^2$, $m\in\N$, $g\in\mathrm{RotComp}(c,m)$, $\sigma$ a permutation of $\{0,1,2\}$, and $t,C$
ordered triples with $C=g\circ t\circ\sigma$ and $\orient_3(t)\ne0$. Then
\[0<\orient_3(t)\iff\operatorname{litVal}(\sgn\sigma,\,C).\]
\end{theorem}
\begin{proof}
By Lemma~\ref{oc:lem:rot-perm}, $\orient_3(C)=\sgn(\sigma)\orient_3(t)$. By (S1), $\sgn\sigma=1$ or
$\sgn\sigma=-1$.
If $\sgn\sigma=1$, then $\operatorname{litVal}(\sgn\sigma,C)$ is ``$0<\orient_3(C)$'' and
$\orient_3(C)=\orient_3(t)$, so the two sides coincide.
If $\sgn\sigma=-1$, then $\sgn\sigma\ne1$, so $\operatorname{litVal}(\sgn\sigma,C)$ is ``not
$0<\orient_3(C)$'', and $\orient_3(C)=-\orient_3(t)$. If $0<\orient_3(t)$, then $-\orient_3(t)<0$, so the
literal holds. Conversely, if the literal holds, then $-\orient_3(t)\le0$, i.e.\ $\orient_3(t)\ge0$, and
$\orient_3(t)\ne0$ gives $0<\orient_3(t)$.
\end{proof}

\begin{theorem}[literal soundness, label-level form modelled on the encoder; \leanref{ErdosSzekeres.OrbitConsistency.encoder_literal_sound}{LEANEXT/orbit/OrbitConsistency.lean:253}]\label{oc:thm:encoder}
Let $\iota$ be any type (of labels), $P:\iota\to\R^2$ and $r:\iota\to\iota$ maps, $c\in\R^2$, $m\in\N$, and
$\rho=R_{c,2\pi/m}$. Assume
\[P(r(i))=\rho(P(i))\qquad\text{for every } i\in\iota .\]
Let $T,T_0,C:\{0,1,2\}\to\iota$, let $\sigma_1,\sigma_2$ be permutations of $\{0,1,2\}$, and $s\in\N$, with
\[T_0=T\circ\sigma_1\qquad\text{and}\qquad C=r^{[s]}\circ T_0\circ\sigma_2 .\]
If $\orient_3(P\circ T)\ne0$, then
\[0<\orient_3(P\circ T)\iff\operatorname{litVal}\bigl(\sgn\sigma_1\cdot\sgn\sigma_2,\;P\circ C\bigr).\]
Nothing else is assumed: $P$ need not be injective, and $m$ is arbitrary (for $m=0$,
$\rho=R_{c,0}=\mathrm{id}$). Distinctness of the labels $T(0),T(1),T(2)$ is not a separate hypothesis; it
follows from $\orient_3(P\circ T)\ne0$, because $\orient$ vanishes when two of its arguments are equal
(Lemma~\ref{fd:lem:orient}(e)).
\end{theorem}
\begin{proof}
\emph{Step 1: $P\circ r^{[s]}=\rho^{[s]}\circ P$.} Induction on $s$: for $s=0$ both sides are $P$; if it
holds for $s$, then for every label $i$,
$P(r^{[s+1]}(i))=P(r(r^{[s]}(i)))=\rho(P(r^{[s]}(i)))=\rho(\rho^{[s]}(P(i)))=\rho^{[s+1]}(P(i))$.

\emph{Step 2: $P\circ C=\rho^{[s]}\circ(P\circ T)\circ(\sigma_1\sigma_2)$.} For $i\in\{0,1,2\}$,
$P(C(i))=P\bigl(r^{[s]}(T_0(\sigma_2(i)))\bigr)=\rho^{[s]}\bigl(P(T(\sigma_1(\sigma_2(i))))\bigr)$ by
Step~1 and $T_0=T\circ\sigma_1$, and $\sigma_1(\sigma_2(i))=(\sigma_1\sigma_2)(i)$.

\emph{Step 3.} $\rho^{[s]}\in\mathrm{RotComp}(c,m)$ by Lemma~\ref{oc:lem:rotcomp-pres}(c). Theorem~\ref{oc:thm:literal}
with $g=\rho^{[s]}$, $\sigma=\sigma_1\sigma_2$, $t=P\circ T$ and the triple $P\circ C$ (Step~2), and
$\orient_3(P\circ T)\ne0$, gives $0<\orient_3(P\circ T)\iff\operatorname{litVal}(\sgn(\sigma_1\sigma_2),P\circ
C)$. Finally $\sgn(\sigma_1\sigma_2)=\sgn\sigma_1\cdot\sgn\sigma_2$ by (S2).
\end{proof}

\begin{remark}[parity by counting inversions is not modelled in Lean]\label{oc:rem:parity}
The encoders compute parities by counting inversions, in Python. For a list $a$ of three labels let
$\mathrm{inv}(a)$ be the number of pairs $i<j$ of positions with $a(i)>a(j)$. The function
\texttt{sortpar} of \dirn{C7UPPER/gen.py} (line 108) bubble-sorts $a$ and flips the parity at each
exchange of adjacent entries; each exchange removes exactly one inversion and the sorted list has none,
so it returns $(-1)^{\mathrm{inv}(a)}$. It is applied to the ordered label triple (line 120) and to the
rotations of the sorted triple (line 78); these give the parities that play the roles of
$\sgn\sigma_1$ and $\sgn\sigma_2$. The function \texttt{perm\_sign} of \dirn{C6UPPER/encb/gen_b2.py}
(line 56) counts the inversions of its argument and also returns $(-1)^{\mathrm{inv}(a)}$; it is
applied (lines 72--73) to the list $\sigma(0)\sigma(1)\sigma(2)$ of the permutation $\sigma$ that sorts
the rotated triple, so this encoder corresponds to $\sigma_1=\mathrm{id}$ and $\sigma_2=\sigma$.
Theorem~\ref{oc:thm:encoder} is stated with Mathlib's $\sgn$ instead. The two agree on lists of three
distinct labels, by a six-case check. Write a permutation $\tau$ of $\{0,1,2\}$ as the list $\tau(0)\tau(1)\tau(2)$; so
$\mathrm{inv}(\tau)$ is the number of pairs $i<j$ with $\tau(i)>\tau(j)$. If $\sigma$ sorts a list $a$ of
three distinct labels, i.e.\ $a\circ\sigma$ is increasing, then $a(i)>a(j)$ iff
$\sigma^{-1}(i)>\sigma^{-1}(j)$, so $\mathrm{inv}(a)=\mathrm{inv}(\sigma^{-1})$. For the six
permutations $\tau$ one has
\begin{center}
\begin{tabular}{cccc}
\toprule
$\tau$ & $\tau$ as a product of transpositions & $\mathrm{inv}(\tau)$ & $\sgn\tau$ by (S2), (S3)\\
\midrule
$012$ & $\mathrm{id}$ & $0$ & $1$\\
$102$ & $\mathrm{swap}(0,1)$ & $1$ & $-1$\\
$021$ & $\mathrm{swap}(1,2)$ & $1$ & $-1$\\
$210$ & $\mathrm{swap}(0,2)$ & $3$ & $-1$\\
$120$ & $\mathrm{swap}(0,1)\,\mathrm{swap}(1,2)$ & $2$ & $1$\\
$201$ & $\mathrm{swap}(1,2)\,\mathrm{swap}(0,1)$ & $2$ & $1$\\
\bottomrule
\end{tabular}
\end{center}
so $(-1)^{\mathrm{inv}(\tau)}=\sgn\tau$ for every $\tau$. Hence \texttt{perm\_sign}, applied to the
list of a permutation $\sigma$, returns $(-1)^{\mathrm{inv}(\sigma)}=\sgn\sigma$; and \texttt{sortpar},
applied to a list $a$ of distinct labels that $\sigma$ sorts, returns
$(-1)^{\mathrm{inv}(a)}=(-1)^{\mathrm{inv}(\sigma^{-1})}=\sgn(\sigma^{-1})=\sgn(\sigma)$ by (S2). This
six-case argument, the reading of the code above and the correspondence between the generator code and
the data $T,T_0,C,\sigma_1,\sigma_2,s$ of Theorem~\ref{oc:thm:encoder} are \emph{not} part of any Lean
statement.
The Python script \dirn{LEANEXT/orbit/check_literals.py} checked (lines 36--39) that \texttt{sortpar}
sorts each of the six orders of $(0,1,2)$ and returns the same parity as the script's own inversion
count \texttt{perm\_sign} (lines 25--31); this compares two inversion counts, not \texttt{sortpar} with
Mathlib's $\sgn$. The same script compared the literal code of \dirn{C7UPPER/gen.py} with
floating-point orientations on 368,814 ordered triples with no mismatch
(\dirn{LEANEXT/orbit/logs/check_literals.log}, lines 1--5). These Python checks are evidence, not
proof. The design notes of the ENCBLEAN route record, in their entry for this theorem, the mismatch
between $\sgn$ and the inversion count and that a six-case bridge would be needed, and give the verdict
``Not used'' (\dirn{ENCBLEAN/DESIGN.md:625}).
\end{remark}

\begin{remark}[how the theorems are used]\label{oc:rem:literal-use}
As described in the doc comment of \lean{encoder_literal_sound}, the encoder keeps one Boolean variable
$V(C)$ per canonical representative $C$ of an orbit of ordered label triples under the label rotation $r$
and the permutations of the positions; $V(C)$ is meant to be $[0<\orient_3(P\circ C)]$. For an ordered
triple $T$ it sorts $T$ ($T_0=T\circ\sigma_1$), maps the sorted triple to the representative
($C=r^{[s]}\circ T_0\circ\sigma_2$), and uses $V(C)$, negated when $\sgn\sigma_1\sgn\sigma_2=-1$. Under the semiconjugacy $P\circ r=\rho\circ P$ and
$\orient_3(P\circ T)\ne0$ (which holds when the three points are distinct points of a set in general
position, by Lemma~\ref{fd:lem:collinear}), Theorem~\ref{oc:thm:encoder} says that this literal is true
exactly when $\orient(P(T(0)),P(T(1)),P(T(2)))>0$, whichever representation is chosen;
Lemma~\ref{oc:lem:two-reps} shows directly that the sign does not depend on the representation.
\end{remark}

\subsection{Even order: antipodes, rays and phases}\label{oc:sub:even}

\begin{lemma}[antipodes; \leanref{ErdosSzekeres.OrbitConsistency.antipode_mem_of_even}{LEANEXT/orbit/OrbitConsistency.lean:273}, \leanref{ErdosSzekeres.LayerLemma.antipode_mem_of_even}{LEANEXT/merged/Merged.lean:1160}]\label{oc:lem:antipode}
Let $S$ be a finite set of points, $c\in\R^2$ and $m\in\N$ with $m$ even and $0<m$, and assume that
$\rho=R_{c,2\pi/m}$ maps $S$ into $S$. If $x\in S$, then the point $(2c_0-x_0,\;2c_1-x_1)=2c-x$ belongs
to $S$.
\end{lemma}
\begin{proof}
By Lemma~\ref{oc:lem:iter-mem}, $\rho^{[m/2]}(x)\in S$, and $\rho^{[m/2]}(x)=(2c_0-x_0,2c_1-x_1)$ by
Corollary~\ref{oc:cor:period}.
\end{proof}

\begin{lemma}[no two points on one open ray from the centre; \leanref{ErdosSzekeres.OrbitConsistency.same_ray_false}{LEANEXT/orbit/OrbitConsistency.lean:295}, \leanref{ErdosSzekeres.LayerLemma.same_ray_false}{LEANEXT/merged/Merged.lean:1181}]\label{oc:lem:same-ray}
Let $S$ be a finite set of points, $c\in\R^2$ and $m\in\N$ with $m$ even and $0<m$; assume that
$\rho=R_{c,2\pi/m}$ maps $S$ into $S$ and that $S$ is in general position. Let $x,y\in S$ with $x\ne y$,
and $l\in\R$ with $l>0$. Then
\[y-c\ne l\,(x-c).\]
(In Lean the equation $y-c=l\,(x-c)$ is a hypothesis and the conclusion is \lean{False}; this is the
same statement.) The condition $c\notin S$ is not a hypothesis; it follows from the others by
Lemma~\ref{lg:lem:centre}, since $x\ne y$ gives a point of $S$ other than $c$, but the proof below does
not use it.
\end{lemma}
\begin{proof}
Suppose $y-c=l(x-c)$, i.e.\ $y=c+l(x-c)$. Let $q=2c-x$; then $q\in S$ by Lemma~\ref{oc:lem:antipode}.
\begin{enumerate}[label=(\roman*)]
\item $x\ne c$: if $x=c$, then $y=c+l\cdot0=c=x$, contrary to $x\ne y$.
\item $x\ne q$: $x=q$ means $2(x-c)=0$, i.e.\ $x=c$, excluded by (i).
\item $y\ne q$: $y=q$ means $c+l(x-c)=2c-x$, i.e.\ $(l+1)(x-c)=0$; as $l+1>0$, this gives $x=c$, excluded by
  (i).
\item $\orient(x,y,q)=0$: $y-x=(l-1)(x-c)$ and $q-x=-2(x-c)$, so
  $\orient(x,y,q)=\det\bigl(y-x\mid q-x\bigr)=-2(l-1)\det\bigl(x-c\mid x-c\bigr)=0$.
\end{enumerate}
Thus $x,y,q$ are pairwise distinct points of $S$ with $\orient(x,y,q)=0$, which contradicts general
position by Lemma~\ref{fd:lem:collinear} (\lean{generalPosition_iff_orient_ne_zero},
\dirn{ErdosSzekeres/OrientationGeometry.lean:66}).
\end{proof}

\begin{remark}[the hypothesis $l>0$]\label{oc:rem:ray}
The proof uses $l>0$ only through $l+1\ne0$. The hypothesis cannot be replaced by $l\ne0$: for any
$x\in S$ with $x\ne c$, the antipode $y=2c-x$ lies in $S$ (Lemma~\ref{oc:lem:antipode}), $y\ne x$, and
$y-c=(-1)(x-c)$. Such $S$ exist, for example $m=2$ and $S=\{x,2c-x\}$ with $x\ne c$ (general position
holds because $S$ has only two points).
\end{remark}

\begin{lemma}[rotating and shifting polar points; \leanref{ErdosSzekeres.OrbitConsistency.rotAbout_polar}{LEANEXT/orbit/OrbitConsistency.lean:355}, \leanref{ErdosSzekeres.OrbitConsistency.polar_add_int}{LEANEXT/orbit/OrbitConsistency.lean:383}, \leanref{ErdosSzekeres.LayerLemma.rotAbout_polar}{LEANEXT/merged/Merged.lean:1241}, \leanref{ErdosSzekeres.LayerLemma.polar_add_int}{LEANEXT/merged/Merged.lean:1269}]\label{oc:lem:polar-rot}
A rotation about $c$ adds its angle to the angle of a polar point about $c$, and polar points are
$2\pi$-periodic in the angle:
\begin{enumerate}[label=(\alph*)]
\item (\lean{rotAbout_polar}) For all $c\in\R^2$ and $r,\theta,\varphi\in\R$:
  $R_{c,\varphi}(\operatorname{polar}(c,r,\theta))=\operatorname{polar}(c,r,\theta+\varphi)$.
\item (\lean{polar_add_int}) For all $c\in\R^2$, $r,\theta\in\R$ and $n\in\Z$:
  $\operatorname{polar}(c,r,\theta+n\cdot2\pi)=\operatorname{polar}(c,r,\theta)$.
\end{enumerate}
\end{lemma}
\begin{proof}
(a) $\operatorname{polar}(c,r,\theta)-c=(r\cos\theta,r\sin\theta)$, so $R_{c,\varphi}$ maps the point to
$c+\bigl(r(\cos\varphi\cos\theta-\sin\varphi\sin\theta),\;r(\sin\varphi\cos\theta+\cos\varphi\sin\theta)\bigr)
=c+\bigl(r\cos(\theta+\varphi),\;r\sin(\theta+\varphi)\bigr)$ by the addition formulas.
(b) $\cos$ and $\sin$ have period $2\pi$, so $\cos(\theta+2\pi n)=\cos\theta$ and
$\sin(\theta+2\pi n)=\sin\theta$ for $n\in\Z$.
\end{proof}

\begin{lemma}[equal phases modulo $2\pi/m$, natural form; \leanref{ErdosSzekeres.OrbitConsistency.eq_iterate_of_phase_eq}{LEANEXT/orbit/OrbitConsistency.lean:363}, \leanref{ErdosSzekeres.LayerLemma.eq_iterate_of_phase_eq}{LEANEXT/merged/Merged.lean:1249}]\label{oc:lem:phase-nat}
Let $S$ be a finite set of points, $c\in\R^2$ and $m\in\N$ with $m$ even and $0<m$; assume that
$\rho=R_{c,2\pi/m}$ maps $S$ into $S$ and that $S$ is in general position. Let $r,r',\alpha,\beta\in\R$
with $r>0$, $r'>0$, $\operatorname{polar}(c,r,\alpha)\in S$ and $\operatorname{polar}(c,r',\beta)\in S$,
and let $k\in\N$ with $\beta=\alpha+k\cdot(2\pi/m)$. Then
\[\operatorname{polar}(c,r',\beta)=\rho^{[k]}\bigl(\operatorname{polar}(c,r,\alpha)\bigr).\]
\end{lemma}
\begin{proof}
Let $z=\rho^{[k]}(\operatorname{polar}(c,r,\alpha))$ and $y=\operatorname{polar}(c,r',\beta)$. By
Lemma~\ref{oc:lem:iterate}, $\rho^{[k]}=R_{c,k(2\pi/m)}$, so by Lemma~\ref{oc:lem:polar-rot}(a),
$z=\operatorname{polar}(c,r,\alpha+k(2\pi/m))=\operatorname{polar}(c,r,\beta)$; and $z\in S$ by
Lemma~\ref{oc:lem:iter-mem}. Suppose $y\ne z$. Then
\[y-c=(r'\cos\beta,\,r'\sin\beta)=\tfrac{r'}{r}\,(r\cos\beta,\,r\sin\beta)=\tfrac{r'}{r}\,(z-c),\]
with $r'/r>0$. So $z,y\in S$, $z\ne y$ and $y-c=l(z-c)$ with $l=r'/r>0$, which
Lemma~\ref{oc:lem:same-ray} (applied with $x:=z$) excludes. Hence $y=z$.
\end{proof}

\begin{lemma}[equal phases modulo $2\pi/m$, integer form; \leanref{ErdosSzekeres.OrbitConsistency.eq_iterate_of_phase_eq_int}{LEANEXT/orbit/OrbitConsistency.lean:389}, \leanref{ErdosSzekeres.LayerLemma.eq_iterate_of_phase_eq_int}{LEANEXT/merged/Merged.lean:1275}]\label{oc:lem:phase-int}
Let $S$ be a finite set of points, $c\in\R^2$ and $m\in\N$ with $m$ even and $0<m$; assume that
$\rho=R_{c,2\pi/m}$ maps $S$ into $S$ and that $S$ is in general position. Let $r,r',\alpha,\beta\in\R$
with $r>0$, $r'>0$, $\operatorname{polar}(c,r,\alpha)\in S$ and $\operatorname{polar}(c,r',\beta)\in S$,
and let $k\in\Z$ with $\beta=\alpha+k\cdot(2\pi/m)$. Then there is $j\in\N$ with
\[\operatorname{polar}(c,r',\beta)=\rho^{[j]}\bigl(\operatorname{polar}(c,r,\alpha)\bigr).\]
(The proof takes for $j$ the remainder of $k$ modulo $m$.)
\end{lemma}
\begin{proof}
Since $m>0$, Euclidean division of $k$ by $m$ gives $q=\lfloor k/m\rfloor\in\Z$ and a remainder
$j=k-mq\in\{0,1,\dots,m-1\}$, which we regard as a natural number (Lean: \texttt{(k \% m).toNat}, with
$0\le k\bmod m$ by \lean{Int.emod_nonneg} and $k=(k\bmod m)+m\lfloor k/m\rfloor$ by
\lean{Int.emod_add_mul_ediv}). Then $k=j+mq$, and, as $m\ne0$,
$k\cdot\frac{2\pi}{m}=j\cdot\frac{2\pi}{m}+q\cdot2\pi$. Put $\beta'=\alpha+j\cdot(2\pi/m)$; then
$\beta=\beta'+q\cdot2\pi$, so $\operatorname{polar}(c,r',\beta)=\operatorname{polar}(c,r',\beta')$ by
Lemma~\ref{oc:lem:polar-rot}(b). In particular $\operatorname{polar}(c,r',\beta')\in S$, and
Lemma~\ref{oc:lem:phase-nat} (with $\beta'$ and the natural number $j$) gives
$\operatorname{polar}(c,r',\beta')=\rho^{[j]}(\operatorname{polar}(c,r,\alpha))$. Hence the claim.
\end{proof}

\begin{theorem}[points outside an orbit have different phases; \leanref{ErdosSzekeres.OrbitConsistency.phase_ne_of_not_mem_orbit}{LEANEXT/orbit/OrbitConsistency.lean:420}, \leanref{ErdosSzekeres.LayerLemma.phase_ne_of_not_mem_orbit}{LEANEXT/merged/Merged.lean:1304}]\label{oc:thm:phase-orbit}
Let $S$ be a finite set of points, $c\in\R^2$ and $m\in\N$ with $m$ even and $0<m$; assume that
$\rho=R_{c,2\pi/m}$ maps $S$ into $S$ and that $S$ is in general position. Let $r,r',\alpha,\beta\in\R$
with $r>0$, $r'>0$, $\operatorname{polar}(c,r,\alpha)\in S$, $\operatorname{polar}(c,r',\beta)\in S$, and
\[\operatorname{polar}(c,r',\beta)\ne\rho^{[j]}\bigl(\operatorname{polar}(c,r,\alpha)\bigr)\qquad\text{for
every } j\in\N .\]
Then $\beta\ne\alpha+k\cdot(2\pi/m)$ for every $k\in\Z$.
\end{theorem}
\begin{proof}
Let $k\in\Z$ and suppose $\beta=\alpha+k(2\pi/m)$. Lemma~\ref{oc:lem:phase-int} gives $j\in\N$ with
$\operatorname{polar}(c,r',\beta)=\rho^{[j]}(\operatorname{polar}(c,r,\alpha))$, contrary to the
hypothesis.
\end{proof}

\begin{remark}[what is and is not claimed]\label{oc:rem:phase-orbit}
The doc comment of the \dirn{OrbitConsistency.lean} version of \lean{phase_ne_of_not_mem_orbit} adds that,
with the encoder's labelling (layer~$0$ has a point of direction~$0$), the phases of distinct layers differ
and the phases of the layers $i\ge1$ are nonzero. That reading needs facts that are not hypotheses of
Theorem~\ref{oc:thm:phase-orbit}: that a point of one layer is not in the $\rho$-orbit of a point of
another layer (this is what the proof of Theorem~\ref{oc:thm:layers-phase} below establishes, from
Lemmas~\ref{oc:lem:disjoint} and~\ref{oc:lem:closed}), and the normalisation that layer~$0$ contains a
point of direction~$0$, together with the definition of the phases relative to it, which is not part of
any Lean statement of this section.
\end{remark}

\subsection{Distinct layers: rays and phases}\label{oc:sub:layers}

We use the hull vertices $\hv(X)=\{x\in X: x\notin\conv(X\setminus\{x\})\}$ of a finite set $X$
(Definition~\ref{ll:def:hv}) and the peels $\peel(S,0)=S$, $\peel(S,n+1)=\peel(S,n)\setminus\hv(\peel(S,n))$
(Definition~\ref{ll:def:peel}); the $n$-th layer of $S$ is $\hv(\peel(S,n))$. We use three facts from the layer-lemma
section:
\begin{itemize}
\item[(L1)] $\hv(X)\subseteq X$ (by definition), and $\peel(S,n')\subseteq\peel(S,n)$ for $n\le n'$
  (Lemma~\ref{ll:lem:antitone}); in particular $\peel(S,n)\subseteq\peel(S,0)=S$.
\item[(L2)] (Lemma~\ref{ll:lem:inv-peel}) If $0<m$, $c\notin S$ and $\rho=R_{c,2\pi/m}$ maps $S$ into
  $S$, then $\rho$ maps every $\peel(S,n)$ into itself.
\item[(L3)] (Theorem~\ref{ll:thm:core-rot}, \lean{hullVerts_eq_orbitFS_rotAbout},
  \dirn{LAYERLEAN/LayerLemma.lean:364}) If $0<m$, $K\le2m$, $X$ is a finite set of points with $c\notin X$,
  $\rho$ maps $X$ into $X$, $X$ has no $K$ points in convex position, and $x\in\hv(X)$, then
  $\hv(X)=\{\rho^{[a]}(x):0\le a<m\}$ (Lean: \lean{orbitFS} of $\rho$, $m$ and $x$).
\end{itemize}
``$S$ has no $K$ points in convex position'' is $\neg\,$\texttt{HasConvexNGon K S} (Definition~\ref{fd:def:ngon}).

\begin{lemma}[distinct layers are disjoint; \leanref{ErdosSzekeres.LayerLemma.not_mem_hullVerts_peel_of_lt}{LEANEXT/merged/Merged.lean:1316}, \leanref{ErdosSzekeres.LayerLemma.not_mem_hullVerts_peel_of_ne}{LEANEXT/merged/Merged.lean:1322}]\label{oc:lem:disjoint}
Let $S$ be a finite set of points (no other hypothesis).
\begin{enumerate}[label=(\alph*)]
\item If $i<j$ and $y\in\hv(\peel(S,j))$, then $y\notin\hv(\peel(S,i))$.
\item If $i\ne j$, then no point $y$ satisfies both $y\in\hv(\peel(S,i))$ and $y\in\hv(\peel(S,j))$.
\end{enumerate}
\end{lemma}
\begin{proof}
(a) By (L1) and $i+1\le j$: $y\in\hv(\peel(S,j))\subseteq\peel(S,j)\subseteq\peel(S,i+1)
=\peel(S,i)\setminus\hv(\peel(S,i))$, so $y\notin\hv(\peel(S,i))$. (b) If $i<j$, apply (a) to $y$ as a
point of layer $j$; if $j<i$, apply (a) with $i$ and $j$ exchanged.
\end{proof}

\begin{lemma}[each layer is closed under the rotation; \leanref{ErdosSzekeres.LayerLemma.iterate_mem_hullVerts_peel}{LEANEXT/merged/Merged.lean:1330}]\label{oc:lem:closed}
Let $S$ be a finite set of points, $c\in\R^2$, and $m,K\in\N$ with $0<m$ and $K\le2m$; assume $c\notin S$,
that $\rho=R_{c,2\pi/m}$ maps $S$ into $S$, and that $S$ has no $K$ points in convex position. If
$i\in\N$ and $x\in\hv(\peel(S,i))$, then $\rho^{[k]}(x)\in\hv(\peel(S,i))$ for every $k\in\N$.
\end{lemma}
\begin{proof}
Put $X=\peel(S,i)$. By (L1), $X\subseteq S$; hence $c\notin X$, and $X$ has no $K$ points in convex
position (such points would be $K$ points of $S$ in convex position; Lemma~\ref{fd:lem:mono}(c)). By (L2),
$\rho$ maps $X$ into $X$. By (L3), $\hv(X)=O:=\{\rho^{[a]}(x):0\le a<m\}$. The set $O$ is mapped into
itself by $\rho$: for $a+1<m$, $\rho(\rho^{[a]}(x))=\rho^{[a+1]}(x)\in O$, and for $a=m-1$,
$\rho(\rho^{[m-1]}(x))=\rho^{[m]}(x)=x=\rho^{[0]}(x)\in O$, because $\rho^{[m]}=\mathrm{id}$
(Corollary~\ref{oc:cor:period}). Since $x=\rho^{[0]}(x)\in O$ (as $0<m$), Lemma~\ref{oc:lem:iter-mem}
gives $\rho^{[k]}(x)\in O=\hv(X)$ for every $k$. (The Lean proof obtains $\rho^{[m]}(x)=x$ from
$x\ne c$ and the minimal period, Lemma~\ref{ll:lem:minperiod}; the direct identity
$\rho^{[m]}=\mathrm{id}$ used here needs neither.)
\end{proof}

\begin{theorem}[distinct layers never share a ray through a rotated point; \leanref{ErdosSzekeres.LayerLemma.not_same_ray_of_distinct_layers}{LEANEXT/merged/Merged.lean:1350}]\label{oc:thm:layers-ray}
Let $S$ be a finite set of points, $c\in\R^2$, and $m,K\in\N$ with $2\le m$, $m$ even and $K\le2m$; assume
$c\notin S$, that $\rho=R_{c,2\pi/m}$ maps $S$ into $S$, that $S$ is in general position, and that $S$
has no $K$ points in convex position. Let $i,j\in\N$ with $i\ne j$, $x\in\hv(\peel(S,i))$,
$y\in\hv(\peel(S,j))$, $k\in\N$ and $l\in\R$ with $l>0$. Then
\[y-c\ne l\,\bigl(\rho^{[k]}(x)-c\bigr).\]
(In Lean the equation is a hypothesis and the conclusion is \lean{False}.)
\end{theorem}
\begin{proof}
We have $0<m$. Put $x'=\rho^{[k]}(x)$; then $x'\in\hv(\peel(S,i))$ by Lemma~\ref{oc:lem:closed}. Also
$x'\ne y$: otherwise $y$ would lie in $\hv(\peel(S,i))$ and in $\hv(\peel(S,j))$, contrary to
Lemma~\ref{oc:lem:disjoint}(b). By (L1), $x',y\in S$. If $y-c=l(x'-c)$, then
Lemma~\ref{oc:lem:same-ray} (with $x:=x'$, using $m$ even, $0<m$, invariance and general position) gives a
contradiction.
\end{proof}

\begin{theorem}[distinct layers have phases incongruent modulo $2\pi/m$; \leanref{ErdosSzekeres.LayerLemma.phase_ne_of_distinct_layers}{LEANEXT/merged/Merged.lean:1367}]\label{oc:thm:layers-phase}
Let $S$ be a finite set of points, $c\in\R^2$, and $m,K\in\N$ with $2\le m$, $m$ even and $K\le2m$; assume
$c\notin S$, that $\rho=R_{c,2\pi/m}$ maps $S$ into $S$, that $S$ is in general position, and that $S$
has no $K$ points in convex position. Let $i,j\in\N$ with $i\ne j$, and $r,r',\alpha,\beta\in\R$ with
$r>0$, $r'>0$,
\[\operatorname{polar}(c,r,\alpha)\in\hv(\peel(S,i))\qquad\text{and}\qquad
  \operatorname{polar}(c,r',\beta)\in\hv(\peel(S,j)).\]
Then $\beta-\alpha\ne k\cdot(2\pi/m)$ for every $k\in\Z$.
\end{theorem}
\begin{proof}
We apply Theorem~\ref{oc:thm:phase-orbit}. Its hypotheses hold: $m$ is even and $0<m$; $\rho$ maps $S$
into $S$; $S$ is in general position; $r,r'>0$; both points lie in $S$ by (L1). For the orbit
hypothesis, suppose $\operatorname{polar}(c,r',\beta)=\rho^{[n]}(\operatorname{polar}(c,r,\alpha))$ for
some $n\in\N$; then this point lies in $\hv(\peel(S,i))$ by Lemma~\ref{oc:lem:closed}, and also in
$\hv(\peel(S,j))$, contrary to Lemma~\ref{oc:lem:disjoint}(b). Now let $k\in\Z$ and suppose
$\beta-\alpha=k(2\pi/m)$. Then $\beta=\alpha+k(2\pi/m)$, which Theorem~\ref{oc:thm:phase-orbit} excludes.
\end{proof}

\begin{remark}[on the hypotheses and the conclusion]\label{oc:rem:layers}
(1) The hypothesis $c\notin S$ of Theorems~\ref{oc:thm:layers-ray} and~\ref{oc:thm:layers-phase} follows
from the others, although Lean keeps it as a hypothesis: the two given points are distinct by
Lemma~\ref{oc:lem:disjoint}(b), so $S$ contains a point $p\ne c$, and Lemma~\ref{lg:lem:centre}
(\lean{centre_not_mem_of_even}: $0<m$, $m$ even, $\rho$ maps $S$ into $S$, $S$ in general position,
$p\in S$, $p\ne c$) gives $c\notin S$.
(2) The points of Theorem~\ref{oc:thm:layers-phase} are given in polar form with positive radii; no
statement of this section produces such a form from a point.
(3) The modulus $2\pi/m$ cannot be replaced by $\pi/m$. Take $m=4$, $c=(0,0)$, $K=8$ and
$S=\{(\pm2,0),(0,\pm2),(\pm\tfrac12,\pm\tfrac12)\}$, with four outer points $(\pm2,0),(0,\pm2)$ and four
inner points $(\pm\frac12,\pm\frac12)$. The rotation by $\pi/2$ maps $S$ onto itself and $c\notin S$.
$S$ is in general position: among any three points of $S$ two are of the same kind (outer or inner), so
it suffices that the line through two points of the same kind contains no third point of $S$. The lines
through two of the outer points are $x\pm y=\pm2$, $x=0$ and $y=0$, and the lines through two of the
inner points are $x=\pm\frac12$, $y=\pm\frac12$ and $x=\pm y$; no third point of $S$ lies on any of them,
since inner points have $\abs{x}=\abs{y}=\frac12$ and outer points have one coordinate $0$ and the other
$\pm2$. For $\varepsilon_1,\varepsilon_2\in\{1,-1\}$,
\[\bigl(\tfrac{\varepsilon_1}2,\tfrac{\varepsilon_2}2\bigr)=\tfrac38(2\varepsilon_1,0)+\tfrac38(0,2\varepsilon_2)
  +\tfrac18(-2\varepsilon_1,0)+\tfrac18(0,-2\varepsilon_2),\]
a convex combination of the outer points; hence no inner point is a hull vertex of $S$, $S$ is not in
convex position (Lemma~\ref{fd:lem:cp-hull}), and $S$ (which has $8$ points) has no $8$ points in convex
position. Each outer point is a hull vertex of $S$, since it is the only point of $S$
maximising one of the linear functions $\pm x$, $\pm y$ (easy direction of
Lemma~\ref{fd:lem:separation}); so layer~$0$ consists of the four outer points. Likewise each inner point
is the only point of $\peel(S,1)$ (the four inner points) maximising one of $\pm x\pm y$, so layer~$1$
consists of the inner points. With
$(2,0)=\operatorname{polar}(c,2,0)$ in layer~$0$ and
$(\frac12,\frac12)=\operatorname{polar}(c,\frac{\sqrt2}2,\frac\pi4)$ in layer~$1$, all hypotheses of
Theorem~\ref{oc:thm:layers-phase} hold and $\beta-\alpha=\frac\pi4=1\cdot\frac{\pi}{m}$.
\end{remark}

\subsection{Key order}\label{oc:sub:keys}

This subsection concerns only real numbers and integers. Throughout, intervals $[a,b)=\{y\in\R:a\le y<b\}$
are half-open: the left end is included and the right end is not.

\begin{definition}[reduction modulo $p$, Euclidean division, lexicographic order, strict monotonicity; Mathlib]\label{oc:def:keys}
We use the following notions, as defined in Mathlib and in the Lean core library.
\begin{enumerate}[label=(\alph*)]
\item Let $p>0$ and $a,b\in\R$. There is exactly one $n\in\Z$ with $b-np\in[a,a+p)$: for existence take
  $n=\lfloor(b-a)/p\rfloor$, so that $np\le b-a<np+p$; for uniqueness, if $b-np$ and $b-n'p$ both lie in
  $[a,a+p)$ then $\abs{(n'-n)p}<p$, so $n=n'$. Mathlib calls this integer \lean{toIcoDiv} and defines
  $\operatorname{toIcoMod}_p(a,b)=b-np$ (\texttt{toIcoMod hp a b}, where \lean{hp} is a proof of $0<p$).
  In \dirn{KeyOrder.lean}, $p=2\pi$ (with \lean{Real.two_pi_pos}) and $a=0$:
  $\operatorname{toIcoMod}_{2\pi}(0,x)$ is the representative of $x$ modulo $2\pi$ in $[0,2\pi)$.
\item For $N\in\Z$ with $N>0$ and $J\in\Z$, Lean's integer quotient \texttt{J / N} (\lean{Int.ediv}) is
  $\lfloor J/N\rfloor$, the floor of the rational number $J/N$, and Lean's remainder \texttt{J \% N}
  (\lean{Int.emod}) is $J\bmod N=J-N\lfloor J/N\rfloor$. Thus $0\le J\bmod N\le N-1$
  (\lean{Int.emod_nonneg}, \lean{Int.emod_lt_of_pos}) and $J=(J\bmod N)+N\lfloor J/N\rfloor$
  (\lean{Int.emod_add_mul_ediv}).
\item On $\Z\times\N$ the lexicographic order (Lean \lean{toLex} into $\Z\times_\ell\N$) is
  $(a,b)<_{\mathrm{lex}}(a',b')$ iff $a<a'$, or $a=a'$ and $b<b'$ (\lean{Prod.Lex.lt_iff}).
\item A map $\operatorname{rank}:\R\to\N$ is \emph{strictly increasing on} $S\subseteq\R$
  (\texttt{StrictMonoOn rank S}) if $\operatorname{rank}(a)<\operatorname{rank}(b)$ for all $a,b\in S$ with
  $a<b$.
\end{enumerate}
\end{definition}

\begin{lemma}[characterisation of $\operatorname{toIcoMod}$; written step (Mathlib \lean{toIcoMod_eq_iff})]\label{oc:lem:toico}
Let $p>0$ and $a,b,y\in\R$. Then $\operatorname{toIcoMod}_p(a,b)=y$ if and only if $a\le y<a+p$ and
$b=y+zp$ for some $z\in\Z$. In particular $\operatorname{toIcoMod}_p(a,b)\in[a,a+p)$.
\end{lemma}
\begin{proof}
Let $n$ be the integer of Definition~\ref{oc:def:keys}(a). If $y=\operatorname{toIcoMod}_p(a,b)=b-np$, then
$y\in[a,a+p)$ and $b=y+np$. Conversely, if $y\in[a,a+p)$ and $b=y+zp$ with $z\in\Z$, then
$b-zp=y\in[a,a+p)$, so $z=n$ by uniqueness, and $\operatorname{toIcoMod}_p(a,b)=b-zp=y$.
\end{proof}

\begin{lemma}[cells: order and equality; \leanref{KeyOrder.cell_lt_iff}{LEANEXT/keys/KeyOrder.lean:23}, \leanref{KeyOrder.cell_eq_iff}{LEANEXT/keys/KeyOrder.lean:45}]\label{oc:lem:cell}
Let $u\in\R$ with $u>0$, $J_1,J_2\in\Z$, and $\psi_1,\psi_2\in\R$ with $0\le\psi_1<u$ and $0\le\psi_2<u$.
Then
\begin{enumerate}[label=(\alph*)]
\item $J_1u+\psi_1<J_2u+\psi_2$ if and only if $J_1<J_2$, or $J_1=J_2$ and $\psi_1<\psi_2$;
\item $J_1u+\psi_1=J_2u+\psi_2$ if and only if $J_1=J_2$ and $\psi_1=\psi_2$.
\end{enumerate}
\end{lemma}
\begin{proof}
(a) ``If'': when $J_1<J_2$, we have $J_1+1\le J_2$, hence
$J_1u+\psi_1<J_1u+u=(J_1+1)u\le J_2u\le J_2u+\psi_2$, using $\psi_1<u$, $u>0$ and $\psi_2\ge0$; when
$J_1=J_2$ and $\psi_1<\psi_2$, add $J_1u$ to $\psi_1<\psi_2$.
``Only if'': assume $J_1u+\psi_1<J_2u+\psi_2$. If $J_1<J_2$ we are done. If $J_1=J_2$, cancelling $J_1u$
gives $\psi_1<\psi_2$. If $J_1>J_2$, then $J_1-J_2-1\ge0$ and
\[(J_1u+\psi_1)-(J_2u+\psi_2)=(J_1-J_2-1)\,u+\psi_1+(u-\psi_2)>0,\]
since the three terms are $\ge0$, $\ge0$ and $>0$; this contradicts the assumption.
(b) ``If'' is clear. ``Only if'': if $J_1u+\psi_1=J_2u+\psi_2$, then neither side is smaller than the
other, so by (a), applied in both orders, $J_1\not<J_2$ and $J_2\not<J_1$, hence $J_1=J_2$; and then
$\psi_1\not<\psi_2$ and $\psi_2\not<\psi_1$, hence $\psi_1=\psi_2$.
\end{proof}

\begin{lemma}[cells and ranks; \leanref{KeyOrder.cell_lt_iff_rank}{LEANEXT/keys/KeyOrder.lean:62}, \leanref{KeyOrder.cell_lt_iff_rank_finset}{LEANEXT/keys/KeyOrder.lean:71}]\label{oc:lem:rank}
Let $u\in\R$ with $u>0$, let $S\subseteq\R$, and let $\operatorname{rank}:\R\to\N$ be strictly increasing
on $S$. Let $J_1,J_2\in\Z$ and $\psi_1,\psi_2\in S$ with $0\le\psi_1<u$ and $0\le\psi_2<u$. Then
\[J_1u+\psi_1<J_2u+\psi_2\iff(J_1,\operatorname{rank}\psi_1)<_{\mathrm{lex}}(J_2,\operatorname{rank}\psi_2)
\quad\text{in }\Z\times_\ell\N.\]
\lean{cell_lt_iff_rank} is this statement for a set $S$ (a Lean \lean{Set} of reals);
\lean{cell_lt_iff_rank_finset} is the same statement for a finite set $S$ (a Lean \lean{Finset} of reals,
with strict monotonicity on the underlying set and $\psi_1,\psi_2\in S$).
\end{lemma}
\begin{proof}
By Lemma~\ref{oc:lem:cell}(a), the left side is equivalent to ``$J_1<J_2$, or $J_1=J_2$ and
$\psi_1<\psi_2$''. By Definition~\ref{oc:def:keys}(c), the right side is ``$J_1<J_2$, or $J_1=J_2$ and
$\operatorname{rank}\psi_1<\operatorname{rank}\psi_2$''. It remains to show
$\operatorname{rank}\psi_1<\operatorname{rank}\psi_2\iff\psi_1<\psi_2$ for $\psi_1,\psi_2\in S$
(\lean{StrictMonoOn.lt_iff_lt}). ``If'' is strict monotonicity. ``Only if'': if $\psi_1\not<\psi_2$, then
either $\psi_1=\psi_2$, and then $\operatorname{rank}\psi_1=\operatorname{rank}\psi_2$, or $\psi_2<\psi_1$,
and then $\operatorname{rank}\psi_2<\operatorname{rank}\psi_1$; in both cases
$\operatorname{rank}\psi_1\not<\operatorname{rank}\psi_2$. The finite-set version is the set version
applied to the set of elements of the finite set.
\end{proof}

\begin{lemma}[reduction modulo $2\pi$; \leanref{KeyOrder.reduce_mod_two_pi}{LEANEXT/keys/KeyOrder.lean:81}]\label{oc:lem:reduce}
Let $N\in\Z$ with $N>0$, $u\in\R$ with $u>0$ and $Nu=2\pi$, $J\in\Z$, and $\psi\in\R$ with $0\le\psi<u$.
Then
\[\operatorname{toIcoMod}_{2\pi}(0,\,Ju+\psi)=(J\bmod N)\,u+\psi .\]
\end{lemma}
\begin{proof}
Let $R=J\bmod N$ and $Q=\lfloor J/N\rfloor$, so that $0\le R\le N-1$ and $J=R+NQ$
(Definition~\ref{oc:def:keys}(b)). By Lemma~\ref{oc:lem:toico} (with $p=2\pi$, $a=0$) it suffices to show
(i) $0\le Ru+\psi$, (ii) $Ru+\psi<2\pi$, and (iii) $Ju+\psi=(Ru+\psi)+Q\cdot2\pi$.
(i) holds as $R\ge0$, $u>0$ and $\psi\ge0$. (ii): $2\pi-(Ru+\psi)=Nu-Ru-\psi=(N-1-R)u+(u-\psi)>0$, since
$N-1-R\ge0$ and $\psi<u$. (iii): $Ju+\psi=(R+NQ)u+\psi=Ru+\psi+Q(Nu)=Ru+\psi+Q\cdot2\pi$.
\end{proof}

\begin{lemma}[cells of antipodes; \leanref{KeyOrder.antipode_odd}{LEANEXT/keys/KeyOrder.lean:100}, \leanref{KeyOrder.antipode_even}{LEANEXT/keys/KeyOrder.lean:115}]\label{oc:lem:antipode-cell}
Adding $\pi$ to a direction and reducing modulo $2\pi$ changes its cell index as follows and keeps its
phase $\psi$:
\begin{enumerate}[label=(\alph*)]
\item (\lean{antipode_odd}, units $\pi/m$) Let $m\in\N$ with $0<m$, $J\in\Z$ and $\psi\in\R$ with
  $0\le\psi<\pi/m$. Then
  \[\operatorname{toIcoMod}_{2\pi}\bigl(0,\;J\cdot\tfrac{\pi}{m}+\psi+\pi\bigr)
    =\bigl((J+m)\bmod 2m\bigr)\cdot\tfrac{\pi}{m}+\psi .\]
\item (\lean{antipode_even}, units $2\pi/m$) Let $m\in\N$ with $0<m$ and $m$ even, $J\in\Z$ and
  $\psi\in\R$ with $0\le\psi<2\pi/m$. Then
  \[\operatorname{toIcoMod}_{2\pi}\bigl(0,\;J\cdot\tfrac{2\pi}{m}+\psi+\pi\bigr)
    =\bigl((J+m/2)\bmod m\bigr)\cdot\tfrac{2\pi}{m}+\psi,\]
  where $m/2$ is the natural-number quotient, cast to $\Z$.
\end{enumerate}
In (a) the integers $J+m$ and $2m$ are formed in $\Z$; oddness of $m$ is not assumed (the name refers to
the encoder's use of units $\pi/m$ for odd $m$).
\end{lemma}
\begin{proof}
(a) As $m\ne0$, $m\cdot\frac{\pi}{m}=\pi$, so $J\frac{\pi}{m}+\psi+\pi=(J+m)\frac{\pi}{m}+\psi$. Apply
Lemma~\ref{oc:lem:reduce} with $N=2m>0$, $u=\pi/m>0$ (indeed $Nu=2m\cdot\frac{\pi}{m}=2\pi$), the integer
$J+m$, and $\psi$, which satisfies $0\le\psi<u$.
(b) Write $m=k+k$ with $k\in\N$; then $k\ge1$ and $m/2=k$. Since $k\cdot\frac{2\pi}{2k}=\pi$, we get
$J\frac{2\pi}{m}+\psi+\pi=(J+k)\frac{2\pi}{m}+\psi$. Apply Lemma~\ref{oc:lem:reduce} with $N=m>0$,
$u=2\pi/m>0$ (so $Nu=2\pi$), the integer $J+k=J+m/2$, and $\psi$, which satisfies $0\le\psi<u$.
\end{proof}

\begin{theorem}[sorting reduced directions by integer keys; \leanref{KeyOrder.reduced_lt_iff_key}{LEANEXT/keys/KeyOrder.lean:139}]\label{oc:thm:key}
Let $N\in\Z$ with $N>0$, $u\in\R$ with $u>0$ and $Nu=2\pi$, $S\subseteq\R$, and $\operatorname{rank}:\R\to\N$
strictly increasing on $S$. Let $J_1,J_2\in\Z$ and $\psi_1,\psi_2\in S$ with $0\le\psi_1<u$ and
$0\le\psi_2<u$. Then
\[\operatorname{toIcoMod}_{2\pi}(0,J_1u+\psi_1)<\operatorname{toIcoMod}_{2\pi}(0,J_2u+\psi_2)
\iff(J_1\bmod N,\operatorname{rank}\psi_1)<_{\mathrm{lex}}(J_2\bmod N,\operatorname{rank}\psi_2).\]
\end{theorem}
\begin{proof}
By Lemma~\ref{oc:lem:reduce}, applied twice,
$\operatorname{toIcoMod}_{2\pi}(0,J_1u+\psi_1)=(J_1\bmod N)u+\psi_1$ and
$\operatorname{toIcoMod}_{2\pi}(0,J_2u+\psi_2)=(J_2\bmod N)u+\psi_2$. Lemma~\ref{oc:lem:rank}, applied to
the integers $J_1\bmod N$ and $J_2\bmod N$ and to $\psi_1,\psi_2$ (with the same $u$, $S$ and rank),
gives the claim.
\end{proof}

\begin{remark}[half-open intervals]\label{oc:rem:halfopen}
(1) The strict bound $\psi<u$ cannot be weakened to $\psi\le u$: for $J_1=0$, $\psi_1=u$, $J_2=1$,
$\psi_2=0$ we get $J_1u+\psi_1=u=J_2u+\psi_2$, so the left side of Lemma~\ref{oc:lem:cell}(a) is false
while its right side ($0<1$) is true, and Lemma~\ref{oc:lem:cell}(b) fails as well.
(2) The reduction lands in the half-open interval $[0,2\pi)$: for instance $J=N$ and $\psi=0$ give
$\operatorname{toIcoMod}_{2\pi}(0,2\pi)=0=(N\bmod N)u+0$, in accordance with Lemma~\ref{oc:lem:reduce}.
\end{remark}

\begin{remark}[what remains outside Lean]\label{oc:rem:keys-outside}
The statements of this subsection are about real numbers already written as $Ju+\psi$ with $0\le\psi<u$,
and about a rank map assumed strictly increasing on the phases. Three steps are not part of any Lean
statement here: that the encoder's directions have this form (encoder~B computes the cell index as
$J=2k+h_i$, reduced modulo $2m$, for odd $m$, and as $J=k$ for even $m$;
\dirn{C6UPPER/encb/gen_b2.py}, lines 106--109); that its ranks (which are local ranks) are strictly
increasing on the phases; and that sorting the keys gives the circular order of the actual points. The
first two are listed as hypotheses in \dirn{LEANEXT/keys/REPORT.md} (lines 22--25; see also
\dirn{ENCBLEAN/DESIGN.md:626}), the third as outside Lean in \dirn{LEANEXT/README.md} (line~25). (Lines
16--17 of \dirn{LEANEXT/keys/REPORT.md} call \lean{reduced_lt_iff_key} the step ``sorting keys gives the
circular order'' in \dirn{C6UPPER/encb/gen_b2.py}; Theorem~\ref{oc:thm:key} states that step only for real numbers
already in the form above, with a rank as above.) For even $m\ge2$, under the hypotheses of
Theorem~\ref{oc:thm:layers-phase} on the point set (it does not contain $c$, $\rho$ maps it into itself,
it is in general position, and it has no $K$ points in convex position for some $K\le2m$), that theorem
gives, for hull vertices of distinct layers given in polar form with positive radii, directions that are
incongruent modulo $2\pi/m$; once these directions are written as $Ju+\psi$ with $u=2\pi/m$ and $0\le\psi<u$, it follows
that their phases $\psi$ differ. This distinctness of the phases of distinct layers is what a strictly
increasing rank needs.
\end{remark}

\subsection{Coverage}\label{oc:sub:coverage}

\emph{Covered.} From \dirn{LEANEXT/orbit/OrbitConsistency.lean}: \lean{orient_rotAbout_inv}
(Lemma~\ref{oc:lem:rot-inv}); \lean{PreservesOrient} (Definition~\ref{oc:def:preserves}); \lean{RotComp}
(Definition~\ref{oc:def:rotcomp}); \lean{preservesOrient_id}, \lean{PreservesOrient.comp},
\lean{PreservesOrient.iterate}, \lean{preservesOrient_rotAbout}, \lean{RotComp.preservesOrient},
\lean{rotComp_iterate} (Lemma~\ref{oc:lem:rotcomp-pres}); \lean{orient3}
(Definition~\ref{oc:def:orient3}); \lean{orient3_swap} (Lemma~\ref{oc:lem:swap}); \lean{orient3_perm}
(Lemma~\ref{oc:lem:perm}); \lean{orient3_rot_perm} (Lemma~\ref{oc:lem:rot-perm});
\lean{sign_eq_of_two_reps}, \lean{orient3_of_common_rep} (Lemma~\ref{oc:lem:two-reps}); \lean{litVal}
(Definition~\ref{oc:def:litval}); \lean{literal_sound} (Theorem~\ref{oc:thm:literal});
\lean{encoder_literal_sound} (Theorem~\ref{oc:thm:encoder}); \lean{antipode_mem_of_even}
(Lemma~\ref{oc:lem:antipode}); \lean{same_ray_false} (Lemma~\ref{oc:lem:same-ray});
\lean{rotAbout_polar}, \lean{polar_add_int} (Lemma~\ref{oc:lem:polar-rot}); \lean{eq_iterate_of_phase_eq}
(Lemma~\ref{oc:lem:phase-nat}); \lean{eq_iterate_of_phase_eq_int} (Lemma~\ref{oc:lem:phase-int});
\lean{phase_ne_of_not_mem_orbit} (Theorem~\ref{oc:thm:phase-orbit}). Of the copied block (lines 39--112):
\lean{Invariant}, \lean{rotAbout}, \lean{polar} (Definition~\ref{oc:def:basic}), \lean{rotAbout_iterate}
(Lemma~\ref{oc:lem:iterate}) and \lean{iterate_mem} (Lemma~\ref{oc:lem:iter-mem}).
From \dirn{LEANEXT/merged/Merged.lean}, section~13: the seven restatements \lean{antipode_mem_of_even},
\lean{same_ray_false}, \lean{rotAbout_polar}, \lean{eq_iterate_of_phase_eq}, \lean{polar_add_int},
\lean{eq_iterate_of_phase_eq_int}, \lean{phase_ne_of_not_mem_orbit} in the namespace
\lean{ErdosSzekeres.LayerLemma} (same labels as their originals); \lean{not_mem_hullVerts_peel_of_lt},
\lean{not_mem_hullVerts_peel_of_ne} (Lemma~\ref{oc:lem:disjoint}); \lean{iterate_mem_hullVerts_peel}
(Lemma~\ref{oc:lem:closed}); \lean{not_same_ray_of_distinct_layers} (Theorem~\ref{oc:thm:layers-ray});
\lean{phase_ne_of_distinct_layers} (Theorem~\ref{oc:thm:layers-phase}).
From \dirn{LEANEXT/keys/KeyOrder.lean}: \lean{cell_lt_iff}, \lean{cell_eq_iff} (Lemma~\ref{oc:lem:cell});
\lean{cell_lt_iff_rank}, \lean{cell_lt_iff_rank_finset} (Lemma~\ref{oc:lem:rank});
\lean{reduce_mod_two_pi} (Lemma~\ref{oc:lem:reduce}); \lean{antipode_odd}, \lean{antipode_even}
(Lemma~\ref{oc:lem:antipode-cell}); \lean{reduced_lt_iff_key} (Theorem~\ref{oc:thm:key}).

\emph{Not covered, deliberately.} (i) The complex-coordinate helpers of the copied block of
\dirn{OrbitConsistency.lean} (lines 49--90 and 97--103: \lean{toC}, \lean{ofC}, \lean{toC_re}, \lean{toC_im},
\lean{ofC_zero}, \lean{ofC_one}, \lean{toC_ofC}, \lean{ofC_toC}, \lean{toC_injective}, \lean{crot},
\lean{toC_crot}, \lean{crot_one}, \lean{crot_comp}, \lean{crot_iterate}, \lean{rotAbout_eq_crot}): copies
of declarations among lines 231--308 of \dirn{LAYERLEAN/LayerLemma.lean}, which belong to the
layer-lemma unit; our proofs do not need them (the written proof of Lemma~\ref{oc:lem:iterate} uses the
addition formulas). (ii) Lines 1--1147 of \dirn{Merged.lean}: a copy of \dirn{LAYERLEAN/LayerLemma.lean},
treated in the layer-lemma and layer-geometry sections. (iii) Lines 151--158 of \dirn{KeyOrder.lean}:
\texttt{\#print axioms} commands, not results. (iv) The verbatim copies listed under ``Sources'' (section~13
of \dirn{Merged.lean} in \dirn{LEANEXT/linked/Linked.lean}, \dirn{LEANEXT/capstone/Capstone.lean},
\dirn{LEANEXT/capstone/Capstone2.lean} and \dirn{ENCBLEAN/geo/ES8Geo/Capstone2.lean}; the byte-identical
\dirn{ENCBLEAN/geo/ES8Geo/OrbitConsistency.lean} and \dirn{ENCBLEAN/geo/ES8Geo/KeyOrder.lean}), which
are the same declarations. (v) The Python scripts and logs of \dirn{LEANEXT/orbit} (evidence, cited in
Remark~\ref{oc:rem:parity}) and the mutant logs (audit only).
\endgroup

\section{Phase patterns and angular rules}\label{pc:sec}
\begingroup\sloppy

Fix a centre $c\in\R^2$ and three directions $\alpha<\beta<\delta<\alpha+2\pi$, read counterclockwise around
$c$. Together with the three antipodal directions $\alpha+\pi$, $\beta+\pi$, $\delta+\pi$ they give six points on
the circle of directions, distinct unless one of the gaps $\beta-\alpha$, $\delta-\beta$, $\alpha+2\pi-\delta$
equals $\pi$, in which case a direction and an antipode coincide (Remark~\ref{pc:rem:reading}). This section gives
written proofs of the Lean-checked results that classify the cyclic pattern of these points, expressed by the arc
predicates of Definition~\ref{pc:def:arc} (which are defined in every case), and connect it with the sign of
$\orient$ for three points in polar form about $c$.
\begin{itemize}
\item The pattern alternates (every open arc between consecutive directions contains an antipode; pattern
  $PAPAPA$) if and only if all three gaps are smaller than $\pi$ (Theorem~\ref{pc:thm:alternating}).
\item A run $PPP$ with middle point $\beta$ (no antipode in the arcs $(\alpha,\beta)$ and $(\beta,\delta)$) is
  equivalent to $\delta-\alpha<\pi$, provided $\alpha+2\pi-\delta\neq\pi$ (Theorem~\ref{pc:thm:runB}).
\item If no gap equals $\pi$, the pattern alternates or has a run, and alternation excludes every run
  (Theorem~\ref{pc:thm:trichotomy}).
\item Rule~A: an alternating triple at positive radii turns counterclockwise (Theorem~\ref{pc:thm:rule-alt}).
  Rule~B: a triple with a run turns counterclockwise if the gap not adjacent to the middle point of the run
  differs from $\pi$, and the other two radii are positive and at most the radius of the middle point
  (Theorem~\ref{pc:thm:rule-run}). The two rules combined (Theorem~\ref{pc:thm:phase}), and a version for
  arbitrary real directions (Theorem~\ref{pc:thm:sign}).
\item The dichotomy for cyclic Boolean words of length $6$ (Theorem~\ref{pc:thm:word}; computer-assisted in
  Lean, with a written proof by hand here).
\item Rule~B on the convex layers of a rotationally symmetric set, where the radius condition follows from the
  layer indices (Theorem~\ref{pc:thm:layers}).
\end{itemize}

\paragraph{Sources.} File paths are relative to the results directory \dirn{cloud/es7-2026-10-19/results} of
the repository. The Lean sources are:
\begin{itemize}
\item \dirn{LEANEXT/phase/PhaseClassification.lean} (545 lines, SHA-256 \texttt{490efcd0a7e4}\dots), namespace
  \lean{PhaseClassification} (\dirn{LEANEXT/phase/REPORT.md} gives 540 lines for this file, with the same hash
  prefix);
\item \dirn{LEANEXT/linked/Linked.lean} (1416 lines, SHA-256 \texttt{f1e18c119276}\dots); its new lines are
  1385--1416 (\lean{dist2_polar} and \lean{orient_pos_of_layers}, namespace \lean{ErdosSzekeres.LayerLemma});
\item \dirn{LEANEXT/capstone/Capstone.lean} (2004 lines, SHA-256 \texttt{dc382395ca55}\dots); its new lines are
  1966--2004 (namespace \lean{ErdosSzekeres.Capstone});
\item \dirn{LEANEXT/capstone/Capstone2.lean} (2132 lines, SHA-256 \texttt{f8c3930c2427}\dots); its new lines are
  2004--2132.
\end{itemize}
The following relations were checked with \texttt{diff}. Lines 40--102 of \dirn{PhaseClassification.lean}
repeat lines 931--996 of \dirn{LAYERLEAN/LayerLemma.lean} (\lean{polar}, \lean{orient_polar},
\lean{orient_polar_centre_inside}, \lean{sin_add_lt}, \lean{orient_polar_halfplane}) as new constants of the
namespace \lean{PhaseClassification}; only the docstrings differ. Lines 1--1147 of \dirn{Linked.lean} are
lines 1--1147 of \dirn{LAYERLEAN/LayerLemma.lean}, and lines 1--1384 of \dirn{Linked.lean} are those of
\dirn{LEANEXT/merged/Merged.lean}. \dirn{Capstone.lean} consists of an 18-line header, lines 11--1416 of
\dirn{Linked.lean} (as its lines 19--1424), a blank line, lines 6--545 of \dirn{PhaseClassification.lean} (as
its lines 1426--1965) and the new lines 1966--2004. Lines 1--2003 of \dirn{Capstone2.lean} are those of
\dirn{Capstone.lean}. Hence a declaration at line $L$ of \dirn{PhaseClassification.lean} reappears at line
$L+1420$ of \dirn{Capstone.lean} and \dirn{Capstone2.lean}, and one at line $L\ge 11$ of \dirn{Linked.lean}
reappears at line $L+8$. We cite every declaration at its first occurrence. \dirn{Capstone2.lean} is
byte-identical to the copy \dirn{ENCBLEAN/geo/ES8Geo/Capstone2.lean} that the ENCBLEAN geometric package
imports (hash list \dirn{ENCBLEAN/geo/SOURCES.sha256}). Further Lean files cited below, with the first 12 hex
digits of their SHA-256: \dirn{LAYERLEAN/LayerLemma.lean} (\texttt{afb8542dceb0}), \dirn{LEANEXT/merged/Merged.lean}
(\texttt{245b7d65407a}), \dirn{ENCBLEAN/geo/ES8Geo/Audit.lean} (\texttt{39f214bbf30b}),
\dirn{ENCBLEAN/geo/ES8Geo/Labelling.lean} (\texttt{489ba89f8d6f}), \dirn{ENCBLEAN/geo/ES8Geo/ArcBridge.lean}
(\texttt{46966b48a101}), \dirn{ENCBLEAN/geo/ES8Geo/ModelOfLabelling.lean} (\texttt{5b532f49cc11}) and
\dirn{ENCBLEAN/EncB/SoundWord.lean} (\texttt{c1bdb4175ee5}). The files were built with Lean 4.24.0 and, except
\dirn{ENCBLEAN/EncB/SoundWord.lean} (part of a package without Mathlib), with Mathlib \texttt{f897ebcf}; these
are the repository's pins.

\paragraph{Rigor and axioms.} Every definition, lemma, proposition, theorem and corollary below whose head or one
of whose parts carries a Lean reference is \emph{Lean-checked}; the proofs given here are written proofs (not
refereed). Remarks are commentary and, even when they name Lean declarations, are not Lean-checked as a whole; the
mathematical claims they add to the Lean statements are written steps and are not Lean-checked. The log
\dirn{LEANEXT/phase/logs/axioms.log} records the output of \texttt{\#print axioms} for 16 theorems of
\dirn{PhaseClassification.lean} (\lean{alternating_iff_gaps_lt_pi}, \lean{alternating_exactly_one},
\lean{runMidB_iff}, \lean{gap_classification}, \lean{phase_trichotomy}, \lean{orient_pos_of_alternating},
\lean{orient_pos_of_runMidA}, \lean{orient_pos_of_runMidB}, \lean{orient_pos_of_runMidD},
\lean{orient_pos_of_phase}, \lean{orient_sign_of_phase}, \lean{centre_inside_of_alternating},
\lean{halfplane_of_runMidB}, \lean{exists_polar}, \lean{polar_toIcoMod}, \lean{word_dichotomy}); each depends
only on \lean{propext}, \lean{Classical.choice} and \lean{Quot.sound}. Every other declaration of that file
covered here is used in the statement or the proof of one of these 16 theorems, so its axioms are among these
three. For the new
declarations of \dirn{Linked.lean}, \dirn{Capstone.lean} and \dirn{Capstone2.lean} no axiom output is saved in
\dirn{LEANEXT}; \dirn{LEANEXT/linked/REPORT.md} and \dirn{LEANEXT/capstone/REPORT.md} state the same three
axioms. Of these declarations, \lean{dist2_polar}, \lean{Capstone.polar_eq}, \lean{Capstone.exists_angle},
\lean{Capstone.sqrt_dist2_pos} and \lean{Capstone.ne_of_mem} (with \lean{Capstone.polar_zero}, used in the
proof of \lean{sqrt_dist2_pos}) are used, together with \lean{PhaseClassification.orient_pos_of_alternating},
\lean{PhaseClassification.orient_pos_of_runMidB} and \lean{PhaseClassification.polar_add_int_two_pi}, through the
byte-identical copy of \dirn{Capstone2.lean}, in proofs in \dirn{ENCBLEAN/geo/ES8Geo/Labelling.lean} and
\dirn{ENCBLEAN/geo/ES8Geo/ArcBridge.lean} on which the declarations \lean{EncB.Geo.labelling_of_counterexample}
(in \dirn{Labelling.lean}) and \lean{EncB.Geo.modelOfLabelling} (in
\dirn{ENCBLEAN/geo/ES8Geo/ModelOfLabelling.lean}) depend. The axioms of these two declarations are pinned to the
same three by \texttt{\#guard\_msgs} in \dirn{ENCBLEAN/geo/ES8Geo/Audit.lean}, lines 45--49. The package compiles
the copy from the same sources, against the same imports, as the original file: besides the byte-identical copy
and the common Lean and Mathlib pins (see Sources), the ten repository modules that \dirn{Capstone2.lean} needs
(its imports \lean{ErdosSzekeres.Definitions}, \lean{ErdosSzekeres.Orientation},
\lean{ErdosSzekeres.OrientationGeometry} and \lean{ErdosSzekeres.Suk.OrientationClique}, with the repository
modules they import) are compiled there from verbatim copies, whose SHA-256 hashes are listed in
\dirn{ENCBLEAN/geo/VENDOR.sha256} and agree with those of the repository files. So the axioms of the declarations
used are among these three as well. The remaining new declarations (\lean{orient_pos_of_layers},
the six layer rules of \lean{ErdosSzekeres.Capstone} and \lean{Capstone.orient_pos_of_alternating_polar}) are not
used in the ENCBLEAN geometric package outside the copy itself, and no \texttt{\#print axioms} output is saved
for them; besides the axiom statements of the two
reports, the evidence for them is the clean builds stated in these reports and the ENCBLEAN build log
\dirn{ENCBLEAN/logs/m4_build.log}, a successful build of the package that imports the byte-identical copy of
\dirn{Capstone2.lean} as a module, in which only the \lean{KeyOrder} module prints messages (its axiom lines). A
text search of the four files finds no \lean{sorry}, \lean{admit}, \lean{axiom}
declaration or \lean{native_decide}; among the declarations of this section, \lean{decide} occurs only in the
proof of \lean{word_dichotomy}.

\paragraph{Dependencies.} We use Definitions~\ref{fd:def:point}, \ref{fd:def:orient}, \ref{fd:def:rot},
\ref{fd:def:inv} and \ref{fd:def:ngon}, and Lemma~\ref{fd:lem:orient}. From the layer sections we use the hull
vertices $\hv$ and the peeling $\peel$ (Definitions~\ref{ll:def:hv} and~\ref{ll:def:peel}), the constant
\lean{LayerLemma.polar} (Definition~\ref{lg:def:polar}), its polar formula and angular rules
(Lemmas~\ref{lg:lem:polar}, \ref{lg:lem:ruleA}, \ref{lg:lem:sinadd} and~\ref{lg:lem:ruleB}), and the radius
order of the layers (Theorem~\ref{lg:thm:radius}).

\subsection{Conventions and polar points}

Points are elements of $\R^2$, written $p=(p_0,p_1)$ (Lean \lean{Point}, Definition~\ref{fd:def:point}); two
points are equal when both coordinates agree. We use
$\orient(a,b,d)=(b_0-a_0)(d_1-a_1)-(b_1-a_1)(d_0-a_0)$ (Definition~\ref{fd:def:orient}), which is positive when
$a,b,d$ turn counterclockwise, and from Lemma~\ref{fd:lem:orient}: $\orient(a,d,b)=-\orient(a,b,d)$
(\lean{orient_swap}), $\orient(b,d,a)=\orient(a,b,d)$ (\lean{orient_cycle}) and
$\orient(a+v,b+v,d+v)=\orient(a,b,d)$ (\lean{orient_translate}). Angles are arbitrary real numbers. The number
$\pi>0$ is Lean's \lean{Real.pi}, and an integer $k$ is identified with its image in $\R$ (Lean writes this cast
explicitly).

\emph{Reduction modulo $2\pi$.} For $\alpha,\theta\in\R$ we write $\operatorname{toIcoMod}(\alpha,\theta)$ for
Mathlib's \lean{toIcoMod} with period $2\pi$ (its positivity argument is \lean{Real.two_pi_pos}), base point
$\alpha$ and argument $\theta$. By definition $\operatorname{toIcoMod}(\alpha,\theta)=\theta-2\pi n$, where
$n=\operatorname{toIcoDiv}(\alpha,\theta)$ is the unique integer with $\theta-2\pi n\in[\alpha,\alpha+2\pi)$,
namely $n=\lfloor(\theta-\alpha)/(2\pi)\rfloor$. Only two properties are used (Mathlib \lean{toIcoMod_mem_Ico}
and \lean{toIcoMod_add_toIcoDiv_zsmul}):
\begin{enumerate}[label=(T\arabic*)]
\item $\alpha\le\operatorname{toIcoMod}(\alpha,\theta)<\alpha+2\pi$;
\item $\theta=\operatorname{toIcoMod}(\alpha,\theta)+n\cdot 2\pi$ with the integer
  $n=\operatorname{toIcoDiv}(\alpha,\theta)$.
\end{enumerate}

\begin{definition}[polar points and squared distance; \leanref{PhaseClassification.polar}{LEANEXT/phase/PhaseClassification.lean:41}, \leanref{ErdosSzekeres.LayerLemma.dist2}{LAYERLEAN/LayerLemma.lean:999}]\label{pc:def:polar}
For $c\in\R^2$ and $r,\theta\in\R$,
\[
  \operatorname{polar}(c,r,\theta)=(c_0+r\cos\theta,\ c_1+r\sin\theta).
\]
No sign condition is imposed on $r$. For $c,y\in\R^2$, $\operatorname{dist2}(c,y)=(y_0-c_0)^2+(y_1-c_1)^2$.
\end{definition}

The constant \lean{ErdosSzekeres.LayerLemma.polar} (Definition~\ref{lg:def:polar};
\dirn{LAYERLEAN/LayerLemma.lean:932}, repeated at \dirn{LEANEXT/linked/Linked.lean:932} and
\dirn{LEANEXT/capstone/Capstone2.lean:940}) is a different Lean constant given by the same formula;
Lemma~\ref{pc:lem:helpers}(a) records that the two agree. We write $\operatorname{polar}$ for both; every
statement below that involves $\operatorname{polar}$ says, directly or through the convention of its subsection,
which of the two constants its Lean version uses. The definition of $\operatorname{dist2}$ is repeated at
\dirn{LEANEXT/linked/Linked.lean:999} and \dirn{LEANEXT/capstone/Capstone2.lean:1007}.

\subsection{The polar orientation formula and the two rules}

\begin{lemma}[polar formula, rule A, rule B; \leanref{PhaseClassification.orient_polar}{LEANEXT/phase/PhaseClassification.lean:44} for (a), and the declarations cited in (b)--(e)]\label{pc:lem:polar-facts}
Let $c\in\R^2$; all other letters denote real numbers, and $\operatorname{polar}$ is
\lean{PhaseClassification.polar}.
\begin{enumerate}[label=(\alph*)]
\item $\orient(\operatorname{polar}(c,r_a,\alpha),\operatorname{polar}(c,r_b,\beta),\operatorname{polar}(c,r_d,\delta))
  =r_ar_b\sin(\beta-\alpha)+r_br_d\sin(\delta-\beta)+r_dr_a\sin(\alpha-\delta)$.
\item \textup{(\leanref{PhaseClassification.orient_polar_centre}{LEANEXT/phase/PhaseClassification.lean:407}.)}
  $\orient(\operatorname{polar}(c,r_a,\alpha),\operatorname{polar}(c,r_b,\beta),c)=r_ar_b\sin(\beta-\alpha)$.
\item \textup{(Rule A; \leanref{PhaseClassification.orient_polar_centre_inside}{LEANEXT/phase/PhaseClassification.lean:51}.)}
  If $0<r_a$, $0<r_b$, $0<r_d$, $0<\beta-\alpha$, $\beta-\alpha<\pi$, $0<\delta-\beta$,
  $\delta-\beta<\pi$, $0<\alpha-\delta+2\pi$ and $\alpha-\delta+2\pi<\pi$, then
  $0<\orient(\operatorname{polar}(c,r_a,\alpha),\operatorname{polar}(c,r_b,\beta),\operatorname{polar}(c,r_d,\delta))$.
\item \textup{(\leanref{PhaseClassification.sin_add_lt}{LEANEXT/phase/PhaseClassification.lean:64}.)}
  If $0<x$, $x<\pi$, $0<y$ and $y<\pi$, then $\sin(x+y)<\sin x+\sin y$.
\item \textup{(Rule B; \leanref{PhaseClassification.orient_polar_halfplane}{LEANEXT/phase/PhaseClassification.lean:80}.)}
  If $0<r_a$, $0<r_d$, $r_a\le r_b$, $r_d\le r_b$, $0<\beta-\alpha$, $0<\delta-\beta$ and
  $\delta-\alpha<\pi$, then
  $0<\orient(\operatorname{polar}(c,r_a,\alpha),\operatorname{polar}(c,r_b,\beta),\operatorname{polar}(c,r_d,\delta))$.
\end{enumerate}
\end{lemma}

Parts (a), (c) and (e) are the statements of Lemmas~\ref{lg:lem:polar}, \ref{lg:lem:ruleA}
and~\ref{lg:lem:ruleB} with the constant \lean{PhaseClassification.polar} in place of \lean{LayerLemma.polar}, and
part (d) is the statement of Lemma~\ref{lg:lem:sinadd}; the Lean proofs are the same. For completeness we repeat
the short arguments.

\begin{proof}
(a) Put $u=(r_a\cos\alpha,r_a\sin\alpha)$, $v=(r_b\cos\beta,r_b\sin\beta)$, $w=(r_d\cos\delta,r_d\sin\delta)$,
so that the three points are $c+u$, $c+v$, $c+w$, and write $[x,y]=x_0y_1-x_1y_0$ for $x,y\in\R^2$. By
translation invariance, $\orient(c+u,c+v,c+w)=\orient(u,v,w)=[v-u,w-u]$, and by bilinearity and antisymmetry of
$[\cdot,\cdot]$,
\[
  [v-u,w-u]=[v,w]-[v,u]-[u,w]+[u,u]=[u,v]+[v,w]+[w,u].
\]
Finally $[u,v]=r_ar_b(\cos\alpha\sin\beta-\sin\alpha\cos\beta)=r_ar_b\sin(\beta-\alpha)$, and in the same way
$[v,w]=r_br_d\sin(\delta-\beta)$ and $[w,u]=r_dr_a\sin(\alpha-\delta)$.

(b) For every $\delta$, $\operatorname{polar}(c,0,\delta)=c$. Hence (b) is (a) with $r_d=0$. (Lean expands the
definitions directly; the computation is the same.)

(c) By (a) the orientation is $r_ar_b\sin(\beta-\alpha)+r_br_d\sin(\delta-\beta)+r_dr_a\sin(\alpha-\delta)$. The
sine is positive on $(0,\pi)$, so $\sin(\beta-\alpha)>0$, $\sin(\delta-\beta)>0$ and
$\sin(\alpha-\delta)=\sin(\alpha-\delta+2\pi)>0$. All radii are positive, so all three terms are positive.

(d) By the addition formula, $\sin x+\sin y-\sin(x+y)=\sin x\,(1-\cos y)+\sin y\,(1-\cos x)$. Here $\sin x>0$
and $\sin y>0$. Moreover $\cos x<1$ and $\cos y<1$, because $\cos$ is strictly decreasing on $[0,\pi]$ with
$\cos 0=1$. So both products are positive.

(e) Put $x=\beta-\alpha>0$ and $y=\delta-\beta>0$. Then $x+y=\delta-\alpha<\pi$, so $x,y\in(0,\pi)$, and
$\alpha-\delta=-(x+y)$. By (a), the orientation $O$ in (e) satisfies
\begin{align*}
  O&=r_ar_b\sin x+r_br_d\sin y-r_dr_a\sin(x+y)\\
   &=r_a(r_b-r_d)\sin x+r_d(r_b-r_a)\sin y+r_ar_d\bigl(\sin x+\sin y-\sin(x+y)\bigr).
\end{align*}
The first two terms are $\ge0$ (as $r_a,r_d>0$, $r_b-r_d\ge0$, $r_b-r_a\ge0$, $\sin x,\sin y>0$), and the last
is $>0$ by (d) and $r_ar_d>0$. (Positivity of $r_b$ is not assumed; it follows from $r_b\ge r_a>0$.)
\end{proof}

\subsection{Periodicity and polar forms}

\begin{lemma}[periodicity; \leanref{PhaseClassification.polar_add_int_two_pi}{LEANEXT/phase/PhaseClassification.lean:106}, \leanref{PhaseClassification.polar_add_two_pi}{LEANEXT/phase/PhaseClassification.lean:111}, \leanref{PhaseClassification.polar_toIcoMod}{LEANEXT/phase/PhaseClassification.lean:485}]\label{pc:lem:periodic}
Let $c\in\R^2$ and $r,\theta,\alpha\in\R$, and let $\operatorname{polar}$ be \lean{PhaseClassification.polar}.
\begin{enumerate}[label=(\alph*)]
\item For every $k\in\Z$, $\operatorname{polar}(c,r,\theta+k\cdot2\pi)=\operatorname{polar}(c,r,\theta)$.
\item $\operatorname{polar}(c,r,\theta+2\pi)=\operatorname{polar}(c,r,\theta)$.
\item $\operatorname{polar}(c,r,\operatorname{toIcoMod}(\alpha,\theta))=\operatorname{polar}(c,r,\theta)$.
\end{enumerate}
\end{lemma}

\begin{proof}
(a) $\cos(\theta+2\pi k)=\cos\theta$ and $\sin(\theta+2\pi k)=\sin\theta$ for $k\in\Z$. (b) is (a) with $k=1$.
(c) By (T2), $\theta=\operatorname{toIcoMod}(\alpha,\theta)+n\cdot2\pi$ with $n\in\Z$, so by (a)
$\operatorname{polar}(c,r,\theta)=\operatorname{polar}(c,r,\operatorname{toIcoMod}(\alpha,\theta)+n\cdot2\pi)
=\operatorname{polar}(c,r,\operatorname{toIcoMod}(\alpha,\theta))$.
\end{proof}

\begin{lemma}[every point other than the centre has a polar form; \leanref{PhaseClassification.exists_polar}{LEANEXT/phase/PhaseClassification.lean:459}]\label{pc:lem:exists-polar}
For all $c,p\in\R^2$ with $p\neq c$ there are $r,\theta\in\R$ with $0<r$ and $p=\operatorname{polar}(c,r,\theta)$,
where $\operatorname{polar}$ is \lean{PhaseClassification.polar}.
\end{lemma}

\begin{proof}
Let $u=p-c\neq0$ and $r=\sqrt{u_0^2+u_1^2}>0$, and put $t=u_0/r\in[-1,1]$; note $1-t^2=u_1^2/r^2$. Recall that
$\arccos t\in[0,\pi]$, $\cos(\arccos t)=t$ and $\sin(\arccos t)=\sqrt{1-t^2}=\abs{u_1}/r$. If $u_1\ge0$, let
$\theta=\arccos t$; then $\cos\theta=u_0/r$ and $\sin\theta=\abs{u_1}/r=u_1/r$. If $u_1<0$, let
$\theta=-\arccos t$; then $\cos\theta=t=u_0/r$ and $\sin\theta=-\abs{u_1}/r=u_1/r$. Hence
$p=c+u=(c_0+r\cos\theta,c_1+r\sin\theta)=\operatorname{polar}(c,r,\theta)$. (The Lean proof takes
$r=\abs{z}$ and $\theta=\arg z$ for the complex number $z=u_0+iu_1$, using $\abs{z}\cos(\arg z)=\operatorname{Re}z$
and $\abs{z}\sin(\arg z)=\operatorname{Im}z$; this is the same polar representation.)
\end{proof}

\begin{lemma}[helpers for layered points; \leanref{ErdosSzekeres.Capstone.polar_eq}{LEANEXT/capstone/Capstone.lean:2001}, \leanref{ErdosSzekeres.Capstone.polar_zero}{LEANEXT/capstone/Capstone2.lean:2051}, \leanref{ErdosSzekeres.LayerLemma.dist2_polar}{LEANEXT/linked/Linked.lean:1390}, \leanref{ErdosSzekeres.Capstone.exists_angle}{LEANEXT/capstone/Capstone2.lean:2056}, \leanref{ErdosSzekeres.Capstone.sqrt_dist2_pos}{LEANEXT/capstone/Capstone2.lean:2064}, \leanref{ErdosSzekeres.Capstone.ne_of_mem}{LEANEXT/capstone/Capstone2.lean:2073}]\label{pc:lem:helpers}
Let $c,p\in\R^2$ and $r,\theta\in\R$. In \textup{(b)--(e)}, $\operatorname{polar}$ is the constant
\lean{LayerLemma.polar}.
\begin{enumerate}[label=(\alph*)]
\item $\text{\lean{LayerLemma.polar}}(c,r,\theta)=\text{\lean{PhaseClassification.polar}}(c,r,\theta)$.
\item $\operatorname{polar}(c,0,\theta)=c$.
\item $\operatorname{dist2}(c,\operatorname{polar}(c,r,\theta))=r^2$.
\item If $p\neq c$, there is $\theta\in\R$ with $p=\operatorname{polar}\bigl(c,\sqrt{\operatorname{dist2}(c,p)},\theta\bigr)$.
\item If $p\neq c$, then $0<\sqrt{\operatorname{dist2}(c,p)}$.
\item If $S$ is a finite set of points with $c\notin S$ and $p\in S$, then $p\neq c$.
\end{enumerate}
\end{lemma}

\begin{proof}
(a) Both constants are defined by the formula of Definition~\ref{pc:def:polar} (in Lean the proof is
\lean{rfl}). (b) $c_0+0\cdot\cos\theta=c_0$ and $c_1+0\cdot\sin\theta=c_1$.
(c) $(r\cos\theta)^2+(r\sin\theta)^2=r^2(\cos^2\theta+\sin^2\theta)=r^2$.
(d) By Lemma~\ref{pc:lem:exists-polar} and (a) there are $r>0$ and $\theta$ with $p=\operatorname{polar}(c,r,\theta)$.
By (c), $\operatorname{dist2}(c,p)=r^2$, and $\sqrt{r^2}=r$ because $r\ge0$. Hence
$p=\operatorname{polar}(c,\sqrt{\operatorname{dist2}(c,p)},\theta)$.
(e) The square root is $\ge0$. If it were $0$, then by (d) and (b) $p=\operatorname{polar}(c,0,\theta)=c$, a
contradiction.
(f) If $p=c$, then $c=p\in S$, contrary to $c\notin S$.
\end{proof}

\subsection{Arcs and phase patterns}

\begin{definition}[arcs and phase patterns; \leanref{PhaseClassification.inArc}{LEANEXT/phase/PhaseClassification.lean:117}, \leanref{PhaseClassification.arcHasAntipode}{LEANEXT/phase/PhaseClassification.lean:120}, \leanref{PhaseClassification.Alternating}{LEANEXT/phase/PhaseClassification.lean:126}, \leanref{PhaseClassification.RunMidB}{LEANEXT/phase/PhaseClassification.lean:132}, \leanref{PhaseClassification.RunMidD}{LEANEXT/phase/PhaseClassification.lean:136}, \leanref{PhaseClassification.RunMidA}{LEANEXT/phase/PhaseClassification.lean:140}]\label{pc:def:arc}
For $s,g,\theta\in\R$,
\[
  \operatorname{inArc}(s,g,\theta)\ :\Longleftrightarrow\ \exists k\in\Z:\ s<\theta+2\pi k\ \text{ and }\ \theta+2\pi k<s+g .
\]
For $\alpha,\beta,\delta,s,g\in\R$,
\[
  \operatorname{arcHasAntipode}(\alpha,\beta,\delta,s,g)\ :\Longleftrightarrow\
  \operatorname{inArc}(s,g,\alpha+\pi)\vee\operatorname{inArc}(s,g,\beta+\pi)\vee\operatorname{inArc}(s,g,\delta+\pi).
\]
For $\alpha,\beta,\delta\in\R$ abbreviate (this notation is not a Lean name)
\begin{align*}
  H_1&=\operatorname{arcHasAntipode}(\alpha,\beta,\delta,\alpha,\beta-\alpha),\\
  H_2&=\operatorname{arcHasAntipode}(\alpha,\beta,\delta,\beta,\delta-\beta),\\
  H_3&=\operatorname{arcHasAntipode}(\alpha,\beta,\delta,\delta,\alpha+2\pi-\delta),
\end{align*}
and write $H_j(\alpha,\beta,\delta)$ when the triple has to be shown. Then
\begin{align*}
  \operatorname{Alternating}(\alpha,\beta,\delta)&:\Longleftrightarrow H_1\wedge H_2\wedge H_3,&
  \operatorname{RunMidB}(\alpha,\beta,\delta)&:\Longleftrightarrow \neg H_1\wedge\neg H_2,\\
  \operatorname{RunMidD}(\alpha,\beta,\delta)&:\Longleftrightarrow \neg H_2\wedge\neg H_3,&
  \operatorname{RunMidA}(\alpha,\beta,\delta)&:\Longleftrightarrow \neg H_3\wedge\neg H_1 .
\end{align*}
\end{definition}

\begin{remark}[reading the definitions]\label{pc:rem:reading}
For $0<g\le2\pi$, $\operatorname{inArc}(s,g,\theta)$ says that the direction $\theta$, taken modulo $2\pi$, lies
on the open counterclockwise arc from $s$ to $s+g$. When $\alpha<\beta<\delta<\alpha+2\pi$, the arcs of
$H_1,H_2,H_3$ are the open arcs $(\alpha,\beta)$, $(\beta,\delta)$, $(\delta,\alpha+2\pi)$ between
counterclockwise-consecutive directions (for $H_3$, $s+g=\delta+(\alpha+2\pi-\delta)=\alpha+2\pi$), and $H_j$ says
that the $j$-th arc contains the direction of an antipode. We use the gaps $g_1=\beta-\alpha$,
$g_2=\delta-\beta$, $g_3=\alpha+2\pi-\delta$; they satisfy $g_1+g_2+g_3=2\pi$, and $g_1,g_2,g_3>0$ when
$\alpha<\beta<\delta<\alpha+2\pi$. Note $g_1+g_2=\delta-\alpha$, $g_2+g_3=\alpha+2\pi-\beta$ and
$g_3+g_1=\beta+2\pi-\delta$. If $\alpha<\beta<\delta<\alpha+2\pi$ and no gap equals $\pi$, the six directions
$\alpha,\beta,\delta,\alpha+\pi,\beta+\pi,\delta+\pi$ are distinct modulo $2\pi$ (a direction and an antipode
coincide modulo $2\pi$ exactly when $\beta-\alpha$, $\delta-\beta$ or $\delta-\alpha$ equals $\pi$, that is, when
$g_1$, $g_2$ or $g_3$ equals $\pi$). In that case, writing $P$ for a direction and $A$ for an antipode,
$\operatorname{Alternating}$ is the cyclic pattern $PAPAPA$, and $\operatorname{RunMidB}$ (resp.\
$\operatorname{RunMidD}$, $\operatorname{RunMidA}$) is a cyclic run $PPP$ whose middle point is $\beta$ (resp.\
$\delta$, $\alpha$). The definitions themselves assume no order; every result below states the order it uses.
\end{remark}

\begin{lemma}[membership and non-membership; \leanref{PhaseClassification.inArc_of}{LEANEXT/phase/PhaseClassification.lean:143}, \leanref{PhaseClassification.not_inArc}{LEANEXT/phase/PhaseClassification.lean:147}]\label{pc:lem:inarc}
Let $s,g,\theta\in\R$.
\begin{enumerate}[label=(\alph*)]
\item For every $k\in\Z$: if $s<\theta+2\pi k$ and $\theta+2\pi k<s+g$, then $\operatorname{inArc}(s,g,\theta)$.
\item For every $m\in\Z$: if $s+g\le\theta+2\pi m$ and $\theta+2\pi m\le s+2\pi$, then
  $\neg\operatorname{inArc}(s,g,\theta)$.
\end{enumerate}
\end{lemma}

\begin{proof}
(a) is the definition, with witness $k$.
(b) Suppose $k\in\Z$ satisfies $s<\theta+2\pi k<s+g$. If $k\ge m$, then $\theta+2\pi k\ge\theta+2\pi m\ge s+g$
(as $\pi>0$), contradicting $\theta+2\pi k<s+g$. If $k\le m-1$, then
$\theta+2\pi k\le\theta+2\pi m-2\pi\le s$, contradicting $s<\theta+2\pi k$. (The Lean proof makes the same case
split on the sign of $k-m$.)
\end{proof}

\begin{lemma}[shift invariance and relabelling; written step]\label{pc:lem:shift}
Let $s,g,\theta,\alpha,\beta,\delta\in\R$ and $n\in\Z$.
\begin{enumerate}[label=(\alph*)]
\item $\operatorname{inArc}(s,g,\theta+2\pi n)\iff\operatorname{inArc}(s,g,\theta)$.
\item $\operatorname{inArc}(s+2\pi n,g,\theta)\iff\operatorname{inArc}(s,g,\theta)$.
\item $\operatorname{RunMidD}(\alpha,\beta,\delta)\iff\operatorname{RunMidB}(\beta,\delta,\alpha+2\pi)$.
\item $\operatorname{RunMidA}(\alpha,\beta,\delta)\iff\operatorname{RunMidB}(\delta,\alpha+2\pi,\beta+2\pi)$.
\end{enumerate}
\end{lemma}

\begin{proof}
(a) If $s<\theta+2\pi n+2\pi k<s+g$, then $k'=n+k$ witnesses $\operatorname{inArc}(s,g,\theta)$; conversely a
witness $k'$ for $\theta$ gives the witness $k'-n$ for $\theta+2\pi n$.
(b) $s+2\pi n<\theta+2\pi k<s+2\pi n+g$ is equivalent to $s<\theta+2\pi(k-n)<s+g$.

(c) For all $s,g$, $\operatorname{arcHasAntipode}(\beta,\delta,\alpha+2\pi,s,g)$ is the disjunction of
$\operatorname{inArc}(s,g,\beta+\pi)$, $\operatorname{inArc}(s,g,\delta+\pi)$ and
$\operatorname{inArc}(s,g,\alpha+\pi+2\pi)$; by (a) with $n=1$ the last one is equivalent to
$\operatorname{inArc}(s,g,\alpha+\pi)$. Hence
$\operatorname{arcHasAntipode}(\beta,\delta,\alpha+2\pi,s,g)\iff\operatorname{arcHasAntipode}(\alpha,\beta,\delta,s,g)$
for all $s,g$. By definition, $\operatorname{RunMidB}(\beta,\delta,\alpha+2\pi)$ is the conjunction of
$\neg\operatorname{arcHasAntipode}(\beta,\delta,\alpha+2\pi,\beta,\delta-\beta)$ and
$\neg\operatorname{arcHasAntipode}(\beta,\delta,\alpha+2\pi,\delta,(\alpha+2\pi)-\delta)$, which are therefore
equivalent to $\neg H_2$ and $\neg H_3$ (of $(\alpha,\beta,\delta)$). So it is equivalent to
$\operatorname{RunMidD}(\alpha,\beta,\delta)$.

(d) As in (c), $\operatorname{arcHasAntipode}(\delta,\alpha+2\pi,\beta+2\pi,s,g)\iff
\operatorname{arcHasAntipode}(\alpha,\beta,\delta,s,g)$ for all $s,g$ (the antipodes
$\alpha+2\pi+\pi$, $\beta+2\pi+\pi$ are shifted by $2\pi$). $\operatorname{RunMidB}(\delta,\alpha+2\pi,\beta+2\pi)$
is the conjunction of $\neg\operatorname{arcHasAntipode}(\delta,\alpha+2\pi,\beta+2\pi,\delta,(\alpha+2\pi)-\delta)$,
equivalent to $\neg H_3$, and of
$\neg\operatorname{arcHasAntipode}(\delta,\alpha+2\pi,\beta+2\pi,\alpha+2\pi,(\beta+2\pi)-(\alpha+2\pi))$. Since
$(\beta+2\pi)-(\alpha+2\pi)=\beta-\alpha$ and, by (b) with $n=1$,
$\operatorname{inArc}(\alpha+2\pi,\beta-\alpha,\theta)\iff\operatorname{inArc}(\alpha,\beta-\alpha,\theta)$ for
every $\theta$, the second conjunct is equivalent to $\neg H_1$. So the conjunction is equivalent to
$\neg H_3\wedge\neg H_1=\operatorname{RunMidA}(\alpha,\beta,\delta)$.

The Lean proofs of \lean{orient_pos_of_runMidD} and \lean{orient_pos_of_runMidA}
(\dirn{LEANEXT/phase/PhaseClassification.lean}, lines 323--335 and 346--362) prove the implications from left to right
inside the proofs, disjunct by disjunct, keeping the witness $k$ or changing it by $+1$ or $-1$, as (a) and (b)
do above.
\end{proof}

\begin{lemma}[empty arcs; \leanref{PhaseClassification.arc1_empty}{LEANEXT/phase/PhaseClassification.lean:161}, \leanref{PhaseClassification.arc2_empty}{LEANEXT/phase/PhaseClassification.lean:169}, \leanref{PhaseClassification.arc2_empty'}{LEANEXT/phase/PhaseClassification.lean:177}, \leanref{PhaseClassification.arc3_empty}{LEANEXT/phase/PhaseClassification.lean:185}, \leanref{PhaseClassification.arc3_empty'}{LEANEXT/phase/PhaseClassification.lean:193}, \leanref{PhaseClassification.arc1_empty'}{LEANEXT/phase/PhaseClassification.lean:201}]\label{pc:lem:empty}
Let $\alpha,\beta,\delta\in\R$, with $H_1,H_2,H_3$ as in Definition~\ref{pc:def:arc}.
\begin{enumerate}[label=(\alph*)]
\item \textup{(\lean{arc1_empty})} If $\alpha<\beta$, $\beta<\delta$ and $\delta-\alpha\le\pi$, then $\neg H_1$.
\item \textup{(\lean{arc2_empty})} If $\alpha<\beta$, $\beta<\delta$ and $\delta-\alpha\le\pi$, then $\neg H_2$.
\item \textup{(\lean{arc2_empty'})} If $\beta<\delta$, $\delta<\alpha+2\pi$ and $\alpha+2\pi-\beta\le\pi$, then $\neg H_2$.
\item \textup{(\lean{arc3_empty})} If $\beta<\delta$, $\delta<\alpha+2\pi$ and $\alpha+2\pi-\beta\le\pi$, then $\neg H_3$.
\item \textup{(\lean{arc3_empty'})} If $\alpha<\beta$, $\delta<\alpha+2\pi$ and $\beta+2\pi-\delta\le\pi$, then $\neg H_3$.
\item \textup{(\lean{arc1_empty'})} If $\alpha<\beta$, $\delta<\alpha+2\pi$ and $\beta+2\pi-\delta\le\pi$, then $\neg H_1$.
\end{enumerate}
\end{lemma}

\begin{proof}
In each part we apply Lemma~\ref{pc:lem:inarc}(b) to each of the three antipodes $\theta\in\{\alpha+\pi,
\beta+\pi,\delta+\pi\}$ with the integer $m$ indicated, and verify the two inequalities
$s+g\le\theta+2\pi m\le s+2\pi$ for the arc in question; this excludes all three disjuncts of
$\operatorname{arcHasAntipode}$. Throughout we use $\pi>0$. (The Lean proof uses the same integers $m$.)

(a) Arc $H_1$: $s=\alpha$, $s+g=\beta$, $s+2\pi=\alpha+2\pi$. From $\beta<\delta$ and $\delta-\alpha\le\pi$ we get
$\beta<\alpha+\pi$ and $\beta-\alpha<\pi$.
$\theta=\alpha+\pi$, $m=0$: $\beta\le\alpha+\pi\le\alpha+2\pi$.
$\theta=\beta+\pi$, $m=0$: $\beta\le\beta+\pi$, and $\beta+\pi\le\alpha+2\pi$ as $\beta-\alpha<\pi$.
$\theta=\delta+\pi$, $m=0$: $\beta\le\delta+\pi$ as $\beta<\delta$, and $\delta+\pi\le\alpha+2\pi$ as $\delta-\alpha\le\pi$.

(b) Arc $H_2$: $s=\beta$, $s+g=\delta$, $s+2\pi=\beta+2\pi$. From $\alpha<\beta$ and $\delta-\alpha\le\pi$ we get
$\delta-\beta<\pi$.
$\theta=\alpha+\pi$, $m=0$: $\delta\le\alpha+\pi$, and $\alpha+\pi\le\beta+2\pi$ as $\alpha<\beta$.
$\theta=\beta+\pi$, $m=0$: $\delta\le\beta+\pi$ as $\delta-\beta<\pi$, and $\beta+\pi\le\beta+2\pi$.
$\theta=\delta+\pi$, $m=0$: $\delta\le\delta+\pi$, and $\delta+\pi\le\beta+2\pi$ as $\delta-\beta<\pi$.

(c) Arc $H_2$ as in (b). The hypothesis says $\alpha+\pi\le\beta$; with $\delta<\alpha+2\pi$ this gives
$\delta<\beta+\pi$.
$\theta=\alpha+\pi$, $m=1$: $\delta\le\alpha+3\pi$ as $\delta<\alpha+2\pi$, and $\alpha+3\pi\le\beta+2\pi$ as
$\alpha+\pi\le\beta$.
$\theta=\beta+\pi$, $m=0$: $\delta\le\beta+\pi$, and $\beta+\pi\le\beta+2\pi$.
$\theta=\delta+\pi$, $m=0$: $\delta\le\delta+\pi$, and $\delta+\pi\le\beta+2\pi$ as $\delta<\beta+\pi$.

(d) Arc $H_3$: $s=\delta$, $s+g=\alpha+2\pi$, $s+2\pi=\delta+2\pi$. The hypothesis says $\alpha+\pi\le\beta$; with
$\beta<\delta$ this gives $\alpha+\pi<\delta$.
$\theta=\alpha+\pi$, $m=1$: $\alpha+2\pi\le\alpha+3\pi$, and $\alpha+3\pi\le\delta+2\pi$ as $\alpha+\pi<\delta$.
$\theta=\beta+\pi$, $m=0$: $\alpha+2\pi\le\beta+\pi$ as $\alpha+\pi\le\beta$, and $\beta+\pi\le\delta+2\pi$ as
$\beta<\delta$.
$\theta=\delta+\pi$, $m=0$: $\alpha+2\pi\le\delta+\pi$ as $\alpha+\pi<\delta$, and $\delta+\pi\le\delta+2\pi$.

(e) Arc $H_3$ as in (d). The hypothesis says $\beta+\pi\le\delta$; with $\alpha<\beta$ this gives
$\alpha+\pi<\delta$.
$\theta=\alpha+\pi$, $m=1$: $\alpha+2\pi\le\alpha+3\pi$, and $\alpha+3\pi\le\delta+2\pi$ as $\alpha+\pi<\delta$.
$\theta=\beta+\pi$, $m=1$: $\alpha+2\pi\le\beta+3\pi$ as $\alpha<\beta$, and $\beta+3\pi\le\delta+2\pi$ as
$\beta+\pi\le\delta$.
$\theta=\delta+\pi$, $m=0$: $\alpha+2\pi\le\delta+\pi$ as $\alpha+\pi<\delta$, and $\delta+\pi\le\delta+2\pi$.

(f) Arc $H_1$ as in (a). The hypothesis says $\beta+\pi\le\delta$; with $\delta<\alpha+2\pi$ this gives
$\beta<\alpha+\pi$ and $\beta+\pi<\alpha+2\pi$.
$\theta=\alpha+\pi$, $m=0$: $\beta\le\alpha+\pi\le\alpha+2\pi$.
$\theta=\beta+\pi$, $m=0$: $\beta\le\beta+\pi$, and $\beta+\pi\le\alpha+2\pi$.
$\theta=\delta+\pi$, $m=-1$: $\beta\le\delta-\pi$ as $\beta+\pi\le\delta$, and $\delta-\pi\le\alpha+2\pi$ as
$\delta<\alpha+2\pi$.
\end{proof}

\begin{lemma}[a long arc meets an antipode; \leanref{PhaseClassification.arc12_nonempty}{LEANEXT/phase/PhaseClassification.lean:229}]\label{pc:lem:arc12}
Let $\alpha,\beta,\delta\in\R$ with $\alpha<\beta$, $\beta<\delta$, $\delta<\alpha+2\pi$ and $\pi<\delta-\alpha$.
Then $H_1\vee H_2$.
\end{lemma}

\begin{proof}
We give a three-case proof; the Lean proof uses a finer case split (on $\delta-\beta$ against $\pi$ first, then
on $\beta-\alpha$), each case again exhibiting one antipode with an explicit witness $k$.

Case $\beta-\alpha<\pi$: the antipode $\alpha+\pi$ lies in the arc of $H_2$ with $k=0$, since
$\beta<\alpha+\pi$ and $\alpha+\pi<\delta$ (the latter is $\pi<\delta-\alpha$). So $H_2$.

Case $\beta-\alpha>\pi$: $\alpha+\pi$ lies in the arc of $H_1$ with $k=0$, since $\alpha<\alpha+\pi<\beta$. So $H_1$.

Case $\beta-\alpha=\pi$: the antipode $\delta+\pi$ lies in the arc of $H_1$ with $k=-1$: indeed
$\alpha<\delta-\pi$ is $\pi<\delta-\alpha$, and $\delta-\pi<\beta=\alpha+\pi$ is $\delta<\alpha+2\pi$. So $H_1$.
\end{proof}

\subsection{Alternation and runs}

\begin{theorem}[alternation iff all gaps are $<\pi$; \leanref{PhaseClassification.alternating_iff_gaps_lt_pi}{LEANEXT/phase/PhaseClassification.lean:210}]\label{pc:thm:alternating}
Let $\alpha,\beta,\delta\in\R$ with $\alpha<\beta$, $\beta<\delta$ and $\delta<\alpha+2\pi$. Then
\[
  \operatorname{Alternating}(\alpha,\beta,\delta)\iff
  \bigl(\beta-\alpha<\pi\ \wedge\ \delta-\beta<\pi\ \wedge\ \alpha+2\pi-\delta<\pi\bigr).
\]
\end{theorem}

\begin{proof}
($\Rightarrow$) Assume $H_1\wedge H_2\wedge H_3$. If $\beta-\alpha\ge\pi$, then $\alpha+2\pi-\beta\le\pi$, and
Lemma~\ref{pc:lem:empty}(c) (with $\beta<\delta<\alpha+2\pi$) gives $\neg H_2$, a contradiction. If
$\delta-\beta\ge\pi$, then $\beta+2\pi-\delta\le\pi$, and Lemma~\ref{pc:lem:empty}(e) (with $\alpha<\beta$,
$\delta<\alpha+2\pi$) gives $\neg H_3$. If $\alpha+2\pi-\delta\ge\pi$, then $\delta-\alpha\le\pi$, and
Lemma~\ref{pc:lem:empty}(a) gives $\neg H_1$. Hence all three gaps are $<\pi$.

($\Leftarrow$) Assume $g_1,g_2,g_3<\pi$. By Lemma~\ref{pc:lem:inarc}(a):
$H_1$ holds through $\delta+\pi$ with $k=-1$, since $\alpha<\delta-\pi$ (that is $g_3<\pi$) and $\delta-\pi<\beta$
(that is $g_2<\pi$);
$H_2$ holds through $\alpha+\pi$ with $k=0$, since $\beta<\alpha+\pi$ ($g_1<\pi$) and $\alpha+\pi<\delta$
($g_3<\pi$);
$H_3$ holds through $\beta+\pi$ with $k=0$, since $\delta<\beta+\pi$ ($g_2<\pi$) and
$\beta+\pi<\delta+(\alpha+2\pi-\delta)=\alpha+2\pi$ ($g_1<\pi$).
\end{proof}

\begin{remark}[a gap equal to $\pi$]\label{pc:rem:alt-boundary}
No genericity is assumed: a gap may equal $\pi$, and then alternation fails. For example, if $\beta-\alpha=\pi$,
the antipode $\alpha+\pi$ is the endpoint $\beta$ of the open arc $(\beta,\delta)$, and
Lemma~\ref{pc:lem:empty}(c) shows that this arc contains no antipode.
\end{remark}

\begin{proposition}[each arc of an alternating pattern holds exactly the opposite antipode; \leanref{PhaseClassification.alternating_exactly_one}{LEANEXT/phase/PhaseClassification.lean:386}]\label{pc:prop:exactly-one}
Let $\alpha,\beta,\delta\in\R$ with $\alpha<\beta$, $\beta<\delta$, $\delta<\alpha+2\pi$ and
$\operatorname{Alternating}(\alpha,\beta,\delta)$. Then
\begin{enumerate}[label=(\roman*)]
\item $\operatorname{inArc}(\alpha,\beta-\alpha,\delta+\pi)$, $\neg\operatorname{inArc}(\alpha,\beta-\alpha,\alpha+\pi)$ and
  $\neg\operatorname{inArc}(\alpha,\beta-\alpha,\beta+\pi)$;
\item $\operatorname{inArc}(\beta,\delta-\beta,\alpha+\pi)$, $\neg\operatorname{inArc}(\beta,\delta-\beta,\beta+\pi)$ and
  $\neg\operatorname{inArc}(\beta,\delta-\beta,\delta+\pi)$;
\item $\operatorname{inArc}(\delta,\alpha+2\pi-\delta,\beta+\pi)$,
  $\neg\operatorname{inArc}(\delta,\alpha+2\pi-\delta,\alpha+\pi)$ and
  $\neg\operatorname{inArc}(\delta,\alpha+2\pi-\delta,\delta+\pi)$.
\end{enumerate}
\end{proposition}

\begin{proof}
By Theorem~\ref{pc:thm:alternating}, $g_1,g_2,g_3<\pi$. The three memberships were shown in the proof of
Theorem~\ref{pc:thm:alternating}, ($\Leftarrow$). For the non-memberships we use Lemma~\ref{pc:lem:inarc}(b):
\begin{itemize}
\item arc $(\alpha,\beta)$, i.e.\ $s+g=\beta$, $s+2\pi=\alpha+2\pi$: for $\alpha+\pi$ with $m=0$,
  $\beta\le\alpha+\pi$ ($g_1<\pi$) and $\alpha+\pi\le\alpha+2\pi$; for $\beta+\pi$ with $m=0$, $\beta\le\beta+\pi$
  and $\beta+\pi\le\alpha+2\pi$ ($g_1<\pi$);
\item arc $(\beta,\delta)$, i.e.\ $s+g=\delta$, $s+2\pi=\beta+2\pi$: for $\beta+\pi$ with $m=0$,
  $\delta\le\beta+\pi$ ($g_2<\pi$) and $\beta+\pi\le\beta+2\pi$; for $\delta+\pi$ with $m=0$, $\delta\le\delta+\pi$
  and $\delta+\pi\le\beta+2\pi$ ($g_2<\pi$);
\item arc $(\delta,\alpha+2\pi)$, i.e.\ $s+g=\alpha+2\pi$, $s+2\pi=\delta+2\pi$: for $\alpha+\pi$ with $m=1$,
  $\alpha+2\pi\le\alpha+3\pi$ and $\alpha+3\pi\le\delta+2\pi$ ($g_3<\pi$); for $\delta+\pi$ with $m=0$,
  $\alpha+2\pi\le\delta+\pi$ ($g_3<\pi$) and $\delta+\pi\le\delta+2\pi$.\qedhere
\end{itemize}
\end{proof}

\begin{theorem}[a run with middle $\beta$ iff an open half-plane; \leanref{PhaseClassification.runMidB_iff}{LEANEXT/phase/PhaseClassification.lean:251}]\label{pc:thm:runB}
Let $\alpha,\beta,\delta\in\R$ with $\alpha<\beta$, $\beta<\delta$, $\delta<\alpha+2\pi$ and
$\alpha+2\pi-\delta\neq\pi$. Then
\[
  \operatorname{RunMidB}(\alpha,\beta,\delta)\iff\delta-\alpha<\pi .
\]
\end{theorem}

\begin{proof}
($\Leftarrow$) If $\delta-\alpha<\pi$, then $\delta-\alpha\le\pi$, and Lemma~\ref{pc:lem:empty}(a),(b) give
$\neg H_1$ and $\neg H_2$.
($\Rightarrow$) Assume $\neg H_1\wedge\neg H_2$ and, for a contradiction, $\delta-\alpha\ge\pi$. Since
$\alpha+2\pi-\delta\neq\pi$ we have $\delta-\alpha\neq\pi$, hence $\pi<\delta-\alpha$, and
Lemma~\ref{pc:lem:arc12} gives $H_1\vee H_2$, a contradiction.
\end{proof}

\begin{remark}[the gap hypothesis is needed]\label{pc:rem:runB-example}
The direction ($\Rightarrow$) fails without $\alpha+2\pi-\delta\neq\pi$. Take any $\alpha<\beta<\delta$ with
$\delta-\alpha=\pi$, for instance $(\alpha,\beta,\delta)=(0,\pi/2,\pi)$. Then $\delta<\alpha+2\pi$, and
$\operatorname{RunMidB}(\alpha,\beta,\delta)$ holds by Lemma~\ref{pc:lem:empty}(a),(b) (their hypothesis is
$\delta-\alpha\le\pi$), but $\delta-\alpha<\pi$ is false. Geometrically, points at positive radii in the
directions $\alpha$ and $\delta=\alpha+\pi$ lie on opposite rays from $c$, so no open half-plane bounded by a line
through $c$ contains both (compare Proposition~\ref{pc:prop:halfplane}). The direction ($\Leftarrow$) does not use
the gap hypothesis.
\end{remark}

\subsection{Classification under genericity}

\begin{proposition}[gap classification; \leanref{PhaseClassification.gap_classification}{LEANEXT/phase/PhaseClassification.lean:268}]\label{pc:prop:gaps}
Let $\alpha,\beta,\delta\in\R$ with $\alpha<\beta$, $\beta<\delta$, $\delta<\alpha+2\pi$, $\beta-\alpha\neq\pi$,
$\delta-\beta\neq\pi$ and $\alpha+2\pi-\delta\neq\pi$. Then
\[
  \bigl(\beta-\alpha<\pi\wedge\delta-\beta<\pi\wedge\alpha+2\pi-\delta<\pi\bigr)\ \vee\ \delta-\alpha<\pi\ \vee\
  \alpha+2\pi-\beta<\pi\ \vee\ \beta+2\pi-\delta<\pi .
\]
\end{proposition}

\begin{proof}
The gaps $g_1,g_2,g_3$ are positive and sum to $2\pi$, and none equals $\pi$, so each is $<\pi$ or $>\pi$. If all
three are $<\pi$, the first disjunct holds. Two gaps cannot both exceed $\pi$: their sum would exceed $2\pi$,
whereas it equals $2\pi$ minus the (positive) third gap. If exactly one gap exceeds $\pi$, the other two sum to
less than $\pi$: if $g_3>\pi$ then $\delta-\alpha=g_1+g_2=2\pi-g_3<\pi$; if $g_1>\pi$ then
$\alpha+2\pi-\beta=g_2+g_3<\pi$; if $g_2>\pi$ then $\beta+2\pi-\delta=g_3+g_1<\pi$. (The Lean proof runs through
the same $2^3$ sign cases, closing the impossible ones by linear arithmetic.)
\end{proof}

\begin{remark}[boundary cases and exclusivity]\label{pc:rem:gaps}
(1) Genericity cannot be dropped here: for $(\alpha,\beta,\delta)=(0,\pi,3\pi/2)$ all hypotheses hold except
$\beta-\alpha\neq\pi$ (here $g_1=\pi$, $g_2=g_3=\pi/2$), and the four disjuncts read $\pi<\pi$ (first conjunct of the first disjunct), $3\pi/2<\pi$, $\pi<\pi$ and $3\pi/2<\pi$, all
false.

(2) The docstring of \lean{gap_classification} speaks of exactly one case; the Lean statement is only the
disjunction above. The disjuncts are in fact mutually exclusive whenever $\alpha<\beta<\delta<\alpha+2\pi$ (a
written step, not part of the Lean statement): if all gaps are $<\pi$, then each sum of two gaps is
$2\pi$ minus the third, hence $>\pi$; and two of the last three disjuncts together give, after adding,
$2\pi+g_j<2\pi$ for some $j$, which contradicts $g_j>0$.
\end{remark}

\begin{theorem}[phase trichotomy; \leanref{PhaseClassification.phase_trichotomy}{LEANEXT/phase/PhaseClassification.lean:285}]\label{pc:thm:trichotomy}
Let $\alpha,\beta,\delta\in\R$ with $\alpha<\beta$, $\beta<\delta$, $\delta<\alpha+2\pi$, $\beta-\alpha\neq\pi$,
$\delta-\beta\neq\pi$ and $\alpha+2\pi-\delta\neq\pi$. Then
\begin{enumerate}[label=(\roman*)]
\item $\operatorname{Alternating}(\alpha,\beta,\delta)\vee\operatorname{RunMidA}(\alpha,\beta,\delta)\vee
  \operatorname{RunMidB}(\alpha,\beta,\delta)\vee\operatorname{RunMidD}(\alpha,\beta,\delta)$, and
\item $\operatorname{Alternating}(\alpha,\beta,\delta)\Rightarrow\bigl(\neg\operatorname{RunMidA}(\alpha,\beta,\delta)
  \wedge\neg\operatorname{RunMidB}(\alpha,\beta,\delta)\wedge\neg\operatorname{RunMidD}(\alpha,\beta,\delta)\bigr)$.
\end{enumerate}
\end{theorem}

\begin{proof}
(i) Apply Proposition~\ref{pc:prop:gaps}. If all gaps are $<\pi$, Theorem~\ref{pc:thm:alternating} gives
$\operatorname{Alternating}$. If $\delta-\alpha<\pi$, Lemma~\ref{pc:lem:empty}(a),(b) give $\neg H_1\wedge\neg H_2$,
i.e.\ $\operatorname{RunMidB}$. If $\alpha+2\pi-\beta<\pi$, Lemma~\ref{pc:lem:empty}(c),(d) give
$\neg H_2\wedge\neg H_3$, i.e.\ $\operatorname{RunMidD}$. If $\beta+2\pi-\delta<\pi$, Lemma~\ref{pc:lem:empty}(e),(f)
give $\neg H_3\wedge\neg H_1$, i.e.\ $\operatorname{RunMidA}$. (Each strict inequality implies the weak one
required by the lemma.)

(ii) $\operatorname{Alternating}$ is $H_1\wedge H_2\wedge H_3$. Its conjunct $H_3$ contradicts the first conjunct
$\neg H_3$ of $\operatorname{RunMidA}$, $H_1$ contradicts the first conjunct $\neg H_1$ of $\operatorname{RunMidB}$,
and $H_2$ contradicts the first conjunct $\neg H_2$ of $\operatorname{RunMidD}$. This part uses none of the
hypotheses.
\end{proof}

\begin{remark}[what Theorem~\ref{pc:thm:trichotomy} does not state]\label{pc:rem:runs}
The Lean statement is the conjunction of the disjunction (i) and the implication (ii). It does \emph{not} state
that the three runs exclude each other. They do (a written step, not Lean-checked): if
$\alpha<\beta<\delta<\alpha+2\pi$, then $H_1\vee H_2\vee H_3$. Indeed $\alpha<\alpha+\pi<\alpha+2\pi$. If
$\alpha+\pi<\beta$, the antipode $\alpha+\pi$ gives $H_1$ (with $k=0$); if $\beta<\alpha+\pi<\delta$, it gives $H_2$;
if $\delta<\alpha+\pi$, it gives $H_3$, since $\alpha+\pi<\alpha+2\pi$. If $\alpha+\pi=\beta$, then
$\delta\in(\alpha+\pi,\alpha+2\pi)$, so $\delta-\pi\in(\alpha,\beta)$, and $\delta+\pi$ gives $H_1$ with $k=-1$. If
$\alpha+\pi=\delta$, then $\beta\in(\alpha,\alpha+\pi)$, so $\beta+\pi\in(\delta,\alpha+2\pi)$, and $\beta+\pi$ gives
$H_3$ with $k=0$. Any two of $\operatorname{RunMidA}$, $\operatorname{RunMidB}$, $\operatorname{RunMidD}$ together
give $\neg H_1\wedge\neg H_2\wedge\neg H_3$, which is therefore impossible. Hence, under the hypotheses of
Theorem~\ref{pc:thm:trichotomy}, exactly one of the four patterns holds.
\end{remark}

\subsection{The rules in pattern form}

In this subsection $c\in\R^2$, $r_a,r_b,r_d,\alpha,\beta,\delta\in\R$, $\operatorname{polar}$ is
\lean{PhaseClassification.polar}, and
\[
  a=\operatorname{polar}(c,r_a,\alpha),\qquad b=\operatorname{polar}(c,r_b,\beta),\qquad d=\operatorname{polar}(c,r_d,\delta).
\]

\begin{theorem}[rule A from the pattern; \leanref{PhaseClassification.orient_pos_of_alternating}{LEANEXT/phase/PhaseClassification.lean:301}]\label{pc:thm:rule-alt}
If $0<r_a$, $0<r_b$, $0<r_d$, $\alpha<\beta$, $\beta<\delta$, $\delta<\alpha+2\pi$ and
$\operatorname{Alternating}(\alpha,\beta,\delta)$, then $0<\orient(a,b,d)$.
\end{theorem}

\begin{proof}
By Theorem~\ref{pc:thm:alternating}, $\beta-\alpha<\pi$, $\delta-\beta<\pi$ and $\alpha+2\pi-\delta<\pi$; the
ordering gives $0<\beta-\alpha$, $0<\delta-\beta$ and $0<\alpha-\delta+2\pi$. Lemma~\ref{pc:lem:polar-facts}(c)
applies.
\end{proof}

\begin{theorem}[rule B from a run; \leanref{PhaseClassification.orient_pos_of_runMidB}{LEANEXT/phase/PhaseClassification.lean:310}, \leanref{PhaseClassification.orient_pos_of_runMidD}{LEANEXT/phase/PhaseClassification.lean:318}, \leanref{PhaseClassification.orient_pos_of_runMidA}{LEANEXT/phase/PhaseClassification.lean:342}]\label{pc:thm:rule-run}
Let $c$, $r_a,r_b,r_d,\alpha,\beta,\delta$ and $a,b,d$ be as above.
\begin{enumerate}[label=(\Alph*)]
\item[\textup{(B)}] If $0<r_a$, $0<r_d$, $\alpha<\beta$, $\beta<\delta$, $\delta<\alpha+2\pi$,
  $\alpha+2\pi-\delta\neq\pi$, $\operatorname{RunMidB}(\alpha,\beta,\delta)$, $r_a\le r_b$ and $r_d\le r_b$, then
  $0<\orient(a,b,d)$.
\item[\textup{(D)}] If $0<r_a$, $0<r_b$, $\alpha<\beta$, $\beta<\delta$, $\delta<\alpha+2\pi$, $\beta-\alpha\neq\pi$,
  $\operatorname{RunMidD}(\alpha,\beta,\delta)$, $r_b\le r_d$ and $r_a\le r_d$, then $0<\orient(a,b,d)$.
\item[\textup{(A)}] If $0<r_b$, $0<r_d$, $\alpha<\beta$, $\beta<\delta$, $\delta<\alpha+2\pi$, $\delta-\beta\neq\pi$,
  $\operatorname{RunMidA}(\alpha,\beta,\delta)$, $r_b\le r_a$ and $r_d\le r_a$, then $0<\orient(a,b,d)$.
\end{enumerate}
In each case the radius of the middle point of the run is not assumed positive.
\end{theorem}

\begin{proof}
(B) Theorem~\ref{pc:thm:runB} turns $\operatorname{RunMidB}(\alpha,\beta,\delta)$ into $\delta-\alpha<\pi$. With
$0<\beta-\alpha$ and $0<\delta-\beta$ from the ordering, Lemma~\ref{pc:lem:polar-facts}(e) gives the claim.

(D) Put $(\alpha',\beta',\delta')=(\beta,\delta,\alpha+2\pi)$. Then $\alpha'<\beta'$ ($\beta<\delta$),
$\beta'<\delta'$ ($\delta<\alpha+2\pi$) and $\delta'<\alpha'+2\pi$ ($\alpha+2\pi<\beta+2\pi$, from $\alpha<\beta$).
Further $\alpha'+2\pi-\delta'=\beta-\alpha\neq\pi$, and $\operatorname{RunMidB}(\alpha',\beta',\delta')$ holds by
Lemma~\ref{pc:lem:shift}(c). Apply (B) to these directions with the radii $(r_b,r_d,r_a)$ in place of
$(r_a,r_b,r_d)$; its radius hypotheses $0<r_b$, $0<r_a$, $r_b\le r_d$, $r_a\le r_d$ are among ours. It gives
$0<\orient(\operatorname{polar}(c,r_b,\beta),\operatorname{polar}(c,r_d,\delta),\operatorname{polar}(c,r_a,\alpha+2\pi))$.
By Lemma~\ref{pc:lem:periodic}(b) the last point is $a$, so $0<\orient(b,d,a)=\orient(a,b,d)$ by cyclic invariance.

(A) Put $(\alpha',\beta',\delta')=(\delta,\alpha+2\pi,\beta+2\pi)$. Then $\delta<\alpha+2\pi$,
$\alpha+2\pi<\beta+2\pi$ and $\beta+2\pi<\delta+2\pi$, so $\alpha'<\beta'<\delta'<\alpha'+2\pi$; further
$\alpha'+2\pi-\delta'=\delta-\beta\neq\pi$, and $\operatorname{RunMidB}(\alpha',\beta',\delta')$ holds by
Lemma~\ref{pc:lem:shift}(d). Apply (B) with the radii $(r_d,r_a,r_b)$ in place of $(r_a,r_b,r_d)$; its radius
hypotheses $0<r_d$, $0<r_b$, $r_d\le r_a$, $r_b\le r_a$ are among ours. It gives
$0<\orient(\operatorname{polar}(c,r_d,\delta),\operatorname{polar}(c,r_a,\alpha+2\pi),
\operatorname{polar}(c,r_b,\beta+2\pi))=\orient(d,a,b)$ by Lemma~\ref{pc:lem:periodic}(b), and
$\orient(d,a,b)=\orient(a,b,d)$ by cyclic invariance.
\end{proof}

\begin{theorem}[sign from the phase pattern, sorted directions; \leanref{PhaseClassification.orient_pos_of_phase}{LEANEXT/phase/PhaseClassification.lean:372}]\label{pc:thm:phase}
Assume $0<r_a$, $0<r_b$, $0<r_d$, $\alpha<\beta$, $\beta<\delta$, $\delta<\alpha+2\pi$, $\beta-\alpha\neq\pi$,
$\delta-\beta\neq\pi$, $\alpha+2\pi-\delta\neq\pi$, and
\begin{itemize}
\item $\operatorname{RunMidA}(\alpha,\beta,\delta)\Rightarrow(r_b\le r_a\wedge r_d\le r_a)$,
\item $\operatorname{RunMidB}(\alpha,\beta,\delta)\Rightarrow(r_a\le r_b\wedge r_d\le r_b)$,
\item $\operatorname{RunMidD}(\alpha,\beta,\delta)\Rightarrow(r_b\le r_d\wedge r_a\le r_d)$.
\end{itemize}
Then $0<\orient(a,b,d)$.
\end{theorem}

\begin{proof}
By Theorem~\ref{pc:thm:trichotomy}(i) one of four cases holds. If $\operatorname{Alternating}(\alpha,\beta,\delta)$,
apply Theorem~\ref{pc:thm:rule-alt}. If $\operatorname{RunMidA}$, the first implication gives $r_b\le r_a$ and
$r_d\le r_a$, and Theorem~\ref{pc:thm:rule-run}(A) applies (with $0<r_b$, $0<r_d$, $\delta-\beta\neq\pi$). If
$\operatorname{RunMidB}$, the second gives $r_a\le r_b$, $r_d\le r_b$, and (B) applies (with $0<r_a$, $0<r_d$,
$\alpha+2\pi-\delta\neq\pi$). If $\operatorname{RunMidD}$, the third gives $r_b\le r_d$, $r_a\le r_d$, and (D)
applies (with $0<r_a$, $0<r_b$, $\beta-\alpha\neq\pi$).
\end{proof}

The two patterns have a direct geometric meaning.

\begin{proposition}[alternation puts the centre inside the triangle; \leanref{PhaseClassification.centre_inside_of_alternating}{LEANEXT/phase/PhaseClassification.lean:414}]\label{pc:prop:centre}
If $0<r_a$, $0<r_b$, $0<r_d$, $\alpha<\beta$, $\beta<\delta$, $\delta<\alpha+2\pi$ and
$\operatorname{Alternating}(\alpha,\beta,\delta)$, then
\[
  0<\orient(a,b,c),\qquad 0<\orient(b,d,c),\qquad 0<\orient(d,a,c).
\]
\end{proposition}

\begin{proof}
By Theorem~\ref{pc:thm:alternating}, $\beta-\alpha$, $\delta-\beta$ and $\alpha+2\pi-\delta$ lie in $(0,\pi)$. By
Lemma~\ref{pc:lem:polar-facts}(b), applied to the pairs $(a,b)$, $(b,d)$ and $(d,a)$,
\[
  \orient(a,b,c)=r_ar_b\sin(\beta-\alpha),\quad \orient(b,d,c)=r_br_d\sin(\delta-\beta),\quad
  \orient(d,a,c)=r_dr_a\sin(\alpha-\delta),
\]
and $\sin(\alpha-\delta)=\sin(\alpha+2\pi-\delta)$. All three sines and all radii are positive.
\end{proof}

\begin{proposition}[a run with middle $\beta$ lies in an open half-plane through the centre; \leanref{PhaseClassification.halfplane_of_runMidB}{LEANEXT/phase/PhaseClassification.lean:431}]\label{pc:prop:halfplane}
If $0<r_a$, $0<r_b$, $0<r_d$, $\alpha<\beta$, $\beta<\delta$, $\delta<\alpha+2\pi$, $\alpha+2\pi-\delta\neq\pi$ and
$\operatorname{RunMidB}(\alpha,\beta,\delta)$, then there is $\varphi\in\R$ such that for every
$p\in\{a,b,d\}$,
\[
  0<(p_0-c_0)\cos\varphi+(p_1-c_1)\sin\varphi .
\]
\end{proposition}

\begin{proof}
By Theorem~\ref{pc:thm:runB}, $\delta-\alpha<\pi$. Take $\varphi=(\alpha+\delta)/2$. For $p=\operatorname{polar}(c,r,\theta)$,
\[
  (p_0-c_0)\cos\varphi+(p_1-c_1)\sin\varphi=r(\cos\theta\cos\varphi+\sin\theta\sin\varphi)=r\cos(\theta-\varphi).
\]
For $\theta\in\{\alpha,\beta,\delta\}\subseteq[\alpha,\delta]$ we have
$\abs{\theta-\varphi}\le(\delta-\alpha)/2<\pi/2$, so $\cos(\theta-\varphi)>0$; the radius $r\in\{r_a,r_b,r_d\}$ is
positive.
\end{proof}

\subsection{Arbitrary real directions}

\begin{theorem}[sign for arbitrary directions; \leanref{PhaseClassification.orient_sign_of_phase}{LEANEXT/phase/PhaseClassification.lean:496}]\label{pc:thm:sign}
Let $c\in\R^2$ and $r_a,r_b,r_d,\alpha,\beta,\delta\in\R$ with $0<r_a$, $0<r_b$, $0<r_d$, and let
$a=\operatorname{polar}(c,r_a,\alpha)$, $b=\operatorname{polar}(c,r_b,\beta)$, $d=\operatorname{polar}(c,r_d,\delta)$,
where $\operatorname{polar}$ is \lean{PhaseClassification.polar}. Put $\beta'=\operatorname{toIcoMod}(\alpha,\beta)$
and $\delta'=\operatorname{toIcoMod}(\alpha,\delta)$, and assume $\beta'\neq\alpha$ and $\delta'\neq\alpha$. Then
both of the following hold.
\begin{enumerate}[label=(\roman*)]
\item If $\beta'<\delta'$, $\beta'-\alpha\neq\pi$, $\delta'-\beta'\neq\pi$, $\alpha+2\pi-\delta'\neq\pi$,
  $\operatorname{RunMidA}(\alpha,\beta',\delta')\Rightarrow(r_b\le r_a\wedge r_d\le r_a)$,
  $\operatorname{RunMidB}(\alpha,\beta',\delta')\Rightarrow(r_a\le r_b\wedge r_d\le r_b)$ and
  $\operatorname{RunMidD}(\alpha,\beta',\delta')\Rightarrow(r_b\le r_d\wedge r_a\le r_d)$, then $0<\orient(a,b,d)$.
\item If $\delta'<\beta'$, $\delta'-\alpha\neq\pi$, $\beta'-\delta'\neq\pi$, $\alpha+2\pi-\beta'\neq\pi$,
  $\operatorname{RunMidA}(\alpha,\delta',\beta')\Rightarrow(r_d\le r_a\wedge r_b\le r_a)$,
  $\operatorname{RunMidB}(\alpha,\delta',\beta')\Rightarrow(r_a\le r_d\wedge r_b\le r_d)$ and
  $\operatorname{RunMidD}(\alpha,\delta',\beta')\Rightarrow(r_d\le r_b\wedge r_a\le r_b)$, then $\orient(a,b,d)<0$.
\end{enumerate}
\end{theorem}

\begin{proof}
By (T1), $\alpha\le\beta'<\alpha+2\pi$ and $\alpha\le\delta'<\alpha+2\pi$; with $\beta'\neq\alpha$ and
$\delta'\neq\alpha$ we get $\alpha<\beta'$ and $\alpha<\delta'$. By Lemma~\ref{pc:lem:periodic}(c),
$b=\operatorname{polar}(c,r_b,\beta')$ and $d=\operatorname{polar}(c,r_d,\delta')$.

(i) The directions $(\alpha,\beta',\delta')$ satisfy $\alpha<\beta'<\delta'<\alpha+2\pi$, and the remaining
hypotheses of (i) are exactly the gap conditions and radius implications of Theorem~\ref{pc:thm:phase} for these
directions and the radii $(r_a,r_b,r_d)$. Hence $0<\orient(a,b,d)$.

(ii) Apply Theorem~\ref{pc:thm:phase} to the directions $(\alpha,\delta',\beta')$, which satisfy
$\alpha<\delta'<\beta'<\alpha+2\pi$, and the radii $(r_a,r_d,r_b)$ in place of $(r_a,r_b,r_d)$. Its gap conditions
are $\delta'-\alpha\neq\pi$, $\beta'-\delta'\neq\pi$, $\alpha+2\pi-\beta'\neq\pi$, and its radius implications,
with $r_b$ and $r_d$ interchanged, are exactly the three implications of (ii). Hence
$0<\orient(a,\operatorname{polar}(c,r_d,\delta'),\operatorname{polar}(c,r_b,\beta'))=\orient(a,d,b)=-\orient(a,b,d)$,
so $\orient(a,b,d)<0$.
\end{proof}

\begin{remark}[scope of Theorem~\ref{pc:thm:sign}]\label{pc:rem:sign-scope}
The case $\beta'=\delta'$ (the points $b$ and $d$ on one ray from $c$) is not covered, and $\beta'=\alpha$ or
$\delta'=\alpha$ is excluded by hypothesis. In Lean this theorem is not combined with the layer theorems of
Subsection~\ref{pc:sub:layers} (\dirn{LEANEXT/capstone/REPORT.md}, open issues).
\end{remark}

\subsection{The discrete dichotomy of cyclic words}

Let $\mathrm{Fin}\,6=\{0,1,\dots,5\}$ with addition modulo $6$, and $\mathrm{Bool}=\{\mathrm{true},\mathrm{false}\}$
with negation $\neg$. A Boolean value $x$ used as a statement means $x=\mathrm{true}$.

\begin{theorem}[discrete dichotomy, computer-assisted in Lean; written proof by hand; \leanref{PhaseClassification.word_dichotomy}{LEANEXT/phase/PhaseClassification.lean:539}]\label{pc:thm:word}
For every $s\colon\mathrm{Fin}\,6\to\mathrm{Bool}$ such that $s(i+3)=\neg s(i)$ for all $i\in\mathrm{Fin}\,6$:
\begin{enumerate}[label=(\roman*)]
\item $\bigl(\forall i:\ s(i+1)=\neg s(i)\bigr)\ \vee\ \bigl(\exists i:\ s(i)\wedge s(i+1)\wedge s(i+2)\bigr)$, and
\item $\neg\Bigl(\bigl(\forall i:\ s(i+1)=\neg s(i)\bigr)\wedge\bigl(\exists i:\ s(i)\wedge s(i+1)\wedge s(i+2)\bigr)\Bigr)$.
\end{enumerate}
\end{theorem}

\paragraph{How Lean checks it.} The statement quantifies over the finite type of functions
$\mathrm{Fin}\,6\to\mathrm{Bool}$, which has $2^6=64$ elements, and inside over $\mathrm{Fin}\,6$; all atoms are
equalities of Booleans. Lean therefore has a decision procedure for it. For each of the $64$ functions $s$, the
procedure decides the implication from the premise (the six equalities $s(i+3)=\neg s(i)$) to the conclusion
(i)~$\wedge$~(ii), whose atoms are the six equalities $s(i+1)=\neg s(i)$ of the alternation condition and the six
triples $s(i),s(i+1),s(i+2)$ of the run condition. Lean's decision procedure for an implication evaluates the
premise first and the conclusion only when the premise holds, so the conclusion is evaluated only for the $8$
functions listed in the proof below. The tactic \lean{decide} evaluates the procedure during elaboration, and the
Lean kernel evaluates it again, to \lean{true}, when it type-checks the resulting proof term
(\lean{of_decide_eq_true} applied to \lean{Eq.refl}~\lean{true}). No
\lean{native_decide} is used, so no \lean{Lean.ofReduceBool} axiom enters; \dirn{LEANEXT/phase/logs/axioms.log}
records the three standard axioms. Mutant M7 of the mutation script \dirn{LEANEXT/phase/mutants.py} replaces the
premise by \lean{True}; the mutation log \dirn{LEANEXT/phase/logs/mutants.log}, which calls this mutant ``drop
antipodal premise in \lean{word_dichotomy}'', records that it fails to compile with the error ``Tactic
\lean{decide} proved that the proposition'' (the script logs only the first line of the first error message;
Lean's message continues with the proposition and ``is false''). The mutated statement is indeed false: for
example, the word $PPAPPA$ (with $P=\mathrm{true}$, $A=\mathrm{false}$; it violates the premise, as $s(3)=s(0)$) is
neither alternating nor contains a cyclic $PPP$.

\begin{proof}[Written proof (by hand)]
Let $s$ satisfy the premise and put $x=s(0)$, $y=s(1)$, $z=s(2)$. The premise gives $s(3)=\neg x$, $s(4)=\neg y$,
$s(5)=\neg z$, so $s=(x,y,z,\neg x,\neg y,\neg z)$ and there are exactly eight such words. Conversely each of
these eight words satisfies the premise, since $s(i+3)=\neg s(i)$ for $i=0,1,2$ gives it for $i=3,4,5$ as well. In
the notation $P=\mathrm{true}$, $A=\mathrm{false}$, and listing positions $0,\dots,5$, they are:
\begin{center}
\begin{tabular}{lll}
\toprule
$(x,y,z)$ & word & property \\
\midrule
$(P,P,P)$ & $PPPAAA$ & run $PPP$ at $i=0$ \\
$(P,P,A)$ & $PPAAAP$ & run at $i=5$ (positions $5,0,1$) \\
$(P,A,P)$ & $PAPAPA$ & alternating \\
$(P,A,A)$ & $PAAAPP$ & run at $i=4$ (positions $4,5,0$) \\
$(A,P,P)$ & $APPPAA$ & run at $i=1$ \\
$(A,P,A)$ & $APAPAP$ & alternating \\
$(A,A,P)$ & $AAPPPA$ & run at $i=2$ \\
$(A,A,A)$ & $AAAPPP$ & run at $i=3$ \\
\bottomrule
\end{tabular}
\end{center}
This proves (i). For (ii): if $s(i+1)=\neg s(i)$ for all $i$, then $s(i)$ and $s(i+1)$ are never both true, so
there is no cyclic run. (A shorter argument for (i): if $s$ is not alternating, then $s(i+1)=s(i)$ for some $i$, and
then also $s(i+4)=\neg s(i+1)=\neg s(i)=s(i+3)$; one of the equal pairs $\{i,i+1\}$, $\{i+3,i+4\}$ consists of two
$P$'s, say at $j,j+1$. If $s(j+2)=P$ there is a run at $j$; otherwise $s(j-1)=s(j+5)=\neg s(j+2)=P$ and there is a
run at $j-1$.)
\end{proof}

\begin{remark}[scope of Theorem~\ref{pc:thm:word}]\label{pc:rem:word-scope}
This is a statement about Boolean words only. The layered encoders read the six directions $\alpha,\beta,\delta,
\alpha+\pi,\beta+\pi,\delta+\pi$, sorted counterclockwise, as such a word ($P$ for a direction, $A$ for an
antipode). That this sorted word satisfies $s(i+3)=\neg s(i)$, and that a cyclic $PPP$ in it is
$\operatorname{RunMidA}$, $\operatorname{RunMidB}$ or $\operatorname{RunMidD}$ of the middle point, is not proved in
\dirn{LEANEXT} (\dirn{LEANEXT/phase/REPORT.md}, ``Not in Lean'', item 2). For encoder~B, the Lean version of the
encoder in the ENCBLEAN package has a Lean-checked counterpart of the second link that does not use
Theorem~\ref{pc:thm:word}: \lean{EncB.wordSpec} (\dirn{ENCBLEAN/EncB/SoundWord.lean:222}) derives the rule-B
pattern \lean{EncB.BlockCCW} from three cyclically consecutive points of the encoder's sorted word, and, in
\dirn{ENCBLEAN/geo/ES8Geo/ArcBridge.lean}, \lean{EncB.Geo.Arc.arcFrame} (line 393) identifies, for a polar labelling
of the layers, the arc tests of this pattern with the tests $\operatorname{arcHasAntipode}$ for true directions
$\alpha<\beta<\delta<\alpha+2\pi$, from which \lean{EncB.Geo.labelling_ruleB} (line 473) obtains
$\operatorname{RunMidB}(\alpha,\beta,\delta)$ and applies Theorem~\ref{pc:thm:rule-run}(B). These three
declarations depend only on \lean{propext}, \lean{Classical.choice} and \lean{Quot.sound}: for \lean{wordSpec} by
the \texttt{\#guard\_msgs} pin that follows it (\dirn{ENCBLEAN/EncB/SoundWord.lean}, lines 228--229; the module is
built in \dirn{ENCBLEAN/logs/m3_build.log}, line 9, and \dirn{ENCBLEAN/logs/axioms.log}, line 446, gives the same
list), and for \lean{arcFrame} and \lean{labelling_ruleB} by their \texttt{\#print axioms} output in
\dirn{ENCBLEAN/logs/m8_fix_build.log}, lines 35 and 29, and by the pin of \lean{EncB.Geo.modelOfLabelling} (which
uses \lean{labelling_ruleB}) described under Rigor and axioms. The case split for real directions that is
Lean-checked here is Theorem~\ref{pc:thm:trichotomy}.
\end{remark}

\subsection{Rule B on the layers of a symmetric set}\label{pc:sub:layers}

In this subsection $S$ is a finite set of points, $c\in\R^2$, $m,K\in\N$, and $\rho=\operatorname{rotAbout}(c,2\pi/m)$
is the counterclockwise rotation about $c$ by $2\pi/m$ (Definition~\ref{fd:def:rot}; $m$ is cast to $\R$). We say that
\emph{the layer hypotheses \textup{(L)} hold} if
\[
  0<m,\qquad K\le 2m,\qquad c\notin S,\qquad \operatorname{Invariant}(\rho,S),\qquad \neg\operatorname{HasConvexNGon}(K,S),
\]
where $\operatorname{Invariant}(\rho,S)$ means $\rho(p)\in S$ for all $p\in S$ (Definition~\ref{fd:def:inv}) and
$\operatorname{HasConvexNGon}(K,S)$ means that some $K$-element subset of $S$ is in convex position
(Definition~\ref{fd:def:ngon}). In this subsection $\operatorname{polar}$ is \lean{LayerLemma.polar}; by
Lemma~\ref{pc:lem:helpers}(a) it is the same function as \lean{PhaseClassification.polar}, so the results of the
previous subsections apply to it.

We use two facts from the layer sections. First, $\peel(S,n)\subseteq S$ and $\hv(T)\subseteq T$ for all $n$ and
all finite $T$ (Lean \lean{peel_subset}, \lean{hullVerts_subset}): $\hv(T)$ is defined as a subset of $T$
(Definition~\ref{ll:def:hv}), and $\peel(S,n+1)=\peel(S,n)\setminus\hv(\peel(S,n))\subseteq\peel(S,n)$
(Definition~\ref{ll:def:peel}), so induction on $n$ gives the first inclusion. Second, the radius order of the
layers (Theorem~\ref{lg:thm:radius}; Lean \lean{ErdosSzekeres.LayerLemma.dist2_le_of_later},
\dirn{LAYERLEAN/LayerLemma.lean:1012}, repeated at \dirn{LEANEXT/linked/Linked.lean:1012} and
\dirn{LEANEXT/capstone/Capstone2.lean:1020}): \emph{if \textup{(L)} holds, $n\le n'$, $x\in\hv(\peel(S,n))$ and
$y\in\peel(S,n')$, then $\operatorname{dist2}(c,y)\le\operatorname{dist2}(c,x)$.} We also use the elementary fact
that $0<r$, $0<r'$ and $r^2\le r'^2$ imply $r\le r'$ (otherwise $r^2-r'^2=(r-r')(r+r')>0$).

\begin{theorem}[rule B on the convex layers; \leanref{ErdosSzekeres.LayerLemma.orient_pos_of_layers}{LEANEXT/linked/Linked.lean:1398} for (a), and the six declarations of \lean{ErdosSzekeres.Capstone} cited in (b)--(e)]\label{pc:thm:layers}
Assume \textup{(L)}. Let $n_a,n_b,n_d\in\N$ and $a,b,d\in\R^2$. For real $r_a,r_b,r_d,\alpha,\beta,\delta$ let
\textup{(P)} denote the conditions
\[
  0<r_a,\quad 0<r_b,\quad 0<r_d,\quad a=\operatorname{polar}(c,r_a,\alpha),\quad b=\operatorname{polar}(c,r_b,\beta),
  \quad d=\operatorname{polar}(c,r_d,\delta),
\]
and for real $\alpha,\beta,\delta$ let \textup{(P$\surd$)} denote the conditions
\[
  a=\operatorname{polar}\bigl(c,\sqrt{\operatorname{dist2}(c,a)},\alpha\bigr),\quad
  b=\operatorname{polar}\bigl(c,\sqrt{\operatorname{dist2}(c,b)},\beta\bigr),\quad
  d=\operatorname{polar}\bigl(c,\sqrt{\operatorname{dist2}(c,d)},\delta\bigr).
\]
\begin{enumerate}[label=(\alph*)]
\item \textup{(\lean{orient_pos_of_layers})} If $a\in\peel(S,n_a)$, $b\in\hv(\peel(S,n_b))$, $d\in\peel(S,n_d)$,
  $n_b\le n_a$, $n_b\le n_d$, and $r_a,r_b,r_d,\alpha,\beta,\delta\in\R$ satisfy \textup{(P)}, $0<\beta-\alpha$,
  $0<\delta-\beta$ and $\delta-\alpha<\pi$, then $0<\orient(a,b,d)$.
\item \textup{(\leanref{ErdosSzekeres.Capstone.orient_pos_of_layers_runMidB}{LEANEXT/capstone/Capstone.lean:1977}.)}
  If $a\in\peel(S,n_a)$, $b\in\hv(\peel(S,n_b))$,
  $d\in\peel(S,n_d)$, $n_b\le n_a$, $n_b\le n_d$, and $r_a,r_b,r_d,\alpha,\beta,\delta\in\R$ satisfy \textup{(P)},
  $\alpha<\beta$, $\beta<\delta$, $\delta<\alpha+2\pi$, $\alpha+2\pi-\delta\neq\pi$ and
  $\operatorname{RunMidB}(\alpha,\beta,\delta)$, then $0<\orient(a,b,d)$.
\item \textup{(\leanref{ErdosSzekeres.Capstone.orient_pos_of_layers_runMidD}{LEANEXT/capstone/Capstone2.lean:2011}.)}
  If $a\in\peel(S,n_a)$, $b\in\peel(S,n_b)$,
  $d\in\hv(\peel(S,n_d))$, $n_d\le n_a$, $n_d\le n_b$, and $r_a,r_b,r_d,\alpha,\beta,\delta\in\R$ satisfy
  \textup{(P)}, $\alpha<\beta$, $\beta<\delta$, $\delta<\alpha+2\pi$, $\beta-\alpha\neq\pi$ and
  $\operatorname{RunMidD}(\alpha,\beta,\delta)$, then $0<\orient(a,b,d)$.
\item \textup{(\leanref{ErdosSzekeres.Capstone.orient_pos_of_layers_runMidA}{LEANEXT/capstone/Capstone2.lean:2031}.)}
  If $a\in\hv(\peel(S,n_a))$, $b\in\peel(S,n_b)$,
  $d\in\peel(S,n_d)$, $n_a\le n_b$, $n_a\le n_d$, and $r_a,r_b,r_d,\alpha,\beta,\delta\in\R$ satisfy \textup{(P)},
  $\alpha<\beta$, $\beta<\delta$, $\delta<\alpha+2\pi$, $\delta-\beta\neq\pi$ and
  $\operatorname{RunMidA}(\alpha,\beta,\delta)$, then $0<\orient(a,b,d)$.
\item \textup{(The \lean{_intrinsic} versions.)} Let $\alpha,\beta,\delta\in\R$. Then $0<\orient(a,b,d)$ in each
  of the following three cases:
  \begin{enumerate}[label=(e\arabic*)]
  \item \textup{(\leanref{ErdosSzekeres.Capstone.orient_pos_of_layers_runMidB_intrinsic}{LEANEXT/capstone/Capstone2.lean:2082}.)}
    $a\in\peel(S,n_a)$, $b\in\hv(\peel(S,n_b))$,
    $d\in\peel(S,n_d)$, $n_b\le n_a$, $n_b\le n_d$, \textup{(P$\surd$)}, $\alpha<\beta$, $\beta<\delta$,
    $\delta<\alpha+2\pi$, $\alpha+2\pi-\delta\neq\pi$ and $\operatorname{RunMidB}(\alpha,\beta,\delta)$;
  \item \textup{(\leanref{ErdosSzekeres.Capstone.orient_pos_of_layers_runMidD_intrinsic}{LEANEXT/capstone/Capstone2.lean:2098}.)}
    $a\in\peel(S,n_a)$, $b\in\peel(S,n_b)$,
    $d\in\hv(\peel(S,n_d))$, $n_d\le n_a$, $n_d\le n_b$, \textup{(P$\surd$)}, $\alpha<\beta$, $\beta<\delta$,
    $\delta<\alpha+2\pi$, $\beta-\alpha\neq\pi$ and $\operatorname{RunMidD}(\alpha,\beta,\delta)$;
  \item \textup{(\leanref{ErdosSzekeres.Capstone.orient_pos_of_layers_runMidA_intrinsic}{LEANEXT/capstone/Capstone2.lean:2115}.)}
    $a\in\hv(\peel(S,n_a))$, $b\in\peel(S,n_b)$,
    $d\in\peel(S,n_d)$, $n_a\le n_b$, $n_a\le n_d$, \textup{(P$\surd$)}, $\alpha<\beta$, $\beta<\delta$,
    $\delta<\alpha+2\pi$, $\delta-\beta\neq\pi$ and $\operatorname{RunMidA}(\alpha,\beta,\delta)$.
  \end{enumerate}
\end{enumerate}
General position of $S$ is not assumed.
\end{theorem}

\begin{proof}
(a) Apply the radius order with $(n,n',x,y)=(n_b,n_a,b,a)$ and with $(n_b,n_d,b,d)$: this gives
$\operatorname{dist2}(c,a)\le\operatorname{dist2}(c,b)$ and $\operatorname{dist2}(c,d)\le\operatorname{dist2}(c,b)$. By
(P) and Lemma~\ref{pc:lem:helpers}(c) these read $r_a^2\le r_b^2$ and $r_d^2\le r_b^2$, hence $r_a\le r_b$ and
$r_d\le r_b$ as all radii are positive. Rule~B, Lemma~\ref{pc:lem:polar-facts}(e) (in Lean the \lean{LayerLemma}
copy, Lemma~\ref{lg:lem:ruleB}), with $0<r_a$, $0<r_d$, $0<\beta-\alpha$, $0<\delta-\beta$ and $\delta-\alpha<\pi$
gives $0<\orient(a,b,d)$.

(b) By Theorem~\ref{pc:thm:runB} (with $\alpha<\beta<\delta<\alpha+2\pi$ and $\alpha+2\pi-\delta\neq\pi$),
$\operatorname{RunMidB}(\alpha,\beta,\delta)$ gives $\delta-\alpha<\pi$. Together with $0<\beta-\alpha$ and
$0<\delta-\beta$, part (a) applies.

(c) The radius order with $(n_d,n_a,d,a)$ and $(n_d,n_b,d,b)$ gives $\operatorname{dist2}(c,a)\le\operatorname{dist2}(c,d)$
and $\operatorname{dist2}(c,b)\le\operatorname{dist2}(c,d)$, so $r_a^2\le r_d^2$ and $r_b^2\le r_d^2$, hence
$r_a\le r_d$ and $r_b\le r_d$. Theorem~\ref{pc:thm:rule-run}(D), with $0<r_a$, $0<r_b$, the ordering,
$\beta-\alpha\neq\pi$, $\operatorname{RunMidD}(\alpha,\beta,\delta)$, $r_b\le r_d$ and $r_a\le r_d$, gives
$0<\orient(a,b,d)$.

(d) The radius order with $(n_a,n_b,a,b)$ and $(n_a,n_d,a,d)$ gives $r_b^2\le r_a^2$ and $r_d^2\le r_a^2$, hence
$r_b\le r_a$ and $r_d\le r_a$. Theorem~\ref{pc:thm:rule-run}(A), with $0<r_b$, $0<r_d$, the ordering,
$\delta-\beta\neq\pi$, $\operatorname{RunMidA}(\alpha,\beta,\delta)$, $r_b\le r_a$ and $r_d\le r_a$, gives
$0<\orient(a,b,d)$.

(e) In each case $a,b,d\in S$: a point of $\peel(S,n)$ lies in $S$, and a point of $\hv(\peel(S,n))$ lies in
$\peel(S,n)\subseteq S$. As $c\notin S$, Lemma~\ref{pc:lem:helpers}(f) gives $a,b,d\neq c$, and
Lemma~\ref{pc:lem:helpers}(e) gives $0<\sqrt{\operatorname{dist2}(c,a)}$, $0<\sqrt{\operatorname{dist2}(c,b)}$ and
$0<\sqrt{\operatorname{dist2}(c,d)}$. So (P) holds with $r_a=\sqrt{\operatorname{dist2}(c,a)}$,
$r_b=\sqrt{\operatorname{dist2}(c,b)}$, $r_d=\sqrt{\operatorname{dist2}(c,d)}$, and (e1), (e2), (e3) follow from (b),
(c), (d) respectively.
\end{proof}

\begin{corollary}[rule A restated; \leanref{ErdosSzekeres.Capstone.orient_pos_of_alternating_polar}{LEANEXT/capstone/Capstone.lean:1992}]\label{pc:cor:alt-polar}
Let $c\in\R^2$ and $r_a,r_b,r_d,\alpha,\beta,\delta\in\R$ with $0<r_a$, $0<r_b$, $0<r_d$, $\alpha<\beta$,
$\beta<\delta$, $\delta<\alpha+2\pi$ and $\operatorname{Alternating}(\alpha,\beta,\delta)$. Then, with
$\operatorname{polar}$ the constant \lean{LayerLemma.polar},
$0<\orient(\operatorname{polar}(c,r_a,\alpha),\operatorname{polar}(c,r_b,\beta),\operatorname{polar}(c,r_d,\delta))$.
No layer hypotheses are involved.
\end{corollary}

\begin{proof}
By Lemma~\ref{pc:lem:helpers}(a) this is Theorem~\ref{pc:thm:rule-alt}. (The Lean proof repeats the proof of
Theorem~\ref{pc:thm:rule-alt} with the \lean{LayerLemma} copy of rule~A, Lemma~\ref{lg:lem:ruleA}.)
\end{proof}

\begin{remark}[what the layer theorems assume]\label{pc:rem:layers-scope}
In Theorem~\ref{pc:thm:layers} the polar forms (with radius variables in (a)--(d), with the intrinsic radii
$\sqrt{\operatorname{dist2}(c,\cdot)}$ in (e)) and the conditions on the given directions are hypotheses: in (a)
the conditions $0<\beta-\alpha$, $0<\delta-\beta$ and $\delta-\alpha<\pi$; in (b)--(e) the strict order
$\alpha<\beta<\delta<\alpha+2\pi$, the gap condition $\neq\pi$ and the pattern. They are not derived from the
symmetry or from general position. Directions with (P$\surd$) exist for every point other than $c$
(Lemma~\ref{pc:lem:helpers}(d)), but choosing them in sorted order and excluding gaps equal to $\pi$ is left to the
user of these theorems; the reduction of arbitrary directions in Theorem~\ref{pc:thm:sign} is not combined with
the layer hypotheses in Lean. The middle point (the point $b$ of the middle direction in (a), the middle point of
the run in (b)--(e)) is required to be a hull vertex of a layer whose index is at most the indices of the layers
of the other two points, and the conclusion is about the fixed order $(a,b,d)$.
\end{remark}

\paragraph{Coverage.} All declarations of \dirn{LEANEXT/phase/PhaseClassification.lean} are covered:
\lean{polar} (Definition~\ref{pc:def:polar}); \lean{orient_polar}, \lean{orient_polar_centre},
\lean{orient_polar_centre_inside}, \lean{sin_add_lt}, \lean{orient_polar_halfplane}
(Lemma~\ref{pc:lem:polar-facts}; except \lean{orient_polar_centre}, these are the copies of lines 931--996 of
\dirn{LAYERLEAN/LayerLemma.lean}, whose originals also have written proofs in the section on layer geometry);
\lean{polar_add_int_two_pi},
\lean{polar_add_two_pi}, \lean{polar_toIcoMod} (Lemma~\ref{pc:lem:periodic}); \lean{exists_polar}
(Lemma~\ref{pc:lem:exists-polar}); \lean{inArc}, \lean{arcHasAntipode}, \lean{Alternating}, \lean{RunMidB},
\lean{RunMidD}, \lean{RunMidA} (Definition~\ref{pc:def:arc}); \lean{inArc_of}, \lean{not_inArc}
(Lemma~\ref{pc:lem:inarc}); \lean{arc1_empty}, \lean{arc2_empty}, \lean{arc2_empty'}, \lean{arc3_empty},
\lean{arc3_empty'}, \lean{arc1_empty'} (Lemma~\ref{pc:lem:empty}); \lean{arc12_nonempty}
(Lemma~\ref{pc:lem:arc12}); \lean{alternating_iff_gaps_lt_pi} (Theorem~\ref{pc:thm:alternating});
\lean{alternating_exactly_one} (Proposition~\ref{pc:prop:exactly-one}); \lean{runMidB_iff}
(Theorem~\ref{pc:thm:runB}); \lean{gap_classification} (Proposition~\ref{pc:prop:gaps}); \lean{phase_trichotomy}
(Theorem~\ref{pc:thm:trichotomy}); \lean{orient_pos_of_alternating} (Theorem~\ref{pc:thm:rule-alt});
\lean{orient_pos_of_runMidB}, \lean{orient_pos_of_runMidD}, \lean{orient_pos_of_runMidA}
(Theorem~\ref{pc:thm:rule-run}); \lean{orient_pos_of_phase} (Theorem~\ref{pc:thm:phase});
\lean{centre_inside_of_alternating} (Proposition~\ref{pc:prop:centre}); \lean{halfplane_of_runMidB}
(Proposition~\ref{pc:prop:halfplane}); \lean{orient_sign_of_phase} (Theorem~\ref{pc:thm:sign});
\lean{word_dichotomy} (Theorem~\ref{pc:thm:word}, computer-assisted in Lean; written proof by hand). The new declarations of
\dirn{LEANEXT/linked/Linked.lean} are covered: \lean{dist2_polar} (Lemma~\ref{pc:lem:helpers}(c)) and
\lean{orient_pos_of_layers} (Theorem~\ref{pc:thm:layers}(a)); so are those of \dirn{LEANEXT/capstone/Capstone.lean}:
\lean{orient_pos_of_layers_runMidB} (Theorem~\ref{pc:thm:layers}(b)), \lean{orient_pos_of_alternating_polar}
(Corollary~\ref{pc:cor:alt-polar}) and \lean{polar_eq} (Lemma~\ref{pc:lem:helpers}(a)); and those of
\dirn{LEANEXT/capstone/Capstone2.lean}: \lean{orient_pos_of_layers_runMidD}, \lean{orient_pos_of_layers_runMidA}
(Theorem~\ref{pc:thm:layers}(c),(d)), \lean{polar_zero}, \lean{exists_angle}, \lean{sqrt_dist2_pos},
\lean{ne_of_mem} (Lemma~\ref{pc:lem:helpers}(b),(d),(e),(f)) and the three \lean{_intrinsic} theorems
(Theorem~\ref{pc:thm:layers}(e)). Deliberately not covered separately: lines 1--1384 of \dirn{Linked.lean} and the
corresponding lines of \dirn{Capstone.lean} and \dirn{Capstone2.lean} (copies of \dirn{LAYERLEAN/LayerLemma.lean}
and of \dirn{LEANEXT/merged/Merged.lean}, treated in the sections on the layer lemma, layer geometry and orbit
literals), and lines 1426--1965 of \dirn{Capstone.lean} and \dirn{Capstone2.lean} (verbatim copies of
\dirn{PhaseClassification.lean}, whose statements and proofs are those covered above, at line offset $+1420$).
The helper Lemma~\ref{pc:lem:shift} is a written step that makes explicit two \lean{have} steps of the Lean proofs
of \lean{orient_pos_of_runMidD} and \lean{orient_pos_of_runMidA}; Remarks~\ref{pc:rem:alt-boundary}, \ref{pc:rem:runB-example}, \ref{pc:rem:gaps} and~\ref{pc:rem:runs} contain
further written steps (with proofs) that are not Lean statements.
\par\endgroup

\section{Small corollaries for \texorpdfstring{$k=7$}{k=7} and \texorpdfstring{$k=8$}{k=8}}\label{ec:sec}
% Local to this section: the statement heads contain long unbreakable Lean names; \sloppy lets TeX break them at spaces.
\begingroup\sloppy

This section gives written proofs for a group of short Lean files about rotationally symmetric point sets.
\begin{enumerate}[label=(\arabic*)]
\item \emph{Conditional $ES(7)$ corollaries} (Subsection~\ref{ec:ssec:es7}). If the case classes
  $\mathrm{Fivefold35}$ and $\mathrm{Fourfold32}$ hold (they are hypotheses of the Lean theorems), then every finite
  set of points in general position that is invariant under the rotation by $2\pi/5$ about some point and has at
  least $32$ points, and every such set for the rotation by $2\pi/4$ with at least $29$ points, has $7$ points in
  convex position.
\item \emph{Rotations of order at least $8$} (Subsection~\ref{ec:ssec:order8}). A subset of $\C$ that is mapped into
  itself by $z\mapsto e^{2\pi i/m}z$ with $m\ge8$ and contains no $8$ points in convex position is contained in
  $\{0\}$. Also the threshold arithmetic for $m=4,5,6,7$: for all $t,e\in\N$ with $e\le1$, $65\le4t\iff t\ge17$,
  $65\le5t+e\iff t\ge13$, $65\le6t\iff t\ge11$ and $65\le7t+e\iff t\ge10$.
\item \emph{Orientations in inverse-radius polar form} (Subsection~\ref{ec:ssec:polar}): the sign of the orientation of
  three points $(\cos\varphi_k,\sin\varphi_k)/u_k$, and its consequences for layered configurations and for nested
  regular $m$-gons.
\item \emph{Size compatibility} (Subsections~\ref{ec:ssec:compat} and~\ref{ec:ssec:l0}): the orbit count
  $\card S\equiv[0\in S]\pmod m$ for finite sets invariant under an element of exact order $m$, its converse under
  explicit hypotheses, the classification of the $m$ with $2^{m-1}\bmod m\in\{0,1\}$, and the consequence for
  $m$-fold symmetric sets of $2^{m-1}$ points.
\end{enumerate}

\paragraph{Sources.} Paths are relative to \dirn{cloud/es7-2026-10-19/results}. We computed the SHA-256 of each file;
the table gives the first $12$ hexadecimal digits. The line ranges are the scope of this section; the Coverage paragraph
at the end lists what in them is not covered.
\begin{center}\footnotesize
\begin{tabular}{@{}lll@{}}
\toprule
Lean file (lines in scope) & SHA-256 & axioms log\\
\midrule
\dirn{LEANEXT/es7corr/ES7Corollaries.lean} (1148--1233) & \texttt{5883c8ebac55} & \dirn{LEANEXT/es7corr/axioms.log}\\
\dirn{ES8SYM/lean/C6check.lean} (1--123) & \texttt{92f15b5cc907} & \dirn{ES8SYM/lean/C6check.out}\\
\dirn{ES8SYM/lean/PolarOrient.lean} (1--132) & \texttt{6e54f5f83a19} & \dirn{ES8SYM/lean/PolarOrient.out}\\
\dirn{PRIMEFOLD/compat/lean/PCompatCore.lean} (1--291) & \texttt{ebc6aacefadb} & \dirn{PRIMEFOLD/compat/lean/PCompatCore.out}\\
\dirn{PRIMEFOLD/compat/lean/L0Check.lean} (1--152) & \texttt{05a67c3b8b16} & \dirn{PRIMEFOLD/compat/lean/L0Check.out}\\
\bottomrule
\end{tabular}
\end{center}
According to the records, all five files were compiled with the repository's pins, Lean~4.24.0 and Mathlib
\texttt{f897ebcf}; nothing was re-run for this document. The hashes of \dirn{LEANEXT/es7corr/ES7Corollaries.lean} and
of the two Lean files in \dirn{PRIMEFOLD/compat/lean} agree with the prefixes recorded in \dirn{LEANEXT/es7corr/REPORT.md}
and \dirn{PRIMEFOLD/README.md}. For the two \dirn{ES8SYM} files, \dirn{ES8SYM/README.md} (lines 258--260) records the
prefixes \texttt{8c090168} and \texttt{4d42e2fb} for the compiled files and states that the present copies differ from
them only in four comment lines; we could not check this, so the logs of these two files refer to the compiled
versions.

\paragraph{Axioms.} The logs list the following theorems, each depending on at most the standard axioms
\lean{propext}, \lean{Classical.choice} and \lean{Quot.sound}: the five $ES(7)$ corollaries,
Theorems~\ref{ec:thm:fivefold-erase}--\ref{ec:thm:fourfold} (\dirn{LEANEXT/es7corr/axioms.log}); \lean{RefuteC6.part_a}
and \lean{RefuteC6.thr_m4}--\lean{thr_m7} (\dirn{ES8SYM/lean/C6check.out}); the six theorems \lean{C7.orient_polar},
\lean{orient_polar_cross}, \lean{orient_layered}, \lean{orient_layered_scale}, \lean{nested_regular_orient},
\lean{orient_centre} (\dirn{ES8SYM/lean/PolarOrient.out}); \lean{PCompat.dvd_card}, \lean{card_mod},
\lean{exists_iff}, \lean{compat_even}, \lean{compat_odd}, \lean{classification}, and \lean{PCompat.ex341}, which
depends on \lean{propext} only (\dirn{PRIMEFOLD/compat/lean/PCompatCore.out}); \lean{L0Check.L0_rot}
(\dirn{PRIMEFOLD/compat/lean/L0Check.out}). Every other declaration covered here occurs in the statement or the proof
of one of these, so the same logs cover it, except \lean{RefuteC6.m4_with_centre_counter}, which has no axiom line (it is a closed
arithmetic statement whose Lean proof is \lean{omega}; the compile log of its file shows no error). No file uses \lean{sorry},
\lean{native_decide} or new axioms; \lean{PCompat.ex341} is checked by kernel evaluation (\lean{decide}).

\paragraph{Rigor and dependencies.} The theorems with a Lean reference are Lean-checked, with the axioms above; for the
two \dirn{ES8SYM} files this rests on the logs of the compiled versions, which according to the records differ from the
present copies only in comment lines (see Sources). The proofs below are written proofs (not refereed). Statements
tagged ``written step'' are not declarations of the five Lean files: Fact~\ref{ec:fact:prim} restates Mathlib lemmas,
and Corollary~\ref{ec:cor:planar} has no Lean counterpart. The theorems of Subsection~\ref{ec:ssec:es7} are
conditional on the propositions $\mathrm{Fivefold35}$ and $\mathrm{Fourfold32}$, and Lemma~\ref{ec:lem:pc-exist} and
Theorem~\ref{ec:thm:compat}(c) take hypotheses on $r$ that Lean does not discharge (Remark~\ref{ec:rem:compat}). We
use Theorem~\ref{ll:thm:reduction} (reduction to exactly $mt$ points) and Lemma~\ref{lg:lem:centre} (even orders
exclude the centre), and the conventions of Definitions~\ref{fd:def:point}, \ref{fd:def:orient}, \ref{fd:def:gp},
\ref{fd:def:cp}, \ref{fd:def:ngon}, \ref{fd:def:inv} and~\ref{fd:def:rot}; remarks also refer to
Lemma~\ref{ll:lem:minperiod}, Theorems~\ref{ll:thm:layers} and~\ref{ll:thm:card-peel}, Lemma~\ref{lg:lem:polar},
Lemmas~\ref{oc:lem:rot-inv}, \ref{oc:lem:perm} and~\ref{oc:lem:same-ray}, Theorem~\ref{oc:thm:phase-orbit} and
Definition~\ref{sr:def:classes}. The other subsections are self-contained apart from standard facts, which we state
where they are used.

%---------------------------------------------------------------------------------------------------------------------
\subsection{Conditional \texorpdfstring{$ES(7)$}{ES(7)} corollaries for fivefold and fourfold symmetry}\label{ec:ssec:es7}

The source is Section~11 of \dirn{LEANEXT/es7corr/ES7Corollaries.lean} (lines 1148--1233, namespace
\lean{ErdosSzekeres.LayerLemma}). Lines 1--1147 of this file coincide line by line with lines 1--1147 of
\dirn{LAYERLEAN/LayerLemma.lean} (we checked this with \texttt{diff}; the latter file has one more line, its closing
\texttt{end}). Hence the declarations \lean{Invariant}, \lean{rotAbout}, \lean{centre_not_mem_of_even} and
\lean{convexNGon_of_layers} used below are verbatim copies of those of \dirn{LAYERLEAN/LayerLemma.lean}, at the same
lines ($105$, $298$, $582$, $1066$).

\emph{Conventions.} $S$ is a finite set of points (Lean \texttt{Finset Point}) and $c$ a point. $\mathrm{rotAbout}(c,\theta)$
is the counterclockwise rotation by $\theta$ about $c$ (Definition~\ref{fd:def:rot}). ``$S$ is invariant under $f$''
means $f(p)\in S$ for every $p\in S$ (Definition~\ref{fd:def:inv}). General position is
Definition~\ref{fd:def:gp}, and $\mathrm{HasConvexNGon}(n,S)$ means that some $n$-element subset of $S$ is in convex
position (Definition~\ref{fd:def:ngon}; no emptiness condition). We write $S\setminus\{c\}$ for Lean's
\texttt{S.erase c}. In the statements below the angles are $2\pi/5$ and $2\pi/4$ with real numerals $5$ and $4$; in
Theorem~\ref{ll:thm:reduction} and Lemma~\ref{lg:lem:centre} (items (R) and (C) below) the angle is $2\pi/m$ with
$m\in\N$ cast to $\R$. For $m=5$ and $m=4$ these are the same real numbers (the Lean proofs pass between the two forms
with the cast-normalizing tactic \lean{exact_mod_cast}).

\begin{definition}[the case classes; \leanref{ErdosSzekeres.LayerLemma.Fivefold35}{LEANEXT/es7corr/ES7Corollaries.lean:1157}, \leanref{ErdosSzekeres.LayerLemma.Fourfold32}{LEANEXT/es7corr/ES7Corollaries.lean:1164}]\label{ec:def:cases}
$\mathrm{Fivefold35}$ is the proposition: for every finite set $S$ of points and every point $c$, if $S$ is in
general position, $\card S=35$, $c\notin S$, and $S$ is invariant under $\mathrm{rotAbout}(c,2\pi/5)$, then
$\mathrm{HasConvexNGon}(7,S)$.

$\mathrm{Fourfold32}$ is the proposition: for every finite set $S$ of points and every point $c$, if $S$ is in general
position, $\card S=32$, $c\notin S$, and $S$ is invariant under $\mathrm{rotAbout}(c,2\pi/4)$, then
$\mathrm{HasConvexNGon}(7,S)$.
\end{definition}

\begin{remark}[status of the case classes]\label{ec:rem:cases}
No Lean theorem proves $\mathrm{Fivefold35}$ or $\mathrm{Fourfold32}$. The theorems of this subsection take them as
hypotheses and are therefore conditional. In the project they are supported by computer-assisted results outside
Lean, which also rest on written arguments. For SAT certificates the project counts only the verdicts of the proof
checker cake\_lpr (\dirn{WRITEUP/SYMMETRIC.md}, Section~3). The records also show acceptances by a second checker,
lrat-check, which the project does not count: it is a secondary check only, and its unpatched upstream build is
unsound on the empty clause (\dirn{E2/logs/lrat_check_soundness.md}; \dirn{LRATCHECK20261105}).
\begin{itemize}
\item $\mathrm{Fivefold35}$: \dirn{SYM5} (\dirn{SYM5/README.md}). Two formulas for the $35$-point fivefold case, an
  angle form and a layer form of one SAT encoder (not independent encodings: audit \dirn{AUDIT3}, finding A4-3), are
  recorded as unsatisfiable; the LRAT proofs of both were accepted by cake\_lpr. The passage from
  these verdicts to $\mathrm{Fivefold35}$ uses the soundness of the encoder (recorded as a written proof, reviewed in
  the audit, \dirn{AUDIT3}, A4), a normal-form lemma for the labelling (for the layer form, the layer lemma), and the
  centre-perturbation lemma, which makes the centre generic (\dirn{C2ENC/perturbation.md}, recorded there as a
  complete written proof, not Lean-checked, refereed within the project on the day it was written:
  \dirn{C2ENC/referee/NOTES.md}, Section~4). \dirn{C2ENC/class_table.md} summarizes this chain for the constant of
  the same text in \dirn{SYMLEAN} (see below).
\item $\mathrm{Fourfold32}$: \dirn{C4CUT} (\dirn{C4CUT/README.md}). The layered formula for the $32$-point fourfold
  case with six realizability cuts, each justified by a Grassmann--Pl\"ucker certificate checked in exact rational
  arithmetic, is recorded as unsatisfiable; its LRAT proof was accepted by cake\_lpr. The final run does not record
  which checker builds it used (\dirn{AUDIT3}, A4-5, which concludes that the claim is carried by the cake\_lpr verdict);
  the replay \dirn{REPLAY}, with hash-checked files and checker binaries, accepted the proof again with cake\_lpr. The
  passage to $\mathrm{Fourfold32}$ uses the soundness of the layered encoder (recorded in \dirn{C4CUT/README.md} as a
  written proof, refereed within the project: the review of 1~November 2026, \dirn{C4CUT/review/README.md}), the
  layer lemma, and the sector normal form: each of the other orbits has exactly one
  point in the open quadrant between the rays from the centre through two consecutive points of the outermost orbit.
  The audit found the argument for this not written down (\dirn{AUDIT3}, A4-1); \dirn{C4CUT/README.md}, Section~1,
  now gives it in one sentence: a point of another orbit on a bounding ray would be collinear with two antipodal
  points of the outermost orbit.
\end{itemize}
The reviews, audits and refereeing named in this list were all made within the project (internal reviews); none of
them is an external peer review. Apart from the layer lemma and the ray fact for even orders behind the sector
normal form, none of these written arguments is checked in Lean. Both have Lean-checked forms. For the layer lemma
it is Theorem~\ref{ll:thm:layers}.
For the ray fact it is Lemma~\ref{oc:lem:same-ray} together with Theorem~\ref{oc:thm:phase-orbit}
(\dirn{LEANEXT/orbit/OrbitConsistency.lean}): for even $m>0$, no two points of a finite set in general position that
is invariant under $\mathrm{rotAbout}(c,2\pi/m)$ lie on one open ray from $c$; and if $x$ and $y$ are points of such a
set, both different from $c$, and $y$ is not in the orbit of $x$ under that rotation, then the directions of $x$ and
$y$ at $c$ differ by no integer multiple of $2\pi/m$. The
remaining step, that each of the other orbits has exactly one point in the open quadrant, is a counting argument and
is not in Lean. The rotation invariance of orientations and their sign rule under permutations, on which the
encoders' single variable per orbit of triples rests, are Lean-checked too (Lemmas~\ref{oc:lem:rot-inv}
and~\ref{oc:lem:perm}). Nothing in Lean connects any of these Lean-checked forms to the formulas, and no Lean theorem
links a formula or a certificate to these propositions. Propositions with the same text occur in
\dirn{SYMLEAN/SymmetricReduction.lean:477} and \dirn{SYMLEAN/SymmetricReduction.lean:496}
(Definition~\ref{sr:def:classes}); they are different Lean constants (namespace \lean{ErdosSzekeres.SymReduction},
built on that file's own copies of \lean{Invariant} and \lean{rotAbout}), and no Lean theorem identifies the two
versions. The docstring of \lean{Fourfold32}, copied from \dirn{SYMLEAN}, says that it is used only by a variant
theorem of \dirn{SYMLEAN}; in \dirn{LEANEXT/es7corr/ES7Corollaries.lean} it is the hypothesis of
Theorems~\ref{ec:thm:fourfold-erase} and~\ref{ec:thm:fourfold}.
\end{remark}

We use two Lean-checked results whose written proofs are in other sections, in the following form.
\begin{itemize}
\item[(R)] Theorem~\ref{ll:thm:reduction} (\lean{convexNGon_of_layers}, \dirn{LAYERLEAN/LayerLemma.lean:1066}). Let
  $m,K,t\in\N$ with $0<m$ and $K\le2m$, and assume: \emph{for every finite set $T$ and every point $c'$, if $T$ is in
  general position, $c'\notin T$, $T$ is invariant under $\mathrm{rotAbout}(c',2\pi/m)$ and $\card T=m\cdot t$, then
  $\mathrm{HasConvexNGon}(K,T)$}. Then for every finite set $S$ and every point $c$ such that $S$ is in general
  position, $S$ is invariant under $\mathrm{rotAbout}(c,2\pi/m)$ and $m\cdot(t\mathbin{\dot{-}}1)<\card{S\setminus\{c\}}$,
  we have $\mathrm{HasConvexNGon}(K,S)$. Here $t\mathbin{\dot{-}}1=\max(t-1,0)$ is Lean's truncated subtraction
  on $\N$ (the notation of the layer-lemma section).
\item[(C)] Lemma~\ref{lg:lem:centre} (\lean{centre_not_mem_of_even}, \dirn{LAYERLEAN/LayerLemma.lean:582}). Let
  $m\in\N$, let $S$ be a finite set of points and $c$ a point. If $0<m$, $m$ is even, $S$ is invariant under
  $\mathrm{rotAbout}(c,2\pi/m)$ and in general position, and some $p\in S$ satisfies $p\ne c$, then $c\notin S$. (In
  this section, for a map $f$ and $n\in\N$, $f^n$ is the $n$-fold composite of $f$, written $f^{[n]}$ in the
  foundations section; so $f^0=\mathrm{id}$. The power $\mathrm{rotAbout}(c,2\pi/m)^{m/2}$ is the half-turn
  $p\mapsto 2c-p$, and $p$, $c$, $2c-p$ would be three distinct collinear points of $S$.)
\end{itemize}

In each of the following five theorems the case class is a hypothesis (Definition~\ref{ec:def:cases} and
Remark~\ref{ec:rem:cases}).

\begin{theorem}[fivefold, layer form; \leanref{ErdosSzekeres.LayerLemma.fivefold_heptagon_of_erase}{LEANEXT/es7corr/ES7Corollaries.lean:1170}]\label{ec:thm:fivefold-erase}
Assume $\mathrm{Fivefold35}$. Let $S$ be a finite set of points and $c$ a point such that $S$ is in general position,
$S$ is invariant under $\mathrm{rotAbout}(c,2\pi/5)$, and $30<\card{S\setminus\{c\}}$. Then
$\mathrm{HasConvexNGon}(7,S)$.
\end{theorem}

\begin{proof}
We follow the Lean proof. Apply (R) with $m=5$, $K=7$, $t=7$; here $0<5$ and $K=7\le10=2m$. The assumption of (R)
holds: let $T$ be finite and $c'$ a point such that $T$ is in general position, $c'\notin T$, $T$ is invariant under
$\mathrm{rotAbout}(c',2\pi/5)$ and $\card T=5\cdot7=35$. These are the four hypotheses of $\mathrm{Fivefold35}$ for
$T$ and $c'$ (listed there in a different order), so $\mathrm{Fivefold35}$ gives $\mathrm{HasConvexNGon}(7,T)$. The
size condition of (R) is $5\cdot(7\mathbin{\dot{-}}1)=5\cdot6=30<\card{S\setminus\{c\}}$, which is our hypothesis. So (R) gives
$\mathrm{HasConvexNGon}(7,S)$.
\end{proof}

\begin{theorem}[fivefold; \leanref{ErdosSzekeres.LayerLemma.fivefold_heptagon}{LEANEXT/es7corr/ES7Corollaries.lean:1183}]\label{ec:thm:fivefold}
Assume $\mathrm{Fivefold35}$. Let $S$ be a finite set of points and $c$ a point such that $S$ is in general position,
$S$ is invariant under $\mathrm{rotAbout}(c,2\pi/5)$, and $32\le\card S$. Then $\mathrm{HasConvexNGon}(7,S)$.
\end{theorem}

\begin{proof}
As in Lean: removing one element lowers the cardinality by at most one, so
$\card{S\setminus\{c\}}\ge\card S-1\ge31>30$, and Theorem~\ref{ec:thm:fivefold-erase} applies.
\end{proof}

\begin{theorem}[centre-free; \leanref{ErdosSzekeres.LayerLemma.fivefold_heptagon_of_centre_not_mem}{LEANEXT/es7corr/ES7Corollaries.lean:1193}]\label{ec:thm:fivefold-nc}
Assume $\mathrm{Fivefold35}$. Let $S$ be a finite set of points and $c$ a point such that $S$ is in general position,
$S$ is invariant under $\mathrm{rotAbout}(c,2\pi/5)$, $c\notin S$, and $31\le\card S$. Then
$\mathrm{HasConvexNGon}(7,S)$.
\end{theorem}

\begin{proof}
As in Lean: since $c\notin S$, $S\setminus\{c\}=S$, so $\card{S\setminus\{c\}}=\card S\ge31>30$, and
Theorem~\ref{ec:thm:fivefold-erase} applies.
\end{proof}

\begin{theorem}[fourfold, layer form; \leanref{ErdosSzekeres.LayerLemma.fourfold_heptagon_of_erase}{LEANEXT/es7corr/ES7Corollaries.lean:1203}]\label{ec:thm:fourfold-erase}
Assume $\mathrm{Fourfold32}$. Let $S$ be a finite set of points and $c$ a point such that $S$ is in general position,
$S$ is invariant under $\mathrm{rotAbout}(c,2\pi/4)$, and $28<\card{S\setminus\{c\}}$. Then
$\mathrm{HasConvexNGon}(7,S)$.
\end{theorem}

\begin{proof}
We follow the Lean proof. Apply (R) with $m=4$, $K=7$, $t=8$; here $0<4$ and $K=7\le8=2m$. The assumption of (R)
holds by $\mathrm{Fourfold32}$, exactly as in the proof of Theorem~\ref{ec:thm:fivefold-erase}, since $4\cdot8=32$.
The size condition of (R) is $4\cdot(8\mathbin{\dot{-}}1)=4\cdot7=28<\card{S\setminus\{c\}}$, our hypothesis.
\end{proof}

\begin{theorem}[fourfold; \leanref{ErdosSzekeres.LayerLemma.fourfold_heptagon}{LEANEXT/es7corr/ES7Corollaries.lean:1216}]\label{ec:thm:fourfold}
Assume $\mathrm{Fourfold32}$. Let $S$ be a finite set of points and $c$ a point such that $S$ is in general position,
$S$ is invariant under $\mathrm{rotAbout}(c,2\pi/4)$, and $29\le\card S$. Then $\mathrm{HasConvexNGon}(7,S)$.
\end{theorem}

\begin{proof}
We follow the Lean proof. Since $\card S\ge29>1$, $S$ is not contained in $\{c\}$, so some $p\in S$ satisfies $p\ne
c$. Apply (C) with $m=4$: $0<4$, $4$ is even, $S$ is in general position, and $S$ is invariant under
$\mathrm{rotAbout}(c,2\pi/4)$, which is the rotation $\mathrm{rotAbout}(c,2\pi/m)$ of (C) for $m=4$ (the cast angle
$2\pi/m$ equals $2\pi/4$ for $m=4$; see the Conventions). So $c\notin S$. Hence $S\setminus\{c\}=S$ and
$\card{S\setminus\{c\}}=\card S\ge29>28$, and
Theorem~\ref{ec:thm:fourfold-erase} applies.
\end{proof}

\begin{remark}[what the theorems give]\label{ec:rem:es7-sharp}
Under $\mathrm{Fivefold35}$, Theorem~\ref{ec:thm:fivefold} says that a finite set in general position that is
invariant under $\mathrm{rotAbout}(c,2\pi/5)$ for some $c$ and has no $7$ points in convex position has at most $31$
points; under $\mathrm{Fourfold32}$, Theorem~\ref{ec:thm:fourfold} gives at most $28$ points for
$\mathrm{rotAbout}(c,2\pi/4)$. These are the upper bounds $N(C_5)\le31$ and $N(C_4)\le28$ of the symmetric $ES(7)$
note (\dirn{WRITEUP/note/symmetric_es7.tex}, the proposition labelled \texttt{prop:C}), where they are derived from
the $35$-point fivefold and the $32$-point fourfold case by orbit counting (the lemma labelled \texttt{lem:sizes}). A set
with symmetry $C_s$ in the sense of the note is invariant under a rotation $\rho$ of order $s$ about some point $c$,
that is $\rho=\mathrm{rotAbout}(c,2\pi k/s)$ for an integer $k$ with $\gcd(k,s)=1$. Choose $j\in\N$ with
$jk\equiv1\pmod s$. By the complex form of Definition~\ref{fd:def:rot},
$\rho^j=\mathrm{rotAbout}(c,2\pi jk/s)=\mathrm{rotAbout}(c,2\pi/s)$, since the two angles differ by a multiple of
$2\pi$; and a set mapped into itself by $\rho$ is mapped into itself by $\rho^j$. So the two theorems apply to every
finite set in general position with symmetry $C_5$, respectively $C_4$. The matching lower-bound examples \dirn{SYMMAX/c5_31points.txt} and
\dirn{SYMMAX/c4_28points.txt} rest on exact computations outside Lean, which were not repeated for this document. The
logs \dirn{LEANEXT/es7corr/mutant31.log} and \dirn{LEANEXT/es7corr/mutant28.log} record that the Lean proofs fail at their
final arithmetic step when the thresholds $32$ and $29$ are lowered to $31$ and $28$; this audits the proof scripts
only.
\end{remark}

\begin{remark}[a route without the layer lemma]\label{ec:rem:orbits}
Theorems~\ref{ec:thm:fivefold-erase} and~\ref{ec:thm:fourfold-erase} also follow by orbit counting, the argument
of the symmetric $ES(7)$ note. Let $\rho=\mathrm{rotAbout}(c,2\pi/m)$ with $m\in\{4,5\}$. The map $\rho$ fixes $c$ and
is injective, so it maps $S'=S\setminus\{c\}$ into itself, hence onto itself ($S'$ is finite). Every $y\ne c$ has
minimal period $m$ under $\rho$ (Lemma~\ref{ll:lem:minperiod}), so the $\rho$-orbits $\{y,\rho y,\dots,\rho^{m-1}y\}$
of the points $y\in S'$ have exactly $m$ elements and partition $S'$. If $m(t-1)<\card{S'}$, with $t=7$ for $m=5$ and
$t=8$ for $m=4$, then $S'$ consists of at least $t$ orbits; the union $T$ of $t$ of them has $mt$ points ($35$,
respectively $32$), satisfies $c\notin T$, is in general position (as a subset of $S$) and is invariant under
$\rho$. The case class applies to $T$, and $7$ points of $T$ in convex position lie in $S$. The Lean proofs use
Theorem~\ref{ll:thm:reduction} instead, which takes for $T$ the $t$ outermost convex layers of $S'$.
\end{remark}

%---------------------------------------------------------------------------------------------------------------------
\subsection{Rotations of order at least eight}\label{ec:ssec:order8}

The source is \dirn{ES8SYM/lean/C6check.lean} (namespace \lean{RefuteC6}; it imports Mathlib only, not the project
library). It is the Lean check of a complex form of part (a) of Lemma~2 of \dirn{ES8SYM/proofs/layers.md} (the planar
form is Corollary~\ref{ec:cor:planar}, a written step) and of the arithmetic of part (b) of that lemma.

\emph{Conventions.} In this subsection points are complex numbers. $\C$ is a real vector space of dimension $2$,
$\abs z$ is the modulus (the Euclidean norm), and for $r\in\R$ we write $B_r=\{z\in\C:\abs z\le r\}$ (Lean
\texttt{closedBall 0 r}); the circle $\{z\in\C:\abs z=r\}$ is Lean's \texttt{sphere 0 r}, and both sets are empty if
$r<0$. $\conv$ is the convex hull over $\R$, the smallest convex set containing the given set. The following notions are
Mathlib's.
\begin{itemize}
\item An \emph{extreme point} of $A\subseteq\C$ (Mathlib \lean{Set.extremePoints}) is a point $x\in A$ such that
  whenever $x_1,x_2\in A$, $a,b\in\R$, $a>0$, $b>0$, $a+b=1$ and $x=ax_1+bx_2$, we have $x_1=x$.
\item A family $f:\iota\to\C$ is \emph{convex independent} over $\R$ (Mathlib \lean{ConvexIndependent}) if for every
  set $s\subseteq\iota$ and every $x\in\iota$, $f(x)\in\conv f(s)$ implies $x\in s$.
\item $\zeta\in\C$ is a \emph{primitive $m$-th root of unity} (Mathlib \lean{IsPrimitiveRoot}) if $\zeta^m=1$ and,
  for every $l\in\N$, $\zeta^l=1$ implies $m\mid l$.
\end{itemize}
$\mathrm{Fin}\,8=\{0,1,\dots,7\}$. A finite set of points is in convex position (Definition~\ref{fd:def:cp}) if the
inclusion family of the set is convex independent; Corollary~\ref{ec:cor:planar} makes the passage between injective
families in $\C$ and finite sets of points precise.

\begin{lemma}[points of a circle are extreme points of the disc; \leanref{RefuteC6.sphere_subset_extremePoints}{ES8SYM/lean/C6check.lean:24}]\label{ec:lem:sphere}
For every $r\in\R$, every $x\in\C$ with $\abs x=r$ is an extreme point of $B_r$.
\end{lemma}

\begin{proof}
We follow the Lean proof, which uses the strict convexity of the norm of $\C$; we make it explicit. For
$x_1,x_2\in\C$ and real $a,b$ with $a+b=1$,
\begin{equation*}
\abs{ax_1+bx_2}^2=a\abs{x_1}^2+b\abs{x_2}^2-ab\,\abs{x_1-x_2}^2 .
\end{equation*}
Indeed, with $\langle z,w\rangle=\operatorname{Re}(z\bar w)$ the left side is
$a^2\abs{x_1}^2+b^2\abs{x_2}^2+2ab\langle x_1,x_2\rangle$, the right side is
$(a-ab)\abs{x_1}^2+(b-ab)\abs{x_2}^2+2ab\langle x_1,x_2\rangle$, and $a-ab=a(1-b)=a^2$, $b-ab=b^2$.

Let $\abs x=r$. Then $x\in B_r$. Let $x_1,x_2\in B_r$ and $a,b>0$ with $a+b=1$ and $x=ax_1+bx_2$. If $x_1=x_2$, then
$x=(a+b)x_1=x_1$. Suppose $x_1\ne x_2$. Since $0\le\abs{x_1}\le r$ we have $r\ge0$ and $\abs{x_1}^2\le r^2$, and
likewise $\abs{x_2}^2\le r^2$; also $ab\abs{x_1-x_2}^2>0$. The identity gives
$r^2=\abs x^2<ar^2+br^2=r^2$, a contradiction. Hence $x_1=x$.
\end{proof}

\begin{lemma}[injective families on a circle are convex independent; \leanref{RefuteC6.convexIndependent_of_norm_eq}{ES8SYM/lean/C6check.lean:43}]\label{ec:lem:ci-circle}
Let $\iota$ be any type, $f:\iota\to\C$ injective, and $r\in\R$ with $\abs{f(i)}=r$ for every $i\in\iota$. Then $f$ is
convex independent over $\R$.
\end{lemma}

\begin{proof}
Step 1. \emph{If $x$ is an extreme point of a convex set $A\subseteq\C$, then $A\setminus\{x\}$ is convex.} Let
$y,z\in A\setminus\{x\}$, $a,b\ge0$ real with $a+b=1$, and $w=ay+bz$; then $w\in A$. If $w=x$: for $b=0$ we get
$w=y\ne x$, for $a=0$ we get $w=z\ne x$, and for $a,b>0$ extremality gives $y=x$; each case is a contradiction. So
$w\in A\setminus\{x\}$.

Step 2. The disc $B_r$ is convex: for $y,z\in B_r$ and $a,b\ge0$ with $a+b=1$, $\abs{ay+bz}\le a\abs y+b\abs
z\le r$. Let $s\subseteq\iota$ and $x\in\iota$ with $f(x)\in\conv f(s)$, and suppose $x\notin s$. For $j\in s$ we have
$\abs{f(j)}=r$, so $f(j)\in B_r$, and $j\ne x$, so $f(j)\ne f(x)$ by injectivity. Thus $f(s)\subseteq
B_r\setminus\{f(x)\}$. Since $\abs{f(x)}=r$, $f(x)$ is an extreme point of $B_r$ (Lemma~\ref{ec:lem:sphere}), so
$B_r\setminus\{f(x)\}$ is convex by Step~1, and therefore contains $\conv f(s)$. This contradicts $f(x)\in\conv f(s)$.
Hence $x\in s$.

(The Lean proof uses Mathlib's reformulation of convex independence, ``$f(i)\notin\conv f(s\setminus\{i\})$ for all
$i,s$'', with the same two steps.)
\end{proof}

\begin{definition}[the generator; \leanref{RefuteC6.zeta}{ES8SYM/lean/C6check.lean:60}]\label{ec:def:zeta}
For $m\in\N$, $\zeta_m=\exp(2\pi i/m)\in\C$, with $m$ cast to $\C$. (For $m=0$, Lean's division gives $2\pi i/0=0$ and
$\zeta_0=1$; this case is not used.) Rotation by $2\pi/m$ about $0$ is $z\mapsto\zeta_mz$.
\end{definition}

\begin{lemma}[primitive root; \leanref{RefuteC6.zeta_prim}{ES8SYM/lean/C6check.lean:62}]\label{ec:lem:zeta}
For $m\in\N$ with $m\ne0$, $\zeta_m$ is a primitive $m$-th root of unity.
\end{lemma}

\begin{proof}
This is Mathlib's \lean{Complex.isPrimitiveRoot_exp}, a standard fact. Proof: $\exp(w)=1$ if and only if
$w\in2\pi i\Z$ (periodicity of the exponential function), and $\zeta_m^l=\exp(2\pi il/m)$ for $l\in\N$. Hence
$\zeta_m^l=1$ if and only if $l/m\in\Z$, that is $m\mid l$; in particular $\zeta_m^m=1$.
\end{proof}

\begin{fact}[properties of primitive roots; written step]\label{ec:fact:prim}
Let $m\ge1$ and let $\zeta\in\C$ be a primitive $m$-th root of unity. Then
\textup{(i)} $\zeta\ne0$; \textup{(ii)} $\abs\zeta=1$; \textup{(iii)} if $i,j<m$ and $\zeta^i=\zeta^j$, then $i=j$;
\textup{(iv)} $\zeta^j\ne1$ for $0<j<m$.
\end{fact}

\begin{proof}
(i) $\zeta^m=1\ne0$ and $m\ge1$. (ii) $\abs\zeta^m=\abs{\zeta^m}=1$ with $\abs\zeta\ge0$, and $x\mapsto x^m$ is
strictly increasing on $[0,\infty)$, so $\abs\zeta=1$. (iii) Say $i\le j$. Then
$\zeta^i\zeta^{j-i}=\zeta^j=\zeta^i$, and $\zeta^i\ne0$ by (i), so $\zeta^{j-i}=1$, hence $m\mid j-i$; as $0\le
j-i<m$, $j=i$. (iv) $\zeta^j=1$ would give $m\mid j$, impossible for $0<j<m$. (In Mathlib these are
\lean{IsPrimitiveRoot.ne_zero}, \lean{IsPrimitiveRoot.norm'_eq_one}, \lean{IsPrimitiveRoot.pow_inj} and
\lean{IsPrimitiveRoot.pow_ne_one_of_pos_of_lt}.)
\end{proof}

\begin{lemma}[eight points of an orbit; \leanref{RefuteC6.orbit_injective}{ES8SYM/lean/C6check.lean:66}, \leanref{RefuteC6.orbit_convexIndependent}{ES8SYM/lean/C6check.lean:75}]\label{ec:lem:orbit8}
Let $m\in\N$ with $8\le m$ and $p\in\C$ with $p\ne0$, and let $F:\mathrm{Fin}\,8\to\C$, $F(k)=p\,\zeta_m^k$. Then
\textup{(a)} $F$ is injective, and \textup{(b)} $F$ is convex independent over $\R$.
\end{lemma}

\begin{proof}
Since $m\ne0$, $\zeta_m$ is a primitive $m$-th root of unity (Lemma~\ref{ec:lem:zeta}).
(a) If $p\zeta_m^j=p\zeta_m^k$, then $\zeta_m^j=\zeta_m^k$ because $p\ne0$, and $j,k<8\le m$, so $j=k$ by
Fact~\ref{ec:fact:prim}(iii).
(b) $\abs{F(k)}=\abs p\,\abs{\zeta_m}^k=\abs p$ by Fact~\ref{ec:fact:prim}(ii). Apply Lemma~\ref{ec:lem:ci-circle}
with $r=\abs p$ and (a).
\end{proof}

\begin{theorem}[order at least $8$; \leanref{RefuteC6.part_a}{ES8SYM/lean/C6check.lean:85}]\label{ec:thm:big-order}
Let $m\in\N$ with $8\le m$, and let $S\subseteq\C$ be any set \textup{(}not necessarily finite\textup{)} such that
$\exp(2\pi i/m)\,z\in S$ for every $z\in S$. Assume that for every injective $f:\mathrm{Fin}\,8\to\C$ with $f(k)\in S$
for all $k$, $f$ is not convex independent over $\R$. Then $S\subseteq\{0\}$.
\end{theorem}

\begin{proof}
We follow the Lean proof. Let $p\in S$ and suppose $p\ne0$. By induction, $p\zeta_m^k\in S$ for every $k\in\N$:
$p\zeta_m^0=p\in S$, and if $p\zeta_m^k\in S$, then $\exp(2\pi i/m)\,(p\zeta_m^k)=p\zeta_m^{k+1}\in S$ by the hypothesis
(note $\exp(2\pi i/m)=\zeta_m$ by definition). Hence the family $F$ of Lemma~\ref{ec:lem:orbit8} takes values in $S$; it
is injective and convex independent by that lemma, contradicting the assumption. So $p=0$.
\end{proof}

\begin{remark}\label{ec:rem:big-order}
The theorem uses no finiteness and no general position, and only the forward invariance $\zeta_mS\subseteq S$. Its
conclusion says $S=\emptyset$ or $S=\{0\}$; both sets satisfy the hypotheses, since an injective family of $8$ values
needs $8$ distinct points. The planar statement used in the project (\dirn{ES8SYM/proofs/layers.md}, Lemma~2(a),
stated there for sets in general position; the row ``order $\ge8$'', labelled Lean-checked, of the $k=8$ tables in
\dirn{ES8SYM/README.md} and \dirn{WRITEUP/note/symmetric_extremal.tex}) follows from the following corollary, which
does not assume general position. Lemma~2(a) is stated for the rotation by $2\pi/m$, and the corollary applies to it
directly. The tables concern any rotation $\rho$ of order $m$ (\dirn{WRITEUP/note/symmetric_extremal.tex}, line~338),
that is $\rho=\mathrm{rotAbout}(c,2\pi k/m)$ with $\gcd(k,m)=1$; as in Remark~\ref{ec:rem:es7-sharp}, the power
$\rho^j$ with $j\in\N$ and $jk\equiv1\pmod m$ is $\mathrm{rotAbout}(c,2\pi/m)$ and maps every $\rho$-invariant set
into itself, so the corollary applies to $\rho^j$. Lean checks only the complex form, Theorem~\ref{ec:thm:big-order};
the passage to the plane is not formalized in Lean, and a written proof follows.
\end{remark}

\begin{corollary}[planar form of Theorem~\ref{ec:thm:big-order}; written step]\label{ec:cor:planar}
Let $m\in\N$ with $m\ge8$, let $c$ be a point, and let $S$ be a finite set of points that is invariant under
$\mathrm{rotAbout}(c,2\pi/m)$ and satisfies $\neg\,\mathrm{HasConvexNGon}(8,S)$. Then $S\subseteq\{c\}$.
\end{corollary}

\begin{proof}
Let $\varphi:\R^2\to\C$, $\varphi(p)=(p_0-c_0)+i(p_1-c_1)$. It is a bijection and an affine map of real vector spaces,
with inverse $\psi(z)=(c_0+\operatorname{Re}z,\;c_1+\operatorname{Im}z)$. By the complex form of
Definition~\ref{fd:def:rot}, $\varphi(\mathrm{rotAbout}(c,\theta)\,p)=e^{i\theta}\varphi(p)$ for all $\theta\in\R$ and
all points $p$. We also use the standard fact that an affine map $\alpha$ between real vector spaces satisfies
$\alpha(\conv Y)=\conv\alpha(Y)$, because $\conv Y$ is the set of finite convex combinations of points of $Y$ and
$\alpha$ preserves convex combinations. Let $\tilde S=\varphi(S)\subseteq\C$.

\emph{Invariance.} Let $z=\varphi(p)$ with $p\in S$. With $\theta=2\pi/m$ we have $e^{i\theta}=\exp(2\pi i/m)$, so
$\exp(2\pi i/m)\,z=\varphi(\mathrm{rotAbout}(c,2\pi/m)\,p)\in\tilde S$, because $\mathrm{rotAbout}(c,2\pi/m)\,p\in S$.

\emph{No convex independent $8$-family.} Let $f:\mathrm{Fin}\,8\to\C$ be injective with values in $\tilde S$, and
suppose $f$ is convex independent. Put $g=\psi\circ f$; it is injective with values in $S$, so $T=g(\mathrm{Fin}\,8)$
is an $8$-element subset of $S$. We show that $T$ is in convex position, i.e.\ (Definition~\ref{fd:def:cp}) that for
every $U\subseteq T$ and every $x\in T$, $x\in\conv U$ implies $x\in U$. Write $x=g(k)$ and
$s=\{j\in\mathrm{Fin}\,8: g(j)\in U\}$; then $U=g(s)$, and $\varphi\circ g=f$. From $x\in\conv U$ we get
$f(k)=\varphi(x)\in\varphi(\conv U)=\conv\varphi(U)=\conv f(s)$, so $k\in s$ by convex independence of $f$, that is
$x=g(k)\in U$. Hence $\mathrm{HasConvexNGon}(8,S)$, a contradiction.

So $\tilde S$ satisfies the hypotheses of Theorem~\ref{ec:thm:big-order}, and $\tilde S\subseteq\{0\}$. Since
$\varphi(p)=0$ only for $p=c$, $S\subseteq\{c\}$.
\end{proof}

\begin{proposition}[threshold arithmetic; \leanref{RefuteC6.thr_m4}{ES8SYM/lean/C6check.lean:108}, \leanref{RefuteC6.thr_m5}{ES8SYM/lean/C6check.lean:109}, \leanref{RefuteC6.thr_m6}{ES8SYM/lean/C6check.lean:110}, \leanref{RefuteC6.thr_m7}{ES8SYM/lean/C6check.lean:111}, \leanref{RefuteC6.m4_with_centre_counter}{ES8SYM/lean/C6check.lean:115}]\label{ec:prop:thr}
In $\N$:
\begin{enumerate}[label=(\roman*)]
\item \textup{(\lean{thr_m4})} for every $t$: $65\le4t\iff17\le t$;
\item \textup{(\lean{thr_m5})} for all $t,e$ with $e\le1$: $65\le5t+e\iff13\le t$;
\item \textup{(\lean{thr_m6})} for every $t$: $65\le6t\iff11\le t$;
\item \textup{(\lean{thr_m7})} for all $t,e$ with $e\le1$: $65\le7t+e\iff10\le t$;
\item \textup{(\lean{m4_with_centre_counter})} $65\le4\cdot16+1$ and $\neg(17\le16)$.
\end{enumerate}
\textup{(}In Lean the variable $e$ is called \texttt{c}.\textup{)}
\end{proposition}

\begin{proof}
(i) If $t\ge17$ then $4t\ge68$; if $t\le16$ then $4t\le64$. (ii) If $t\ge13$ then $5t+e\ge65$; if $t\le12$ then
$5t+e\le60+1=61$. (iii) If $t\ge11$ then $6t\ge66$; if $t\le10$ then $6t\le60$. (iv) If $t\ge10$ then $7t+e\ge70$; if
$t\le9$ then $7t+e\le63+1=64$. (v) $4\cdot16+1=65$, and $17>16$. (Lean proves each line with \lean{omega}, a decision
procedure for linear arithmetic over $\N$ and $\Z$ whose proof term is checked by the kernel.)
\end{proof}

\begin{remark}[the geometric part of Lemma~2(b), not in this file]\label{ec:rem:thr}
In \dirn{ES8SYM/proofs/layers.md}, Lemma~2(b), the proposition is applied to a finite set $S$ in general position,
invariant under $\mathrm{rotAbout}(c,2\pi/m)$ with $m\in\{4,5,6,7\}$ and without $8$ points in convex position. The
set $S\setminus\{c\}$ is again invariant (the rotation is injective and fixes $c$), and the layer results applied to it
(Theorems~\ref{ll:thm:layers} and~\ref{ll:thm:card-peel}, with $K=8\le2m$) give $\card{S\setminus\{c\}}=mt$, where
$t$ is the number of convex layers of $S\setminus\{c\}$. If moreover $S\not\subseteq\{c\}$ (for instance if
$\card S\ge2$), then $c\notin S$ for even $m$ (Lemma~\ref{lg:lem:centre}). So $\card S=mt+e$ with $e=[c\in S]\le1$,
and $e=0$ for $m=4,6$ when $\card S\ge2$. Then Proposition~\ref{ec:prop:thr} turns $\card S\ge65$ into
$t\ge17,13,11,10$. This geometric input is not part of \dirn{ES8SYM/lean/C6check.lean}, so Lemma~2(b) as a whole is
not Lean-checked there. Item (v) records the arithmetic behind the
remark of \dirn{ES8SYM/proofs/layers.md} (lines 41--42) that general position is needed: with the centre allowed for
$m=4$, $16$ orbits and the centre would already give $65$ points.
\end{remark}

%---------------------------------------------------------------------------------------------------------------------
\subsection{Orientations in inverse-radius polar form}\label{ec:ssec:polar}

The source is \dirn{ES8SYM/lean/PolarOrient.lean} (namespace \lean{C7}; Mathlib only). It is the Lean check of
Lemma~1 of \dirn{ES8SYM/proofs/layers.md} (lines 15--23) and of part of the consequence stated after it (lines 25--29;
see Remark~\ref{ec:rem:polar}(iii)). For $x\in\R$, $\operatorname{sign}x\in\{-1,0,1\}$ is
Mathlib's \lean{SignType.sign}: $1$ if $x>0$, $0$ if $x=0$, $-1$ if $x<0$. We use
$\operatorname{sign}(xy)=\operatorname{sign}(x)\operatorname{sign}(y)$ and $\operatorname{sign}(x)=1$ for $x>0$; hence
$\operatorname{sign}(xP)=\operatorname{sign}(x)$ whenever $P>0$. For $p\in\R\times\R$ we write $p=(p_1,p_2)$ (Lean
\lean{p.1}, \lean{p.2}).

\begin{definition}[cross product, inverse-radius points, layered coefficients; \leanref{C7.cross}{ES8SYM/lean/PolarOrient.lean:45}, \leanref{C7.pt}{ES8SYM/lean/PolarOrient.lean:48}, \leanref{C7.coef}{ES8SYM/lean/PolarOrient.lean:63}, \leanref{C7.phase}{ES8SYM/lean/PolarOrient.lean:94}, \leanref{C7.ngonPt}{ES8SYM/lean/PolarOrient.lean:97}]\label{ec:def:polar}
The following objects are used in this subsection.
\begin{enumerate}[label=(\roman*)]
\item For $p,q,r\in\R\times\R$:
  $\mathrm{cross}(p,q,r)=(q_1-p_1)(r_2-p_2)-(q_2-p_2)(r_1-p_1)$.
\item For $\varphi,u\in\R$: $\mathrm{pt}(\varphi,u)=(\cos\varphi/u,\ \sin\varphi/u)\in\R\times\R$, the point with
  direction $\varphi$ and inverse radius $u$ (Lean's division gives $x/0=0$; only $u>0$ is used).
\item For a type $\iota$, $t\in\N$, $L:\iota\to\mathrm{Fin}\,t$, $\varphi:\iota\to\R$ and $i,j,k\in\iota$, the vector
  $\mathrm{coef}_{L,\varphi}(i,j,k)\in\R^{\mathrm{Fin}\,t}$ is
  \begin{equation*}
  \mathrm{coef}_{L,\varphi}(i,j,k)=\sin(\varphi_k-\varphi_j)\,e_{L(i)}+\sin(\varphi_i-\varphi_k)\,e_{L(j)}
  +\sin(\varphi_j-\varphi_i)\,e_{L(k)},
  \end{equation*}
  where $\varphi_i=\varphi(i)$ and $e_o$ is the $o$-th unit vector (Lean \lean{Pi.single}); when two of
  $L(i),L(j),L(k)$ coincide, the coefficients add.
\item For $m,t\in\N$, $\theta:\mathrm{Fin}\,t\to\R$ and $(o,k)\in\mathrm{Fin}\,t\times\mathrm{Fin}\,m$:
  $\mathrm{phase}_{m,\theta}(o,k)=\theta_o+2\pi k/m$, computed in $\R$ as $(2\pi k)/m$ with $k$ and $m$ cast to $\R$.
\item For moreover $r:\mathrm{Fin}\,t\to\R$: $\mathrm{ngonPt}_{m,\theta,r}(o,k)=\bigl(r_o\cos\mathrm{phase}_{m,\theta}(o,k),\
  r_o\sin\mathrm{phase}_{m,\theta}(o,k)\bigr)$, vertex $k$ of the regular $m$-gon with centre $0$, radius $r_o$ and
  phase $\theta_o$.
\end{enumerate}
\end{definition}

For the rest of this subsection, for real $\varphi_1,\varphi_2,\varphi_3,u_1,u_2,u_3$ let
\begin{equation*}
D=\det\begin{pmatrix}\cos\varphi_1/u_1&\sin\varphi_1/u_1&1\\ \cos\varphi_2/u_2&\sin\varphi_2/u_2&1\\
\cos\varphi_3/u_3&\sin\varphi_3/u_3&1\end{pmatrix},\qquad
\Sigma=u_1\sin(\varphi_3-\varphi_2)+u_2\sin(\varphi_1-\varphi_3)+u_3\sin(\varphi_2-\varphi_1).
\end{equation*}

\begin{lemma}[the value identity; \leanref{C7.det_mul_eq}{ES8SYM/lean/PolarOrient.lean:24}]\label{ec:lem:det-mul}
For all real $\varphi_1,\varphi_2,\varphi_3,u_1,u_2,u_3$ with $u_1\ne0$, $u_2\ne0$, $u_3\ne0$: $D\cdot(u_1u_2u_3)=\Sigma$.
\end{lemma}

\begin{proof}
Write $x_k=\cos\varphi_k/u_k$ and $y_k=\sin\varphi_k/u_k$. Expanding the determinant along its third column,
\begin{equation*}
D=(x_2y_3-x_3y_2)+(x_3y_1-x_1y_3)+(x_1y_2-x_2y_1).
\end{equation*}
For indices $a\ne b$, the subtraction formula $\sin(\beta-\alpha)=\sin\beta\cos\alpha-\cos\beta\sin\alpha$ gives
$x_ay_b-x_by_a=(\cos\varphi_a\sin\varphi_b-\cos\varphi_b\sin\varphi_a)/(u_au_b)=\sin(\varphi_b-\varphi_a)/(u_au_b)$.
Multiplying by $u_1u_2u_3$ turns the three terms into $u_1\sin(\varphi_3-\varphi_2)$, $u_2\sin(\varphi_1-\varphi_3)$
and $u_3\sin(\varphi_2-\varphi_1)$. (The Lean proof expands the determinant by \lean{Matrix.det_fin_three}, applies
the subtraction formula, clears denominators and normalizes; it is the same computation.)
\end{proof}

\begin{lemma}[polar orientation formula; \leanref{C7.orient_polar}{ES8SYM/lean/PolarOrient.lean:37}]\label{ec:lem:polar-u}
For all real $\varphi_1,\varphi_2,\varphi_3,u_1,u_2,u_3$ with $u_1>0$, $u_2>0$, $u_3>0$:
$\operatorname{sign}D=\operatorname{sign}\Sigma$.
\end{lemma}

\begin{proof}
Let $P=u_1u_2u_3>0$. By Lemma~\ref{ec:lem:det-mul}, $\operatorname{sign}\Sigma=\operatorname{sign}(DP)=\operatorname{sign}D$.
\end{proof}

\begin{lemma}[cross-product form; \leanref{C7.orient_polar_cross}{ES8SYM/lean/PolarOrient.lean:50}]\label{ec:lem:polar-cross}
For all real $\varphi_1,\varphi_2,\varphi_3,u_1,u_2,u_3$ with $u_1>0$, $u_2>0$, $u_3>0$:
\begin{equation*}
\operatorname{sign}\mathrm{cross}\bigl(\mathrm{pt}(\varphi_1,u_1),\mathrm{pt}(\varphi_2,u_2),\mathrm{pt}(\varphi_3,u_3)\bigr)
=\operatorname{sign}\Sigma .
\end{equation*}
\end{lemma}

\begin{proof}
With $x_k,y_k$ as before, $\mathrm{pt}(\varphi_k,u_k)=(x_k,y_k)$, and expanding,
\begin{equation*}
(x_2-x_1)(y_3-y_1)-(y_2-y_1)(x_3-x_1)=(x_2y_3-x_3y_2)+(x_3y_1-x_1y_3)+(x_1y_2-x_2y_1),
\end{equation*}
the same three-term expression as $D$. So $\mathrm{cross}(\dots)\cdot u_1u_2u_3=\Sigma$ by the computation in the proof
of Lemma~\ref{ec:lem:det-mul}, and the sign argument of Lemma~\ref{ec:lem:polar-u} applies.
\end{proof}

\begin{proposition}[layered configurations; \leanref{C7.orient_layered}{ES8SYM/lean/PolarOrient.lean:68}]\label{ec:prop:layered}
Let $\iota$ be a type, $t\in\N$, $L:\iota\to\mathrm{Fin}\,t$, $\varphi:\iota\to\R$, and $u:\mathrm{Fin}\,t\to\R$ with
$u(o)>0$ for every $o$. Then for all $i,j,k\in\iota$ \textup{(}not necessarily distinct\textup{)}:
\begin{equation*}
\operatorname{sign}\mathrm{cross}\bigl(\mathrm{pt}(\varphi_i,u(L(i))),\mathrm{pt}(\varphi_j,u(L(j))),\mathrm{pt}(\varphi_k,u(L(k)))\bigr)
=\operatorname{sign}\sum_{o\in\mathrm{Fin}\,t}\mathrm{coef}_{L,\varphi}(i,j,k)_o\,u(o).
\end{equation*}
\end{proposition}

\begin{proof}
For $a\in\mathrm{Fin}\,t$ and $v\in\R$, $\sum_o(v\,e_a)_o\,u(o)=v\,u(a)$, and the sum is linear in the coefficient
vector, so
\begin{equation*}
\sum_o\mathrm{coef}_{L,\varphi}(i,j,k)_o\,u(o)=u(L(i))\sin(\varphi_k-\varphi_j)+u(L(j))\sin(\varphi_i-\varphi_k)
+u(L(k))\sin(\varphi_j-\varphi_i),
\end{equation*}
also when some of $L(i),L(j),L(k)$ coincide. This is $\Sigma$ for $(\varphi_1,\varphi_2,\varphi_3)=(\varphi_i,\varphi_j,\varphi_k)$
and $(u_1,u_2,u_3)=(u(L(i)),u(L(j)),u(L(k)))$, all positive; apply Lemma~\ref{ec:lem:polar-cross}.
\end{proof}

\begin{corollary}[scaling; \leanref{C7.orient_layered_scale}{ES8SYM/lean/PolarOrient.lean:80}]\label{ec:cor:scale}
In the setting of Proposition~\ref{ec:prop:layered}, let $\lambda\in\R$ with $\lambda>0$. Then for all $i,j,k\in\iota$,
\begin{multline*}
\operatorname{sign}\mathrm{cross}\bigl(\mathrm{pt}(\varphi_i,\lambda u(L(i))),\mathrm{pt}(\varphi_j,\lambda u(L(j))),\mathrm{pt}(\varphi_k,\lambda u(L(k)))\bigr)\\
=\operatorname{sign}\mathrm{cross}\bigl(\mathrm{pt}(\varphi_i,u(L(i))),\mathrm{pt}(\varphi_j,u(L(j))),\mathrm{pt}(\varphi_k,u(L(k)))\bigr).
\end{multline*}
\end{corollary}

\begin{proof}
The function $o\mapsto\lambda u(o)$ is positive, so Proposition~\ref{ec:prop:layered} applies to it and to $u$. With
$c_o=\mathrm{coef}_{L,\varphi}(i,j,k)_o$, the left side is
$\operatorname{sign}\sum_oc_o\lambda u(o)=\operatorname{sign}\bigl(\lambda\sum_oc_ou(o)\bigr)=\operatorname{sign}\sum_oc_ou(o)$,
which is the right side. (Directly: $\mathrm{pt}(\varphi,\lambda u)=\lambda^{-1}\mathrm{pt}(\varphi,u)$ and
$\mathrm{cross}$ is homogeneous of degree $2$.)
\end{proof}

\begin{lemma}[polygon vertices in inverse-radius form; \leanref{C7.ngonPt_eq_pt}{ES8SYM/lean/PolarOrient.lean:100}]\label{ec:lem:ngon-pt}
Let $m,t\in\N$, $\theta,r:\mathrm{Fin}\,t\to\R$ with $r_o>0$ for every $o$, and $p\in\mathrm{Fin}\,t\times\mathrm{Fin}\,m$.
Then $\mathrm{ngonPt}_{m,\theta,r}(p)=\mathrm{pt}\bigl(\mathrm{phase}_{m,\theta}(p),\,1/r_{p_1}\bigr)$, where $p_1$ is
the first component of $p$.
\end{lemma}

\begin{proof}
For $r_o\ne0$ and every real $\phi$, $\cos\phi/(1/r_o)=r_o\cos\phi$ and $\sin\phi/(1/r_o)=r_o\sin\phi$.
\end{proof}

\begin{theorem}[nested regular $m$-gons; \leanref{C7.nested_regular_orient}{ES8SYM/lean/PolarOrient.lean:108}]\label{ec:thm:nested}
Let $m,t\in\N$ and $\theta,r:\mathrm{Fin}\,t\to\R$ with $r_o>0$ for every $o$. Then for all
$a,b,c\in\mathrm{Fin}\,t\times\mathrm{Fin}\,m$:
\begin{equation*}
\operatorname{sign}\mathrm{cross}\bigl(\mathrm{ngonPt}_{m,\theta,r}(a),\mathrm{ngonPt}_{m,\theta,r}(b),\mathrm{ngonPt}_{m,\theta,r}(c)\bigr)
=\operatorname{sign}\sum_{o\in\mathrm{Fin}\,t}\mathrm{coef}_{\pi_1,\mathrm{phase}_{m,\theta}}(a,b,c)_o\cdot\frac1{r_o},
\end{equation*}
where $\pi_1:\mathrm{Fin}\,t\times\mathrm{Fin}\,m\to\mathrm{Fin}\,t$ is the first projection \textup{(}Lean
\lean{Prod.fst}\textup{)}. The coefficient vector depends only on $m$ and $\theta$, not on $r$.
\end{theorem}

\begin{proof}
By Lemma~\ref{ec:lem:ngon-pt}, $\mathrm{ngonPt}_{m,\theta,r}(a)=\mathrm{pt}(\mathrm{phase}_{m,\theta}(a),u(\pi_1(a)))$
with $u(o)=1/r_o>0$, and likewise for $b$ and $c$. Apply Proposition~\ref{ec:prop:layered} with
$\iota=\mathrm{Fin}\,t\times\mathrm{Fin}\,m$, $L=\pi_1$, $\varphi=\mathrm{phase}_{m,\theta}$ and this $u$.
\end{proof}

\begin{lemma}[triples with the centre; \leanref{C7.orient_centre}{ES8SYM/lean/PolarOrient.lean:116}]\label{ec:lem:centre-u}
For all real $\varphi_1,\varphi_2,u_1,u_2$ with $u_1>0$ and $u_2>0$:
$\operatorname{sign}\mathrm{cross}\bigl(\mathrm{pt}(\varphi_1,u_1),\mathrm{pt}(\varphi_2,u_2),(0,0)\bigr)=\operatorname{sign}\sin(\varphi_2-\varphi_1)$.
\end{lemma}

\begin{proof}
With $x_k,y_k$ as before, $\mathrm{cross}((x_1,y_1),(x_2,y_2),(0,0))=(x_2-x_1)(-y_1)-(y_2-y_1)(-x_1)=x_1y_2-x_2y_1
=\sin(\varphi_2-\varphi_1)/(u_1u_2)$. Multiply by $u_1u_2>0$ and argue as in Lemma~\ref{ec:lem:polar-u}.
\end{proof}

\begin{remark}[scope of the polar lemmas]\label{ec:rem:polar}
(i) If $\hat p$ denotes the point with coordinates $\hat p_0=p_1$, $\hat p_1=p_2$, then
$\mathrm{cross}(p,q,r)=\orient(\hat p,\hat q,\hat r)$ by the formulas (Definition~\ref{fd:def:orient}), so a positive
value means counterclockwise; Lean has no link between \lean{C7.cross} and \lean{ErdosSzekeres.orient}. By the
expansions in the proofs of Lemmas~\ref{ec:lem:det-mul} and~\ref{ec:lem:polar-cross}, $\mathrm{cross}(p,q,r)$ is
the determinant of the matrix with rows $(p_1,p_2,1)$, $(q_1,q_2,1)$, $(r_1,r_2,1)$.
(ii) Lemma~\ref{ec:lem:polar-u} is not the theorem \lean{orient_polar} of \dirn{LAYERLEAN/LayerLemma.lean:935}
(Lemma~\ref{lg:lem:polar}), which is a radius form about a centre $c$:
$\orient(\mathrm{polar}(c,r_a,\alpha),\mathrm{polar}(c,r_b,\beta),\mathrm{polar}(c,r_d,\delta))=r_ar_b\sin(\beta-\alpha)
+r_br_d\sin(\delta-\beta)+r_dr_a\sin(\alpha-\delta)$, where $\mathrm{polar}(c,r,\alpha)=(c_0+r\cos\alpha,\,c_1+r\sin\alpha)$.
For $u\ne0$ the point $\mathrm{pt}(\varphi,u)$ corresponds as in (i) to $\mathrm{polar}(0,1/u,\varphi)$, and for
$r_a=1/u_1$, $r_b=1/u_2$, $r_d=1/u_3$ this right side times $u_1u_2u_3$ is $\Sigma$ with
$(\varphi_1,\varphi_2,\varphi_3)=(\alpha,\beta,\delta)$, so the two agree in substance,
but they are distinct declarations; the name \lean{orient_polar} in \dirn{WRITEUP/P1/p1.tex} and
\dirn{WRITEUP/note/k8_soundness.tex} refers to the radius form.
(iii) Theorem~\ref{ec:thm:nested} assumes only $r_o>0$: no nesting, no order of the radii, no distinctness of the
points and no general position. It says that for fixed phases every orientation sign is the sign of a linear form in
$(1/r_o)_o$ whose coefficients depend only on $\theta$ and $m$; Corollary~\ref{ec:cor:scale} adds that the signs do
not change when all inverse radii are multiplied by the same $\lambda>0$. Lemma~\ref{ec:lem:centre-u} treats triples
that contain the centre $(0,0)$, which have no finite inverse radius. The further statement of
\dirn{ES8SYM/proofs/layers.md} (lines 26--27) that for fixed phases the order type is constant on the cells of a
central hyperplane arrangement is not formalized as such.
\end{remark}

%---------------------------------------------------------------------------------------------------------------------
\subsection{Size compatibility in a group with zero}\label{ec:ssec:compat}

The source is \dirn{PRIMEFOLD/compat/lean/PCompatCore.lean} (namespace \lean{PCompat}). Here the plane is replaced
by an arbitrary commutative group with zero $K$ (Mathlib \lean{CommGroupWithZero}: a commutative monoid with zero in
which $0\ne1$, with an inverse map such that $aa^{-1}=1$ for every $a\ne0$ and $0^{-1}=0$; the structure also carries a
division and integer powers, which are not used here), with decidable equality (Lean needs it to form finite sets by
image, difference, erasure and insertion, and for the test $0\in S$ in Theorem~\ref{ec:thm:compat}(b); mathematically
it is no restriction). In the application $K=\C$, the centre is $0$ and the rotation is
$z\mapsto\zeta z$; a rotation about another centre reduces to this by a translation. We use only these consequences of
the axioms: $0\cdot a=0$; if $a\ne0$ and $ab=ac$, then $b=c$ (multiply by $a^{-1}$); if $a\ne0$ and $b\ne0$, then
$ab\ne0$ (otherwise $b=a^{-1}(ab)=0$); hence $a^j\ne0$ for $a\ne0$. For a finite $T\subseteq K$, ``$T$ is
$\zeta$-invariant'' means $\zeta z\in T$ for every $z\in T$. We write $[P]$ for $1$ if $P$ holds and $0$ otherwise,
and $N\bmod m$ for the remainder (Lean \texttt{N \% m}). The hypotheses on $m\in\N$, $\zeta\in K$ and $r:\N\to K$ are:
\begin{itemize}
\item[(O)] $0<m$, $\zeta^m=1$, and $\zeta^j\ne1$ for every $j$ with $0<j<m$ \ ($\zeta$ has exact order $m$);
\item[(R0)] $r_i\ne0$ for every $i\in\N$;
\item[(RD)] for all $i,i',j,j'\in\N$ with $j<m$ and $j'<m$: if $\zeta^jr_i=\zeta^{j'}r_{i'}$, then $i=i'$ \ (the
  orbits of the $r_i$ are pairwise disjoint).
\end{itemize}

\begin{lemma}[powers; \leanref{PCompat.ne_zero_of_order}{PRIMEFOLD/compat/lean/PCompatCore.lean:36}, \leanref{PCompat.pow_inj_of_order}{PRIMEFOLD/compat/lean/PCompatCore.lean:43}, \leanref{PCompat.pow_mul_mem}{PRIMEFOLD/compat/lean/PCompatCore.lean:58}]\label{ec:lem:pc-basic}
Let $K$ be a commutative group with zero, $m\in\N$ and $\zeta\in K$.
\begin{enumerate}[label=(\alph*)]
\item If $0<m$ and $\zeta^m=1$, then $\zeta\ne0$.
\item If \textup{(O)} holds, $i<m$, $j<m$ and $\zeta^i=\zeta^j$, then $i=j$.
\item If $T\subseteq K$ is finite and $\zeta$-invariant and $z\in T$, then $\zeta^jz\in T$ for every $j\in\N$.
\end{enumerate}
\end{lemma}

\begin{proof}
(a) If $\zeta=0$, then $\zeta^m=0$ since $m>0$, so $0=1$, which is false in $K$.
(b) By (a), $\zeta\ne0$. Suppose $i<j$. Then $\zeta^i\zeta^{j-i}=\zeta^j=\zeta^i=\zeta^i\cdot1$, and $\zeta^i\ne0$, so
$\zeta^{j-i}=1$ with $0<j-i<m$, contradicting (O). The case $j<i$ is symmetric. Hence $i=j$.
(c) Induction on $j$: $\zeta^0z=z\in T$, and $\zeta^{j+1}z=\zeta(\zeta^jz)\in T$ by invariance.
\end{proof}

\begin{lemma}[the converse constructions; \leanref{PCompat.exists_card0}{PRIMEFOLD/compat/lean/PCompatCore.lean:129}, \leanref{PCompat.exists_card1}{PRIMEFOLD/compat/lean/PCompatCore.lean:160}]\label{ec:lem:pc-exist}
Let $K$ be a commutative group with zero, $m\in\N$, $\zeta\in K$ and $r:\N\to K$, and assume \textup{(O)},
\textup{(R0)} and \textup{(RD)}. Then for every $t\in\N$:
\begin{enumerate}[label=(\alph*)]
\item there is a finite $\zeta$-invariant $S\subseteq K$ with $0\notin S$ and $\card S=mt$;
\item there is a finite $\zeta$-invariant $S\subseteq K$ with $\card S=mt+1$.
\end{enumerate}
\end{lemma}

\begin{proof}
(a) Let $P=\{0,\dots,m-1\}\times\{0,\dots,t-1\}$ and $S=\{\zeta^jr_i:(j,i)\in P\}$. \emph{Invariance:} for $(j,i)\in
P$, if $j+1<m$ then $\zeta\cdot\zeta^jr_i=\zeta^{j+1}r_i$ with $(j+1,i)\in P$; if $j+1=m$ then
$\zeta\cdot\zeta^jr_i=\zeta^mr_i=r_i=\zeta^0r_i$ with $(0,i)\in P$ (as $0<m$). \emph{$0\notin S$:} $\zeta\ne0$ by
Lemma~\ref{ec:lem:pc-basic}(a), so $\zeta^j\ne0$, and $r_i\ne0$ by (R0), so $\zeta^jr_i\ne0$. \emph{Cardinality:} the
map $(j,i)\mapsto\zeta^jr_i$ is injective on $P$: if $\zeta^jr_i=\zeta^{j'}r_{i'}$ with $(j,i),(j',i')\in P$, then
$i=i'$ by (RD), then $\zeta^j=\zeta^{j'}$ by cancelling $r_i\ne0$, then $j=j'$ by Lemma~\ref{ec:lem:pc-basic}(b). Hence
$\card S=\card P=mt$.

(b) Take $S$ from (a) and $S'=S\cup\{0\}$. It is $\zeta$-invariant ($\zeta\cdot0=0$), and $\card{S'}=mt+1$ because
$0\notin S$.
\end{proof}

\begin{definition}[size condition; \leanref{PCompat.Compat}{PRIMEFOLD/compat/lean/PCompatCore.lean:205}]\label{ec:def:compat}
For $m,N\in\N$, $\mathrm{Compat}(m,N)$ means: $N\bmod m=0$ or $N\bmod m=1\bmod m$. (For every $m\ge1$ this is the
same condition as $N\bmod m\in\{0,1\}$: for $m\ge2$, $1\bmod m=1$, and for $m=1$ both conditions hold, since
$N\bmod1=0$. The form $1\bmod m$ matches the right side of Theorem~\ref{ec:thm:compat}(b). For $m=0$, Lean's
convention $N\bmod0=N$ applies, but $m=0$ is not used.)
\end{definition}

\begin{proposition}[parity; \leanref{PCompat.compat_even}{PRIMEFOLD/compat/lean/PCompatCore.lean:208}, \leanref{PCompat.compat_odd}{PRIMEFOLD/compat/lean/PCompatCore.lean:235}]\label{ec:prop:parity}
Both parts are statements about natural numbers.
\begin{enumerate}[label=(\alph*)]
\item For every even $m\in\N$ with $2\le m$: $\mathrm{Compat}(m,2^{m-1})$ holds if and only if $m=2^a$ for some
  $a\ge1$.
\item For every odd $m\in\N$ with $3\le m$: $\mathrm{Compat}(m,2^{m-1})$ holds if and only if $2^{m-1}\bmod m=1$.
\end{enumerate}
\end{proposition}

\begin{proof}
We use the standard fact (unique factorization; Mathlib \lean{Nat.dvd_prime_pow}) that the divisors of $2^n$ are
exactly $2^0,2^1,\dots,2^n$. In both parts $m\ge2$, so $1\bmod m=1$ and $\mathrm{Compat}(m,2^{m-1})$ says
$2^{m-1}\bmod m\in\{0,1\}$.

(a) Suppose $2^{m-1}\bmod m=0$. Then $m\mid2^{m-1}$, so $m=2^a$ with $a\le m-1$, and $a\ne0$ because $m\ge2$. Suppose
$2^{m-1}\bmod m=1$. Write $2^{m-1}=qm+1$; since $2\mid m$, $2^{m-1}$ is odd. But $m-1\ge1$, so $2^{m-1}$ is even, a
contradiction. Conversely, let $m=2^a$ with $a\ge1$. Since $a<2^a=m$, we have $a\le m-1$, so $m=2^a\mid2^{m-1}$ and
$2^{m-1}\bmod m=0$.

(b) Suppose $2^{m-1}\bmod m=0$. Then $m=2^a$ for some $a$; $a=0$ gives $m=1<3$, and $a\ge1$ makes $m$ even; both are
impossible. So $\mathrm{Compat}(m,2^{m-1})$ holds exactly when $2^{m-1}\bmod m=1$.
\end{proof}

\begin{theorem}[compatibility; \leanref{PCompat.dvd_card}{PRIMEFOLD/compat/lean/PCompatCore.lean:67}, \leanref{PCompat.card_mod}{PRIMEFOLD/compat/lean/PCompatCore.lean:112}, \leanref{PCompat.exists_iff}{PRIMEFOLD/compat/lean/PCompatCore.lean:174}, \leanref{PCompat.classification}{PRIMEFOLD/compat/lean/PCompatCore.lean:253}]\label{ec:thm:compat}
Let $m\in\N$, and in \textup{(a)}--\textup{(c)} let $K$ be a commutative group with zero and $\zeta\in K$.
\begin{enumerate}[label=(\alph*)]
\item \textup{(\lean{dvd_card})} If \textup{(O)} holds, then every finite $\zeta$-invariant $T\subseteq K$ with
  $0\notin T$ satisfies $m\mid\card T$.
\item \textup{(\lean{card_mod})} If \textup{(O)} holds, then every finite $\zeta$-invariant $S\subseteq K$ satisfies
  $\card S\bmod m=[0\in S]\bmod m$.
\item \textup{(\lean{exists_iff})} Let $r:\N\to K$. If \textup{(O)}, \textup{(R0)} and \textup{(RD)} hold, then for
  every $N\in\N$: there is a finite $\zeta$-invariant $S\subseteq K$ with $\card S=N$ if and only if $N\bmod m=0$ or
  $N\bmod m=1\bmod m$ \textup{(}that is, $\mathrm{Compat}(m,N)$\textup{)}.
\item \textup{(\lean{classification}; this part does not involve $K$ and $\zeta$)} If $1\le m$, then
  $\mathrm{Compat}(m,2^{m-1})$ holds if and only if $m=1$, or $m=2^a$ for some $a\ge1$, or \textup{(}$m$ is odd,
  $3\le m$ and $2^{m-1}\bmod m=1$\textup{)}.
\end{enumerate}
\end{theorem}

\begin{proof}
We follow the Lean proofs.

(a) By Lemma~\ref{ec:lem:pc-basic}(a), $\zeta\ne0$. We use strong induction on $\card T$. If $T=\emptyset$, then
$m\mid0$. Otherwise pick $z\in T$; $z\ne0$ because $0\notin T$. Let $O=\{\zeta^jz:0\le j<m\}$. By
Lemma~\ref{ec:lem:pc-basic}(c), $O\subseteq T$. The map $j\mapsto\zeta^jz$ is injective on $\{0,\dots,m-1\}$: if
$\zeta^iz=\zeta^jz$, then $\zeta^i=\zeta^j$ (cancel $z\ne0$) and $i=j$ by Lemma~\ref{ec:lem:pc-basic}(b). So
$\card O=m$, and $\card T=\card{T\setminus O}+m$. The set $T\setminus O$ omits $0$ and is $\zeta$-invariant: let
$w\in T\setminus O$; then $\zeta w\in T$. Suppose $\zeta w\in O$, say $\zeta w=\zeta^jz$ with $j<m$. If $j=0$, then
$\zeta^{m-1}z=\zeta^{m-1}\zeta w=\zeta^mw=w$, so $w\in O$ (as $m-1<m$). If $j\ge1$, then
$\zeta w=\zeta\cdot\zeta^{j-1}z$, so $w=\zeta^{j-1}z\in O$ (cancel $\zeta\ne0$). Both contradict $w\notin O$, so
$\zeta w\in T\setminus O$. Since $z\in O$, $\card{T\setminus O}<\card T$, and the induction hypothesis gives
$m\mid\card{T\setminus O}$. Hence $m\mid\card T$.

(b) Let $T_0=S\setminus\{0\}$. Then $0\notin T_0$, and $T_0$ is $\zeta$-invariant: for $z\in T_0$, $\zeta z\in S$ and
$\zeta z\ne0$ since $\zeta\ne0$ and $z\ne0$. By (a), $\card{T_0}=mt$ for some $t$. If $0\in S$, then $\card S=mt+1$
and $\card S\bmod m=1\bmod m=[0\in S]\bmod m$. If $0\notin S$, then $S=T_0$ and $\card S\bmod m=0=[0\in S]\bmod m$.

(c) ($\Rightarrow$) Let $S$ be finite, $\zeta$-invariant, with $\card S=N$. By (b), $N\bmod m=1\bmod m$ if $0\in S$,
and $N\bmod m=0\bmod m=0$ if $0\notin S$. ($\Leftarrow$) Write $N=m\lfloor N/m\rfloor+(N\bmod m)$. If $N\bmod m=0$,
Lemma~\ref{ec:lem:pc-exist}(a) with $t=\lfloor N/m\rfloor$ gives a $\zeta$-invariant $S$ with $\card S=N$. Otherwise
$N\bmod m=1\bmod m$. Then $m\ne1$ (for $m=1$ every remainder is $0$), so $m>1$, $1\bmod m=1$, $N=m\lfloor N/m\rfloor+1$,
and Lemma~\ref{ec:lem:pc-exist}(b) with $t=\lfloor N/m\rfloor$ gives the set.

(d) If $m$ is even, then $m\ge2$, and by Proposition~\ref{ec:prop:parity}(a) $\mathrm{Compat}(m,2^{m-1})$ holds iff
$m=2^a$ with $a\ge1$; the other two alternatives of the right side are false for even $m$. If $m$ is odd, then $m=1$ or
$m\ge3$. For $m=1$ both sides hold: $2^0\bmod1=0$, and $m=1$. For odd $m\ge3$, by Proposition~\ref{ec:prop:parity}(b)
the left side holds iff $2^{m-1}\bmod m=1$; on the right side $m=1$ is false and $m=2^a$ with $a\ge1$ is impossible
($m$ is odd), so the right side holds iff $2^{m-1}\bmod m=1$.
\end{proof}

\begin{proposition}[$m=341$; computer-assisted in Lean; written proof by hand; \leanref{PCompat.ex341}{PRIMEFOLD/compat/lean/PCompatCore.lean:277}]\label{ec:prop:ex341}
$\mathrm{Compat}(341,2^{341-1})$; $341=11\cdot31$; $341$ is not prime; $2^{10}\bmod11=1$; $2^{10}\bmod31=1$;
$2^{10}\bmod341=1$.
\end{proposition}

\begin{proof}
This is a finite computation with closed numerals. Lean checks the five numerical claims by kernel evaluation
(\lean{decide}, with \lean{exponentiation.threshold} raised to $400$; no \lean{native_decide}) and non-primality by
Mathlib's \lean{Nat.not_prime_mul} with $11\ne1$ and $31\ne1$; the log lists only \lean{propext}. The computation is
small enough to do by hand, as follows. $11\cdot31=341$, and $11$ is a divisor with
$1<11<341$, so $341$ is not prime. Further $2^{10}=1024=93\cdot11+1=33\cdot31+1=3\cdot341+1$, which gives the three
remainders. Finally $2^{340}=(2^{10})^{34}\equiv1^{34}=1\pmod{341}$, so $2^{340}\bmod341=1=1\bmod341$ and
$\mathrm{Compat}(341,2^{340})$ holds by its second alternative.
\end{proof}

\begin{remark}[what is and is not formalized]\label{ec:rem:compat}
(i) For $K=\C$ and $\zeta=\zeta_m=e^{2\pi i/m}$ with $m\ge1$, (O) holds by Lemma~\ref{ec:lem:zeta} and
Fact~\ref{ec:fact:prim}(iv). (ii) In Theorem~\ref{ec:thm:compat}(c) the hypotheses (R0) and (RD) on $r$ are not
discharged in Lean for $\C$. They hold for $K=\C$, $\zeta=\zeta_m$ and $r_i=i+1$ (a written step): (R0) is clear, and
$\abs{\zeta_m^jr_i}=\abs{\zeta_m}^j(i+1)=i+1$ by Fact~\ref{ec:fact:prim}(ii), so $\zeta_m^jr_i=\zeta_m^{j'}r_{i'}$ forces
$i+1=i'+1$. The sets of Lemma~\ref{ec:lem:pc-exist} need not be in general position (for $m=2$ and $r_i=i+1$ they lie
on the real axis); the existence of size-compatible sets in general position
(\dirn{WRITEUP/note/symmetric_extremal.tex}, the proposition labelled \texttt{prop:compat}) is not formalized.
(iii) Fermat's little theorem is not formalized: the third alternative of Theorem~\ref{ec:thm:compat}(d) is the
arithmetic condition $2^{m-1}\bmod m=1$; by Fermat's little theorem it holds for every odd prime $m$, and for odd
composite $m$ it says that $m$ is a base-$2$ Fermat pseudoprime, as for $m=341$.
\end{remark}

%---------------------------------------------------------------------------------------------------------------------
\subsection{Symmetric sets of \texorpdfstring{$2^{m-1}$}{2\^{}(m-1)} points}\label{ec:ssec:l0}

The source is \dirn{PRIMEFOLD/compat/lean/L0Check.lean} (namespace \lean{L0Check}), an independent check of Lemma~0
of \dirn{PRIMEFOLD/cluster/PROOFS.md}. The plane is $\C$, the centre is $0$, and the rotation by $2\pi/m$ is
multiplication by a primitive $m$-th root of unity. Mathlib's base-$2$ Fermat pseudoprimes are defined by
\begin{equation*}
\mathrm{FermatPsp}(m,2)\iff m\mid2^{m-1}-1,\ m\text{ is not prime, and }1<m
\end{equation*}
(\lean{Nat.FermatPsp}, with natural subtraction; $2^{m-1}\ge1$, so $2^{m-1}-1$ is the usual difference).

\begin{lemma}[arithmetic core; \leanref{L0Check.arith}{PRIMEFOLD/compat/lean/L0Check.lean:21}]\label{ec:lem:arith}
Let $m,t,e\in\N$ with $2\le m$, $e\le1$ and $mt+e=2^{m-1}$. Then either \textup{(}$m$ is odd, $e=1$, and $m$ is prime
or $\mathrm{FermatPsp}(m,2)$\textup{)}, or \textup{(}$e=0$ and $m=2^j$ for some $j\in\N$\textup{)}.
\end{lemma}

\begin{proof}
\emph{$m$ even.} Then $mt$ is even, and $2^{m-1}$ is even because $m-1\ge1$. Hence $e=2^{m-1}-mt$ is even, and $e\le1$
gives $e=0$. Then $mt=2^{m-1}$, so $m\mid2^{m-1}$ and $m=2^j$ for some $j$ (the divisors of $2^n$ are powers of $2$).
The second alternative holds.

\emph{$m$ odd.} If $e\ne1$, then $e=0$, $m\mid2^{m-1}$ and $m=2^j$; $j=0$ gives $m=1<2$, and $j\ge1$ makes $m$ even;
both are impossible. So $e=1$ and $mt+1=2^{m-1}$, hence $m\mid mt=2^{m-1}-1$. If $m$ is prime, the first
alternative holds; otherwise $m\mid2^{m-1}-1$, $m$ is not prime and $1<m$, which is $\mathrm{FermatPsp}(m,2)$.
\end{proof}

\begin{lemma}[orbit counting in $\C$; \leanref{L0Check.orbit_dvd}{PRIMEFOLD/compat/lean/L0Check.lean:60}]\label{ec:lem:orbit-dvd}
Let $m\in\N$ with $0<m$, and let $\zeta\in\C$ be a primitive $m$-th root of unity. Then every finite $T\subseteq\C$
with $0\notin T$ and $\zeta z\in T$ for all $z\in T$ satisfies $m\mid\card T$.
\end{lemma}

\begin{proof}
$\C$ is a commutative group with zero, $\zeta^m=1$, and $\zeta^j\ne1$ for $0<j<m$ by Fact~\ref{ec:fact:prim}(iv); so
(O) holds and Theorem~\ref{ec:thm:compat}(a) applies. (The Lean file proves this statement again, independently, by
the same strong induction, using Mathlib's \lean{IsPrimitiveRoot.pow_inj} for the size of an orbit.)
\end{proof}

\begin{theorem}[sets of $2^{m-1}$ points; \leanref{L0Check.L0}{PRIMEFOLD/compat/lean/L0Check.lean:122}, \leanref{L0Check.L0_rot}{PRIMEFOLD/compat/lean/L0Check.lean:144}]\label{ec:thm:l0}
Let $m\in\N$ with $2\le m$.
\begin{enumerate}[label=(\alph*)]
\item \textup{(\lean{L0})} Let $\zeta\in\C$ be a primitive $m$-th root of unity, and let $S\subseteq\C$ be finite with
  $\zeta z\in S$ for all $z\in S$ and $\card S=2^{m-1}$. Then either \textup{(}$m$ is odd, $0\in S$, and $m$ is prime
  or $\mathrm{FermatPsp}(m,2)$\textup{)}, or \textup{(}$0\notin S$ and $m=2^j$ for some $j\in\N$\textup{)}.
\item \textup{(\lean{L0_rot})} Let $S\subseteq\C$ be finite with $\exp(2\pi i/m)\,z\in S$ for all $z\in S$ and
  $\card S=2^{m-1}$. Then the same conclusion holds.
\end{enumerate}
\end{theorem}

\begin{proof}
(a) By Fact~\ref{ec:fact:prim}(i), $\zeta\ne0$. The set $T=S\setminus\{0\}$ omits $0$ and is invariant: for $z\in T$,
$\zeta z\in S$ and $\zeta z\ne0$. By Lemma~\ref{ec:lem:orbit-dvd}, $\card T=mt$ for some $t\in\N$. If $0\in S$, then
$mt+1=\card S=2^{m-1}$, and Lemma~\ref{ec:lem:arith} with $e=1$ gives the first alternative of the lemma (the second
needs $e=0$); together with $0\in S$ this is the first alternative of the theorem. If $0\notin S$, then $T=S$ and
$mt+0=2^{m-1}$; Lemma~\ref{ec:lem:arith} with $e=0$ excludes its first alternative (which needs $e=1$), so $m=2^j$ for
some $j$, and with $0\notin S$ this is the second alternative of the theorem.

(b) $\exp(2\pi i/m)=\zeta_m$ is a primitive $m$-th root of unity by Lemma~\ref{ec:lem:zeta} ($m\ge2>0$); apply (a).
\end{proof}

\begin{remark}[on the hypotheses]\label{ec:rem:l0}
No general-position or convexity hypothesis is assumed. For even $m$ the centre is excluded by parity, as in the proof
of Lemma~\ref{ec:lem:arith}: $\card S=2^{m-1}$ is even and $\card{S\setminus\{0\}}$ is a multiple of $m$. (Lemma~0 of
\dirn{PRIMEFOLD/cluster/PROOFS.md} excludes the centre by general position instead; the Lean statement does not
need it.) For odd $m$ the theorem only says that $m$ is prime or a base-$2$ Fermat pseudoprime; that every odd prime
actually satisfies $2^{m-1}\equiv1\pmod m$ is Fermat's little theorem, which is not used and not formalized here. For
a rotation about a centre $c\ne0$, translate by $-c$ first. For $k=8$, i.e.\ $m=7$ and $2^{m-1}=64=7\cdot9+1$: a
finite $S\subseteq\C$ with $e^{2\pi i/7}S\subseteq S$ and $\card S=64$ contains $0$, by part (b), since $7$ is odd and
not a power of $2$.
\end{remark}

%---------------------------------------------------------------------------------------------------------------------
\paragraph{Coverage.} Covered, by label:
\begin{itemize}
\item \dirn{LEANEXT/es7corr/ES7Corollaries.lean}, Section~11: \lean{Fivefold35} and \lean{Fourfold32}
  (Definition~\ref{ec:def:cases}); \lean{fivefold_heptagon_of_erase} (Theorem~\ref{ec:thm:fivefold-erase});
  \lean{fivefold_heptagon} (Theorem~\ref{ec:thm:fivefold}); \lean{fivefold_heptagon_of_centre_not_mem}
  (Theorem~\ref{ec:thm:fivefold-nc}); \lean{fourfold_heptagon_of_erase} (Theorem~\ref{ec:thm:fourfold-erase});
  \lean{fourfold_heptagon} (Theorem~\ref{ec:thm:fourfold}). The three fivefold theorems are conditional on
  \lean{Fivefold35}, the two fourfold theorems on \lean{Fourfold32}.
\item \dirn{ES8SYM/lean/C6check.lean}: \lean{sphere_subset_extremePoints} (Lemma~\ref{ec:lem:sphere});
  \lean{convexIndependent_of_norm_eq} (Lemma~\ref{ec:lem:ci-circle}); \lean{zeta} (Definition~\ref{ec:def:zeta});
  \lean{zeta_prim} (Lemma~\ref{ec:lem:zeta}); \lean{orbit_injective} and \lean{orbit_convexIndependent}
  (Lemma~\ref{ec:lem:orbit8}); \lean{part_a} (Theorem~\ref{ec:thm:big-order}); \lean{thr_m4}--\lean{thr_m7} and
  \lean{m4_with_centre_counter} (Proposition~\ref{ec:prop:thr}).
\item \dirn{ES8SYM/lean/PolarOrient.lean}: \lean{cross}, \lean{pt}, \lean{coef}, \lean{phase}, \lean{ngonPt}
  (Definition~\ref{ec:def:polar}); \lean{det_mul_eq} (Lemma~\ref{ec:lem:det-mul}); \lean{orient_polar}
  (Lemma~\ref{ec:lem:polar-u}); \lean{orient_polar_cross} (Lemma~\ref{ec:lem:polar-cross}); \lean{orient_layered}
  (Proposition~\ref{ec:prop:layered}); \lean{orient_layered_scale} (Corollary~\ref{ec:cor:scale});
  \lean{ngonPt_eq_pt} (Lemma~\ref{ec:lem:ngon-pt}); \lean{nested_regular_orient} (Theorem~\ref{ec:thm:nested});
  \lean{orient_centre} (Lemma~\ref{ec:lem:centre-u}).
\item \dirn{PRIMEFOLD/compat/lean/PCompatCore.lean}: \lean{ne_zero_of_order}, \lean{pow_inj_of_order},
  \lean{pow_mul_mem} (Lemma~\ref{ec:lem:pc-basic}); \lean{exists_card0}, \lean{exists_card1}
  (Lemma~\ref{ec:lem:pc-exist}); \lean{Compat} (Definition~\ref{ec:def:compat}); \lean{compat_even},
  \lean{compat_odd} (Proposition~\ref{ec:prop:parity}); \lean{dvd_card}, \lean{card_mod}, \lean{exists_iff},
  \lean{classification} (Theorem~\ref{ec:thm:compat}); \lean{ex341} (Proposition~\ref{ec:prop:ex341}, a finite
  computation).
\item \dirn{PRIMEFOLD/compat/lean/L0Check.lean}: \lean{arith} (Lemma~\ref{ec:lem:arith}); \lean{orbit_dvd}
  (Lemma~\ref{ec:lem:orbit-dvd}); \lean{L0}, \lean{L0_rot} (Theorem~\ref{ec:thm:l0}).
\end{itemize}
Not covered, deliberately: lines 1--1147 of \dirn{LEANEXT/es7corr/ES7Corollaries.lean} (outside the scope of this
section), a verbatim copy of lines 1--1147 of \dirn{LAYERLEAN/LayerLemma.lean}, which belongs to the layer-lemma
sections (we use Theorem~\ref{ll:thm:reduction} and Lemma~\ref{lg:lem:centre} from there); the \texttt{\#print axioms}
commands at the ends of the four other files (\dirn{ES8SYM/lean/C6check.lean}, lines 119--123;
\dirn{ES8SYM/lean/PolarOrient.lean}, lines 127--132; \dirn{PRIMEFOLD/compat/lean/PCompatCore.lean}, lines 285--291;
\dirn{PRIMEFOLD/compat/lean/L0Check.lean}, line 152), which produce the logs cited above and are not results. Out of
scope for this section: the Lean files of \dirn{ES8SLACK} (\dirn{ES8SLACK/review/C6_early_vertex/EarlyVertex.lean},
\dirn{ES8SLACK/review/C2_identity/Identity.lean}, and the two files that according to \dirn{ES8SLACK/README.md} have
not been compiled, \dirn{ES8SLACK/review/C1_certificates/Composition311.lean} and
\dirn{ES8SLACK/review/C3_residual/TemplateRowsC3.lean}); written proofs of the \dirn{ES8SLACK} lemmas are in
\dirn{ES8SLACK/proofs.md}. Fact~\ref{ec:fact:prim}, Corollary~\ref{ec:cor:planar} and the example in
Remark~\ref{ec:rem:compat}(ii) are written steps and not declarations of the five Lean files; only the Fact has Lean
counterparts, the Mathlib lemmas named in its proof.
\endgroup

\section{Reduction of symmetric sets to case classes}\label{sr:sec}

This section gives written proofs of the declarations of \dirn{SYMLEAN/SymmetricReduction.lean}. Its main result,
Theorem~\ref{sr:thm:main} (\lean{heptagon_of_nontrivial_isometry}), is a conditional statement. Five propositions
about point sets of an exact size with one prescribed symmetry, the \emph{case classes} \lean{HalfTurn34},
\lean{Threefold33}, \lean{Fivefold35}, \lean{MirrorOneAxis33} and \lean{MirrorNoAxis34}
(Definition~\ref{sr:def:classes}), imply the following: every finite set of at least $33$ points in general position
that is mapped into itself by a Euclidean isometry of the plane other than the identity contains $7$ points in convex
position. The classes are \emph{hypotheses} of the theorem. Nothing in this section proves them;
Remark~\ref{sr:rem:records} says which project records supply them and with which rigor label. The section also proves
the variant Theorem~\ref{sr:thm:variant} (\lean{halfTurn_of_symmetric_counterexample}), in which a fourfold class
\lean{Fourfold32} replaces \lean{HalfTurn34}, and the bridge Lemma~\ref{sr:lem:mirror-line} to a mirror statement of
the directory \dirn{H3}.

The proof of Theorem~\ref{sr:thm:main} has three parts.
\begin{enumerate}
\item Classification (Theorem~\ref{sr:thm:structure}). A Euclidean isometry $g\ne\mathrm{id}$ of the plane that maps
  a finite set with at least two points into itself is a rotation $z\mapsto c+a(z-c)$ with $a\ne1$ a root of unity, or a
  reflection in a line. Translations and glide reflections are excluded because the set is finite.
\item Rotations. A power of $g$ is the half-turn, or a rotation whose factor is a primitive $p$-th root of unity for an
  odd prime $p$ (Lemma~\ref{sr:lem:root-cases}). Whole orbits are selected (Lemma~\ref{sr:lem:orbit}) to reach the
  exact sizes of the classes: $17$ antipodal pairs; $11$ orbits of size $3$; $7$ orbits of size $5$. For $p\ge7$ one
  orbit consists of $p$ concyclic points, which are in convex position (Lemma~\ref{sr:lem:concyclic}).
\item Reflections. At most two points lie on the mirror and the others come in pairs; $17$ pairs, or $16$ pairs and one
  point of the mirror, are kept.
\end{enumerate}

\paragraph{Sources.}
Paths of research records are relative to \dirn{cloud/es7-2026-10-19/results}; paths of the main library are relative
to the repository root. The first 12 hex digits of the SHA-256 of each Lean source used are:
\begin{itemize}
\item \dirn{SYMLEAN/SymmetricReduction.lean} (1059 lines; every declaration of this file is covered here):
  \texttt{6d5ec90a4215};
\item \dirn{SYMLEAN/SymmetricReductionAxioms.lean}: \texttt{7bf809cbdc77}. Its lines 1--1059 are byte-identical to
  the previous file; two \texttt{\#print axioms} commands follow;
\item the eight repository modules that the file imports, directly or through other imports:
  \begin{itemize}
  \item \dirn{ErdosSzekeres/Definitions.lean}: \texttt{2425500eefd1};
  \item \dirn{ErdosSzekeres/Basic.lean}: \texttt{13667aa2d3dc};
  \item \dirn{ErdosSzekeres/Orientation.lean}: \texttt{f67955ac7091};
  \item \dirn{ErdosSzekeres/OrientationGeometry.lean}: \texttt{c7db9bbbb65e};
  \item \dirn{ErdosSzekeres/Suk/SupportingLines.lean}: \texttt{6a8e12677685};
  \item \dirn{ErdosSzekeres/Suk/BaseOrder.lean}: \texttt{d21becd29e84};
  \item \dirn{ErdosSzekeres/Suk/BaseChart.lean}: \texttt{08d0e856f005};
  \item \dirn{ErdosSzekeres/Suk/TriangleOrder.lean}: \texttt{6447cd4833d4}.
  \end{itemize}
  The last three are imported only through \dirn{ErdosSzekeres/Suk/SupportingLines.lean}, and this section cites
  nothing from them.
\end{itemize}
The files were checked with Lean 4.24.0 and Mathlib \texttt{f897ebcf}.

\paragraph{Axioms and rigor.}
\begin{itemize}
\item The compile log \dirn{SYMLEAN/logs/SymmetricReduction.out} is empty, and the exit code is 0
  (\dirn{SYMLEAN/logs/runs.tsv}). So no declaration of the file uses \lean{sorry}.
\item The log \dirn{SYMLEAN/logs/SymmetricReductionAxioms.out} shows that the two theorems
  \lean{heptagon_of_nontrivial_isometry} and \lean{halfTurn_of_symmetric_counterexample} depend only on the axioms
  \lean{propext}, \lean{Classical.choice} and \lean{Quot.sound}.
\item The bridge lemma \lean{mirrorOneAxis33_of_exclusionLine} is used in the proof of the review theorem
  \lean{Review.mirrorOneAxis33_iff_exclusionLine} at \dirn{SYMLEAN/review/ReviewSanity.lean:1629}. That file's lines 1--1059
  are byte-identical to \dirn{SYMLEAN/SymmetricReduction.lean}, and its log \dirn{SYMLEAN/logs/ReviewSanity.out} shows
  the same three axioms for this theorem.
\item The file and the eight imported repository modules listed above contain no \lean{axiom} declaration, no
  \lean{sorry}, no \lean{decide} and no \lean{native_decide} (the word ``axiom'' occurs there only in comments). A
  declaration of the file that the printed theorems do not use (such as
  \lean{IsEucIsometry.injective}) is covered by the empty compile log and by this scan: it uses no \lean{sorry}, and
  neither the file nor these modules add an axiom. No statement of this section is a finite computation.
\end{itemize}

Every statement below whose head carries a Lean reference is \emph{Lean-checked} in this sense. The proofs written here
are \emph{written proofs (not refereed)}. Unless a proof says otherwise, it follows the Lean proof step by step.
Statements tagged ``written step'' are introduced only for the written proofs; most of them are standard facts that
the Lean proofs take from Mathlib, and each is proved here.

\paragraph{Dependencies.}
The section uses the foundations: Definitions~\ref{fd:def:point}, \ref{fd:def:orient}, \ref{fd:def:gp},
\ref{fd:def:cp}, \ref{fd:def:ngon}, \ref{fd:def:inv} and \ref{fd:def:rot}, and Lemmas~\ref{fd:lem:collinear},
\ref{fd:lem:cp-hull} and~\ref{fd:lem:mono}. Lemma~\ref{sr:lem:supports} is the lemma titled ``Strict supporting
lines'' in the foundations section (the same declaration of the main library), repeated here with its short proof.
Remarks~\ref{sr:rem:records}
and~\ref{sr:rem:main} refer to Theorems~\ref{cb:thm:c2} and~\ref{cb:thm:c3} (bookkeeping of two case splits; both
theorems are computer-assisted, as their Lean proofs rest on kernel evaluation, and Theorem~\ref{cb:thm:c2} is
Lean-checked according to the record, Theorem~\ref{cb:thm:c3} Lean-checked), to Theorem~\ref{mb:thm:classes}, which
gives written proofs of the two mirror classes, to Theorem~\ref{mb:thm:main3} and to Theorem~\ref{ar:thm:main}; no
proof of this section uses them.

\subsection{Setting and definitions}\label{sr:ssec:defs}

\paragraph{Conventions.}
Points are elements of $\R^2$, written $p=(p_0,p_1)$ (Definition~\ref{fd:def:point}; Lean type \lean{Point}). Point
sets are finite (Lean \texttt{Finset Point}), and $\card S$ is the number of elements of $S$. The function $\orient$ is
that of Definition~\ref{fd:def:orient}: $\orient(a,b,c)=(b_0-a_0)(c_1-a_1)-(b_1-a_1)(c_0-a_0)$. General position,
convex position and \texttt{HasConvexNGon n S} are as in Definitions~\ref{fd:def:gp}, \ref{fd:def:cp}
and~\ref{fd:def:ngon}. We use them in the following forms.
\begin{itemize}
\item $S$ is in general position if and only if $\orient(a,b,c)\ne0$ for all pairwise distinct $a,b,c\in S$
  (Lemma~\ref{fd:lem:collinear}).
\item $S$ is in convex position if and only if $p\notin\conv(S\setminus\{p\})$ for every $p\in S$
  (Lemma~\ref{fd:lem:cp-hull}).
\item \texttt{HasConvexNGon n S} means: some $T\subseteq S$ with $\card T=n$ is in convex position.
\item Monotonicity (Lemma~\ref{fd:lem:mono}): subsets of a set in general position are in general position; subsets of
  a set in convex position are in convex position; \texttt{HasConvexNGon n T} and $T\subseteq S$ imply
  \texttt{HasConvexNGon n S}.
\item If $T$ is in convex position and $n\le\card T$, then \texttt{HasConvexNGon n T}: any $n$-element subset of $T$ is
  in convex position by monotonicity (Lean \lean{ConvexPosition.hasConvexNGon}, \dirn{ErdosSzekeres/Basic.lean:31}).
\end{itemize}
For a map $f\colon X\to X$ the iterates are $f^{[0]}=\mathrm{id}$ and $f^{[k+1]}=f\circ f^{[k]}$; then also
$f^{[k+1]}=f^{[k]}\circ f$ and $f^{[i+j]}=f^{[i]}\circ f^{[j]}$. Natural numbers include $0$, and every natural number
divides $0$. For $z\in\C$, $\bar z$ is the complex conjugate (Mathlib \lean{starRingEnd}) and $\abs z$ the modulus;
Mathlib's $\mathrm{normSq}(z)=(\operatorname{Re}z)^2+(\operatorname{Im}z)^2$ equals $\abs z^2$ and is multiplicative.
We write $\langle z,w\rangle=\operatorname{Re}(z\bar w)=\operatorname{Re}z\operatorname{Re}w+\operatorname{Im}z
\operatorname{Im}w$ for the Euclidean inner product of $\C=\R^2$, so that $\abs{z}^2=\langle z,z\rangle$.

Lean names in the text are read relative to the namespace \lean{ErdosSzekeres.SymReduction}: a name $n$ stands for the
first of \lean{ErdosSzekeres.SymReduction.}$n$, \lean{ErdosSzekeres.}$n$ and $n$ that exists. So \lean{orbit_select}
is \lean{ErdosSzekeres.SymReduction.orbit_select}, \lean{Review.} and \lean{VariantChecks.} denote sub-namespaces of
\lean{ErdosSzekeres.SymReduction}, \lean{Point} is \lean{ErdosSzekeres.Point}, \lean{Suk.} denotes
\lean{ErdosSzekeres.Suk}, and the names of Mathlib and those with the prefixes \lean{Mir.} (the Lean package in
\dirn{H3}) and \lean{V2L.} (the Lean package in \dirn{V2-L}) are full names. Statement heads give full names.

\begin{definition}[Squared distance, Euclidean isometry; \leanref{ErdosSzekeres.SymReduction.eDist2}{SYMLEAN/SymmetricReduction.lean:265}, \leanref{ErdosSzekeres.SymReduction.IsEucIsometry}{SYMLEAN/SymmetricReduction.lean:268}]\label{sr:def:iso}
For $p,q\in\R^2$ put $\operatorname{eDist2}(p,q)=(p_0-q_0)^2+(p_1-q_1)^2$. A map $g\colon\R^2\to\R^2$ is a
\emph{Euclidean isometry} (Lean \texttt{IsEucIsometry g}) if
\[
  \operatorname{eDist2}(g(p),g(q))=\operatorname{eDist2}(p,q)\qquad\text{for all }p,q\in\R^2 .
\]
\end{definition}

No further property of $g$ is assumed: $g$ is not assumed continuous, surjective or affine, but it is defined on the
whole plane. That $g$, read in complex coordinates, has the form $z\mapsto az+b$ or $z\mapsto a\bar z+b$ with
$\abs a=1$ follows from Lemma~\ref{sr:lem:toC-iso}(b) and Theorem~\ref{sr:thm:structure}(a).

\begin{remark}[Why an explicit Euclidean distance]\label{sr:rem:sup}
Mathlib equips the type \lean{Point} of functions $\{0,1\}\to\R$ with the product (sup) metric
$\mathrm{dist}(p,q)=\max(\abs{p_0-q_0},\abs{p_1-q_1})$; for instance $\mathrm{dist}((0,0),(3,4))=4$. Mathlib's predicate
\lean{Isometry} for maps of \lean{Point} therefore means preservation of the sup distance, which a rotation need not
have. For example, the rotation $\rho$ by $2\pi/3$ about a point $c$ fixes $c$ and maps $c+(1,0)$, at sup distance $1$
from $c$, to $c+(-\tfrac12,\tfrac{\sqrt3}{2})$, at sup distance $\tfrac{\sqrt3}{2}$ from $c$. So $\rho$ does not preserve
the sup distance, and a theorem stated with Mathlib's \lean{Isometry} would not cover the rotations used in
\lean{Threefold33}. This is why the file states isometries with the explicit squared Euclidean distance
$\operatorname{eDist2}$. General position, convex position and \lean{HasConvexNGon} are affine notions (Mathlib
\lean{Collinear} and \lean{ConvexIndependent}) and involve no metric. Lean checks of such sup-metric facts are in the
review file \dirn{SYMLEAN/review/SupMetric.lean}, which is outside the scope of this section.
\end{remark}

\begin{definition}[Invariance; \leanref{ErdosSzekeres.SymReduction.Invariant}{SYMLEAN/SymmetricReduction.lean:271}]\label{sr:def:inv}
For $f\colon\R^2\to\R^2$ and a finite set $S$, $\operatorname{Invariant}(f,S)$ means $f(p)\in S$ for every $p\in S$.
\end{definition}

This is the notion of Definition~\ref{fd:def:inv}: $f$ maps $S$ \emph{into} $S$; equality $f(S)=S$ is not required.
(It follows when $f$ is injective, because $S$ is finite, but it is never used.)

\begin{definition}[Complex coordinates; \leanref{ErdosSzekeres.SymReduction.toC}{SYMLEAN/SymmetricReduction.lean:274}, \leanref{ErdosSzekeres.SymReduction.ofC}{SYMLEAN/SymmetricReduction.lean:277}]\label{sr:def:complex}
For $p\in\R^2$ let $\operatorname{toC}(p)=p_0+ip_1\in\C$, and for $z\in\C$ let $\operatorname{ofC}(z)=(\operatorname{Re}
z,\operatorname{Im}z)\in\R^2$. In proofs we abbreviate $\hat p=\operatorname{toC}(p)$.
\end{definition}

\begin{lemma}[Coordinates; \leanref{ErdosSzekeres.SymReduction.toC_re}{SYMLEAN/SymmetricReduction.lean:279}, \leanref{ErdosSzekeres.SymReduction.toC_im}{SYMLEAN/SymmetricReduction.lean:280}, \leanref{ErdosSzekeres.SymReduction.ofC_zero}{SYMLEAN/SymmetricReduction.lean:281}, \leanref{ErdosSzekeres.SymReduction.ofC_one}{SYMLEAN/SymmetricReduction.lean:282}, \leanref{ErdosSzekeres.SymReduction.toC_ofC}{SYMLEAN/SymmetricReduction.lean:284}, \leanref{ErdosSzekeres.SymReduction.ofC_toC}{SYMLEAN/SymmetricReduction.lean:287}, \leanref{ErdosSzekeres.SymReduction.toC_injective}{SYMLEAN/SymmetricReduction.lean:291}, \leanref{ErdosSzekeres.SymReduction.eDist2_eq_normSq}{SYMLEAN/SymmetricReduction.lean:294}, \leanref{ErdosSzekeres.SymReduction.eDist2_pos_of_ne}{SYMLEAN/SymmetricReduction.lean:297}, \leanref{ErdosSzekeres.SymReduction.point_ext}{SYMLEAN/SymmetricReduction.lean:502}]\label{sr:lem:complex}
Let $p,q\in\R^2$ and $z\in\C$.
\begin{enumerate}[label=(\alph*)]
\item (\lean{toC_re}, \lean{toC_im}, \lean{ofC_zero}, \lean{ofC_one}) $\operatorname{Re}\operatorname{toC}(p)=p_0$,
  $\operatorname{Im}\operatorname{toC}(p)=p_1$,
  $\operatorname{ofC}(z)_0=\operatorname{Re}z$ and $\operatorname{ofC}(z)_1=\operatorname{Im}z$.
\item (\lean{point_ext}) If $p_0=q_0$ and $p_1=q_1$, then $p=q$.
\item (\lean{toC_ofC}, \lean{ofC_toC}, \lean{toC_injective}) $\operatorname{toC}(\operatorname{ofC}(z))=z$ and
  $\operatorname{ofC}(\operatorname{toC}(p))=p$; the map
  $\operatorname{toC}$ is injective.
\item (\lean{eDist2_eq_normSq}) $\operatorname{eDist2}(p,q)=\mathrm{normSq}(\operatorname{toC}(p)-\operatorname{toC}(q))$.
\item (\lean{eDist2_pos_of_ne}) If $p\ne q$, then $0<\operatorname{eDist2}(p,q)$.
\end{enumerate}
\end{lemma}

\begin{proof}
(a) holds by the definitions. (b) A point is a function on $\{0,1\}$, determined by its two values. (c) By (a),
$\operatorname{toC}(\operatorname{ofC}(z))=\operatorname{Re}z+i\operatorname{Im}z=z$, and
$\operatorname{ofC}(\operatorname{toC}(p))$ has coordinates $p_0,p_1$, so it equals $p$ by (b). If
$\operatorname{toC}(p)=\operatorname{toC}(q)$, then $p=\operatorname{ofC}(\operatorname{toC}(p))=
\operatorname{ofC}(\operatorname{toC}(q))=q$. (d) $\operatorname{toC}(p)-\operatorname{toC}(q)=(p_0-q_0)+i(p_1-q_1)$,
whose $\mathrm{normSq}$ is $(p_0-q_0)^2+(p_1-q_1)^2$. (e) If $p\ne q$, then $\hat p-\hat q\ne0$ by (c), and a nonzero
complex number has positive $\mathrm{normSq}$; conclude by (d).
\end{proof}

\begin{definition}[Rotations, half-turn, line reflection; \leanref{ErdosSzekeres.SymReduction.crot}{SYMLEAN/SymmetricReduction.lean:349}, \leanref{ErdosSzekeres.SymReduction.rotAbout}{SYMLEAN/SymmetricReduction.lean:400}, \leanref{ErdosSzekeres.SymReduction.halfTurn}{SYMLEAN/SymmetricReduction.lean:413}, \leanref{ErdosSzekeres.SymReduction.reflLine}{SYMLEAN/SymmetricReduction.lean:423}]\label{sr:def:maps}
Let $p\in\R^2$.
\begin{enumerate}[label=(\alph*)]
\item For $c,w\in\C$ (Lean \texttt{crot c w}):
  $\operatorname{crot}_{c,w}(p)=\operatorname{ofC}\bigl(c+w(\operatorname{toC}(p)-c)\bigr)$.
\item For $c\in\R^2$ and $\theta\in\R$ (Lean \lean{rotAbout}; the formula of Definition~\ref{fd:def:rot}):
  \[
  \operatorname{rotAbout}_{c,\theta}(p)=\bigl(c_0+\cos\theta\,(p_0-c_0)-\sin\theta\,(p_1-c_1),\;
  c_1+\sin\theta\,(p_0-c_0)+\cos\theta\,(p_1-c_1)\bigr).
  \]
\item For $c\in\R^2$ (Lean \texttt{halfTurn c}): $\operatorname{halfTurn}_c(p)=(2c_0-p_0,\,2c_1-p_1)$.
\item For $a,u\in\R^2$ (Lean \texttt{reflLine a u}): with $d=(p_0-a_0)u_0+(p_1-a_1)u_1$,
  \[
  \operatorname{reflLine}_{a,u}(p)=\bigl(a_0+2du_0-(p_0-a_0),\;a_1+2du_1-(p_1-a_1)\bigr).
  \]
\end{enumerate}
\end{definition}

For $\abs w=1$, $\operatorname{crot}_{c,w}$ is the rotation about the point $\operatorname{ofC}(c)$ through the angle
$\theta$ with $w=e^{i\theta}$ (Lemma~\ref{sr:lem:rot-half}(a)), and $\operatorname{rotAbout}_{c,\theta}$ is the
counterclockwise rotation by $\theta$ about $c$. For $u_0^2+u_1^2=1$, $\operatorname{reflLine}_{a,u}$ is the reflection
in the line $a+\R u$: by Lemma~\ref{sr:lem:refl}(a) it is $z\mapsto\hat a+\hat u^{\,2}\,\overline{(z-\hat a)}$ in
complex coordinates, and as $\hat u\,\overline{\hat u}=1$, this map fixes $\hat a+t\hat u$ and sends $\hat a+it\hat u$
to $\hat a-it\hat u$ for every $t\in\R$. For $u=0$ it is the half-turn about $a$. For every other $u$ it is not even a
Euclidean isometry: with $s=u_0^2+u_1^2\notin\{0,1\}$, it fixes $a$ and sends $a+u$ to $a+(2s-1)u$, and
$\operatorname{eDist2}(a+(2s-1)u,a)=(2s-1)^2s\ne s=\operatorname{eDist2}(a+u,a)$. So every statement below that uses
it as a reflection assumes $u_0^2+u_1^2=1$. The definition has the same text as \lean{Mir.reflLine} in
\dirn{H3/MirrorEncoding/Line.lean:18}.

\begin{definition}[The case classes; \leanref{ErdosSzekeres.SymReduction.HalfTurn34}{SYMLEAN/SymmetricReduction.lean:467}, \leanref{ErdosSzekeres.SymReduction.Threefold33}{SYMLEAN/SymmetricReduction.lean:472}, \leanref{ErdosSzekeres.SymReduction.Fivefold35}{SYMLEAN/SymmetricReduction.lean:477}, \leanref{ErdosSzekeres.SymReduction.MirrorOneAxis33}{SYMLEAN/SymmetricReduction.lean:482}, \leanref{ErdosSzekeres.SymReduction.MirrorNoAxis34}{SYMLEAN/SymmetricReduction.lean:488}, \leanref{ErdosSzekeres.SymReduction.Fourfold32}{SYMLEAN/SymmetricReduction.lean:496}]\label{sr:def:classes}
The following six propositions are defined (Lean: definitions of type \lean{Prop}). In each, $S$ ranges over all finite
subsets of $\R^2$, and $c$, $a$, $u$ range over all points of $\R^2$.
\begin{enumerate}[label=(\arabic*)]
\item \lean{HalfTurn34}: for all $S$ and $c$, if $S$ is in general position, $\card S=34$, $c\notin S$ and
  $\operatorname{Invariant}(\operatorname{halfTurn}_c,S)$, then \texttt{HasConvexNGon 7 S}.
\item \lean{Threefold33}: for all $S$ and $c$, if $S$ is in general position, $\card S=33$, $c\notin S$ and
  $\operatorname{Invariant}(\operatorname{rotAbout}_{c,2\pi/3},S)$, then \texttt{HasConvexNGon 7 S}.
\item \lean{Fivefold35}: for all $S$ and $c$, if $S$ is in general position, $\card S=35$, $c\notin S$ and
  $\operatorname{Invariant}(\operatorname{rotAbout}_{c,2\pi/5},S)$, then \texttt{HasConvexNGon 7 S}.
\item \lean{MirrorOneAxis33}: for all $S$, $a$ and $u$, if $u_0^2+u_1^2=1$, $S$ is in general position, $\card S=33$,
  $\operatorname{Invariant}(\operatorname{reflLine}_{a,u},S)$, and there is $x\in S$ with
  $\operatorname{reflLine}_{a,u}(x)=x$ such that every $y\in S$ with $\operatorname{reflLine}_{a,u}(y)=y$ equals $x$, then
  \texttt{HasConvexNGon 7 S}.
\item \lean{MirrorNoAxis34}: for all $S$, $a$ and $u$, if $u_0^2+u_1^2=1$, $S$ is in general position, $\card S=34$,
  $\operatorname{Invariant}(\operatorname{reflLine}_{a,u},S)$, and $\operatorname{reflLine}_{a,u}(p)\ne p$ for every
  $p\in S$, then \texttt{HasConvexNGon 7 S}.
\item \lean{Fourfold32}: for all $S$ and $c$, if $S$ is in general position, $\card S=32$, $c\notin S$ and
  $\operatorname{Invariant}(\operatorname{rotAbout}_{c,2\pi/4},S)$, then \texttt{HasConvexNGon 7 S}.
\end{enumerate}
The angles are the real numbers $2\pi/3$, $2\pi/5$ and $2\pi/4$ (Lean \texttt{2 * Real.pi / 3} etc.); the rotations are
counterclockwise. In \lean{MirrorOneAxis33} the last hypothesis says that exactly one point of $S$ is fixed by the
reflection; in \lean{MirrorNoAxis34} no point of $S$ is fixed. For $u_0^2+u_1^2=1$ the fixed points of
$\operatorname{reflLine}_{a,u}$ are exactly the points of the mirror $a+\R u$. Indeed, with $d$ as in
Definition~\ref{sr:def:maps}(d), the equation $\operatorname{reflLine}_{a,u}(p)=p$ says $2du_0=2(p_0-a_0)$ and
$2du_1=2(p_1-a_1)$, so a fixed point is $p=a+du$; and for $p=a+tu$ one has $d=t(u_0^2+u_1^2)=t$ and
$\operatorname{reflLine}_{a,u}(p)=a+2tu-tu=p$.
\end{definition}

These are definitions, not assertions. No class is proved in this section. Theorem~\ref{sr:thm:main} takes the first
five as hypotheses, and Theorem~\ref{sr:thm:variant} takes \lean{Threefold33}, \lean{Fourfold32}, \lean{Fivefold35} and
the two mirror classes. Remark~\ref{sr:rem:redundant} records which hypotheses of the classes are implied by the others.

\begin{remark}[Where the project supplies the classes]\label{sr:rem:records}
None of the following is checked in this section; it follows \dirn{C2ENC/class_table.md}, \S5 and \S9 of
\dirn{SYMLEAN/README.md}, \S3 of \dirn{WRITEUP/SYMMETRIC.md} and the records named below. For no class is the step from
a SAT certificate to the Lean class Lean-checked. On the certificate routes of \lean{HalfTurn34}, \lean{Threefold33},
\lean{Fivefold35} and \lean{Fourfold32}, and for the \dirn{SYMZ} certificates of the two mirror classes, that step rests
on written proofs (the soundness of the encoders and the lemmas named below) and on the review of the class
definitions in \dirn{SYMLEAN/README.md} \S6, which compared each class with its certificate by reading. On the route
through \dirn{H3} and \dirn{V2-L} to \lean{MirrorOneAxis33}, the soundness of the encodings and the reduction of
\lean{Mir.MirrorExclusionLine} to CNF facts are Lean-checked. The cube covers in that reduction are checked by kernel
evaluation (computer-assisted), with one exception: the cover of the top two levels of the hull-pair case split of
sub-case (ii), where the point of $S$ on the mirror lies in the convex hull of the other points
(\dirn{WRITEUP/MIRROR.md} \S4). That cover is itself given by $1+121$ of those CNF facts. What remains outside Lean on
this route is described in Remark~\ref{sr:rem:mirror-line}. The
project trusts one proof checker, \texttt{cake\_lpr}: where a certificate below is said to be accepted by both
checkers, \texttt{cake\_lpr} accepted it, and \texttt{lrat-check} is only a secondary check (\dirn{WRITEUP/SYMMETRIC.md}
\S3, revised on 6 November 2026). The trust statement in \S9 of \dirn{SYMLEAN/README.md}, which this remark follows on
other points, predates that revision and still lists both checkers as trusted.
\begin{itemize}
\item \lean{HalfTurn34}: claimed in \dirn{C2FINAL} as computer-assisted. It rests on SAT certificates (\dirn{C2C34}:
  hull cover and $120$ hexagonal hulls, proofs retained; \dirn{C2P}--\dirn{C2P5}: parallelogram subcases, most of which
  are verdicts of both checkers recorded at solve time for proofs that were then deleted; re-checks in
  \dirn{C2RECHECK}); on the written soundness proof of the encoder, \dirn{C2ENC/proof.md}; on the reductions A and B,
  which pair equivalent parallelogram cases and subcases (written proofs, reviewed in \dirn{C2P/review} and checked at
  the CNF level in \dirn{AUDIT3}); and on the bookkeeping of the subcase orbits, Lean-checked according to the record
  (\dirn{C2FINAL/README.md} \S6), by kernel evaluation, and so computer-assisted (Theorem~\ref{cb:thm:c2}; that the
  listed subcases are refuted rests on the records). A table of the evidence is in \dirn{C2ENC/class_table.md} \S1.
\item \lean{Threefold33}: claimed in the records as computer-assisted. It rests on \dirn{C3FINAL}, a composition of
  $116{,}883$ pieces whose proofs were not retained, so that their verdicts are recorded ones (a random sample of
  $2{,}120$ pieces was
  re-solved in \dirn{C3SPOT}, and the $204$ indispensable $\sigma$-cube pieces were re-solved with retained proofs in
  \dirn{C3PIECES}); on the cover of the $910$ roots, Lean-checked by kernel evaluation and so computer-assisted
  (Theorem~\ref{cb:thm:c3}; that each leaf names a record with recorded checker verdicts is checked outside Lean); on
  the written normal form \dirn{C3-S/soundness.md} and the written Lemma H6M (\dirn{C3H6M/soundness_h6m.md}, checked
  adversarially in \dirn{C3AUDIT2}); on the written exhaustiveness of the layer splits and $\sigma$-splits (Lemmas 1--4
  and the cover theorem of \dirn{C3-S/soundness.md}, its version for the reduced cut in
  \dirn{C3H6M/soundness_h6m.md}, and rules 2 and 3 of \dirn{C3U2/sigma_split.md}); on the encoding \dirn{C3a}; and on
  the written centre-perturbation lemma \dirn{C2ENC/perturbation.md}. The perturbation lemma is needed because the
  class admits centres that lie on a line through two points of different orbits, which the encoder excludes.
\item \lean{Fivefold35}: claimed in the records as computer-assisted (\dirn{SYM5/README.md}). It rests on \dirn{SYM5}:
  two normal forms of one encoder (an angle form and a layer form; not independent encodings), each with a retained
  proof that both checkers accepted; either proof would suffice for the class without the other, given the written
  proofs listed next. It rests further on written proofs of the encoder's soundness and of the two normal forms, and on
  the same perturbation lemma. An earlier computation reported by B.~Dumitru (arXiv:2512.24061, \S7), which has
  priority, also implies the fivefold bound; no proof files of that run are public (\dirn{SYM5/README.md}).
\item \lean{MirrorOneAxis33} and \lean{MirrorNoAxis34}: Lean-checked in \dirn{SYMLEAN/mirror}
  (\lean{mirrorOneAxis33_proved}, \lean{mirrorNoAxis34_proved}), with no SAT input, from the repository's formalization
  of Baek's theorem (J.~Baek, arXiv:2206.04260); written proofs are in Theorem~\ref{mb:thm:classes}. The file
  \dirn{SYMLEAN/mirror/ReductionThreeClasses.lean} does not import \dirn{SYMLEAN/SymmetricReduction.lean}. Its lines
  38--1084 are byte-identical to lines 12--1058 of \dirn{SYMLEAN/SymmetricReduction.lean}, so the classes defined in
  \dirn{SYMLEAN/mirror/ReductionThreeClasses.lean} are copies of those of Definition~\ref{sr:def:classes}, with the same
  full names and the same text. Both classes are also supplied by SAT certificates in \dirn{SYMZ} (Theorems 2a and 2b
  of \dirn{WRITEUP/SYMMETRIC.md}, claimed in the records as computer-assisted), and \lean{MirrorOneAxis33} also,
  conditionally, by the Lean chain of \dirn{H3} and \dirn{V2-L} (Remark~\ref{sr:rem:mirror-line}).
\item \lean{Fourfold32} (among the statements of this section, a hypothesis only of Theorem~\ref{sr:thm:variant} and
  of Lemma~\ref{sr:lem:fourfold}, which its proof uses): claimed in \dirn{C4CUT}, which states it without a
  rigor label. It rests on a final certificate accepted by both checkers, on six cuts checked by exact computation and
  on a written proof of the encoder's soundness. The statement is due to Subercaseaux, Mackey, Qian and Heule
  (arXiv:2506.00224, \S6.2).
\end{itemize}
\end{remark}

\subsection{Finite dynamics}\label{sr:ssec:dyn}

\begin{lemma}[Iterates stay in an invariant set; \leanref{ErdosSzekeres.SymReduction.iterate_mem}{SYMLEAN/SymmetricReduction.lean:69}, \leanref{ErdosSzekeres.SymReduction.Invariant.iterate}{SYMLEAN/SymmetricReduction.lean:309}]\label{sr:lem:iterate}
Iterates of a map stay in a set that the map sends into itself.
\begin{enumerate}[label=(\alph*)]
\item (\lean{iterate_mem}) Let $X$ be a set, $f\colon X\to X$, and $S\subseteq X$ finite with $f(y)\in S$ for every
  $y\in S$. If $x\in S$, then $f^{[i]}(x)\in S$ for every $i\in\N$.
\item (\lean{Invariant.iterate}) If $f\colon\R^2\to\R^2$, $S\subseteq\R^2$ is finite and $\operatorname{Invariant}(f,S)$,
  then $\operatorname{Invariant}(f^{[k]},S)$ for every $k\in\N$.
\end{enumerate}
\end{lemma}

\begin{proof}
(a) Induction on $i$: $f^{[0]}(x)=x\in S$, and $f^{[i+1]}(x)=f(f^{[i]}(x))\in S$ because $f^{[i]}(x)\in S$. (b) is (a)
applied to each $x\in S$.
\end{proof}

\begin{lemma}[Pigeonhole for iterates; written step]\label{sr:lem:pigeon}
Let $S\subseteq\R^2$ be finite, $f\colon\R^2\to\R^2$ with $\operatorname{Invariant}(f,S)$, and $q\in S$. Then there are
$i<j$ in $\N$ with $f^{[i]}(q)=f^{[j]}(q)$.
\end{lemma}

\begin{proof}
By Lemma~\ref{sr:lem:iterate}, the $\card S+1$ points $f^{[k]}(q)$, $0\le k\le\card S$, lie in $S$. By the pigeonhole
principle two of them with different indices coincide; call the smaller index $i$ and the larger $j$. (In Lean this is
the local statement \lean{pigeon} in the proof of \lean{isometry_structure}, obtained from Mathlib's
\lean{Set.Finite.exists_lt_map_eq_of_forall_mem}.)
\end{proof}

\begin{definition}[Minimal period; written step]\label{sr:def:minper}
Let $X$ be a set, $f\colon X\to X$ and $x\in X$. If $f^{[n]}(x)=x$ for some $n>0$, then $\operatorname{minimalPeriod}(f,x)$
is the least such $n$; otherwise $\operatorname{minimalPeriod}(f,x)=0$. (This is Mathlib's
\lean{Function.minimalPeriod}.)
\end{definition}

\begin{lemma}[Minimal periods; written step]\label{sr:lem:minper}
Let $f\colon X\to X$, $x\in X$ and $d=\operatorname{minimalPeriod}(f,x)$.
\begin{enumerate}[label=(MP\arabic*)]
\item If $f^{[n]}(x)=x$ for some $n\in\N$, then $d\mid n$. (Mathlib \lean{Function.IsPeriodicPt.minimalPeriod_dvd}.)
\item $d=1$ if and only if $f(x)=x$. (Mathlib \lean{Function.minimalPeriod_eq_one_iff_isFixedPt}.)
\item $f^{[d]}(x)=x$. (Mathlib \lean{Function.iterate_minimalPeriod}.)
\item If $i<j<d$, then $f^{[i]}(x)\ne f^{[j]}(x)$. (Mathlib \lean{Function.iterate_injOn_Iio_minimalPeriod}.)
\end{enumerate}
\end{lemma}

\begin{proof}
(MP3) If some $n>0$ has $f^{[n]}(x)=x$, this holds for the least one, $d$. Otherwise $d=0$ and $f^{[0]}(x)=x$.

(MP1) If $n=0$ there is nothing to prove. Let $n>0$. Then $d>0$ is the least positive $m$ with $f^{[m]}(x)=x$. By
induction on $t$, $f^{[td]}(x)=x$ for all $t\in\N$. Write $n=td+r$ with $0\le r<d$. Then
$x=f^{[n]}(x)=f^{[r]}(f^{[td]}(x))=f^{[r]}(x)$. If $r>0$, this contradicts the minimality of $d$; so $r=0$ and $d\mid n$.

(MP2) If $f(x)=x$, then $1$ is the least positive integer and $f^{[1]}(x)=x$, so $d=1$. Conversely, if $d=1\ne0$, then by
definition $f^{[1]}(x)=x$.

(MP4) Suppose $f^{[i]}(x)=f^{[j]}(x)$ with $i<j<d$. Then $d>0$, and applying $f^{[d-j]}$ to both sides gives
$f^{[d-j+i]}(x)=f^{[d]}(x)=x$ by (MP3). Since $j<d$ and $i<j$, we have $0<d-j+i<d$, which contradicts the minimality
of $d$.
\end{proof}

\begin{definition}[Forward orbit segment; \leanref{ErdosSzekeres.SymReduction.orbitFS}{SYMLEAN/SymmetricReduction.lean:65}]\label{sr:def:orbitFS}
For $f\colon X\to X$, $p\in\N$ and $x\in X$, let $\operatorname{orbitFS}(f,p,x)=\{f^{[i]}(x):0\le i<p\}$, a finite set
(Lean: the image of $\{0,\dots,p-1\}$ under $i\mapsto f^{[i]}(x)$; Lean needs decidable equality on $X$ to form it).
\end{definition}

\begin{lemma}[Orbit segments; \leanref{ErdosSzekeres.SymReduction.card_orbitFS_of_minimalPeriod}{SYMLEAN/SymmetricReduction.lean:111}, \leanref{ErdosSzekeres.SymReduction.mem_orbitFS_self}{SYMLEAN/SymmetricReduction.lean:90}, \leanref{ErdosSzekeres.SymReduction.orbitFS_subset}{SYMLEAN/SymmetricReduction.lean:93}, \leanref{ErdosSzekeres.SymReduction.orbitFS_invariant}{SYMLEAN/SymmetricReduction.lean:99}, \leanref{ErdosSzekeres.SymReduction.card_orbitFS}{SYMLEAN/SymmetricReduction.lean:76}]\label{sr:lem:orbitFS}
Let $X$ be a set, $f\colon X\to X$, $p\in\N$ and $x\in X$, and write $O=\operatorname{orbitFS}(f,p,x)$.
\begin{enumerate}[label=(\alph*)]
\item (\lean{mem_orbitFS_self}) If $0<p$, then $x\in O$.
\item (\lean{orbitFS_subset}) If $S\subseteq X$ is finite, $f(y)\in S$ for every $y\in S$, and $x\in S$, then
  $O\subseteq S$.
\item (\lean{orbitFS_invariant}) If $0<p$ and $f^{[p]}(x)=x$, then $f(y)\in O$ for every $y\in O$.
\item (\lean{card_orbitFS}) If $p$ is prime, $f^{[p]}(x)=x$ and $f(x)\ne x$, then $\card O=p$.
\item (\lean{card_orbitFS_of_minimalPeriod}) If $\operatorname{minimalPeriod}(f,x)=p$, then $\card O=p$.
\end{enumerate}
\end{lemma}

\begin{proof}
(a) $x=f^{[0]}(x)$ and $0<p$. (b) Each $f^{[i]}(x)$ lies in $S$ by Lemma~\ref{sr:lem:iterate}(a). (c) Let
$y=f^{[i]}(x)$ with $i<p$. Then $f(y)=f^{[i+1]}(x)$. If $i+1<p$, this lies in $O$. Otherwise $i+1=p$ and
$f(y)=f^{[p]}(x)=x\in O$ by (a). (e) By (MP4), $i\mapsto f^{[i]}(x)$ is injective on $\{0,\dots,p-1\}$, so $O$ has $p$
elements. (This also holds for $p=0$, when $O=\varnothing$.) (d) By (MP1), $d=\operatorname{minimalPeriod}(f,x)$ divides
the prime $p$, so $d=1$ or $d=p$. By (MP2), $d\ne1$ because $f(x)\ne x$. Hence $d=p$, and (e) applies.
\end{proof}

The next statement contains the induction shared by the two orbit selection lemmas. The Lean file carries out the same
induction twice, once in each lemma; we do it once.

\begin{lemma}[Orbit selection, common form; written step]\label{sr:lem:orbit-core}
Let $X$ be a set, $f\colon X\to X$ and $p\in\N$ with $0<p$. Let $S\subseteq X$ be finite, and assume: $f(y)\in S$ for
every $y\in S$; $f$ is injective on $S$; and every $y\in S$ satisfies $f^{[p]}(y)=y$ and
$\card{\operatorname{orbitFS}(f,p,y)}=p$. Then $p\mid\card S$, and for every $k\in\N$ with $kp\le\card S$ there is
$T\subseteq S$ with $f(y)\in T$ for every $y\in T$ and $\card T=kp$.
\end{lemma}

\begin{proof}
Strong induction on $\card S$. The hypotheses on single points pass to subsets of $S$, and so does injectivity.

If $S=\varnothing$, then $p\mid0$; and $kp\le0$ forces $kp=0$, so $T=\varnothing$ works.

Let $S\ne\varnothing$ and pick $x\in S$. Put $O=\operatorname{orbitFS}(f,p,x)$. By Lemma~\ref{sr:lem:orbitFS} and the
hypotheses, $O\subseteq S$, $\card O=p$, $x\in O$, and $f(y)\in O$ for every $y\in O$. Let $S'=S\setminus O$; then
$\card{S'}=\card S-p<\card S$.

\emph{$S'$ is mapped into itself.} Let $y\in S'$. Then $f(y)\in S$. Suppose $f(y)\in O$, say $f(y)=f^{[i]}(x)$ with
$i<p$. If $i=0$, then $f(y)=x=f^{[p]}(x)=f(f^{[p-1]}(x))$ (here $p\ge1$). Both $y$ and $f^{[p-1]}(x)$ lie in $S$
(Lemma~\ref{sr:lem:iterate}(a)), so injectivity on $S$ gives $y=f^{[p-1]}(x)\in O$. If $i\ge1$, then
$f(y)=f(f^{[i-1]}(x))$ and injectivity gives $y=f^{[i-1]}(x)\in O$. Both contradict $y\notin O$. Hence $f(y)\in S'$.

So $S'$ satisfies all the hypotheses, and by induction $p\mid\card{S'}$, and $S'$ has $f$-invariant subsets of every
size $k'p\le\card{S'}$. Since $\card S=\card{S'}+p$, we get $p\mid\card S$. Let $kp\le\card S$. If $k=0$, take
$T=\varnothing$. If $k=k'+1$, then $k'p\le\card S-p=\card{S'}$. Take an $f$-invariant $T'\subseteq S'$ with
$\card{T'}=k'p$, and put $T=T'\cup O\subseteq S$. The union of two sets that $f$ maps into themselves is mapped into
itself, and $T'\cap O=\varnothing$ because $T'\subseteq S\setminus O$. So $\card T=k'p+p=kp$.
\end{proof}

\begin{lemma}[Orbit selection; \leanref{ErdosSzekeres.SymReduction.orbit_select_minimalPeriod}{SYMLEAN/SymmetricReduction.lean:123}, \leanref{ErdosSzekeres.SymReduction.orbit_select}{SYMLEAN/SymmetricReduction.lean:194}]\label{sr:lem:orbit}
Let $X$ be a set and $f\colon X\to X$.
\begin{enumerate}[label=(\alph*)]
\item (\lean{orbit_select_minimalPeriod}) Let $p\in\N$ with $0<p$. For every finite $S\subseteq X$: if $f(y)\in S$ for
  every $y\in S$, $f$ is injective on $S$, and $\operatorname{minimalPeriod}(f,y)=p$ for every $y\in S$, then
  $p\mid\card S$, and for every $k\in\N$ with $kp\le\card S$ there is $T\subseteq S$ with $f(y)\in T$ for every $y\in
  T$ and $\card T=kp$.
\item (\lean{orbit_select}) Let $p$ be a prime. For every finite $S\subseteq X$: if $f(y)\in S$ for every $y\in S$,
  $f$ is injective on $S$, $f^{[p]}(y)=y$ for every $y\in S$, and $f(y)\ne y$ for every $y\in S$, then $p\mid\card S$,
  and for every $k\in\N$ with $kp\le\card S$ there is $T\subseteq S$ with $f(y)\in T$ for every $y\in T$ and $\card
  T=kp$.
\end{enumerate}
\end{lemma}

\begin{proof}
Both follow from Lemma~\ref{sr:lem:orbit-core}; we check its hypotheses. (a) For $y\in S$, $f^{[p]}(y)=y$ by (MP3) and
$\card{\operatorname{orbitFS}(f,p,y)}=p$ by Lemma~\ref{sr:lem:orbitFS}(e). (b) $p>0$ because $p$ is prime, and
$\card{\operatorname{orbitFS}(f,p,y)}=p$ for $y\in S$ by Lemma~\ref{sr:lem:orbitFS}(d).
\end{proof}

In (b) the injectivity hypothesis is redundant: if $f(y)=f(y')$ with $y,y'\in S$, then applying $f^{[p-1]}$ gives
$y=f^{[p]}(y)=f^{[p]}(y')=y'$. It is part of the Lean statement, and it is available in all applications.

\subsection{Rotations and reflections in coordinates}\label{sr:ssec:maps}

\begin{lemma}[The maps $\operatorname{crot}$; \leanref{ErdosSzekeres.SymReduction.toC_crot}{SYMLEAN/SymmetricReduction.lean:351}, \leanref{ErdosSzekeres.SymReduction.crot_one}{SYMLEAN/SymmetricReduction.lean:354}, \leanref{ErdosSzekeres.SymReduction.crot_comp}{SYMLEAN/SymmetricReduction.lean:359}, \leanref{ErdosSzekeres.SymReduction.crot_iterate}{SYMLEAN/SymmetricReduction.lean:364}, \leanref{ErdosSzekeres.SymReduction.crot_eq_self_iff}{SYMLEAN/SymmetricReduction.lean:370}, \leanref{ErdosSzekeres.SymReduction.crot_injective}{SYMLEAN/SymmetricReduction.lean:384}, \leanref{ErdosSzekeres.SymReduction.eDist2_crot_center}{SYMLEAN/SymmetricReduction.lean:394}]\label{sr:lem:crot}
Let $c,w,v\in\C$ and $p\in\R^2$.
\begin{enumerate}[label=(\alph*)]
\item (\lean{toC_crot}) $\operatorname{toC}(\operatorname{crot}_{c,w}(p))=c+w(\operatorname{toC}(p)-c)$.
\item (\lean{crot_one}) $\operatorname{crot}_{c,1}=\mathrm{id}$.
\item (\lean{crot_comp}) $\operatorname{crot}_{c,w}(\operatorname{crot}_{c,v}(p))=\operatorname{crot}_{c,wv}(p)$.
\item (\lean{crot_iterate}) $(\operatorname{crot}_{c,w})^{[k]}=\operatorname{crot}_{c,w^k}$ for every $k\in\N$.
\item (\lean{crot_eq_self_iff}) If $w\ne1$, then $\operatorname{crot}_{c,w}(p)=p$ if and only if
  $\operatorname{toC}(p)=c$.
\item (\lean{crot_injective}) If $w\ne0$, then $\operatorname{crot}_{c,w}$ is injective.
\item (\lean{eDist2_crot_center}) If $\mathrm{normSq}(w)=1$, then
  $\operatorname{eDist2}(\operatorname{crot}_{c,w}(p),\operatorname{ofC}(c))=
  \operatorname{eDist2}(p,\operatorname{ofC}(c))$.
\end{enumerate}
\end{lemma}

\begin{proof}
Write $z=\hat p$. Equalities of points are proved by applying the injective map $\operatorname{toC}$
(Lemma~\ref{sr:lem:complex}(c)). (a) is $\operatorname{toC}(\operatorname{ofC}(\zeta))=\zeta$. (b) $c+1\cdot(z-c)=z$.
(c) By (a), $c+w\bigl((c+v(z-c))-c\bigr)=c+wv(z-c)$. (d) Induction on $k$. For $k=0$ both sides are the identity, by
(b). If it holds for $k$, then $(\operatorname{crot}_{c,w})^{[k+1]}(p)=\operatorname{crot}_{c,w}(\operatorname{crot}_{c,w^k}(p))
=\operatorname{crot}_{c,w\cdot w^k}(p)=\operatorname{crot}_{c,w^{k+1}}(p)$ by (c). (e) By (a),
$\operatorname{toC}(\operatorname{crot}_{c,w}(p))-z=(w-1)(z-c)$. For $w\ne1$ this vanishes if and only if $z=c$. (f) If
$\operatorname{crot}_{c,w}(p)=\operatorname{crot}_{c,w}(q)$, then by (a) $w(\hat p-\hat q)=0$, so $\hat p=\hat q$ and
$p=q$. (g) By Lemma~\ref{sr:lem:complex}(c),(d) and (a), the left side is $\mathrm{normSq}(c+w(z-c)-c)=
\mathrm{normSq}(w)\,\mathrm{normSq}(z-c)=\mathrm{normSq}(z-c)$, which is the right side.
\end{proof}

\begin{lemma}[Rotations and the half-turn as $\operatorname{crot}$; \leanref{ErdosSzekeres.SymReduction.rotAbout_eq_crot}{SYMLEAN/SymmetricReduction.lean:404}, \leanref{ErdosSzekeres.SymReduction.halfTurn_eq_crot}{SYMLEAN/SymmetricReduction.lean:415}]\label{sr:lem:rot-half}
For $c\in\R^2$ and $\theta\in\R$:
\begin{enumerate}[label=(\alph*)]
\item (\lean{rotAbout_eq_crot}) $\operatorname{rotAbout}_{c,\theta}=\operatorname{crot}_{\operatorname{toC}(c),\,\exp(\theta i)}$ (equality of maps
  of $\R^2$);
\item (\lean{halfTurn_eq_crot}) $\operatorname{halfTurn}_c=\operatorname{crot}_{\operatorname{toC}(c),\,-1}$.
\end{enumerate}
\end{lemma}

\begin{proof}
Let $p\in\R^2$, $x=p_0-c_0$ and $y=p_1-c_1$, so that $\hat p-\hat c=x+iy$. (a) By Euler's formula
$\exp(\theta i)=\cos\theta+i\sin\theta$, hence $\exp(\theta i)(x+iy)=(x\cos\theta-y\sin\theta)+i(x\sin\theta+y\cos\theta)$.
Adding $\hat c=c_0+ic_1$ gives $\operatorname{toC}(\operatorname{rotAbout}_{c,\theta}(p))$; conclude by
Lemma~\ref{sr:lem:crot}(a) and the injectivity of $\operatorname{toC}$. (b) $\hat c+(-1)(\hat p-\hat c)=2\hat c-\hat p=
\operatorname{toC}(\operatorname{halfTurn}_c(p))$.
\end{proof}

\begin{lemma}[Line reflections; \leanref{ErdosSzekeres.SymReduction.toC_reflLine}{SYMLEAN/SymmetricReduction.lean:427}, \leanref{ErdosSzekeres.SymReduction.reflLine_involutive}{SYMLEAN/SymmetricReduction.lean:435}, \leanref{ErdosSzekeres.SymReduction.cross_of_reflLine_fixed}{SYMLEAN/SymmetricReduction.lean:447}, \leanref{ErdosSzekeres.SymReduction.orient_eq_zero_of_cross}{SYMLEAN/SymmetricReduction.lean:455}]\label{sr:lem:refl}
Let $a,u,p,q,r\in\R^2$.
\begin{enumerate}[label=(\alph*)]
\item (\lean{toC_reflLine}) If $u_0^2+u_1^2=1$, then $\operatorname{toC}(\operatorname{reflLine}_{a,u}(p))=\hat a+\hat u^{\,2}\,
  \overline{(\hat p-\hat a)}$.
\item (\lean{reflLine_involutive}) If $u_0^2+u_1^2=1$, then
  $\operatorname{reflLine}_{a,u}(\operatorname{reflLine}_{a,u}(p))=p$.
\item (\lean{cross_of_reflLine_fixed}) If $\operatorname{reflLine}_{a,u}(p)=p$, then $(p_0-a_0)u_1-(p_1-a_1)u_0=0$.
  (No hypothesis on $u$.)
\item (\lean{orient_eq_zero_of_cross}) If $u_0^2+u_1^2=1$ and $(z_0-a_0)u_1-(z_1-a_1)u_0=0$ for each
  $z\in\{p,q,r\}$, then $\orient(p,q,r)=0$.
\end{enumerate}
\end{lemma}

\begin{proof}
Let $x=p_0-a_0$, $y=p_1-a_1$ and $d=xu_0+yu_1$, so that $\operatorname{reflLine}_{a,u}(p)-a=(2du_0-x,\,2du_1-y)$.

(a) Expanding, $\hat u^{\,2}\,\overline{(x+iy)}=(u_0^2-u_1^2+2iu_0u_1)(x-iy)=\bigl((u_0^2-u_1^2)x+2u_0u_1y\bigr)+
i\bigl(2u_0u_1x-(u_0^2-u_1^2)y\bigr)$. On the other hand $2du_0-x=(2u_0^2-1)x+2u_0u_1y$ and
$2du_1-y=2u_0u_1x+(2u_1^2-1)y$. Since $u_0^2+u_1^2=1$, we have $2u_0^2-1=u_0^2-u_1^2$ and $2u_1^2-1=-(u_0^2-u_1^2)$, so the
two expressions agree. (In other words, the real and imaginary parts of the difference are $x(u_0^2+u_1^2-1)$ and
$y(u_0^2+u_1^2-1)$.)

(b) Put $\zeta=\hat p-\hat a$; note $\hat u\,\overline{\hat u}=u_0^2+u_1^2=1$. By (a), applied twice,
$\operatorname{toC}(\operatorname{reflLine}_{a,u}(\operatorname{reflLine}_{a,u}(p)))-\hat a=\hat u^{\,2}\,
\overline{\hat u^{\,2}\,\bar\zeta}=(\hat u\,\overline{\hat u})^2\zeta=\zeta$. Hence the two points have the same
$\operatorname{toC}$, and they are equal.

(c) The fixed-point equations read $a_0+2du_0-x=p_0=a_0+x$ and $a_1+2du_1-y=a_1+y$, that is, $du_0=x$ and $du_1=y$. Hence
$xu_1-yu_0=du_0u_1-du_1u_0=0$.

(d) Write $v\times w=v_0w_1-v_1w_0$ and $v\cdot w=v_0w_0+v_1w_1$ for $v,w\in\R^2$, so that
$\orient(p,q,r)=(q-p)\times(r-p)$. For all $v,w,u\in\R^2$ one has the identity
\[
  (v\times w)(u_0^2+u_1^2)=(w\cdot u)(v\times u)-(v\cdot u)(w\times u),
\]
checked by expanding both sides. Take $v=q-p$ and $w=r-p$. The hypotheses say $(z-a)\times u=0$ for $z=p,q,r$, so
$v\times u=(q-a)\times u-(p-a)\times u=0$, and likewise $w\times u=0$. Hence $\orient(p,q,r)\cdot1=0$.
\end{proof}

\subsection{Isometries of the plane}\label{sr:ssec:iso}

\begin{lemma}[Isometries in coordinates; \leanref{ErdosSzekeres.SymReduction.IsEucIsometry.injective}{SYMLEAN/SymmetricReduction.lean:301}, \leanref{ErdosSzekeres.SymReduction.isometry_toC_conj}{SYMLEAN/SymmetricReduction.lean:313}]\label{sr:lem:toC-iso}
Let $g\colon\R^2\to\R^2$ be a Euclidean isometry.
\begin{enumerate}[label=(\alph*)]
\item (\lean{IsEucIsometry.injective}) $g$ is injective.
\item (\lean{isometry_toC_conj}) The map $G\colon\C\to\C$, $G(z)=\operatorname{toC}(g(\operatorname{ofC}(z)))$, is
  an isometry of $\C$:
  $\abs{G(z)-G(w)}=\abs{z-w}$ for all $z,w\in\C$. (Lean: Mathlib's \texttt{Isometry G}; for the metric space $\C$, whose
  distance is $\abs{z-w}$, this is equivalent to the displayed equation.)
\end{enumerate}
\end{lemma}

\begin{proof}
(a) If $g(p)=g(q)$, then $\operatorname{eDist2}(p,q)=\operatorname{eDist2}(g(p),g(q))=0$, so $p=q$ by
Lemma~\ref{sr:lem:complex}(e). (b) By Lemma~\ref{sr:lem:complex}(c),(d),
\[
\abs{G(z)-G(w)}^2=\operatorname{eDist2}\bigl(g(\operatorname{ofC}z),g(\operatorname{ofC}w)\bigr)=
\operatorname{eDist2}(\operatorname{ofC}z,\operatorname{ofC}w)=\abs{z-w}^2,
\]
and both moduli are nonnegative.
\end{proof}

\begin{theorem}[Isometries of the plane; \leanref{ErdosSzekeres.SymReduction.isometry_complex_cases}{SYMLEAN/SymmetricReduction.lean:321}, \leanref{ErdosSzekeres.SymReduction.isometry_structure}{SYMLEAN/SymmetricReduction.lean:869}]\label{sr:thm:structure}
Part (a) describes the isometries of $\C$; part (b) the Euclidean isometries of the plane that map a finite set with at
least two points into itself.
\begin{enumerate}[label=(\alph*)]
\item (\lean{isometry_complex_cases}) Let $G\colon\C\to\C$ be an isometry: $\abs{G(z)-G(w)}=\abs{z-w}$ for all
  $z,w\in\C$. Then there are $a,b\in\C$ with $\abs a=1$ such that either $G(z)=az+b$ for all $z\in\C$, or
  $G(z)=a\bar z+b$ for all $z\in\C$.
\item (\lean{isometry_structure}) Let $S\subseteq\R^2$ be finite with $1<\card S$, and let $g\colon\R^2\to\R^2$ be a
  Euclidean isometry with $g\ne\mathrm{id}$ (as maps of $\R^2$) and $\operatorname{Invariant}(g,S)$. Then one of the
  following holds:
  \begin{enumerate}[label=(\roman*)]
  \item there are $c,a\in\C$ with $g=\operatorname{crot}_{c,a}$, $a\ne1$, and $a^m=1$ for some $m\in\N$ with $0<m$;
  \item there are $A,u\in\R^2$ with $u_0^2+u_1^2=1$ and $g=\operatorname{reflLine}_{A,u}$.
  \end{enumerate}
\end{enumerate}
\end{theorem}

\begin{proof}
\emph{(a).} This proof differs from the Lean proof. The Lean proof applies Mathlib's construction of an affine isometry
from an isometry between real normed affine spaces with strictly convex codomain
(\lean{Isometry.affineIsometryOfStrictConvexSpace}, a definition whose underlying map is the given isometry; it is a
result of Mazur--Ulam type that does not need surjectivity), and then Mathlib's theorem describing the real-linear
isometric equivalences of $\C$ (\lean{linear_isometry_complex}). We give a self-contained argument.

Let $b=G(0)$ and $H(z)=G(z)-b$. Then $H(0)=0$, $\abs{H(z)-H(w)}=\abs{z-w}$, and so $\abs{H(z)}=\abs{H(z)-H(0)}=\abs z$.
For all $z,w$ we have $\abs{z-w}^2=\abs z^2-2\langle z,w\rangle+\abs w^2$, and the same for $H(z),H(w)$. Hence
\[
2\langle H(z),H(w)\rangle=\abs{H(z)}^2+\abs{H(w)}^2-\abs{H(z)-H(w)}^2=\abs z^2+\abs w^2-\abs{z-w}^2=2\langle z,w\rangle,
\]
so $H$ preserves the inner product. Put $a=H(1)$ and $e=H(i)$. Then $\abs a=\abs1=1$, $\abs e=1$ and $\langle
e,a\rangle=\langle i,1\rangle=0$. Let $v=e\bar a$. Then $\abs v=1$ and $\operatorname{Re}v=\langle e,a\rangle=0$, so
$v=i$ or $v=-i$; and $e=e\bar a a=va$, so $e=ia$ or $e=-ia$.

Next, every $\xi\in\C$ satisfies $\xi=\langle\xi,a\rangle a+\langle\xi,e\rangle e$. Indeed, write $\bar a\xi=\alpha+i\beta$
with $\alpha,\beta\in\R$; then $\xi=a\bar a\xi=\alpha a+\beta(ia)$, $\langle\xi,a\rangle=\operatorname{Re}(\xi\bar a)=\alpha$
and $\langle\xi,ia\rangle=\operatorname{Re}(-i\xi\bar a)=\beta$; and $\langle\xi,e\rangle e=\langle\xi,ia\rangle\,ia$ for
both choices $e=\pm ia$.

Now let $z=x+iy$ with $x,y\in\R$. Since $H$ preserves the inner product, $\langle H(z),a\rangle=\langle
H(z),H(1)\rangle=\langle z,1\rangle=x$ and $\langle H(z),e\rangle=\langle z,i\rangle=y$. Hence $H(z)=xa+ye$. If $e=ia$,
this is $a(x+iy)=az$; if $e=-ia$, it is $a(x-iy)=a\bar z$. The choice between the two cases does not depend on $z$. So
$G(z)=az+b$ for all $z$, or $G(z)=a\bar z+b$ for all $z$, and $\abs a=1$.

\emph{(b).} This part follows the Lean proof. Let $G(z)=\operatorname{toC}(g(\operatorname{ofC}z))$. By
Lemma~\ref{sr:lem:toC-iso}(b) and part (a) there are $a,b\in\C$ with $\abs a=1$ such that either
\textup{(R)} $G(z)=az+b$ for all $z$, or \textup{(C)} $G(z)=a\bar z+b$ for all $z$. Since $\operatorname{ofC}(\hat
p)=p$, we have $\operatorname{toC}(g(p))=G(\hat p)$ for every $p\in\R^2$. The set $S$ is nonempty because $\card S>1$.

\emph{Case (R), $a=1$.} Then $\operatorname{toC}(g(p))=\hat p+b$ for all $p$, and by induction
$\operatorname{toC}(g^{[k]}(q))=\hat q+kb$ for all $k\in\N$ and $q\in\R^2$. Take $q\in S$. By Lemma~\ref{sr:lem:pigeon}
there are $i<j$ with $g^{[i]}(q)=g^{[j]}(q)$, so $\hat q+ib=\hat q+jb$ and $(j-i)b=0$. Since $j-i$ is a nonzero
integer, $b=0$. Then $\operatorname{toC}(g(p))=\hat p$ for all $p$, so $g=\mathrm{id}$ by the injectivity of
$\operatorname{toC}$. This contradicts $g\ne\mathrm{id}$; so this case does not occur.

\emph{Case (R), $a\ne1$.} Put $c=b/(1-a)$. For every $p$, $c+a(\hat p-c)=a\hat p+(1-a)c=a\hat p+b=
\operatorname{toC}(g(p))$, so $g=\operatorname{crot}_{c,a}$ by Lemma~\ref{sr:lem:crot}(a). Also $a\ne0$ because
$\abs a=1$. Since $\card S\ge2$, some $q\in S$ differs from $\operatorname{ofC}(c)$. By Lemma~\ref{sr:lem:pigeon},
applied to $g=\operatorname{crot}_{c,a}$, there are $i<j$ with $(\operatorname{crot}_{c,a})^{[i]}(q)=
(\operatorname{crot}_{c,a})^{[j]}(q)$. By Lemma~\ref{sr:lem:crot}(d),(a), $c+a^i(\hat q-c)=c+a^j(\hat q-c)$. Here
$\hat q-c\ne0$, since otherwise $q=\operatorname{ofC}(\hat q)=\operatorname{ofC}(c)$. Hence $a^i=a^j=a^ia^{j-i}$, and as
$a^i\ne0$, $a^{j-i}=1$. With $m=j-i>0$ we obtain alternative (i).

\emph{Case (C).} We prove alternative (ii); the hypothesis $g\ne\mathrm{id}$ is not needed here. Since $\abs a=1$,
$a\bar a=1$. For every $p$,
\[
\operatorname{toC}(g(g(p)))=a\,\overline{(a\bar{\hat p}+b)}+b=a\bar a\hat p+a\bar b+b=\hat p+t,\qquad t:=a\bar b+b .
\]
By induction, $\operatorname{toC}((g\circ g)^{[k]}(p))=\hat p+kt$ for all $k$ and $p$. The map $g\circ g$ maps $S$ into
$S$. Take $q\in S$; Lemma~\ref{sr:lem:pigeon} applied to $g\circ g$ gives $i<j$ with $(j-i)t=0$, hence $t=0$, that is,
$a\bar b=-b$. Choose $u\in\C$ with $u^2=a$ (every complex number has a square root; the Lean proof takes the principal
value $a^{1/2}$). Then $\abs u^2=\abs a=1$, so the point $U=\operatorname{ofC}(u)$ satisfies $U_0^2+U_1^2=1$. Let
$A=\operatorname{ofC}(b/2)$. By Lemma~\ref{sr:lem:refl}(a), for every $p$,
\[
\operatorname{toC}(\operatorname{reflLine}_{A,U}(p))=\tfrac b2+u^2\,\overline{\bigl(\hat p-\tfrac b2\bigr)}
=a\bar{\hat p}+\tfrac b2-\tfrac{a\bar b}{2}=a\bar{\hat p}+\tfrac b2+\tfrac b2=a\bar{\hat p}+b=\operatorname{toC}(g(p)).
\]
Hence $g=\operatorname{reflLine}_{A,U}$, which is alternative (ii).
\end{proof}

In (b), the pigeonhole steps are where finiteness enters: they exclude translations ($a=1$) and glide reflections
($t\ne0$), and they make $a$ a root of unity.

\subsection{Elementary geometric lemmas}\label{sr:ssec:geom}

\begin{lemma}[Strict supporting lines; \leanref{ErdosSzekeres.Suk.convexPosition_of_strict_supports}{ErdosSzekeres/Suk/SupportingLines.lean:28}]\label{sr:lem:supports}
Let $s\subseteq\R^2$ be finite. If for every $p\in s$ there is $q\in\R^2$ with $\orient(p,q,r)>0$ for every $r\in s$
with $r\ne p$, then $s$ is in convex position.
\end{lemma}

\begin{proof}
This is the lemma titled ``Strict supporting lines'' in the foundations section, the same declaration of the main
library; the short proof is repeated here because Lemma~\ref{sr:lem:concyclic} uses it. By Lemma~\ref{fd:lem:cp-hull} it
suffices to show $p\notin\conv(s\setminus\{p\})$ for every $p\in s$. Fix $p$ and a point $q$ as in the hypothesis, and
let $\varphi(x)=\orient(p,q,x)=(q_0-p_0)(x_1-p_1)-(q_1-p_1)(x_0-p_0)$. This is an affine function of $x$:
$\varphi(\lambda x+\mu y)=\lambda\varphi(x)+\mu\varphi(y)$ whenever $\lambda+\mu=1$. So the open half-plane
$H=\{x:\varphi(x)>0\}$ is convex: for $x,y\in H$ and $\lambda,\mu\ge0$ with $\lambda+\mu=1$, both terms
$\lambda\varphi(x)$ and $\mu\varphi(y)$ are $\ge0$ and at least one is $>0$. By hypothesis $s\setminus\{p\}\subseteq
H$, hence $\conv(s\setminus\{p\})\subseteq H$. But $\varphi(p)=\orient(p,q,p)=0$, so $p\notin H$, and therefore
$p\notin\conv(s\setminus\{p\})$.
\end{proof}

\begin{lemma}[Concyclic points; \leanref{ErdosSzekeres.SymReduction.convexPosition_of_concyclic}{SYMLEAN/SymmetricReduction.lean:509}]\label{sr:lem:concyclic}
Let $s\subseteq\R^2$ be finite, $c\in\R^2$ and $r\in\R$, and assume $\operatorname{eDist2}(p,c)=r$ for every $p\in s$.
Then $s$ is in convex position.
\end{lemma}

No general position is assumed, and $r$ is arbitrary (for $r=0$ the set has at most one point).

\begin{proof}
We apply Lemma~\ref{sr:lem:supports} with the tangent direction. For $p\in s$ let
$q=\bigl(p_0-(p_1-c_1),\;p_1+(p_0-c_0)\bigr)$, that is, $p$ plus the vector $p-c$ turned by $+90^\circ$. For every
$x\in\R^2$, expanding gives
\[
\orient(p,q,x)=-\bigl[(p_0-c_0)(x_0-p_0)+(p_1-c_1)(x_1-p_1)\bigr]
=\tfrac12\bigl(\operatorname{eDist2}(x,p)+\operatorname{eDist2}(p,c)-\operatorname{eDist2}(x,c)\bigr).
\]
(The second equality is $\abs{x-c}^2=\abs{x-p}^2+2\langle x-p,p-c\rangle+\abs{p-c}^2$.) For $x\in s$ we have
$\operatorname{eDist2}(x,c)=r=\operatorname{eDist2}(p,c)$, so $\orient(p,q,x)=\frac12\operatorname{eDist2}(x,p)$, which is
positive when $x\ne p$ (Lemma~\ref{sr:lem:complex}(e)). Lemma~\ref{sr:lem:supports} gives the claim.
\end{proof}

\begin{corollary}[A concyclic subset gives convex polygons; \leanref{ErdosSzekeres.SymReduction.hasConvexNGon_of_concyclic}{SYMLEAN/SymmetricReduction.lean:522}]\label{sr:cor:concyclic}
Let $T\subseteq S$ be finite sets of points, $c\in\R^2$ and $r\in\R$ with $\operatorname{eDist2}(p,c)=r$ for every $p\in
T$, and let $n\in\N$ with $n\le\card T$. Then \texttt{HasConvexNGon n S}.
\end{corollary}

\begin{proof}
$T$ is in convex position by Lemma~\ref{sr:lem:concyclic}, hence \texttt{HasConvexNGon n T} since $n\le\card T$, hence
\texttt{HasConvexNGon n S} by monotonicity (Lemma~\ref{fd:lem:mono}).
\end{proof}

\begin{lemma}[The centre of a half-turn; \leanref{ErdosSzekeres.SymReduction.center_not_mem_of_halfTurn}{SYMLEAN/SymmetricReduction.lean:528}]\label{sr:lem:center}
Let $S$ be in general position with $2\le\card S$, let $c\in\R^2$, and assume
$\operatorname{Invariant}(\operatorname{halfTurn}_c,S)$. Then $c\notin S$.
\end{lemma}

\begin{proof}
Suppose $c\in S$. Since $\card S\ge2$, there is $q\in S$ with $q\ne c$, and $q'=\operatorname{halfTurn}_c(q)=2c-q$ lies
in $S$. We have $q'\ne q$, since $2c-q=q$ would give $q=c$, and $q'\ne c$, since $2c-q=c$ would give $q=c$. So $q,c,q'$
are pairwise distinct points of $S$. But
\[
\orient(q,c,2c-q)=(c_0-q_0)\cdot2(c_1-q_1)-(c_1-q_1)\cdot2(c_0-q_0)=0,
\]
contradicting general position (Lemma~\ref{fd:lem:collinear}).
\end{proof}

\subsection{Roots of unity}\label{sr:ssec:roots}

\begin{definition}[Primitive roots of unity, order; written step]\label{sr:def:prim}
For $\zeta\in\C$ and $k\in\N$, $\zeta$ is a \emph{primitive $k$-th root of unity} (Mathlib \lean{IsPrimitiveRoot}) if
$\zeta^k=1$ and every $l\in\N$ with $\zeta^l=1$ is a multiple of $k$. If $\zeta^m=1$ for some $m>0$, the \emph{order}
$\operatorname{ord}(\zeta)$ (Mathlib \lean{orderOf}) is the least such $m$; otherwise it is $0$.
\end{definition}

\begin{lemma}[Facts on roots of unity; written step]\label{sr:lem:prim}
Let $\zeta,\xi\in\C$ and $n,k,d\in\N$.
\begin{enumerate}[label=(P\arabic*)]
\item If $\zeta^m=1$ for some $m>0$, then $\operatorname{ord}(\zeta)>0$, $\zeta$ is a primitive
  $\operatorname{ord}(\zeta)$-th root of unity, and $\operatorname{ord}(\zeta)=1$ if and only if $\zeta=1$. (Mathlib
  \lean{IsPrimitiveRoot.orderOf}, \lean{orderOf_eq_one_iff}.)
\item If $\zeta$ is a primitive $n$-th root of unity, $n>0$ and $n=k\,d$, then $\zeta^k$ is a primitive $d$-th root of
  unity. (\lean{IsPrimitiveRoot.pow}.)
\item A primitive $2$nd root of unity equals $-1$. (\lean{IsPrimitiveRoot.eq_neg_one_of_two_right}.)
\item If $\zeta$ is a primitive $k$-th root of unity and $k>1$, then $\zeta^k=1$, $\zeta\ne1$ and $\zeta\ne0$.
  (\lean{IsPrimitiveRoot.pow_eq_one}, \lean{IsPrimitiveRoot.ne_one}, \lean{IsPrimitiveRoot.ne_zero}.)
\item If $\zeta$ is a primitive $k$-th root of unity, $k>0$ and $\xi^k=1$, then $\xi=\zeta^e$ for some $e<k$.
  (\lean{IsPrimitiveRoot.eq_pow_of_pow_eq_one}.)
\item If $n>0$, then $\exp(2\pi i/n)$ is a primitive $n$-th root of unity. (\lean{Complex.isPrimitiveRoot_exp}.)
\end{enumerate}
\end{lemma}

\begin{proof}
(P1) Put $s=\operatorname{ord}(\zeta)$. Since $\zeta^m=1$ with $m>0$, Definition~\ref{sr:def:prim} makes $s$ the least
positive integer with $\zeta^s=1$; in particular $s>0$. If $\zeta^l=1$, write $l=ts+r$ with $0\le r<s$; then
$1=\zeta^l=(\zeta^s)^t\zeta^r=\zeta^r$, so $r=0$ by minimality, and $s\mid l$. So $\zeta$ is a primitive $s$-th root of
unity. Finally $s=1$ if and only if $\zeta^1=1$, that is, $\zeta=1$.

(P2) $(\zeta^k)^d=\zeta^n=1$. If $(\zeta^k)^l=1$, then $\zeta^{kl}=1$, so $kd=n\mid kl$; since $n>0$ we have $k>0$, and
so $d\mid l$.

(P3) Let $\zeta$ be a primitive $2$nd root. Then $\zeta^2=1$, and $\zeta\ne1$ because $\zeta^1=1$ would give $2\mid1$.
From $(\zeta-1)(\zeta+1)=\zeta^2-1=0$ we get $\zeta=-1$.

(P4) $\zeta^k=1$ by definition. If $\zeta=1$, then $\zeta^1=1$ gives $k\mid1$, so $k=1$. If $\zeta=0$, then
$\zeta^k=0\ne1$ since $k\ge1$.

(P5) The numbers $\zeta^0,\dots,\zeta^{k-1}$ are pairwise distinct: if $\zeta^i=\zeta^j$ with $i<j<k$, then
$\zeta^i\ne0$ (as $\zeta^k=1$), so $\zeta^{j-i}=1$ and $k\mid j-i$ with $0<j-i<k$, which is impossible. They are $k$
distinct roots of the polynomial $X^k-1$, which has degree $k$ and therefore at most $k$ roots in $\C$. Since $\xi$ is
also a root, it is one of them.

(P6) Let $\zeta=\exp(2\pi i/n)$. For $l\in\N$, $\zeta^l=\exp(2\pi il/n)$, and $\exp(w)=1$ holds exactly when $w\in2\pi
i\Z$. So $\zeta^l=1$ if and only if $l/n\in\Z$, that is, $n\mid l$. In particular $\zeta^n=1$.
\end{proof}

\begin{lemma}[The angles $2\pi/3$, $2\pi/4$, $2\pi/5$; \leanref{ErdosSzekeres.SymReduction.isPrimitiveRoot_rot3}{SYMLEAN/SymmetricReduction.lean:584}, \leanref{ErdosSzekeres.SymReduction.isPrimitiveRoot_rot5}{SYMLEAN/SymmetricReduction.lean:591}, \leanref{ErdosSzekeres.SymReduction.isPrimitiveRoot_rot4}{SYMLEAN/SymmetricReduction.lean:598}]\label{sr:lem:rot345}
For $n\in\{3,5,4\}$, $\exp\bigl((2\pi/n)\,i\bigr)$ is a primitive $n$-th root of unity. Here $2\pi/n$ is the real
number $2\pi/n$ (Lean \texttt{2 * Real.pi / n}), viewed as a complex number.
\end{lemma}

\begin{proof}
$(2\pi/n)i=2\pi i/n$; apply (P6).
\end{proof}

\begin{lemma}[Even order or an odd prime; \leanref{ErdosSzekeres.SymReduction.root_cases}{SYMLEAN/SymmetricReduction.lean:552}]\label{sr:lem:root-cases}
Let $a\in\C$ and $m\in\N$ with $0<m$, $a^m=1$ and $a\ne1$. Then either $a^k=-1$ for some $k\in\N$, or there are
$p,k\in\N$ such that $p$ is prime, $p\ne2$, and $a^k$ is a primitive $p$-th root of unity.
\end{lemma}

\begin{proof}
Let $n=\operatorname{ord}(a)$. By (P1), $n>0$, $a$ is a primitive $n$-th root of unity, and $n\ne1$ because $a\ne1$.
If $2\mid n$, write $n=2k=k\cdot2$; by (P2), $a^k$ is a primitive $2$nd root of unity, so $a^k=-1$ by (P3). Otherwise
let $p$ be the least prime factor of $n$ (it exists since $n>1$). Then $p\ne2$ because $n$ is odd, and $n=(n/p)\,p$, so
$a^{n/p}$ is a primitive $p$-th root of unity by (P2).
\end{proof}

In the proof the first alternative is taken for \emph{every} even order (so orders $6,10,12,\dots$ lead to the
half-turn), and in the second $p$ is the least prime factor of the order; the statement itself does not record this.

\begin{lemma}[Odd primes; \leanref{ErdosSzekeres.SymReduction.prime_cases}{SYMLEAN/SymmetricReduction.lean:572}]\label{sr:lem:prime-cases}
If $p$ is prime and $p\ne2$, then $p=3$, $p=5$ or $7\le p$.
\end{lemma}

\begin{proof}
$p\ge2$ and $p\ne2$, so $p\in\{3,4,5,6\}$ or $p\ge7$; and $4=2\cdot2$ and $6=2\cdot3$ are not prime.
\end{proof}

\begin{lemma}[Order $2$, a multiple of $4$, or an odd prime; \leanref{ErdosSzekeres.SymReduction.root_cases2}{SYMLEAN/SymmetricReduction.lean:607}]\label{sr:lem:root-cases2}
Let $a\in\C$ and $m\in\N$ with $0<m$, $a^m=1$ and $a\ne1$. Then $a=-1$, or $a^k$ is a primitive $4$th root of unity for
some $k\in\N$, or there are $p,k\in\N$ such that $p$ is prime, $p\ne2$, and $a^k$ is a primitive $p$-th root of unity.
\end{lemma}

\begin{proof}
Let $n=\operatorname{ord}(a)$; as before, $n>1$ and $a$ is a primitive $n$-th root of unity. If some prime $p\ne2$
divides $n$, then $a^{n/p}$ is a primitive $p$-th root of unity by (P2). Otherwise $2$ is the only prime factor of $n$.
If $4\mid n$, then $a^{n/4}$ is a primitive $4$th root of unity by (P2). Otherwise $n$ is even (its least prime factor
is $2$), say $n=2m'$ with $m'\ge1$. If $m'>1$, a prime factor of $m'$ divides $n$, so it equals $2$, and then $4\mid
n$, which is excluded. So $m'=1$, $n=2$, and $a=-1$ by (P3).
\end{proof}

\subsection{The case lemmas}\label{sr:ssec:cases}

\begin{lemma}[Half-turn invariant sets; \leanref{ErdosSzekeres.SymReduction.halfTurn_facts}{SYMLEAN/SymmetricReduction.lean:647}]\label{sr:lem:halfturn-facts}
Let $S$ be in general position with $2\le\card S$, let $c\in\R^2$, and assume
$\operatorname{Invariant}(\operatorname{halfTurn}_c,S)$. Then $c\notin S$, $2\mid\card S$, and for every $k\in\N$ with
$2k\le\card S$ there is $T\subseteq S$ with $\operatorname{Invariant}(\operatorname{halfTurn}_c,T)$ and $\card T=2k$.
\end{lemma}

\begin{proof}
$c\notin S$ by Lemma~\ref{sr:lem:center}. Let $f=\operatorname{halfTurn}_c=\operatorname{crot}_{\hat c,-1}$
(Lemma~\ref{sr:lem:rot-half}(b)). By Lemma~\ref{sr:lem:crot}(d),(b), $f^{[2]}=\operatorname{crot}_{\hat c,1}=\mathrm{id}$.
If $f(y)=y$ for some $y\in S$, then $\hat y=\hat c$ by Lemma~\ref{sr:lem:crot}(e) (as $-1\ne1$), so $y=c\in S$, a
contradiction. $f$ is injective by Lemma~\ref{sr:lem:crot}(f). Lemma~\ref{sr:lem:orbit}(b) with the prime $p=2$ gives
the claims.
\end{proof}

\begin{lemma}[Case of a half-turn; \leanref{ErdosSzekeres.SymReduction.case_halfTurn}{SYMLEAN/SymmetricReduction.lean:670}; conditional on \lean{HalfTurn34}]\label{sr:lem:case-halfturn}
Assume \lean{HalfTurn34} (a hypothesis). Let $S$ be in general position with $33\le\card S$, let $c\in\R^2$, and assume
$\operatorname{Invariant}(\operatorname{halfTurn}_c,S)$. Then \texttt{HasConvexNGon 7 S}.
\end{lemma}

\begin{proof}
By Lemma~\ref{sr:lem:halfturn-facts}, $c\notin S$ and $\card S$ is even; with $\card S\ge33$ this gives $\card S\ge34=17\cdot2$.
So there is $T\subseteq S$, invariant under $\operatorname{halfTurn}_c$, with $\card T=34$. $T$ is in general position
(monotonicity) and $c\notin T$. \lean{HalfTurn34} applied to $T$ and $c$ gives \texttt{HasConvexNGon 7 T}, hence
\texttt{HasConvexNGon 7 S} by monotonicity.
\end{proof}

\begin{lemma}[Rotations of prime order off the centre; \leanref{ErdosSzekeres.SymReduction.rot_erase_facts}{SYMLEAN/SymmetricReduction.lean:676}]\label{sr:lem:rot-erase}
Let $S$ be finite, $c,w\in\C$, and $p$ a prime with $w^p=1$ and $w\ne1$; assume
$\operatorname{Invariant}(\operatorname{crot}_{c,w},S)$. Let $S'=S\setminus\{\operatorname{ofC}(c)\}$. Then
$p\mid\card{S'}$, and for every $k\in\N$ with $kp\le\card{S'}$ there is $T\subseteq S$ with
$\operatorname{ofC}(c)\notin T$, $\operatorname{Invariant}(\operatorname{crot}_{c,w},T)$ and $\card T=kp$.
\end{lemma}

\begin{proof}
Let $f=\operatorname{crot}_{c,w}$. Since $p>0$ and $w^p=1$, $w\ne0$, so $f$ is injective (Lemma~\ref{sr:lem:crot}(f)).
The centre is fixed: $\operatorname{toC}(f(\operatorname{ofC}c))=c+w(c-c)=c$, so $f(\operatorname{ofC}c)=\operatorname{ofC}c$.
\emph{$S'$ is mapped into itself:} for $y\in S'$, $f(y)\in S$, and $f(y)=\operatorname{ofC}(c)=f(\operatorname{ofC}c)$
would give $y=\operatorname{ofC}(c)$ by injectivity. Further $f^{[p]}=\operatorname{crot}_{c,w^p}=\operatorname{crot}_{c,1}=
\mathrm{id}$ (Lemma~\ref{sr:lem:crot}(d),(b)). If $f(y)=y$ for some $y\in S'$, then $\hat y=c$ by
Lemma~\ref{sr:lem:crot}(e), so $y=\operatorname{ofC}(c)\notin S'$, a contradiction. Lemma~\ref{sr:lem:orbit}(b) applied
to $S'$ gives $p\mid\card{S'}$ and, for $kp\le\card{S'}$, an $f$-invariant $T\subseteq S'$ with $\card T=kp$; then
$T\subseteq S$ and $\operatorname{ofC}(c)\notin T$.
\end{proof}

\begin{lemma}[Rotations of prime order at least $7$; \leanref{ErdosSzekeres.SymReduction.case_rot_big}{SYMLEAN/SymmetricReduction.lean:711}]\label{sr:lem:rot-big}
Let $S$ be finite, $c,w\in\C$, and $p$ a prime with $7\le p$, $w^p=1$ and $w\ne1$. Assume
$\operatorname{Invariant}(\operatorname{crot}_{c,w},S)$ and $2\le\card S$. Then \texttt{HasConvexNGon 7 S}.
\end{lemma}

No general position is assumed.

\begin{proof}
Let $f=\operatorname{crot}_{c,w}$. Since $\card S\ge2$, some $q\in S$ differs from $\operatorname{ofC}(c)$. We have
$\mathrm{normSq}(w)^p=\mathrm{normSq}(w^p)=1$ and $\mathrm{normSq}(w)\ge0$; as $t\mapsto t^p$ is strictly increasing on
$[0,\infty)$, $\mathrm{normSq}(w)=1$. Let $O=\operatorname{orbitFS}(f,p,q)$. Then $O\subseteq S$
(Lemma~\ref{sr:lem:orbitFS}(b)); $f^{[p]}(q)=\operatorname{crot}_{c,w^p}(q)=q$; and $f(q)\ne q$, since otherwise $\hat
q=c$ by Lemma~\ref{sr:lem:crot}(e) and $q=\operatorname{ofC}(c)$. By Lemma~\ref{sr:lem:orbitFS}(d), $\card O=p$. Every
$x\in O$ has the form $x=f^{[i]}(q)=\operatorname{crot}_{c,w^i}(q)$ with $\mathrm{normSq}(w^i)=\mathrm{normSq}(w)^i=1$, so
$\operatorname{eDist2}(x,\operatorname{ofC}c)=\operatorname{eDist2}(q,\operatorname{ofC}c)$ by
Lemma~\ref{sr:lem:crot}(g). Corollary~\ref{sr:cor:concyclic} with $T=O$, centre $\operatorname{ofC}(c)$, radius
$\operatorname{eDist2}(q,\operatorname{ofC}c)$ and $n=7\le p=\card O$ gives \texttt{HasConvexNGon 7 S}.
\end{proof}

\begin{lemma}[Case of an odd prime order; \leanref{ErdosSzekeres.SymReduction.case_odd_prime}{SYMLEAN/SymmetricReduction.lean:736}; conditional on \lean{Threefold33} and \lean{Fivefold35}]\label{sr:lem:odd-prime}
Assume \lean{Threefold33} and \lean{Fivefold35} (hypotheses). Let $S$ be in general position with $33\le\card S$, let
$c,w\in\C$, let $p$ be a prime with $p\ne2$, let $w$ be a primitive $p$-th root of unity, and assume
$\operatorname{Invariant}(\operatorname{crot}_{c,w},S)$. Then \texttt{HasConvexNGon 7 S}.
\end{lemma}

\begin{proof}
By (P4), $w^p=1$ and $w\ne1$. Let $S'=S\setminus\{\operatorname{ofC}(c)\}$; then $\card{S'}\ge\card S-1\ge32$. By
Lemma~\ref{sr:lem:prime-cases}, $p=3$, $p=5$ or $p\ge7$.

\emph{$p=3$.} By Lemma~\ref{sr:lem:rot-erase}, $3\mid\card{S'}$; with $\card{S'}\ge32$ this gives $\card{S'}\ge33=11\cdot3$.
So there is $T\subseteq S$ with $\operatorname{ofC}(c)\notin T$, $\operatorname{Invariant}(\operatorname{crot}_{c,w},T)$
and $\card T=33$. Let $\zeta=\exp((2\pi/3)i)$. By Lemma~\ref{sr:lem:rot345}, $\zeta^3=1$, so by (P5) $\zeta=w^e$ for
some $e$. By Lemmas~\ref{sr:lem:rot-half}(a), \ref{sr:lem:complex}(c) and~\ref{sr:lem:crot}(d),
\[
\operatorname{rotAbout}_{\operatorname{ofC}(c),2\pi/3}=\operatorname{crot}_{\operatorname{toC}(\operatorname{ofC}c),\,\zeta}
=\operatorname{crot}_{c,w^e}=(\operatorname{crot}_{c,w})^{[e]},
\]
so $\operatorname{Invariant}(\operatorname{rotAbout}_{\operatorname{ofC}(c),2\pi/3},T)$ by Lemma~\ref{sr:lem:iterate}(b).
$T$ is in general position (monotonicity), $\card T=33$ and $\operatorname{ofC}(c)\notin T$. \lean{Threefold33}, applied
to $T$ and the centre $\operatorname{ofC}(c)$, gives \texttt{HasConvexNGon 7 T}, hence \texttt{HasConvexNGon 7 S}.

\emph{$p=5$.} In the same way, $5\mid\card{S'}$ and $\card{S'}\ge32$ give $\card{S'}\ge35=7\cdot5$; take $T\subseteq S$
with $\operatorname{ofC}(c)\notin T$, invariant under $\operatorname{crot}_{c,w}$, with $\card T=35$; with
$\zeta=\exp((2\pi/5)i)=w^e$, $T$ is invariant under $\operatorname{rotAbout}_{\operatorname{ofC}(c),2\pi/5}$; and
\lean{Fivefold35} gives \texttt{HasConvexNGon 7 T}, hence \texttt{HasConvexNGon 7 S}.

\emph{$p\ge7$.} Lemma~\ref{sr:lem:rot-big} applies, since $\card S\ge2$.
\end{proof}

\begin{lemma}[Case of a primitive fourth root; \leanref{ErdosSzekeres.SymReduction.case_fourfold}{SYMLEAN/SymmetricReduction.lean:763}; conditional on \lean{Fourfold32}]\label{sr:lem:fourfold}
Assume \lean{Fourfold32} (a hypothesis). Let $S$ be in general position with $33\le\card S$, let $c,w\in\C$, let $w$ be
a primitive $4$th root of unity, and assume $\operatorname{Invariant}(\operatorname{crot}_{c,w},S)$. Then
\texttt{HasConvexNGon 7 S}.
\end{lemma}

\begin{proof}
By (P4), $w^4=1$, $w\ne1$ and $w\ne0$; by (P2) with $4=2\cdot2$ and (P3), $w^2=-1$. Let $f=\operatorname{crot}_{c,w}$.
By Lemmas~\ref{sr:lem:iterate}(b) and~\ref{sr:lem:crot}(d), $S$ is invariant under $f^{[2]}=\operatorname{crot}_{c,-1}$,
which equals $\operatorname{halfTurn}_{\operatorname{ofC}(c)}$ by Lemma~\ref{sr:lem:rot-half}(b) and
$\operatorname{toC}(\operatorname{ofC}c)=c$. By Lemma~\ref{sr:lem:center}, $\operatorname{ofC}(c)\notin S$. Hence $\hat
y\ne c$ for every $y\in S$.

\emph{Every $y\in S$ has minimal period $4$ under $f$.} First, $f^{[4]}(y)=\operatorname{crot}_{c,w^4}(y)=y$, so by (MP1)
$d=\operatorname{minimalPeriod}(f,y)$ divides $4$, and $d\in\{1,2,4\}$. If $d=1$, then $f(y)=y$ by (MP2), and $\hat y=c$
by Lemma~\ref{sr:lem:crot}(e), which is impossible. If $d=2$, then $f^{[2]}(y)=y$ by (MP3), that is,
$\operatorname{crot}_{c,-1}(y)=y$, and again $\hat y=c$ (as $-1\ne1$). So $d=4$.

$f$ is injective (Lemma~\ref{sr:lem:crot}(f)). Lemma~\ref{sr:lem:orbit}(a) with $p=4$ gives $4\mid\card S$ and invariant
subsets of every size $4k\le\card S$. Since $8\cdot4=32\le\card S$, there is $T\subseteq S$ with
$\operatorname{Invariant}(f,T)$ and $\card T=32$. (Also $\card S\ge36$, but this is not needed.) Let
$\zeta=\exp((2\pi/4)i)$; by Lemma~\ref{sr:lem:rot345} and (P5), $\zeta=w^e$ for some $e$, and as in the proof of
Lemma~\ref{sr:lem:odd-prime}, $\operatorname{rotAbout}_{\operatorname{ofC}(c),2\pi/4}=f^{[e]}$, so $T$ is invariant under
it. $T$ is in general position, $\card T=32$, and $\operatorname{ofC}(c)\notin T$ because $T\subseteq S$.
\lean{Fourfold32} gives \texttt{HasConvexNGon 7 T}, hence \texttt{HasConvexNGon 7 S}.
\end{proof}

\begin{lemma}[Case of a reflection; \leanref{ErdosSzekeres.SymReduction.case_mirror}{SYMLEAN/SymmetricReduction.lean:811}; conditional on the two mirror classes]\label{sr:lem:mirror}
Assume \lean{MirrorOneAxis33} and \lean{MirrorNoAxis34} (hypotheses). Let $a,u\in\R^2$ with $u_0^2+u_1^2=1$, let $S$ be in
general position with $33\le\card S$, and assume $\operatorname{Invariant}(\operatorname{reflLine}_{a,u},S)$. Then
\texttt{HasConvexNGon 7 S}.
\end{lemma}

\begin{proof}
Let $\sigma=\operatorname{reflLine}_{a,u}$, $F=\{p\in S:\sigma(p)=p\}$ and $M=\{p\in S:\sigma(p)\ne p\}$, so that
$\card F+\card M=\card S$.

\emph{$\card F\le2$.} Otherwise $F$ contains three pairwise distinct points $x,y,z$. By Lemma~\ref{sr:lem:refl}(c),
each $t\in\{x,y,z\}$ satisfies $(t_0-a_0)u_1-(t_1-a_1)u_0=0$, so $\orient(x,y,z)=0$ by Lemma~\ref{sr:lem:refl}(d),
contradicting general position.

\emph{$M$ is mapped into itself.} By Lemma~\ref{sr:lem:refl}(b), $\sigma\circ\sigma=\mathrm{id}$; in particular $\sigma$
is injective. Let $y\in M$. Then $\sigma(y)\in S$, and if $\sigma(\sigma(y))=\sigma(y)$ then $y=\sigma(y)$, which is
false; so $\sigma(y)\in M$. Moreover $\sigma^{[2]}(y)=y$ and $\sigma(y)\ne y$ for $y\in M$. Lemma~\ref{sr:lem:orbit}(b)
with the prime $2$ gives: $2\mid\card M$, and for every $k$ with $2k\le\card M$ there is $T\subseteq M$, mapped into
itself by $\sigma$, with $\card T=2k$.

\emph{Case $\card F=0$.} Then $\card M=\card S\ge33$ and $\card M$ is even, so $\card M\ge34$. Take $T\subseteq M$
invariant with $\card T=34$. $T\subseteq S$ is in general position, and no point of $T$ is fixed by $\sigma$ (as
$T\subseteq M$). \lean{MirrorNoAxis34} applied to $T$, $a$, $u$ gives \texttt{HasConvexNGon 7 T}, hence
\texttt{HasConvexNGon 7 S}.

\emph{Case $\card F\ge1$.} Pick $x\in F$. Since $\card F\le2$, $\card M\ge31$, and $\card M$ is even, so $\card M\ge32$.
Take $T\subseteq M$ invariant with $\card T=32$, and put $T'=T\cup\{x\}\subseteq S$. Then $x\notin T$ (points of $T$ are
not fixed), so $\card{T'}=33$. $T'$ is invariant: $\sigma(x)=x$, and $T$ is invariant. The point $x\in T'$ is fixed, and
every fixed $y\in T'$ equals $x$, because no point of $T$ is fixed. $T'$ is in general position. \lean{MirrorOneAxis33}
applied to $T'$, $a$, $u$ gives \texttt{HasConvexNGon 7 T'}, hence \texttt{HasConvexNGon 7 S}.
\end{proof}

\subsection{The reduction theorems}\label{sr:ssec:main}

\begin{theorem}[Main reduction; \leanref{ErdosSzekeres.SymReduction.heptagon_of_nontrivial_isometry}{SYMLEAN/SymmetricReduction.lean:997}; conditional on five case classes]\label{sr:thm:main}
Assume the propositions \lean{HalfTurn34}, \lean{Threefold33}, \lean{Fivefold35}, \lean{MirrorOneAxis33} and
\lean{MirrorNoAxis34} of Definition~\ref{sr:def:classes} (hypotheses; none is proved in this section). Let $S$ be a
finite set of points and $g\colon\R^2\to\R^2$ a map. If $S$ is in general position, $33\le\card S$, $g$ is a Euclidean
isometry, $g\ne\mathrm{id}$ and $\operatorname{Invariant}(g,S)$, then \texttt{HasConvexNGon 7 S}.
\end{theorem}

\begin{proof}
Since $\card S\ge33>1$, Theorem~\ref{sr:thm:structure}(b) applies.

\emph{Rotation: $g=\operatorname{crot}_{c,a}$ with $a\ne1$ and $a^m=1$, $m>0$.} Then
$\operatorname{Invariant}(\operatorname{crot}_{c,a},S)$, and for every $k$,
$\operatorname{Invariant}(\operatorname{crot}_{c,a^k},S)$ by Lemmas~\ref{sr:lem:iterate}(b) and~\ref{sr:lem:crot}(d).
By Lemma~\ref{sr:lem:root-cases} there are two cases.
\begin{itemize}
\item $a^k=-1$ for some $k$. Then $S$ is invariant under $\operatorname{crot}_{c,-1}=
  \operatorname{halfTurn}_{\operatorname{ofC}(c)}$ (Lemma~\ref{sr:lem:rot-half}(b) and
  $\operatorname{toC}(\operatorname{ofC}c)=c$), and Lemma~\ref{sr:lem:case-halfturn} gives the claim.
\item $a^k$ is a primitive $p$-th root of unity for a prime $p\ne2$. Then Lemma~\ref{sr:lem:odd-prime}, applied with
  $w=a^k$, gives the claim.
\end{itemize}

\emph{Reflection: $g=\operatorname{reflLine}_{A,u}$ with $u_0^2+u_1^2=1$.} Lemma~\ref{sr:lem:mirror} gives the claim.
\end{proof}

\begin{remark}[Status of Theorem~\ref{sr:thm:main}]\label{sr:rem:main}
The theorem is Lean-checked as an implication: the five classes are explicit hypotheses (Lean binders
\texttt{h2 h3 h5 hm1 hm0}), not axioms, and the axioms log is \dirn{SYMLEAN/logs/SymmetricReductionAxioms.out}. The
unconditional statement, Theorem~\ref{sr:thm:main} without its five hypotheses, is claimed in the project as
computer-assisted (Theorem~A of \dirn{WRITEUP/SYMMETRIC.md}, \S0, in its form for Euclidean isometries), on the
evidence summarized in Remark~\ref{sr:rem:records}; it is not Lean-checked as a whole. In this document the theorem is
established only as the stated implication. Theorem~\ref{mb:thm:main3} restates the theorem with the two mirror
classes discharged, so that only \lean{HalfTurn34}, \lean{Threefold33} and \lean{Fivefold35} remain hypotheses; it is
stated in \dirn{SYMLEAN/mirror/ReductionThreeClasses.lean}, for the byte-identical copies of the definitions described
in Remark~\ref{sr:rem:records}. The
symmetry $g$ must be a Euclidean isometry of the whole plane; a distance-preserving permutation of $S$ alone is not
treated. Affine symmetries are treated in Theorem~\ref{ar:thm:main}.
\end{remark}

\begin{theorem}[Variant; \leanref{ErdosSzekeres.SymReduction.halfTurn_of_symmetric_counterexample}{SYMLEAN/SymmetricReduction.lean:1024}; conditional on five case classes]\label{sr:thm:variant}
Assume the propositions \lean{Threefold33}, \lean{Fourfold32}, \lean{Fivefold35}, \lean{MirrorOneAxis33} and
\lean{MirrorNoAxis34} of Definition~\ref{sr:def:classes} (hypotheses; none is proved in this section). Let $S$ be a
finite set of points and $g\colon\R^2\to\R^2$ a map. Assume that $S$ is in general position, $33\le\card S$, $g$ is a
Euclidean isometry, $g\ne\mathrm{id}$, $\operatorname{Invariant}(g,S)$, and $\neg$\,\texttt{HasConvexNGon 7 S}. Then
there is $c\in\R^2$ such that $g=\operatorname{halfTurn}_c$ (as maps of $\R^2$), $c\notin S$, $2\mid\card S$ and
$34\le\card S$.
\end{theorem}

\begin{proof}
By Theorem~\ref{sr:thm:structure}(b), as $\card S>1$, $g$ is a rotation or a reflection.

\emph{Rotation: $g=\operatorname{crot}_{c,a}$ with $a\ne1$ and $a^m=1$, $m>0$.} By Lemma~\ref{sr:lem:root-cases2} there
are three cases.
\begin{itemize}
\item $a=-1$. Then $g=\operatorname{crot}_{c,-1}=\operatorname{halfTurn}_{\operatorname{ofC}(c)}$
  (Lemma~\ref{sr:lem:rot-half}(b) and $\operatorname{toC}(\operatorname{ofC}c)=c$). By Lemma~\ref{sr:lem:halfturn-facts}, $\operatorname{ofC}(c)\notin S$ and
  $2\mid\card S$; with $\card S\ge33$ this gives $\card S\ge34$. So the point $\operatorname{ofC}(c)$ has all the
  required properties.
\item $a^k$ is a primitive $4$th root of unity. $S$ is invariant under $\operatorname{crot}_{c,a^k}$
  (Lemmas~\ref{sr:lem:iterate}(b) and~\ref{sr:lem:crot}(d)), and Lemma~\ref{sr:lem:fourfold} gives
  \texttt{HasConvexNGon 7 S}, contradicting the hypothesis.
\item $a^k$ is a primitive $p$-th root of unity for a prime $p\ne2$. Likewise Lemma~\ref{sr:lem:odd-prime} gives a
  contradiction.
\end{itemize}

\emph{Reflection.} Lemma~\ref{sr:lem:mirror} gives \texttt{HasConvexNGon 7 S}, a contradiction.
\end{proof}

\begin{remark}[Scope of Theorem~\ref{sr:thm:variant}]\label{sr:rem:variant}
The theorem concerns one symmetry at a time and Euclidean isometries only. Corollary~3 of
\dirn{WRITEUP/SYMMETRIC.md} has the same per-symmetry form (every symmetry of such a set is a half-turn about a point
outside the set), with the case theorems and the fourfold exclusion as assumptions, but there a symmetry may also be
an affine bijection; the theorem is the isometry case of that corollary. A group form, stated as the corollary labelled
\texttt{cor:only} in \dirn{WRITEUP/note/symmetric_es7.tex}, is Lean-checked for Euclidean isometries, with the same five
classes as hypotheses, in the review file \dirn{SYMLEAN/review/VariantChecks.lean} (\lean{VariantChecks.group_form},
line 1168): for $S$ as in the theorem, the Euclidean isometries that map $S$ into itself are the identity and at most
one half-turn, whose centre is not in $S$. The evidence is the log \dirn{SYMLEAN/logs/VariantChecks.out}, which prints
the axioms \lean{propext}, \lean{Classical.choice} and \lean{Quot.sound} for \lean{VariantChecks.group_form}, and the
row \texttt{review/VariantChecks.lean} of \dirn{SYMLEAN/logs/runs.tsv}, which records exit code~$0$ for the file with
SHA-256 prefix \texttt{4b823bb6}; the file at the commit used here has the prefix \texttt{4b823bb6cf85}. The group form
follows from the theorem. Let $g$ and $g'$ be Euclidean isometries other than the identity that map $S$ into itself.
The theorem gives $g=\operatorname{halfTurn}_c$ and $g'=\operatorname{halfTurn}_d$ with $c,d\notin S$. If $c\ne d$,
then $g\circ g'$ is the translation $p\mapsto p+2(c-d)$. It is a Euclidean isometry other than the identity that maps
$S$ into itself, so by the theorem it equals $\operatorname{halfTurn}_e$ for some $e\in\R^2$; evaluating at $p=e$ gives
$e+2(c-d)=e$, so $c=d$, a contradiction. The group form is not restated in this document as a numbered statement, and
no statement of this document uses it. The variant does not use \lean{HalfTurn34}; it uses \lean{Fourfold32}, which
rests on \dirn{C4CUT} (Remark~\ref{sr:rem:records}).
\end{remark}

\begin{definition}[The \dirn{H3} mirror statement; \leanref{ErdosSzekeres.SymReduction.MirrorExclusionLineCopy}{SYMLEAN/SymmetricReduction.lean:1052}]\label{sr:def:exclusion-line}
\lean{MirrorExclusionLineCopy} is the proposition: for every finite $S\subseteq\R^2$ and all $a,u\in\R^2$, if
$u_0^2+u_1^2=1$, $S$ is in general position, $\card S=33$, and $\operatorname{reflLine}_{a,u}(q)\in S$ for every $q\in
S$, then \texttt{HasConvexNGon 7 S}.
\end{definition}

\begin{lemma}[Bridge; \leanref{ErdosSzekeres.SymReduction.mirrorOneAxis33_of_exclusionLine}{SYMLEAN/SymmetricReduction.lean:1056}]\label{sr:lem:mirror-line}
\lean{MirrorExclusionLineCopy} implies \lean{MirrorOneAxis33}.
\end{lemma}

\begin{proof}
Let $S$, $a$, $u$ satisfy the hypotheses of \lean{MirrorOneAxis33}. The condition
$\operatorname{Invariant}(\operatorname{reflLine}_{a,u},S)$ is, by definition, the condition
$\operatorname{reflLine}_{a,u}(q)\in S$ for every $q\in S$. So all hypotheses of \lean{MirrorExclusionLineCopy} hold
(the hypothesis about the fixed point is not used), and it gives \texttt{HasConvexNGon 7 S}.
\end{proof}

\begin{remark}[About the bridge]\label{sr:rem:mirror-line}
\lean{MirrorExclusionLineCopy} has the same text as \lean{Mir.MirrorExclusionLine}
(\dirn{H3/MirrorEncoding/Line.lean:94}), and \lean{reflLine} has the same text as \lean{Mir.reflLine}. They are
different Lean constants of separately compiled developments, and no Lean theorem in \dirn{SYMLEAN} identifies them;
\dirn{SYMLEAN/README.md} \S5 reports a separate check of both identities by \texttt{rfl}, whose file is not kept in
\dirn{SYMLEAN}. In \dirn{H3}, \lean{Mir.mirrorExclusionLine_of_unsat} and
\lean{Mir.mirrorExclusionLine_of_cubes} (\dirn{H3/MirrorEncoding/Line.lean:139} and \dirn{:147}) derive
\lean{Mir.MirrorExclusionLine} from hypotheses saying that certain CNF formulas are unsatisfiable (the second also from
two hypotheses that the Boolean cover check \texttt{coverCheck} of a cube list returns \texttt{true}). In \dirn{V2-L},
\lean{V2L.mirrorExclusionLine_final} (\dirn{V2-L/V2L/FinalTree.lean:28}; \dirn{V2-L/README.md} \S9.1 and \S9.6)
derives it from hypotheses that all carry recorded verdicts. It is Lean-checked with four hypotheses, \texttt{h1},
\texttt{hroot}, \texttt{hup} and \texttt{hleaves}, which together assert $1{,}198+1+121+74{,}000$ facts of the form
``this CNF formula is unsatisfiable''; the $1+121$ facts of \texttt{hroot} and \texttt{hup} are the cover obligations
of the top two levels of the hull-pair case split of sub-case (ii) (Remark~\ref{sr:rem:records};
\dirn{WRITEUP/MIRROR.md} \S4). Each fact has a recorded \texttt{cake\_lpr} verdict on a file whose proof was
not retained, and the step from a verdict on a file to the Lean hypothesis is made outside Lean
(\dirn{WRITEUP/MIRROR.md} \S6). So
Lemma~\ref{sr:lem:mirror-line} supplies \lean{MirrorOneAxis33} only conditionally on those CNF facts and on the
identification of the two texts; the Lean proof without SAT input is in \dirn{SYMLEAN/mirror}
(Remark~\ref{sr:rem:records}). The converse implication also holds, and is Lean-checked in the review file
\dirn{SYMLEAN/review/ReviewSanity.lean} (\lean{Review.exclusionLine_of_mirrorOneAxis33}, line 1593, and the
equivalence \lean{Review.mirrorOneAxis33_iff_exclusionLine}, line 1629, whose axioms the log
\dirn{SYMLEAN/logs/ReviewSanity.out} prints): for $33$ points the fixed-point hypothesis of \lean{MirrorOneAxis33} is
automatic. Indeed, let $S$, $a$, $u$ satisfy the hypotheses of \lean{MirrorExclusionLineCopy}. In the notation of the
proof of Lemma~\ref{sr:lem:mirror}, whose first two steps use only these hypotheses, $\card F\le2$, $\card M$ is even
and $\card F+\card M=33$, so $\card F=1$. The point of $F$ is then the only point of $S$ fixed by
$\operatorname{reflLine}_{a,u}$, and \lean{MirrorOneAxis33} gives \texttt{HasConvexNGon 7 S}. Neither review theorem
is restated in this document as a numbered statement, and no statement of this document uses them.
\end{remark}

\begin{remark}[Hypotheses of the classes that follow from the others]\label{sr:rem:redundant}
Some hypotheses of the classes are implied by the others; this is not used anywhere.
\begin{itemize}
\item In \lean{HalfTurn34}, $c\notin S$ follows from general position and invariance (Lemma~\ref{sr:lem:center}).
\item In \lean{Fourfold32}, $c\notin S$ follows as well: by Lemma~\ref{sr:lem:rot-half}(a) the rotation is
  $f=\operatorname{crot}_{\hat c,\zeta}$ with $\zeta=\exp((2\pi/4)i)$, a primitive $4$th root of unity
  (Lemma~\ref{sr:lem:rot345}), so $\zeta^2=-1$ by (P2) and (P3). By Lemmas~\ref{sr:lem:crot}(d)
  and~\ref{sr:lem:rot-half}(b), $f^{[2]}=\operatorname{crot}_{\hat c,-1}=\operatorname{halfTurn}_c$; $S$ is invariant under
  it (Lemma~\ref{sr:lem:iterate}(b)), and Lemma~\ref{sr:lem:center} applies.
\item In \lean{Threefold33} and \lean{Fivefold35}, $c\notin S$ follows from the size. Let $p\in\{3,5\}$ and
  $w=\exp((2\pi/p)i)$, so that $w^p=1$ and $w\ne1$ (Lemma~\ref{sr:lem:rot345} and (P4)), and the rotation is
  $\operatorname{crot}_{\hat c,w}$ (Lemma~\ref{sr:lem:rot-half}(a)). If $c\in S$, then by Lemma~\ref{sr:lem:rot-erase}
  (with $\operatorname{ofC}(\hat c)=c$), $p$ divides $\card{(S\setminus\{c\})}$, which is $32$ for $p=3$ and $34$ for
  $p=5$; but $3\nmid32$ and $5\nmid34$.
\item In \lean{MirrorOneAxis33}, the fixed-point hypothesis follows from $\card S=33$ (Remark~\ref{sr:rem:mirror-line}).
  In \lean{MirrorNoAxis34} the hypothesis that no point is fixed is not implied in this way: for $34$ points the same
  argument gives only $\card F\in\{0,2\}$.
\end{itemize}
Dropping one of the implied hypotheses above gives an equivalent class. In general, adding a hypothesis to a class can
only make it weaker (the class without the hypothesis implies the class with it), that is, easier to supply.
\end{remark}

\paragraph{Coverage.}
All $72$ declarations of \dirn{SYMLEAN/SymmetricReduction.lean} are covered.
\begin{itemize}
\item Definitions: \lean{eDist2}, \lean{IsEucIsometry} (Definition~\ref{sr:def:iso}); \lean{Invariant}
  (Definition~\ref{sr:def:inv}); \lean{toC}, \lean{ofC} (Definition~\ref{sr:def:complex}); \lean{crot}, \lean{rotAbout},
  \lean{halfTurn}, \lean{reflLine} (Definition~\ref{sr:def:maps}); \lean{orbitFS} (Definition~\ref{sr:def:orbitFS}); the
  six classes \lean{HalfTurn34}, \lean{Threefold33}, \lean{Fivefold35}, \lean{MirrorOneAxis33}, \lean{MirrorNoAxis34},
  \lean{Fourfold32} (Definition~\ref{sr:def:classes}); \lean{MirrorExclusionLineCopy}
  (Definition~\ref{sr:def:exclusion-line}).
\item Finite dynamics: \lean{iterate_mem}, \lean{Invariant.iterate} (Lemma~\ref{sr:lem:iterate});
  \lean{mem_orbitFS_self}, \lean{orbitFS_subset}, \lean{orbitFS_invariant}, \lean{card_orbitFS},
  \lean{card_orbitFS_of_minimalPeriod} (Lemma~\ref{sr:lem:orbitFS}); \lean{orbit_select_minimalPeriod},
  \lean{orbit_select} (Lemma~\ref{sr:lem:orbit}).
\item Coordinates and maps: \lean{toC_re}, \lean{toC_im}, \lean{ofC_zero}, \lean{ofC_one}, \lean{toC_ofC},
  \lean{ofC_toC}, \lean{toC_injective}, \lean{eDist2_eq_normSq}, \lean{eDist2_pos_of_ne}, \lean{point_ext}
  (Lemma~\ref{sr:lem:complex}); \lean{toC_crot}, \lean{crot_one}, \lean{crot_comp}, \lean{crot_iterate},
  \lean{crot_eq_self_iff}, \lean{crot_injective}, \lean{eDist2_crot_center} (Lemma~\ref{sr:lem:crot});
  \lean{rotAbout_eq_crot}, \lean{halfTurn_eq_crot} (Lemma~\ref{sr:lem:rot-half}); \lean{toC_reflLine},
  \lean{reflLine_involutive}, \lean{cross_of_reflLine_fixed}, \lean{orient_eq_zero_of_cross}
  (Lemma~\ref{sr:lem:refl}).
\item Isometries: \lean{IsEucIsometry.injective}, \lean{isometry_toC_conj} (Lemma~\ref{sr:lem:toC-iso});
  \lean{isometry_complex_cases}, \lean{isometry_structure} (Theorem~\ref{sr:thm:structure}).
\item Geometry: \lean{convexPosition_of_concyclic} (Lemma~\ref{sr:lem:concyclic}), \lean{hasConvexNGon_of_concyclic}
  (Corollary~\ref{sr:cor:concyclic}), \lean{center_not_mem_of_halfTurn} (Lemma~\ref{sr:lem:center}); also the library
  lemma \lean{Suk.convexPosition_of_strict_supports} (Lemma~\ref{sr:lem:supports}). This is the lemma titled ``Strict
  supporting lines'' in the foundations section; the short proof is repeated here.
\item Roots of unity: \lean{root_cases} (Lemma~\ref{sr:lem:root-cases}); \lean{prime_cases}
  (Lemma~\ref{sr:lem:prime-cases}); \lean{root_cases2} (Lemma~\ref{sr:lem:root-cases2}); and
  \lean{isPrimitiveRoot_rot3}, \lean{isPrimitiveRoot_rot5}, \lean{isPrimitiveRoot_rot4} (Lemma~\ref{sr:lem:rot345}).
\item Cases: \lean{halfTurn_facts} (Lemma~\ref{sr:lem:halfturn-facts}); \lean{case_halfTurn}
  (Lemma~\ref{sr:lem:case-halfturn}, conditional on \lean{HalfTurn34}); \lean{rot_erase_facts}
  (Lemma~\ref{sr:lem:rot-erase}); \lean{case_rot_big} (Lemma~\ref{sr:lem:rot-big}); \lean{case_odd_prime}
  (Lemma~\ref{sr:lem:odd-prime}, conditional on \lean{Threefold33} and \lean{Fivefold35}); \lean{case_fourfold}
  (Lemma~\ref{sr:lem:fourfold}, conditional on \lean{Fourfold32}); \lean{case_mirror} (Lemma~\ref{sr:lem:mirror},
  conditional on \lean{MirrorOneAxis33} and \lean{MirrorNoAxis34}).
\item Main results: \lean{heptagon_of_nontrivial_isometry} (Theorem~\ref{sr:thm:main}, conditional on
  \lean{HalfTurn34}, \lean{Threefold33}, \lean{Fivefold35}, \lean{MirrorOneAxis33} and \lean{MirrorNoAxis34});
  \lean{halfTurn_of_symmetric_counterexample} (Theorem~\ref{sr:thm:variant}, conditional on the same classes with
  \lean{Fourfold32} in place of \lean{HalfTurn34});
  \lean{mirrorOneAxis33_of_exclusionLine} (Lemma~\ref{sr:lem:mirror-line}, an implication between two propositions,
  so conditional on its hypothesis \lean{MirrorExclusionLineCopy}).
\end{itemize}
The local statement \lean{pigeon} in the proof of \lean{isometry_structure} is Lemma~\ref{sr:lem:pigeon}. The Mathlib
facts on minimal periods and on roots of unity that the Lean proofs use are restated and proved in
Lemmas~\ref{sr:lem:minper} and~\ref{sr:lem:prim}; the Mathlib results on isometries that the Lean proof uses (the
construction \lean{Isometry.affineIsometryOfStrictConvexSpace} and the theorem \lean{linear_isometry_complex}) are
replaced by the self-contained proof of Theorem~\ref{sr:thm:structure}(a); other standard facts (the pigeonhole
principle, least prime factors, square roots of complex numbers, elementary facts on finite sets) are stated where they
are used. Not covered,
deliberately:
\dirn{SYMLEAN/SymmetricReductionAxioms.lean} (a verbatim copy with two \texttt{\#print axioms} lines); lines 1--1058 of
\dirn{SYMLEAN/AffineReduction.lean} (a verbatim copy; its affine part belongs to the section on affine symmetries);
the review files \dirn{SYMLEAN/review/ReviewSanity.lean}, \dirn{SYMLEAN/review/VariantChecks.lean} and
\dirn{SYMLEAN/review/SupMetric.lean} (review and audit files, cited only in the paragraph on axioms and rigor and in
Remarks~\ref{sr:rem:sup}, \ref{sr:rem:variant} and~\ref{sr:rem:mirror-line}); and \dirn{SYMLEAN/mirror} (the section
on mirror bounds). Two results of the review files go beyond audit checks and are not restated as numbered
statements: the group form \lean{VariantChecks.group_form} (\dirn{SYMLEAN/review/VariantChecks.lean:1168}, conditional
on the five classes of Theorem~\ref{sr:thm:variant}), and the converse bridge
\lean{Review.exclusionLine_of_mirrorOneAxis33} (\dirn{SYMLEAN/review/ReviewSanity.lean:1593}, an implication between
two propositions, so conditional on its hypothesis \lean{MirrorOneAxis33}) with the equivalence
\lean{Review.mirrorOneAxis33_iff_exclusionLine} of the two propositions (line~1629 of the same file), which asserts
neither of them. The reason is that they lie in review files and that no statement of this document uses them.
Remarks~\ref{sr:rem:variant} and~\ref{sr:rem:mirror-line} state them, name
their logs and give short written arguments.

\section{Affine symmetries}\label{ar:sec}
\begingroup\sloppy\emergencystretch=\linewidth

This section extends the reduction of Theorem~\ref{sr:thm:main} from Euclidean isometries to affine symmetries. Its
central statement is a transfer principle (Theorem~\ref{ar:thm:transfer}). Let $N\ge3$ and $n$ be natural numbers, and
suppose that every finite set in general position with at least $N$ points that is mapped into itself by a Euclidean
isometry other than the identity contains $n$ points in convex position. Then the same holds with an injective affine
map $p\mapsto Lp+b$ in place of the isometry. The Euclidean statement is a \emph{hypothesis} of the transfer (a
parameter of the theorem), so the transfer is an implication and asserts neither statement on its own.
Corollary~\ref{ar:cor:transfer-forms} restates the transfer for Mathlib affine maps, for affine equivalences and for
maps given in coordinates. With $N=33$, $n=7$ and the Euclidean statement supplied by Theorem~\ref{sr:thm:main}, the
transfer gives five forms of the statement ``a set of at least $33$ points in general
position with a non-trivial affine symmetry contains $7$ points in convex position'' (Theorem~\ref{ar:thm:main}). Like
Theorem~\ref{sr:thm:main}, these five theorems take five of the six case classes of Definition~\ref{sr:def:classes}
(all except \lean{Fourfold32}) as hypotheses, so they are conditional (Remark~\ref{ar:rem:classes}).

The idea is as follows. The symmetry $g(p)=Lp+b$ permutes the finite set $S$, and $S$ contains three non-collinear
points, because it is in general position and has at least $N\ge3$ points. So some power $g^m$ with $m\ge1$ fixes these
three points, and then its linear part satisfies $L^m=1$. The averaged quadratic form $z\mapsto\sum_{k<m}\operatorname{nsq}(L^kz)$, where $\operatorname{nsq}$ is the squared
Euclidean norm, is $L$-invariant and positive definite. Writing it as $\operatorname{nsq}(Tz)$ with an explicit upper
triangular matrix $T$ (a Cholesky factor) turns $g$ into the Euclidean isometry $T\circ g\circ T^{-1}$ of the set
$T(S)$. Finally, $T$ preserves general position and cardinality, and $T^{-1}$ carries $n$ points of $T(S)$ in convex
position back to $n$ points of $S$ in convex position.

\paragraph{Sources.}
All declarations of this section are in the namespace \lean{ErdosSzekeres.SymReduction} of the file
\dirn{SYMLEAN/AffineReduction.lean} (SHA-256 \texttt{41cfc0b60aec}\dots), lines 1059--1533 (its \S\S9--10).
\begin{itemize}
\item Lines 1--1058 of that file are byte-identical to lines 1--1058 of \dirn{SYMLEAN/SymmetricReduction.lean}
(SHA-256 \texttt{6d5ec90a4215}\dots). Their declarations, among them \lean{eDist2}, \lean{IsEucIsometry},
\lean{Invariant}, the case classes and \lean{heptagon_of_nontrivial_isometry} (line 997 in both files), are treated in
the section of Theorem~\ref{sr:thm:main}; the theorems of this section use the copies in
\dirn{SYMLEAN/AffineReduction.lean}.
\item The repository definitions \lean{Point}, \lean{GeneralPosition}, \lean{ConvexPosition} and \lean{HasConvexNGon}
are in \dirn{ErdosSzekeres/Definitions.lean} (SHA-256 \texttt{2425500eefd1}\dots), and \lean{orient} is in
\dirn{ErdosSzekeres/Orientation.lean:9} (SHA-256 \texttt{f67955ac7091}\dots).
\item Two reformulations of these definitions are used. The first is \lean{convexPosition_iff}, at
\dirn{ErdosSzekeres/Basic.lean:15} (file SHA-256 \texttt{13667aa2d3dc}\dots). The second is
\lean{generalPosition_iff_orient_ne_zero}, at \dirn{ErdosSzekeres/OrientationGeometry.lean:66} (file SHA-256
\texttt{c7db9bbbb65e}\dots).
\item Toolchain: Lean 4.24.0, Mathlib \texttt{f897ebcf}.
\end{itemize}

\paragraph{Rigor and axioms.}
Every statement below that carries a Lean reference is \emph{Lean-checked}. The evidence is as follows.
\begin{itemize}
\item The file \dirn{SYMLEAN/AffineReductionAxioms.lean} (SHA-256 \texttt{f6cb9372b174}\dots) consists of the 1533
lines of \dirn{SYMLEAN/AffineReduction.lean}, followed by a blank line, nine \texttt{\#print axioms} lines and one
\texttt{\#check} line.
\item The run table \dirn{SYMLEAN/logs/runs.tsv} records that \dirn{SYMLEAN/AffineReduction.lean} and
\dirn{SYMLEAN/AffineReductionAxioms.lean} were each compiled on 1 November 2026 with exit code 0.
\item The recorded output \dirn{SYMLEAN/logs/AffineReductionAxioms.out} (SHA-256 \texttt{38d0b6f59398}\dots) of the
axioms file states
that each of the following nine theorems depends on exactly the axioms \texttt{propext}, \texttt{Classical.choice} and
\texttt{Quot.sound}: \lean{heptagon_of_nontrivial_affineEquiv}, \lean{heptagon_of_nontrivial_affineEquiv_image},
\lean{heptagon_of_nontrivial_affineMap}, \lean{heptagon_of_nontrivial_affine},
\lean{heptagon_of_nontrivial_affine_coords}, \lean{hasConvexNGon_of_affine_symmetry},
\lean{generalPosition_image_affineEquiv}, \lean{hasConvexNGon_image_affine} and
\lean{heptagon_of_nontrivial_isometry}.
\item Every other declaration of this section is used, directly or through other declarations, in the proof of one of
these nine theorems. So it depends on no further axiom. For \lean{lin2_fst} and \lean{lin2_snd}
(Lemma~\ref{ar:lem:lin2}(a)) this inference needs a supplement. Their Lean proofs are \lean{rfl}, and they are used only inside
\texttt{simp only} (\dirn{SYMLEAN/AffineReduction.lean}, lines 1180, 1187, 1193 and 1197), which may apply such lemmas as definitional rewrites without
recording them in the proof term. So they need not occur in the proof terms from which \texttt{\#print axioms}
collects the axioms of the nine theorems. Apart from equality and \lean{rfl}, which use no axiom, their statements and
proofs mention only \lean{lin2} applied to points and the coordinates, additions and multiplications in the definition
of \lean{lin2}. The axioms check of \lean{hasConvexNGon_of_affine_symmetry} reaches all of these, because its proof applies
the map $T$ of Proposition~\ref{ar:prop:averaging}(b), a value of \lean{lin2}, to points (Lean lines 1373--1376), and the
check also follows definitions. So these two lemmas depend on no
further axiom either.
\item The compilation log \dirn{SYMLEAN/logs/AffineReduction.out} of \dirn{SYMLEAN/AffineReduction.lean} is empty
(no message). So no declaration of the file uses \lean{sorry}.
\item The five case classes are definitions of type \lean{Prop} that occur as explicit premises, not axioms. The
output of the \texttt{\#check} line, in the same log as the axioms, displays them as premises of
\lean{heptagon_of_nontrivial_affineEquiv}.
\item No proof uses \lean{decide} or \lean{native_decide}. Nothing in this section is a computation.
\end{itemize}
The proofs below are \emph{written proofs (not refereed)}. They follow the Lean proofs step by step unless a proof says
otherwise. Statements tagged ``written step'' are not Lean declarations.

\paragraph{Dependencies.}
From the foundations: Definitions~\ref{fd:def:point}, \ref{fd:def:orient}, \ref{fd:def:gp}, \ref{fd:def:cp},
\ref{fd:def:ngon} and~\ref{fd:def:inv}; Lemma~\ref{fd:lem:collinear} (general position in terms of $\orient$); and
Lemma~\ref{fd:lem:cp-hull} (convex position in terms of convex hulls). The lemmas below re-prove, in the exact form
used here, the facts summarised in Lemmas~\ref{fd:lem:affine-orient} and~\ref{fd:lem:affine-cp}. From the reduction
for isometries: Definition~\ref{sr:def:iso} (\lean{eDist2} and \lean{IsEucIsometry}),
Definition~\ref{sr:def:classes} (the case classes) and Theorem~\ref{sr:thm:main}.

\subsection{Conventions}

\paragraph{Points, matrices, iterates.}
Points are elements of $\R^2$ (Lean \lean{Point}, the functions from the index type \lean{Fin}~$2$ to the reals),
written
$p=(p_0,p_1)$. Points and vectors are the same objects, and $p\pm q$ and $\lambda p$ are taken coordinatewise.
\begin{itemize}
\item A \emph{matrix} is a real $2\times2$ matrix $M=(M_{ij})_{i,j\in\{0,1\}}$ (Lean: a
\lean{Matrix} with rows and columns indexed by \lean{Fin}~$2$ and real entries).
\item $Mz$ is the matrix--vector product, $(Mz)_i=\sum_{j\in\{0,1\}}M_{ij}z_j$ (Lean \lean{Matrix.mulVec}). It is
linear in $z$, and $M(M'z)=(MM')z$ for matrices $M,M'$.
\item $M^k$ is the $k$-th matrix power, with $M^0=1$, the identity matrix, and $M^{k+1}=M^kM=MM^k$.
\item $\det M=M_{00}M_{11}-M_{01}M_{10}$ (Mathlib \lean{Matrix.det_fin_two}).
\item For a map $g:\R^2\to\R^2$, $g^k$ is the $k$-fold iterate (Lean \texttt{g\textasciicircum[k]}, from \lean{Nat.iterate}): $g^0=\mathrm{id}$ and
$g^{k+1}=g\circ g^k$. (The definition of \lean{Nat.iterate} recurses on the other side, $g^{k+1}=g^k\circ g$. The
form used here, $g^{k+1}(p)=g(g^k(p))$, is the lemma \lean{Function.iterate_succ_apply'}, which the Lean proof of
Lemma~\ref{ar:lem:iter} uses.)
\item ``$g\ne\mathrm{id}$'' is always an inequality of maps of the whole plane.
\end{itemize}

\paragraph{Finite sets.}
A finite set is a Lean \lean{Finset} of points, and $\card S$ is its cardinality. For a finite set $S$ and a map $f$,
$f(S)=\{f(p):p\in S\}$ is again a finite set (Lean \lean{Finset.image}); if $f$ is injective, then $\card{f(S)}=\card S$.
``$f$ maps $S$ into $S$'' means \lean{Invariant}~$f$~$S$: $f(p)\in S$ for all $p\in S$ (Definition~\ref{fd:def:inv};
Lean \dirn{SYMLEAN/AffineReduction.lean:271}). General position, convex position and \lean{HasConvexNGon} are as in
Definitions~\ref{fd:def:gp}, \ref{fd:def:cp} and~\ref{fd:def:ngon}. By Lemma~\ref{fd:lem:collinear}, a finite set $S$
is in general position if and only if $\orient(a,b,c)\ne0$ for all pairwise distinct $a,b,c\in S$. By
Lemma~\ref{fd:lem:cp-hull}, $S$ is in convex position if and only if $p\notin\conv(S\setminus\{p\})$ for every $p\in
S$. Here $\conv(A)$ is the convex hull, the smallest convex set containing $A$. Finally, \lean{HasConvexNGon}~$n$~$S$
means that some $U\subseteq S$ with $\card U=n$ is in convex position.

\paragraph{Euclidean isometries.}
$\operatorname{eDist2}(p,q)=(p_0-q_0)^2+(p_1-q_1)^2$, and a map $h:\R^2\to\R^2$ is a \emph{Euclidean isometry}
(\lean{IsEucIsometry}~$h$) if $\operatorname{eDist2}(h(p),h(q))=\operatorname{eDist2}(p,q)$ for \emph{all}
$p,q\in\R^2$ (Definition~\ref{sr:def:iso}; Lean \dirn{SYMLEAN/AffineReduction.lean:265} and~\dirn{:268}). The squared
distance is written out because the metric that Mathlib puts on \lean{Point} is the sup metric, not the Euclidean
one.

\paragraph{Affine maps (Mathlib).}
\begin{itemize}
\item An \emph{affine map} $f$ of $\R^2$ (Mathlib \lean{AffineMap}) is a map $f:\R^2\to\R^2$ together with a linear
map $f_{\mathrm{lin}}:\R^2\to\R^2$ (Lean \lean{f.linear}) such that $f(v+p)=f_{\mathrm{lin}}(v)+f(p)$ for all
$v,p\in\R^2$. The same letter denotes the map. A linear map $\varphi$ is an affine map with
$\varphi_{\mathrm{lin}}=\varphi$ (Lean \lean{LinearMap.toAffineMap}).
\item An \emph{affine equivalence} $f$ of $\R^2$ (Mathlib \lean{AffineEquiv}) is a bijection $f$ of $\R^2$ together
with a linear bijection $f_{\mathrm{lin}}$ with the same property. Its underlying affine map
(\lean{AffineEquiv.toAffineMap}) has the same map and the same linear part.
\item $\mathrm{refl}$ denotes the identity affine equivalence of $\R^2$, whose map is $\mathrm{id}$ (Lean
\lean{AffineEquiv.refl}, for the reals and \lean{Point}).
\item Two affine equivalences with the same map are equal (\lean{AffineEquiv.ext}), because the linear part is
determined by the map: $f_{\mathrm{lin}}(v)=f(v)-f(0)$, by the defining property with $p=0$. Hence, for an affine
equivalence $g$, $g\ne\mathrm{refl}$ holds if and only if $g\ne\mathrm{id}$ as maps.
\item For a linear map $\varphi$ of $\R^2$, $M_\varphi$ is the matrix whose $j$-th column is $\varphi(e_j)$, where
$e_0=(1,0)$ and $e_1=(0,1)$ (Lean \lean{LinearMap.toMatrix'}). Then $M_\varphi z=\varphi(z)$ for every $z$, since
$\varphi(z)=z_0\varphi(e_0)+z_1\varphi(e_1)$ (\lean{LinearMap.toMatrix'_mulVec}). The matrix of a composition is the
product of the matrices.
\end{itemize}

\subsection{The averaged quadratic form}

\begin{definition}[squared norm; \leanref{ErdosSzekeres.SymReduction.nsq}{SYMLEAN/AffineReduction.lean:1088}]\label{ar:def:nsq}
For $v\in\R^2$, $\operatorname{nsq}(v)=v_0^2+v_1^2$.
\end{definition}

\begin{lemma}[coordinates of $Az$; \leanref{ErdosSzekeres.SymReduction.mulVec2_fst}{SYMLEAN/AffineReduction.lean:1097}, \leanref{ErdosSzekeres.SymReduction.mulVec2_snd}{SYMLEAN/AffineReduction.lean:1101}]\label{ar:lem:mulvec}
For every matrix $A$ and every $z\in\R^2$: $(Az)_0=A_{00}z_0+A_{01}z_1$ \textup{(}\lean{mulVec2_fst}\textup{)} and
$(Az)_1=A_{10}z_0+A_{11}z_1$ \textup{(}\lean{mulVec2_snd}\textup{)}.
\end{lemma}

\begin{proof}
This is the definition $(Az)_i=\sum_{j\in\{0,1\}}A_{ij}z_j$, written out for $i=0$ and $i=1$.
\end{proof}

\begin{lemma}[elementary facts on $\operatorname{nsq}$; \leanref{ErdosSzekeres.SymReduction.nsq_nonneg}{SYMLEAN/AffineReduction.lean:1090}, \leanref{ErdosSzekeres.SymReduction.eDist2_eq_nsq}{SYMLEAN/AffineReduction.lean:1094}]\label{ar:lem:nsq}
The following hold.
\begin{enumerate}[label=\textup{(\alph*)}]
\item \textup{(}\lean{nsq_nonneg}\textup{)} $0\le\operatorname{nsq}(v)$ for every $v\in\R^2$.
\item \textup{(}\lean{eDist2_eq_nsq}\textup{)} $\operatorname{eDist2}(p,q)=\operatorname{nsq}(p-q)$ for all
$p,q\in\R^2$.
\end{enumerate}
\end{lemma}

\begin{proof}
(a) $\operatorname{nsq}(v)$ is a sum of two squares. (b) $(p-q)_i=p_i-q_i$, so
$\operatorname{nsq}(p-q)=(p_0-q_0)^2+(p_1-q_1)^2=\operatorname{eDist2}(p,q)$.
\end{proof}

\begin{definition}[Gram entries; \leanref{ErdosSzekeres.SymReduction.gA}{SYMLEAN/AffineReduction.lean:1106}, \leanref{ErdosSzekeres.SymReduction.gB}{SYMLEAN/AffineReduction.lean:1107}, \leanref{ErdosSzekeres.SymReduction.gC}{SYMLEAN/AffineReduction.lean:1108}]\label{ar:def:gram}
For a matrix $A$,
\[
g_A(A)=A_{00}^2+A_{10}^2,\qquad g_B(A)=A_{00}A_{01}+A_{10}A_{11},\qquad g_C(A)=A_{01}^2+A_{11}^2
\]
(Lean \lean{gA}, \lean{gB}, \lean{gC}). These are the entries of the Gram matrix of the two columns of $A$.
\end{definition}

\begin{lemma}[squared norm of $Az$; \leanref{ErdosSzekeres.SymReduction.nsq_mulVec}{SYMLEAN/AffineReduction.lean:1110}]\label{ar:lem:nsq-mulvec}
For every matrix $A$ and every $z\in\R^2$,
$\operatorname{nsq}(Az)=g_A(A)\,z_0^2+2\,g_B(A)\,z_0z_1+g_C(A)\,z_1^2$.
\end{lemma}

\begin{proof}
By Lemma~\ref{ar:lem:mulvec}, $\operatorname{nsq}(Az)=(A_{00}z_0+A_{01}z_1)^2+(A_{10}z_0+A_{11}z_1)^2$. Expanding gives
$(A_{00}^2+A_{10}^2)z_0^2+2(A_{00}A_{01}+A_{10}A_{11})z_0z_1+(A_{01}^2+A_{11}^2)z_1^2$.
\end{proof}

\begin{definition}[the averaged form; \leanref{ErdosSzekeres.SymReduction.avgQ}{SYMLEAN/AffineReduction.lean:1116}, \leanref{ErdosSzekeres.SymReduction.qA}{SYMLEAN/AffineReduction.lean:1120}, \leanref{ErdosSzekeres.SymReduction.qB}{SYMLEAN/AffineReduction.lean:1121}, \leanref{ErdosSzekeres.SymReduction.qC}{SYMLEAN/AffineReduction.lean:1122}]\label{ar:def:avgq}
For a matrix $L$, $n\in\N$ and $z\in\R^2$,
\[
\operatorname{avgQ}(L,n,z)=\sum_{k=0}^{n-1}\operatorname{nsq}(L^kz),
\]
\[
q_A(L,n)=\sum_{k=0}^{n-1}g_A(L^k),\qquad q_B(L,n)=\sum_{k=0}^{n-1}g_B(L^k),\qquad q_C(L,n)=\sum_{k=0}^{n-1}g_C(L^k)
\]
(Lean \lean{avgQ}, \lean{qA}, \lean{qB}, \lean{qC}; the sums run over \lean{Finset.range}~$n$, and they are $0$ for
$n=0$).
\end{definition}

\begin{lemma}[coefficients of the averaged form; \leanref{ErdosSzekeres.SymReduction.avgQ_eq}{SYMLEAN/AffineReduction.lean:1124}]\label{ar:lem:avgq-eq}
For every matrix $L$, every $n\in\N$ and every $z\in\R^2$,
\[
\operatorname{avgQ}(L,n,z)=q_A(L,n)\,z_0^2+2\,q_B(L,n)\,z_0z_1+q_C(L,n)\,z_1^2 .
\]
\end{lemma}

\begin{proof}
Apply Lemma~\ref{ar:lem:nsq-mulvec} to $A=L^k$ for each $k<n$ and add the $n$ equalities. On the left this gives
$\operatorname{avgQ}(L,n,z)$. On the right, the coefficients of $z_0^2$, $2z_0z_1$ and $z_1^2$ add up to $q_A(L,n)$,
$q_B(L,n)$ and $q_C(L,n)$. (Lean argues by induction on $n$; each induction step is the addition of one such
equality.)
\end{proof}

\begin{lemma}[positivity of the averaged form; \leanref{ErdosSzekeres.SymReduction.avgQ_ge}{SYMLEAN/AffineReduction.lean:1133}, \leanref{ErdosSzekeres.SymReduction.qA_pos}{SYMLEAN/AffineReduction.lean:1152}, \leanref{ErdosSzekeres.SymReduction.qdisc_pos}{SYMLEAN/AffineReduction.lean:1158}]\label{ar:lem:avgq-pos}
Let $L$ be a matrix and $n\in\N$ with $0<n$. Then:
\begin{enumerate}[label=\textup{(\alph*)}]
\item \textup{(}\lean{avgQ_ge}\textup{)} $\operatorname{nsq}(z)\le\operatorname{avgQ}(L,n,z)$ for every $z\in\R^2$;
\item \textup{(}\lean{qA_pos}\textup{)} $0<q_A(L,n)$;
\item \textup{(}\lean{qdisc_pos}\textup{)} $0<q_A(L,n)\,q_C(L,n)-q_B(L,n)^2$.
\end{enumerate}
\end{lemma}

\begin{proof}
(a) Since $n>0$, the sum defining $\operatorname{avgQ}(L,n,z)$ contains the term $k=0$, which is
$\operatorname{nsq}(L^0z)=\operatorname{nsq}(z)$. The other terms are $\ge0$ by Lemma~\ref{ar:lem:nsq}(a).

Write $\alpha=q_A(L,n)$, $\beta=q_B(L,n)$, $\gamma=q_C(L,n)$.

(b) Apply (a) to $z=(1,0)$ and use Lemma~\ref{ar:lem:avgq-eq}: $1=\operatorname{nsq}(1,0)\le\operatorname{avgQ}(L,n,(1,0))
=\alpha$. Hence $\alpha\ge1>0$.

(c) Apply (a) to $z=(-\beta,\alpha)$ and use Lemma~\ref{ar:lem:avgq-eq}:
\[
\alpha^2+\beta^2=\operatorname{nsq}(-\beta,\alpha)\le\alpha\beta^2+2\beta(-\beta)\alpha+\gamma\alpha^2
=\alpha(\alpha\gamma-\beta^2).
\]
Hence $\alpha(\alpha\gamma-\beta^2)\ge\alpha^2+\beta^2\ge\alpha^2>0$, where the last step uses (b). Since $\alpha>0$, we
get $\alpha\gamma-\beta^2>0$.
\end{proof}

\begin{lemma}[invariance of the averaged form; \leanref{ErdosSzekeres.SymReduction.avgQ_mulVec}{SYMLEAN/AffineReduction.lean:1139}]\label{ar:lem:avgq-inv}
Let $L$ be a matrix and $n\in\N$ with $L^n=1$. Then $\operatorname{avgQ}(L,n,Lz)=\operatorname{avgQ}(L,n,z)$ for every
$z\in\R^2$.
\end{lemma}

\begin{proof}
Put $f(k)=\operatorname{nsq}(L^kz)$ for $k\in\N$. Since $L^k(Lz)=(L^kL)z=L^{k+1}z$, we have
$\operatorname{avgQ}(L,n,Lz)=\sum_{k=0}^{n-1}f(k+1)$. Splitting off the last term, respectively the first term, of
$\sum_{k=0}^{n}f(k)$ gives
\[
\sum_{k=0}^{n-1}f(k)+f(n)=\sum_{k=0}^{n}f(k)=\sum_{k=0}^{n-1}f(k+1)+f(0).
\]
Here $f(n)=\operatorname{nsq}(L^nz)=\operatorname{nsq}(z)=\operatorname{nsq}(L^0z)=f(0)$, because $L^n=1=L^0$. Hence
$\sum_{k=0}^{n-1}f(k+1)=\sum_{k=0}^{n-1}f(k)=\operatorname{avgQ}(L,n,z)$. No positivity of $n$ is needed; for $n=0$
both sides are $0$.
\end{proof}

\subsection{Linear maps with explicit matrices; the Cholesky factor}

\begin{definition}[the linear map $\operatorname{lin2}$; \leanref{ErdosSzekeres.SymReduction.lin2}{SYMLEAN/AffineReduction.lean:1167}]\label{ar:def:lin2}
For $a,b,c,d\in\R$, $\operatorname{lin2}(a,b,c,d):\R^2\to\R^2$ is the $\R$-linear map
$p\mapsto(ap_0+bp_1,\ cp_0+dp_1)$, that is, multiplication by the matrix
$\begin{pmatrix}a&b\\c&d\end{pmatrix}$. In Lean it is a bundled linear map; its additivity and homogeneity are
discharged inside the definition by coordinate computations.
\end{definition}

\begin{lemma}[coordinates and orientation for $\operatorname{lin2}$; \leanref{ErdosSzekeres.SymReduction.lin2_fst}{SYMLEAN/AffineReduction.lean:1174}, \leanref{ErdosSzekeres.SymReduction.lin2_snd}{SYMLEAN/AffineReduction.lean:1175}, \leanref{ErdosSzekeres.SymReduction.orient_lin2}{SYMLEAN/AffineReduction.lean:1177}]\label{ar:lem:lin2}
Let $a,b,c,d\in\R$ and $\Lambda=\operatorname{lin2}(a,b,c,d)$.
\begin{enumerate}[label=\textup{(\alph*)}]
\item \textup{(}\lean{lin2_fst}, \lean{lin2_snd}\textup{)} $(\Lambda p)_0=ap_0+bp_1$ and $(\Lambda p)_1=cp_0+dp_1$ for
every $p\in\R^2$.
\item \textup{(}\lean{orient_lin2}\textup{)} $\orient(\Lambda x,\Lambda y,\Lambda z)=(ad-bc)\orient(x,y,z)$ for all
$x,y,z\in\R^2$.
\end{enumerate}
\end{lemma}

\begin{proof}
(a) This is the definition. (b) By definition, $\orient(X,Y,Z)=(Y-X)_0(Z-X)_1-(Y-X)_1(Z-X)_0$. Put $u=y-x$ and $w=z-x$.
By linearity, $\Lambda y-\Lambda x=\Lambda u$ and $\Lambda z-\Lambda x=\Lambda w$. Hence, by (a),
\[
\orient(\Lambda x,\Lambda y,\Lambda z)=(au_0+bu_1)(cw_0+dw_1)-(cu_0+du_1)(aw_0+bw_1)=(ad-bc)(u_0w_1-u_1w_0),
\]
and $u_0w_1-u_1w_0=\orient(x,y,z)$. This is the case $v=0$ of Lemma~\ref{fd:lem:affine-orient}.
\end{proof}

\begin{lemma}[the Cholesky factor and its inverse; \leanref{ErdosSzekeres.SymReduction.nsq_lin2_chol}{SYMLEAN/AffineReduction.lean:1184}, \leanref{ErdosSzekeres.SymReduction.lin2_chol_inv_left}{SYMLEAN/AffineReduction.lean:1191}, \leanref{ErdosSzekeres.SymReduction.lin2_chol_inv_right}{SYMLEAN/AffineReduction.lean:1195}]\label{ar:lem:chol}
Let $s,t,\beta\in\R$. Put $T=\operatorname{lin2}(s,\beta/s,0,t)$ and $T'=\operatorname{lin2}(1/s,\,-\beta/(s^2t),\,0,\,1/t)$.
\textup(Division follows the Lean convention $x/0=0$, so $T$ and $T'$ are defined for all $s,t,\beta$; each part
below states the non-vanishing it assumes.\textup)
\begin{enumerate}[label=\textup{(\alph*)}]
\item \textup{(}\lean{nsq_lin2_chol}\textup{)} If $s\ne0$, then for every $z\in\R^2$,
\[
\operatorname{nsq}(Tz)=s^2z_0^2+2\beta z_0z_1+\Bigl(\frac{\beta^2}{s^2}+t^2\Bigr)z_1^2 .
\]
\item \textup{(}\lean{lin2_chol_inv_left}\textup{)} If $s\ne0$ and $t\ne0$, then $T'(Tx)=x$ for every $x\in\R^2$.
\item \textup{(}\lean{lin2_chol_inv_right}\textup{)} If $s\ne0$ and $t\ne0$, then $T(T'x)=x$ for every $x\in\R^2$.
\end{enumerate}
\end{lemma}

\begin{proof}
By Lemma~\ref{ar:lem:lin2}(a), $Tz=(sz_0+(\beta/s)z_1,\ tz_1)$ and $T'x=(x_0/s-\beta x_1/(s^2t),\ x_1/t)$.

(a) $\operatorname{nsq}(Tz)=s^2z_0^2+2s(\beta/s)z_0z_1+(\beta/s)^2z_1^2+t^2z_1^2$, and $s(\beta/s)=\beta$ because
$s\ne0$.

(b) The first coordinate of $T'(Tx)$ is $\frac1s\bigl(sx_0+\frac\beta sx_1\bigr)-\frac{\beta}{s^2t}\,tx_1
=x_0+\frac{\beta x_1}{s^2}-\frac{\beta x_1}{s^2}=x_0$. The second is $\frac1t\,tx_1=x_1$.

(c) The first coordinate of $T(T'x)$ is $s\bigl(\frac{x_0}s-\frac{\beta x_1}{s^2t}\bigr)+\frac\beta s\cdot\frac{x_1}t
=x_0-\frac{\beta x_1}{st}+\frac{\beta x_1}{st}=x_0$. The second is $t\cdot\frac{x_1}t=x_1$.
\end{proof}

\subsection{Affine maps in matrix form}

In this subsection, $g:\R^2\to\R^2$ is a map, $L$ a matrix and $b\in\R^2$; ``$g(p)=Lp+b$ for all $p$'' is the Lean
hypothesis \lean{hg}.

\begin{lemma}[differences of iterates; \leanref{ErdosSzekeres.SymReduction.iterate_sub_of_affine}{SYMLEAN/AffineReduction.lean:1199}]\label{ar:lem:iter}
If $g(p)=Lp+b$ for all $p\in\R^2$, then $g^k(p)-g^k(q)=L^k(p-q)$ for all $k\in\N$ and $p,q\in\R^2$.
\end{lemma}

\begin{proof}
Induction on $k$. For $k=0$ both sides equal $p-q$, as $L^0=1$. If the claim holds for $k$, then
\[
g^{k+1}(p)-g^{k+1}(q)=\bigl(Lg^k(p)+b\bigr)-\bigl(Lg^k(q)+b\bigr)=L\bigl(g^k(p)-g^k(q)\bigr)=L\bigl(L^k(p-q)\bigr)
=L^{k+1}(p-q).\qedhere
\]
\end{proof}

\begin{lemma}[periodic points; \leanref{ErdosSzekeres.SymReduction.exists_periodic_of_injective}{SYMLEAN/AffineReduction.lean:1208}]\label{ar:lem:periodic}
Let $g:\R^2\to\R^2$ be injective, let $S$ be a finite set that $g$ maps into itself, and let $q\in S$. Then there is
$m\in\N$ with $0<m$ and $g^m(q)=q$.
\end{lemma}

\begin{proof}
By induction on $i$, $g^i(q)\in S$ for every $i\in\N$: $g^0(q)=q\in S$, and $g^{i+1}(q)=g(g^i(q))\in S$ because $g$
maps $S$ into itself. (In Lean this is \lean{iterate_mem}, \dirn{SYMLEAN/AffineReduction.lean:69}.) Since $\N$ is
infinite and $S$ is finite, the map $i\mapsto g^i(q)$ from $\N$ to $S$ is not injective (pigeonhole). So there are
$i<j$ with $g^i(q)=g^j(q)$. Put $m=j-i>0$. Then $g^i(g^m(q))=g^{i+m}(q)=g^j(q)=g^i(q)$. The map $g^i$ is injective, as
a composition of injective maps, so $g^m(q)=q$.
\end{proof}

\begin{lemma}[a matrix fixing two independent vectors; \leanref{ErdosSzekeres.SymReduction.matrix_eq_one_of_fix}{SYMLEAN/AffineReduction.lean:1217}]\label{ar:lem:fix}
Let $M$ be a matrix and $u,w\in\R^2$ with $u_0w_1-u_1w_0\ne0$, $Mu=u$ and $Mw=w$. Then $M=1$.
\end{lemma}

\begin{proof}
Put $\delta=u_0w_1-u_1w_0$. By Lemma~\ref{ar:lem:mulvec}, the hypotheses say that the four quantities
\[
E_1=M_{00}u_0+M_{01}u_1-u_0,\quad E_2=M_{10}u_0+M_{11}u_1-u_1,\quad
E_3=M_{00}w_0+M_{01}w_1-w_0,\quad E_4=M_{10}w_0+M_{11}w_1-w_1
\]
vanish. Expanding gives the polynomial identities
\[
w_1E_1-u_1E_3=(M_{00}-1)\delta,\quad -w_0E_1+u_0E_3=M_{01}\delta,\quad
w_1E_2-u_1E_4=M_{10}\delta,\quad -w_0E_2+u_0E_4=(M_{11}-1)\delta .
\]
So all four right-hand sides vanish. As $\delta\ne0$, we get $M_{00}=1$, $M_{01}=0$, $M_{10}=0$ and $M_{11}=1$, that
is, $M=1$.
\end{proof}

\begin{lemma}[orientation under affine maps; \leanref{ErdosSzekeres.SymReduction.orient_affine}{SYMLEAN/AffineReduction.lean:1241}]\label{ar:lem:orient-affine}
If $g(p)=Lp+b$ for all $p\in\R^2$, then $\orient(g(x),g(y),g(z))=\det(L)\,\orient(x,y,z)$ for all $x,y,z\in\R^2$.
\end{lemma}

\begin{proof}
We have $g(y)-g(x)=L(y-x)$ and $g(z)-g(x)=L(z-x)$, since the translation $b$ cancels. So the computation in the proof
of Lemma~\ref{ar:lem:lin2}(b), with the entries $L_{00},L_{01},L_{10},L_{11}$ in place of the four parameters of
$\operatorname{lin2}$, gives the factor $L_{00}L_{11}-L_{01}L_{10}=\det L$. This is Lemma~\ref{fd:lem:affine-orient}.
\end{proof}

\begin{lemma}[injectivity; \leanref{ErdosSzekeres.SymReduction.injective_of_affine}{SYMLEAN/AffineReduction.lean:1305}]\label{ar:lem:inj}
If $g(p)=Lp+b$ for all $p\in\R^2$ and $\det L\ne0$, then $g$ is injective.
\end{lemma}

\begin{proof}
Let $g(p)=g(q)$ and put $v=p-q$. Then $Lv=Lp-Lq=(Lp+b)-(Lq+b)=g(p)-g(q)=0$. So
$h_0:=L_{00}v_0+L_{01}v_1=0$ and $h_1:=L_{10}v_0+L_{11}v_1=0$ by Lemma~\ref{ar:lem:mulvec}. Then
\[
(\det L)\,v_0=L_{11}h_0-L_{01}h_1=0,\qquad (\det L)\,v_1=L_{00}h_1-L_{10}h_0=0,
\]
so $v=0$ because $\det L\ne0$. Hence $p=q$.
\end{proof}

\begin{lemma}[matrix form of a Mathlib affine map; \leanref{ErdosSzekeres.SymReduction.affineMap_eq_mulVec}{SYMLEAN/AffineReduction.lean:1284}]\label{ar:lem:matrix-form}
Let $f$ be an affine map of $\R^2$ \textup(Mathlib \lean{AffineMap}\textup) and let $M_f$ be the matrix of its linear
part $f_{\mathrm{lin}}$ \textup(Lean \lean{LinearMap.toMatrix'} applied to \lean{f.linear}\textup). Then $f(p)=M_fp+f(0)$ for every
$p\in\R^2$.
\end{lemma}

\begin{proof}
By the defining property of an affine map with $v=p$ and base point $0$, $f(p)=f(p+0)=f_{\mathrm{lin}}(p)+f(0)$ (in
Mathlib, \lean{AffineMap.decomp}). By the conventions, $f_{\mathrm{lin}}(p)=M_fp$.
\end{proof}

\subsection{Transport of general position and convex position}

\begin{lemma}[general position under orientation-scaling maps; \leanref{ErdosSzekeres.SymReduction.generalPosition_image_of_orient}{SYMLEAN/AffineReduction.lean:1249}]\label{ar:lem:gp-orient}
Let $f:\R^2\to\R^2$ be a map and $D\in\R$ with $D\ne0$ and $\orient(f(x),f(y),f(z))=D\,\orient(x,y,z)$ for all
$x,y,z\in\R^2$. If a finite set $S$ is in general position, then so is $f(S)$.
\end{lemma}

\begin{proof}
By Lemma~\ref{fd:lem:collinear} it suffices to show $\orient(X,Y,Z)\ne0$ for pairwise distinct $X,Y,Z\in f(S)$. Write
$X=f(x)$, $Y=f(y)$, $Z=f(z)$ with $x,y,z\in S$. Then $x,y,z$ are pairwise distinct: for instance $x=y$ would give
$X=Y$. By Lemma~\ref{fd:lem:collinear} again, $\orient(x,y,z)\ne0$, and so
$\orient(X,Y,Z)=D\,\orient(x,y,z)\ne0$. Injectivity of $f$ is not assumed and not needed.
\end{proof}

\begin{lemma}[general position under affine maps; \leanref{ErdosSzekeres.SymReduction.generalPosition_image_affine}{SYMLEAN/AffineReduction.lean:1278}, \leanref{ErdosSzekeres.SymReduction.generalPosition_image_affineEquiv}{SYMLEAN/AffineReduction.lean:1290}]\label{ar:lem:gp-affine}
The following hold.
\begin{enumerate}[label=\textup{(\alph*)}]
\item \textup{(}\lean{generalPosition_image_affine}\textup{)} Let $g:\R^2\to\R^2$, $L$ a matrix and $b\in\R^2$ with
$g(p)=Lp+b$ for all $p$, and $\det L\ne0$. If a finite set $S$ is in general position, then so is $g(S)$.
\item \textup{(}\lean{generalPosition_image_affineEquiv}\textup{)} If a finite set $S$ is in general position and $f$
is an affine equivalence of $\R^2$ \textup(Mathlib \lean{AffineEquiv}\textup), then $f(S)$ is in general position.
\end{enumerate}
\end{lemma}

\begin{proof}
(a) By Lemma~\ref{ar:lem:orient-affine}, the map $g$ satisfies the hypothesis of Lemma~\ref{ar:lem:gp-orient} with
$D=\det L\ne0$.

(b) By Lemma~\ref{ar:lem:matrix-form}, applied to the underlying affine map, $f(p)=Mp+f(0)$ for all $p$, where $M$ is
the matrix of $f_{\mathrm{lin}}$. The linear bijection $f_{\mathrm{lin}}$ has a linear inverse $\psi$. The product $MM_\psi$
is the matrix of $f_{\mathrm{lin}}\circ\psi=\mathrm{id}$, so $MM_\psi=1$ and $\det M\cdot\det M_\psi=1$; hence $\det
M\ne0$. (Lean uses \lean{LinearMap.det_toMatrix'} and \lean{LinearEquiv.isUnit_det'}.) Now apply (a).
\end{proof}

\begin{lemma}[convex position under injective affine maps; \leanref{ErdosSzekeres.SymReduction.convexPosition_image_affine}{SYMLEAN/AffineReduction.lean:1263}]\label{ar:lem:cp-affine}
Let $S$ be a finite set in convex position and let $f$ be an affine map of $\R^2$ \textup(Mathlib
\lean{AffineMap}\textup) that is injective. Then $f(S)$ is in convex position.
\end{lemma}

\begin{proof}
\emph{Step 1: $f(\conv A)=\conv(f(A))$ for every $A\subseteq\R^2$} (in Mathlib, \lean{AffineMap.image_convexHull}). By
Lemma~\ref{ar:lem:matrix-form}, $f(p)=Mp+c$ with $M=M_f$ and $c=f(0)$. So for $x,y\in\R^2$ and $\lambda\in[0,1]$,
\[
f\bigl((1-\lambda)x+\lambda y\bigr)=(1-\lambda)Mx+\lambda My+\bigl((1-\lambda)+\lambda\bigr)c
=(1-\lambda)f(x)+\lambda f(y).
\]
Hence images and preimages of convex sets under $f$ are convex. Let $C$ be convex. If $X=f(x)$ and $Y=f(y)$ with
$x,y\in C$, then $(1-\lambda)x+\lambda y\in C$, so $(1-\lambda)X+\lambda Y=f((1-\lambda)x+\lambda y)\in f(C)$. Let
$C'$ be convex. If $f(x),f(y)\in C'$, then $f((1-\lambda)x+\lambda y)=(1-\lambda)f(x)+\lambda f(y)\in C'$, so
$(1-\lambda)x+\lambda y\in f^{-1}(C')$. Now $f(\conv A)$ is a convex set containing $f(A)$, so
it contains $\conv(f(A))$. And $f^{-1}(\conv(f(A)))$ is a convex set containing $A$, so it contains $\conv A$; that is,
$f(\conv A)\subseteq\conv(f(A))$.

\emph{Step 2.} By Lemma~\ref{fd:lem:cp-hull} we must show $P\notin\conv(f(S)\setminus\{P\})$ for every $P\in f(S)$.
Write $P=f(q)$ with $q\in S$. Since $f$ is injective, $f(S)\setminus\{f(q)\}=f(S\setminus\{q\})$. Indeed, if $r\in S$
and $f(r)\ne f(q)$ then $r\ne q$; conversely, if $r\in S$ and $r\ne q$ then $f(r)\ne f(q)$. Suppose
$f(q)\in\conv(f(S)\setminus\{f(q)\})$. By this identity and Step~1, this set is
$\conv(f(S\setminus\{q\}))=f(\conv(S\setminus\{q\}))$. So $f(q)=f(r)$ for some $r\in\conv(S\setminus\{q\})$.
Injectivity gives $q=r\in\conv(S\setminus\{q\})$, which contradicts the convex position of $S$
(Lemma~\ref{fd:lem:cp-hull}).
\end{proof}

\begin{lemma}[convex $n$-gons under injective affine maps; \leanref{ErdosSzekeres.SymReduction.hasConvexNGon_image_affine}{SYMLEAN/AffineReduction.lean:1298}]\label{ar:lem:ngon-affine}
Let $n\in\N$, let $S$ be a finite set with \lean{HasConvexNGon}~$n$~$S$, and let $f$ be an injective affine map of $\R^2$
\textup(Mathlib \lean{AffineMap}\textup). Then \lean{HasConvexNGon}~$n$ holds for $f(S)$.
\end{lemma}

\begin{proof}
Take $U\subseteq S$ with $\card U=n$ in convex position. Then $f(U)\subseteq f(S)$, $\card{f(U)}=\card U=n$ since $f$ is
injective, and $f(U)$ is in convex position by Lemma~\ref{ar:lem:cp-affine}.
\end{proof}

\subsection{The transfer theorem}

\begin{definition}[the Euclidean statement $\mathcal E(N,n)$; written step]\label{ar:def:E}
For $N,n\in\N$, $\mathcal E(N,n)$ is the following proposition. \emph{For every finite set $S'\subseteq\R^2$ and every
map $h:\R^2\to\R^2$: if $S'$ is in general position, $N\le\card{S'}$, $h$ is a Euclidean isometry \textup(that is,
$\operatorname{eDist2}(h(p),h(q))=\operatorname{eDist2}(p,q)$ for all $p,q\in\R^2$\textup), $h\ne\mathrm{id}$, and
$h(p)\in S'$ for all $p\in S'$, then \lean{HasConvexNGon}~$n$~$S'$.} This is the type of the hypothesis \lean{hE} of
Theorem~\ref{ar:thm:transfer} and Corollary~\ref{ar:cor:transfer-forms}, up to two changes of notation: the bound
variables, \lean{S} and \lean{g} in Lean, are called $S'$ and $h$ here, and \lean{IsEucIsometry} and \lean{Invariant} are
unfolded. The quantifiers, the five premises and their order, and the conclusion are those of Lean. It has no name in
Lean.
\end{definition}

The next proposition collects the steps of the proof of Theorem~\ref{ar:thm:transfer} that construct the Euclidean
isometry. It is stated separately because the remark on affine symmetries in \dirn{WRITEUP/P1/p1.tex} (Remark
\texttt{rem:affine}) uses this construction.

\begin{proposition}[averaging an affine symmetry; written step]\label{ar:prop:averaging}
Let $S$ be a finite set in general position with $\card S\ge3$. Let $L$ be a matrix, $b\in\R^2$ and
$g:\R^2\to\R^2$ with $g(p)=Lp+b$ for all $p$. Assume that $g$ is injective and maps $S$ into itself.
\begin{enumerate}[label=\textup{(\alph*)}]
\item There is an integer $m\ge1$ with $L^m=1$ and $g^m=\mathrm{id}$.
\item Fix any integer $m\ge1$ with $L^m=1$, and put $\alpha=q_A(L,m)$, $\beta=q_B(L,m)$, $\gamma=q_C(L,m)$. Then
$\alpha>0$ and $\alpha\gamma-\beta^2>0$. Put
\[
s=\sqrt\alpha,\qquad t=\sqrt{(\alpha\gamma-\beta^2)/\alpha},\qquad
T=\operatorname{lin2}(s,\beta/s,0,t),\qquad T'=\operatorname{lin2}(1/s,\,-\beta/(s^2t),\,0,\,1/t).
\]
Then $s>0$ and $t>0$. The maps $T$ and $T'$ are mutually inverse linear bijections of $\R^2$. Moreover
$\orient(Tx,Ty,Tz)=st\cdot\orient(x,y,z)$ for all $x,y,z$, with $st>0$, and
$\operatorname{nsq}(Tz)=\operatorname{avgQ}(L,m,z)$ for every $z$.
\item With $m$, $T$, $T'$ as in \textup{(b)}, the map $h=T\circ g\circ T'$ is a Euclidean isometry, and $h$ maps
$T(S)$ into $T(S)$.
\item $h^k=T\circ g^k\circ T'$ for every $k\in\N$. Hence $h^k=\mathrm{id}$ if and only if $g^k=\mathrm{id}$. In
particular, $h\ne\mathrm{id}$ if and only if $g\ne\mathrm{id}$, and $h$ and $g$ have the same order (the least $k\ge1$
with $g^k=\mathrm{id}$, which exists by \textup{(a)}).
\end{enumerate}
\end{proposition}

\begin{proof}
(a) Since $\card S\ge3$, $S$ contains pairwise distinct points $p_0,p_1,p_2$. By general position and
Lemma~\ref{fd:lem:collinear}, $\delta:=\orient(p_0,p_1,p_2)\ne0$. By Lemma~\ref{ar:lem:periodic} there are
$m_0,m_1,m_2\ge1$ with $g^{m_i}(p_i)=p_i$. Put $m=m_0m_1m_2\ge1$. For each $i$, $m=r_im_i$ for an integer $r_i\ge1$. By
induction on $r$, $g^{rm_i}(p_i)=p_i$: for $r=0$ this is $g^0(p_i)=p_i$, and
$g^{(r+1)m_i}(p_i)=g^{m_i}(g^{rm_i}(p_i))=g^{m_i}(p_i)=p_i$. So $g^m(p_i)=p_i$ for $i=0,1,2$.

Put $u=p_1-p_0$ and $w=p_2-p_0$. By Lemma~\ref{ar:lem:iter}, $L^mu=g^m(p_1)-g^m(p_0)=p_1-p_0=u$, and likewise
$L^mw=w$. By the definition of $\orient$, $u_0w_1-u_1w_0=\orient(p_0,p_1,p_2)=\delta\ne0$. So Lemma~\ref{ar:lem:fix}
gives $L^m=1$.

Finally, for every $p\in\R^2$, Lemma~\ref{ar:lem:iter} gives $g^m(p)-g^m(p_0)=L^m(p-p_0)=p-p_0$. Since
$g^m(p_0)=p_0$, this gives $g^m(p)=p$. So $g^m=\mathrm{id}$. (The Lean proof derives $L^m=1$ in the same way, as its
step \lean{hLm}. It does not derive or use $g^m=\mathrm{id}$.)

(b) By Lemma~\ref{ar:lem:avgq-pos}(b),(c) with $n=m\ge1$, $\alpha>0$ and $\alpha\gamma-\beta^2>0$. So
$(\alpha\gamma-\beta^2)/\alpha>0$; hence $s>0$, $t>0$, $s^2=\alpha$ and $t^2=(\alpha\gamma-\beta^2)/\alpha$. As $s\ne0$
and $t\ne0$, Lemma~\ref{ar:lem:chol}(b),(c) gives $T'(Tx)=x$ and $T(T'x)=x$ for all $x$. So $T$ and $T'$ are mutually
inverse bijections; they are linear by Definition~\ref{ar:def:lin2}. By Lemma~\ref{ar:lem:lin2}(b), applied to
$T=\operatorname{lin2}(s,\beta/s,0,t)$, the orientation factor is $st-(\beta/s)\cdot0=st>0$. Finally, by
Lemma~\ref{ar:lem:chol}(a),
$\operatorname{nsq}(Tz)=s^2z_0^2+2\beta z_0z_1+(\beta^2/s^2+t^2)z_1^2$. Here $s^2=\alpha$ and
\[
\frac{\beta^2}{s^2}+t^2=\frac{\beta^2}{\alpha}+\frac{\alpha\gamma-\beta^2}{\alpha}=\gamma .
\]
So $\operatorname{nsq}(Tz)=\alpha z_0^2+2\beta z_0z_1+\gamma z_1^2=\operatorname{avgQ}(L,m,z)$ by
Lemma~\ref{ar:lem:avgq-eq}.

(c) Let $x,y\in\R^2$ and put $w=T'x-T'y$. Since $g(T'x)-g(T'y)=(LT'x+b)-(LT'y+b)=Lw$ and $T$ is linear,
\begin{align*}
\operatorname{eDist2}(h(x),h(y))
&=\operatorname{nsq}\bigl(T(g(T'x))-T(g(T'y))\bigr) &&\text{Lemma~\ref{ar:lem:nsq}(b)}\\
&=\operatorname{nsq}\bigl(T(Lw)\bigr)=\operatorname{avgQ}(L,m,Lw) &&\text{linearity of $T$; part (b)}\\
&=\operatorname{avgQ}(L,m,w)=\operatorname{nsq}(Tw) &&\text{Lemma~\ref{ar:lem:avgq-inv} ($L^m=1$); part (b)}\\
&=\operatorname{nsq}\bigl(T(T'x)-T(T'y)\bigr)=\operatorname{nsq}(x-y) &&\text{linearity of $T$; $TT'=\mathrm{id}$}\\
&=\operatorname{eDist2}(x,y) &&\text{Lemma~\ref{ar:lem:nsq}(b).}
\end{align*}
So $h$ is a Euclidean isometry. If $X\in T(S)$, write $X=Tp$ with $p\in S$. Then $h(X)=T(g(T'(Tp)))=T(g(p))\in T(S)$,
because $g(p)\in S$.

(d) Induction on $k$. For $k=0$: $h^0=\mathrm{id}=T\circ T'=T\circ g^0\circ T'$. If $h^k=T\circ g^k\circ T'$, then
$h^{k+1}=h\circ h^k=T\circ g\circ(T'\circ T)\circ g^k\circ T'=T\circ g^{k+1}\circ T'$, because $T'\circ T=\mathrm{id}$.

If $g^k=\mathrm{id}$, then $h^k=T\circ T'=\mathrm{id}$. If $h^k=\mathrm{id}$, then for every $p$,
$T(g^k(p))=T(g^k(T'(Tp)))=h^k(Tp)=Tp$, so $g^k(p)=p$ because $T$ is injective. The case $k=1$ gives $h\ne\mathrm{id}$
if and only if $g\ne\mathrm{id}$. Since $h^k=\mathrm{id}\Leftrightarrow g^k=\mathrm{id}$ for every $k$, the least
$k\ge1$ with this property is the same for $g$ and $h$.

\emph{Relation to Lean.} Parts (b) and (c) are steps of the Lean proof of
\lean{hasConvexNGon_of_affine_symmetry}, where $T'$ is called \lean{Ti}. These steps are: the positivity of $\alpha$,
of $\alpha\gamma-\beta^2$, of $s$ and of $t$; the inverse identities \lean{hTiT} and \lean{hTTi}, and the injectivity of
$T$ and $T'$ derived from them (\lean{hTinj}, \lean{hTiinj}); the identity \lean{hnT}; the orientation factor
$st-(\beta/s)\cdot0$ of $T$, which Lean uses only inside the step \lean{hgp'} (general position of $T(S)$, through
\lean{generalPosition_image_of_orient} and \lean{orient_lin2}); and the steps \lean{hiso} and \lean{hinv'}. Of
part (d), only the implication $g\ne\mathrm{id}\Rightarrow h\ne\mathrm{id}$ is in Lean (step \lean{hne'}).
\end{proof}

\begin{theorem}[affine-to-Euclidean transfer; \leanref{ErdosSzekeres.SymReduction.hasConvexNGon_of_affine_symmetry}{SYMLEAN/AffineReduction.lean:1330}]\label{ar:thm:transfer}
Let $N,n\in\N$. Assume
\begin{itemize}
\item[\textup{(hE)}] for every finite set $S'\subseteq\R^2$ and every map $h:\R^2\to\R^2$: if $S'$ is in general
position, $N\le\card{S'}$, $h$ is a Euclidean isometry \textup(\lean{IsEucIsometry}~$h$: $\operatorname{eDist2}(h(p),
h(q))=\operatorname{eDist2}(p,q)$ for all $p,q\in\R^2$\textup), $h\ne\mathrm{id}$, and $h(p)\in S'$ for all $p\in S'$,
then \lean{HasConvexNGon}~$n$~$S'$
\end{itemize}
\textup(that is, assume $\mathcal E(N,n)$ of Definition~\textup{\ref{ar:def:E}}\textup), and assume $3\le N$. Let $S$
be a finite set, $g:\R^2\to\R^2$ a map, $L$ a matrix and $b\in\R^2$, such that $g(p)=Lp+b$ for every $p\in\R^2$, $g$ is
injective, $S$ is in general position, $N\le\card S$, $g\ne\mathrm{id}$, and $g(p)\in S$ for every $p\in S$. Then
\lean{HasConvexNGon}~$n$~$S$: some $n$ points of $S$ are in convex position.
\end{theorem}

\begin{proof}
Since $\card S\ge N\ge3$, Proposition~\ref{ar:prop:averaging} applies. Fix $m\ge1$ with $L^m=1$ (part (a)), and let
$T$, $T'$ and $h=T\circ g\circ T'$ be as in parts (b) and (c). We check the hypotheses of (hE) for $S'=T(S)$ and $h$.
\begin{itemize}
\item $T(S)$ is in general position, by Lemma~\ref{ar:lem:gp-orient} with $f=T$ and $D=st\ne0$
(Proposition~\ref{ar:prop:averaging}(b)).
\item $N\le\card S=\card{T(S)}$, because $T$ is injective ($T'\circ T=\mathrm{id}$).
\item $h$ is a Euclidean isometry and maps $T(S)$ into $T(S)$ (Proposition~\ref{ar:prop:averaging}(c)).
\item $h\ne\mathrm{id}$, by Proposition~\ref{ar:prop:averaging}(d), since $g\ne\mathrm{id}$.
\end{itemize}
So (hE) gives $U\subseteq T(S)$ with $\card U=n$ in convex position. Put $V=T'(U)$.
\begin{itemize}
\item $V\subseteq S$: every $y\in V$ is $y=T'(x)$ with $x\in U\subseteq T(S)$, so $x=Tp$ with $p\in S$, and
$y=T'(Tp)=p\in S$.
\item $\card V=\card U=n$, because $T'$ is injective ($T\circ T'=\mathrm{id}$).
\item $V$ is in convex position, by Lemma~\ref{ar:lem:cp-affine} applied to $T'$. This is a linear map, hence an
affine map (Lean \lean{Ti.toAffineMap}), and it is injective.
\end{itemize}
Hence \lean{HasConvexNGon}~$n$~$S$.
\end{proof}

\begin{remark}[the hypothesis (hE), and what the statement forgets]\label{ar:rem:hE}
The following three points concern the hypotheses of Theorem~\ref{ar:thm:transfer} and what its statement does not
record.
\begin{enumerate}[label=(\arabic*)]
\item The hypothesis (hE) is a parameter of Theorem~\ref{ar:thm:transfer}: the Euclidean statement $\mathcal E(N,n)$
for the \emph{same} $N$ and $n$. The theorem is the implication from $\mathcal E(N,n)$ to the corresponding statement
for injective affine maps. It involves no unproved claim of the project.
\item (hE) is used once, for the set $T(S)$ and the isometry $h=T\circ g\circ T'$. Since (hE) quantifies over
\emph{all} Euclidean isometries other than the identity, the statement does not record which isometry $h$ is. The
written Proposition~\ref{ar:prop:averaging}(d) adds that $h$ has the same order as $g$. This is not part of any Lean
statement. An order-preserving form of the transfer, whose Euclidean hypothesis concerns only isometries of the order
of $g$, is not formalized. The remark on affine symmetries in \dirn{WRITEUP/P1/p1.tex} (Remark \texttt{rem:affine})
needs such a form, and says that it is not formalized.
\item The hypothesis $3\le N$ is used only to find three distinct points of $S$. Injectivity of $g$ is a separate
hypothesis. It follows from $\det L\ne0$ (Lemma~\ref{ar:lem:inj}), and that is how
Corollary~\ref{ar:cor:transfer-forms}(c) supplies it.
\end{enumerate}
\end{remark}

\begin{corollary}[the transfer for other presentations of the map; \leanref{ErdosSzekeres.SymReduction.hasConvexNGon_of_affineMap_symmetry}{SYMLEAN/AffineReduction.lean:1423}, \leanref{ErdosSzekeres.SymReduction.hasConvexNGon_of_affineEquiv_symmetry}{SYMLEAN/AffineReduction.lean:1433}, \leanref{ErdosSzekeres.SymReduction.hasConvexNGon_of_affine_coords_symmetry}{SYMLEAN/AffineReduction.lean:1444}]\label{ar:cor:transfer-forms}
Let $N,n\in\N$. Assume \textup{(hE)}, exactly as in Theorem~\textup{\ref{ar:thm:transfer}} \textup(the statement
$\mathcal E(N,n)$\textup), and assume $3\le N$. Let $S$ be a finite set in general position with $N\le\card S$.
\begin{enumerate}[label=\textup{(\alph*)}]
\item \textup{(}\lean{hasConvexNGon_of_affineMap_symmetry}\textup{)} If $f$ is an affine map of $\R^2$ \textup(Mathlib
\lean{AffineMap}\textup) that is injective, $f\ne\mathrm{id}$ as maps, and $f(p)\in S$ for all $p\in S$, then
\lean{HasConvexNGon}~$n$~$S$.
\item \textup{(}\lean{hasConvexNGon_of_affineEquiv_symmetry}\textup{)} If $f$ is an affine equivalence of $\R^2$
\textup(Mathlib \lean{AffineEquiv}\textup) with $f\ne\mathrm{refl}$ and $f(p)\in S$ for all $p\in S$, then
\lean{HasConvexNGon}~$n$~$S$.
\item \textup{(}\lean{hasConvexNGon_of_affine_coords_symmetry}\textup{)} If $g:\R^2\to\R^2$ is a map and
$a,b,c,d,e,f\in\R$ \textup(here $f$ is a real number\textup) satisfy $ad-bc\ne0$,
$g(p)=(ap_0+bp_1+e,\ cp_0+dp_1+f)$ for every $p\in\R^2$, $g\ne\mathrm{id}$, and $g(p)\in S$ for all $p\in S$, then
\lean{HasConvexNGon}~$n$~$S$.
\end{enumerate}
\end{corollary}

\begin{proof}
(a) By Lemma~\ref{ar:lem:matrix-form}, $f(p)=M_fp+f(0)$ for every $p$. Apply Theorem~\ref{ar:thm:transfer} with $g=f$,
$L=M_f$ and $b=f(0)$; the remaining hypotheses are those of (a).

(b) The underlying affine map of $f$ has the same map, which is injective because $f$ is a bijection. It differs from
$\mathrm{id}$: if the map $f$ were $\mathrm{id}$, then $f$ and $\mathrm{refl}$ would have the same map, so
$f=\mathrm{refl}$ by \lean{AffineEquiv.ext}, a contradiction. Now apply (a).

(c) Let $L=\begin{pmatrix}a&b\\c&d\end{pmatrix}$ and $b'=(e,f)$. By Lemma~\ref{ar:lem:mulvec},
$Lp+b'=(ap_0+bp_1+e,\ cp_0+dp_1+f)=g(p)$ for every $p$. Moreover $\det L=ad-bc\ne0$ (Mathlib
\lean{Matrix.det_fin_two_of}), so $g$ is injective by Lemma~\ref{ar:lem:inj}. Apply Theorem~\ref{ar:thm:transfer}.
\end{proof}

\subsection{Affine symmetries and convex heptagons}

\begin{theorem}[convex heptagons from affine symmetries, conditional on the five case classes; \leanref{ErdosSzekeres.SymReduction.heptagon_of_nontrivial_affineEquiv}{SYMLEAN/AffineReduction.lean:1501}, \leanref{ErdosSzekeres.SymReduction.heptagon_of_nontrivial_affineEquiv_image}{SYMLEAN/AffineReduction.lean:1510}, \leanref{ErdosSzekeres.SymReduction.heptagon_of_nontrivial_affineMap}{SYMLEAN/AffineReduction.lean:1489}, \leanref{ErdosSzekeres.SymReduction.heptagon_of_nontrivial_affine}{SYMLEAN/AffineReduction.lean:1477}, \leanref{ErdosSzekeres.SymReduction.heptagon_of_nontrivial_affine_coords}{SYMLEAN/AffineReduction.lean:1521}]\label{ar:thm:main}
Assume, as hypotheses, the five propositions \lean{HalfTurn34}, \lean{Threefold33}, \lean{Fivefold35},
\lean{MirrorOneAxis33} and \lean{MirrorNoAxis34} of Definition~\textup{\ref{sr:def:classes}}. Let $S\subseteq\R^2$ be a
finite set in general position with $33\le\card S$.
\begin{enumerate}[label=\textup{(\alph*)}]
\item \textup{(}\lean{heptagon_of_nontrivial_affineEquiv}\textup{)} If $g$ is an affine equivalence of $\R^2$
\textup(Mathlib \lean{AffineEquiv}\textup) with $g\ne\mathrm{refl}$ and $g(p)\in S$ for all $p\in S$, then
\lean{HasConvexNGon}~$7$~$S$.
\item \textup{(}\lean{heptagon_of_nontrivial_affineEquiv_image}\textup{)} If $g$ is an affine equivalence of $\R^2$
with $g\ne\mathrm{refl}$ and $g(S)\subseteq S$ \textup(an inclusion of sets: the image of $S$ under $g$ is contained in
$S$\textup), then \lean{HasConvexNGon}~$7$~$S$.
\item \textup{(}\lean{heptagon_of_nontrivial_affineMap}\textup{)} If $f$ is an affine map of $\R^2$ \textup(Mathlib
\lean{AffineMap}\textup) that is injective, $f\ne\mathrm{id}$ as maps, and $f(p)\in S$ for all $p\in S$, then
\lean{HasConvexNGon}~$7$~$S$.
\item \textup{(}\lean{heptagon_of_nontrivial_affine}\textup{)} If $g:\R^2\to\R^2$ is a map, $L$ a matrix and
$b\in\R^2$, with $g(p)=Lp+b$ for every $p$, $g$ injective, $g\ne\mathrm{id}$, and $g(p)\in S$ for all $p\in S$, then
\lean{HasConvexNGon}~$7$~$S$.
\item \textup{(}\lean{heptagon_of_nontrivial_affine_coords}\textup{)} If $g:\R^2\to\R^2$ is a map and
$a,b,c,d,e,f\in\R$ \textup(in this part $b$ and $f$ are real numbers\textup) satisfy $ad-bc\ne0$,
$g(p)=(ap_0+bp_1+e,\ cp_0+dp_1+f)$ for every $p\in\R^2$, $g\ne\mathrm{id}$, and $g(p)\in S$ for all $p\in S$, then
\lean{HasConvexNGon}~$7$~$S$.
\end{enumerate}
In each case, \lean{HasConvexNGon}~$7$~$S$ means that some $7$ points of $S$ are in convex position.
\end{theorem}

\begin{proof}
By Theorem~\ref{sr:thm:main} (\lean{heptagon_of_nontrivial_isometry}, applied to the five hypotheses), the following
holds. \emph{For every finite set $S'$ and every map $h:\R^2\to\R^2$: if $S'$ is in general position, $33\le\card{S'}$,
$h$ is a Euclidean isometry, $h\ne\mathrm{id}$ and $h$ maps $S'$ into $S'$, then \lean{HasConvexNGon}~$7$~$S'$.} This is
exactly $\mathcal E(33,7)$, the hypothesis (hE) with $N=33$ and $n=7$. Moreover $3\le33$.

(a) Apply Corollary~\ref{ar:cor:transfer-forms}(b).

(b) For $p\in S$ we have $g(p)\in g(S)\subseteq S$, so $g$ maps $S$ into $S$. Now apply (a).

(c) Apply Corollary~\ref{ar:cor:transfer-forms}(a).

(d) Apply Theorem~\ref{ar:thm:transfer}.

(e) Apply Corollary~\ref{ar:cor:transfer-forms}(c).

The Lean proofs are the same applications, each a two-line term. For (a), (c), (d) and (e) they pass
\lean{heptagon_of_nontrivial_isometry}, applied to the five hypotheses \lean{h2}, \lean{h3}, \lean{h5}, \lean{hm1} and
\lean{hm0}, as \lean{hE}, and prove $3\le33$ by \lean{norm_num}. For (b) the Lean proof passes no \lean{hE}: it applies
\lean{heptagon_of_nontrivial_affineEquiv}, that is (a), after turning the inclusion into maps-into.
\end{proof}

\begin{remark}[the five hypotheses]\label{ar:rem:classes}
Theorem~\ref{ar:thm:main} is a conditional statement. Its five hypotheses enter only through
Theorem~\ref{sr:thm:main}. In Lean they are definitions of type \lean{Prop} that occur as explicit premises
(\lean{h2}, \lean{h3}, \lean{h5}, \lean{hm1}, \lean{hm0}), not axioms. In short, each class says that a finite set in
general position contains $7$ points in convex position if it has one of the following: exactly $34$ points and a
half-turn symmetry whose centre is not in the set; exactly $33$ points and a symmetry under the rotation by $2\pi/3$
about a centre not in the set; exactly $35$ points and a symmetry under the rotation by $2\pi/5$ about a centre not in
the set; exactly $33$ points, a reflection symmetry in a line, and exactly one point of the set on that line; exactly
$34$ points, a reflection symmetry, and no point of the set on the line. See Definition~\ref{sr:def:classes} for the
exact statements. In the project they are supplied as follows (\dirn{SYMLEAN/README.md}, \S\S5, 6 and~9;
\dirn{WRITEUP/SYMMETRIC.md}, \S\S1, 3 and~4).
\begin{itemize}
\item \lean{MirrorOneAxis33} and \lean{MirrorNoAxis34} are Lean-checked theorems: \lean{mirrorOneAxis33_proved}
(line 4274) and \lean{mirrorNoAxis34_proved} (line 4282) of \dirn{SYMLEAN/mirror/ReductionThreeClasses.lean}
(Theorem~\ref{mb:thm:classes}). These two theorems are about the copies of the class definitions in that file. It does
not import \dirn{SYMLEAN/AffineReduction.lean}; lines 38--1559 of \dirn{SYMLEAN/mirror/ReductionThreeClasses.lean} are
byte-identical to lines 12--1533 of \dirn{SYMLEAN/AffineReduction.lean}. So the two theorems discharge the premises
\lean{hm1} and \lean{hm0} of Theorem~\ref{ar:thm:main} only through the identity of the two texts. Their axioms log
\dirn{SYMLEAN/mirror/axioms.out} lists only \texttt{propext}, \texttt{Classical.choice} and \texttt{Quot.sound}. They
are obtained without SAT input from the repository's Lean formalization of Baek's theorem
(Theorem~\ref{mb:thm:baek}) and the mirror reduction. The same file restates part (a) with
only the three rotation classes as hypotheses (\lean{heptagon_of_nontrivial_affineEquiv_3}, line 4303;
Theorem~\ref{mb:thm:main3}). Both mirror classes also have SAT certificates (record \dirn{SYMZ}).
\item \lean{HalfTurn34}, \lean{Threefold33} and \lean{Fivefold35} are not Lean-checked. They are computer-assisted
claims of the records \dirn{C2FINAL} (with \dirn{C2C34} and \dirn{C2P}--\dirn{C2P5}), \dirn{C3FINAL} and \dirn{SYM5}.
They rest on four kinds of evidence.
\begin{itemize}
\item \emph{SAT certificates.} Every certificate that these claims use, including those of the pieces of their case
splits, carries a recorded acceptance by the proof checker \texttt{cake\_lpr}. Since a revision of 6 November 2026 this is the
only checker in the project's trust base; the second checker, \texttt{lrat-check}, is a secondary check only
(\dirn{WRITEUP/SYMMETRIC.md}, \S3; \dirn{SYMLEAN/README.md}, \S9, which predates that revision, still names both). For
the parallelogram pieces of \lean{HalfTurn34} and the $116{,}883$ pieces of \lean{Threefold33}, the proof files were
deleted after checking, so for most of these pieces only the recorded verdicts remain. Some of these pieces were
re-solved (\dirn{C2FINAL}, \S\S4, 5 and~9; \dirn{C2RECHECK}; \dirn{C3SPOT}; \dirn{C3PIECES}).
\item \emph{Written proofs, reviewed or audited within the project} (listed in \dirn{WRITEUP/SYMMETRIC.md}, \S3).
Among them are the soundness proofs of the encoders; the reductions A and B of the parallelogram cases of
\lean{HalfTurn34}; the normal form of \lean{Threefold33} (\dirn{C3-S/soundness.md}) and Lemma~5 and Lemma~H6M of its
case split; and, for \lean{Threefold33} and \lean{Fivefold35}, a perturbation argument for centres that lie on a line
through two points of different orbits.
\item \emph{Exact computations} that compose the pieces (\dirn{C2FINAL}, \dirn{C3FINAL}). Two bookkeeping steps are
computer-assisted because their finite checks are kernel evaluations (\texttt{decide +kernel}): the cover of the $910$
threefold roots, which is Lean-checked (Theorem~\ref{cb:thm:c3}), and the closure of the central parallelogram orbits,
which is Lean-checked according to the record (\dirn{C2FINAL/README.md}, \S6; Theorem~\ref{cb:thm:c2}). Lean does not
see the records: that every leaf named by the keys of the certificates of Theorem~\ref{cb:thm:c3}
and every subcase listed in the data of Theorem~\ref{cb:thm:c2} has a record with the checker verdicts and a CNF hash equal to that of the rebuilt formula is an exact computation outside
Lean.
\item \emph{The match between the classes and the formulas.} No formal link joins a Lean class to its CNF formulas.
That each class states what its certificates prove was checked by reading (\dirn{SYMLEAN/README.md}, \S6).
\end{itemize}
The records list the remaining caveats and open checks (\dirn{WRITEUP/SYMMETRIC.md}, \S4).
\end{itemize}
The unconditional statement, Theorem~\ref{ar:thm:main} without its five hypotheses, is claimed in the project as
computer-assisted (Theorem~A of \dirn{WRITEUP/SYMMETRIC.md}, \S0, in its form for affine bijections), with the evidence
listed above; it is not Lean-checked as a whole. This document establishes it only conditionally: as
Theorem~\ref{ar:thm:main} and, for part~(a) with three hypotheses, as Theorem~\ref{mb:thm:main3}.
\end{remark}

\begin{remark}[what the statements of this section do not say]\label{ar:rem:scope}
The statements of this section are limited in four ways.
\begin{enumerate}[label=(\arabic*)]
\item The symmetry is a map of the whole plane: an affine equivalence, an injective affine map, or an injective map
$p\mapsto Lp+b$. A permutation of $S$ that is not the restriction of such a map is not treated.
\item One symmetry is treated at a time. Nothing is stated about the full symmetry group of $S$.
\item The variant without \lean{HalfTurn34} (Theorem~\ref{sr:thm:variant}) has no affine version in the Lean files
(\dirn{SYMLEAN/README.md}, \S9).
\item The order of the symmetry is not tracked (Remark~\ref{ar:rem:hE}(2)).
\end{enumerate}
\end{remark}

\paragraph{Coverage.}
This section covers every declaration of \dirn{SYMLEAN/AffineReduction.lean}, lines 1059--1533:
\begin{itemize}
\item \lean{nsq} (Definition~\ref{ar:def:nsq}); \lean{mulVec2_fst} and \lean{mulVec2_snd}
(Lemma~\ref{ar:lem:mulvec}); \lean{nsq_nonneg} and \lean{eDist2_eq_nsq} (Lemma~\ref{ar:lem:nsq});
\item \lean{gA}, \lean{gB} and \lean{gC} (Definition~\ref{ar:def:gram}); \lean{nsq_mulVec}
(Lemma~\ref{ar:lem:nsq-mulvec});
\item \lean{avgQ}, \lean{qA}, \lean{qB} and \lean{qC} (Definition~\ref{ar:def:avgq}); \lean{avgQ_eq}
(Lemma~\ref{ar:lem:avgq-eq}); \lean{avgQ_ge}, \lean{qA_pos} and \lean{qdisc_pos} (Lemma~\ref{ar:lem:avgq-pos});
\lean{avgQ_mulVec} (Lemma~\ref{ar:lem:avgq-inv});
\item \lean{lin2} (Definition~\ref{ar:def:lin2}); \lean{lin2_fst}, \lean{lin2_snd} and \lean{orient_lin2}
(Lemma~\ref{ar:lem:lin2}); \lean{nsq_lin2_chol}, \lean{lin2_chol_inv_left} and \lean{lin2_chol_inv_right}
(Lemma~\ref{ar:lem:chol});
\item \lean{iterate_sub_of_affine} (Lemma~\ref{ar:lem:iter}); \lean{exists_periodic_of_injective}
(Lemma~\ref{ar:lem:periodic}); \lean{matrix_eq_one_of_fix} (Lemma~\ref{ar:lem:fix}); \lean{orient_affine}
(Lemma~\ref{ar:lem:orient-affine}); \lean{injective_of_affine} (Lemma~\ref{ar:lem:inj}); \lean{affineMap_eq_mulVec}
(Lemma~\ref{ar:lem:matrix-form});
\item \lean{generalPosition_image_of_orient} (Lemma~\ref{ar:lem:gp-orient}); \lean{generalPosition_image_affine} and
\lean{generalPosition_image_affineEquiv} (Lemma~\ref{ar:lem:gp-affine}); \lean{convexPosition_image_affine}
(Lemma~\ref{ar:lem:cp-affine}); \lean{hasConvexNGon_image_affine} (Lemma~\ref{ar:lem:ngon-affine});
\item \lean{hasConvexNGon_of_affine_symmetry} (Theorem~\ref{ar:thm:transfer}, conditional on (hE));
\lean{hasConvexNGon_of_affineMap_symmetry}, \lean{hasConvexNGon_of_affineEquiv_symmetry} and
\lean{hasConvexNGon_of_affine_coords_symmetry} (Corollary~\ref{ar:cor:transfer-forms}, conditional on (hE));
\item \lean{heptagon_of_nontrivial_affineEquiv}, \lean{heptagon_of_nontrivial_affineEquiv_image},
\lean{heptagon_of_nontrivial_affineMap}, \lean{heptagon_of_nontrivial_affine} and
\lean{heptagon_of_nontrivial_affine_coords} (Theorem~\ref{ar:thm:main}, conditional on the case classes
\lean{HalfTurn34}, \lean{Threefold33}, \lean{Fivefold35}, \lean{MirrorOneAxis33} and \lean{MirrorNoAxis34}).
\end{itemize}
Definition~\ref{ar:def:E} and Proposition~\ref{ar:prop:averaging} are written steps, not Lean declarations.

Deliberately not covered here:
\begin{itemize}
\item Lines 1--1058 of \dirn{SYMLEAN/AffineReduction.lean}. They are a verbatim duplicate of lines 1--1058 of
\dirn{SYMLEAN/SymmetricReduction.lean} and are treated in the section of Theorem~\ref{sr:thm:main}. Of them, \lean{iterate_mem}
(line 69) is re-derived inside the proof of Lemma~\ref{ar:lem:periodic}, and \lean{point_ext} (line 502: two points
with equal coordinates are equal) is used tacitly.
\item \dirn{SYMLEAN/AffineReductionAxioms.lean}. It is the same file followed by a blank line, nine
\texttt{\#print axioms} lines and one \texttt{\#check} line, and it is used here only as evidence for the axioms.
\item The verbatim copy of lines 12--1533 inside \dirn{SYMLEAN/mirror/ReductionThreeClasses.lean} (the declarations of
this section sit there at their line number plus 26). The own content of that file is treated
in the section of Theorem~\ref{mb:thm:main3}.
\end{itemize}
\endgroup

\section{The mirror lemma and the mirror reduction (Theorem~A of \texorpdfstring{\dirn{MIRRORETV}}{MIRRORETV})}\label{mm:sec}
% Local typesetting setting, closed by the \endgroup at the end of this file: long Lean names in
% statement heads cannot be hyphenated, so allow loose lines instead of lines running into the margin.
\begingroup\setlength{\emergencystretch}{0.5\linewidth}

This section gives written proofs for the part of the Lean development \lean{MirrorOdd}
(file \dirn{MIRRORETV/lean/MirrorMain.lean}) that concerns point sets symmetric under a reflection:
the mirror lemma, the reduction of a mirror-symmetric set to an $x$-sorted configuration (``Theorem~A''
of \dirn{MIRRORETV/README.md}, in its Lean form), the transport from an arbitrary mirror line to the
$x$-axis, and the bounds for $k=4,5,6$. In this section ``Theorem~A'' always means this mirror reduction.
The main theorem of \dirn{WRITEUP/SYMMETRIC.md}, also called Theorem~A there, is a different statement,
about sets of at least $33$ points with a symmetry; this document states it only in conditional forms,
whose hypotheses are case classes (Theorems~\ref{sr:thm:main}, \ref{ar:thm:main} and~\ref{mb:thm:main3}).
\begin{itemize}
\item \emph{Mirror lemma} (Theorems~\ref{mm:thm:mirror-polygon} and~\ref{mm:thm:axis-polygon}). Let
  $c_0,\dots,c_{m-1}$ have strictly increasing abscissae, and let every triple $c_i,c_j,c_l$ with $i<j<l$
  be clockwise. If $m\ge2$ and all $c_i$ lie strictly above the $x$-axis, then the points $c_i$ and their
  mirror images $\operatorname{mir}(c_i)$ contain a convex $2m$-gon. If $m\ge3$ and the first (or the
  last) point lies on the axis while the others lie strictly above it, they contain a convex
  $(2m-1)$-gon. As in Lean, the proof lists the points in an explicit cyclic order and checks that
  \emph{every} triple of the list, taken in list order, is counterclockwise.
\item \emph{Theorem~A} (Theorems~\ref{mm:thm:odd}, \ref{mm:thm:even} and~\ref{mm:thm:line}). Let $j\ge2$,
  and let $M$ be a natural number such that every $x$-sorted configuration with no $(j+1)$-cap and no
  convex $(2j+1)$-gon has at most $M$ points. Then every finite set in general position that is
  symmetric under the reflection in a line and has no convex $(2j+1)$-gon has at most $2M$ points. In
  the even case ($j$-caps, convex $2j$-gons) the bound is $2M+2$. Here $M$ is an \emph{assumed} upper
  bound: the bound on configurations is a hypothesis of these theorems, and it is not proved in this
  section.
\item \emph{Small $k$} (Theorem~\ref{mm:thm:small}). Without further hypotheses, a finite set in general
  position that is symmetric in a line and has no convex $k$-gon has at most $4$, $8$, $12$ points for
  $k=4,5,6$.
\end{itemize}

\paragraph{Sources.}
{\raggedright
The files are the following (SHA-256, first 12 hexadecimal digits).
\begin{itemize}
\item The declarations treated here: \dirn{MIRRORETV/lean/MirrorMain.lean} (2677 lines,
  \texttt{ca623e55a1b9}\dots, namespace \lean{MirrorOdd}) and, for configurations and caps,
  \dirn{ErdosSzekeres/Omission/Realization.lean} (\texttt{ae4970e0808e}\dots).
\item The positive-cycle lemma (Lemma~\ref{fd:lem:pos-cycle}): \dirn{ErdosSzekeres/ConvexCycleCuts.lean}
  (\texttt{b4ad0ba9e583}\dots).
\item The basic definitions recalled in the foundations section: \dirn{ErdosSzekeres/Definitions.lean}
  (\texttt{2425500eefd1}\dots), \dirn{ErdosSzekeres/Basic.lean} (\texttt{13667aa2d3dc}\dots),
  \dirn{ErdosSzekeres/Orientation.lean} (\texttt{f67955ac7091}\dots) and
  \dirn{ErdosSzekeres/OrientationGeometry.lean} (\texttt{c7db9bbbb65e}\dots).
\end{itemize}
Lines 22--1506 of
\dirn{MIRRORETV/lean/MirrorMain.lean} hold five positive-cycle helpers and nine small polygons
(Lemma~\ref{mm:lem:cyc} and Proposition~\ref{mm:prop:small-polygons}). The file header (lines 4--10)
describes the polygons for caps of four points, says that every orientation inequality is closed by
\lean{linarith} with a certificate found offline by a script \dirn{gen_certs.py}, and lists
\dirn{MirrorPolygons.lean} and \dirn{small_polygons.lean}, generated by a script \dirn{gen_lean.py},
among the source files of the assembled file; the comment at line 993 marks the polygons for caps of two
and three points as generated. Neither script is included in \dirn{MIRRORETV}, and the header lists the
other source files without saying how they were written. Toolchain: Lean~4.24.0, Mathlib \texttt{f897ebcf}. The text of \dirn{MIRRORETV/lean/MirrorMain.lean},
lines 3--2677, occurs verbatim as lines 1562--4236 of \dirn{SYMLEAN/mirror/ReductionThreeClasses.lean}
(\texttt{0cc83994fe8c}\dots): a declaration at line $L$ of \dirn{MIRRORETV/lean/MirrorMain.lean} is at
line $L+1559$ there. That copy has its own record: lines 1--16 of \dirn{SYMLEAN/mirror/axioms.out} print
the same three axioms for the sixteen theorems listed in the Rigor paragraph below, and
\dirn{SYMLEAN/mirror/README.md} reports a compile without errors. With this record in place of
\dirn{MIRRORETV/lean/MirrorMain.out}, everything below applies to both copies.
\dirn{MIRRORETV/lean/FarFlat.lean} (\texttt{08a79f881eea}\dots) and the four Lean files of
\dirn{MIRRORETV/review/} are not used here (see the Coverage paragraph).\par}

\paragraph{Rigor.}
Every statement below that carries a Lean reference is \emph{Lean-checked}, except the two parts marked
as written steps (part~(e) of Lemma~\ref{mm:lem:orient-mir} and part~(b) of Lemma~\ref{mm:lem:slope});
Lemma~\ref{mm:lem:upper-caps} is a written step and carries no Lean reference. The recorded compile log
\dirn{MIRRORETV/lean/MirrorMain.out}, lines 848--863, prints for the sixteen theorems \lean{mirror_four},
\lean{mirror_five}, \lean{mirror_six}, \lean{mirror_seven}, \lean{mirror_eight}, their five
\lean{_line} versions, \lean{mirror_odd_general}, \lean{mirror_even_general},
\lean{mirror_odd_general_line}, \lean{mirror_even_general_line}, \lean{mirror_seven'} and
\lean{mirror_eight'} (all in \lean{MirrorOdd}) exactly the axioms \lean{propext},
\lean{Classical.choice} and \lean{Quot.sound}; the log reports no error, only 121 linter warnings about
unused simp arguments. Every other declaration treated in this section is used, directly or indirectly,
by one of these sixteen theorems, so it depends on no further axiom. The file contains no \lean{sorry},
\lean{admit}, \lean{native_decide} or \lean{axiom} declaration; in the declarations treated here
\lean{decide} only compares explicit elements of finite index types (kernel evaluation). This section
relies on the recorded log; Lean was not re-run for it. The proofs below are written proofs (not
refereed). Theorems~\ref{mm:thm:odd} and~\ref{mm:thm:even} and parts (b), (c) of Theorem~\ref{mm:thm:line} take a
parametric hypothesis (hM): an assumed bound $M$, not an unproved claim of the project; several instances of
(hM) are discharged by Lean-checked theorems (Remark~\ref{mm:rem:hm}).

\paragraph{Dependencies.}
We use the conventions and facts of the foundations section: Definitions~\ref{fd:def:point},
\ref{fd:def:orient}, \ref{fd:def:gp}, \ref{fd:def:cp}, \ref{fd:def:ngon} and~\ref{fd:def:inv}, and
Lemmas~\ref{fd:lem:orient}, \ref{fd:lem:affine-orient}, \ref{fd:lem:collinear}, \ref{fd:lem:cp-hull},
\ref{fd:lem:mono}, \ref{fd:lem:affine-cp} and~\ref{fd:lem:pos-cycle}. The bounds for $k=7,8,9$, which
combine the results of this section with the repository's Lean formalization of Baek's theorem (for
$k=9$ also with the vendored cups--caps theorem), are Theorems~\ref{mb:thm:seven}, \ref{mb:thm:eight}
and~\ref{mb:thm:nine}.

\paragraph{Conventions.}
Points are $p=(p_0,p_1)\in\R^2$ (Lean \lean{Point}, Definition~\ref{fd:def:point}); $p_0$ is the abscissa
and $p_1$ the ordinate. For $m\in\N$ we write $[m]=\{0,1,\dots,m-1\}$ with its usual order; it stands for
Lean's \texttt{Fin m}. All point sets are finite (Lean \texttt{Finset Point}). A finite set $X\subset\R^2$
\emph{has a convex $n$-gon} if \lean{HasConvexNGon} $n$ $X$ holds (Definition~\ref{fd:def:ngon}): some
$T\subseteq X$ with $\card T=n$ is in convex position. A finite set $T$ is \emph{$x$-injective} if
$p\mapsto p_0$ is injective on $T$ (Lean \texttt{Set.InjOn (fun p => p 0) T}). In Lean, $2m-1$, $m-1$
and $K-1$ in the statements below are truncated subtractions of natural numbers; wherever they occur,
the hypotheses ($m\ge3$, $K\ge4$) make them equal to the ordinary differences.

\subsection{Reflection, configurations and sorting}

\begin{definition}[reflection, upper part, axis part; \leanref{MirrorOdd.mir}{MIRRORETV/lean/MirrorMain.lean:16}, \leanref{MirrorOdd.Rset}{MIRRORETV/lean/MirrorMain.lean:1598}, \leanref{MirrorOdd.Aset}{MIRRORETV/lean/MirrorMain.lean:1599}]\label{mm:def:mir}
For $p\in\R^2$ let $\operatorname{mir}(p)=(p_0,-p_1)$, the reflection in the $x$-axis. For a finite set
$S\subset\R^2$ let
\[ R(S)=\{p\in S: p_1>0\},\qquad A(S)=\{p\in S: p_1=0\} \]
(Lean \lean{Rset} and \lean{Aset}). In all theorems below, symmetry of $S$ in the $x$-axis is the
hypothesis ``$\operatorname{mir}(p)\in S$ for every $p\in S$'' (``maps into'', as in
Definition~\ref{fd:def:inv}). Since $\operatorname{mir}$ is injective and $S$ is finite, this is
equivalent to $\operatorname{mir}(S)=S$, but only the stated form is used.
\end{definition}

\begin{lemma}[identities for $\operatorname{mir}$ and $\orient$; \leanref{MirrorOdd.mir_zero}{MIRRORETV/lean/MirrorMain.lean:18}, \leanref{MirrorOdd.mir_one}{MIRRORETV/lean/MirrorMain.lean:19}, \leanref{MirrorOdd.mir_mir}{MIRRORETV/lean/MirrorMain.lean:1543}, \leanref{MirrorOdd.orient_mir3}{MIRRORETV/lean/MirrorMain.lean:2187}, \leanref{MirrorOdd.orient_mir2}{MIRRORETV/lean/MirrorMain.lean:2190}, \leanref{MirrorOdd.orient_cycl}{MIRRORETV/lean/MirrorMain.lean:2193}, \leanref{MirrorOdd.orient_swap12}{MIRRORETV/lean/MirrorMain.lean:2196}, \leanref{MirrorOdd.orient_rev}{MIRRORETV/lean/MirrorMain.lean:2199}]\label{mm:lem:orient-mir}
For all $p,a,b,d\in\R^2$:
\begin{enumerate}[label=(\alph*)]
\item $\operatorname{mir}(p)_0=p_0$, $\operatorname{mir}(p)_1=-p_1$ and $\operatorname{mir}(\operatorname{mir}(p))=p$
  (\lean{mir_zero}, \lean{mir_one}, \lean{mir_mir});
\item $\orient(\operatorname{mir}a,\operatorname{mir}b,\operatorname{mir}d)=-\orient(a,b,d)$ (\lean{orient_mir3});
\item $\orient(\operatorname{mir}a,\operatorname{mir}b,d)=-\orient(a,b,\operatorname{mir}d)$ (\lean{orient_mir2});
\item $\orient(a,b,d)=\orient(b,d,a)$, $\orient(a,b,d)=-\orient(b,a,d)$ and $\orient(a,b,d)=-\orient(d,b,a)$
  (\lean{orient_cycl}, \lean{orient_swap12}, \lean{orient_rev});
\item (written step; not a separate Lean declaration)
  $\orient(a,b,\operatorname{mir}d)=\orient(a,b,d)-2d_1(b_0-a_0)$.
\end{enumerate}
\end{lemma}

\begin{proof}
(a) is the definition of $\operatorname{mir}$. For (b) and (e) expand
$\orient(u,v,w)=(v_0-u_0)(w_1-u_1)-(v_1-u_1)(w_0-u_0)$ (Definition~\ref{fd:def:orient}):
\begin{align*}
\orient(\operatorname{mir}a,\operatorname{mir}b,\operatorname{mir}d)
  &=(b_0-a_0)(-d_1+a_1)-(-b_1+a_1)(d_0-a_0)=-\orient(a,b,d),\\
\orient(a,b,\operatorname{mir}d)
  &=(b_0-a_0)(-d_1-a_1)-(b_1-a_1)(d_0-a_0)=\orient(a,b,d)-2d_1(b_0-a_0).
\end{align*}
(c) By (a), $d=\operatorname{mir}(\operatorname{mir}d)$, so (b) applied to $a,b,\operatorname{mir}d$ gives
$\orient(\operatorname{mir}a,\operatorname{mir}b,d)
=\orient(\operatorname{mir}a,\operatorname{mir}b,\operatorname{mir}(\operatorname{mir}d))
=-\orient(a,b,\operatorname{mir}d)$.
(d) These are the cases of Lemma~\ref{fd:lem:orient}(d) for the cyclic permutation
$(a,b,d)\mapsto(b,d,a)$ (sign $+1$), the transposition of the first two arguments and the transposition
of the first and the third argument (sign $-1$). Each can also be checked by expanding.
\end{proof}

\begin{definition}[configurations, caps and sorting; \leanref{ErdosSzekeres.Omission.Configuration}{ErdosSzekeres/Omission/Realization.lean:9}, \leanref{ErdosSzekeres.Omission.Configuration.HasCap}{ErdosSzekeres/Omission/Realization.lean:64}, \leanref{MirrorOdd.xset}{MIRRORETV/lean/MirrorMain.lean:1549}, \leanref{MirrorOdd.emb}{MIRRORETV/lean/MirrorMain.lean:1553}, \leanref{MirrorOdd.sortPt}{MIRRORETV/lean/MirrorMain.lean:1558}, \leanref{MirrorOdd.cfg}{MIRRORETV/lean/MirrorMain.lean:1584}]\label{mm:def:config}
The following notions concern finite sequences of points sorted by abscissa.
\begin{enumerate}[label=(\alph*)]
\item Let $n\in\N$. A \emph{configuration of $n$ points} is a map $C:[n]\to\R^2$, $i\mapsto C_i$ (Lean
  field \lean{point}), such that $i\mapsto(C_i)_0$ is strictly increasing (field \lean{xOrder}) and
  $\orient(C_i,C_j,C_k)\ne0$ for all $i<j<k$ in $[n]$ (field \lean{generalPosition}). Its \emph{point
  set} is $\mathrm{pts}(C)=\{C_i:i\in[n]\}$ (Lean \texttt{Finset.univ.image c.point}).
\item Let $k\in\N$. A configuration $C$ of $n$ points \emph{has a $k$-cap} (Lean \texttt{c.HasCap k}) if
  there is a strictly increasing map $f:[k]\to[n]$ such that $\orient(C_{f(i)},C_{f(j)},C_{f(l)})<0$ for
  \emph{all} $i<j<l$ in $[k]$.
\item For every finite set $T\subset\R^2$ let $\mathrm{xset}(T)=\{p_0:p\in T\}\subset\R$. From now on
  let $T$ also be $x$-injective. By Lemma~\ref{mm:lem:sort}(a) the set $\mathrm{xset}(T)$ has $\card T$
  elements, so there is a unique strictly increasing map $\mathrm{emb}(T):[\card T]\to\R$ with image
  $\mathrm{xset}(T)$ (Mathlib \lean{Finset.orderEmbOfFin}); $\mathrm{emb}(T)(i)$ is the $(i+1)$-st
  smallest abscissa of a point of $T$. For $i\in[\card T]$, $\mathrm{sortPt}(T)(i)$ is a point $p\in T$
  with $p_0=\mathrm{emb}(T)(i)$. Such a point exists by Lemma~\ref{mm:lem:sort}(b), and Lean selects one
  with the axiom of choice (\lean{Classical.choose}). By $x$-injectivity it is unique, but this is not
  used. (Parts (a) and (b) of Lemma~\ref{mm:lem:sort} concern only $\mathrm{xset}$ and $\mathrm{emb}$,
  so these forward references are not circular.)
\item If, in addition, $T$ is in general position, $\mathrm{cfg}(T)$ is the configuration of $\card T$ points
  $i\mapsto\mathrm{sortPt}(T)(i)$. Its two defining properties are Lemma~\ref{mm:lem:sort}(e) and~(h).
\end{enumerate}
\end{definition}

\begin{remark}
Caps in Definition~\ref{mm:def:config}(b) are \emph{all-triples} caps: every increasing triple of the
selected points is clockwise. The written sources \dirn{MIRRORETV/README.md} and
\dirn{MIRRORETV/PROOFS.md} define caps by consecutive triples; for points with increasing abscissae the
two notions agree (\dirn{MIRRORETV/PROOFS.md}, Lemma~2.1(b)), but only the all-triples notion occurs in
this section. For $k\le2$ the condition on triples is vacuous, so a configuration of $n$ points has a
$k$-cap if and only if $k\le n$.
\end{remark}

\begin{lemma}[sorting; \leanref{MirrorOdd.xset_card}{MIRRORETV/lean/MirrorMain.lean:1551}, \leanref{MirrorOdd.emb_mem}{MIRRORETV/lean/MirrorMain.lean:1555}, \leanref{MirrorOdd.sortPt_mem}{MIRRORETV/lean/MirrorMain.lean:1560}, \leanref{MirrorOdd.sortPt_x}{MIRRORETV/lean/MirrorMain.lean:1563}, \leanref{MirrorOdd.sortPt_strictMono}{MIRRORETV/lean/MirrorMain.lean:1566}, \leanref{MirrorOdd.sortPt_injective}{MIRRORETV/lean/MirrorMain.lean:1572}, \leanref{MirrorOdd.sortPt_image}{MIRRORETV/lean/MirrorMain.lean:1576}]\label{mm:lem:sort}
Let $T\subset\R^2$ be finite and $x$-injective, and write $s_i=\mathrm{sortPt}(T)(i)$ for
$i\in[\card T]$.
\begin{enumerate}[label=(\alph*)]
\item $\card{\mathrm{xset}(T)}=\card T$ (\lean{xset_card}).
\item For every $i\in[\card T]$ there is $p\in T$ with $p_0=\mathrm{emb}(T)(i)$ (\lean{emb_mem}).
\item $s_i\in T$ for every $i$ (\lean{sortPt_mem}).
\item $(s_i)_0=\mathrm{emb}(T)(i)$ for every $i$ (\lean{sortPt_x}).
\item $i\mapsto(s_i)_0$ is strictly increasing on $[\card T]$ (\lean{sortPt_strictMono}).
\item $i\mapsto s_i$ is injective (\lean{sortPt_injective}).
\item $\{s_i:i\in[\card T]\}=T$ (\lean{sortPt_image}).
\item (the field \lean{generalPosition} of \lean{cfg}, \dirn{MIRRORETV/lean/MirrorMain.lean:1587})
  If $T$ is in general position, then $\orient(s_i,s_j,s_k)\ne0$ for all $i<j<k$ in $[\card T]$.
\end{enumerate}
Consequently, if $T$ is in general position, $\mathrm{cfg}(T)$ is a configuration of $\card T$ points with
$\mathrm{pts}(\mathrm{cfg}(T))=T$.
\end{lemma}

\begin{proof}
(a) The map $p\mapsto p_0$ is injective on $T$, so its image $\mathrm{xset}(T)$ has $\card T$ elements.
(b) $\mathrm{emb}(T)(i)$ lies in the image $\mathrm{xset}(T)$ of $\mathrm{emb}(T)$, so it equals $p_0$ for
some $p\in T$.
(c), (d) These are the two properties of the point chosen in Definition~\ref{mm:def:config}(c).
(e) $\mathrm{emb}(T)$ is strictly increasing; use (d).
(f) If $s_i=s_j$, then $(s_i)_0=(s_j)_0$, so $i=j$ by (e).
(g) By (c) the set $I=\{s_i:i\in[\card T]\}$ is contained in $T$, and by (f) it has $\card T$ elements. A
subset of the finite set $T$ with $\card T$ elements is $T$ itself.
(h) By (f) the points $s_i,s_j,s_k$ are pairwise distinct, and by (c) they lie in $T$. Since $T$ is in
general position, $\orient(s_i,s_j,s_k)\ne0$ by Lemma~\ref{fd:lem:collinear}.
Finally, (e) and (h) are the two conditions of Definition~\ref{mm:def:config}(a) for $i\mapsto s_i$, and
(g) says that the point set is $T$.
\end{proof}

\begin{lemma}[two facts on convex polygons; \leanref{MirrorOdd.hasConvexNGon_of_le}{MIRRORETV/lean/MirrorMain.lean:2182}, \leanref{MirrorOdd.hasConvexNGon_affineMap}{MIRRORETV/lean/MirrorMain.lean:1859}]\label{mm:lem:ngon}
Let $S\subset\R^2$ be finite.
\begin{enumerate}[label=(\alph*)]
\item (\lean{hasConvexNGon_of_le}) If $S$ has a convex $m$-gon and $n\le m$, then $S$ has a convex
  $n$-gon.
\item (\lean{hasConvexNGon_affineMap}) If $S$ has a convex $n$-gon and $f:\R^2\to\R^2$ is an injective
  affine map, then $f(S)$ has a convex $n$-gon.
\end{enumerate}
\end{lemma}

\begin{proof}
(a) This is Lemma~\ref{fd:lem:mono}(d). Directly: take $T\subseteq S$ with $\card T=m$ in convex
position and choose $T'\subseteq T$ with $\card{T'}=n$. Then $T'\subseteq S$, and $T'$ is in convex
position by Lemma~\ref{fd:lem:mono}(a).

(b) Take $T\subseteq S$ with $\card T=n$ in convex position. Then $f(T)\subseteq f(S)$ and
$\card{f(T)}=n$ because $f$ is injective. An affine map preserves affine combinations:
$f(\sum_i\lambda_iy_i)=\sum_i\lambda_if(y_i)$ whenever $\sum_i\lambda_i=1$. Since the convex hull of a set
consists of the convex combinations of its finite subsets, $f(\conv Y)=\conv f(Y)$ for every
$Y\subseteq\R^2$. Let $q\in f(T)$, say $q=f(p)$ with $p\in T$. By injectivity,
$f(T)\setminus\{f(p)\}=f(T\setminus\{p\})$. If $f(p)\in\conv f(T\setminus\{p\})=f(\conv(T\setminus\{p\}))$,
then $f(p)=f(x)$ for some $x\in\conv(T\setminus\{p\})$, so $p=x\in\conv(T\setminus\{p\})$ by injectivity,
which contradicts the convex position of $T$ (Lemma~\ref{fd:lem:cp-hull}). Hence no point of $f(T)$ lies
in the convex hull of the other points of $f(T)$, and $f(T)$ is in convex position by
Lemma~\ref{fd:lem:cp-hull}. (This is also contained in Lemma~\ref{fd:lem:affine-cp}.)
\end{proof}

\subsection{The mirror lemma}

\begin{lemma}[cap facts; \leanref{MirrorOdd.cap_below_lines}{MIRRORETV/lean/MirrorMain.lean:2204}]\label{mm:lem:slope}
Both parts concern points that lie on or above the $x$-axis.
\begin{enumerate}[label=(\alph*)]
\item (\lean{cap_below_lines}) Let $m\in\N$ and let $c:[m]\to\R^2$ be such that $i\mapsto(c_i)_0$ is
  strictly increasing, $(c_i)_1\ge0$ for all $i\in[m]$, and $\orient(c_i,c_j,c_l)<0$ for all $i<j<l$ in
  $[m]$. Then for all $i,j,l\in[m]$ with $i<j$ and $(c_l)_1>0$,
  \[ \orient(c_i,c_j,\operatorname{mir}(c_l))<0. \]
\item (Written step. It is not a separate Lean declaration: the same argument is carried out inline in
  the proof of \lean{mirror_odd_general}, \dirn{MIRRORETV/lean/MirrorMain.lean:2542}--2557, and as the
  local claim \lean{interior} in \lean{mirror_five}, line~2060, and in \lean{mirror_seven_of},
  line~1716.) Let $p,a,q\in\R^2$ with $p_0<a_0<q_0$, $p_1\ge0$, $q_1\ge0$ and $a_1=0$. Then
  \[ \orient(p,a,q)=(a_0-p_0)\,q_1+p_1\,(q_0-a_0)\ \ge\ 0 . \]
  In particular $\orient(p,a,q)<0$ is impossible: a point on the axis cannot be the middle point of a
  clockwise triple with increasing abscissae whose other two points lie on or above the axis.
\end{enumerate}
\end{lemma}

In words, (a) says that the mirror image of a cap point lying strictly above the axis lies strictly to the
right of every directed chord $c_i\to c_j$ of the cap ($i<j$), that is, strictly below the line through
$c_i$ and $c_j$.

\begin{proof}
(a) Write $x_k=(c_k)_0$ and $y_k=(c_k)_1$. By Lemma~\ref{mm:lem:orient-mir}(e),
\[ \orient(c_i,c_j,\operatorname{mir}c_l)=\orient(c_i,c_j,c_l)-2y_l(x_j-x_i), \tag{$*$} \]
and $2y_l(x_j-x_i)>0$ because $y_l>0$ and $x_j>x_i$. We distinguish the position of $l$.
\begin{itemize}
\item $l<i$: $\orient(c_i,c_j,c_l)=\orient(c_l,c_i,c_j)$ by Lemma~\ref{mm:lem:orient-mir}(d), and this is
  negative by the cap hypothesis for $l<i<j$. So $(*)$ is negative.
\item $l=i$ or $l=j$: $\orient(c_i,c_j,c_l)=0$ because two arguments coincide, so $(*)$ equals
  $-2y_l(x_j-x_i)<0$.
\item $i<l<j$: expanding directly,
  \[ \orient(c_i,c_j,\operatorname{mir}c_l)=-(x_j-x_i)\,y_l-y_i\,(x_j-x_l)-y_j\,(x_l-x_i). \]
  The first term is negative, and the other two are $\le0$ because $y_i,y_j\ge0$ and $x_i<x_l<x_j$.
  (The cap hypothesis is not needed in this case.)
\item $l>j$: $\orient(c_i,c_j,c_l)<0$ by the cap hypothesis for $i<j<l$, so $(*)$ is negative.
\end{itemize}
(b) Expanding with $a_1=0$:
$\orient(p,a,q)=(a_0-p_0)(q_1-p_1)+p_1(q_0-p_0)=(a_0-p_0)q_1+p_1(q_0-a_0)$, and both products are $\ge0$.
\end{proof}

\begin{theorem}[mirror lemma, no point on the axis; \leanref{MirrorOdd.cap_mirror_polygon}{MIRRORETV/lean/MirrorMain.lean:2233}]\label{mm:thm:mirror-polygon}
Let $m\in\N$ with $2\le m$, and let $c:[m]\to\R^2$ be such that $i\mapsto(c_i)_0$ is strictly increasing,
$(c_i)_1>0$ for all $i\in[m]$, and $\orient(c_i,c_j,c_l)<0$ for all $i<j<l$ in $[m]$. Let $P\subset\R^2$
be a finite set with $c_i\in P$ and $\operatorname{mir}(c_i)\in P$ for all $i\in[m]$. Then $P$ has a
convex $2m$-gon.
\end{theorem}

\begin{proof}
We follow the Lean proof. Write $x_i=(c_i)_0$ and $y_i=(c_i)_1>0$. Define $q:[2m]\to\R^2$ by
\[ q_r=\operatorname{mir}(c_r)\quad(0\le r\le m-1),\qquad q_r=c_{2m-1-r}\quad(m\le r\le 2m-1). \]
Thus $q$ lists $\operatorname{mir}c_0,\dots,\operatorname{mir}c_{m-1}$ (the reflected cap, from left to
right) followed by $c_{m-1},\dots,c_0$ (the cap, from right to left). For $m\le r\le 2m-1$ the index
$2m-1-r$ lies in $[m]$, and $r\mapsto2m-1-r$ is strictly decreasing. Since all $y_i>0$, the hypotheses
of Lemma~\ref{mm:lem:slope}(a) hold for $c$.

\emph{Claim~1: $\orient(q_r,q_s,q_t)>0$ for all $r<s<t$ in $[2m]$.} We distinguish how many of $r,s,t$
are smaller than $m$.
\begin{itemize}
\item Case~1, $t<m$. By Lemma~\ref{mm:lem:orient-mir}(b),
  $\orient(\operatorname{mir}c_r,\operatorname{mir}c_s,\operatorname{mir}c_t)=-\orient(c_r,c_s,c_t)$, which
  is positive by the cap hypothesis for $r<s<t$.
\item Case~2, $s<m\le t$. Put $u=2m-1-t\in[m]$, so $q_t=c_u$. By Lemma~\ref{mm:lem:orient-mir}(c),
  $\orient(\operatorname{mir}c_r,\operatorname{mir}c_s,c_u)=-\orient(c_r,c_s,\operatorname{mir}c_u)$, and
  $\orient(c_r,c_s,\operatorname{mir}c_u)<0$ by Lemma~\ref{mm:lem:slope}(a) with $(i,j,l)=(r,s,u)$, since
  $r<s$ and $y_u>0$.
\item Case~3, $r<m\le s$. Put $u'=2m-1-s$ and $u=2m-1-t$. Then $u<u'$ (because $s<t$), $q_s=c_{u'}$ and
  $q_t=c_u$. By the first two identities of Lemma~\ref{mm:lem:orient-mir}(d),
  \[ \orient(\operatorname{mir}c_r,c_{u'},c_u)=\orient(c_{u'},c_u,\operatorname{mir}c_r)
     =-\orient(c_u,c_{u'},\operatorname{mir}c_r), \]
  and $\orient(c_u,c_{u'},\operatorname{mir}c_r)<0$ by Lemma~\ref{mm:lem:slope}(a) with $(i,j,l)=(u,u',r)$,
  since $u<u'$ and $y_r>0$.
\item Case~4, $m\le r$. Put $\alpha=2m-1-r$, $\beta=2m-1-s$ and $\delta=2m-1-t$; then
  $\delta<\beta<\alpha$, $q_r=c_\alpha$, $q_s=c_\beta$ and $q_t=c_\delta$. By the third identity of
  Lemma~\ref{mm:lem:orient-mir}(d), $\orient(c_\alpha,c_\beta,c_\delta)=-\orient(c_\delta,c_\beta,c_\alpha)$,
  which is positive by the cap hypothesis for $\delta<\beta<\alpha$.
\end{itemize}

\emph{Claim~2: $q$ is injective.} Let $r\ne s$ in $[2m]$. If $r,s<m$ and $q_r=q_s$, then comparing
abscissae gives $x_r=x_s$ (Lemma~\ref{mm:lem:orient-mir}(a)), so $r=s$ because $i\mapsto x_i$ is strictly
increasing; a contradiction. If $r<m\le s$, then $(q_r)_1=-y_r<0<y_{2m-1-s}=(q_s)_1$, so $q_r\ne q_s$; the
case $s<m\le r$ is symmetric. If $r,s\ge m$ and $q_r=q_s$, then $x_{2m-1-r}=x_{2m-1-s}$, so
$2m-1-r=2m-1-s$ and $r=s$; a contradiction.

\emph{Conclusion.} Every $q_r$ lies in $P$, being some $\operatorname{mir}(c_i)$ or some $c_i$. Apply
Lemma~\ref{fd:lem:pos-cycle} with $N=k=2m$ (note $2m\ge4$ because $m\ge2$), the injective map $p=q$
(Claim~2), $e$ the identity of $[2m]$ and the finite set $P$; its positivity hypothesis is Claim~1. Hence
$P$ has a convex $2m$-gon.
\end{proof}

\begin{remark}
(1) Claim~1 checks the sign of \emph{every} triple $q_r,q_s,q_t$ with $r<s<t$, that is, of every triple
of the list taken in list order, which is what Lemma~\ref{fd:lem:pos-cycle} requires. It would not suffice to check that the closed polygon
$q_0q_1\cdots q_{2m-1}$ turns the same way at every vertex: such a polygon can wind around twice, and then
a vertex can lie inside the convex hull of the others (\dirn{MIRRORETV/PROOFS.md}, remark after Lemma~3.1,
gives a five-point example).
(2) \dirn{MIRRORETV/PROOFS.md}, Lemma~3.1, gives a written argument by a different method (a region
argument) for a version with caps $C$ defined by consecutive triples: if the interior points of $C$ have
positive ordinates and its endpoints have nonnegative ordinates (positive if $C$ has only one point), then
$C\cup\operatorname{mir}(C)$ is in convex position. That argument is not used here; the Lean statements
are Theorems~\ref{mm:thm:mirror-polygon} and~\ref{mm:thm:axis-polygon}, and these are what we prove.
\end{remark}

\begin{theorem}[mirror lemma, one endpoint on the axis; \leanref{MirrorOdd.axis_left_polygon}{MIRRORETV/lean/MirrorMain.lean:2309}, \leanref{MirrorOdd.axis_right_polygon}{MIRRORETV/lean/MirrorMain.lean:2403}]\label{mm:thm:axis-polygon}
Let $m\in\N$ with $3\le m$, and let $c:[m]\to\R^2$ be such that $i\mapsto(c_i)_0$ is strictly increasing
and $\orient(c_i,c_j,c_l)<0$ for all $i<j<l$ in $[m]$. Let $P\subset\R^2$ be a finite set with $c_i\in P$
and $\operatorname{mir}(c_i)\in P$ for all $i\in[m]$.
\begin{enumerate}[label=(\alph*)]
\item (\lean{axis_left_polygon}) If $(c_0)_1=0$ and $(c_i)_1>0$ for every $i\in[m]$ with $0<i$, then $P$
  has a convex $(2m-1)$-gon.
\item (\lean{axis_right_polygon}) If $(c_{m-1})_1=0$ and $(c_i)_1>0$ for every $i\in[m]$ with $i<m-1$,
  then $P$ has a convex $(2m-1)$-gon.
\end{enumerate}
\end{theorem}

\begin{proof}
We follow the Lean proofs. Write $x_i=(c_i)_0$ and $y_i=(c_i)_1$. In both parts $y_i\ge0$ for all $i$,
so Lemma~\ref{mm:lem:slope}(a) applies to $c$; and $2m-1\ge5$.

(a) Here $\operatorname{mir}(c_0)=c_0$ because $(c_0)_1=0$. Define $q:[2m-1]\to\R^2$ by
\[ q_r=\operatorname{mir}(c_r)\quad(0\le r\le m-1),\qquad q_r=c_{2m-1-r}\quad(m\le r\le 2m-2). \]
So $q$ lists $\operatorname{mir}c_0=c_0,\operatorname{mir}c_1,\dots,\operatorname{mir}c_{m-1}$ and then
$c_{m-1},\dots,c_1$. For $m\le r\le 2m-2$ the index $2m-1-r$ lies in $\{1,\dots,m-1\}$, hence
$y_{2m-1-r}>0$.

\emph{Claim~1: $\orient(q_r,q_s,q_t)>0$ for all $r<s<t$ in $[2m-1]$.}
\begin{itemize}
\item Case~1 ($t<m$) and Case~4 ($m\le r$) are exactly as in the proof of
  Theorem~\ref{mm:thm:mirror-polygon}; they use only the cap hypothesis (in Case~4 the indices are
  $2m-1-r>2m-1-s>2m-1-t$).
\item Case~2, $s<m\le t$. With $u=2m-1-t\in\{1,\dots,m-1\}$ we have $q_t=c_u$ and $y_u>0$, and Case~2 of
  the proof of Theorem~\ref{mm:thm:mirror-polygon} applies verbatim.
\item Case~3, $r<m\le s$. Put $u'=2m-1-s$ and $u=2m-1-t$, so $1\le u<u'\le m-1$, $q_s=c_{u'}$, $q_t=c_u$.
  \begin{itemize}
  \item Case~3a, $r=0$. Then $q_r=\operatorname{mir}c_0=c_0$, and by Lemma~\ref{fd:lem:orient}(a)
    (exchange of the last two arguments) $\orient(c_0,c_{u'},c_u)=-\orient(c_0,c_u,c_{u'})$, which is
    positive by the cap hypothesis for $0<u<u'$.
  \item Case~3b, $r\ge1$. Then $y_r>0$, and Case~3 of the proof of Theorem~\ref{mm:thm:mirror-polygon}
    applies verbatim:
    $\orient(\operatorname{mir}c_r,c_{u'},c_u)=-\orient(c_u,c_{u'},\operatorname{mir}c_r)>0$.
  \end{itemize}
\end{itemize}

\emph{Claim~2: $q$ is injective.} Two indices $r\ne s$ that are both $<m$, or both $\ge m$, are separated by
the abscissae as in Theorem~\ref{mm:thm:mirror-polygon}. If $r<m\le s$, then $(q_s)_1=y_{2m-1-s}>0$
because $2m-1-s\ge1$, while $(q_r)_1=-y_r\le0$. The case $s<m\le r$ is symmetric.

\emph{Conclusion.} Every $q_r$ lies in $P$, and Lemma~\ref{fd:lem:pos-cycle} with $N=k=2m-1\ge4$, $p=q$
and $e$ the identity gives a convex $(2m-1)$-gon in $P$.

(b) Here $\operatorname{mir}(c_{m-1})=c_{m-1}$. Define $q:[2m-1]\to\R^2$ by
\[ q_r=\operatorname{mir}(c_r)\quad(0\le r\le m-1),\qquad q_r=c_{2m-2-r}\quad(m\le r\le 2m-2). \]
So $q$ lists $\operatorname{mir}c_0,\dots,\operatorname{mir}c_{m-2},\operatorname{mir}c_{m-1}=c_{m-1}$ and
then $c_{m-2},\dots,c_0$. For $m\le r\le 2m-2$ the index $2m-2-r$ lies in $\{0,\dots,m-2\}$, hence
$y_{2m-2-r}>0$.

\emph{Claim~1.}
\begin{itemize}
\item Cases~1 and~4 are as in Theorem~\ref{mm:thm:mirror-polygon}, with $2m-2-\cdot$ in place of
  $2m-1-\cdot$.
\item Case~2, $s<m\le t$. With $u=2m-2-t\le m-2$ we have $q_t=c_u$ and $y_u>0$, and Case~2 of
  Theorem~\ref{mm:thm:mirror-polygon} applies.
\item Case~3, $r<m\le s$. Put $u'=2m-2-s$ and $u=2m-2-t$, so $0\le u<u'\le m-2$.
  \begin{itemize}
  \item Case~3a, $r=m-1$. Then $q_r=c_{m-1}$, and by the third identity of
    Lemma~\ref{mm:lem:orient-mir}(d), $\orient(c_{m-1},c_{u'},c_u)=-\orient(c_u,c_{u'},c_{m-1})$, which is
    positive by the cap hypothesis for $u<u'<m-1$.
  \item Case~3b, $r<m-1$. Then $y_r>0$, and Case~3 of Theorem~\ref{mm:thm:mirror-polygon} applies.
  \end{itemize}
\end{itemize}

\emph{Claim~2.} Two indices $r\ne s$ that are both $<m$, or both $\ge m$, are separated by the abscissae
as in (a). If $r<m\le s$, then $(q_s)_1=y_{2m-2-s}>0$ because $2m-2-s\le m-2$, while
$(q_r)_1=-y_r\le0$. The case $s<m\le r$ is symmetric.

\emph{Conclusion} as in (a).
\end{proof}

\begin{remark}
The hypothesis $m\ge3$ makes the list long enough for Lemma~\ref{fd:lem:pos-cycle}, which needs at least
four points. In (a) the point $c_0$ occurs in the list only once, as $\operatorname{mir}c_0=c_0$; in (b)
the same holds for $c_{m-1}$. This is why one point is lost compared with Theorem~\ref{mm:thm:mirror-polygon}.
\end{remark}

\subsection{The generated small polygons}

Lines 22--1506 of \dirn{MIRRORETV/lean/MirrorMain.lean} contain five helper lemmas $\mathrm{cyc}_k$
($k=4,\dots,8$) and nine lemmas on small mirror polygons. In Lean each polygon lemma is proved directly:
the pairwise distinctness of its $k$ vertices by comparing one coordinate, and each of the $\binom k3$
orientation inequalities, after expanding $\orient$ and $\operatorname{mir}$, by linear arithmetic
(\lean{linarith}) from the hypotheses (with the cap orientations expanded) and auxiliary products
$(c_i)_1\,(u_0-v_0)>0$, where $u$ and $v$ are two of the lemma's points with $v_0<u_0$ (in the axis
lemmas one of them can be the axis point $a$), with certificates found offline; then $\mathrm{cyc}_k$ is
applied. These Lean proofs do not use the
general lemmas, which come later in the file. Our written proofs are different: we derive the nine lemmas
from Theorems~\ref{mm:thm:mirror-polygon} and~\ref{mm:thm:axis-polygon}.

\begin{lemma}[positive cycles of length $4$ to $8$; \leanref{MirrorOdd.cyc4}{MIRRORETV/lean/MirrorMain.lean:995}, \leanref{MirrorOdd.cyc5}{MIRRORETV/lean/MirrorMain.lean:1026}, \leanref{MirrorOdd.cyc6}{MIRRORETV/lean/MirrorMain.lean:1068}, \leanref{MirrorOdd.cyc7}{MIRRORETV/lean/MirrorMain.lean:22}, \leanref{MirrorOdd.cyc8}{MIRRORETV/lean/MirrorMain.lean:103}]\label{mm:lem:cyc}
Let $k\in\{4,5,6,7,8\}$, let $q_0,\dots,q_{k-1}\in\R^2$, and let $P\subset\R^2$ be finite. Suppose that
$q_r\in P$ for all $r\in[k]$, that $q_r\ne q_s$ for all $r<s$ in $[k]$, and that
$\orient(q_r,q_s,q_t)>0$ for all $r<s<t$ in $[k]$. Then $P$ has a convex $k$-gon. (The Lean lemma
$\mathrm{cyc}_k$ lists these $k$ memberships, $\binom k2$ inequalities $q_r\ne q_s$ and $\binom k3$
orientation inequalities as separate hypotheses.)
\end{lemma}

\begin{proof}
The map $r\mapsto q_r$ on $[k]$ is injective by the distinctness hypotheses. Lemma~\ref{fd:lem:pos-cycle}
with $N=k\ge4$, $p=q$, $e$ the identity of $[k]$ and the set $P$ gives the claim; its remaining hypotheses
are exactly the memberships and the orientation inequalities. This is also the Lean proof.
\end{proof}

\begin{proposition}[generated small mirror polygons; \leanref{MirrorOdd.cap_mirror_quadrilateral}{MIRRORETV/lean/MirrorMain.lean:1126}, \leanref{MirrorOdd.cap_mirror_pentagon}{MIRRORETV/lean/MirrorMain.lean:1282}, \leanref{MirrorOdd.cap_mirror_hexagon}{MIRRORETV/lean/MirrorMain.lean:1168}, \leanref{MirrorOdd.cap_mirror_heptagon}{MIRRORETV/lean/MirrorMain.lean:212}, \leanref{MirrorOdd.cap_mirror_octagon}{MIRRORETV/lean/MirrorMain.lean:392}, \leanref{MirrorOdd.axis_left_pentagon}{MIRRORETV/lean/MirrorMain.lean:1358}, \leanref{MirrorOdd.axis_right_pentagon}{MIRRORETV/lean/MirrorMain.lean:1433}, \leanref{MirrorOdd.axis_left_heptagon}{MIRRORETV/lean/MirrorMain.lean:639}, \leanref{MirrorOdd.axis_right_heptagon}{MIRRORETV/lean/MirrorMain.lean:816}]\label{mm:prop:small-polygons}
Let $a,c_0,c_1,c_2,c_3\in\R^2$ and let $P\subset\R^2$ be finite. Write $x_i=(c_i)_0$ and $y_i=(c_i)_1$.
\begin{enumerate}[label=(\alph*)]
\item (\lean{cap_mirror_quadrilateral}) If $x_0<x_1$, $y_0>0$, $y_1>0$, and
  $\operatorname{mir}c_0,\operatorname{mir}c_1,c_1,c_0\in P$, then $P$ has a convex $4$-gon.
\item (\lean{cap_mirror_pentagon}) If $x_0<x_1<x_2$, $y_0,y_1,y_2>0$, $\orient(c_0,c_1,c_2)<0$, and
  $\operatorname{mir}c_0,\operatorname{mir}c_1,\operatorname{mir}c_2,c_2,c_1\in P$, then $P$ has a convex
  $5$-gon.
\item (\lean{cap_mirror_hexagon}) If $x_0<x_1<x_2$, $y_0,y_1,y_2>0$, $\orient(c_0,c_1,c_2)<0$, and
  $\operatorname{mir}c_0,\operatorname{mir}c_1,\operatorname{mir}c_2,c_2,c_1,c_0\in P$, then $P$ has a
  convex $6$-gon.
\item (\lean{cap_mirror_heptagon}) If $x_0<x_1<x_2<x_3$, $y_0,y_1,y_2,y_3>0$, the four orientations
  $\orient(c_0,c_1,c_2)$, $\orient(c_0,c_1,c_3)$, $\orient(c_0,c_2,c_3)$, $\orient(c_1,c_2,c_3)$ are $<0$,
  and $c_0,\operatorname{mir}c_1,\operatorname{mir}c_2,\operatorname{mir}c_3,c_3,c_2,c_1\in P$, then $P$
  has a convex $7$-gon.
\item (\lean{cap_mirror_octagon}) Under the hypotheses on $c_0,\dots,c_3$ of (d), if
  $\operatorname{mir}c_0,\operatorname{mir}c_1,\operatorname{mir}c_2,\operatorname{mir}c_3,c_3,c_2,c_1,c_0\in P$,
  then $P$ has a convex $8$-gon.
\item (\lean{axis_left_pentagon}) If $a_1=0$, $a_0<x_0<x_1$, $y_0,y_1>0$, $\orient(a,c_0,c_1)<0$, and
  $a,\operatorname{mir}c_0,\operatorname{mir}c_1,c_1,c_0\in P$, then $P$ has a convex $5$-gon.
\item (\lean{axis_right_pentagon}) If $a_1=0$, $x_0<x_1<a_0$, $y_0,y_1>0$, $\orient(c_0,c_1,a)<0$, and
  $\operatorname{mir}c_0,\operatorname{mir}c_1,a,c_1,c_0\in P$, then $P$ has a convex $5$-gon.
\item (\lean{axis_left_heptagon}) If $a_1=0$, $a_0<x_0<x_1<x_2$, $y_0,y_1,y_2>0$, the four orientations
  $\orient(a,c_0,c_1)$, $\orient(a,c_0,c_2)$, $\orient(a,c_1,c_2)$, $\orient(c_0,c_1,c_2)$ are $<0$, and
  $a,\operatorname{mir}c_0,\operatorname{mir}c_1,\operatorname{mir}c_2,c_2,c_1,c_0\in P$, then $P$ has a
  convex $7$-gon.
\item (\lean{axis_right_heptagon}) If $a_1=0$, $x_0<x_1<x_2<a_0$, $y_0,y_1,y_2>0$, the four orientations
  $\orient(c_0,c_1,c_2)$, $\orient(c_0,c_1,a)$, $\orient(c_0,c_2,a)$, $\orient(c_1,c_2,a)$ are $<0$, and
  $\operatorname{mir}c_0,\operatorname{mir}c_1,\operatorname{mir}c_2,a,c_2,c_1,c_0\in P$, then $P$ has a
  convex $7$-gon.
\end{enumerate}
In each item the memberships are listed in the vertex order that the Lean proof passes to
$\mathrm{cyc}_k$ (Lemma~\ref{mm:lem:cyc}).
\end{proposition}

\begin{proof}
We use the general theorems; this differs from the generated Lean proofs described above.

(a), (c), (e): apply Theorem~\ref{mm:thm:mirror-polygon} with $m=2,3,4$ to $c=(c_0,\dots,c_{m-1})$ and to
$P$. The abscissae increase strictly because consecutive ones do, and all $y_i>0$. The cap hypothesis
concerns the triples $i<j<l$ in $[m]$: there are none for $m=2$, only $(0,1,2)$ for $m=3$, and
$(0,1,2)$, $(0,1,3)$, $(0,2,3)$, $(1,2,3)$ for $m=4$; these are exactly the listed orientation hypotheses.
$P$ contains all $c_i$ and $\operatorname{mir}c_i$. The conclusion is a convex $2m$-gon, that is, a
$4$-, $6$- or $8$-gon.

(f), (h): note that $\operatorname{mir}(a)=(a_0,-0)=a$. Apply Theorem~\ref{mm:thm:axis-polygon}(a) with
$m=3$ to $(a,c_0,c_1)$, respectively with $m=4$ to $(a,c_0,c_1,c_2)$. The abscissae increase strictly, the
first point $a$ has ordinate $0$ and the others have positive ordinates. For $m=3$ the only triple is
$(a,c_0,c_1)$; for $m=4$ the triples $(0,1,2)$, $(0,1,3)$, $(0,2,3)$, $(1,2,3)$ of $(a,c_0,c_1,c_2)$ are
$(a,c_0,c_1)$, $(a,c_0,c_2)$, $(a,c_1,c_2)$, $(c_0,c_1,c_2)$. These are the listed orientation hypotheses.
$P$ contains the listed points and their mirror images (recall $\operatorname{mir}a=a$). The conclusion is
a convex $(2m-1)$-gon, that is, a $5$- or $7$-gon.

(g), (i): the same with Theorem~\ref{mm:thm:axis-polygon}(b), applied to $(c_0,c_1,a)$, respectively
$(c_0,c_1,c_2,a)$, whose last point $a$ has ordinate $0$. For $m=4$ the four triples are
$(c_0,c_1,c_2)$, $(c_0,c_1,a)$, $(c_0,c_2,a)$, $(c_1,c_2,a)$.

(b): let $Q=\{\operatorname{mir}c_0,\operatorname{mir}c_1,\operatorname{mir}c_2,c_2,c_1,c_0\}$, a set with
at most $6$ elements. By (c) applied to the finite set $Q$ in place of $P$, there is $T\subseteq Q$ with
$\card T=6$ in convex position. As $\card Q\le6$, we get $T=Q$: the six listed points are pairwise distinct
and $Q$ is in convex position. Then $Q\setminus\{c_0\}=\{\operatorname{mir}c_0,\operatorname{mir}c_1,
\operatorname{mir}c_2,c_2,c_1\}$ has $5$ elements, is in convex position (Lemma~\ref{fd:lem:mono}(a)), and
is contained in $P$ by hypothesis. So $P$ has a convex $5$-gon.

(d): in the same way, (e) applied to the set $Q$ of its eight points shows that these eight points are
pairwise distinct and in convex position. The set $Q\setminus\{\operatorname{mir}c_0\}=\{c_0,
\operatorname{mir}c_1,\operatorname{mir}c_2,\operatorname{mir}c_3,c_3,c_2,c_1\}$ has $7$ elements, is in
convex position and is contained in $P$.
\end{proof}

In Lean, items (a), (b), (c), (f), (g) are used in the proofs of Theorem~\ref{mm:thm:small}, and items
(d), (e), (h), (i) in those of the bounds for $k=7,8$ (Theorems~\ref{mb:thm:seven}
and~\ref{mb:thm:eight}).

\subsection{Counting and distinct abscissae}

\begin{lemma}[counting; \leanref{MirrorOdd.card_split}{MIRRORETV/lean/MirrorMain.lean:1601}, \leanref{MirrorOdd.card_A_le}{MIRRORETV/lean/MirrorMain.lean:1636}]\label{mm:lem:count}
Let $S\subset\R^2$ be finite.
\begin{enumerate}[label=(\alph*)]
\item (\lean{card_split}) If $\operatorname{mir}(p)\in S$ for every $p\in S$, then
  $\card S=2\card{R(S)}+\card{A(S)}$. (General position is not assumed.)
\item (\lean{card_A_le}) If $S$ is in general position, then $\card{A(S)}\le2$.
\end{enumerate}
\end{lemma}

\begin{proof}
(a) Let $N=\{p\in S:p_1<0\}$. By trichotomy of $p_1$, $S$ is the disjoint union of $R(S)$, $A(S)$ and $N$.
We show $\operatorname{mir}(R(S))=N$. If $q\in R(S)$, then $\operatorname{mir}q\in S$ by hypothesis and
$(\operatorname{mir}q)_1=-q_1<0$, so $\operatorname{mir}q\in N$. Conversely, if $p\in N$, then
$\operatorname{mir}p\in S$ and $(\operatorname{mir}p)_1=-p_1>0$, so $\operatorname{mir}p\in R(S)$, and
$p=\operatorname{mir}(\operatorname{mir}p)$ by Lemma~\ref{mm:lem:orient-mir}(a). Since $\operatorname{mir}$
is injective (it is its own inverse), $\card N=\card{R(S)}$, and therefore
$\card S=\card{R(S)}+\card{A(S)}+\card N=2\card{R(S)}+\card{A(S)}$.

(b) Suppose that $A(S)$ contains three pairwise distinct points $a,b,d$. Since $a_1=b_1=d_1=0$,
$\orient(a,b,d)=(b_0-a_0)(d_1-a_1)-(b_1-a_1)(d_0-a_0)=0$, which contradicts general position
(Lemma~\ref{fd:lem:collinear}). Hence $\card{A(S)}\le2$.
\end{proof}

\begin{lemma}[distinct abscissae; \leanref{MirrorOdd.injOn_x}{MIRRORETV/lean/MirrorMain.lean:1647}]\label{mm:lem:injx}
Let $S\subset\R^2$ be finite and in general position, with $\operatorname{mir}(p)\in S$ for every $p\in S$.
Let $T\subseteq S$ be a set such that $p_1\ge0$ for every $p\in T$, and such that $p=q$ whenever $p,q\in T$
and $p_1=q_1=0$ (at most one point of $T$ lies on the axis). Then $T$ is $x$-injective.
\end{lemma}

\begin{proof}
First, if $p,q\in T$, $p_1>0$ and $p_0=q_0$, then $p=q$. Indeed, suppose $p\ne q$. The points $p$, $q$
and $\operatorname{mir}p$ lie in $S$ (as $T\subseteq S$ and $S$ is $\operatorname{mir}$-symmetric) and are
pairwise distinct: $p\ne\operatorname{mir}p$ because $p_1>0>-p_1$, and $q\ne\operatorname{mir}p$ because
$q_1\ge0>-p_1$. But
$\orient(p,q,\operatorname{mir}p)=(q_0-p_0)(-p_1-p_1)-(q_1-p_1)(p_0-p_0)=0$ because $q_0=p_0$, which
contradicts general position (Lemma~\ref{fd:lem:collinear}).

Now let $p,q\in T$ with $p_0=q_0$. If $p_1>0$, then $p=q$ by the first step. Otherwise $p_1=0$. If
$q_1>0$, the first step with the roles of $p$ and $q$ exchanged gives $q=p$. Otherwise $q_1=0$ as well,
and $p=q$ by the hypothesis on axis points.
\end{proof}

\begin{lemma}[caps in the closed upper half; written step]\label{mm:lem:upper-caps}
Let $S\subset\R^2$ be finite and in general position, with $\operatorname{mir}(p)\in S$ for every $p\in S$.
Let $T\subseteq S$ satisfy $p_1\ge0$ for every $p\in T$, and $p=q$ whenever $p,q\in T$ and $p_1=q_1=0$.
Then $T$ is $x$-injective and in general position, so $C=\mathrm{cfg}(T)$ is a configuration of
$\card T$ points with $\mathrm{pts}(C)=T$. Let $m\ge2$ and suppose that $C$ has an $m$-cap.
\begin{enumerate}[label=(\alph*)]
\item If $p_1>0$ for every $p\in T$, then $S$ has a convex $2m$-gon.
\item If $m\ge3$, then $S$ has a convex $(2m-1)$-gon.
\end{enumerate}
(This lemma isolates the cap argument that the Lean proofs of Theorems~\ref{mm:thm:odd}
and~\ref{mm:thm:even} and of parts (b) and~(c) of Theorem~\ref{mm:thm:small} carry out inside their local
claims \lean{keyT} (odd case and $k=5$) and \lean{hR} (even case and $k=6$).)
\end{lemma}

\begin{proof}
$T$ is $x$-injective by Lemma~\ref{mm:lem:injx} and in general position by Lemma~\ref{fd:lem:mono}(b); the
statements about $C$ are Lemma~\ref{mm:lem:sort}. Let $f:[m]\to[\card T]$ be strictly increasing with
$\orient(C_{f(i)},C_{f(i')},C_{f(i'')})<0$ for all $i<i'<i''$ in $[m]$, and put $c_i=C_{f(i)}$ for
$i\in[m]$. Then $i\mapsto(c_i)_0$ is strictly increasing (a composition of strictly increasing maps);
$\orient(c_i,c_{i'},c_{i''})<0$ for all $i<i'<i''$; $c_i\in T\subseteq S$, hence
$\operatorname{mir}(c_i)\in S$; and $(c_i)_1\ge0$.

(a) Here $(c_i)_1>0$ for all $i$, because $c_i\in T$. Theorem~\ref{mm:thm:mirror-polygon} with $P=S$ gives
a convex $2m$-gon in $S$.

(b) If $(c_i)_1>0$ for all $i$, Theorem~\ref{mm:thm:mirror-polygon} gives a convex $2m$-gon in $S$, hence a
convex $(2m-1)$-gon by Lemma~\ref{mm:lem:ngon}(a). Otherwise there is $i_0$ with $(c_{i_0})_1\le0$, hence
$(c_{i_0})_1=0$. For $i\ne i_0$ we have $(c_i)_1>0$: otherwise $(c_i)_1=0$, the hypothesis on axis points
gives $c_i=c_{i_0}$, so $(c_i)_0=(c_{i_0})_0$ and $i=i_0$. There are three cases.
\begin{itemize}
\item $i_0=0$: Theorem~\ref{mm:thm:axis-polygon}(a) with $P=S$ gives a convex $(2m-1)$-gon in $S$.
\item $i_0=m-1$: Theorem~\ref{mm:thm:axis-polygon}(b) with $P=S$ gives a convex $(2m-1)$-gon in $S$.
\item $0<i_0<m-1$: put $p=c_{i_0-1}$, $a=c_{i_0}$ and $q=c_{i_0+1}$. Then $p_0<a_0<q_0$, $p_1,q_1\ge0$ and
  $a_1=0$, so $\orient(p,a,q)\ge0$ by Lemma~\ref{mm:lem:slope}(b); but $\orient(p,a,q)<0$ by the cap
  property for $i_0-1<i_0<i_0+1$. So this case does not occur. \qedhere
\end{itemize}
\end{proof}

\subsection{The mirror reduction}

\begin{theorem}[Theorem A, odd case; \leanref{MirrorOdd.mirror_odd_general}{MIRRORETV/lean/MirrorMain.lean:2498}]\label{mm:thm:odd}
Let $j,M\in\N$ with $2\le j$, and assume the following parametric hypothesis on $M$, which is part of the
statement (Remark~\ref{mm:rem:hm}):
\begin{itemize}
\item[(hM)] for every $n\in\N$ and every configuration $C$ of $n$ points: if $C$ has no $(j+1)$-cap and
  $\mathrm{pts}(C)$ has no convex $(2j+1)$-gon, then $n\le M$.
\end{itemize}
Let $S\subset\R^2$ be finite and in general position, with $\operatorname{mir}(p)\in S$ for every $p\in S$,
and suppose that $S$ has no convex $(2j+1)$-gon. Then $\card S\le2M$.
\end{theorem}

\begin{proof}
We follow the Lean proof.

\emph{Step~1.} Let $T\subseteq S$ be such that $p_1\ge0$ for all $p\in T$, and $p=q$ whenever $p,q\in T$
and $p_1=q_1=0$. Then $\card T\le M$. Indeed, let $C=\mathrm{cfg}(T)$, a configuration of $\card T$ points
with $\mathrm{pts}(C)=T$ (Lemma~\ref{mm:lem:upper-caps}). If $C$ had a $(j+1)$-cap,
Lemma~\ref{mm:lem:upper-caps}(b) with $m=j+1\ge3$ would give a convex $(2(j+1)-1)$-gon, that is, a convex
$(2j+1)$-gon, in $S$; so $C$ has no $(j+1)$-cap. If $\mathrm{pts}(C)=T$ had a convex $(2j+1)$-gon, so would
$S\supseteq T$ (Lemma~\ref{fd:lem:mono}(c)). Hypothesis (hM) with $n=\card T$ gives $\card T\le M$.

\emph{Step~2.} The set $T=R(S)$ satisfies the assumptions of Step~1: all its points have positive
ordinate, so none of them lies on the axis. Hence $\card{R(S)}\le M$.

\emph{Step~3.} If $A(S)=\emptyset$, then $\card S=2\card{R(S)}\le2M$ by Lemma~\ref{mm:lem:count}(a).
Otherwise pick $a\in A(S)$ and let $T=R(S)\cup\{a\}$. Then $a\notin R(S)$ (as $a_1=0$), $T\subseteq S$,
$p_1\ge0$ on $T$, and $a$ is the only point of $T$ on the axis. Step~1 gives
$\card{R(S)}+1=\card T\le M$. By Lemma~\ref{mm:lem:count}(a) and~(b),
$\card S=2\card{R(S)}+\card{A(S)}\le2\card{R(S)}+2\le2M$.
\end{proof}

\begin{theorem}[Theorem A, even case; \leanref{MirrorOdd.mirror_even_general}{MIRRORETV/lean/MirrorMain.lean:2599}]\label{mm:thm:even}
Let $j,M\in\N$ with $2\le j$, and assume the following parametric hypothesis on $M$, which is part of the
statement (Remark~\ref{mm:rem:hm}):
\begin{itemize}
\item[(hM)] for every $n\in\N$ and every configuration $C$ of $n$ points: if $C$ has no $j$-cap and
  $\mathrm{pts}(C)$ has no convex $2j$-gon, then $n\le M$.
\end{itemize}
Let $S\subset\R^2$ be finite and in general position, with $\operatorname{mir}(p)\in S$ for every $p\in S$,
and suppose that $S$ has no convex $2j$-gon. Then $\card S\le2M+2$.
\end{theorem}

\begin{proof}
We follow the Lean proof. The set $T=R(S)$ satisfies the assumptions of Lemma~\ref{mm:lem:upper-caps},
with $p_1>0$ for all $p\in T$. Let $C=\mathrm{cfg}(R(S))$. If $C$ had a $j$-cap,
Lemma~\ref{mm:lem:upper-caps}(a) with $m=j\ge2$ would give a convex $2j$-gon in $S$; so $C$ has no
$j$-cap. The set $\mathrm{pts}(C)=R(S)\subseteq S$ has no convex $2j$-gon (Lemma~\ref{fd:lem:mono}(c)).
By (hM), $\card{R(S)}\le M$, and by Lemma~\ref{mm:lem:count},
$\card S=2\card{R(S)}+\card{A(S)}\le2M+2$.
\end{proof}

\begin{remark}[the hypothesis (hM), and indexing]\label{mm:rem:hm}
We comment on the hypothesis of Theorems~\ref{mm:thm:odd} and~\ref{mm:thm:even} and of parts (b), (c) of Theorem~\ref{mm:thm:line}.
\begin{enumerate}[label=(\arabic*)]
\item \emph{(hM) is a parametric hypothesis.} $M$ is any natural number for which (hM) holds; it is not
  defined as a maximum. Theorems~\ref{mm:thm:odd}, \ref{mm:thm:even} and~\ref{mm:thm:line}(b),(c) are
  implications from (hM) to a bound on $\card S$. Thus (hM) is an assumption on the parameter $M$, not an
  unproved claim of the project: for a given pair $(j,M)$ the bound holds as soon as (hM) is proved for that
  pair. In the project, (hM) is discharged by Lean-checked theorems in the following instances. Odd case:
  $(j,M)=(3,16)$ by \lean{MirrorOdd.capacity16} (in \lean{MirrorOdd.mirror_seven'},
  \dirn{MIRRORETV/lean/MirrorMain.lean:2631}) and $(j,M)=(4,113)$ by \lean{C12Check.capacity113} (in
  \lean{C12Check.mirror_nine}, \dirn{MIRRORETV/lean/MirrorNine_append.lean:131}). Even case:
  $(j,M)=(4,22)$ by \lean{MirrorOdd.capacity22g} (in \lean{MirrorOdd.mirror_eight'},
  \dirn{MIRRORETV/lean/MirrorMain.lean:2636}). The axioms of \lean{mirror_seven'} and \lean{mirror_eight'}
  are recorded in \dirn{MIRRORETV/lean/MirrorMain.out} (lines 860--861), those of \lean{C12Check.mirror_nine}
  in \dirn{MIRRORETV/lean/MirrorNine.out} (line 881; that log checks
  \dirn{MIRRORETV/lean/MirrorNine_append.lean} appended to \dirn{MIRRORETV/lean/MirrorMain.lean}); in each
  case they are the three standard axioms. These instances rest on the repository's Lean formalization of
  Baek's theorem (the vendored theorem \lean{main}, Theorem~\ref{mb:thm:baek}), and the instance
  $(4,113)$ also on the vendored cups--caps theorem \lean{Config.cap_cup}. They are treated in the section
  of Theorems~\ref{mb:thm:seven}, \ref{mb:thm:eight} and~\ref{mb:thm:nine}: the three capacity bounds
  appear there as statements used in the proofs of these theorems, \lean{mirror_seven'} and
  \lean{mirror_eight'} in remarks after Theorems~\ref{mb:thm:seven} and~\ref{mb:thm:eight}, and
  \lean{C12Check.mirror_nine} is part~(b) of Theorem~\ref{mb:thm:nine}. For small $j$:
  Lemma~\ref{mm:lem:cup}(b) is exactly (hM) for the odd case
  $(j,M)=(2,4)$ (with $K=5$) and for the even case $(j,M)=(3,5)$ (with $K=6$); and (hM) for the even case
  $(j,M)=(2,1)$ holds because any two points of a configuration form a $2$-cap. So the bounds $8$, $12$
  and $4$ of Theorem~\ref{mm:thm:small} also follow from Theorems~\ref{mm:thm:odd}
  and~\ref{mm:thm:even}; the Lean proofs of Theorem~\ref{mm:thm:small} do not take this route.
\item \emph{Relation to the Erd\H{o}s--Tuza--Valtr function.} \dirn{MIRRORETV/README.md} and
  \dirn{MIRRORETV/PROOFS.md} state Theorem~A with Baek's number $N(n,a,b)$ in the place of $M$: the
  largest size of an $x$-distinct set in general position with no $n$ points in convex position, no
  $a$-cap and no $b$-cup, where caps and cups are defined by consecutive triples. In their own index $j'$
  (see (3)) the odd case uses $N(2j'-1,j',2j'-1)$ and the even case $N(2j',j',2j')$; in the indexing of
  Theorems~\ref{mm:thm:odd} and~\ref{mm:thm:even} these are $N(2j+1,j+1,2j+1)$ (with $j'=j+1$) and
  $N(2j,j,2j)$ (with $j'=j$). These numbers are finite: by the cups--caps theorem of Erd\H{o}s and
  Szekeres (1935), for $a,b\ge2$ every set in the class that defines $N(n,a,b)$ has at most
  $\binom{a+b-4}{a-2}$ points (\dirn{MIRRORETV/PROOFS.md}, \S1), and the two numbers above have
  $a,b\ge j\ge2$.
  That (hM) holds for $M$ equal to the corresponding number then follows, since the point set of a
  configuration of $n$ points is an $x$-distinct set of $n$ points in general position, from two facts
  stated in writing in \dirn{MIRRORETV/PROOFS.md}: consecutive clockwise triples imply that all
  increasing triples are clockwise (Lemma~2.1(b)), and a cup is in convex position (\S1; see also
  Lemma~2.3). No declaration of \lean{MirrorOdd} states this, and it is not used in this section. A
  separate Lean formalization of Theorem~A in this form, written by a review with its own definitions
  and importing only Mathlib, is
  \dirn{MIRRORETV/review/MirrorA.lean} (for example \lean{theoremA_odd_bound}, whose hypothesis is an
  upper bound on the sets of the class that defines $N(2j'-1,j',2j'-1)$). It is not restated in this
  document; the Coverage paragraph at the end of this section describes it and gives the reason.
\item \emph{Index shift.} In Theorem~\ref{mm:thm:odd}, $k=2j+1$ and the caps have $j+1$ points. The
  written sources state the odd case as $k=2j'-1$ with $j'$-caps, so the Lean $j$ equals $j'-1$; and
  $2\le j$ means $k\ge5$, so the case $k=3$ is not covered by Theorem~\ref{mm:thm:odd}. In
  Theorem~\ref{mm:thm:even}, $k=2j$ with $j$-caps, the same $j$ as in the written sources.
\item In the odd case, Step~3 adds an axis point to $R(S)$ at no cost: by
  Lemma~\ref{mm:lem:upper-caps}(b) a $(j+1)$-cap through it would still give a convex $(2j+1)$-gon. This
  yields $2M$ instead of $2M+2$. In the even case a $j$-cap with an endpoint on the axis gives only
  $2j-1$ points, and the (at most two) axis points are simply counted.
\end{enumerate}
\end{remark}

\subsection{Arbitrary mirror lines}

\begin{definition}[line reflection and frame; \leanref{MirrorOdd.lineRefl}{MIRRORETV/lean/MirrorMain.lean:1879}, \leanref{MirrorOdd.frame}{MIRRORETV/lean/MirrorMain.lean:1883}, \leanref{MirrorOdd.unframeLin}{MIRRORETV/lean/MirrorMain.lean:1886}, \leanref{MirrorOdd.unframe}{MIRRORETV/lean/MirrorMain.lean:1892}]\label{mm:def:frame}
Let $n_0,n_1,c\in\R$. For $p,q\in\R^2$ put
\begin{align*}
\mathrm{lineRefl}(p)&=\bigl(p_0-2(n_0p_0+n_1p_1-c)\,n_0,\ \ p_1-2(n_0p_0+n_1p_1-c)\,n_1\bigr),\\
\mathrm{frame}(p)&=\bigl(n_1(p_0-cn_0)-n_0(p_1-cn_1),\ \ n_0(p_0-cn_0)+n_1(p_1-cn_1)\bigr),\\
\mathrm{unframeLin}(q)&=\bigl(n_1q_0+n_0q_1,\ \ -n_0q_0+n_1q_1\bigr),\\
\mathrm{unframe}(q)&=\bigl(n_1q_0+n_0q_1+cn_0,\ \ -n_0q_0+n_1q_1+cn_1\bigr).
\end{align*}
These maps depend on $n_0,n_1,c$ (Lean \lean{lineRefl}, \lean{frame} and \lean{unframe} take the
arguments $n_0,n_1,c$; \lean{unframeLin} takes $n_0,n_1$). The map $\mathrm{unframeLin}$ is linear, and
$\mathrm{unframe}$ is the affine map with linear part $\mathrm{unframeLin}$ and translation part
$(cn_0,cn_1)$.
\end{definition}

\begin{remark}
Let $n_0^2+n_1^2=1$ and $\ell=\{p:n_0p_0+n_1p_1=c\}$; every line in $\R^2$ has this form. With
$\lambda(p)=n_0p_0+n_1p_1-c$ we have $\mathrm{lineRefl}(p)=p-2\lambda(p)(n_0,n_1)$, so
$\lambda(\mathrm{lineRefl}(p))=-\lambda(p)$ while $n_1p_0-n_0p_1$ is unchanged: $\mathrm{lineRefl}$ is the
orthogonal reflection in $\ell$. Moreover $\mathrm{frame}(p)=(n_1p_0-n_0p_1,\ \lambda(p))$ is a rotation by
the matrix $\begin{psmallmatrix}n_1&-n_0\\ n_0&n_1\end{psmallmatrix}$ followed by a translation, and it maps
$\ell$ onto the $x$-axis. This explains the definitions; only Lemma~\ref{mm:lem:frame} is used below.
\end{remark}

\begin{lemma}[frame identities; \leanref{MirrorOdd.unframe_apply}{MIRRORETV/lean/MirrorMain.lean:1897}, \leanref{MirrorOdd.unframe_frame}{MIRRORETV/lean/MirrorMain.lean:1900}, \leanref{MirrorOdd.frame_unframe}{MIRRORETV/lean/MirrorMain.lean:1907}, \leanref{MirrorOdd.frame_lineRefl}{MIRRORETV/lean/MirrorMain.lean:1914}, \leanref{MirrorOdd.orient_frame}{MIRRORETV/lean/MirrorMain.lean:1920}, \leanref{MirrorOdd.frame_injective}{MIRRORETV/lean/MirrorMain.lean:1925}, \leanref{MirrorOdd.unframe_injective}{MIRRORETV/lean/MirrorMain.lean:1930}]\label{mm:lem:frame}
Let $n_0,n_1,c\in\R$. For every $q\in\R^2$, $\mathrm{unframe}(q)$ is given by the formula of
Definition~\ref{mm:def:frame} (\lean{unframe_apply}; no hypothesis). If $n_0^2+n_1^2=1$, then for all
$p,q,a,b,d\in\R^2$:
\begin{enumerate}[label=(\alph*)]
\item $\mathrm{unframe}(\mathrm{frame}(p))=p$ (\lean{unframe_frame});
\item $\mathrm{frame}(\mathrm{unframe}(q))=q$ (\lean{frame_unframe});
\item $\mathrm{frame}(\mathrm{lineRefl}(p))=\operatorname{mir}(\mathrm{frame}(p))$ (\lean{frame_lineRefl});
\item $\orient(\mathrm{frame}(a),\mathrm{frame}(b),\mathrm{frame}(d))=\orient(a,b,d)$ (\lean{orient_frame});
\item $\mathrm{frame}$ and $\mathrm{unframe}$ are injective (\lean{frame_injective},
  \lean{unframe_injective}).
\end{enumerate}
\end{lemma}

\begin{proof}
The formula for $\mathrm{unframe}$ holds by definition (in Lean, \lean{unframe_apply} is proved by
\lean{rfl}). Put $s=n_0^2+n_1^2$ and $\lambda=n_0p_0+n_1p_1-c$. Expanding the
definition gives, with no hypothesis on $n_0,n_1,c$,
\[ \mathrm{frame}(p)=\bigl(n_1p_0-n_0p_1,\ n_0p_0+n_1p_1-cs\bigr)=Rp-(0,cs),\qquad
   R=\begin{pmatrix}n_1&-n_0\\ n_0&n_1\end{pmatrix}. \]
We prove four polynomial identities, valid for all $n_0,n_1,c$.
\begin{itemize}
\item $\mathrm{unframe}(\mathrm{frame}(p))$ has first coordinate
  $n_1(n_1p_0-n_0p_1)+n_0(n_0p_0+n_1p_1-cs)+cn_0=sp_0-csn_0+cn_0=p_0+(s-1)(p_0-cn_0)$, and second coordinate
  $-n_0(n_1p_0-n_0p_1)+n_1(n_0p_0+n_1p_1-cs)+cn_1=sp_1-csn_1+cn_1=p_1+(s-1)(p_1-cn_1)$.
\item With $X=n_1q_0+n_0q_1+cn_0$ and $Y=-n_0q_0+n_1q_1+cn_1$, we get
  $\mathrm{frame}(\mathrm{unframe}(q))=(n_1X-n_0Y,\ n_0X+n_1Y-cs)=(sq_0,\ sq_1+cs-cs)=q+(s-1)q$.
\item $\mathrm{lineRefl}(p)=p-2\lambda\,(n_0,n_1)$ and $R\,(n_0,n_1)^{\mathsf T}=(0,s)^{\mathsf T}$, so
  $\mathrm{frame}(\mathrm{lineRefl}(p))=\mathrm{frame}(p)-2\lambda\,(0,s)$. Its first coordinate equals that
  of $\mathrm{frame}(p)$, which is also that of $\operatorname{mir}(\mathrm{frame}(p))$. Its second
  coordinate minus that of $\operatorname{mir}(\mathrm{frame}(p))$ is
  $2(n_0p_0+n_1p_1-cs)-2\lambda s=-2(s-1)(n_0p_0+n_1p_1)$.
\item $\mathrm{frame}$ is the affine map $x\mapsto Rx-(0,cs)$ with $\det R=s$, so
  $\orient(\mathrm{frame}(a),\mathrm{frame}(b),\mathrm{frame}(d))=s\,\orient(a,b,d)$ by
  Lemma~\ref{fd:lem:affine-orient}.
\end{itemize}
With $s=1$ these identities give (a)--(d). (They are the certificates of the Lean proofs: the first
coordinate in (c) is closed by \lean{ring}, and every other coordinate, and (d), by
\lean{linear_combination} with the multiples $p_0-cn_0$, $p_1-cn_1$, $q_0$, $q_1$, $-2(n_0p_0+n_1p_1)$ and
$\orient(a,b,d)$ of the hypothesis $s=1$.)
(e) By (a), $\mathrm{unframe}$ is a left inverse of $\mathrm{frame}$, so $\mathrm{frame}$ is injective; by
(b), $\mathrm{frame}$ is a left inverse of $\mathrm{unframe}$, so $\mathrm{unframe}$ is injective.
\end{proof}

\begin{theorem}[transport, and Theorem A for an arbitrary mirror line; \leanref{MirrorOdd.transport}{MIRRORETV/lean/MirrorMain.lean:1937}, \leanref{MirrorOdd.mirror_odd_general_line}{MIRRORETV/lean/MirrorMain.lean:2641}, \leanref{MirrorOdd.mirror_even_general_line}{MIRRORETV/lean/MirrorMain.lean:2650}]\label{mm:thm:line}
Part (a) moves a set that is symmetric in a line to a set that is symmetric in the $x$-axis; parts (b)
and~(c) are Theorems~\ref{mm:thm:odd} and~\ref{mm:thm:even} for an arbitrary mirror line.
\begin{enumerate}[label=(\alph*)]
\item (\lean{transport}) Let $n_0,n_1,c\in\R$ with $n_0^2+n_1^2=1$, let $S\subset\R^2$ be finite, and let
  $m\in\N$. Suppose that $S$ is in general position, that $\mathrm{lineRefl}(p)\in S$ for every $p\in S$,
  and that $S$ has no convex $m$-gon. Let $S'=\mathrm{frame}(S)$. Then (i) $S'$ is in general position;
  (ii) $\operatorname{mir}(q)\in S'$ for every $q\in S'$; (iii) $S'$ has no convex $m$-gon; and
  (iv) $\card{S'}=\card S$.
\item (\lean{mirror_odd_general_line}) Let $j,M\in\N$ with $2\le j$, and assume the parametric
  hypothesis (hM) of Theorem~\ref{mm:thm:odd}. Let $n_0,n_1,c\in\R$ with $n_0^2+n_1^2=1$, and let
  $S\subset\R^2$ be finite and in general position, with $\mathrm{lineRefl}(p)\in S$ for every $p\in S$
  and with no convex $(2j+1)$-gon. Then $\card S\le2M$.
\item (\lean{mirror_even_general_line}) Let $j,M\in\N$ with $2\le j$, and assume the parametric
  hypothesis (hM) of Theorem~\ref{mm:thm:even}. Let $n_0,n_1,c\in\R$ with $n_0^2+n_1^2=1$, and let
  $S\subset\R^2$ be finite and in general position, with $\mathrm{lineRefl}(p)\in S$ for every $p\in S$
  and with no convex $2j$-gon. Then $\card S\le2M+2$.
\end{enumerate}
Here $\mathrm{lineRefl}$ and $\mathrm{frame}$ are the maps of Definition~\ref{mm:def:frame} for the given
$n_0,n_1,c$.
\end{theorem}

\begin{proof}
(a) (i) Let $a',b',d'\in S'$ be pairwise distinct, say $a'=\mathrm{frame}(a)$, $b'=\mathrm{frame}(b)$,
$d'=\mathrm{frame}(d)$ with $a,b,d\in S$. Then $a,b,d$ are pairwise distinct (equal points have equal
images), so $\orient(a,b,d)\ne0$ by Lemma~\ref{fd:lem:collinear}, and
$\orient(a',b',d')=\orient(a,b,d)\ne0$ by Lemma~\ref{mm:lem:frame}(d). By Lemma~\ref{fd:lem:collinear},
$S'$ is in general position.
(ii) If $q=\mathrm{frame}(p)$ with $p\in S$, then $\operatorname{mir}(q)=\mathrm{frame}(\mathrm{lineRefl}(p))$
by Lemma~\ref{mm:lem:frame}(c), and $\mathrm{lineRefl}(p)\in S$.
(iii) Suppose that $S'$ has a convex $m$-gon. The map $\mathrm{unframe}$ is affine and injective
(Lemma~\ref{mm:lem:frame}(e)), so $\mathrm{unframe}(S')$ has a convex $m$-gon by
Lemma~\ref{mm:lem:ngon}(b). But $\mathrm{unframe}(S')=\{\mathrm{unframe}(\mathrm{frame}(p)):p\in S\}=S$ by
Lemma~\ref{mm:lem:frame}(a), a contradiction.
(iv) $\mathrm{frame}$ is injective (Lemma~\ref{mm:lem:frame}(e)).

(b) By (a) with $m=2j+1$, the finite set $S'=\mathrm{frame}(S)$ is in general position, is
$\operatorname{mir}$-symmetric, has no convex $(2j+1)$-gon, and $\card{S'}=\card S$.
Theorem~\ref{mm:thm:odd} (with the same $j$, $M$ and (hM)) gives $\card S=\card{S'}\le2M$.

(c) The same with $m=2j$ and Theorem~\ref{mm:thm:even}. This is also the structure of the Lean proofs.
\end{proof}

\subsection{Small \texorpdfstring{$k$}{k}}

\begin{lemma}[configurations without a $3$-cap; \leanref{MirrorOdd.hasCap3_of}{MIRRORETV/lean/MirrorMain.lean:1991}, \leanref{MirrorOdd.capacity_cup}{MIRRORETV/lean/MirrorMain.lean:2002}]\label{mm:lem:cup}
Both parts concern configurations in the sense of Definition~\ref{mm:def:config}.
\begin{enumerate}[label=(\alph*)]
\item (\lean{hasCap3_of}) Let $C$ be a configuration of $n$ points, and let $a<b<d$ in $[n]$ with
  $\orient(C_a,C_b,C_d)<0$. Then $C$ has a $3$-cap.
\item (\lean{capacity_cup}) Let $n,K\in\N$ with $4\le K$, and let $C$ be a configuration of $n$ points
  that has no $3$-cap and such that $\mathrm{pts}(C)$ has no convex $K$-gon. Then $n\le K-1$.
\end{enumerate}
\end{lemma}

\begin{proof}
(a) The map $f:[3]\to[n]$ with $f(0)=a$, $f(1)=b$, $f(2)=d$ is strictly increasing, and the only triple
$i<j<l$ in $[3]$ is $(0,1,2)$, for which the required inequality is the hypothesis.

(b) Suppose $n>K-1$, that is, $n\ge K$. The map $i\mapsto C_i$ is injective because the abscissae
increase strictly. Let $e:[K]\to[n]$ be the inclusion $i\mapsto i$ (Lean \lean{Fin.castLEEmb}). For
$i<j<l$ in $[K]$ we have $\orient(C_i,C_j,C_l)\ne0$ (Definition~\ref{mm:def:config}(a)), and
$\orient(C_i,C_j,C_l)<0$ is impossible because by (a) it would give a $3$-cap; hence
$\orient(C_i,C_j,C_l)>0$. Lemma~\ref{fd:lem:pos-cycle} with $N=n$, $k=K\ge4$, $p=C$, this $e$ and the set
$\mathrm{pts}(C)$ (which contains every $C_{e(i)}$) gives a convex $K$-gon in $\mathrm{pts}(C)$, a
contradiction.
\end{proof}

\begin{theorem}[$k=4,5,6$; \leanref{MirrorOdd.mirror_four}{MIRRORETV/lean/MirrorMain.lean:2019}, \leanref{MirrorOdd.mirror_five}{MIRRORETV/lean/MirrorMain.lean:2040}, \leanref{MirrorOdd.mirror_six}{MIRRORETV/lean/MirrorMain.lean:2129}, \leanref{MirrorOdd.mirror_four_line}{MIRRORETV/lean/MirrorMain.lean:2159}, \leanref{MirrorOdd.mirror_five_line}{MIRRORETV/lean/MirrorMain.lean:2165}, \leanref{MirrorOdd.mirror_six_line}{MIRRORETV/lean/MirrorMain.lean:2171}]\label{mm:thm:small}
In (a)--(c) let $S\subset\R^2$ be finite and in general position, with $\operatorname{mir}(p)\in S$ for
every $p\in S$. In (d)--(f) let $n_0,n_1,c\in\R$ with $n_0^2+n_1^2=1$, and let $S\subset\R^2$ be finite
and in general position, with $\mathrm{lineRefl}(p)\in S$ for every $p\in S$.
\begin{enumerate}[label=(\alph*)]
\item (\lean{mirror_four}) If $S$ has no convex $4$-gon, then $\card S\le4$.
\item (\lean{mirror_five}) If $S$ has no convex $5$-gon, then $\card S\le8$.
\item (\lean{mirror_six}) If $S$ has no convex $6$-gon, then $\card S\le12$.
\item (\lean{mirror_four_line}) If $S$ has no convex $4$-gon, then $\card S\le4$.
\item (\lean{mirror_five_line}) If $S$ has no convex $5$-gon, then $\card S\le8$.
\item (\lean{mirror_six_line}) If $S$ has no convex $6$-gon, then $\card S\le12$.
\end{enumerate}
These statements have no hypotheses beyond the displayed ones.
\end{theorem}

\begin{proof}
We follow the Lean proofs.

(a) First, $\card{R(S)}\le1$. Otherwise there are $p\ne q$ in $R(S)$. Lemma~\ref{mm:lem:injx}, applied with
$T=R(S)$ (all its points have positive ordinate, so the condition on axis points is vacuous), gives
$p_0\ne q_0$; exchanging the names if necessary, $p_0<q_0$. Proposition~\ref{mm:prop:small-polygons}(a)
with $(c_0,c_1)=(p,q)$ and $P=S$ applies ($p_1,q_1>0$, and $\operatorname{mir}p,\operatorname{mir}q,q,p\in S$)
and gives a convex $4$-gon in $S$, a contradiction. By Lemma~\ref{mm:lem:count},
$\card S=2\card{R(S)}+\card{A(S)}\le2+2=4$.

(b) \emph{Step~1:} if $T\subseteq S$ has $p_1\ge0$ for all $p\in T$ and at most one point on the axis,
then $\card T\le4$. Let $C=\mathrm{cfg}(T)$ (Lemma~\ref{mm:lem:upper-caps}). $C$ has no $3$-cap, because
otherwise Lemma~\ref{mm:lem:upper-caps}(b) with $m=3$ gives a convex $5$-gon in $S$. (The Lean proof makes
the same case distinction as the proof of Lemma~\ref{mm:lem:upper-caps}(b) for the cap $c_0,c_1,c_2$, but
applies Proposition~\ref{mm:prop:small-polygons}(b) when the three points lie strictly above the axis;
Proposition~\ref{mm:prop:small-polygons}(g), with the points $a,c_0,c_1$ of the Proposition taken to be
$c_2,c_0,c_1$, when $c_2$ lies on the axis; and Proposition~\ref{mm:prop:small-polygons}(f), with
$a,c_0,c_1$ taken to be $c_0,c_1,c_2$, when $c_0$ lies on the axis. The case that $c_1$ lies on the axis
is excluded by Lemma~\ref{mm:lem:slope}(b).) Moreover $\mathrm{pts}(C)=T\subseteq S$ has no convex $5$-gon.
Lemma~\ref{mm:lem:cup}(b) with $K=5$ gives $\card T\le4$.
\emph{Step~2:} $\card{R(S)}\le4$ by Step~1.
\emph{Step~3:} if $A(S)=\emptyset$, then $\card S=2\card{R(S)}\le8$. Otherwise, for $a\in A(S)$, Step~1
applied to $R(S)\cup\{a\}$ (as in Step~3 of Theorem~\ref{mm:thm:odd}) gives $\card{R(S)}+1\le4$, and
$\card S\le2\cdot3+2=8$ by Lemma~\ref{mm:lem:count}.

(c) Let $C=\mathrm{cfg}(R(S))$ (Lemma~\ref{mm:lem:upper-caps} with $T=R(S)$). Suppose $C$ has a $3$-cap,
given by $f$, and put $c_i=C_{f(i)}$; then $(c_0)_0<(c_1)_0<(c_2)_0$, all $(c_i)_1>0$,
$\orient(c_0,c_1,c_2)<0$, and $S$ contains $c_i$ and $\operatorname{mir}c_i$. Then
Proposition~\ref{mm:prop:small-polygons}(c) with $P=S$ (equivalently, Lemma~\ref{mm:lem:upper-caps}(a)
with $m=3$) gives a convex $6$-gon in $S$, a contradiction. Also $\mathrm{pts}(C)=R(S)\subseteq S$ has no
convex $6$-gon. Lemma~\ref{mm:lem:cup}(b) with $K=6$ gives $\card{R(S)}\le5$, so
$\card S\le2\cdot5+2=12$ by Lemma~\ref{mm:lem:count}.

(d)--(f) By Theorem~\ref{mm:thm:line}(a) with $m=4,5,6$, the set $S'=\mathrm{frame}(S)$ satisfies the
hypotheses of (a)--(c), respectively, and $\card{S'}=\card S$.
\end{proof}

\paragraph{Coverage.}
{\raggedright
Covered, all in namespace \lean{MirrorOdd} of \dirn{MIRRORETV/lean/MirrorMain.lean} unless stated
otherwise:
\lean{mir}, \lean{Rset}, \lean{Aset} (Definition~\ref{mm:def:mir});
\lean{mir_zero}, \lean{mir_one}, \lean{mir_mir}, \lean{orient_mir3}, \lean{orient_mir2},
\lean{orient_cycl}, \lean{orient_swap12}, \lean{orient_rev} (Lemma~\ref{mm:lem:orient-mir});
\lean{ErdosSzekeres.Omission.Configuration} and \lean{Configuration.HasCap}
(\dirn{ErdosSzekeres/Omission/Realization.lean}), \lean{xset}, \lean{emb}, \lean{sortPt}, \lean{cfg}
(Definition~\ref{mm:def:config});
\lean{xset_card}, \lean{emb_mem}, \lean{sortPt_mem}, \lean{sortPt_x}, \lean{sortPt_strictMono},
\lean{sortPt_injective}, \lean{sortPt_image} (Lemma~\ref{mm:lem:sort});
\lean{hasConvexNGon_of_le}, \lean{hasConvexNGon_affineMap} (Lemma~\ref{mm:lem:ngon});
\lean{cap_below_lines} (Lemma~\ref{mm:lem:slope});
\lean{cap_mirror_polygon} (Theorem~\ref{mm:thm:mirror-polygon});
\lean{axis_left_polygon}, \lean{axis_right_polygon} (Theorem~\ref{mm:thm:axis-polygon});
\lean{cyc4}--\lean{cyc8} (Lemma~\ref{mm:lem:cyc});
\lean{cap_mirror_quadrilateral}, \lean{cap_mirror_pentagon}, \lean{cap_mirror_hexagon},
\lean{cap_mirror_heptagon}, \lean{cap_mirror_octagon}, \lean{axis_left_pentagon},
\lean{axis_right_pentagon}, \lean{axis_left_heptagon}, \lean{axis_right_heptagon}
(Proposition~\ref{mm:prop:small-polygons});
\lean{card_split}, \lean{card_A_le} (Lemma~\ref{mm:lem:count}); \lean{injOn_x}
(Lemma~\ref{mm:lem:injx});
\lean{mirror_odd_general} (Theorem~\ref{mm:thm:odd}); \lean{mirror_even_general}
(Theorem~\ref{mm:thm:even});
\lean{lineRefl}, \lean{frame}, \lean{unframeLin}, \lean{unframe} (Definition~\ref{mm:def:frame});
\lean{unframe_apply}, \lean{unframe_frame}, \lean{frame_unframe}, \lean{frame_lineRefl},
\lean{orient_frame}, \lean{frame_injective}, \lean{unframe_injective} (Lemma~\ref{mm:lem:frame});
\lean{transport}, \lean{mirror_odd_general_line}, \lean{mirror_even_general_line} (Theorem~\ref{mm:thm:line});
\lean{hasCap3_of}, \lean{capacity_cup} (Lemma~\ref{mm:lem:cup});
\lean{mirror_four}, \lean{mirror_five}, \lean{mirror_six} and their \lean{_line} versions
(Theorem~\ref{mm:thm:small}).
Three items are written steps with no Lean declaration of their own: Lemma~\ref{mm:lem:upper-caps}
isolates the cap argument of the local claims \lean{keyT} and \lean{hR}; part~(b) of
Lemma~\ref{mm:lem:slope} states the local claim \lean{interior} and the inline interior-point argument of
the Lean proofs; and part~(e) of Lemma~\ref{mm:lem:orient-mir} is an identity that the Lean proof of
\lean{cap_below_lines} uses implicitly when it expands $\orient$ and closes each case by
\lean{nlinarith}.

Not covered here, deliberately, from \dirn{MIRRORETV/lean/MirrorMain.lean}:
\begin{itemize}
\item \lean{hasNGon_heptagon} (line 1512), \lean{capacity16} (1529), \lean{mirror_seven_of} (1676),
  \lean{mirror_seven} (1803), \lean{capacity22g} (1812), \lean{mirror_eight} (1820),
  \lean{mirror_seven_line} (1969) and \lean{mirror_eight_line} (1977). \lean{hasNGon_heptagon} is the
  bridge, used by \lean{capacity16}, from a $7$-gon of the combinatorial model of a configuration
  (\lean{colorConfig}) to a convex $7$-gon of its point set. The others rest on Baek's theorem (the
  vendored theorem \lean{main}): \lean{capacity16} applies it directly, \lean{capacity22g} through
  \lean{geometric_capacity22}, and the remaining five through these two. They are treated in the section
  of Theorems~\ref{mb:thm:seven} and~\ref{mb:thm:eight}: \lean{mirror_seven}, \lean{mirror_eight} and
  their \lean{_line} versions are these two theorems, and the other four are statements used in their
  proofs.
\item \lean{mirror_seven'} (2631) and \lean{mirror_eight'} (2636): consistency re-derivations of the bounds
  $32$ and $46$. They are Theorem~\ref{mm:thm:odd} with $(j,M)=(3,16)$ and Theorem~\ref{mm:thm:even} with
  $(j,M)=(4,22)$, with (hM) supplied by \lean{capacity16} and \lean{capacity22g} (see
  Remark~\ref{mm:rem:hm}); the bounds themselves are Theorems~\ref{mb:thm:seven}
  and~\ref{mb:thm:eight}.
\item The \texttt{\#print axioms} lines 2662--2677: audit commands; their output is the axioms evidence
  cited at the beginning of this section.
\end{itemize}
Not covered here, deliberately, from the other Lean files of \dirn{MIRRORETV}:
\begin{itemize}
\item The five declarations of \dirn{MIRRORETV/lean/MirrorNine_append.lean} (namespace
  \lean{C12Check}). They are treated in the section of Theorem~\ref{mb:thm:nine}, and
  \lean{config_capacity113}, \lean{mirror_nine} and \lean{mirror_nine_line} are that theorem.
\item The independent formalizations \dirn{MIRRORETV/review/MirrorA.lean} and
  \dirn{MIRRORETV/review/MirrorCap.lean}, written by reviews. Each imports only Mathlib and has its own
  definitions, and no Lean file of the project imports either of them. No statement or proof of this
  document uses them. Their content overlaps that of
  Theorems~\ref{mm:thm:mirror-polygon}, \ref{mm:thm:axis-polygon}, \ref{mm:thm:odd}, \ref{mm:thm:even}
  and~\ref{mm:thm:line}, but in another formulation: points in $\R\times\R$, caps and cups defined by
  consecutive triples, and an upper bound on the class that defines $N(n,a,b)$
  (Remark~\ref{mm:rem:hm}(2)) in place of (hM). Beyond this overlap they contain the odd case $k=3$,
  which Theorem~\ref{mm:thm:odd} excludes (Remark~\ref{mm:rem:hm}(3)); the mirror lemma for caps with
  both endpoints on the axis; the convex position of caps and of cups; and the coincidence statement
  \lean{MirrorCap.coincidences}.
  \begin{itemize}
  \item \dirn{MirrorA.lean} (891 lines, \texttt{1a8a62cb4385}\dots; namespace \lean{MirrorA}; 16
    definitions, 55 theorems and lemmas, and one \texttt{example}). The definitions are \lean{orient}, the
    reflection \lean{mu} in the $x$-axis and its linear form \lean{muL}, \lean{ConvexPos}, the caps and cups
    \lean{IsCapSeq} and \lean{IsCupSeq}, the set \lean{V} of the points of a cap and their mirror images,
    \lean{GenPos}, \lean{XDistinct}, \lean{HasGon}, \lean{HasCap}, \lean{HasCup}, the class \lean{InClass}
    that defines $N(n,a,b)$, \lean{MirrorSym}, \lean{refl}, the reflection in the line through a point
    with a given unit direction, and \lean{isoT}, the rigid motion that maps this line onto the $x$-axis.
    The main results are \lean{lemma31} (282: a cap, in the
    sense of \lean{IsCapSeq}, all of whose points have ordinates $\ge0$, together with its mirror image,
    is in convex position; the review's form of \dirn{MIRRORETV/PROOFS.md}, Lemma~3.1),
    \lean{cap_convexPos} (352) and \lean{cup_convexPos} (417) (the points of a cap, resp.\ of a cup, are in
    convex position), \lean{theoremA_even} (593) and \lean{theoremA_odd} (615) (in the indexing of the
    written sources, for $j\ge2$: if $S$ is symmetric in the $x$-axis, in general position, and has no $2j$, resp.\
    $2j-1$, points in convex position, then some set $X$ in the class of $N(2j,j,2j)$, resp.\
    $N(2j-1,j,2j-1)$, satisfies $\card S\le2\card X+2$, resp.\ $\card S\le2\card X$),
    \lean{theoremA_even_bound} (685) and \lean{theoremA_odd_bound} (694) (the bounds $2N_b+2$ and $2N_b$
    from an upper bound $N_b$ on that class), \lean{transport} (771; the transport from an arbitrary
    mirror line to the $x$-axis by the rigid motion \lean{isoT}), and \lean{theoremA_even_line} (810)
    and \lean{theoremA_odd_line} (821) (the two bounds for the reflection in an arbitrary line). With
    $j=2$ the odd case is $k=3$. The other 45 lemmas are auxiliary, among them three checks of the
    definition of \lean{refl} (\lean{refl_refl}, \lean{refl_fix_line}, \lean{refl_normal}); the
    \texttt{example} (852) checks that \lean{ConvexPos} is not vacuous. The recorded log
    \dirn{MIRRORETV/review/MirrorA.out} prints the three standard axioms for eight of these theorems.
  \item \dirn{MirrorCap.lean} (354 lines, \texttt{a5c6dadb5ac6}\dots; namespace \lean{MirrorCap}; 16
    declarations): the definitions \lean{orient}, \lean{sl} (the slope of an edge) and \lean{V} (the
    points of the cap and their mirror images), ten auxiliary lemmas and theorems (\lean{edge_eq},
    \lean{orient_eq_slopes}, \lean{slope_dec_of_orient}, \lean{sl_anti}, \lean{x_mono}, \lean{chord_right},
    \lean{chord_left}, \lean{y_nonneg}, \lean{exposed_c}, \lean{exposed_mu}), and \lean{convex_position}
    (292), \lean{convexIndependent} (316) and \lean{coincidences} (327). For a cap $c_0,\dots,c_{m-1}$ given
    by coordinates (strictly increasing abscissae, consecutive triples clockwise) whose interior points
    have positive ordinates and whose endpoints have nonnegative ordinates (positive if $m=1$),
    \lean{convex_position} says that no point of $V$ lies in the convex hull of the other points of $V$
    (the form of \dirn{MIRRORETV/PROOFS.md}, Lemma~3.1), and \lean{convexIndependent} says the same in
    Mathlib's terms. From the monotonicity of the abscissae and the positivity of the interior ordinates,
    \lean{coincidences} derives that $c_i=c_k$ or $\operatorname{mir}c_i=\operatorname{mir}c_k$ forces
    $i=k$, and that $c_i=\operatorname{mir}c_k$ forces $i=k$, $(c_i)_1=0$ and $i\in\{0,m-1\}$; by the
    file's comment, this gives $\card V=2m-e$, where $e$ is the number of endpoints of the cap on the axis.
    The recorded log \dirn{MIRRORETV/review/MirrorCap.out} prints the three standard axioms for these three
    theorems.
  \end{itemize}
\item The far-flat files \dirn{MIRRORETV/lean/FarFlat.lean} (\texttt{08a79f881eea}\dots) and
  \dirn{MIRRORETV/review/FarFlatExtra.lean} (144 lines, \texttt{d30a9b874ba5}\dots). Both import only
  Mathlib, and no Lean file of the project imports them. They hold estimates for the far-flat lower-bound
  construction (Theorem~B of \dirn{MIRRORETV/README.md}). This document treats no part of Theorem~B, so
  no section restates either file.
  \begin{itemize}
  \item \dirn{FarFlat.lean}: the root-namespace definition \lean{orient} on six reals (line 12),
    \lean{orient_wzq_neg} (15), \lean{orient_uqz_neg} (34), \lean{orient_cyclic} (55), \lean{orient_sum}
    (60), \lean{inside_triangle} (67) and \lean{line_misses} (89): three orientation estimates
    (\lean{orient_wzq_neg}, \lean{orient_uqz_neg}, \lean{line_misses}), two identities
    (\lean{orient_cyclic}, \lean{orient_sum}) and a barycentric representation (\lean{inside_triangle}).
    The recorded log \dirn{MIRRORETV/lean/FarFlat.out} prints the axioms of five of these six theorems
    (all but \lean{orient_cyclic}), and they are the standard three.
  \item \dirn{FarFlatExtra.lean} (namespace \lean{FFX}): the definition \lean{orient} on six reals (line
    20), \lean{orient_sum} (22), \lean{bary_of_neg} (28), \lean{bary_of_pos} (44), \lean{gp_cross_axis}
    (60), \lean{a1_inside} (69) and \lean{lemma51_core} (91). Here \lean{orient_sum} and
    \lean{bary_of_neg} have the same statements as \lean{orient_sum} and \lean{inside_triangle} of
    \dirn{FarFlat.lean}; \lean{bary_of_pos} is the barycentric representation for three positive
    orientations; \lean{gp_cross_axis} is a general-position estimate for a point above the axis, a point
    below it and a point of the axis to the right of both; \lean{a1_inside} is the counting step of
    Theorem~B; and \lean{lemma51_core} is the composed core of Lemma~5.1 of \dirn{MIRRORETV/PROOFS.md}.
    The recorded log \dirn{MIRRORETV/review/FarFlatExtra.out} prints the axioms of five of these six
    theorems (all but \lean{orient_sum}), the standard three, and four warnings about unused variables.
  \end{itemize}
\item \dirn{MIRRORETV/review/ReplayRepo.lean} (118 lines, \texttt{20e479e8b3ab}\dots) contains no
  definition or theorem, only a command \texttt{\#replay\_repo}, defined by \texttt{elab} at line~51 and
  run at line~118. By its header comment, it imports only library modules, reads the compiled repository
  modules without importing them, and sends the repository constants in the dependency closure of those
  that \dirn{MIRRORETV/lean/MirrorMain.lean} references directly to Lean's kernel again. The recorded log
  \dirn{MIRRORETV/review/ReplayRepo.out} reports $497$ such constants, none of them an axiom, and a
  successful replay. It is an audit and states no theorem, so there is nothing to restate.
\end{itemize}
\par}

% End of the local typesetting group opened after the section heading.
\endgroup

\section{Mirror bounds via Baek's theorem and the three-class reduction}\label{mb:sec}

\paragraph{What this section proves.}
Let $S\subset\R^2$ be a finite set in general position that is mapped into itself by the reflection in a line. We prove
that $\card S\le 32$ if $S$ contains no $7$ points in convex position, $\card S\le 46$ if it contains no $8$ points in
convex position, and $\card S\le 226$ if it contains no $9$ points in convex position
(Theorems~\ref{mb:thm:seven}, \ref{mb:thm:eight} and~\ref{mb:thm:nine}; each has a version for the $x$-axis and a version
for an arbitrary line). The bounds for $7$ and $8$ rest on Baek's theorem on $4$-caps, the bound for $9$ in addition on
the cups--caps theorem of Erd\H{o}s and Szekeres. These two theorems are \emph{cited}, in the exact form in which the Lean
code uses them (Theorems~\ref{mb:thm:baek} and~\ref{mb:thm:capcup}). Everything between their abstract, purely
combinatorial statements and the geometric notions is proved here in full (written proofs, not refereed): the bridge
from Baek's configurations to points of the plane (Subsection~\ref{mb:ssec:bridge}), the capacity bounds $16$ and $22$
(Subsection~\ref{mb:ssec:cap}) and $113$ (Subsection~\ref{mb:ssec:nine}) and the mirror reductions
(Subsections~\ref{mb:ssec:red}--\ref{mb:ssec:nine}); the latter use the general mirror lemma
and the general reductions of Theorems~\ref{mm:thm:mirror-polygon}, \ref{mm:thm:axis-polygon}, \ref{mm:thm:odd}
and~\ref{mm:thm:line}, which are proved elsewhere in this document. Finally (Subsection~\ref{mb:ssec:three})
we identify the two parametrisations of line reflections that occur in the Lean files of this section
(Theorem~\ref{mb:thm:refl32}), prove the two mirror case classes of Definition~\ref{sr:def:classes}
(Theorem~\ref{mb:thm:classes}), and restate the reduction Theorem~\ref{sr:thm:main} and part~(a) of
Theorem~\ref{ar:thm:main} (the \lean{AffineEquiv} form) with the three remaining case classes as their only case-class
hypotheses (Theorem~\ref{mb:thm:main3}, which is
\emph{conditional} on these three statements; they are not proved in Lean or in this document, see
Remark~\ref{mb:rem:main3}).

\paragraph{Sources.}
Paths beginning with \dirn{ErdosSzekeres/} or \dirn{docs/} are relative to the repository root; all other paths (such
as \dirn{MIRRORETV/...}, \dirn{SYMLEAN/...} and the record directories named in Remark~\ref{mb:rem:main3}) are relative
to the results directory of the project. A file name without a directory, such as \dirn{MirrorMain.lean}, is short for
the file of that name whose full path is given in this section. The declarations treated here are in the following
files (SHA-256, first $12$ hexadecimal digits):
\begin{itemize}
\item \dirn{MIRRORETV/lean/MirrorMain.lean} (\texttt{ca623e55a1b9}), namespace \lean{MirrorOdd};
\item \dirn{MIRRORETV/lean/MirrorNine_append.lean} (\texttt{c2afca380211}), namespace \lean{C12Check}; this file is not
  standalone: it was compiled appended to \dirn{MirrorMain.lean}; line numbers below refer to the file itself;
\item \dirn{SYMLEAN/mirror/ReductionThreeClasses.lean} (\texttt{0cc83994fe8c}), called RTC below; its Part~C (lines
  4238--4322) is treated in Subsection~\ref{mb:ssec:three};
\item the vendored formalization of Baek's theorem (adapted from Baek's Lean code, Apache-2.0 licence):
  \begin{itemize}
  \item \dirn{ErdosSzekeres/Vendor/ETV/Main/Main.lean} (\texttt{83dd0a214204}),
  \item \dirn{ErdosSzekeres/Vendor/ETV/Main/CapCup.lean} (\texttt{fb8fbefac260}),
  \item \dirn{ErdosSzekeres/Vendor/ETV/Config/Defs.lean} (\texttt{ce2ccf7c1d28}),
  \item \dirn{ErdosSzekeres/Vendor/ETV/Config/Lemmas.lean} (\texttt{5395e27a82cc}),
  \item \dirn{ErdosSzekeres/Vendor/ETV/Lib/List/Defs.lean} (\texttt{46d750763528});
  \end{itemize}
\item the bridge files: \dirn{ErdosSzekeres/Omission/Realization.lean} (\texttt{ae4970e0808e}), which is the project's
  own code (it has no adaptation header; \dirn{docs/THIRD_PARTY.md} says that its connections were added in the
  project, using the strategy of M.~Schroeder's archive~\cite{Sch26}), and six files adapted from that archive (MIT
  licence):
  \begin{itemize}
  \item \dirn{ErdosSzekeres/Omission/ActualRanks.lean} (\texttt{3788f5149b3a}),
  \item \dirn{ErdosSzekeres/Omission/PlanarChains.lean} (\texttt{21161d053dc0}),
  \item \dirn{ErdosSzekeres/Omission/ListGeometry.lean} (\texttt{1a03c283176b}),
  \item \dirn{ErdosSzekeres/Omission/BoundaryGeometry.lean} (\texttt{f55a711e7aa9}),
  \item \dirn{ErdosSzekeres/Omission/BaekGeometry.lean} (\texttt{db8ac3c30dc0}),
  \item \dirn{ErdosSzekeres/Omission/Baek.lean} (\texttt{3433c3c2124d}).
  \end{itemize}
\end{itemize}

\paragraph{Lean version and axioms.}
The files were compiled with Lean 4.24.0 and Mathlib \texttt{f897ebcf}. The recorded \texttt{\#print axioms} outputs
report exactly the three standard axioms \lean{propext}, \lean{Classical.choice}, \lean{Quot.sound} for the following
declarations.
\begin{itemize}
\item \dirn{MIRRORETV/lean/MirrorMain.out} (\texttt{c6024155067e}), lines 848--863: the $16$ main theorems of
  \lean{MirrorOdd}, among them \lean{mirror_seven} (line 848), \lean{mirror_eight} (849), \lean{mirror_seven_line} (850),
  \lean{mirror_eight_line} (851), \lean{mirror_odd_general} (858) and \lean{mirror_odd_general_line} (862). The log
  records no error.
\item \dirn{MIRRORETV/lean/MirrorNine.out} (\texttt{976fda168e70}), lines 879--882: \lean{C12Check.config_capacity113},
  \lean{C12Check.capacity113}, \lean{C12Check.mirror_nine}, \lean{C12Check.mirror_nine_line}. This log records the
  compilation of the concatenation of \dirn{MirrorMain.lean} and \dirn{MirrorNine_append.lean}; in it, the lines of the
  appended file are shifted by $2677$.
\item \dirn{SYMLEAN/mirror/axioms.out} (\texttt{ab92257c6f17}): the six Part~C theorems \lean{reflLine_eq_lineRefl},
  \lean{card_le_32_of_reflLine}, \lean{mirrorOneAxis33_proved}, \lean{mirrorNoAxis34_proved},
  \lean{heptagon_of_nontrivial_isometry_3}, \lean{heptagon_of_nontrivial_affineEquiv_3}, and the $16$ \lean{MirrorOdd}
  theorems in their copy inside RTC. This file is a filtered excerpt of the compile output: some lines are repeated,
  and of the five \texttt{\#check} commands of RTC (lines 4312--4316) only the two for the \lean{_3} theorems have
  output in it (lines 17--20, the first one again at lines 39--40), truncated after the binders \lean{HalfTurn34} and
  \lean{Threefold33}.
\end{itemize}
Since \texttt{\#print axioms} reports the axioms of the whole dependency cone, these lines also cover the declarations
that are used but not printed individually. Every other Lean declaration named in the head of a statement below is
used by one of the printed theorems, for example \lean{capacity16}, \lean{capacity22g}, \lean{hasNGon_heptagon}, the
bridge lemmas of Subsection~\ref{mb:ssec:bridge}, and the vendored theorems \lean{main} and \lean{Config.cap_cup}. In
the three files of this section the tactic
\lean{decide} occurs only on comparisons of small indices and on the binomial values $\binom62=15$, $\binom82=28$,
$\binom93=84$, which the kernel evaluates; no \lean{native_decide}, \lean{sorry} or added axiom occurs in them or in the
vendored and bridge files listed above. The logs were produced by the project and were not
re-run for this document.

\paragraph{Copies of the same code.}
RTC lines 1562--4236 are byte-identical to \dirn{MIRRORETV/lean/MirrorMain.lean} lines 3--2677, so line $L$ of
\dirn{MirrorMain.lean} is line $L+1559$ of RTC; we cite \lean{MirrorOdd} declarations by their lines in
\dirn{MirrorMain.lean}. RTC lines 38--1559 are byte-identical to \dirn{SYMLEAN/AffineReduction.lean} lines 12--1533, and
RTC lines 38--1084 to \dirn{SYMLEAN/SymmetricReduction.lean} lines 12--1058 (so a line $L$ of
\dirn{SymmetricReduction.lean} is line $L+26$ of RTC). These identities can be checked with \texttt{sed} and \texttt{diff}.

\paragraph{Dependencies and rigour.}
We use the foundations (Definitions~\ref{fd:def:point}, \ref{fd:def:orient}, \ref{fd:def:gp}, \ref{fd:def:cp},
\ref{fd:def:ngon}, \ref{fd:def:inv} and Lemmas~\ref{fd:lem:orient}, \ref{fd:lem:affine-orient}, \ref{fd:lem:collinear},
\ref{fd:lem:cp-hull}, \ref{fd:lem:mono}, \ref{fd:lem:affine-cp}, \ref{fd:lem:separation}), the mirror lemma and the
general reductions of Definition~\ref{mm:def:config} and Theorems~\ref{mm:thm:mirror-polygon},
\ref{mm:thm:axis-polygon}, \ref{mm:thm:odd}, \ref{mm:thm:line}, and in
Subsection~\ref{mb:ssec:three} Definitions~\ref{sr:def:iso}, \ref{sr:def:classes} and Theorems~\ref{sr:thm:main},
\ref{ar:thm:main}. Three remarks refer to further results that no proof of this section uses:
Remark~\ref{mb:rem:generated} (on the Lean proofs of Lemma~\ref{mb:lem:polygons}) to Lemma~\ref{fd:lem:pos-cycle},
Remark~\ref{mb:rem:eight} (on the second Lean proof \texttt{mirror\_eight\textquotesingle} of
Theorem~\ref{mb:thm:eight}(a)) to Theorem~\ref{mm:thm:even},
and Remark~\ref{mb:rem:main3} (on the hypotheses of Theorem~\ref{mb:thm:main3}) to Theorems~\ref{ll:thm:layers},
\ref{lg:thm:nesting}, \ref{cb:thm:c2} and~\ref{cb:thm:c3}. Every Lean declaration stated below is
\emph{Lean-checked} (axioms evidence above); the proofs given here are \emph{written proofs (not refereed)}.
Theorems~\ref{mb:thm:baek} and~\ref{mb:thm:capcup} are cited published theorems; their vendored formalizations are
Lean-checked as part of the dependency cone above, but this document does not reprove them. Theorem~\ref{mb:thm:main3}
takes as hypotheses three propositions that are not proved in Lean or in this document.

\paragraph{Conventions.}
Points are $p=(p_0,p_1)\in\R^2$ (Lean \lean{Point}, Definition~\ref{fd:def:point}); we also write $x(p):=p_0$ and
$y(p):=p_1$, which avoids double subscripts for points that carry an index. $\orient$ is as in
Definition~\ref{fd:def:orient}. We write
$[n]:=\{0,1,\dots,n-1\}$ (Lean \texttt{Finset.range n}). Where $[m]$ is the index set of a configuration of $m$
points, or the domain or codomain of a cap map (such as $f\colon[k]\to[n]$ in Subsection~\ref{mb:ssec:bridge}), it is
the same set viewed as Lean's \texttt{Fin m}, as in Definition~\ref{mm:def:config}.
``$S$ contains $k$ points in convex position'' means
$\mathsf{HasConvexNGon}(k,S)$ of Definition~\ref{fd:def:ngon}: some $k$-element subset of $S$ is in convex position
(Definition~\ref{fd:def:cp}); no emptiness is required. We use without comment the identities
$\orient(a,c,b)=-\orient(a,b,c)$, $\orient(b,c,a)=\orient(a,b,c)$, $\orient(b,a,c)=-\orient(a,b,c)$,
$\orient(a,b,a)=\orient(a,b,b)=0$ (Lemma~\ref{fd:lem:orient}; the last two are immediate from the formula of
Definition~\ref{fd:def:orient}), and the fact that $\orient(a,b,w)$ is an affine function
of $w$ for fixed $a,b$ (it equals $(x(b)-x(a))(y(w)-y(a))-(y(b)-y(a))(x(w)-x(a))$).

%======================================================================
\subsection{Baek's configurations and the two cited theorems}\label{mb:ssec:baek}

\begin{definition}[configurations, caps, cups and gons;\\
  \leanref{Config}{ErdosSzekeres/Vendor/ETV/Config/Defs.lean:20},\\
  \leanref{Config.Cap3}{ErdosSzekeres/Vendor/ETV/Config/Defs.lean:30},\\
  \leanref{Config.Cap}{ErdosSzekeres/Vendor/ETV/Config/Defs.lean:38},\\
  \leanref{Config.Cup}{ErdosSzekeres/Vendor/ETV/Config/Defs.lean:41},\\
  \leanref{Config.Gon}{ErdosSzekeres/Vendor/ETV/Config/Defs.lean:45},\\
  \leanref{Config.NGon}{ErdosSzekeres/Vendor/ETV/Config/Defs.lean:57},\\
  \leanref{Config.HasNCap}{ErdosSzekeres/Vendor/ETV/Config/Defs.lean:60},\\
  \leanref{Config.HasNCup}{ErdosSzekeres/Vendor/ETV/Config/Defs.lean:63},\\
  \leanref{Config.HasNGon}{ErdosSzekeres/Vendor/ETV/Config/Defs.lean:66}]\label{mb:def:config}
Let $(\alpha,<)$ be a linearly ordered set. A \emph{configuration} $C$ on $\alpha$ is a ternary relation
$\mathrm{Cup3}_C$ on $\alpha$ (in Lean a structure with the field \lean{Cup3} and a decidability witness, which plays no
mathematical role). Put $\mathrm{Cap3}_C(a,b,c):\Leftrightarrow\neg\,\mathrm{Cup3}_C(a,b,c)$ (\lean{Config.Cap3}).
Let $l=(l_0,\dots,l_{r-1})$ be a finite list of elements of $\alpha$, of length $\card l=r$.
\begin{enumerate}[label=(\alph*)]
\item $l$ is a \emph{$C$-cap} if $l_0<l_1<\dots<l_{r-1}$ and $\mathrm{Cap3}_C(l_i,l_{i+1},l_{i+2})$ for every $i$ with
  $i+2<r$; $l$ is a \emph{$C$-cup} if $l_0<l_1<\dots<l_{r-1}$ and $\mathrm{Cup3}_C(l_i,l_{i+1},l_{i+2})$ for every $i$
  with $i+2<r$. Only consecutive triples are constrained; every strictly increasing list of length at most $2$ is both a
  $C$-cap and a $C$-cup.
\item A pair $(l,l')$ of lists is a \emph{$C$-gon} if $\card l\ge2$, $l$ is a $C$-cap, $\card{l'}\ge2$, $l'$ is a
  $C$-cup, and $l$ and $l'$ have the same first entry and the same last entry. It is an \emph{$n$-gon of $C$} if moreover
  $\card l+\card{l'}=n+2$.
\item For a finite set $S\subseteq\alpha$, $\mathsf{HasNCap}_C(n,S)$ means that some $C$-cap of length $n$ has all its
  entries in $S$; $\mathsf{HasNCup}_C(n,S)$ means the same for $C$-cups; and $\mathsf{HasNGon}_C(n,S)$ means that some
  $n$-gon $(l,l')$ of $C$ has all entries of $l$ and of $l'$ in $S$.
\end{enumerate}
\end{definition}

\begin{remark}\label{mb:rem:config}
In Lean, ``strictly increasing'' is \lean{List.IsChain} for the relation $<$ (consecutive entries increase, hence all
entries are in increasing order), the condition on consecutive triples is \texttt{List.Chain3\textquotesingle}
(\dirn{ErdosSzekeres/Vendor/ETV/Lib/List/Defs.lean:46}), ``same first entry'' is the equality
\texttt{l.head? = l\textquotesingle.head?} of the optional first entries and ``same last entry'' is \texttt{l.getLast? = l\textquotesingle.getLast?};
since both lists are nonempty these are the stated equalities. ``All entries in $S$'' is
\lean{List.In} (\dirn{ErdosSzekeres/Vendor/ETV/Lib/List/Defs.lean:22}). An $n$-gon $(l,l')$ consists of $n$ entries if
the two lists share only their end entries; this is the reason for the normalisation $\card l+\card{l'}=n+2$. Nothing
else is required of a gon (Baek calls such gons \emph{weak}): for an arbitrary relation $\mathrm{Cup3}_C$ there is no
geometric content. The geometric meaning enters only in Subsection~\ref{mb:ssec:bridge}.
\end{remark}

\begin{lemma}[monotonicity, and cups as gons;\\
  \leanref{hasNCap_supset}{ErdosSzekeres/Vendor/ETV/Config/Lemmas.lean:293},\\
  \leanref{hasNGon_supset}{ErdosSzekeres/Vendor/ETV/Config/Lemmas.lean:305},\\
  \leanref{ncup_is_ngon}{ErdosSzekeres/Vendor/ETV/Config/Lemmas.lean:275}]\label{mb:lem:config-mono}
Let $C$ be a configuration on a linearly ordered set $\alpha$ and $n\in\N$.
\begin{enumerate}[label=(\alph*)]
\item If $S_1\subseteq S_2\subseteq\alpha$ are finite, then $\mathsf{HasNCap}_C(n,S_1)$ implies
  $\mathsf{HasNCap}_C(n,S_2)$, and $\mathsf{HasNGon}_C(n,S_1)$ implies $\mathsf{HasNGon}_C(n,S_2)$.
\item If $S\subseteq\alpha$ is finite, $2\le n$ and $\mathsf{HasNCup}_C(n,S)$, then $\mathsf{HasNGon}_C(n,S)$.
\end{enumerate}
\end{lemma}

\begin{proof}
(a) The same lists are witnesses, since their entries lie in $S_1\subseteq S_2$.

(b) Let $l=(l_0,\dots,l_{n-1})$ be a $C$-cup with all entries in $S$. As $n\ge2$ and $l$ is strictly increasing,
$l_0<l_{n-1}$. The list $(l_0,l_{n-1})$ is strictly increasing and has no consecutive triple, so it is a $C$-cap of
length $2$; its entries lie in $S$; it has the same first entry $l_0$ and the same last entry $l_{n-1}$ as $l$; and
$2+n=n+2$. Hence $((l_0,l_{n-1}),l)$ is an $n$-gon of $C$ with all entries in $S$.
\end{proof}

\begin{theorem}[Baek's theorem, cited;\\
  \leanref{main}{ErdosSzekeres/Vendor/ETV/Main/Main.lean:87}]\label{mb:thm:baek}
Let $(\gamma,<)$ be a linearly ordered set, $n\in\N$, $C$ a configuration on $\gamma$, and $S\subseteq\gamma$ a finite
set with $\binom{n+2}{2}+2\le\card S$. Then $\mathsf{HasNCap}_C(4,S)$ or $\mathsf{HasNGon}_C(n+3,S)$.
\end{theorem}

\begin{remark}[source, status and index shift]\label{mb:rem:baek}
(1) \emph{Source.} J.~Baek, \emph{On the Erd\H{o}s--Tuza--Valtr conjecture}, Eur.\ J.\ Combin.\ 124 (2025) 104085;
arXiv:2206.04260. Its geometric main theorem (Theorem~1.6 in the arXiv numbering, quoted in
\dirn{MIRRORETV/PROOFS.md}, proof of Corollary~7.2) says: for every $N\ge3$, every set of $\binom{N-1}{2}+2$ points in
the plane in general position with pairwise distinct $x$-coordinates contains a $4$-cap or $N$ points in convex
position. Theorem~\ref{mb:thm:baek} is a combinatorial statement, for an arbitrary ternary relation and weak gons. It
is the theorem \lean{main} of the vendored file \dirn{ErdosSzekeres/Vendor/ETV/Main/Main.lean}, which is adapted from
Baek's own Lean code (according to its header, the changes concern only imports and compatibility with Lean/Mathlib
4.24.0); the vendored file is compiled with the repository, so the axioms evidence above covers it.
\dirn{MIRRORETV/README.md}, \S8, states that Baek's Theorem~2.7 holds for all colourings with weak gons; we have not
compared the wording of that theorem with Theorem~\ref{mb:thm:baek}. In this document Theorem~\ref{mb:thm:baek} is
cited, not reproved.

(2) \emph{Index shift.} The Lean theorem \lean{main} at $n$ is Baek's statement for gons of size $N=n+3$: indeed
$\binom{n+2}{2}+2=\binom{N-1}{2}+2$. We use $n=4$ ($N=7$, threshold $\binom62+2=17$), $n=5$ ($N=8$, threshold
$\binom72+2=23$) and $n=6$ ($N=9$, threshold $\binom82+2=30$). Equivalently: if a finite set $S$ contains no $4$-cap and
no $N$-gon of $C$, then $\card S\le\binom{N-1}{2}+1$, which is $16$, $22$, $29$ for $N=7,8,9$.

(3) \emph{Geometric meaning.} For the configuration $C_c$ of a point configuration $c$ of $n$ points
(Definition~\ref{mb:def:color}), $C_c$-caps and $C_c$-cups with entries in $[n]$ are geometric caps and cups
(Lemma~\ref{mb:lem:lists}), and for $N\in\{7,8,9\}$ every $N$-gon of $C_c$ with entries in $[n]$ yields $N$ of the
points of $c$ in convex position (Proposition~\ref{mb:prop:gon}; its proof does not use the value of $N$).
\end{remark}

\begin{theorem}[cups--caps theorem, cited;\\
  \leanref{Config.cap_cup}{ErdosSzekeres/Vendor/ETV/Main/CapCup.lean:46}]\label{mb:thm:capcup}
Let $C$ be a configuration on a linearly ordered set $\alpha$, let $a,b\in\N$, and let $S\subseteq\alpha$ be finite with
$\binom{a+b}{a}<\card S$. Then $\mathsf{HasNCap}_C(a+2,S)$ or $\mathsf{HasNCup}_C(b+2,S)$.
\end{theorem}

\begin{remark}\label{mb:rem:capcup}
This is the cups--caps theorem of P.~Erd\H{o}s and G.~Szekeres, \emph{A combinatorial problem in geometry}, Compositio
Math.\ 2 (1935) 463--470. Its classical proof (induction on $a+b$; split $S$ into the first entries of $(a+1)$-caps,
that is, of caps one entry shorter than the target $(a+2)$-cap, and the rest) uses only consecutive triples, hence holds
for every configuration; the vendored file contains a Lean proof of this form, which we cite. The proof of
Theorem~\ref{mb:thm:nine}(a) below repeats its induction step once (target $5$-cap, split by the first entries of
$4$-caps).
\end{remark}

\begin{lemma}[Baek's theorem for $8$-gons;\\
  \leanref{ErdosSzekeres.Omission.Baek.four_cap_or_eight_gon}{ErdosSzekeres/Omission/Baek.lean:12},\\
  \leanref{ErdosSzekeres.Omission.Baek.capacity22}{ErdosSzekeres/Omission/Baek.lean:16}]\label{mb:lem:baek8}
Let $C$ be a configuration on a linearly ordered set $\alpha$ and $S\subseteq\alpha$ finite.
\begin{enumerate}[label=(\alph*)]
\item If $23\le\card S$, then $\mathsf{HasNCap}_C(4,S)$ or $\mathsf{HasNGon}_C(8,S)$.
\item If $\neg\mathsf{HasNCap}_C(4,S)$ and $\neg\mathsf{HasNGon}_C(8,S)$, then $\card S\le22$.
\end{enumerate}
\end{lemma}

\begin{proof}
(a) is Theorem~\ref{mb:thm:baek} with $n=5$, because $\binom{7}{2}+2=23$. (b) If $\card S\ge23$, then (a) contradicts the
two hypotheses; hence $\card S\le22$.
\end{proof}

%======================================================================
\subsection{From configurations to points of the plane}\label{mb:ssec:bridge}

We use the $x$-sorted point configurations of Definition~\ref{mm:def:config}. A \emph{configuration of $n$ points}
(Lean \lean{ErdosSzekeres.Omission.Configuration}, \dirn{ErdosSzekeres/Omission/Realization.lean:9}) is a map
$c\colon[n]\to\R^2$, $i\mapsto c_i$, with $x(c_i)<x(c_j)$ for all $i<j$ and $\orient(c_i,c_j,c_k)\ne0$ for all
$i<j<k$. For $k\in\N$, $c.\mathsf{HasCap}(k)$ (\lean{Configuration.HasCap}, \dirn{ErdosSzekeres/Omission/Realization.lean:64})
means that there is a strictly increasing $f\colon[k]\to[n]$ with $\orient(c_{f(i)},c_{f(j)},c_{f(l)})<0$ for \emph{all}
$i<j<l$ in $[k]$. We write $\mathcal P(c):=\{c_0,\dots,c_{n-1}\}$ for the set of points of $c$ (Lean
\texttt{Finset.univ.image c.point}; it is the point set $\mathrm{pts}(c)$ of Definition~\ref{mm:def:config}(a)); since
$x$ is strictly increasing along $c$, $\card{\mathcal P(c)}=n$. (A
configuration of $n$ points is not a configuration in the sense of Definition~\ref{mb:def:config}; the link between the
two notions is the configuration $C_c$ of Definition~\ref{mb:def:color}.)

\begin{definition}[the configuration $C_c$ of a point configuration;\\
  \leanref{ErdosSzekeres.Omission.Configuration.cup}{ErdosSzekeres/Omission/Realization.lean:33},\\
  \leanref{ErdosSzekeres.Omission.ActualRanks.colorConfig}{ErdosSzekeres/Omission/ActualRanks.lean:28},\\
  \leanref{ErdosSzekeres.Omission.ListGeometry.pointAt}{ErdosSzekeres/Omission/ListGeometry.lean:13}]\label{mb:def:color}\raggedright
Let $c$ be a configuration of $n$ points. For $i,j,k\in\N$ put $c.\mathrm{cup}(i,j,k)=\mathrm{true}$ if $i,j,k<n$ and
$\orient(c_i,c_j,c_k)>0$, and $c.\mathrm{cup}(i,j,k)=\mathrm{false}$ otherwise. Let $C_c$ be the configuration on
$(\N,<)$ with $\mathrm{Cup3}_{C_c}(i,j,k)$ if and only if $c.\mathrm{cup}(i,j,k)=\mathrm{true}$ (Lean
\texttt{colorConfig c.cup true}). Finally $\mathrm{pointAt}(c,i):=c_i$ for $i<n$ and $\mathrm{pointAt}(c,i):=(0,0)$ for
$i\ge n$; for indices $i<n$, which are the only ones occurring below, we simply write $c_i$.
\end{definition}

\begin{lemma}[signs of the colouring;\\
  \leanref{ErdosSzekeres.Omission.Configuration.cup_true_iff}{ErdosSzekeres/Omission/Realization.lean:39},\\
  \leanref{ErdosSzekeres.Omission.ListGeometry.cup_sign}{ErdosSzekeres/Omission/ListGeometry.lean:19},\\
  \leanref{ErdosSzekeres.Omission.ListGeometry.cap_sign}{ErdosSzekeres/Omission/ListGeometry.lean:24}]\label{mb:lem:signs}\raggedright
Let $c$ be a configuration of $n$ points and $a,b,d\in\N$ with $a,b,d<n$. Then $c.\mathrm{cup}(a,b,d)=\mathrm{true}$,
that is $\mathrm{Cup3}_{C_c}(a,b,d)$, holds if and only if $\orient(c_a,c_b,c_d)>0$. If moreover $a<b<d$, then
$c.\mathrm{cup}(a,b,d)\ne\mathrm{true}$, that is $\mathrm{Cap3}_{C_c}(a,b,d)$, holds if and only if
$\orient(c_a,c_b,c_d)<0$.
\end{lemma}

\begin{proof}
The first claim is Definition~\ref{mb:def:color}, since $a,b,d<n$. If $a<b<d$, then $\orient(c_a,c_b,c_d)\ne0$ because
$c$ is a configuration; so $\mathrm{Cap3}_{C_c}(a,b,d)$, which is $\neg(\orient(c_a,c_b,c_d)>0)$, is equivalent to
$\orient(c_a,c_b,c_d)<0$. (In Lean, \lean{cup_true_iff} states the first claim with indices $i,j,k$ in \texttt{Fin n}
and the points \texttt{c.point i}; \lean{cup_sign} and \lean{cap_sign} state the two claims for natural numbers
$a,b,d$, with the bounds $a,b,d<n$ as hypotheses and the points $\mathrm{pointAt}(c,a)$, which equal $c_a$ for $a<n$.
The companion \lean{Configuration.cup_false_iff} (\dirn{ErdosSzekeres/Omission/Realization.lean:44}), the second claim
with indices in \texttt{Fin n} and the equation $c.\mathrm{cup}(i,j,k)=\mathrm{false}$, is not used on the proof paths
of this section and is not covered here.)
\end{proof}

\begin{lemma}[four $x$-sorted points;\\
  \leanref{ErdosSzekeres.Omission.PlanarChains.four_point_identity_left}{ErdosSzekeres/Omission/PlanarChains.lean:15},\\
  \leanref{ErdosSzekeres.Omission.PlanarChains.four_point_identity_right}{ErdosSzekeres/Omission/PlanarChains.lean:19},\\
  \leanref{ErdosSzekeres.Omission.PlanarChains.positive_four}{ErdosSzekeres/Omission/PlanarChains.lean:23},\\
  \leanref{ErdosSzekeres.Omission.PlanarChains.negative_four}{ErdosSzekeres/Omission/PlanarChains.lean:35}]\label{mb:lem:four}\raggedright
For all $a,b,c,d\in\R^2$,
\begin{align}
(x(c)-x(b))\,\orient(a,b,d)&=(x(d)-x(b))\,\orient(a,b,c)+(x(b)-x(a))\,\orient(b,c,d),\label{mb:eq:four-left}\\
(x(c)-x(b))\,\orient(a,c,d)&=(x(d)-x(c))\,\orient(a,b,c)+(x(c)-x(a))\,\orient(b,c,d).\label{mb:eq:four-right}
\end{align}
Let moreover $x(a)<x(b)<x(c)<x(d)$. If $\orient(a,b,c)>0$ and $\orient(b,c,d)>0$, then $\orient(a,b,d)>0$ and
$\orient(a,c,d)>0$. If $\orient(a,b,c)<0$ and $\orient(b,c,d)<0$, then $\orient(a,b,d)<0$ and $\orient(a,c,d)<0$.
\end{lemma}

\begin{proof}
Both identities are polynomial identities in the eight coordinates; expanding both sides gives the same polynomial (we
checked this symbolically; it is also the Lean proof, \texttt{unfold orient; ring}). Under $x(a)<x(b)<x(c)<x(d)$ the
coefficients $x(d)-x(b)$, $x(b)-x(a)$, $x(d)-x(c)$, $x(c)-x(a)$ on the right are positive, so each right-hand side is
positive (negative) when $\orient(a,b,c)$ and $\orient(b,c,d)$ are both positive (negative). Dividing by
$x(c)-x(b)>0$ gives the claims.
\end{proof}

\begin{lemma}[closure of chain conditions;\\
  \leanref{ErdosSzekeres.Omission.PlanarChains.chain3_all}{ErdosSzekeres/Omission/PlanarChains.lean:58}]\label{mb:lem:chain}\raggedright
Let $(\alpha,<)$ be a linearly ordered set, $R$ a ternary relation on $\alpha$, and $D\subseteq\alpha$ such that for all
$a,b,c,d\in D$ with $a<b<c<d$,
\begin{equation}\label{mb:eq:trans}
R(a,b,c)\ \text{and}\ R(b,c,d)\ \Longrightarrow\ R(a,b,d)\ \text{and}\ R(a,c,d).
\end{equation}
Let $l=(l_0,\dots,l_{r-1})$ be a list of elements of $D$ with $l_0<l_1<\dots<l_{r-1}$ and
$R(l_i,l_{i+1},l_{i+2})$ for all $i$ with $i+2<r$. Then $R(a,b,c)$ holds for all entries $a,b,c$ of $l$ with $a<b<c$
(Lean: \texttt{OrderedTriples R l}).
\end{lemma}

\begin{proof}
Since $l$ is strictly increasing, the claim says $R(l_i,l_j,l_k)$ for all $0\le i<j<k<r$. We prove by downward induction
on $i$ the statement $Q(i)$: $R(l_i,l_j,l_k)$ for all $j,k$ with $i<j<k<r$. For $i\ge r-2$ there is nothing to prove.
Let $i\le r-3$ and assume $Q(i')$ for all $i'>i$.

\emph{Step A: $R(l_i,l_{i+1},l_k)$ for $i+2\le k<r$.} Induction on $k$. For $k=i+2$ this is a hypothesis on consecutive
triples. For $k>i+2$ we have $R(l_i,l_{i+1},l_{k-1})$ by the inner induction hypothesis and $R(l_{i+1},l_{k-1},l_k)$ by
$Q(i+1)$; the first conclusion of \eqref{mb:eq:trans} for $l_i<l_{i+1}<l_{k-1}<l_k$ gives $R(l_i,l_{i+1},l_k)$.

\emph{Step B: $R(l_i,l_j,l_k)$ for $i+1<j<k<r$.} By Step~A, $R(l_i,l_{i+1},l_j)$; by $Q(i+1)$, $R(l_{i+1},l_j,l_k)$;
the second conclusion of \eqref{mb:eq:trans} for $l_i<l_{i+1}<l_j<l_k$ gives $R(l_i,l_j,l_k)$.

Steps A ($j=i+1$) and B ($j>i+1$) give $Q(i)$. All points used lie in $D$. (The Lean proof is an induction on the list,
with the same two steps.)
\end{proof}

\begin{lemma}[$C_c$-caps and $C_c$-cups are geometric;\\
  \leanref{ErdosSzekeres.Omission.ListGeometry.cup_list_all}{ErdosSzekeres/Omission/ListGeometry.lean:39},\\
  \leanref{ErdosSzekeres.Omission.ListGeometry.cap_list_all}{ErdosSzekeres/Omission/ListGeometry.lean:53},\\
  \leanref{ErdosSzekeres.Omission.ListGeometry.hasCap_of_list}{ErdosSzekeres/Omission/ListGeometry.lean:85}]\label{mb:lem:lists}\raggedright
Let $c$ be a configuration of $n$ points and $l$ a list all of whose entries are $<n$.
\begin{enumerate}[label=(\alph*)]
\item If $l$ is a $C_c$-cup, then $\orient(c_a,c_b,c_d)>0$ for all entries $a<b<d$ of $l$.
\item If $l$ is a $C_c$-cap, then $\orient(c_a,c_b,c_d)<0$ for all entries $a<b<d$ of $l$.
\item For every $k\in\N$, $\mathsf{HasNCap}_{C_c}(k,[n])$ implies $c.\mathsf{HasCap}(k)$.
\end{enumerate}
\end{lemma}

\begin{proof}
(b) Apply Lemma~\ref{mb:lem:chain} with $\alpha=\N$, $D=[n]$ and $R(a,b,d):\Leftrightarrow\mathrm{Cap3}_{C_c}(a,b,d)$.
Hypothesis \eqref{mb:eq:trans}: let $a<b<d<e$ be in $[n]$ with $R(a,b,d)$ and $R(b,d,e)$. By Lemma~\ref{mb:lem:signs},
$\orient(c_a,c_b,c_d)<0$ and $\orient(c_b,c_d,c_e)<0$; since $c$ is a configuration,
$x(c_a)<x(c_b)<x(c_d)<x(c_e)$; Lemma~\ref{mb:lem:four} gives $\orient(c_a,c_b,c_e)<0$ and $\orient(c_a,c_d,c_e)<0$, that
is $R(a,b,e)$ and $R(a,d,e)$ by Lemma~\ref{mb:lem:signs}. The list $l$ is strictly increasing with entries in $[n]$, and
its consecutive triples satisfy $\mathrm{Cap3}_{C_c}=R$ because $l$ is a $C_c$-cap. Lemma~\ref{mb:lem:chain} gives
$R(a,b,d)$ for all entries $a<b<d$ of $l$, i.e.\ $\orient(c_a,c_b,c_d)<0$ by Lemma~\ref{mb:lem:signs}.

(a) The same argument with $R=\mathrm{Cup3}_{C_c}$ and the positive half of Lemma~\ref{mb:lem:four}.

(c) Let $l=(l_0,\dots,l_{k-1})$ be a $C_c$-cap of length $k$ with entries in $[n]$, and put $f(i):=l_i$. Then
$f\colon[k]\to[n]$ is strictly increasing because $l$ is, and for $i<j<m$ in $[k]$, (b) applied to the entries
$l_i<l_j<l_m$ gives $\orient(c_{f(i)},c_{f(j)},c_{f(m)})<0$. Hence $c.\mathsf{HasCap}(k)$.
\end{proof}

The next lemma is the geometric heart of the bridge: a geometric cup and a geometric cap with the same two end points
form a convex polygon. (This is the direction ``$\Leftarrow$'' of \dirn{MIRRORETV/PROOFS.md}, Lemma~2.3.)

\begin{lemma}[cup and cap with common ends;\\
  \leanref{ErdosSzekeres.Omission.BoundaryGeometry.boundary_polygon}{ErdosSzekeres/Omission/BoundaryGeometry.lean:403},\\
  \leanref{ErdosSzekeres.Omission.BoundaryGeometry.Boundary.card}{ErdosSzekeres/Omission/BoundaryGeometry.lean:359},\\
  \leanref{ErdosSzekeres.Omission.BoundaryGeometry.boundary_convex_position}{ErdosSzekeres/Omission/BoundaryGeometry.lean:317}]\label{mb:lem:boundary}\raggedright
Let $s,t\in\N$. Let $u_0,\dots,u_{s+1}\in\R^2$ satisfy $x(u_0)<x(u_1)<\dots<x(u_{s+1})$ and $\orient(u_i,u_j,u_m)>0$ for
all $i<j<m$, and let $v_0,\dots,v_{t+1}\in\R^2$ satisfy $x(v_0)<\dots<x(v_{t+1})$ and $\orient(v_i,v_j,v_m)<0$ for all
$i<j<m$. Assume $u_0=v_0$ and $u_{s+1}=v_{t+1}$, and let $B:=\{u_0,\dots,u_{s+1}\}\cup\{v_0,\dots,v_{t+1}\}$. Then
\begin{enumerate}[label=(\alph*)]
\item $\card B=s+t+2$;
\item $p\notin\conv(B\setminus\{p\})$ for every $p\in B$; equivalently (Lemma~\ref{fd:lem:cp-hull}), $B$ is in convex
  position;
\item for every finite $S\subset\R^2$ with $B\subseteq S$, $S$ contains $s+t+2$ points in convex position.
\end{enumerate}
\end{lemma}

\begin{remark}[the Lean encoding]\label{mb:rem:boundary}
Lean stores this data as a \texttt{Boundary s t} (\dirn{ErdosSzekeres/Omission/BoundaryGeometry.lean:120}): the lower
arc \lean{lower} is $(u_i)$, a \texttt{CupArc s} (strictly increasing $x$, all ordered triples positive, line 56), and the
upper arc is stored through its mirror image \lean{upperReflected}${}=(\rho(v_j))$, a \texttt{CupArc t}, where
$\rho(p)=(p_0,-p_1)$ is \lean{BoundaryGeometry.reflect} (line 69); the fields \lean{first_eq} and \lean{last_eq} say
$u_0=\rho(\rho(v_0))=v_0$ and $u_{s+1}=v_{t+1}$, and \lean{Boundary.points} is $B$. Since $\rho$ preserves $x$ and
$\orient(\rho a,\rho b,\rho d)=-\orient(a,b,d)$, $(\rho(v_j))$ is a \lean{CupArc} if and only if $(v_j)$ has strictly
increasing $x$ and all ordered triples negative. So Lemma~\ref{mb:lem:boundary} is \lean{boundary_polygon} (part (c)),
\lean{Boundary.card} (part (a)) and \lean{boundary_convex_position} (part (b)) after the change of variables
$v_j=\rho(w_j)$, where $(w_j)$ is \lean{upperReflected}. This change of variables is a bijection between the data
$(u_i)$, $(v_j)$ of the lemma and the values of type \texttt{Boundary s t}, and \lean{Boundary.card} counts the finite
set \lean{Boundary.finite}, which is $B$; so the lemma is the Lean statement in this encoding, not a variant of it.
\end{remark}

\begin{proof}[Proof of Lemma~\ref{mb:lem:boundary}]
Put $L:=u_0=v_0$, $R:=u_{s+1}=v_{t+1}$, $U:=\{u_0,\dots,u_{s+1}\}$ and $V:=\{v_0,\dots,v_{t+1}\}$. The $u_i$ are pairwise
distinct and the $v_j$ are pairwise distinct (their $x$-coordinates are), so $\card U=s+2$, $\card V=t+2$, and
$x(L)<x(R)$, in particular $L\ne R$. The proof follows the Lean proof.

\emph{Step 1 (the chord $LR$).} For $0<i<s+1$, $\orient(L,R,u_i)=-\orient(u_0,u_i,u_{s+1})<0$; for $i\in\{0,s+1\}$,
$\orient(L,R,u_i)=0$. For $0<j<t+1$, $\orient(L,R,v_j)=-\orient(v_0,v_j,v_{t+1})>0$; for $j\in\{0,t+1\}$,
$\orient(L,R,v_j)=0$. (Lean: \lean{chord_negative}, \lean{chord_nonpositive}.)

\emph{Step 2 (only the ends are shared; part (a)).} If $u_i=v_j$, then $\orient(L,R,u_i)=\orient(L,R,v_j)\ge0$ by
Step~1, so $i\in\{0,s+1\}$ by Step~1 again, i.e.\ $u_i\in\{L,R\}$. Hence $U\cap V=\{L,R\}$ (both $L$ and $R$ lie in
$U\cap V$), and $\card B=\card U+\card V-\card{U\cap V}=(s+2)+(t+2)-2=s+t+2$. (Lean: \lean{shared_only_endpoints},
\lean{Boundary.card}.)

\emph{Step 3 (the lower edges support $B$).} We claim: for $0\le i\le s$ and $w\in B$,
\begin{equation}\label{mb:eq:edge}
\orient(u_i,u_{i+1},w)\ge0,\quad\text{with strict inequality if } w\notin\{u_i,u_{i+1}\}.
\end{equation}
If $w=u_j$: for $j\in\{i,i+1\}$ the orientation is $0$; for $j<i$ it equals $\orient(u_j,u_i,u_{i+1})>0$ (cyclic
invariance); for $j>i+1$ it is $\orient(u_i,u_{i+1},u_j)>0$. If $w=v_j$ with $j\in\{0,t+1\}$, then $w\in\{L,R\}\subseteq
U$ and the previous case applies. If $w=v_j$ with $0<j<t+1$, we use the identity (valid for all points, checked
symbolically; Lean \lean{chord_identity}, line 64)
\begin{equation}\label{mb:eq:chord}
\begin{split}
(x(R)-x(L))\,\orient(a,b,w)&=(x(R)-x(w))\,\orient(a,b,L)+(x(w)-x(L))\,\orient(a,b,R)\\
&\quad+(x(b)-x(a))\,\orient(L,R,w)
\end{split}
\end{equation}
with $a=u_i$, $b=u_{i+1}$. Here $x(L)<x(w)<x(R)$ because the $x(v_j)$ increase, $\orient(u_i,u_{i+1},L)\ge0$ and
$\orient(u_i,u_{i+1},R)\ge0$ by the case $w\in U$ (as $L=u_0$, $R=u_{s+1}$), $x(u_{i+1})-x(u_i)>0$, and
$\orient(L,R,w)>0$ by Step~1. So the right-hand side of \eqref{mb:eq:chord} is positive, and dividing by $x(R)-x(L)>0$
gives $\orient(u_i,u_{i+1},w)>0$. (In this last case $w\notin U$ by Step~2.) (Lean: \lean{edge_nonnegative},
\lean{edge_positive}, \lean{cross_edge_positive}, \lean{lower_edge_support}.)

\emph{Step 4 (the upper edges support $B$).} We claim: for $0\le j\le t$ and $w\in B$,
\begin{equation}\label{mb:eq:upper-edge}
\orient(v_{j+1},v_j,w)\ge0,\quad\text{with strict inequality if } w\notin\{v_j,v_{j+1}\}.
\end{equation}
Let $\rho(p)=(p_0,-p_1)$; then $x(\rho p)=x(p)$ and $\orient(\rho a,\rho b,\rho d)=-\orient(a,b,d)$
(Lemma~\ref{fd:lem:affine-orient}, $\det\rho=-1$). The points $\rho(v_0),\dots,\rho(v_{t+1})$ therefore satisfy the
hypotheses on the $u$'s (with $t$ in place of $s$), the points $\rho(u_0),\dots,\rho(u_{s+1})$ satisfy the hypotheses on
the $v$'s, and $\rho(v_0)=\rho(u_0)$, $\rho(v_{t+1})=\rho(u_{s+1})$; the set of all these points is $\rho(B)$. Step~3 for
this reflected data, applied to the edge $(\rho v_j,\rho v_{j+1})$ and the point $\rho w\in\rho(B)$, gives
$\orient(\rho v_j,\rho v_{j+1},\rho w)\ge0$, strictly if $\rho w\notin\{\rho v_j,\rho v_{j+1}\}$, i.e.\ (as $\rho$ is
injective) if $w\notin\{v_j,v_{j+1}\}$. Finally
$\orient(\rho v_j,\rho v_{j+1},\rho w)=-\orient(v_j,v_{j+1},w)=\orient(v_{j+1},v_j,w)$. (Lean: \lean{Boundary.flip},
\lean{upper_edge_support}.)

\emph{Step 5 (exposing functions; part (b)).} We use the following elementary fact (the easy direction of
Lemma~\ref{fd:lem:separation}; Lean \lean{exposed_not_in_hull}, line 30): \emph{if $g\colon\R^2\to\R$ is affine, $g(p)=0$
and $g(w)>0$ for every $w$ in a finite set $F$, then $p\notin\conv F$.} Indeed every point of $\conv F$ is
$\sum_k\lambda_kw_k$ with $w_k\in F$, $\lambda_k\ge0$, $\sum_k\lambda_k=1$, and $g(\sum_k\lambda_kw_k)=\sum_k\lambda_k
g(w_k)>0$. So it suffices to find, for every $p\in B$, an affine $g$ with $g(p)=0$ and $g>0$ on $B\setminus\{p\}$. Every $p\in B$ is $L$, $R$, some $u_i$ with $0<i<s+1$, or some $v_j$ with $0<j<t+1$.
\begin{itemize}
\item $p=L$: $g(w)=x(w)-x(L)$. Every $w\in B\setminus\{L\}$ is $u_i$ or $v_j$ with $i\ge1$ or $j\ge1$, so $x(w)>x(L)$.
  (Lean \lean{boundary_first_exposed}.)
\item $p=R$: $g(w)=x(R)-x(w)$, symmetrically. (Lean \lean{boundary_last_exposed}.)
\item $p=u_i$, $0<i<s+1$: $g(w)=\orient(u_{i-1},u_i,w)+\orient(u_i,u_{i+1},w)$, which is affine with $g(u_i)=0$. Let
  $w\in B\setminus\{u_i\}$. If $w=u_{i-1}$, the first term is $0$ and the second is $>0$ by \eqref{mb:eq:edge}, since
  $u_{i-1}\notin\{u_i,u_{i+1}\}$. If $w\ne u_{i-1}$, the first term is $>0$ and the second $\ge0$ by \eqref{mb:eq:edge}.
  (Lean \lean{two_edges_expose}, \lean{lower_interior_exposed}.)
\item $p=v_j$, $0<j<t+1$: $g(w)=\orient(v_{j+1},v_j,w)+\orient(v_j,v_{j-1},w)$, affine with $g(v_j)=0$. For
  $w\in B\setminus\{v_j\}$: if $w=v_{j+1}$ the first term is $0$ and the second is $>0$ by \eqref{mb:eq:upper-edge}
  (since $v_{j+1}\notin\{v_{j-1},v_j\}$); otherwise the first term is $>0$ and the second $\ge0$. (Lean
  \lean{upper_interior_exposed}.)
\end{itemize}
So $p\notin\conv(B\setminus\{p\})$ for every $p\in B$, which is (b); by Lemma~\ref{fd:lem:cp-hull}, $B$ is in convex
position.

\emph{Step 6 (part (c)).} $B\subseteq S$, $\card B=s+t+2$ and $B$ is in convex position, so $B$ witnesses
$\mathsf{HasConvexNGon}(s+t+2,S)$.
\end{proof}

\begin{proposition}[gons of $C_c$ are convex polygons;\\
  \leanref{MirrorOdd.hasNGon_heptagon}{MIRRORETV/lean/MirrorMain.lean:1512},\\
  \leanref{ErdosSzekeres.Omission.BaekGeometry.hasNGon_octagon}{ErdosSzekeres/Omission/BaekGeometry.lean:94},\\
  \leanref{C12Check.hasNGon_nonagon}{MIRRORETV/lean/MirrorNine_append.lean:99}]\label{mb:prop:gon}\raggedright
Let $c$ be a configuration of $n$ points and $m\in\{7,8,9\}$. If $\mathsf{HasNGon}_{C_c}(m,[n])$, then $\mathcal P(c)$
contains $m$ points in convex position. (For $m=7,8,9$ this is \lean{hasNGon_heptagon}, \lean{hasNGon_octagon},
\lean{hasNGon_nonagon} respectively.)
\end{proposition}

\begin{proof}
The three Lean proofs are identical up to the number $m$, which plays no role. Let $(l,l')$ be an $m$-gon of $C_c$ with
all entries in $[n]$: $l$ is a $C_c$-cap, $l'$ a $C_c$-cup, $\card l,\card{l'}\ge2$, the first entries agree, the last
entries agree, and $\card l+\card{l'}=m+2$. Put $s:=\card{l'}-2$ and $t:=\card l-2$ (Lean calls these numbers \lean{l} and
\lean{k}), so that $s+t+2=m$, and put $u_i:=c_{l'_i}$ ($0\le i\le s+1$) and $v_j:=c_{l_j}$ ($0\le j\le t+1$). The lists
are strictly increasing and $c$ is $x$-sorted, so $x(u_0)<\dots<x(u_{s+1})$ and $x(v_0)<\dots<x(v_{t+1})$. By
Lemma~\ref{mb:lem:lists}(a), $\orient(u_i,u_j,u_k)>0$ for $i<j<k$; by Lemma~\ref{mb:lem:lists}(b),
$\orient(v_i,v_j,v_k)<0$ for $i<j<k$. Equal first entries give $u_0=v_0$ and equal last entries give $u_{s+1}=v_{t+1}$.
All $u_i$ and $v_j$ lie in $\mathcal P(c)$. Lemma~\ref{mb:lem:boundary}(c) with $S=\mathcal P(c)$ gives $s+t+2=m$ points
of $\mathcal P(c)$ in convex position. (In Lean, \lean{boundaryOfGon}, \dirn{ErdosSzekeres/Omission/BaekGeometry.lean:62},
builds the \lean{Boundary} from the two lists, and \lean{boundaryOfGon_inside}, line 77, says that its points lie in
$\mathcal P(c)$.)
\end{proof}

%======================================================================
\subsection{The capacity bounds 16 and 22}\label{mb:ssec:cap}

\begin{proposition}[Baek at $N=7$ for point configurations;\\
  \leanref{MirrorOdd.capacity16}{MIRRORETV/lean/MirrorMain.lean:1529}]\label{mb:prop:cap16}
Let $c$ be a configuration of $n$ points with $\neg\,c.\mathsf{HasCap}(4)$ such that $\mathcal P(c)$ does not contain $7$
points in convex position. Then $n\le16$.
\end{proposition}

\begin{proof}
Suppose $n\ge17$. Then $\binom{4+2}{2}+2=17\le n=\card{[n]}$ (the value $\binom62=15$ is evaluated by \lean{decide}), so
Theorem~\ref{mb:thm:baek} with the value $4$ for its parameter (called $n$ there), $C=C_c$ and $S=[n]$ gives
$\mathsf{HasNCap}_{C_c}(4,[n])$ or
$\mathsf{HasNGon}_{C_c}(7,[n])$. In the first case $c.\mathsf{HasCap}(4)$ by Lemma~\ref{mb:lem:lists}(c); in the second,
$\mathcal P(c)$ contains $7$ points in convex position by Proposition~\ref{mb:prop:gon}. Both contradict the hypotheses.
\end{proof}

\begin{lemma}[Baek at $N=8$ on index sets;\\
  \leanref{ErdosSzekeres.Omission.BaekGeometry.geometric_capacity22}{ErdosSzekeres/Omission/BaekGeometry.lean:111}]\label{mb:lem:gc22}\raggedright
Let $c$ be a configuration of $n$ points and $S\subseteq[n]$ such that $\neg\mathsf{HasNCap}_{C_c}(4,S)$ and
$\mathcal P(c)$ does not contain $8$ points in convex position. Then $\card S\le22$.
\end{lemma}

\begin{proof}
By Lemma~\ref{mb:lem:baek8}(b) it suffices to show $\neg\mathsf{HasNGon}_{C_c}(8,S)$. An $8$-gon of $C_c$ with entries in
$S$ is one with entries in $[n]$ (Lemma~\ref{mb:lem:config-mono}(a)), so it would give $8$ points of $\mathcal P(c)$ in
convex position by Proposition~\ref{mb:prop:gon}.
\end{proof}

\begin{proposition}[Baek at $N=8$ for point configurations;\\
  \leanref{MirrorOdd.capacity22g}{MIRRORETV/lean/MirrorMain.lean:1812}]\label{mb:prop:cap22}
Let $c$ be a configuration of $n$ points with $\neg\,c.\mathsf{HasCap}(4)$ such that $\mathcal P(c)$ does not contain $8$
points in convex position. Then $n\le22$.
\end{proposition}

\begin{proof}
Apply Lemma~\ref{mb:lem:gc22} with $S=[n]$: a $C_c$-cap of length $4$ in $[n]$ would give $c.\mathsf{HasCap}(4)$ by
Lemma~\ref{mb:lem:lists}(c). Hence $n=\card{[n]}\le22$.
\end{proof}

%======================================================================
\subsection{Mirror-symmetric sets: notation and reduction lemmas}\label{mb:ssec:red}

The Lean declarations stated in this subsection, and those of Definition~\ref{mb:def:frame} and
Lemmas~\ref{mb:lem:frame}, \ref{mb:lem:affine} and~\ref{mb:lem:transport} below, are also stated (and, where they are
lemmas, proved) in the section of Theorem~\ref{mm:thm:odd}; we repeat them here, with proofs, so that this section is
self-contained.

\begin{definition}[the reflection in the $x$-axis; upper and axis parts;\\
  \leanref{MirrorOdd.mir}{MIRRORETV/lean/MirrorMain.lean:16},\\
  \leanref{MirrorOdd.Rset}{MIRRORETV/lean/MirrorMain.lean:1598},\\
  \leanref{MirrorOdd.Aset}{MIRRORETV/lean/MirrorMain.lean:1599}]\label{mb:def:mir}
$\operatorname{mir}(p):=(p_0,-p_1)$. For a finite $S\subset\R^2$, $R(S):=\{p\in S:y(p)>0\}$ and
$A(S):=\{p\in S:y(p)=0\}$. We say that $S$ is \emph{$\operatorname{mir}$-symmetric} if $\operatorname{mir}(p)\in S$ for
all $p\in S$ (this is the hypothesis $\forall p\in S,\ \operatorname{mir}(p)\in S$, named \lean{sym}, of the Lean
statements; for finite $S$ it is equivalent to $\operatorname{mir}(S)=S$, because $\operatorname{mir}$ is injective).
\end{definition}

Note that $\operatorname{mir}(\operatorname{mir}(p))=p$ (\lean{MirrorOdd.mir_mir}, line 1543), that
$\operatorname{mir}(a)=a$ if $y(a)=0$, and that $\orient(\operatorname{mir}a,\operatorname{mir}b,\operatorname{mir}d)
=-\orient(a,b,d)$ (Lemma~\ref{fd:lem:affine-orient}).

\begin{lemma}[splitting a mirror-symmetric set;\\
  \leanref{MirrorOdd.card_split}{MIRRORETV/lean/MirrorMain.lean:1601},\\
  \leanref{MirrorOdd.card_A_le}{MIRRORETV/lean/MirrorMain.lean:1636},\\
  \leanref{MirrorOdd.injOn_x}{MIRRORETV/lean/MirrorMain.lean:1647}]\label{mb:lem:split}
Let $S\subset\R^2$ be finite.
\begin{enumerate}[label=(\alph*)]
\item If $S$ is $\operatorname{mir}$-symmetric, then $\card S=2\card{R(S)}+\card{A(S)}$.
\item If $S$ is in general position, then $\card{A(S)}\le2$.
\item Let $S$ be in general position and $\operatorname{mir}$-symmetric, and let $T\subseteq S$ be such that $y(p)\ge0$
  for all $p\in T$ and $p=q$ for all $p,q\in T$ with $y(p)=0=y(q)$. Then $p\mapsto x(p)$ is injective on $T$.
\end{enumerate}
\end{lemma}

\begin{proof}
(a) Let $N:=\{p\in S:y(p)<0\}$; $S$ is the disjoint union of $R(S)$, $A(S)$ and $N$. For $p\in R(S)$,
$\operatorname{mir}(p)\in S$ and $y(\operatorname{mir}p)=-y(p)<0$, so $\operatorname{mir}$ maps $R(S)$ into $N$; likewise
it maps $N$ into $R(S)$. As $\operatorname{mir}\circ\operatorname{mir}=\mathrm{id}$, these two maps are mutually inverse,
so $\card N=\card{R(S)}$ and $\card S=\card{R(S)}+\card{A(S)}+\card N=2\card{R(S)}+\card{A(S)}$.

(b) If $A(S)$ had three distinct points $a,b,d$, then $\orient(a,b,d)=(x(b)-x(a))\cdot0-0\cdot(x(d)-x(a))=0$, contradicting
general position (Lemma~\ref{fd:lem:collinear}).

(c) First let $p,q\in T$ with $y(p)>0$ and $x(p)=x(q)$, and suppose $p\ne q$. Then $p\ne\operatorname{mir}(p)$ (as
$y(p)\ne-y(p)$) and $q\ne\operatorname{mir}(p)$ (as $y(q)\ge0>-y(p)$), so $p,q,\operatorname{mir}(p)$ are three pairwise
distinct points of $S$. Since $x(q)=x(p)=x(\operatorname{mir}p)$,
$\orient(p,q,\operatorname{mir}p)=(x(q)-x(p))(-2y(p))-(y(q)-y(p))\cdot0=0$, contradicting general position. So $p=q$. Now
let $p,q\in T$ with $x(p)=x(q)$. If $y(p)>0$ or $y(q)>0$, the first part (with the roles of $p,q$ exchanged if needed)
gives $p=q$; otherwise $y(p)=y(q)=0$ and $p=q$ by hypothesis.
\end{proof}

\begin{lemma}[sorting by abscissa;\\
  \leanref{MirrorOdd.cfg}{MIRRORETV/lean/MirrorMain.lean:1584},\\
  \leanref{MirrorOdd.sortPt_image}{MIRRORETV/lean/MirrorMain.lean:1576}]\label{mb:lem:sort}
Let $T\subset\R^2$ be finite and in general position, and let $p\mapsto x(p)$ be injective on $T$. Then there is a
configuration $\mathrm{cfg}(T)$ of $\card T$ points (Lean \texttt{cfg T hx hgp}, with points \texttt{sortPt T hx}) all
of whose points lie in $T$, and $\mathcal P(\mathrm{cfg}(T))=T$.
\end{lemma}

\begin{proof}
This is the sorting of Definition~\ref{mm:def:config}. As $x$ is injective on $T$, the set $\{x(p):p\in T\}$ has $\card T$
elements $\xi_0<\dots<\xi_{\card T-1}$ (Lean \lean{Finset.orderEmbOfFin}); let $t_i$ be the point of $T$ with
$x(t_i)=\xi_i$. Then $x(t_i)$ is strictly increasing in $i$, so $i\mapsto t_i$ is injective and its image is a subset of
$T$ with $\card T$ elements, i.e.\ all of $T$. For $i<j<k$ the points $t_i,t_j,t_k$ are pairwise distinct points of $T$,
so $\orient(t_i,t_j,t_k)\ne0$ by general position (Lemma~\ref{fd:lem:collinear}). Hence $(t_i)$ is a configuration of
$\card T$ points with $\mathcal P=T$.
\end{proof}

\begin{lemma}[an axis point is not an interior point of a cap; written step]\label{mb:lem:interior}
Let $p,a,q\in\R^2$ with $x(p)<x(a)<x(q)$, $y(p)\ge0$, $y(q)\ge0$ and $y(a)=0$. Then $\orient(p,a,q)\ge0$.
\end{lemma}

\begin{proof}
If $y(a)=0$, then $\orient(p,a,q)=(x(a)-x(p))\,y(q)+y(p)\,(x(q)-x(a))$ (expand; checked symbolically), and both products
are $\ge0$. (In Lean this is the local claim \lean{interior} inside the proof of \lean{mirror_seven_of}.)
\end{proof}

The next lemma consists of four machine-generated Lean statements. We state them exactly and derive them from the
general mirror lemma of Theorems~\ref{mm:thm:mirror-polygon} and~\ref{mm:thm:axis-polygon}; this is a different proof
from the Lean one (see Remark~\ref{mb:rem:generated}).

\begin{lemma}[mirror heptagons and octagon;\\
  \leanref{MirrorOdd.cap_mirror_heptagon}{MIRRORETV/lean/MirrorMain.lean:212},\\
  \leanref{MirrorOdd.cap_mirror_octagon}{MIRRORETV/lean/MirrorMain.lean:392},\\
  \leanref{MirrorOdd.axis_left_heptagon}{MIRRORETV/lean/MirrorMain.lean:639},\\
  \leanref{MirrorOdd.axis_right_heptagon}{MIRRORETV/lean/MirrorMain.lean:816}]\label{mb:lem:polygons}
Let $P\subset\R^2$ be finite.
\begin{enumerate}[label=(\alph*)]
\item Let $c_0,c_1,c_2,c_3\in\R^2$ with $x(c_0)<x(c_1)<x(c_2)<x(c_3)$,
  $y(c_0),y(c_1),y(c_2),y(c_3)>0$ and $\orient(c_0,c_1,c_2)$, $\orient(c_0,c_1,c_3)$, $\orient(c_0,c_2,c_3)$,
  $\orient(c_1,c_2,c_3)<0$. If $P$ contains $c_0$, $\operatorname{mir}c_1$, $\operatorname{mir}c_2$,
  $\operatorname{mir}c_3$, $c_3$, $c_2$, $c_1$, then $P$ contains $7$ points in convex position.
  (\lean{cap_mirror_heptagon})
\item Under the same hypotheses on $c_0,\dots,c_3$: if $P$ contains $\operatorname{mir}c_0$,
  $\operatorname{mir}c_1$, $\operatorname{mir}c_2$, $\operatorname{mir}c_3$, $c_3$, $c_2$, $c_1$, $c_0$, then $P$
  contains $8$ points in convex position. (\lean{cap_mirror_octagon})
\item Let $a,c_0,c_1,c_2\in\R^2$ with $y(a)=0$, $x(a)<x(c_0)<x(c_1)<x(c_2)$,
  $y(c_0),y(c_1),y(c_2)>0$ and $\orient(a,c_0,c_1)$, $\orient(a,c_0,c_2)$, $\orient(a,c_1,c_2)$,
  $\orient(c_0,c_1,c_2)<0$. If $P$ contains $a$, $\operatorname{mir}c_0$, $\operatorname{mir}c_1$,
  $\operatorname{mir}c_2$, $c_2$, $c_1$, $c_0$, then $P$ contains $7$ points in convex position.
  (\lean{axis_left_heptagon})
\item Let $a,c_0,c_1,c_2\in\R^2$ with $y(a)=0$, $x(c_0)<x(c_1)<x(c_2)<x(a)$,
  $y(c_0),y(c_1),y(c_2)>0$ and $\orient(c_0,c_1,c_2)$, $\orient(c_0,c_1,a)$, $\orient(c_0,c_2,a)$,
  $\orient(c_1,c_2,a)<0$. If $P$ contains $\operatorname{mir}c_0$, $\operatorname{mir}c_1$, $\operatorname{mir}c_2$, $a$,
  $c_2$, $c_1$, $c_0$, then $P$ contains $7$ points in convex position. (\lean{axis_right_heptagon})
\end{enumerate}
\end{lemma}

\begin{proof}
We use the Lean statements of the general mirror lemma. Theorem~\ref{mm:thm:mirror-polygon} (\lean{cap_mirror_polygon},
\dirn{MIRRORETV/lean/MirrorMain.lean:2233}) says: if $m\ge2$, $(e_0,\dots,e_{m-1})$ has strictly increasing $x$,
$y(e_i)>0$ for all $i$, $\orient(e_i,e_j,e_l)<0$ for all $i<j<l$, and the finite set $P$ contains every $e_i$ and every
$\operatorname{mir}e_i$, then $P$ contains $2m$ points in convex position. Theorem~\ref{mm:thm:axis-polygon}
(\lean{axis_left_polygon}, line 2309; \lean{axis_right_polygon}, line 2403) says: if $m\ge3$, $(e_0,\dots,e_{m-1})$ has
strictly increasing $x$, $\orient(e_i,e_j,e_l)<0$ for all $i<j<l$, $P$ contains every $e_i$ and every
$\operatorname{mir}e_i$, and either $y(e_0)=0$ and $y(e_i)>0$ for $i>0$ (left), or $y(e_{m-1})=0$ and $y(e_i)>0$ for
$i<m-1$ (right), then $P$ contains $2m-1$ points in convex position.

(b) Apply the first statement with $m=4$ and $e=(c_0,c_1,c_2,c_3)$: the $x$-coordinates increase, all $y(c_i)>0$, the
four triples $i<j<l$ of $\{0,1,2,3\}$ are exactly the four given negative orientations, and $P$ contains all $c_i$ and
all $\operatorname{mir}c_i$. So $P$ contains $2\cdot4=8$ points in convex position.

(a) Let $P':=P\cup\{\operatorname{mir}c_0\}$. By (b) applied to $P'$ there is $T\subseteq P'$ with $\card T=8$ in convex
position. Then $T':=T\setminus\{\operatorname{mir}c_0\}\subseteq P$, $\card{T'}\ge7$, and $T'$ is in convex position
(Lemma~\ref{fd:lem:mono}(a)). So $P$ contains $\card{T'}$ points in convex position, hence $7$
(Lemma~\ref{fd:lem:mono}(d)).

(c) Apply the left form of the second statement with $m=4$ and $e=(a,c_0,c_1,c_2)$: the $x$-coordinates increase,
$y(e_0)=y(a)=0$, $y(e_i)>0$ for $i=1,2,3$, the four triples of $\{0,1,2,3\}$ are the four given negative orientations,
and $P$ contains $e_0=a$, $\operatorname{mir}e_0=\operatorname{mir}a=a$ (as $y(a)=0$), and $c_i$, $\operatorname{mir}c_i$
for $i=0,1,2$. So $P$ contains $2\cdot4-1=7$ points in convex position.

(d) Apply the right form with $m=4$ and $e=(c_0,c_1,c_2,a)$: $y(e_3)=y(a)=0$, $y(e_i)>0$ for $i<3$, the four triples are
$(c_0,c_1,c_2)$, $(c_0,c_1,a)$, $(c_0,c_2,a)$, $(c_1,c_2,a)$, and $\operatorname{mir}a=a\in P$. So $P$ contains $7$ points
in convex position.
\end{proof}

\begin{remark}[the Lean proofs of Lemma~\ref{mb:lem:polygons}]\label{mb:rem:generated}
In Lean these four statements are proved by generated code (\dirn{MIRRORETV/lean/MirrorMain.lean}, lines 22--991,
including the wrappers \lean{cyc7} and \lean{cyc8}; the rest of the generated block, up to line 1506, contains the
smaller mirror polygons used for $k=4,5,6$ and their wrappers): the listed points are taken in the cyclic order in which
they appear in the hypotheses
(for example $c_0,\operatorname{mir}c_1,\operatorname{mir}c_2,\operatorname{mir}c_3,c_3,c_2,c_1$ in (a)); pairwise
distinctness is read off from $x$- or $y$-coordinates; each of the $\binom73=35$ (for (b): $\binom83=56$) orientation
inequalities $\orient(q_i,q_j,q_k)>0$, $i<j<k$, is proved by \lean{linarith}, with hints found offline, from the
hypotheses and from products of a positive ordinate and a positive difference of abscissae, $y(c_i)\,(x(q)-x(q'))>0$
with $q,q'$ two of the four given points and $x(q')<x(q)$ (for example $y(c_0)\,(x(c_0)-x(a))$ in (c)); and the
positive-cycle criterion of Lemma~\ref{fd:lem:pos-cycle} (through the wrappers \lean{cyc7}, \lean{cyc8}) concludes. The
general lemmas used above (Theorems~\ref{mm:thm:mirror-polygon} and~\ref{mm:thm:axis-polygon}) are proved in Lean by
the same method in general form. Statements (b), (c) and (d) are special cases of them ($m=4$), and (a) follows from (b)
applied to $P\cup\{\operatorname{mir}c_0\}$ by deleting the point $\operatorname{mir}c_0$; so nothing is lost by the
derivation above.
\end{remark}

%======================================================================
\subsection{Seven points: the bound 32}\label{mb:ssec:seven}

\begin{lemma}[the reduction for $k=7$;\\
  \leanref{MirrorOdd.mirror_seven_of}{MIRRORETV/lean/MirrorMain.lean:1676}]\label{mb:lem:seven-of}
Let (Hcap), (HL) and (HR) denote the statements of Lemma~\ref{mb:lem:polygons}(a), (c) and (d) respectively, each
universally quantified over the four points ($c_0,\dots,c_3$, resp.\ $a,c_0,c_1,c_2$) and over all finite $P\subset\R^2$.
Assume (Hcap), (HL) and (HR). Let $S\subset\R^2$ be finite, in general position and $\operatorname{mir}$-symmetric, and
suppose that $S$ does not contain $7$ points in convex position. Then $\card S\le32$.
\end{lemma}

\begin{proof}
\emph{Key claim.} Let $T\subseteq S$ with $y(p)\ge0$ for all $p\in T$ and with $p=q$ for all $p,q\in T$ such that
$y(p)=0=y(q)$ (at most one point of $T$ on the axis). Then $\card T\le16$.

By Lemma~\ref{mb:lem:split}(c), $x$ is injective on $T$; $T$ is in general position (Lemma~\ref{fd:lem:mono}(b)). Let
$c:=\mathrm{cfg}(T)$ be the configuration of $\card T$ points of Lemma~\ref{mb:lem:sort}, so $\mathcal P(c)=T$. We apply
Proposition~\ref{mb:prop:cap16} to $c$. Its second hypothesis holds because $\mathcal P(c)=T\subseteq S$ and $S$ does not
contain $7$ points in convex position (Lemma~\ref{fd:lem:mono}(c)). For the first, suppose $c.\mathsf{HasCap}(4)$, with
$f\colon[4]\to[\card T]$ strictly increasing, and put $d_i:=c_{f(i)}$. Then $d_0,\dots,d_3\in T\subseteq S$,
$x(d_0)<x(d_1)<x(d_2)<x(d_3)$, $\orient(d_i,d_j,d_l)<0$ for all $i<j<l$, $\operatorname{mir}(d_i)\in S$ by symmetry, and
$y(d_i)\ge0$. We distinguish cases.
\begin{itemize}
\item $y(d_0),y(d_1),y(d_2),y(d_3)>0$. Then (Hcap) with $P=S$ shows that $S$ contains $7$ points in convex position:
  contradiction.
\item $y(d_0)>0$, $y(d_3)>0$, $y(d_1)>0$, $y(d_2)=0$. Lemma~\ref{mb:lem:interior} for $(d_1,d_2,d_3)$ gives
  $\orient(d_1,d_2,d_3)\ge0$: contradiction.
\item $y(d_0)>0$, $y(d_3)>0$, $y(d_1)=0$. Lemma~\ref{mb:lem:interior} for $(d_0,d_1,d_2)$ gives
  $\orient(d_0,d_1,d_2)\ge0$: contradiction.
\item $y(d_0)>0$, $y(d_3)=0$. If $y(d_1)=0$, then $d_1=d_3$ (at most one axis point in $T$), contradicting
  $x(d_1)<x(d_3)$; so $y(d_1)>0$, and likewise $y(d_2)>0$. Then (HR) with $a=d_3$ and $(c_0,c_1,c_2)=(d_0,d_1,d_2)$, and
  $P=S$, gives $7$ points of $S$ in convex position: contradiction. (The hypotheses of (HR) are $y(d_3)=0$, the order
  $x(d_0)<x(d_1)<x(d_2)<x(d_3)$, the signs $y(d_i)>0$ for $i<3$, the four negative orientations, and $\operatorname{mir}
  d_0,\operatorname{mir}d_1,\operatorname{mir}d_2,d_3,d_2,d_1,d_0\in S$.)
\item $y(d_0)=0$. As before, $y(d_i)>0$ for $i=1,2,3$ (otherwise $d_i=d_0$, contradicting $x(d_0)<x(d_i)$). Then (HL) with
  $a=d_0$ and $(c_0,c_1,c_2)=(d_1,d_2,d_3)$, and $P=S$ (which contains $d_0,\operatorname{mir}d_1,\operatorname{mir}d_2,
  \operatorname{mir}d_3,d_3,d_2,d_1$), gives a contradiction.
\end{itemize}
These cases are exhaustive (since every $y(d_i)\ge0$). Hence $\neg\,c.\mathsf{HasCap}(4)$, and
Proposition~\ref{mb:prop:cap16} gives $\card T=\card{\mathcal P(c)}\le16$.

\emph{Counting.} Write $R=R(S)$ and $A=A(S)$. By Lemma~\ref{mb:lem:split}(a),(b), $\card S=2\card R+\card A$ and
$\card A\le2$. The key claim applies to $T=R$ (all $y>0$, no axis point), so $\card R\le16$. If $A=\emptyset$, then
$\card S=2\card R\le32$. Otherwise pick $a\in A$; then $a\notin R$, and $T=R\cup\{a\}\subseteq S$ has all $y\ge0$ and
exactly one axis point, so the key claim gives $\card R+1\le16$, and $\card S=2\card R+\card A\le2\cdot15+2=32$.
\end{proof}

\begin{definition}[reflection in a line given by a unit normal; the frame;\\
  \leanref{MirrorOdd.lineRefl}{MIRRORETV/lean/MirrorMain.lean:1879},\\
  \leanref{MirrorOdd.frame}{MIRRORETV/lean/MirrorMain.lean:1883},\\
  \leanref{MirrorOdd.unframe}{MIRRORETV/lean/MirrorMain.lean:1892}]\label{mb:def:frame}
For $n_0,n_1,c\in\R$ and $p,q\in\R^2$ put
\begin{align*}
\mathrm{lineRefl}_{n_0,n_1,c}(p)&:=\bigl(p_0-2(n_0p_0+n_1p_1-c)\,n_0,\ p_1-2(n_0p_0+n_1p_1-c)\,n_1\bigr),\\
\mathrm{frame}_{n_0,n_1,c}(p)&:=\bigl(n_1(p_0-cn_0)-n_0(p_1-cn_1),\ n_0(p_0-cn_0)+n_1(p_1-cn_1)\bigr),\\
\mathrm{unframe}_{n_0,n_1,c}(q)&:=\bigl(n_1q_0+n_0q_1+cn_0,\ -n_0q_0+n_1q_1+cn_1\bigr).
\end{align*}
$\mathrm{unframe}$ is an affine map (Lean: a Mathlib \lean{AffineMap} of \lean{Point}) with linear part
$(q_0,q_1)\mapsto(n_1q_0+n_0q_1,-n_0q_0+n_1q_1)$ (\lean{MirrorOdd.unframeLin}, line 1886). When $n_0^2+n_1^2=1$,
$\mathrm{lineRefl}$ is the Euclidean reflection in the line $\{p:n_0p_0+n_1p_1=c\}$ (every line has this form), and
$\mathrm{frame}$ is a rotation composed with a translation which takes this line onto the $x$-axis; no condition on
$n_0,n_1$ is built into the definitions.
\end{definition}

\begin{lemma}[identities of the frame;\\
  \leanref{MirrorOdd.frame_lineRefl}{MIRRORETV/lean/MirrorMain.lean:1914},\\
  \leanref{MirrorOdd.orient_frame}{MIRRORETV/lean/MirrorMain.lean:1920},\\
  \leanref{MirrorOdd.unframe_frame}{MIRRORETV/lean/MirrorMain.lean:1900},\\
  \leanref{MirrorOdd.frame_unframe}{MIRRORETV/lean/MirrorMain.lean:1907},\\
  \leanref{MirrorOdd.frame_injective}{MIRRORETV/lean/MirrorMain.lean:1925},\\
  \leanref{MirrorOdd.unframe_injective}{MIRRORETV/lean/MirrorMain.lean:1930}]\label{mb:lem:frame}
Let $n_0^2+n_1^2=1$ and $c\in\R$; write $\mathrm{frame}$, $\mathrm{unframe}$, $\mathrm{lineRefl}$ for the maps with these
parameters. Then for all $p,q,a,b,d\in\R^2$:
\begin{enumerate}[label=(\alph*)]
\item $\mathrm{frame}(\mathrm{lineRefl}(p))=\operatorname{mir}(\mathrm{frame}(p))$;
\item $\orient(\mathrm{frame}(a),\mathrm{frame}(b),\mathrm{frame}(d))=\orient(a,b,d)$;
\item $\mathrm{unframe}(\mathrm{frame}(p))=p$ and $\mathrm{frame}(\mathrm{unframe}(q))=q$; in particular $\mathrm{frame}$
  and $\mathrm{unframe}$ are injective.
\end{enumerate}
\end{lemma}

\begin{proof}
Put $\nu:=n_0^2+n_1^2$. Expanding the definitions gives the following polynomial identities, valid for all values of the
variables (we checked each symbolically; they are the \lean{linear_combination} certificates of the Lean proofs):
\begin{align*}
\mathrm{frame}(\mathrm{lineRefl}\,p)-\operatorname{mir}(\mathrm{frame}\,p)&=\bigl(0,\ -2(n_0p_0+n_1p_1)(\nu-1)\bigr),\\
\orient(\mathrm{frame}\,a,\mathrm{frame}\,b,\mathrm{frame}\,d)-\orient(a,b,d)&=(\nu-1)\,\orient(a,b,d),\\
\mathrm{unframe}(\mathrm{frame}\,p)-p&=(\nu-1)\,(p_0-cn_0,\ p_1-cn_1),\\
\mathrm{frame}(\mathrm{unframe}\,q)-q&=(\nu-1)\,(q_0,\ q_1).
\end{align*}
With $\nu=1$ all right-hand sides vanish, which proves (a), (b) and the two equalities in (c); a map with a left inverse
is injective, and $\mathrm{frame}$ and $\mathrm{unframe}$ are left inverses of each other. (Conceptually,
$\mathrm{frame}(p)=M(p-c\,n)$ and $\mathrm{unframe}(q)=M^{\mathsf T}q+c\,n$ with $n=(n_0,n_1)$ and
$M=\left(\begin{smallmatrix}n_1&-n_0\\ n_0&n_1\end{smallmatrix}\right)$, $\det M=\nu$, $M^{\mathsf T}M=MM^{\mathsf T}=\nu I$;
(b) is then Lemma~\ref{fd:lem:affine-orient}, and the second coordinate of $\mathrm{frame}(p)$ is $n_0p_0+n_1p_1-c\nu$,
which for $\nu=1$ is the signed distance from $p$ to the line.)
\end{proof}

\begin{lemma}[convex polygons under injective affine maps;\\
  \leanref{MirrorOdd.hasConvexNGon_affineMap}{MIRRORETV/lean/MirrorMain.lean:1859}]\label{mb:lem:affine}
Let $S\subset\R^2$ be finite, $n\in\N$, and $g\colon\R^2\to\R^2$ an injective affine map. If $S$ contains $n$ points in
convex position, then so does $g(S)$.
\end{lemma}

\begin{proof}
Let $T\subseteq S$ with $\card T=n$ be in convex position. Then $g(T)\subseteq g(S)$, $\card{g(T)}=n$ because $g$ is
injective, and $g(T)$ is in convex position by Lemma~\ref{fd:lem:affine-cp}. (Lean argues directly: if
$g(p)\in\conv(g(T)\setminus\{g(p)\})=\conv(g(T\setminus\{p\}))=g(\conv(T\setminus\{p\}))$, injectivity gives
$p\in\conv(T\setminus\{p\})$.)
\end{proof}

\begin{lemma}[transport to the $x$-axis;\\
  \leanref{MirrorOdd.transport}{MIRRORETV/lean/MirrorMain.lean:1937}]\label{mb:lem:transport}
Let $n_0,n_1,c\in\R$ with $n_0^2+n_1^2=1$, let $S\subset\R^2$ be finite and $m\in\N$, and assume that $S$ is in general
position, that $\mathrm{lineRefl}_{n_0,n_1,c}(p)\in S$ for all $p\in S$, and that $S$ does not contain $m$ points in
convex position. Put $S':=\mathrm{frame}_{n_0,n_1,c}(S)$. Then $S'$ is in general position, $S'$ is
$\operatorname{mir}$-symmetric, $S'$ does not contain $m$ points in convex position, and $\card{S'}=\card S$.
\end{lemma}

\begin{proof}
The same statement is part (a) of Theorem~\ref{mm:thm:line}; we include the short proof to keep this section
self-contained. We use Lemma~\ref{mb:lem:frame} throughout. $\card{S'}=\card S$ because $\mathrm{frame}$ is injective. General position:
let $q_1,q_2,q_3\in S'$ be pairwise distinct, $q_k=\mathrm{frame}(p_k)$ with $p_k\in S$; the $p_k$ are pairwise distinct,
so $\orient(q_1,q_2,q_3)=\orient(p_1,p_2,p_3)\ne0$ by Lemma~\ref{mb:lem:frame}(b) and Lemma~\ref{fd:lem:collinear}; by
Lemma~\ref{fd:lem:collinear} again $S'$ is in general position. Symmetry: for $q=\mathrm{frame}(p)\in S'$,
$\operatorname{mir}(q)=\mathrm{frame}(\mathrm{lineRefl}(p))\in S'$ because $\mathrm{lineRefl}(p)\in S$. No $m$ points in
convex position: if $S'$ contained $m$ points in convex position, then so would $\mathrm{unframe}(S')$ by
Lemma~\ref{mb:lem:affine} ($\mathrm{unframe}$ is affine and injective); but $\mathrm{unframe}(S')=S$ by
Lemma~\ref{mb:lem:frame}(c).
\end{proof}

\begin{theorem}[mirror-symmetric sets without convex heptagons;\\
  \leanref{MirrorOdd.mirror_seven}{MIRRORETV/lean/MirrorMain.lean:1803},\\
  \leanref{MirrorOdd.mirror_seven_line}{MIRRORETV/lean/MirrorMain.lean:1969}]\label{mb:thm:seven}
The following two statements hold, (a) for the reflection in the $x$-axis and (b) for the reflection in an arbitrary
line.
\begin{enumerate}[label=(\alph*)]
\item Let $S\subset\R^2$ be finite, in general position and $\operatorname{mir}$-symmetric. If $S$ does not contain $7$
  points in convex position, then $\card S\le32$.
\item Let $n_0,n_1,c\in\R$ with $n_0^2+n_1^2=1$, and let $S\subset\R^2$ be finite, in general position, with
  $\mathrm{lineRefl}_{n_0,n_1,c}(p)\in S$ for every $p\in S$. If $S$ does not contain $7$ points in convex position, then
  $\card S\le32$.
\end{enumerate}
\end{theorem}

Status: Lean-checked (\dirn{MIRRORETV/lean/MirrorMain.out}, lines 848 and 850; also \dirn{SYMLEAN/mirror/axioms.out}).
The proof below is a written proof (not refereed); it rests on the cited Theorem~\ref{mb:thm:baek} (through
Proposition~\ref{mb:prop:cap16}) and on Theorems~\ref{mm:thm:mirror-polygon} and~\ref{mm:thm:axis-polygon}.

\begin{proof}
(a) By Lemma~\ref{mb:lem:polygons}(a), (c), (d), the statements (Hcap), (HL), (HR) of Lemma~\ref{mb:lem:seven-of} hold;
that lemma gives $\card S\le32$. (Lean: \texttt{mirror\_seven\_of cap\_mirror\_heptagon axis\_left\_heptagon
axis\_right\_heptagon}.)

(b) By Lemma~\ref{mb:lem:transport} with $m=7$, $S'=\mathrm{frame}_{n_0,n_1,c}(S)$ is in general position,
$\operatorname{mir}$-symmetric, does not contain $7$ points in convex position, and $\card{S'}=\card S$. By (a),
$\card S=\card{S'}\le32$.
\end{proof}

\begin{remark}\label{mb:rem:seven}
Lean also contains \texttt{MirrorOdd.}\allowbreak\texttt{mirror\_seven\textquotesingle} (\dirn{MIRRORETV/lean/MirrorMain.lean:2631}),
which proves the statement of Theorem~\ref{mb:thm:seven}(a) again as the instance $j=3$, $M=16$ of
Theorem~\ref{mm:thm:odd} (\lean{mirror_odd_general}), with the hypothesis on $M$ supplied by
Proposition~\ref{mb:prop:cap16}. This is a second Lean proof: in place of \lean{mirror_seven_of} and the generated
polygon lemmas of Lemma~\ref{mb:lem:polygons} it uses the general reduction \lean{mirror_odd_general}, which rests on
the general mirror lemma (Theorems~\ref{mm:thm:mirror-polygon} and~\ref{mm:thm:axis-polygon}).
\end{remark}

%======================================================================
\subsection{Eight points: the bound 46}\label{mb:ssec:eight}

\begin{theorem}[mirror-symmetric sets without convex octagons;\\
  \leanref{MirrorOdd.mirror_eight}{MIRRORETV/lean/MirrorMain.lean:1820},\\
  \leanref{MirrorOdd.mirror_eight_line}{MIRRORETV/lean/MirrorMain.lean:1977}]\label{mb:thm:eight}
The following two statements hold, (a) for the reflection in the $x$-axis and (b) for the reflection in an arbitrary
line.
\begin{enumerate}[label=(\alph*)]
\item Let $S\subset\R^2$ be finite, in general position and $\operatorname{mir}$-symmetric. If $S$ does not contain $8$
  points in convex position, then $\card S\le46$.
\item Let $n_0,n_1,c\in\R$ with $n_0^2+n_1^2=1$, and let $S\subset\R^2$ be finite, in general position, with
  $\mathrm{lineRefl}_{n_0,n_1,c}(p)\in S$ for every $p\in S$. If $S$ does not contain $8$ points in convex position, then
  $\card S\le46$.
\end{enumerate}
\end{theorem}

Status: Lean-checked (\dirn{MIRRORETV/lean/MirrorMain.out}, lines 849 and 851). The proof below is a written proof (not
refereed); it rests on the cited Theorem~\ref{mb:thm:baek} (through Proposition~\ref{mb:prop:cap22}) and on
Theorem~\ref{mm:thm:mirror-polygon}.

\begin{proof}
(a) Let $R=R(S)$ and $A=A(S)$. Then $R\subseteq S$ and $y>0$ on $R$, so $x$ is injective on $R$ by
Lemma~\ref{mb:lem:split}(c) (the condition on axis points is vacuous), and $R$ is in general position
(Lemma~\ref{fd:lem:mono}(b)). Let $c:=\mathrm{cfg}(R)$ (Lemma~\ref{mb:lem:sort}), so $\mathcal P(c)=R$. We apply
Proposition~\ref{mb:prop:cap22} to $c$. Its second hypothesis holds since $\mathcal P(c)=R\subseteq S$
(Lemma~\ref{fd:lem:mono}(c)). For the first, suppose $c.\mathsf{HasCap}(4)$ with $f\colon[4]\to[\card R]$ and put
$d_i:=c_{f(i)}$: then $x(d_0)<\dots<x(d_3)$, all four ordered triples have negative orientation, $y(d_i)>0$ (as
$d_i\in R$), and $d_i,\operatorname{mir}(d_i)\in S$. Lemma~\ref{mb:lem:polygons}(b) with $P=S$ shows that $S$ contains $8$
points in convex position: a contradiction. Hence $\card R\le22$, and by Lemma~\ref{mb:lem:split}(a),(b)
$\card S=2\card R+\card A\le44+2=46$.

(b) By Lemma~\ref{mb:lem:transport} with $m=8$ and part (a), $\card S=\card{\mathrm{frame}(S)}\le46$.
\end{proof}

\begin{remark}\label{mb:rem:eight}
Lean also contains \texttt{MirrorOdd.}\allowbreak\texttt{mirror\_eight\textquotesingle}
(\dirn{MIRRORETV/lean/MirrorMain.lean:2636}), which re-derives Theorem~\ref{mb:thm:eight}(a) as the instance $j=4$,
$M=22$ of Theorem~\ref{mm:thm:even} (\lean{mirror_even_general}) with Proposition~\ref{mb:prop:cap22}.
Only the points strictly above the axis are used in (a); the two possible axis points are simply counted.
\end{remark}

%======================================================================
\subsection{Nine points: the bound 226}\label{mb:ssec:nine}

\begin{theorem}[the capacity bound 113; mirror-symmetric sets without convex nonagons;\\
  \leanref{C12Check.config_capacity113}{MIRRORETV/lean/MirrorNine_append.lean:24},\\
  \leanref{C12Check.mirror_nine}{MIRRORETV/lean/MirrorNine_append.lean:131},\\
  \leanref{C12Check.mirror_nine_line}{MIRRORETV/lean/MirrorNine_append.lean:136}]\label{mb:thm:nine}\raggedright
The following three statements hold: (a) is a capacity bound in Baek's model of Definition~\ref{mb:def:config} (no
$5$-cap and no $9$-gon), (b) and (c) are the mirror bounds for the reflection in the $x$-axis and in an arbitrary line.
\begin{enumerate}[label=(\alph*)]
\item Let $(\alpha,<)$ be a linearly ordered set, $C$ a configuration on $\alpha$ and $S\subseteq\alpha$ finite with
  $\neg\mathsf{HasNCap}_C(5,S)$ and $\neg\mathsf{HasNGon}_C(9,S)$. Then $\card S\le113$.
\item Let $S\subset\R^2$ be finite, in general position and $\operatorname{mir}$-symmetric. If $S$ does not contain $9$
  points in convex position, then $\card S\le226$.
\item Let $n_0,n_1,c\in\R$ with $n_0^2+n_1^2=1$, and let $S\subset\R^2$ be finite, in general position, with
  $\mathrm{lineRefl}_{n_0,n_1,c}(p)\in S$ for every $p\in S$. If $S$ does not contain $9$ points in convex position, then
  $\card S\le226$.
\end{enumerate}
\end{theorem}

Status: Lean-checked (\dirn{MIRRORETV/lean/MirrorNine.out}, lines 879, 881 and 882). The proofs below are written proofs
(not refereed); they rest on the cited Theorems~\ref{mb:thm:baek} and~\ref{mb:thm:capcup} and, for (b) and (c), on
Theorems~\ref{mm:thm:odd} and~\ref{mm:thm:line}. Part (a) is proved first; Proposition~\ref{mb:prop:cap113} transfers
it to point configurations; (b) and (c) follow.

\begin{proof}[Proof of Theorem~\ref{mb:thm:nine}(a)]
We follow the Lean proof. Let $T$ be the set of $p\in S$ that are the first entry of some $C$-cap of length $4$ with all
entries in $S$. Then $T\subseteq S$ and $\card S=\card{S\setminus T}+\card T$.

\emph{Claim 1: $\card{S\setminus T}\le29$.} Suppose $\card{S\setminus T}\ge30=\binom{6+2}{2}+2$ (with $\binom82=28$).
Theorem~\ref{mb:thm:baek} with $n=6$ gives $\mathsf{HasNCap}_C(4,S\setminus T)$ or $\mathsf{HasNGon}_C(9,S\setminus T)$.
In the first case let $(h,\dots)$ be a $C$-cap of length $4$ with entries in $S\setminus T$; its entries lie in $S$, so
its first entry $h$ lies in $T$ by definition of $T$, while $h\in S\setminus T$: a contradiction. In the second case
Lemma~\ref{mb:lem:config-mono}(a) gives $\mathsf{HasNGon}_C(9,S)$, contrary to the hypothesis.

\emph{Claim 2: $\card T\le84$.} Suppose $\card T>84=\binom{3+6}{3}$. Theorem~\ref{mb:thm:capcup} with $a=3$, $b=6$ gives
$\mathsf{HasNCap}_C(5,T)$ or $\mathsf{HasNCup}_C(8,T)$. The first is impossible, since a $5$-cap in $T$ is one in $S$
(Lemma~\ref{mb:lem:config-mono}(a)). In the second case let $(z_0,\dots,z_7)$ be a $C$-cup with entries in $T$, and put
$p:=z_6$, $q:=z_7$. As $q\in T$, there is a $C$-cap $(q,r,r_2,r_3)$ with entries in $S$. Note $p<q<r$, since cups and
caps are strictly increasing.
\begin{itemize}
\item If $\mathrm{Cup3}_C(p,q,r)$, then $(z_0,\dots,z_7,r)$ is strictly increasing and all its consecutive triples
  satisfy $\mathrm{Cup3}_C$ (those of the cup, and $(z_6,z_7,r)=(p,q,r)$), so it is a $C$-cup of length $9$ with entries
  in $S$ (as $T\subseteq S$). By Lemma~\ref{mb:lem:config-mono}(b) (with $9\ge2$), $\mathsf{HasNGon}_C(9,S)$: a
  contradiction.
\item If $\mathrm{Cap3}_C(p,q,r)$, then $(p,q,r,r_2,r_3)$ is strictly increasing and all its consecutive triples satisfy
  $\mathrm{Cap3}_C$ (the triple $(p,q,r)$ and those of the cap), so it is a $C$-cap of length $5$ with entries in $S$
  ($p\in T\subseteq S$): a contradiction with $\neg\mathsf{HasNCap}_C(5,S)$.
\end{itemize}

Hence $\card S=\card{S\setminus T}+\card T\le29+84=113$.
\end{proof}

\begin{proposition}[the capacity bound 113 for point configurations;\\
  \leanref{C12Check.capacity113}{MIRRORETV/lean/MirrorNine_append.lean:116}]\label{mb:prop:cap113}
Let $c$ be a configuration of $n$ points with $\neg\,c.\mathsf{HasCap}(5)$ such that $\mathcal P(c)$ does not contain $9$
points in convex position. Then $n\le113$.
\end{proposition}

\begin{proof}
Apply Theorem~\ref{mb:thm:nine}(a) (proved above) to $C=C_c$ and $S=[n]$. A $C_c$-cap of length $5$ in $[n]$ would give
$c.\mathsf{HasCap}(5)$ (Lemma~\ref{mb:lem:lists}(c)), and a $9$-gon of $C_c$ in $[n]$ would give $9$ points of
$\mathcal P(c)$ in convex position (Proposition~\ref{mb:prop:gon}). Hence $n=\card{[n]}\le113$.
\end{proof}

\begin{proof}[Proof of Theorem~\ref{mb:thm:nine}(b) and~(c)]
Theorem~\ref{mm:thm:odd} (\lean{mirror_odd_general}, \dirn{MIRRORETV/lean/MirrorMain.lean:2498}) states: let $j\ge2$ and
$M$ be natural numbers such that every configuration $e$ of $r$ points (for every $r\in\N$) with
$\neg\,e.\mathsf{HasCap}(j+1)$ and without $2j+1$ points of $\mathcal P(e)$ in convex position has $r\le M$; then every
finite $\operatorname{mir}$-symmetric set in general position without $2j+1$ points in convex position has at most $2M$
points. Theorem~\ref{mm:thm:line} contains the same statement for $\mathrm{lineRefl}_{n_0,n_1,c}$-symmetric sets with
$n_0^2+n_1^2=1$ (\lean{mirror_odd_general_line}, line 2641).

Take $j=4$ and $M=113$. Then $j\ge2$, $j+1=5$, $2j+1=9$, and the hypothesis on $M$ is exactly
Proposition~\ref{mb:prop:cap113}. So (b) follows from \lean{mirror_odd_general} and (c) from
\lean{mirror_odd_general_line}, with $2M=226$. (Alternatively, (c) follows from (b) by Lemma~\ref{mb:lem:transport}.)
\end{proof}

%======================================================================
\subsection{Reflections given by a point and a direction; the three-class reduction}\label{mb:ssec:three}

The case classes of Definition~\ref{sr:def:classes} use a second parametrisation of reflections.

\begin{definition}[reflection through a point with a unit direction;\\
  \leanref{ErdosSzekeres.SymReduction.reflLine}{SYMLEAN/mirror/ReductionThreeClasses.lean:449}]\label{mb:def:reflline}
For $a,u,p\in\R^2$ put $d:=(p_0-a_0)u_0+(p_1-a_1)u_1$ and
\[
\mathrm{reflLine}_{a,u}(p):=\bigl(a_0+2du_0-(p_0-a_0),\ a_1+2du_1-(p_1-a_1)\bigr).
\]
For $u_0^2+u_1^2=1$ this is the reflection in the line $a+\R u$. (The same definition is
\dirn{SYMLEAN/SymmetricReduction.lean:423}; the class statements use it with the hypothesis $u_0^2+u_1^2=1$.)
\end{definition}

\begin{theorem}[the two reflections agree; the bound 32 for $\mathrm{reflLine}$;\\
  \leanref{ErdosSzekeres.SymReduction.reflLine_eq_lineRefl}{SYMLEAN/mirror/ReductionThreeClasses.lean:4253},\\
  \leanref{ErdosSzekeres.SymReduction.card_le_32_of_reflLine}{SYMLEAN/mirror/ReductionThreeClasses.lean:4264}]\label{mb:thm:refl32}\raggedright
Let $a,u\in\R^2$ with $u_0^2+u_1^2=1$.
\begin{enumerate}[label=(\alph*)]
\item $\mathrm{reflLine}_{a,u}=\mathrm{lineRefl}_{n_0,n_1,c}$ as maps $\R^2\to\R^2$, where $n_0=-u_1$, $n_1=u_0$ and
  $c=-u_1a_0+u_0a_1$.
\item Let $S\subset\R^2$ be finite and in general position with $\mathrm{reflLine}_{a,u}(p)\in S$ for every $p\in S$
  (Lean \texttt{Invariant (reflLine a u) S}, Definition~\ref{fd:def:inv}). If $S$ does not contain $7$ points in convex
  position, then $\card S\le32$.
\end{enumerate}
\end{theorem}

Status: Lean-checked (\dirn{SYMLEAN/mirror/axioms.out}). Part (b) puts no condition on the number of points of $S$ on the
mirror. The written proof of (b) rests on Theorem~\ref{mb:thm:seven}(b), hence on the cited Theorem~\ref{mb:thm:baek}
and on Theorems~\ref{mm:thm:mirror-polygon} and~\ref{mm:thm:axis-polygon}.

\begin{proof}
(a) Fix $p$ and put $w:=p-a$. In the first coordinate,
\[
\mathrm{reflLine}_{a,u}(p)_0-\mathrm{lineRefl}_{n_0,n_1,c}(p)_0=(2a_0-p_0+2du_0)-\bigl(p_0+2(-u_1p_0+u_0p_1-c)u_1\bigr)
=2w_0\,(u_0^2+u_1^2-1),
\]
using $d=w_0u_0+w_1u_1$ and $-u_1p_0+u_0p_1-c=-u_1w_0+u_0w_1$; in the second coordinate the difference is
$(2a_1-p_1+2du_1)-\bigl(p_1-2(-u_1w_0+u_0w_1)u_0\bigr)=2w_1\,(u_0^2+u_1^2-1)$. (We checked both identities symbolically;
they are the Lean certificates \texttt{linear\_combination (2 * (p i - a i)) * hu}.) Both vanish since $u_0^2+u_1^2=1$.
Geometrically: $n=(-u_1,u_0)$ is $u$ turned by $90^\circ$, $(u,n)$ is an orthonormal basis, $w=\langle w,u\rangle u+\langle
w,n\rangle n$, $\langle n,p\rangle-c=\langle n,w\rangle$, and both maps send $p$ to $a+\langle w,u\rangle u-\langle w,n
\rangle n$.

(b) $n_0^2+n_1^2=u_1^2+u_0^2=1$, and for $p\in S$, $\mathrm{lineRefl}_{n_0,n_1,c}(p)=\mathrm{reflLine}_{a,u}(p)\in S$ by
(a). Theorem~\ref{mb:thm:seven}(b) gives $\card S\le32$.
\end{proof}

\begin{theorem}[the two mirror case classes hold;\\
  \leanref{ErdosSzekeres.SymReduction.mirrorOneAxis33_proved}{SYMLEAN/mirror/ReductionThreeClasses.lean:4274},\\
  \leanref{ErdosSzekeres.SymReduction.mirrorNoAxis34_proved}{SYMLEAN/mirror/ReductionThreeClasses.lean:4282}]\label{mb:thm:classes}\raggedright
The propositions \lean{MirrorOneAxis33} and \lean{MirrorNoAxis34} of Definition~\ref{sr:def:classes} hold. Explicitly:
\begin{enumerate}[label=(\alph*)]
\item Let $S\subset\R^2$ be finite and $a,u\in\R^2$ with $u_0^2+u_1^2=1$. If $S$ is in general position, $\card S=33$,
  $\mathrm{reflLine}_{a,u}(p)\in S$ for all $p\in S$, and there is $x\in S$ with $\mathrm{reflLine}_{a,u}(x)=x$ such that
  every $y\in S$ with $\mathrm{reflLine}_{a,u}(y)=y$ equals $x$, then $S$ contains $7$ points in convex position.
\item Let $S\subset\R^2$ be finite and $a,u\in\R^2$ with $u_0^2+u_1^2=1$. If $S$ is in general position, $\card S=34$,
  $\mathrm{reflLine}_{a,u}(p)\in S$ for all $p\in S$, and $\mathrm{reflLine}_{a,u}(p)\ne p$ for all $p\in S$, then $S$
  contains $7$ points in convex position.
\end{enumerate}
\end{theorem}

Status: Lean-checked (\dirn{SYMLEAN/mirror/axioms.out}); no SAT certificate and no unproved hypothesis is involved; the
written proof rests, through Theorems~\ref{mb:thm:refl32}(b) and~\ref{mb:thm:seven}(b), on the cited
Theorem~\ref{mb:thm:baek} and on Theorems~\ref{mm:thm:mirror-polygon} and~\ref{mm:thm:axis-polygon}.

\begin{proof}
In both cases, if $S$ did not contain $7$ points in convex position, Theorem~\ref{mb:thm:refl32}(b) would give
$\card S\le32$, contradicting $\card S=33$, resp.\ $\card S=34$. The last hypothesis of (a) (exactly one point on the
mirror) and of (b) (no point on the mirror) is not used.
\end{proof}

\begin{theorem}[the reduction with three case classes; \emph{conditional};\\
  \leanref{ErdosSzekeres.SymReduction.heptagon_of_nontrivial_isometry_3}{SYMLEAN/mirror/ReductionThreeClasses.lean:4292},\\
  \leanref{ErdosSzekeres.SymReduction.heptagon_of_nontrivial_affineEquiv_3}{SYMLEAN/mirror/ReductionThreeClasses.lean:4303}]\label{mb:thm:main3}\raggedright
Assume the three propositions \lean{HalfTurn34}, \lean{Threefold33} and \lean{Fivefold35} of
Definition~\ref{sr:def:classes} (they are explicit hypotheses of both statements).
\begin{enumerate}[label=(\alph*)]
\item Let $S\subset\R^2$ be finite and in general position with $33\le\card S$, and let $g\colon\R^2\to\R^2$ be a map
  with $\mathrm{eDist2}(g(p),g(q))=\mathrm{eDist2}(p,q)$ for all $p,q\in\R^2$ (\texttt{IsEucIsometry g},
  Definition~\ref{sr:def:iso}; $\mathrm{eDist2}$ is the squared Euclidean distance; no continuity or surjectivity is
  assumed), $g\ne\mathrm{id}$, and $g(p)\in S$ for all $p\in S$. Then $S$ contains $7$ points in convex position.
\item Let $S\subset\R^2$ be finite and in general position with $33\le\card S$, and let $g$ be an affine bijection of
  $\R^2$ (Lean: $g$ is an \lean{AffineEquiv} of \lean{Point}) with $g\ne\mathrm{id}$ (Lean: $g$ differs from
  \lean{AffineEquiv.refl}; two affine equivalences are equal iff they agree at every point) and $g(p)\in S$ for all
  $p\in S$ (the invariance uses $g$ as a map). Then $S$ contains $7$ points in convex position.
\end{enumerate}
\end{theorem}

Status: Lean-checked \emph{as a conditional statement} (\dirn{SYMLEAN/mirror/axioms.out}: only the three standard axioms;
the three case classes are binders of the theorems, not axioms). The theorem does not establish its conclusion
unconditionally; see Remark~\ref{mb:rem:main3}.

\begin{proof}
(a) Theorem~\ref{sr:thm:main} (\lean{heptagon_of_nontrivial_isometry}) gives the conclusion under the five hypotheses
\lean{HalfTurn34}, \lean{Threefold33}, \lean{Fivefold35}, \lean{MirrorOneAxis33}, \lean{MirrorNoAxis34}, for exactly the
$S$ and $g$ of (a). The last two hold by Theorem~\ref{mb:thm:classes}. (Lean: \texttt{heptagon\_of\_nontrivial\_isometry
h2 h3 h5 mirrorOneAxis33\_proved mirrorNoAxis34\_proved S g hgp hcard hg hne hinv}.)

(b) Likewise from Theorem~\ref{ar:thm:main}(a) (\lean{heptagon_of_nontrivial_affineEquiv}) and Theorem~\ref{mb:thm:classes}.

The theorems used are the copies in RTC (lines 1023 and 1527), which are byte-identical with
\dirn{SYMLEAN/SymmetricReduction.lean:997} and \dirn{SYMLEAN/AffineReduction.lean:1501} (see the paragraph on copies at
the beginning of this section), so the written proofs of Theorems~\ref{sr:thm:main} and~\ref{ar:thm:main} apply
verbatim.
\end{proof}

\begin{remark}[status of the hypotheses of Theorem~\ref{mb:thm:main3}]\label{mb:rem:main3}
\lean{HalfTurn34}, \lean{Threefold33} and \lean{Fivefold35} are not proved in Lean. In the project they are
computer-assisted claims, supplied by the following records (\dirn{SYMLEAN/mirror/README.md}, \S4, and the table in
\dirn{SYMLEAN/README.md}, \S5). The project trusts one proof checker, \texttt{cake\_lpr}: where a proof is said below to
have been accepted by both proof checkers, \texttt{cake\_lpr} accepted it, and that is the verdict that counts;
\texttt{lrat-check} is only a secondary check, whose verdicts the project does not count (\dirn{WRITEUP/SYMMETRIC.md},
\S3; an unpatched build of it accepts invalid proofs, \dirn{LRATCHECK20261105}).
\begin{itemize}
\item \lean{HalfTurn34}: claimed in \dirn{C2FINAL}. The hull cover and the $120$ hexagonal hulls of \dirn{C2C34} have
  retained proofs; most parallelogram pieces (\dirn{C2P}--\dirn{C2P5}) rest on verdicts recorded at solve time (both
  proof checkers accepted the proof, which was then deleted), with re-solved samples and re-checks (status paragraph
  of \dirn{C2FINAL/README.md}; \dirn{C2RECHECK}). The bookkeeping of the subcase orbits, namely that the listed
  subcases of the eight canonical parallelogram hull cases $P_1,\dots,P_8$ (to which reduction~A reduces the others)
  meet every orbit of second-layer subcases, is Lean-checked according to the record (\dirn{C2FINAL/README.md},
  \S6), by kernel evaluation, so computer-assisted (Theorem~\ref{cb:thm:c2}); that the listed subcases are refuted
  rests on the records. The encoder and the two reductions of the case split have written proofs.
\item \lean{Threefold33}: \dirn{C3FINAL}. Its $116{,}883$ piece records carry verdicts recorded at solve time (both
  proof checkers accepted the proof); the proofs were not kept (\dirn{C3LEAN/README.md}, \S4), and $2{,}120$ randomly
  chosen pieces were re-solved (\dirn{C3SPOT}). \dirn{C3LEAN} Lean-checks, by kernel evaluation and so
  computer-assisted (Theorem~\ref{cb:thm:c3}), only that the leaves named in it cover the $910$ roots of the reduced
  case split (that every leaf has such a record is checked outside Lean); the normal form
  \dirn{C3-S/soundness.md}, the encoding \dirn{C3a}, the exhaustiveness of the splits and the reductions of the case
  split are written proofs.
\item \lean{Fivefold35}: \dirn{SYM5}, two normal forms of one encoder, each with a retained proof accepted by both proof
  checkers; the encoder has a written proof. The angle normal form (proof P1) needs no layer lemma; the layer normal
  form (proof L1) uses the layer lemma of \dirn{SYM5/README.md}, \S2, whose statement that every convex layer is one
  orbit, and the strict nesting of the layers, are Lean-checked (Theorems~\ref{ll:thm:layers}
  and~\ref{lg:thm:nesting}).
\end{itemize}
\lean{Threefold33} and \lean{Fivefold35} also rest on the written centre-perturbation lemma
\dirn{C2ENC/perturbation.md}. For none of the three is the step from a certificate to the Lean proposition formal: that
each class states what its certificate proves was checked by reading (\dirn{SYMLEAN/README.md}, \S6 and \S9). So the
three hypotheses are not Lean-checked, and Theorem~\ref{mb:thm:main3} reduces the statement ``every finite set of at
least $33$ points in general position with a non-identity isometric (or affine) symmetry contains $7$ points in convex
position'' to them and proves it only under them. Only the isometric form and the
\lean{AffineEquiv} form have three-class versions in RTC; the other affine variants
\lean{heptagon_of_nontrivial_affine}, \lean{heptagon_of_nontrivial_affineMap},
\lean{heptagon_of_nontrivial_affineEquiv_image}, \lean{heptagon_of_nontrivial_affine_coords} (RTC lines 1503, 1515,
1536, 1547) and \lean{halfTurn_of_symmetric_counterexample} (RTC line 1050) still take all five classes (the last one with
\lean{Fourfold32} in place of \lean{HalfTurn34}) as hypotheses.
\end{remark}

\begin{remark}[the axis hypotheses]\label{mb:rem:axis}
Theorem~\ref{mb:thm:refl32}(b) shows more than Theorem~\ref{mb:thm:classes}: for $u_0^2+u_1^2=1$, every finite
$\mathrm{reflLine}_{a,u}$-invariant set in general position with at least $33$ points, whatever the number of its points on the mirror, contains $7$ points
in convex position. In particular the proposition \lean{MirrorExclusionLineCopy} (RTC line 1078, a copy of the text of
the statement \lean{Mir.MirrorExclusionLine} of \dirn{H3}: for all $S,a,u$ with $u_0^2+u_1^2=1$, $S$ in general position,
$\card S=33$ and $\mathrm{reflLine}_{a,u}(q)\in S$ for all $q\in S$, the set $S$ contains $7$ points in convex position)
follows at once from Theorem~\ref{mb:thm:refl32}(b); RTC does not state this implication as a theorem.
\end{remark}

\paragraph{Coverage.}
Covered, by label: the cited vendored theorems \lean{main} (Theorem~\ref{mb:thm:baek}) and \lean{Config.cap_cup}
(Theorem~\ref{mb:thm:capcup}); the definitions \lean{Config}, \lean{Config.Cap3}, \lean{Config.Cap}, \lean{Config.Cup},
\lean{Config.Gon}, \lean{Config.NGon}, \lean{Config.HasNCap}, \lean{Config.HasNCup}, \lean{Config.HasNGon}
(Definition~\ref{mb:def:config}, with \texttt{List.Chain3\textquotesingle} and \lean{List.In} in Remark~\ref{mb:rem:config});
\lean{hasNCap_supset}, \lean{hasNGon_supset}, \lean{ncup_is_ngon}
(Lemma~\ref{mb:lem:config-mono}); \lean{Baek.four_cap_or_eight_gon}, \lean{Baek.capacity22} (Lemma~\ref{mb:lem:baek8});
\lean{Configuration} and \lean{Configuration.HasCap} (described at the beginning of Subsection~\ref{mb:ssec:bridge} and
stated as Definition~\ref{mm:def:config});
\lean{Configuration.cup}, \lean{colorConfig}, \lean{pointAt} (Definition~\ref{mb:def:color}); \lean{cup_true_iff},
\lean{cup_sign}, \lean{cap_sign} (Lemma~\ref{mb:lem:signs}); \lean{four_point_identity_left},
\lean{four_point_identity_right}, \lean{positive_four}, \lean{negative_four} (Lemma~\ref{mb:lem:four});
\lean{chain3_all} (Lemma~\ref{mb:lem:chain}); \lean{cup_list_all}, \lean{cap_list_all}, \lean{hasCap_of_list}
(Lemma~\ref{mb:lem:lists}); \lean{boundary_polygon}, \lean{Boundary.card}, \lean{boundary_convex_position} with the
helpers named in its proof (Lemma~\ref{mb:lem:boundary}); \lean{hasNGon_heptagon}, \lean{hasNGon_octagon},
\lean{hasNGon_nonagon} with \lean{boundaryOfGon}, \lean{boundaryOfGon_inside} (Proposition~\ref{mb:prop:gon});
\lean{capacity16} (Proposition~\ref{mb:prop:cap16}); \lean{geometric_capacity22} (Lemma~\ref{mb:lem:gc22});
\lean{capacity22g} (Proposition~\ref{mb:prop:cap22}); \lean{config_capacity113} (Theorem~\ref{mb:thm:nine}(a));
\lean{capacity113} (Proposition~\ref{mb:prop:cap113}); \lean{mir}, \lean{Rset}, \lean{Aset}
(Definition~\ref{mb:def:mir}); \lean{card_split}, \lean{card_A_le}, \lean{injOn_x} (Lemma~\ref{mb:lem:split});
\lean{cfg}, \lean{sortPt_image} (Lemma~\ref{mb:lem:sort}); \lean{cap_mirror_heptagon}, \lean{cap_mirror_octagon},
\lean{axis_left_heptagon}, \lean{axis_right_heptagon} (Lemma~\ref{mb:lem:polygons}, by a different proof);
\lean{mirror_seven_of} (Lemma~\ref{mb:lem:seven-of}); \lean{lineRefl}, \lean{frame}, \lean{unframe}
(Definition~\ref{mb:def:frame}); \lean{frame_lineRefl}, \lean{orient_frame}, \lean{unframe_frame}, \lean{frame_unframe},
\lean{frame_injective}, \lean{unframe_injective} (Lemma~\ref{mb:lem:frame}); \lean{hasConvexNGon_affineMap}
(Lemma~\ref{mb:lem:affine}); \lean{transport} (Lemma~\ref{mb:lem:transport}); \lean{mirror_seven},
\lean{mirror_seven_line} (Theorem~\ref{mb:thm:seven}); \lean{mirror_eight}, \lean{mirror_eight_line}
(Theorem~\ref{mb:thm:eight}); \lean{mirror_nine}, \lean{mirror_nine_line} (Theorem~\ref{mb:thm:nine}(b),(c)); \lean{reflLine}
(Definition~\ref{mb:def:reflline}); \lean{reflLine_eq_lineRefl}, \lean{card_le_32_of_reflLine}
(Theorem~\ref{mb:thm:refl32}); \lean{mirrorOneAxis33_proved}, \lean{mirrorNoAxis34_proved}
(Theorem~\ref{mb:thm:classes}); \lean{heptagon_of_nontrivial_isometry_3}, \lean{heptagon_of_nontrivial_affineEquiv_3}
(Theorem~\ref{mb:thm:main3}, conditional on \lean{HalfTurn34}, \lean{Threefold33} and \lean{Fivefold35}). The
declarations stated in Subsection~\ref{mb:ssec:red}, in
Definition~\ref{mb:def:frame} and in Lemmas~\ref{mb:lem:frame}--\ref{mb:lem:transport} are also stated in the section
of Theorem~\ref{mm:thm:odd} (see the beginning of Subsection~\ref{mb:ssec:red}).

Not covered here, deliberately:
\begin{itemize}
\item the Lean proofs of \lean{main} and \lean{Config.cap_cup} (cited published theorems);
\item the generated Lean proofs of the four polygon lemmas and the wrappers \lean{cyc7}, \lean{cyc8} that they use
  (replaced by the derivation from Theorems~\ref{mm:thm:mirror-polygon} and~\ref{mm:thm:axis-polygon},
  Remark~\ref{mb:rem:generated}), and the rest of the generated block (the wrappers \lean{cyc4}--\lean{cyc6} and the
  smaller mirror polygons, which serve the cases $k=4,5,6$);
\item \texttt{mirror\_seven\textquotesingle} and \texttt{mirror\_eight\textquotesingle} (consistency duplicates,
  Remarks~\ref{mb:rem:seven} and~\ref{mb:rem:eight});
\item the general mirror lemma \lean{cap_mirror_polygon}, \lean{axis_left_polygon}, \lean{axis_right_polygon},
  \lean{cap_below_lines}, the general reductions \lean{mirror_odd_general}, \lean{mirror_even_general} and their
  \lean{_line} versions, and \lean{mirror_four}, \lean{mirror_five}, \lean{mirror_six} with their \lean{_line} versions
  and their helpers \lean{hasCap3_of}, \lean{capacity_cup} (they are the subject of Lemma~\ref{mm:lem:slope} and of
  Theorems~\ref{mm:thm:mirror-polygon}, \ref{mm:thm:axis-polygon}, \ref{mm:thm:odd}, \ref{mm:thm:even},
  \ref{mm:thm:line} and~\ref{mm:thm:small}, or are used in their proofs; Part~B of RTC is \dirn{MirrorMain.lean}
  verbatim after its two import lines);
\item the Part~A content of RTC (the subject of Definitions~\ref{sr:def:iso}, \ref{sr:def:classes} and
  Theorems~\ref{sr:thm:main}, \ref{ar:thm:main} and the statements around them);
\item helpers absorbed into the proofs above, among them: \lean{mir_zero}, \lean{mir_one}, \lean{mir_mir}, \lean{xset},
  \lean{xset_card}, \lean{emb}, \lean{emb_mem}, \lean{sortPt}, \lean{sortPt_mem}, \lean{sortPt_x},
  \lean{sortPt_strictMono}, \lean{sortPt_injective}, \lean{unframeLin}, \lean{unframe_apply}; \lean{pointAt_eq},
  \lean{pointAt_order}, \lean{OrderedTriples}, \lean{arcOfList}, \lean{cupArc}, \lean{capArc}, \lean{same_first},
  \lean{same_last}; from \dirn{BoundaryGeometry.lean}, besides the lemmas named in the proof of
  Lemma~\ref{mb:lem:boundary}, \lean{CupArc}, \lean{first}, \lean{last}, \lean{reflect}, \lean{reflect_reflect},
  \lean{orient_reflect}, \lean{Boundary}, \lean{Boundary.points}, \lean{Boundary.flip}, \lean{reflect_mem_flip},
  \lean{Boundary.finite}, \lean{Boundary.finite_coe}, \lean{affine}, \lean{affine_mix}, \lean{affine_positive_convex},
  \lean{supporting_edges_expose}, \lean{point_x_bounds}, \lean{first_x_strict}, \lean{last_x_strict}; and, from the
  vendored library, \lean{Config.NCap}, \lean{Config.NCup} (caps and cups of a given length), the decidability instance
  \lean{Config.DecidableCap3}, and list lemmas such as \lean{List.takeHead}, \lean{List.takeHead2},
  \lean{List.takeLast2}, \lean{List.in_superset}, \lean{List.cons_in};
\item the other declarations of the bridge files, which these proofs do not use: \lean{Configuration.cup_false_iff},
  \lean{hasCup_of_list}, \lean{hasNCup_octagon}, \lean{chain3_of_all}, \lean{listCup_of_hasCup},
  \lean{listCap_of_hasCap}, \lean{PlanarChains.cup_all}, and the declarations of \dirn{Realization.lean} and
  \dirn{ActualRanks.lean} not named in this section (they belong to the project's \dirn{ErdosSzekeres/Omission}
  development and are outside this section);
\item the \texttt{\#print axioms} and \texttt{\#check} lines.
\end{itemize}

\section{Case-split bookkeeping: the threefold cover and the central 34-point split}\label{cb:sec}

\subsection{Overview}\label{cb:sec:overview}

This section is about two Lean files that check the \emph{bookkeeping} of two case splits of the project, that is, the
purely combinatorial statement that an explicitly named family of nodes covers the split.

\begin{enumerate}[label=(\arabic*)]
\item The file \dirn{C3LEAN/C3Cover.lean} models the split of threefold rotationally symmetric sets of 33 points
  (11 orbits of the rotation by $120^\circ$, indexed by the residues $0,\dots,10$) by their convex layers; such sets
  are the subject of the case class \lean{Threefold33} of Definition~\ref{sr:def:classes}. Its headline is
  Theorem~\ref{cb:thm:c3} (\lean{C3Cover.all_roots_covered}): each of the 910 depth-two \emph{roots} of the reduced
  split is \emph{covered}, through layer splits and $\sigma$-splits, by the nodes named by the keys of 910 generated
  certificates. The proof combines a soundness lemma for a certificate checker (Lemma~\ref{cb:lem:run}, with a
  complete written proof here) with a kernel evaluation of the checker on the 910 certificates
  (Fact~\ref{cb:fact:batches}), the computer-assisted match of the codes of the roots with the roots of the
  certificates (Fact~\ref{cb:fact:roots-codes}), and the decoding of each root from its code
  (Fact~\ref{cb:fact:roots-match}, which also has a written proof here).
\item The file \dirn{C2FINAL/Central34Orbits.lean} models the second-layer split of the parallelogram hull cases of
  centrally symmetric sets of 34 points (17 antipodal classes $0,\dots,16$; the case class \lean{HalfTurn34} of
  Definition~\ref{sr:def:classes}). Its headline is Theorem~\ref{cb:thm:c2} (\lean{Central34.all_cases_closed}): for
  each hull case $P_a$, $1\le a\le 8$, each of the 560 second-layer subcases, or its image under the reflection that
  sends $i$ to the residue of $a-i$ modulo $17$, belongs (as a set) to an explicit generated list.
\end{enumerate}

\paragraph{What is not established here.} Both headline theorems are statements about finite lists of natural numbers.
They mention no points, no convexity, no CNF formula and no SAT solver, and they take no hypotheses. The following are
\emph{not} established by them, and are not proved in this document:
\begin{enumerate}[label=(\roman*)]
\item the geometric meaning of the splits: that every set of the class lies below some root (the top cases, and the
  reduction of the second layer by Lemma~5 and Lemma~H6M of the project's written arguments \dirn{C3-S/soundness.md}
  and \dirn{C3H6M/soundness_h6m.md}; for the central split, the reductions A and B reviewed in \dirn{C2P/review}),
  and that a layer split or a $\sigma$-split loses no set (Lemma~4 of \dirn{C3-S/soundness.md}, rules~2 and~3 of
  \dirn{C3U2/sigma_split.md}; the second-layer covers of the central split);
\item the CNF encodings of the nodes;
\item the SAT verdicts of the leaves, and even the fact that the leaves are nodes for which the project holds a run
  record in which both proof checkers, \texttt{cake\_lpr} and \texttt{lrat-check}, are recorded as accepting the proof
  (the project's term for such a record: ``certified''). In Theorem~\ref{cb:thm:c3} the set of leaves is, by
  definition, the set of nodes named by the keys of the certificates themselves, and in Theorem~\ref{cb:thm:c2} the
  lists $\mathsf{closedP}_a$ are just
  lists. Which nodes were put there was decided outside Lean from the run records (by
  \dirn{C3LEAN/tools/gen_c3lean.py}, audited by \dirn{C3LEAN/tools/audit_keys.py} and
  \dirn{C3LEAN/review/indep_audit.py}; by \dirn{C2FINAL/gen_lean.py} with \dirn{C2FINAL/hashmatch.py}).
\end{enumerate}
So neither theorem implies \lean{Threefold33} or \lean{HalfTurn34}. They check one step, ``the named nodes form a cut
of the split'', which the project combines with (i)--(iii); see Remarks~\ref{cb:rem:use} and~\ref{cb:rem:c2-scope}.

\paragraph{Sources and toolchain.} Paths are relative to the results directory \dirn{cloud/es7-2026-10-19/results}
of the repository.
\begin{itemize}
\item \dirn{C3LEAN/C3Cover.lean} (2{,}327 lines, SHA-256 \texttt{4b64293df2ea}\dots). Its lines 29--329 (the
  definitions of the model and of the checker, and the conceptual proofs) are byte-identical to lines 1--301 of
  \dirn{C3LEAN/tools/prelude.lean} (SHA-256 \texttt{cdcf96de7739}\dots), from which the generator
  \dirn{C3LEAN/tools/gen_c3lean.py} copies them. Lines 330--2299 hold the auxiliary definitions \lean{beqLayer},
  \lean{isOption}, \lean{distinctCodes}, \lean{codeSet} and \lean{reflectionFacts} (lines 335--365), the facts about
  the roots (lines 367--409 and 1720--1733), the generated data (lines 413, 423--1717 and 1737), the batch theorems
  (lines 1740--2245), the theorem \lean{all_checked} (line 2247), the main theorem (line 2251) and seven negative
  controls (lines 2266--2297); lines 2301--2327 are \texttt{\#print axioms} commands.
\item \dirn{C2FINAL/Central34Orbits.lean} (298 lines, SHA-256 \texttt{b396a26b2430}\dots), written by
  \dirn{C2FINAL/gen_lean.py}.
\end{itemize}
Both files are checked with Lean~4.24.0 according to their READMEs; for \dirn{C3LEAN/C3Cover.lean} a compile log is
saved, for \dirn{C2FINAL/Central34Orbits.lean} no Lean output (paragraph ``Axioms'' below). Neither has an \texttt{import}: only the core library is used, and the
repository's Mathlib (pinned at \texttt{f897ebcf}) plays no role. Every finite computation is closed by the tactic
\texttt{decide +kernel}: the Lean kernel itself evaluates the Boolean expression to \texttt{true} by definitional
unfolding, using its built-in arbitrary-precision arithmetic on natural-number literals (part of the kernel, not an
axiom). No \texttt{native\_decide} is used, so the axiom \lean{Lean.ofReduceBool} does not occur. The option
\texttt{Elab.async false} set at line~27 of \dirn{C3LEAN/C3Cover.lean} only affects elaboration.

\paragraph{Axioms.} For \dirn{C3LEAN/C3Cover.lean}, the compile log \dirn{C3LEAN/logs/compile.log} (and the identical
axiom lines of the review's recompilation, \dirn{C3LEAN/review/logs/main_compile.log}) records: \lean{all_roots_covered},
\lean{run_sound} and \lean{checkRoot_sound} depend on \lean{propext} and \lean{Quot.sound}; the eleven theorems
\lean{reflection_facts_H3}, \lean{reflection_facts_H6_1_0}, \dots, \lean{reflection_facts_H6_5_1}, and
\lean{roots_eq_tsv} and \lean{all_checked} on \lean{propext} only; \lean{topCases_eq}, \lean{strata},
\lean{roots_match} and \lean{roots_codes_match} on no axiom. The lemmas \lean{Covered.mono}, \lean{beqList_eq},
\lean{beqTok_eq}, \lean{beqNode_eq}, \lean{all_append_of} and the 127 batch theorems are not printed separately; each is
used by a printed theorem (\lean{run_sound}, \lean{all_checked} or \lean{all_roots_covered}), so its axioms are among
\lean{propext} and \lean{Quot.sound}. No theorem of the file depends on \lean{Classical.choice}, \lean{sorryAx} or
\lean{Lean.ofReduceBool}. For \dirn{C2FINAL/Central34Orbits.lean} no compile or axioms log is committed: the file ends
with three \texttt{\#print axioms} commands (lines 296--298) whose output is not in the repository, and
\dirn{C2FINAL/README.md} (\S6) records ``axioms \texttt{propext} only, about 6 s''. As no Lean output of any kind
(compile or build log, axiom output, pinned message) is saved for this file, the labels of its declarations rest on
that record alone: its three theorems and its definitions are \emph{Lean-checked according to the record}, the
qualified label defined with the rigor labels at the beginning of this document.

\paragraph{Rigor labels and dependencies.} The Lean theorems of \dirn{C3LEAN/C3Cover.lean} keep the label
\emph{Lean-checked}, and the three theorems of \dirn{C2FINAL/Central34Orbits.lean} (Theorem~\ref{cb:thm:c2-refl},
Fact~\ref{cb:fact:c2-sub} and Theorem~\ref{cb:thm:c2}) carry the label \emph{Lean-checked according to the record},
as do its definitions (paragraph ``Axioms'' above); the proofs written
here are \emph{written proofs (not refereed)}. Every statement that Lean proves by kernel evaluation is marked
\emph{computer-assisted} in its statement, and so is every statement whose proof uses one of the computer-assisted
statements on generated data (Lemma~\ref{cb:lem:all-checked}, Corollary~\ref{cb:cor:roots},
Theorem~\ref{cb:thm:c3} and Corollary~\ref{cb:cor:c2}). For the kernel evaluations on explicit generated data
(Facts~\ref{cb:fact:batches}, \ref{cb:fact:roots-codes}, \ref{cb:fact:roots-tsv} and~\ref{cb:fact:c2-sub}, and
Theorem~\ref{cb:thm:c2}) we state exactly what is evaluated, and we give no hand proof. For the kernel evaluations of
the definitions alone (Facts~\ref{cb:fact:refl-h3} to~\ref{cb:fact:roots-match} and Theorem~\ref{cb:thm:c2-refl}) we
also give complete written proofs that do not use the computation; the statement names the written proof. Each Lean
declaration of the two files other than the seven negative controls of \dirn{C3LEAN/C3Cover.lean} (tests, listed in
\S\ref{cb:sec:coverage}), or each uniform family of generated declarations (\lean{certs0}, \dots, \lean{certs126};
\lean{batch0}, \dots, \lean{batch126}; the ten theorems \lean{reflection_facts_H6_1_0}, \dots,
\lean{reflection_facts_H6_5_1}; \lean{closedP1}, \dots, \lean{closedP8}), has a statement below whose head carries
its Lean reference; for a family the head cites the first member, and the statement locates the others.
The head of Lemma~\ref{cb:lem:run} carries the two declarations \lean{run_sound} and \lean{checkRoot_sound}. The
proofs of this section use no other section.

\subsection{Notation}\label{cb:sec:notation}

\paragraph{Numbers.} $\N=\{0,1,2,\dots\}$. For $x,y\in\N$ with $y>0$, $\lfloor x/y\rfloor$ and $x\bmod y$ are the
quotient and the remainder of Euclidean division (Lean's division and remainder on $\N$), and
$x\mathbin{\dot-}y=\max(x-y,0)$ is Lean's truncated subtraction on $\N$. For an integer $x$ and $y\in\N$ with
$y>0$, $x\bmod y$ denotes the least nonnegative residue of $x$ modulo $y$; for $x\in\N$ this is the remainder above.
For $\sigma,f\in\{0,1\}$,
$\sigma\oplus f=(\sigma+f)\bmod 2$.

\paragraph{Lists.} We write $[x_0,\dots,x_{k-1}]$ for a list, $[]$ for the empty list, $x\mathbin{::}l$ for the list
with head $x$ and tail $l$, $l_1\mathbin{+\!\!+}l_2$ for concatenation, $\card{l}$ for the length, and $x\in l$ if $x$
is an entry of $l$; $\mathrm{set}(l)$ is the set of entries of $l$. A list is \emph{duplicate-free} if entries at
different positions are different; then $\card{l}=\card{\mathrm{set}(l)}$. The list functions are those of Lean's core
library: $\mathsf{range}(n)=[0,1,\dots,n-1]$; $\mathsf{map}(g,l)=[g(x_0),\dots,g(x_{k-1})]$;
$\mathsf{flatMap}(g,l)=g(x_0)\mathbin{+\!\!+}\cdots\mathbin{+\!\!+}g(x_{k-1})$; $\mathsf{filter}(p,l)$ is the list of
the entries $x$ of $l$ with $p(x)=\mathsf{true}$, in their order; $\mathsf{take}(K,l)$ is the list of the first
$\min(K,\card l)$ entries; $\mathsf{zip}(l,l')=[(x_0,x'_0),(x_1,x'_1),\dots]$ has length $\min(\card l,\card{l'})$;
$\mathsf{foldl}(g,z,l)=g(\cdots g(g(z,x_0),x_1)\cdots,x_{k-1})$; $\mathsf{replicate}(k,x)$ is the list of $k$ copies of
$x$. The Boolean-valued functions are: $\mathsf{contains}(l,x)=\mathsf{true}$ iff $x\in l$;
$\mathsf{all}(l,p)=\mathsf{true}$ iff $p(x)=\mathsf{true}$ for every $x\in l$; $\mathsf{any}(l,p)=\mathsf{true}$ iff
$p(x)=\mathsf{true}$ for some $x\in l$; $\mathsf{isEmpty}(l)=\mathsf{true}$ iff $l=[]$. We use without comment the
following facts, each immediate by induction on $l$ from the defining equations:
\begin{enumerate}[label=(L\arabic*)]
\item $\mathsf{all}(l_1\mathbin{+\!\!+}l_2,p)=\mathsf{all}(l_1,p)\wedge\mathsf{all}(l_2,p)$;
\item $y\in\mathsf{flatMap}(g,l)$ iff $y\in g(x)$ for some $x\in l$;
\item $\card{\mathsf{map}(g,l)}=\card l$, and $\card{\mathsf{flatMap}(g,l)}$ is the sum of $\card{g(x)}$ over the
  positions of $l$;
\item $y\in\mathsf{map}(g,l)$ iff $y=g(x)$ for some $x\in l$.
\end{enumerate}

\paragraph{Booleans.} The Booleans are $\mathsf{true}$ and $\mathsf{false}$. In Boolean expressions $\wedge$, $\vee$,
$\neg$ denote Lean's Boolean conjunction, disjunction and negation, and a decidable proposition $P$ used where a Boolean
is expected stands for $\mathsf{decide}(P)$, the Boolean that is $\mathsf{true}$ iff $P$ holds (Lean inserts this
coercion). For $x,y\in\N$, Lean's Boolean equality tests (written \texttt{x == y} or \texttt{Nat.beq x y}) are
$\mathsf{true}$ iff $x=y$, and \texttt{x != y} is the negation of \texttt{x == y}. A statement ``$b=\mathsf{true}$'' is a
Lean proposition about a Boolean $b$.

\paragraph{Binary digits.} For $n\in\N$ let $\beta(n)=\{i\in\N:\lfloor n/2^i\rfloor\text{ is odd}\}$, the set of
positions of the digit $1$ in the binary expansion of $n$. Then $n=\sum_{i\in\beta(n)}2^i$, and $n\mapsto\beta(n)$ is a
bijection from $\N$ onto the set of finite subsets of $\N$ (uniqueness of the binary expansion). In Lean, shifting $1$
left by $c$ places gives $2^c$, so $\beta(2^c)=\{c\}$; the bitwise or of $x$ and $y$, written here
$x\mathbin{\mathsf{or}}y$, is the number with $\beta(x\mathbin{\mathsf{or}}y)=\beta(x)\cup\beta(y)$; and
$\mathsf{testBit}(n,i)$ (Lean's \texttt{Nat.testBit n i}) is $\mathsf{true}$ iff $i\in\beta(n)$.

\begin{lemma}[Or-folds; written step]\label{cb:lem:bits}
Let $z\in\N$, let $l$ be a list of elements of a set $X$, and let $g:X\to\N$. Then
\[
\beta\big(\mathsf{foldl}((acc,x)\mapsto acc\mathbin{\mathsf{or}}2^{g(x)},\ z,\ l)\big)=\beta(z)\cup\{g(x):x\in l\}.
\]
\end{lemma}

\begin{proof}
Induction on $l$, for all $z$ at once. For $l=[]$ the fold is $z$. For $l=x\mathbin{::}l'$ the fold equals
$\mathsf{foldl}((acc,y)\mapsto acc\mathbin{\mathsf{or}}2^{g(y)},\ z\mathbin{\mathsf{or}}2^{g(x)},\ l')$, whose set of
binary digits is, by induction, $\beta(z)\cup\{g(x)\}\cup\{g(y):y\in l'\}$.
\end{proof}

\subsection{The threefold split: the model}\label{cb:sec:model}

Every Lean reference in a statement head of \S\S\ref{cb:sec:model}--\ref{cb:sec:c3} is to
\dirn{C3LEAN/C3Cover.lean}, namespace \lean{C3Cover}; the statements tagged ``written step'' have no Lean
counterpart. The definitions transcribe the project's Python implementation of the split
(\dirn{C3U2/tools/hullcase.py}); that the transcription agrees with the Python code is tested in the project
(\dirn{C3LEAN/README.md}, \S3), not proved, and plays no role in the theorems.

\begin{definition}[Orbits; \leanref{C3Cover.m}{C3LEAN/C3Cover.lean:34}]\label{cb:def:m}
$m=11$. Below we write $11$ for $m$ wherever it occurs.
\end{definition}

\begin{definition}[Tokens; \leanref{C3Cover.Tok}{C3LEAN/C3Cover.lean:38}]\label{cb:def:tok}
A \emph{token} is a pair $t=(r,s)\in\N\times\N$; we write $t.1=r$ and $t.2=s$.
\end{definition}

\begin{definition}[Layers; \leanref{C3Cover.Layer}{C3LEAN/C3Cover.lean:41}]\label{cb:def:layer}
A \emph{layer} is a list of tokens.
\end{definition}

\begin{definition}[Nodes; \leanref{C3Cover.Node}{C3LEAN/C3Cover.lean:45}]\label{cb:def:node}
A \emph{node} $n=\langle\Lambda,C\rangle$ is a pair consisting of a list $n.\mathsf{layers}=\Lambda$ of layers and a
list $n.\mathsf{cube}=C$ of pairs $(r,v)\in\N\times\N$. Two nodes are equal iff both components are equal.
\end{definition}

\begin{definition}[Used residues; \leanref{C3Cover.Node.used}{C3LEAN/C3Cover.lean:51}]\label{cb:def:used}
For a node $n$,
\[
\mathsf{used}(n)=\mathsf{flatMap}\big(L\mapsto\mathsf{map}(t\mapsto t.1,\ L),\ n.\mathsf{layers}\big),
\]
the first components of the tokens of all layers, in order.
\end{definition}

\begin{remark}[Intended reading]\label{cb:rem:reading}
In the project, a token $(r,s)$ stands for the orbit with residue $r\in\{0,\dots,10\}$ together with its $\sigma$
value $s\in\{0,1\}$ (the side of a fixed line through the centre on which the point with label $r$ lies), and $s=2$
stands for ``no $\sigma$'', which is used for a layer of one orbit other than the top layer (the one-orbit top
layer $[(0,0)]$ carries $s=0$). A layer of one or two tokens is a convex layer of
one or two orbits; the layers of a node are listed from the convex hull inwards, and the cube fixes the $\sigma$ values
$v$ of some further residues $r$. The top layer is $\mathsf{H3}=[(0,0)]$ (a hull of one orbit) or
$\mathsf{H6}(a,s)=[(0,0),(a,s)]$ with $1\le a\le5$ and $s\in\{0,1\}$ (a hull of the orbits $0$ and $a$, with
$\sigma_a=s$); see Lemma~2 of \dirn{C3-S/soundness.md}. Nothing in the Lean types restricts $r\le10$ or $s\le2$: all such
bounds used below are consequences of how particular lists are built. This reading is not used in any proof.
\end{remark}

\begin{definition}[Pair options; \leanref{C3Cover.pairOptions}{C3LEAN/C3Cover.lean:54}]\label{cb:def:pairoptions}
For a list $l$ of naturals,
\[
\mathsf{pairOptions}([])=[],\qquad
\mathsf{pairOptions}(b\mathbin{::}l)=\mathsf{flatMap}\big(c\mapsto Q_{b,c},\ l\big)\mathbin{+\!\!+}\mathsf{pairOptions}(l),
\]
where $Q_{b,c}=\big[\,[(b,0),(c,0)],\ [(b,0),(c,1)],\ [(b,1),(c,0)],\ [(b,1),(c,1)]\,\big]$.
\end{definition}

\begin{definition}[Layer options; \leanref{C3Cover.layerOptions}{C3LEAN/C3Cover.lean:61}]\label{cb:def:layeroptions}
For a list $U$ of naturals let $\mathsf{rem}(U)=\mathsf{filter}\big(r\mapsto\neg\mathsf{contains}(U,r),\
\mathsf{range}(11)\big)$ (a local abbreviation inside \lean{layerOptions}), and
\[
\mathsf{layerOptions}(U)=\mathsf{map}\big(r\mapsto[(r,2)],\ \mathsf{rem}(U)\big)\mathbin{+\!\!+}
\mathsf{pairOptions}(\mathsf{rem}(U)).
\]
\end{definition}

\begin{definition}[Insertion; \leanref{C3Cover.insTok}{C3LEAN/C3Cover.lean:66}]\label{cb:def:instok}
For a pair $t\in\N\times\N$ and a list of such pairs: $\mathsf{insTok}(t,[])=[t]$;
$\mathsf{insTok}(t,u\mathbin{::}us)=t\mathbin{::}u\mathbin{::}us$ if $t.1\le u.1$, and
$\mathsf{insTok}(t,u\mathbin{::}us)=u\mathbin{::}\mathsf{insTok}(t,us)$ if $t.1>u.1$.
\end{definition}

\begin{definition}[Sorting; \leanref{C3Cover.sortTok}{C3LEAN/C3Cover.lean:71}]\label{cb:def:sorttok}
$\mathsf{sortTok}([])=[]$ and $\mathsf{sortTok}(t\mathbin{::}ts)=\mathsf{insTok}(t,\mathsf{sortTok}(ts))$.
\end{definition}

\begin{definition}[Top cases; \leanref{C3Cover.topCases}{C3LEAN/C3Cover.lean:77}]\label{cb:def:topcases}
The list of top cases is
\[
\mathsf{topCases}=[(0,0)]\mathbin{::}\mathsf{flatMap}\big(i\mapsto\big[\,[(0,0),(i+1,0)],\ [(0,0),(i+1,1)]\,\big],\
\mathsf{range}(5)\big).
\]
We write $\mathsf{H3}=[(0,0)]$ and $\mathsf{H6}(a,s)=[(0,0),(a,s)]$.
\end{definition}

\begin{definition}[Option keys; \leanref{C3Cover.okey}{C3LEAN/C3Cover.lean:82}]\label{cb:def:okey}
For a layer $L$, $\mathsf{okey}(L)=\mathsf{map}(\kappa,\mathsf{sortTok}(L))$, where $\kappa(r,s)=3r$ if $s=2$ and
$\kappa(r,s)=3r+s+1$ if $s\ne2$.
\end{definition}

\begin{definition}[Lexicographic order; \leanref{C3Cover.lexLe}{C3LEAN/C3Cover.lean:85}]\label{cb:def:lexle}
For lists of naturals: $\mathsf{lexLe}([],y)=\mathsf{true}$; $\mathsf{lexLe}(x\mathbin{::}xs,[])=\mathsf{false}$; and
$\mathsf{lexLe}(x\mathbin{::}xs,\,y\mathbin{::}ys)=(x<y)\vee\big((x=y)\wedge\mathsf{lexLe}(xs,ys)\big)$.
\end{definition}

\begin{definition}[Mirror; \leanref{C3Cover.mirror}{C3LEAN/C3Cover.lean:91}]\label{cb:def:mirror}
For a layer $L$,
\[
\mathsf{mirror}(L)=\mathsf{sortTok}(\mathsf{map}(\mu,L)),\qquad
\mu(r,s)=\big((11\mathbin{\dot-}r)\bmod11,\ s'\big),\quad s'=\begin{cases}2&\text{if }s=2,\\
1\mathbin{\dot-}s&\text{if }s\ne2.\end{cases}
\]
\end{definition}

\begin{definition}[The map phi; \leanref{C3Cover.phi}{C3LEAN/C3Cover.lean:95}]\label{cb:def:phi}
For $a,s\in\N$ and a layer $L$,
\[
\mathsf{phi}(a,s,L)=\mathsf{sortTok}(\mathsf{map}(\varphi_{a,s},L)),\qquad
\varphi_{a,s}(r,\sigma)=\big(((a+11)\mathbin{\dot-}r)\bmod11,\ \sigma'\big),
\]
where $\sigma'=2$ if $\sigma=2$; $\sigma'=1\mathbin{\dot-}\sigma$ if $\sigma\ne2$ and exactly one of the conditions
$r<a$, $s=0$ holds; and $\sigma'=\sigma$ if $\sigma\neq2$ and both or neither of them hold.
\end{definition}

\begin{definition}[Reflection; \leanref{C3Cover.refl}{C3LEAN/C3Cover.lean:101}]\label{cb:def:refl}
For layers $top$ and $o$: $\mathsf{refl}(top,o)=\mathsf{mirror}(o)$ if $top=[(0,0)]$;
$\mathsf{refl}(top,o)=\mathsf{phi}(a,s,o)$ if $top=[(0,0),(a,s)]$ for some $a,s\in\N$; and
$\mathsf{refl}(top,o)=o$ for every other list $top$.
\end{definition}

\begin{definition}[Kept options; \leanref{C3Cover.keptOptions}{C3LEAN/C3Cover.lean:109}]\label{cb:def:kept}
For a layer $top$,
\begin{align*}
\mathsf{keptOptions}(top)&=\mathsf{filter}\big(\mathsf{kept}_{top},\ \mathsf{layerOptions}(\mathsf{map}(u\mapsto
u.1,\ top))\big),\quad\text{where}\\
\mathsf{kept}_{top}(o)&=\mathsf{lexLe}\big(\mathsf{okey}(o),\ \mathsf{okey}(\mathsf{refl}(top,o))\big).
\end{align*}
\end{definition}

\begin{definition}[Roots; \leanref{C3Cover.roots}{C3LEAN/C3Cover.lean:113}]\label{cb:def:roots}
The list of roots is
\[
\mathsf{roots}=\mathsf{flatMap}\big(t\mapsto\mathsf{map}(o\mapsto\langle[t,o],[]\rangle,\ \mathsf{keptOptions}(t)),\
\mathsf{topCases}\big).
\]
\end{definition}

\begin{definition}[Free residues; \leanref{C3Cover.Node.free}{C3LEAN/C3Cover.lean:116}]\label{cb:def:free}
For a node $n$,
\[
\mathsf{free}(n)=\mathsf{filter}\big(r\mapsto(r\ne0)\wedge\neg\mathsf{contains}(\mathsf{used}(n),r)\wedge
\neg\mathsf{contains}(\mathsf{map}(c\mapsto c.1,\ n.\mathsf{cube}),r),\ \mathsf{range}(11)\big).
\]
\end{definition}

\begin{definition}[Bits; \leanref{C3Cover.bitsOf}{C3LEAN/C3Cover.lean:120}]\label{cb:def:bitsof}
For $K,i\in\N$,
\[
\mathsf{bitsOf}(K,i)=\mathsf{map}\big(j\mapsto\lfloor i/2^{(K\mathbin{\dot-}1)\mathbin{\dot-}j}\rfloor\bmod2,\
\mathsf{range}(K)\big).
\]
\end{definition}

\begin{definition}[Adding bits; \leanref{C3Cover.addBits}{C3LEAN/C3Cover.lean:123}]\label{cb:def:addbits}
For lists $C$ and $new$ of pairs of naturals,
\[
\mathsf{addBits}(C,new)=\mathsf{foldl}\big((c,t)\mapsto\mathsf{insTok}(t,c),\ C,\ new\big).
\]
\end{definition}

\begin{definition}[Sigma children; \leanref{C3Cover.Node.sigmaChildren}{C3LEAN/C3Cover.lean:127}]\label{cb:def:sigmachildren}
For a node $n$ and $K\in\N$, with the local abbreviation $bits=\mathsf{take}(K,\mathsf{free}(n))$,
\[
\mathsf{sigmaChildren}(n,K)=\mathsf{map}\Big(i\mapsto\big\langle n.\mathsf{layers},\
\mathsf{addBits}\big(n.\mathsf{cube},\ \mathsf{zip}(bits,\mathsf{bitsOf}(K,i))\big)\big\rangle,\
\mathsf{range}(2^K)\Big).
\]
\end{definition}

\begin{definition}[Conflicts; \leanref{C3Cover.conflicts}{C3LEAN/C3Cover.lean:132}]\label{cb:def:conflicts}
For a layer $L$ and a list $C$ of pairs of naturals,
\[
\mathsf{conflicts}(L,C)=\mathsf{any}\big(L,\ t\mapsto(t.2\ne2)\wedge\mathsf{any}(C,\ c\mapsto(c.1=t.1)\wedge(c.2\ne
t.2))\big).
\]
\end{definition}

\begin{definition}[Layer children; \leanref{C3Cover.Node.layerChildren}{C3LEAN/C3Cover.lean:137}]\label{cb:def:layerchildren}
For a node $n$,
\[
\mathsf{layerChildren}(n)=\mathsf{map}\Big(o\mapsto\langle n.\mathsf{layers}\mathbin{+\!\!+}[o],\
n.\mathsf{cube}\rangle,\ \mathsf{filter}\big(o\mapsto\neg\mathsf{conflicts}(o,n.\mathsf{cube}),\
\mathsf{layerOptions}(\mathsf{used}(n))\big)\Big).
\]
\end{definition}

\begin{remark}[Reading of the children]\label{cb:rem:children}
If $K\le\card{\mathsf{free}(n)}$ and $b_0<\dots<b_{K-1}$ are the first $K$ entries of $\mathsf{free}(n)$ (the $K$
smallest residues in $\{1,\dots,10\}$ that occur neither in a layer nor in the cube of $n$), then the $i$-th
$\sigma$-child, $0\le i<2^K$, has the layers of $n$ and the cube of $n$ with the pairs $(b_j,v_j)$ inserted, where
$v_j=\lfloor i/2^{K-1-j}\rfloor\bmod2$ is the binary digit of $i$ at position $K-1-j$; as $i$ runs through
$0,\dots,2^K-1$, $(v_0,\dots,v_{K-1})$ runs through $\{0,1\}^K$. A layer child of $n$ appends to the layers of $n$ one
layer option on the residues not used by $n$, except the options that give some residue of the cube an explicit
$\sigma$ value (not $2$) different from the cube's value. No reduction by a reflection is applied to layer children.
This reading is not used in any proof.
\end{remark}

\subsection{The cover predicate}\label{cb:sec:covered}

\begin{definition}[Covered nodes; \leanref{C3Cover.Covered}{C3LEAN/C3Cover.lean:146}]\label{cb:def:covered}
Let $S$ be a predicate on nodes. A predicate $P$ on nodes is \emph{$S$-closed} if for every node $n$ the following
three rules hold.
\begin{description}
\item[(cert)] If $S(n)$, then $P(n)$.
\item[(layer)] If $2\le\card{n.\mathsf{layers}}$, $\mathsf{layerChildren}(n)\ne[]$, and $P(c)$ for every
  $c\in\mathsf{layerChildren}(n)$, then $P(n)$.
\item[(sigma)] If $K\in\N$, $2\le\card{n.\mathsf{layers}}$, $0<K$, $K\le\card{\mathsf{free}(n)}$, and $P(c)$ for every
  $c\in\mathsf{sigmaChildren}(n,K)$, then $P(n)$.
\end{description}
$\mathsf{Covered}(S)$ is the least $S$-closed predicate: $\mathsf{Covered}(S)(n)$ holds iff $P(n)$ holds for every
$S$-closed predicate $P$.
\end{definition}

In each rule the predicate occurs in the premises only positively, so an intersection of $S$-closed predicates is
$S$-closed; hence $\mathsf{Covered}(S)$ is itself $S$-closed. We use the two resulting principles:
\emph{closure} ($\mathsf{Covered}(S)$ satisfies (cert), (layer), (sigma)) and \emph{induction} (if $P$ is $S$-closed,
then $\mathsf{Covered}(S)(n)$ implies $P(n)$ for every node $n$). In Lean, \texttt{Covered S} is an inductive predicate
whose three constructors \lean{Covered.cert}, \lean{Covered.layer}, \lean{Covered.sigma} are exactly these three rules,
with these premises in this order, and whose recursor yields the induction principle. The Lean predicate is therefore
$S$-closed and contained in every $S$-closed predicate, so it coincides with the predicate defined above.

\begin{lemma}[Monotonicity; \leanref{C3Cover.Covered.mono}{C3LEAN/C3Cover.lean:153}]\label{cb:lem:mono}
Let $S,T$ be predicates on nodes such that $S(n)$ implies $T(n)$ for every node $n$. Then for every node $n$,
$\mathsf{Covered}(S)(n)$ implies $\mathsf{Covered}(T)(n)$.
\end{lemma}

\begin{proof}
As in Lean, by induction on the derivation. The predicate $P=\mathsf{Covered}(T)$ is $S$-closed: for (cert), $S(n)$
gives $T(n)$, hence $P(n)$ by the rule (cert) for $T$; the rules (layer) and (sigma) have the same side conditions for
$S$ and for $T$, so they hold for $P$ because $\mathsf{Covered}(T)$ satisfies the rules (layer) and (sigma) for $T$. By
the induction principle, $\mathsf{Covered}(S)(n)$ implies $P(n)$.
\end{proof}

\subsection{Keys, codes and Boolean equality}\label{cb:sec:keys}

\begin{definition}[Token decoding; \leanref{C3Cover.decodeTok}{C3LEAN/C3Cover.lean:165}]\label{cb:def:decodetok}
For $x\in\N$, $\mathsf{decodeTok}(x)=(\lfloor x/3\rfloor,\ x\bmod3)$.
\end{definition}

\begin{definition}[Layer decoding; \leanref{C3Cover.decodeLayer}{C3LEAN/C3Cover.lean:167}]\label{cb:def:decodelayer}
For $c\in\N$,
\[
\mathsf{decodeLayer}(c)=\begin{cases}[\mathsf{decodeTok}(c\mathbin{\dot-}1)]&\text{if }c\le33,\\
\big[\mathsf{decodeTok}(\lfloor(c\mathbin{\dot-}34)/33\rfloor),\ \mathsf{decodeTok}((c\mathbin{\dot-}34)\bmod33)\big]
&\text{if }c>33.\end{cases}
\]
\end{definition}

\begin{definition}[Decoding layers; \leanref{C3Cover.decodeLayers}{C3LEAN/C3Cover.lean:172}]\label{cb:def:decodelayers}
For $f,lk\in\N$ and a list $acc$ of layers: $\mathsf{decodeLayers}(0,lk,acc)=acc$;
$\mathsf{decodeLayers}(f+1,lk,acc)=acc$ if $lk=0$; and $\mathsf{decodeLayers}(f+1,lk,acc)=\mathsf{decodeLayers}\big(f,\
\lfloor lk/2048\rfloor,\ \mathsf{decodeLayer}(lk\bmod2048)\mathbin{::}acc\big)$ if $lk\ne0$.
\end{definition}

\begin{definition}[Cube decoding; \leanref{C3Cover.decodeCube}{C3LEAN/C3Cover.lean:180}]\label{cb:def:decodecube}
For $f,r,cm,cv\in\N$: $\mathsf{decodeCube}(0,r,cm,cv)=[]$; $\mathsf{decodeCube}(f+1,r,cm,cv)=[]$ if $cm=0$; and if
$cm\neq0$, $\mathsf{decodeCube}(f+1,r,cm,cv)$ is $(r,\ cv\bmod2)\mathbin{::}D$ when $cm\bmod2=1$ and $D$ when
$cm\bmod2\ne1$, where $D=\mathsf{decodeCube}(f,\ r+1,\ \lfloor cm/2\rfloor,\ \lfloor cv/2\rfloor)$.
\end{definition}

\begin{definition}[Key decoding; \leanref{C3Cover.decodeKey}{C3LEAN/C3Cover.lean:191}]\label{cb:def:decodekey}
For $k\in\N$,
\[
\mathsf{decodeKey}(k)=\big\langle\mathsf{decodeLayers}(12,\ \lfloor k/2^{22}\rfloor,\ []),\
\mathsf{decodeCube}(11,\ 0,\ \lfloor k/2^{11}\rfloor\bmod2^{11},\ k\bmod2^{11})\big\rangle.
\]
(In Lean the constants are written $4194304=2^{22}$ and $2048=2^{11}$.)
\end{definition}

\begin{remark}[Keys]\label{cb:rem:keys}
Read in base $2048$, the number $\lfloor k/2^{22}\rfloor$ lists the layer codes of $\mathsf{decodeKey}(k)$, its most
significant digit giving the top layer (each step of $\mathsf{decodeLayers}$ puts the layer of the least significant
remaining digit in front of the layers decoded so far); the next $11$ binary digits of $k$ are the set of cube residues
and the last $11$ their values. For example (the review's theorem \lean{C3Cover.rev_parent_decodes},
\dirn{C3LEAN/review/RevCtrl_pass.lean:315}, whose Lean proof is \texttt{decide +kernel}; its file compiles with exit status~0,
\dirn{C3LEAN/review/logs/revctrl_pass.log}, but the axioms of this theorem are not printed),
$\mathsf{decodeKey}(8615616512)=\langle[[(0,0)],[(1,2)]],\ [(2,0),(3,0),(4,0),(5,0),(6,0),(7,0)]\rangle$. The map
$\mathsf{decodeKey}$ is not injective: value digits outside the residue mask are ignored, and a zero digit inside the
layer part decodes to the layer $[(0,0)]$. Nothing below needs injectivity: the leaves are defined as decoded keys, and
the checker compares nodes, not keys.
\end{remark}

\begin{definition}[Token code; \leanref{C3Cover.tokCode}{C3LEAN/C3Cover.lean:223}]\label{cb:def:tokcode}
For a token $t$, $\mathsf{tokCode}(t)=3\,t.1+t.2$.
\end{definition}

\begin{definition}[Layer code; \leanref{C3Cover.layerCode}{C3LEAN/C3Cover.lean:226}]\label{cb:def:layercode}
For tokens $t,u$: $\mathsf{layerCode}([t])=1+\mathsf{tokCode}(t)$,
$\mathsf{layerCode}([t,u])=34+33\,\mathsf{tokCode}(t)+\mathsf{tokCode}(u)$, and $\mathsf{layerCode}(L)=0$ for every
layer $L$ with no or more than two tokens.
\end{definition}

\begin{definition}[Root code; \leanref{C3Cover.rootCode}{C3LEAN/C3Cover.lean:232}]\label{cb:def:rootcode}
For a node $n$: $\mathsf{rootCode}(n)=\mathsf{layerCode}(A)\cdot1123+\mathsf{layerCode}(B)$ if
$n.\mathsf{layers}=[A,B]$, and $\mathsf{rootCode}(n)=0$ otherwise.
\end{definition}

\begin{definition}[Root decoding; \leanref{C3Cover.decodeRoot}{C3LEAN/C3Cover.lean:238}]\label{cb:def:decoderoot}
For $c\in\N$, $\mathsf{decodeRoot}(c)=\big\langle[\mathsf{decodeLayer}(\lfloor c/1123\rfloor),\
\mathsf{decodeLayer}(c\bmod1123)],\ []\big\rangle$.
\end{definition}

\begin{definition}[List equality; \leanref{C3Cover.beqList}{C3LEAN/C3Cover.lean:194}]\label{cb:def:beqlist}
For a type $\alpha$ and a Boolean-valued function $eq$ of two arguments in $\alpha$:
$\mathsf{beqList}(eq,[],[])=\mathsf{true}$;
$\mathsf{beqList}(eq,\,a\mathbin{::}as,\,b\mathbin{::}bs)=eq(a,b)\wedge\mathsf{beqList}(eq,as,bs)$; and
$\mathsf{beqList}(eq,l_1,l_2)=\mathsf{false}$ if exactly one of $l_1,l_2$ is empty.
\end{definition}

\begin{definition}[Token equality; \leanref{C3Cover.beqTok}{C3LEAN/C3Cover.lean:199}]\label{cb:def:beqtok}
For pairs $a,b$ of naturals, $\mathsf{beqTok}(a,b)=\mathsf{Nat.beq}(a.1,b.1)\wedge\mathsf{Nat.beq}(a.2,b.2)$.
\end{definition}

\begin{definition}[Node equality; \leanref{C3Cover.beqNode}{C3LEAN/C3Cover.lean:202}]\label{cb:def:beqnode}
For nodes $a,b$,
\[
\mathsf{beqNode}(a,b)=\mathsf{beqList}\big(\mathsf{beqList}(\mathsf{beqTok},\cdot,\cdot),\
a.\mathsf{layers},\ b.\mathsf{layers}\big)\wedge\mathsf{beqList}(\mathsf{beqTok},\ a.\mathsf{cube},\ b.\mathsf{cube}).
\]
\end{definition}

\begin{definition}[Layer equality; \leanref{C3Cover.beqLayer}{C3LEAN/C3Cover.lean:335}]\label{cb:def:beqlayer}
For layers $L,L'$, $\mathsf{beqLayer}(L,L')=\mathsf{beqList}(\mathsf{beqTok},L,L')$.
\end{definition}

\begin{lemma}[List equality; \leanref{C3Cover.beqList_eq}{C3LEAN/C3Cover.lean:204}]\label{cb:lem:beqlist}
Let $\alpha$ be a type and $eq$ a Boolean-valued function on $\alpha\times\alpha$ such that $eq(a,b)=\mathsf{true}$
implies $a=b$ for all $a,b\in\alpha$. Then for all lists $l_1,l_2$ over $\alpha$,
$\mathsf{beqList}(eq,l_1,l_2)=\mathsf{true}$ implies $l_1=l_2$.
\end{lemma}

\begin{proof}
As in Lean, by induction on $l_1$, for all $l_2$. If both lists are empty they are equal. If exactly one is empty,
$\mathsf{beqList}$ is $\mathsf{false}$, contrary to the assumption. If $l_1=a\mathbin{::}as$ and $l_2=b\mathbin{::}bs$,
the assumption gives $eq(a,b)=\mathsf{true}$ and $\mathsf{beqList}(eq,as,bs)=\mathsf{true}$; so $a=b$ by the hypothesis
on $eq$, and $as=bs$ by induction; hence $l_1=l_2$.
\end{proof}

\begin{lemma}[Token equality; \leanref{C3Cover.beqTok_eq}{C3LEAN/C3Cover.lean:213}]\label{cb:lem:beqtok}
For all tokens $a,b$: $\mathsf{beqTok}(a,b)=\mathsf{true}$ implies $a=b$.
\end{lemma}

\begin{proof}
$\mathsf{Nat.beq}(x,y)=\mathsf{true}$ implies $x=y$, so both components of $a$ and $b$ agree.
\end{proof}

\begin{lemma}[Node equality; \leanref{C3Cover.beqNode_eq}{C3LEAN/C3Cover.lean:218}]\label{cb:lem:beqnode}
For all nodes $a,b$: $\mathsf{beqNode}(a,b)=\mathsf{true}$ implies $a=b$.
\end{lemma}

\begin{proof}
By Lemmas~\ref{cb:lem:beqtok} and~\ref{cb:lem:beqlist}, $\mathsf{beqList}(\mathsf{beqTok},\cdot,\cdot)$ satisfies the
hypothesis of Lemma~\ref{cb:lem:beqlist} on the type of layers, and $\mathsf{beqTok}$ satisfies it on pairs of
naturals. Applying Lemma~\ref{cb:lem:beqlist} to the two conjuncts of $\mathsf{beqNode}(a,b)=\mathsf{true}$ gives
$a.\mathsf{layers}=b.\mathsf{layers}$ and $a.\mathsf{cube}=b.\mathsf{cube}$, hence $a=b$.
\end{proof}

\begin{lemma}[Converses; written step]\label{cb:lem:beq-refl}
$\mathsf{Nat.beq}(x,x)$, $\mathsf{beqTok}(t,t)$, $\mathsf{beqLayer}(L,L)$ and $\mathsf{beqNode}(n,n)$ are
$\mathsf{true}$ for every $x\in\N$, pair $t$, layer $L$ and node $n$. Consequently $\mathsf{beqTok}$,
$\mathsf{beqLayer}$, $\mathsf{beqNode}$ and $\mathsf{beqList}(\mathsf{Nat.beq},\cdot,\cdot)$ return $\mathsf{true}$
exactly on equal arguments.
\end{lemma}

\begin{proof}
If $eq(x,x)=\mathsf{true}$ for all $x$, then $\mathsf{beqList}(eq,l,l)=\mathsf{true}$ for every list $l$, by induction on
$l$. Apply this to $eq=\mathsf{Nat.beq}$ (reflexive), to $\mathsf{beqTok}$ (reflexive since $\mathsf{Nat.beq}$ is) and to
$\mathsf{beqList}(\mathsf{beqTok},\cdot,\cdot)$; for nodes both conjuncts are then $\mathsf{true}$. The other direction
is Lemmas~\ref{cb:lem:beqlist} (with $eq=\mathsf{Nat.beq}$ or $eq=\mathsf{beqTok}$), \ref{cb:lem:beqtok}
and~\ref{cb:lem:beqnode}.
\end{proof}

\subsection{Certificates and the checker}\label{cb:sec:checker}

\begin{definition}[Certificates; \leanref{C3Cover.Cert}{C3LEAN/C3Cover.lean:245}]\label{cb:def:cert}
A \emph{certificate} is a triple $c=\langle c.\mathsf{root},\ c.\mathsf{ds},\ c.\mathsf{ks}\rangle$ of natural numbers.
\end{definition}

\begin{definition}[Unpacking; \leanref{C3Cover.unpack}{C3LEAN/C3Cover.lean:251}]\label{cb:def:unpack}
For all $f,x\in\N$ we have $\mathsf{unpack}(0,x)=[]$; $\mathsf{unpack}(f+1,x)=[]$ if $x=0$; and
$\mathsf{unpack}(f+1,x)=(x\bmod2^{64})\mathbin{::}\mathsf{unpack}(f,\lfloor x/2^{64}\rfloor)$ if $x\ne0$ (in Lean
$2^{64}$ is written $18446744073709551616$).
\end{definition}

\begin{definition}[Keys of a certificate; \leanref{C3Cover.Cert.keys}{C3LEAN/C3Cover.lean:259}]\label{cb:def:keys}
For a certificate $c$, $\mathsf{keys}(c)=\mathsf{unpack}(100000,\ c.\mathsf{ks})$.
\end{definition}

\begin{definition}[The checker; \leanref{C3Cover.run}{C3LEAN/C3Cover.lean:262}]\label{cb:def:run}
For $f,ds\in\N$, a list $stk$ of nodes and a list $ks$ of naturals, the Boolean $\mathsf{run}(f,stk,ds,ks)$ is
defined by recursion on $f$ (the first applicable clause is used):
\begin{enumerate}[label=(R\arabic*)]
\item $\mathsf{run}(f,[],ds,ks)=(ds=0)\wedge\mathsf{isEmpty}(ks)$, for every $f$;
\item $\mathsf{run}(0,\ n\mathbin{::}rest,\ ds,ks)=\mathsf{false}$;
\item $\mathsf{run}(f+1,\ n\mathbin{::}rest,\ ds,ks)$ is computed from $d=ds\bmod16$ and $ds'=\lfloor ds/16\rfloor$:
  \begin{itemize}
  \item if $d=0$: it is $\mathsf{false}$ if $ks=[]$, and it is
    $\mathsf{beqNode}(n,\mathsf{decodeKey}(k))\wedge\mathsf{run}(f,rest,ds',ks')$ if $ks=k\mathbin{::}ks'$;
  \item if $d=1$: it is $\mathsf{decide}(2\le\card{n.\mathsf{layers}})\wedge
    \neg\mathsf{isEmpty}(\mathsf{layerChildren}(n))\wedge
    \mathsf{run}\big(f,\ \mathsf{layerChildren}(n)\mathbin{+\!\!+}rest,\ ds',\ ks\big)$;
  \item if $d\ge2$: it is $\mathsf{decide}(2\le\card{n.\mathsf{layers}})\wedge
    \mathsf{decide}(d\mathbin{\dot-}1\le\card{\mathsf{free}(n)})\wedge
    \mathsf{run}\big(f,\ \mathsf{sigmaChildren}(n,d\mathbin{\dot-}1)\mathbin{+\!\!+}rest,\ ds',\ ks\big)$.
  \end{itemize}
\end{enumerate}
\end{definition}

\begin{definition}[Root check; \leanref{C3Cover.checkRoot}{C3LEAN/C3Cover.lean:277}]\label{cb:def:checkroot}
For a node $\rho$ and a certificate $c$,
$\mathsf{checkRoot}(\rho,c)=\mathsf{run}(10^6,\ [\rho],\ c.\mathsf{ds},\ \mathsf{keys}(c))$ (in Lean $10^6$ is written
$1000000$).
\end{definition}

\begin{remark}[How a certificate is read]\label{cb:rem:cert}
The stack starts as $[\rho]$. At each step the first node $n$ is removed and the least significant base-$16$ digit $d$ of
$ds$ is consumed: $d=0$ declares $n$ a \emph{leaf}, which must be equal, as a node, to the decoded next key, and that
key is consumed; $d=1$ declares a layer split, and $d=1+K$ with $1\le K\le14$ a $\sigma$-split on the $K$ smallest
free residues of $n$ ($K$ $\sigma$ bits, $2^K$ children). For a
split the checker computes the children itself and puts them in front of the remaining stack. A certificate thus only
chooses the kind of split at each node and names the leaves, in depth-first order; it never supplies children. When
the stack is empty, the check succeeds iff no nonzero digit and no key is left. This makes the check strict (the
negative controls listed in \S\ref{cb:sec:coverage} exercise it), but it is not used in the soundness proof below.
Running out of fuel makes the check fail, so the bound $10^6$ is irrelevant for soundness; each step treats one node
of the cover tree, and the largest certificate has 849 of them (\dirn{C3LEAN/README.md}, \S1).
\end{remark}

\begin{lemma}[Checker soundness; \leanref{C3Cover.run_sound}{C3LEAN/C3Cover.lean:279}, \leanref{C3Cover.checkRoot_sound}{C3LEAN/C3Cover.lean:323}]\label{cb:lem:run}
The following hold.
\begin{enumerate}[label=(\alph*)]
\item (\lean{run_sound}) Let $S$ be a predicate on nodes and $f\in\N$. For every list $stk$ of nodes, every $ds\in\N$
  and every list $ks$ of naturals: if $\mathsf{run}(f,stk,ds,ks)=\mathsf{true}$ and $S(\mathsf{decodeKey}(k))$ holds
  for every $k\in ks$, then $\mathsf{Covered}(S)(n)$ holds for every $n\in stk$.
\item (\lean{checkRoot_sound}) Let $\rho$ be a node and $c$ a certificate with
  $\mathsf{checkRoot}(\rho,c)=\mathsf{true}$. Then $\mathsf{Covered}(S_c)(\rho)$, where $S_c(n)$ means: there is
  $k\in\mathsf{keys}(c)$ with $\mathsf{decodeKey}(k)=n$.
\end{enumerate}
\end{lemma}

\begin{proof}
We follow the Lean proof.

(a) Fix $S$. We prove by induction on $f$ the statement $\Phi(f)$: \emph{for all $stk$, $ds$, $ks$, if
$\mathsf{run}(f,stk,ds,ks)=\mathsf{true}$ and $S(\mathsf{decodeKey}(k))$ for every $k\in ks$, then
$\mathsf{Covered}(S)(n)$ for every $n\in stk$.} If $stk=[]$ there is nothing to show, for every $f$; so let
$stk=n_0\mathbin{::}rest$.

\emph{Case $f=0$.} By (R2), $\mathsf{run}(0,n_0\mathbin{::}rest,ds,ks)=\mathsf{false}$; the hypothesis fails, and
$\Phi(0)$ holds.

\emph{Step from $f$ to $f+1$.} Assume $\Phi(f)$, $\mathsf{run}(f+1,n_0\mathbin{::}rest,ds,ks)=\mathsf{true}$ and
$S(\mathsf{decodeKey}(k))$ for every $k\in ks$. Let $d=ds\bmod16$ and $ds'=\lfloor ds/16\rfloor$. We show that $n_0$
and every $n\in rest$ are covered, i.e.\ satisfy $\mathsf{Covered}(S)$.
\begin{itemize}
\item Case $d=0$. If $ks=[]$, the value would be $\mathsf{false}$; so $ks=k\mathbin{::}ks'$, with
  $\mathsf{beqNode}(n_0,\mathsf{decodeKey}(k))=\mathsf{true}$ and $\mathsf{run}(f,rest,ds',ks')=\mathsf{true}$. By
  Lemma~\ref{cb:lem:beqnode}, $n_0=\mathsf{decodeKey}(k)$; as $k\in ks$, $S(n_0)$ holds, and $n_0$ is covered by the
  rule (cert). Every $k'\in ks'$ lies in $ks$, so $S(\mathsf{decodeKey}(k'))$; hence $\Phi(f)$, applied to
  $(rest,ds',ks')$, shows that every $n\in rest$ is covered.
\item Case $d=1$. The three conjuncts are $\mathsf{true}$: $2\le\card{n_0.\mathsf{layers}}$;
  $\mathsf{isEmpty}(\mathsf{layerChildren}(n_0))=\mathsf{false}$, i.e.\ $\mathsf{layerChildren}(n_0)\ne[]$; and
  $\mathsf{run}(f,\mathsf{layerChildren}(n_0)\mathbin{+\!\!+}rest,ds',ks)=\mathsf{true}$. By $\Phi(f)$, applied to
  $(\mathsf{layerChildren}(n_0)\mathbin{+\!\!+}rest,ds',ks)$ with the same keys, every entry of
  $\mathsf{layerChildren}(n_0)$ and every entry of $rest$ is covered. The rule (layer) shows that $n_0$ is covered.
\item Case $d\ge2$. Put $K=d\mathbin{\dot-}1=d-1$, so $K\ge1>0$. The conjuncts give $2\le\card{n_0.\mathsf{layers}}$,
  $K\le\card{\mathsf{free}(n_0)}$ and $\mathsf{run}(f,\mathsf{sigmaChildren}(n_0,K)\mathbin{+\!\!+}rest,ds',ks)=
  \mathsf{true}$. By $\Phi(f)$ every entry of $\mathsf{sigmaChildren}(n_0,K)$ and of $rest$ is covered, and the rule
  (sigma) with this $K$ shows that $n_0$ is covered.
\end{itemize}
This proves $\Phi(f+1)$, and (a) follows. (Keys are consumed only at leaves; this is why the key hypothesis passes
unchanged to the recursive call in the two split cases.)

(b) Apply (a) with $S=S_c$, $f=10^6$, $stk=[\rho]$, $ds=c.\mathsf{ds}$ and $ks=\mathsf{keys}(c)$. The first hypothesis
of (a) is $\mathsf{checkRoot}(\rho,c)=\mathsf{true}$, and for every $k\in\mathsf{keys}(c)$ the statement
$S_c(\mathsf{decodeKey}(k))$ holds with witness $k$. Since $\rho\in[\rho]$, $\mathsf{Covered}(S_c)(\rho)$.
\end{proof}

\begin{lemma}[Concatenation; \leanref{C3Cover.all_append_of}{C3LEAN/C3Cover.lean:327}]\label{cb:lem:append}
Let $\alpha$ be a type, $f$ a Boolean-valued function on $\alpha$, and $l_1,l_2$ lists over $\alpha$. If
$\mathsf{all}(l_1,f)=\mathsf{true}$ and $\mathsf{all}(l_2,f)=\mathsf{true}$, then
$\mathsf{all}(l_1\mathbin{+\!\!+}l_2,f)=\mathsf{true}$.
\end{lemma}

\begin{proof}
By (L1), $\mathsf{all}(l_1\mathbin{+\!\!+}l_2,f)=\mathsf{all}(l_1,f)\wedge\mathsf{all}(l_2,f)=\mathsf{true}$. (Lean:
simplification with \lean{List.all_append}.)
\end{proof}

\subsection{Generated data and the computer-assisted facts}\label{cb:sec:data}

\begin{definition}[TSV root codes; \leanref{C3Cover.tsvRootCodes}{C3LEAN/C3Cover.lean:413}]\label{cb:def:tsv}
$\mathsf{tsvRootCodes}$ is an explicit list of natural numbers. According to its docstring and
\dirn{C3LEAN/README.md} (\S4) it is the root column of \dirn{C3FINAL/C3FINAL_mirrorh6_roots.tsv}, in file order,
written as root codes; this identification is checked by a Python script (\dirn{C3LEAN/tools/audit_keys.py}), not in
Lean.
\end{definition}

\begin{definition}[Certificate lists; \leanref{C3Cover.certs0}{C3LEAN/C3Cover.lean:423}]\label{cb:def:certsN}
$\mathsf{certs}_0,\dots,\mathsf{certs}_{126}$ (Lean \lean{certs0}, \dots, \lean{certs126}, at lines 423--1711; the
last, \lean{certs126}, at line 1668) are explicit lists of certificates, each certificate written as a triple of
numerals.
\end{definition}

\begin{definition}[All certificates; \leanref{C3Cover.certs}{C3LEAN/C3Cover.lean:1714}]\label{cb:def:certs}
The list of all certificates is, nested to the right,
\[
\mathsf{certs}=\mathsf{certs}_0\mathbin{+\!\!+}\big(\mathsf{certs}_1\mathbin{+\!\!+}\big(\cdots
\mathbin{+\!\!+}(\mathsf{certs}_{125}\mathbin{+\!\!+}\mathsf{certs}_{126})\cdots\big)\big).
\]
\end{definition}

\begin{definition}[Batch sizes; \leanref{C3Cover.bounds}{C3LEAN/C3Cover.lean:1717}]\label{cb:def:bounds}
$\mathsf{bounds}$ is an explicit list of 127 numbers, the sizes of the lists $\mathsf{certs}_N$ (between 1 and 85, with
sum 910). No theorem uses it.
\end{definition}

\begin{definition}[Key list; \leanref{C3Cover.certifiedKeys}{C3LEAN/C3Cover.lean:1737}]\label{cb:def:certifiedkeys}
$\mathsf{certifiedKeys}=\mathsf{flatMap}(\mathsf{keys},\mathsf{certs})$, the concatenation of the key lists of all
certificates. The \emph{leaf predicate} of the main theorem is
\[
S^\star(n)\quad:\Longleftrightarrow\quad\text{there is }k\in\mathsf{certifiedKeys}\text{ with }\mathsf{decodeKey}(k)=n.
\]
\end{definition}

The data were written by \dirn{C3LEAN/tools/gen_c3lean.py} from the run records of the project. Despite the name,
$\mathsf{certifiedKeys}$ is just the list of the keys of the certificates; nothing in Lean says which records exist.

\begin{definition}[Options; \leanref{C3Cover.isOption}{C3LEAN/C3Cover.lean:340}]\label{cb:def:isoption}
For a list $U$ of naturals and a layer $L$:
$\mathsf{isOption}(U,[(r,s)])=(r<11)\wedge\neg\mathsf{contains}(U,r)\wedge(s=2)$;
$\mathsf{isOption}(U,[(b,\sigma),(c,\tau)])=(b<c)\wedge(c<11)\wedge\neg\mathsf{contains}(U,b)\wedge
\neg\mathsf{contains}(U,c)\wedge(\sigma<2)\wedge(\tau<2)$; and $\mathsf{isOption}(U,L)=\mathsf{false}$ if $L$ has no
or more than two tokens.
\end{definition}

\begin{definition}[Distinct codes; \leanref{C3Cover.distinctCodes}{C3LEAN/C3Cover.lean:346}]\label{cb:def:distinctcodes}
For a list $l$ of naturals and $acc\in\N$: $\mathsf{distinctCodes}([],acc)=\mathsf{true}$ and
$\mathsf{distinctCodes}(x\mathbin{::}l,acc)=\neg\mathsf{testBit}(acc,x)\wedge\mathsf{distinctCodes}(l,\
acc\mathbin{\mathsf{or}}2^x)$.
\end{definition}

\begin{definition}[Code sets; \leanref{C3Cover.codeSet}{C3LEAN/C3Cover.lean:351}]\label{cb:def:codeset}
For a list $l$ of naturals, $\mathsf{codeSet}(l)=\mathsf{foldl}((acc,x)\mapsto acc\mathbin{\mathsf{or}}2^x,\ 0,\ l)$.
\end{definition}

\begin{definition}[Reflection facts; \leanref{C3Cover.reflectionFacts}{C3LEAN/C3Cover.lean:356}]\label{cb:def:reflectionfacts}
For a layer $t$, let $U=\mathsf{map}(u\mapsto u.1,t)$, $O=\mathsf{layerOptions}(U)$, $\iota(o)=\mathsf{refl}(t,o)$, and
let $\mathsf{kept}_t$ be as in Definition~\ref{cb:def:kept}. Then $\mathsf{reflectionFacts}(t)$ is the $6$-tuple of
the following values:
\begin{enumerate}[label=(F\arabic*)]
\item $\card{O}$;
\item $\mathsf{all}(O,\ \mathsf{isOption}(U,\cdot))\wedge\mathsf{distinctCodes}(\mathsf{map}(\mathsf{layerCode},O),\
  0)$;
\item $\mathsf{all}\big(O,\ o\mapsto\mathsf{isOption}(U,\iota(o))\wedge\mathsf{beqLayer}(\iota(\iota(o)),o)\big)$;
\item $\card{\mathsf{filter}(o\mapsto\mathsf{beqLayer}(\iota(o),o),\ O)}$;
\item $\card{\mathsf{keptOptions}(t)}$;
\item $\mathsf{all}\big(O,\ o\mapsto(\mathsf{kept}_t(o)\vee\mathsf{kept}_t(\iota(o)))\wedge
  (\mathsf{beqLayer}(\iota(o),o)\vee\neg(\mathsf{kept}_t(o)\wedge\mathsf{kept}_t(\iota(o))))\big)$.
\end{enumerate}
\end{definition}

\begin{lemma}[Distinctness and code sets; written step]\label{cb:lem:codes}
Let $l$ be a list of naturals and $acc\in\N$. Then $\mathsf{distinctCodes}(l,acc)=\mathsf{true}$ iff $l$ is
duplicate-free and no entry of $l$ lies in $\beta(acc)$. In particular $\mathsf{distinctCodes}(l,0)=\mathsf{true}$ iff $l$ is
duplicate-free. Moreover $\beta(\mathsf{codeSet}(l))=\mathrm{set}(l)$, so for lists $l,l'$ of naturals
$\mathsf{codeSet}(l)=\mathsf{codeSet}(l')$ iff $\mathrm{set}(l)=\mathrm{set}(l')$.
\end{lemma}

\begin{proof}
Induction on $l$, for all $acc$. For $l=[]$ both sides hold. For $l=x\mathbin{::}l'$ the value is $\mathsf{true}$ iff
$x\notin\beta(acc)$ and $\mathsf{distinctCodes}(l',acc\mathbin{\mathsf{or}}2^x)=\mathsf{true}$; by induction and
$\beta(acc\mathbin{\mathsf{or}}2^x)=\beta(acc)\cup\{x\}$, this holds iff $x\notin\beta(acc)$, $l'$ is duplicate-free,
and no entry of $l'$ lies in $\beta(acc)\cup\{x\}$; that is exactly the statement that $x\mathbin{::}l'$ is
duplicate-free with no entry in $\beta(acc)$. For $\beta(0)=\emptyset$ this gives the second sentence. The claim on
$\mathsf{codeSet}$ is Lemma~\ref{cb:lem:bits} with $z=0$ and $g$ the identity, together with the injectivity of
$\beta$.
\end{proof}

\begin{fact}[Batches; \leanref{C3Cover.batch0}{C3LEAN/C3Cover.lean:1740}]\label{cb:fact:batches}
(Computer-assisted.) For every $N\in\{0,\dots,126\}$,
\[
\mathsf{all}\big(\mathsf{certs}_N,\ c\mapsto\mathsf{checkRoot}(\mathsf{decodeRoot}(c.\mathsf{root}),c)\big)=
\mathsf{true}.
\]
These are the 127 theorems \lean{C3Cover.batch0}, \lean{C3Cover.batch1}, \dots,
\lean{C3Cover.batch126}; the one for $N$ is at line $1740+4N$ (the last, \lean{batch126}, at line 2244).
\end{fact}

Each of the 127 theorems is Lean-checked by kernel evaluation (\texttt{decide +kernel}); no hand
proof is given. What is evaluated: for each of the 910 certificates $c$ in the lists $\mathsf{certs}_N$, the node
$\mathsf{decodeRoot}(c.\mathsf{root})$ and the list $\mathsf{keys}(c)$ are computed, and $\mathsf{run}$ is evaluated
step by step as in Definition~\ref{cb:def:run}; at every split node the kernel computes $\mathsf{layerChildren}$, or
$\mathsf{free}$ and $\mathsf{sigmaChildren}$, from the definitions, and at every leaf it compares the node with the
decoded key by $\mathsf{beqNode}$. Over all 910 certificates this treats 76{,}256 nodes: 73{,}567 leaves, 1{,}747
$\sigma$-splits and 942 layer splits (counts recorded in \dirn{C3LEAN/c3cover_summary.json}; the review's independent
transcription of the checker into Python obtained the same counts, \dirn{C3LEAN/review/logs/indep_audit.json}). For
554 of the 910 roots the decision field $c.\mathsf{ds}$ of the certificate $c$ is the numeral $0$ and
$\mathsf{keys}(c)$ has one entry, so the certificate consists of the root itself as the only leaf. This number is our
own count from the certificate data of \dirn{C3LEAN/C3Cover.lean}; it is not a Lean statement, and no project record
states it.

\begin{lemma}[All checked; \leanref{C3Cover.all_checked}{C3LEAN/C3Cover.lean:2247}]\label{cb:lem:all-checked}
(Computer-assisted: the proof uses Fact~\ref{cb:fact:batches}.) Every certificate passes:
\[
\mathsf{all}\big(\mathsf{certs},\ c\mapsto\mathsf{checkRoot}(\mathsf{decodeRoot}(c.\mathsf{root}),c)\big)=\mathsf{true}.
\]
\end{lemma}

\begin{proof}
As in Lean. Let $p(c)=\mathsf{checkRoot}(\mathsf{decodeRoot}(c.\mathsf{root}),c)$, and let $T_{126}=\mathsf{certs}_{126}$
and $T_N=\mathsf{certs}_N\mathbin{+\!\!+}T_{N+1}$ for $N=125,\dots,0$, so that $T_0=\mathsf{certs}$ by
Definition~\ref{cb:def:certs}. By Fact~\ref{cb:fact:batches}, $\mathsf{all}(T_{126},p)=\mathsf{true}$; and if
$\mathsf{all}(T_{N+1},p)=\mathsf{true}$, then Lemma~\ref{cb:lem:append} together with Fact~\ref{cb:fact:batches} for
$\mathsf{certs}_N$ gives $\mathsf{all}(T_N,p)=\mathsf{true}$. Descending from $N=125$ to $N=0$ (126 applications of
\lean{all_append_of}, as in the Lean term) gives the claim.
\end{proof}

\begin{fact}[Reflection facts below H3; \leanref{C3Cover.reflection_facts_H3}{C3LEAN/C3Cover.lean:367}]\label{cb:fact:refl-h3}
(Computer-assisted in Lean, by kernel evaluation; written proof in Proposition~\ref{cb:prop:refl}.)
$\mathsf{reflectionFacts}(\mathsf{H3})=(190,\mathsf{true},\mathsf{true},10,100,\mathsf{true})$.
\end{fact}

\begin{fact}[Reflection facts below H6; \leanref{C3Cover.reflection_facts_H6_1_0}{C3LEAN/C3Cover.lean:370}]\label{cb:fact:refl-h6}
(Computer-assisted in Lean, by kernel evaluation; written proof in Proposition~\ref{cb:prop:refl}.)
For each $a\in\{1,\dots,5\}$ and $s\in\{0,1\}$:
$\mathsf{reflectionFacts}(\mathsf{H6}(a,s))=(153,\mathsf{true},\mathsf{true},9,81,\mathsf{true})$. These are the ten
theorems \lean{reflection_facts_H6_1_0}, \lean{reflection_facts_H6_1_1}, \lean{reflection_facts_H6_2_0}, \dots,
\lean{reflection_facts_H6_5_1}; the one for $(a,s)$ is at line $370+3(2a+s-2)$ (the last at line 397).
\end{fact}

\begin{fact}[Top cases; \leanref{C3Cover.topCases_eq}{C3LEAN/C3Cover.lean:401}]\label{cb:fact:topcases}
(Computer-assisted in Lean, by kernel evaluation; written proof in Proposition~\ref{cb:prop:refl}.)
The list of top cases is
\[
\mathsf{topCases}=[\mathsf{H3},\mathsf{H6}(1,0),\mathsf{H6}(1,1),\mathsf{H6}(2,0),\mathsf{H6}(2,1),\mathsf{H6}(3,0),
\mathsf{H6}(3,1),\mathsf{H6}(4,0),\mathsf{H6}(4,1),\mathsf{H6}(5,0),\mathsf{H6}(5,1)].
\]
\end{fact}

\begin{fact}[Strata; \leanref{C3Cover.strata}{C3LEAN/C3Cover.lean:406}]\label{cb:fact:strata}
(Computer-assisted in Lean, by kernel evaluation; written proof in Proposition~\ref{cb:prop:refl}.)
The numbers of kept options with one and with two tokens are
\begin{align*}
\mathsf{map}\big(t\mapsto\card{\mathsf{filter}(o\mapsto(\card o=1),\ \mathsf{keptOptions}(t))},\
\mathsf{topCases}\big)&=\mathsf{replicate}(11,5),\\
\mathsf{map}\big(t\mapsto\card{\mathsf{filter}(o\mapsto(\card o=2),\ \mathsf{keptOptions}(t))},\
\mathsf{topCases}\big)&=95\mathbin{::}\mathsf{replicate}(10,76).
\end{align*}
\end{fact}

Each of the 13 theorems of Facts~\ref{cb:fact:refl-h3}--\ref{cb:fact:strata} is Lean-checked by kernel evaluation of
the definitions; Proposition~\ref{cb:prop:refl} below gives a complete written proof of all of them. They are not used
in the proof of Theorem~\ref{cb:thm:c3}; they record that $\mathsf{roots}$ has the properties that the reductions of
the second layer need (Remark~\ref{cb:rem:use}).

\begin{fact}[Roots; \leanref{C3Cover.roots_match}{C3LEAN/C3Cover.lean:1720}]\label{cb:fact:roots-match}
(Computer-assisted in Lean, by kernel evaluation; written proof in Proposition~\ref{cb:prop:roots-match}.)
The list of roots has $910$ entries, and each is recovered from its code:
\[
\card{\mathsf{roots}}=910\quad\text{and}\quad
\mathsf{all}\big(\mathsf{roots},\ \rho\mapsto\mathsf{beqNode}(\mathsf{decodeRoot}(\mathsf{rootCode}(\rho)),\rho)\big)
=\mathsf{true}.
\]
\end{fact}

Lean-checked by kernel evaluation of the definitions (the kernel computes the list $\mathsf{roots}$);
Proposition~\ref{cb:prop:roots-match} below gives a complete written proof that does not use the computation.

\begin{fact}[Root codes; \leanref{C3Cover.roots_codes_match}{C3LEAN/C3Cover.lean:1725}]\label{cb:fact:roots-codes}
(Computer-assisted.) The codes of the roots match the roots of the certificates:
\[
\mathsf{beqList}\big(\mathsf{Nat.beq},\ \mathsf{map}(\mathsf{rootCode},\mathsf{roots}),\
\mathsf{map}(c\mapsto c.\mathsf{root},\mathsf{certs})\big)=\mathsf{true}.
\]
\end{fact}

\begin{fact}[TSV codes; \leanref{C3Cover.roots_eq_tsv}{C3LEAN/C3Cover.lean:1730}]\label{cb:fact:roots-tsv}
(Computer-assisted.) The following four statements hold (as one conjunction):
\begin{gather*}
\card{\mathsf{tsvRootCodes}}=910,\qquad \mathsf{distinctCodes}(\mathsf{tsvRootCodes},0)=\mathsf{true},\\
\mathsf{distinctCodes}(\mathsf{map}(c\mapsto c.\mathsf{root},\mathsf{certs}),0)=\mathsf{true},\\
\mathsf{codeSet}(\mathsf{map}(c\mapsto c.\mathsf{root},\mathsf{certs}))=\mathsf{codeSet}(\mathsf{tsvRootCodes}).
\end{gather*}
\end{fact}

Facts~\ref{cb:fact:roots-codes} and~\ref{cb:fact:roots-tsv} are computer-assisted and Lean-checked by kernel
evaluation: the kernel compares the 910 codes $\mathsf{rootCode}(\rho)$, computed from the definitions, entry by entry
and in order with the 910 generated numerals $c.\mathsf{root}$, and it tests, with the bit-set accumulator of
$\mathsf{distinctCodes}$ and the or-fold $\mathsf{codeSet}$, that the latter and the 910 numerals of
$\mathsf{tsvRootCodes}$ are duplicate-free lists with equal sets. No hand proof of them is given.

\subsection{Written proofs of the facts about the roots}\label{cb:sec:written}

Call a token $(r,s)$ \emph{valid} if $r\le10$ and $s\le2$, and a layer \emph{valid} if it consists of one or two valid
tokens.

\begin{lemma}[Sorting short lists; written step]\label{cb:lem:sort}
$\mathsf{sortTok}([t])=[t]$; $\mathsf{sortTok}([t,u])=[t,u]$ if $t.1\le u.1$, and $\mathsf{sortTok}([t,u])=[u,t]$ if
$t.1>u.1$. In particular a two-token list with distinct residues is put into increasing order of residue, and a layer
of one token, or of two tokens with increasing residues, is fixed by $\mathsf{sortTok}$.
\end{lemma}

\begin{proof}
$\mathsf{sortTok}([t])=\mathsf{insTok}(t,[])=[t]$, and $\mathsf{sortTok}([t,u])=\mathsf{insTok}(t,[u])$, which is
$[t,u]$ if $t.1\le u.1$ and $u\mathbin{::}\mathsf{insTok}(t,[])=[u,t]$ otherwise.
\end{proof}

\begin{lemma}[Layer options; written step]\label{cb:lem:options}
The following hold.
\begin{enumerate}[label=(\alph*)]
\item Let $l$ be a strictly increasing list of naturals. Then $\mathsf{pairOptions}(l)$ is duplicate-free, has
  $4\binom{\card l}{2}$ entries, and its entries are exactly the lists $[(b,\sigma),(c,\tau)]$ with $b,c\in l$, $b<c$
  and $\sigma,\tau\in\{0,1\}$.
\item Let $U$ be a list of naturals, $R=\{r\in\{0,\dots,10\}:r\notin U\}$ and $f=\card R$. Then
  $\mathsf{layerOptions}(U)$ is duplicate-free, has $f+4\binom f2$ entries, and its set of entries is
  \[
  \mathcal O(U)=\{[(r,2)]:r\in R\}\cup\{[(b,\sigma),(c,\tau)]:b,c\in R,\ b<c,\ \sigma,\tau\in\{0,1\}\}.
  \]
  Moreover, for every layer $L$: $\mathsf{isOption}(U,L)=\mathsf{true}$ iff $L\in\mathcal O(U)$.
\item Every element of $\mathcal O(U)$ is a valid layer and is fixed by $\mathsf{sortTok}$.
\end{enumerate}
\end{lemma}

\begin{proof}
(a) Induction on $l$. For $l=[]$ the list is empty and $\binom02=0$. Let $l=b\mathbin{::}l'$; then $l'$ is strictly
increasing and all its entries exceed $b$. The first part $\mathsf{flatMap}(c\mapsto Q_{b,c},l')$ consists of the
$4\card{l'}$ lists $[(b,\sigma),(c,\tau)]$ with $c\in l'$ and $\sigma,\tau\in\{0,1\}$; they are pairwise distinct
(distinct triples $(c,\sigma,\tau)$, as $l'$ is duplicate-free) and all have first residue $b$. By induction,
$\mathsf{pairOptions}(l')$ is duplicate-free with $4\binom{\card{l'}}2$ entries, all with first residue in $l'$, hence
different from $b$. So $\mathsf{pairOptions}(l)$ is duplicate-free with
$4\card{l'}+4\binom{\card{l'}}2=4\binom{\card{l'}+1}2$ entries, and each entry has the stated form ($b<c$ in the first
part since $c\in l'$; by induction in the second). Conversely, let $b',c'\in l$ with $b'<c'$ and
$\sigma,\tau\in\{0,1\}$. If $b'=b$, then $c'\ne b$, so $c'\in l'$, and $[(b,\sigma),(c',\tau)]$ is in the first part.
If $b'\ne b$, then $b'\in l'$, so $c'>b'>b$ and $c'\in l'$; by induction the list is in $\mathsf{pairOptions}(l')$.

(b) $\mathsf{rem}(U)$ is the list of the elements of $R$ in increasing order (filtering keeps the order of
$\mathsf{range}(11)$). The first part $\mathsf{map}(r\mapsto[(r,2)],\mathsf{rem}(U))$ is duplicate-free with $f$
one-token entries; the second part $\mathsf{pairOptions}(\mathsf{rem}(U))$ is described by (a), and its entries have
two tokens. So $\mathsf{layerOptions}(U)$ is duplicate-free with $f+4\binom f2$ entries, and its set of entries is
$\mathcal O(U)$. By Definition~\ref{cb:def:isoption}, $\mathsf{isOption}(U,[(r,s)])=\mathsf{true}$ iff $r\le10$,
$r\notin U$ and $s=2$; $\mathsf{isOption}(U,[(b,\sigma),(c,\tau)])=\mathsf{true}$ iff $b<c\le10$, $b,c\notin U$ and
$\sigma,\tau\in\{0,1\}$; and $\mathsf{isOption}(U,L)=\mathsf{false}$ for the other layers. This is membership in
$\mathcal O(U)$.

(c) The residues are at most $10$, the second components lie in $\{0,1,2\}$, and a two-token element has increasing
residues; apply Lemma~\ref{cb:lem:sort}.
\end{proof}

\begin{lemma}[Decoding; written step]\label{cb:lem:decode}
The following hold.
\begin{enumerate}[label=(\alph*)]
\item $\mathsf{decodeTok}(\mathsf{tokCode}(r,s))=(r,s)$ for all $r\in\N$ and $s\le2$.
\item For every valid layer $L$: $1\le\mathsf{layerCode}(L)\le1122$ and
  $\mathsf{decodeLayer}(\mathsf{layerCode}(L))=L$. In particular $\mathsf{layerCode}$ is injective on valid layers.
\item For valid layers $A,B$:
  $\mathsf{decodeRoot}\big(\mathsf{rootCode}(\langle[A,B],[]\rangle)\big)=\langle[A,B],[]\rangle$.
\end{enumerate}
\end{lemma}

\begin{proof}
(a) $\mathsf{tokCode}(r,s)=3r+s$ with $0\le s<3$, so $\lfloor(3r+s)/3\rfloor=r$ and $(3r+s)\bmod3=s$.

(b) A valid token $t=(r,s)$ has $\mathsf{tokCode}(t)=3r+s\le30+2=32$. If $L=[t]$, then
$c=\mathsf{layerCode}(L)=1+\mathsf{tokCode}(t)\in[1,33]$; as $c\le33$,
$\mathsf{decodeLayer}(c)=[\mathsf{decodeTok}(c\mathbin{\dot-}1)]=[\mathsf{decodeTok}(\mathsf{tokCode}(t))]=[t]$ by (a).
If $L=[t,u]$, let $x=\mathsf{tokCode}(t)$ and $y=\mathsf{tokCode}(u)$, both in $[0,32]$. Then
$c=34+33x+y\in[34,\ 34+33\cdot32+32]=[34,1122]$; as $c>33$, $\mathsf{decodeLayer}(c)$ is computed from
$c\mathbin{\dot-}34=33x+y$, and since $0\le y<33$, $\lfloor(33x+y)/33\rfloor=x$ and $(33x+y)\bmod33=y$; by (a),
$\mathsf{decodeLayer}(c)=[\mathsf{decodeTok}(x),\mathsf{decodeTok}(y)]=[t,u]$.

(c) $\mathsf{rootCode}(\langle[A,B],[]\rangle)=q\cdot1123+r$ with $q=\mathsf{layerCode}(A)$ and
$r=\mathsf{layerCode}(B)\in[1,1122]$ by (b). As $0\le r<1123$, $\lfloor(q\cdot1123+r)/1123\rfloor=q$ and
$(q\cdot1123+r)\bmod1123=r$. Hence $\mathsf{decodeRoot}$ returns
$\langle[\mathsf{decodeLayer}(q),\mathsf{decodeLayer}(r)],[]\rangle=\langle[A,B],[]\rangle$ by (b).
\end{proof}

\begin{lemma}[The order lexLe; written step]\label{cb:lem:lex}
For all lists $x,y$ of naturals: (a) $\mathsf{lexLe}(x,x)=\mathsf{true}$; (b) $\mathsf{lexLe}(x,y)=\mathsf{true}$ or
$\mathsf{lexLe}(y,x)=\mathsf{true}$; (c) if $\mathsf{lexLe}(x,y)=\mathsf{lexLe}(y,x)=\mathsf{true}$, then $x=y$.
\end{lemma}

\begin{proof}
Induction on $x$, for all $y$.
(a) $\mathsf{lexLe}([],[])=\mathsf{true}$, and
$\mathsf{lexLe}(a\mathbin{::}as,a\mathbin{::}as)=(a<a)\vee((a=a)\wedge\mathsf{lexLe}(as,as))=\mathsf{lexLe}(as,as)$,
which is $\mathsf{true}$ by induction.
(b) If $x=[]$ or $y=[]$, the empty list is $\mathsf{lexLe}$ the other one. Let $x=a\mathbin{::}as$ and
$y=b\mathbin{::}bs$. If $a<b$, then $\mathsf{lexLe}(x,y)=\mathsf{true}$; if $b<a$, then $\mathsf{lexLe}(y,x)=\mathsf{true}$;
if $a=b$, then $\mathsf{lexLe}(x,y)=\mathsf{lexLe}(as,bs)$ and $\mathsf{lexLe}(y,x)=\mathsf{lexLe}(bs,as)$, and one of
them is $\mathsf{true}$ by induction.
(c) If $x=[]$, then $\mathsf{lexLe}(y,[])=\mathsf{true}$ forces $y=[]$. If $x=a\mathbin{::}as$, then
$\mathsf{lexLe}(x,y)=\mathsf{true}$ forces $y=b\mathbin{::}bs$, and the hypotheses say: $a<b$ or ($a=b$ and
$\mathsf{lexLe}(as,bs)=\mathsf{true}$); and $b<a$ or ($b=a$ and $\mathsf{lexLe}(bs,as)=\mathsf{true}$). Since $a<b$
excludes both $b<a$ and $b=a$, and $b<a$ excludes $a=b$, we get $a=b$ and
$\mathsf{lexLe}(as,bs)=\mathsf{lexLe}(bs,as)=\mathsf{true}$; so $as=bs$ by induction, and $x=y$.
\end{proof}

\begin{lemma}[Counting kept options; written step]\label{cb:lem:orbits}
Let $X$ be a duplicate-free list of valid layers, each fixed by $\mathsf{sortTok}$, and let $\iota$ be a map from
layers to layers with $\iota(x)\in X$ and $\iota(\iota(x))=x$ for every $x\in X$. For $x\in X$ let
$k(x)=\mathsf{lexLe}(\mathsf{okey}(x),\mathsf{okey}(\iota(x)))$, and let $F$ be the number of $x\in X$ with
$\iota(x)=x$. Then:
\begin{enumerate}[label=(\alph*)]
\item if $\iota(x)=x$, then $k(x)=\mathsf{true}$;
\item if $\iota(x)\ne x$, then exactly one of $k(x)$, $k(\iota(x))$ is $\mathsf{true}$;
\item $\card{\mathsf{filter}(k,X)}=F+(\card X-F)/2$.
\end{enumerate}
\end{lemma}

\begin{proof}
(a) $k(x)=\mathsf{lexLe}(\mathsf{okey}(x),\mathsf{okey}(x))=\mathsf{true}$ by Lemma~\ref{cb:lem:lex}(a).

(b) For $y\in X$ we have $\mathsf{sortTok}(y)=y$, so $\mathsf{okey}(y)=\mathsf{map}(\kappa,y)$. On tokens with second
component $s\le2$, $\kappa(r,s)=3r+e(s)$ with $e(2)=0$, $e(0)=1$, $e(1)=2$, and $e$ is a bijection of $\{0,1,2\}$;
so $r=\lfloor\kappa(r,s)/3\rfloor$ and $e(s)=\kappa(r,s)\bmod3$ determine $(r,s)$, and $\kappa$ is injective on valid
tokens. Hence $\mathsf{map}(\kappa,\cdot)$ is injective on lists of valid tokens (equal images have equal lengths and
equal entries), and $\mathsf{okey}$ is injective on $\mathrm{set}(X)$. Let $\iota(x)\ne x$; both lie in $X$, so
$\mathsf{okey}(x)\ne\mathsf{okey}(\iota(x))$. Moreover
$k(\iota(x))=\mathsf{lexLe}(\mathsf{okey}(\iota(x)),\mathsf{okey}(\iota(\iota(x))))
=\mathsf{lexLe}(\mathsf{okey}(\iota(x)),\mathsf{okey}(x))$. By Lemma~\ref{cb:lem:lex}(b) at least one of $k(x)$,
$k(\iota(x))$ is $\mathsf{true}$, and by Lemma~\ref{cb:lem:lex}(c) not both.

(c) As $X$ is duplicate-free, $\card{\mathsf{filter}(k,X)}$ is the number of elements $x$ of $\mathrm{set}(X)$ with
$k(x)=\mathsf{true}$. The map $\iota$ restricted to $\mathrm{set}(X)$ is an involution of $\mathrm{set}(X)$, so
$\mathrm{set}(X)$ is the disjoint union of its $F$ fixed points and of $(\card X-F)/2$ two-element orbits
$\{x,\iota(x)\}$. By (a) every fixed point is counted, and by (b) exactly one element of each two-element orbit.
\end{proof}

\begin{proposition}[Proof of the reflection facts; written step]\label{cb:prop:refl}
Facts~\ref{cb:fact:refl-h3}, \ref{cb:fact:refl-h6}, \ref{cb:fact:topcases} and~\ref{cb:fact:strata} hold. More
precisely, let $t$ be $\mathsf{H3}$, with $R=\{1,\dots,10\}$, or $\mathsf{H6}(a,s)$ with $1\le a\le5$, $s\in\{0,1\}$,
with $R=\{1,\dots,10\}\setminus\{a\}$; let $O=\mathsf{layerOptions}(\mathsf{map}(u\mapsto u.1,t))$ and
$\iota=\mathsf{refl}(t,\cdot)$. Then $\iota$ maps $\mathrm{set}(O)$ into itself, is an involution of it, preserves the
number of tokens, and its fixed points are:
\begin{itemize}
\item for $\mathsf{H3}$: the ten layers $[(b,\sigma),(11-b,1-\sigma)]$ with $1\le b\le5$ and $\sigma\in\{0,1\}$;
\item for $\mathsf{H6}(a,s)$: the layer $[(r_0,2)]$, where $r_0=a/2$ if $a$ is even and $r_0=(a+11)/2$ if $a$ is odd
  (so $r_0=6,1,7,2,8$ for $a=1,\dots,5$), and the eight layers $[(b,\sigma),(b^*,\sigma\oplus f_b)]$ with $b\in R$,
  $b<b^*$ and $\sigma\in\{0,1\}$; here $r^*=((a+11)\mathbin{\dot-}r)\bmod11$ for $r\in R$ (that is, $r^*=a-r$ if $r<a$
  and $r^*=a+11-r$ if $r>a$), and $f_b=1$ if exactly one of $b<a$, $s=0$ holds and $f_b=0$ otherwise.
\end{itemize}
\end{proposition}

\begin{proof}
Lean checks these facts by evaluation; the following proof does not use the evaluation.

\emph{The used residues and the options.} For $t=\mathsf{H3}$, $\mathsf{map}(u\mapsto u.1,t)=[0]$; for
$t=\mathsf{H6}(a,s)$ it is $[0,a]$. So $R$ is the set of Lemma~\ref{cb:lem:options}(b), with $f=\card R=10$, resp.\
$9$. By that lemma, $O$ is duplicate-free with $\mathrm{set}(O)=\mathcal O$ (the set $\mathcal O(U)$ of the lemma),
and $\card O=10+4\binom{10}2=190$, resp.\ $9+4\binom92=153$. Let $X_1$ and $X_2$ be the sublists of one-token and of
two-token entries of $O$; they have $10$ and $180$, resp.\ $9$ and $144$, entries.

\emph{The map below $\mathsf{H3}$.} Here $\iota=\mathsf{mirror}$ (first case of Definition~\ref{cb:def:refl}). For
$r\in\{1,\dots,10\}$, $(11\mathbin{\dot-}r)\bmod11=11-r\in\{1,\dots,10\}$, and for $\sigma\in\{0,1\}$,
$1\mathbin{\dot-}\sigma=1-\sigma$. Hence, by Lemma~\ref{cb:lem:sort},
\[
\mathsf{mirror}([(r,2)])=[(11-r,2)],\qquad
\mathsf{mirror}([(b,\sigma),(c,\tau)])=[(11-c,1-\tau),(11-b,1-\sigma)]\quad(b<c),
\]
the second because $11-b>11-c$. Both results lie in $\mathcal O$, and applying the formulas twice returns the original
layer ($11-(11-r)=r$, $1-(1-\sigma)=\sigma$). So $\iota$ is an involution of $\mathcal O$ preserving the number of
tokens. A one-token option is never fixed ($11-r=r$ is impossible). A two-token option $[(b,\sigma),(c,\tau)]$, $b<c$,
is fixed iff $11-c=b$, $1-\tau=\sigma$, $11-b=c$ and $1-\sigma=\tau$, i.e.\ iff $c=11-b$ and $\tau=1-\sigma$; with
$b<c=11-b$ this means $1\le b\le5$. So $F=10$: none in $X_1$ and ten in $X_2$.

\emph{The map below $\mathsf{H6}(a,s)$.} Here $\iota=\mathsf{phi}(a,s,\cdot)$ (second case of
Definition~\ref{cb:def:refl}, with this $a$ and $s$). For $r\in R$ let $r^*=((a+11)\mathbin{\dot-}r)\bmod11$, as
in the statement. Since $r\le10<a+11$, $(a+11)\mathbin{\dot-}r=a+11-r$. If $r<a$, then $a+11-r\in[12,a+10]$, so $r^*=a-r\in[1,a-1]$; if $r>a$,
then $a+11-r\in[a+1,10]$, so $r^*=a+11-r\in[a+1,10]$. In both cases $r^*\in R$, $[r^*<a]=[r<a]$ and $(r^*)^*=r$ (indeed
$a-(a-r)=r$, resp.\ $a+11-(a+11-r)=r$). So $r\mapsto r^*$ is an involution of $R$. For $\sigma\in\{0,1\}$ the second
component of $\varphi_{a,s}(r,\sigma)$ is $\sigma\oplus f_r$, where $f_r=1$ if exactly one of $r<a$, $s=0$ holds and
$f_r=0$ otherwise (because $1\mathbin{\dot-}\sigma=\sigma\oplus1$); for $\sigma=2$ it is $2$. As $[r^*<a]=[r<a]$, we
have $f_{r^*}=f_r$. Hence, by Lemma~\ref{cb:lem:sort},
\[
\mathsf{phi}(a,s,[(r,2)])=[(r^*,2)],\qquad
\mathsf{phi}(a,s,[(b,\sigma),(c,\tau)])=\mathsf{sortTok}\big([(b^*,\sigma\oplus f_b),(c^*,\tau\oplus f_c)]\big),
\]
the latter being these two tokens in increasing order of residue ($b^*\ne c^*$ since $r\mapsto r^*$ is injective).
Both results lie in $\mathcal O$. The token map $h(r,\sigma)=(r^*,\sigma\oplus f_r)$, $h(r,2)=(r^*,2)$, satisfies
$h(h(x))=x$ on the tokens of options, since $(r^*)^*=r$ and $\sigma\oplus f_r\oplus f_{r^*}=\sigma$. For a two-token
option $o$, $\mathsf{sortTok}(\mathsf{map}(h,o))$ is $\mathsf{map}(h,o)$ or its reversal; applying $\mathsf{map}(h,\cdot)$
gives $o$ or its reversal, and $\mathsf{sortTok}$ of that is $o$, since $o$ has increasing residues. So $\iota$ is an
involution of $\mathcal O$ preserving the number of tokens.

Fixed points: $[(r,2)]$ is fixed iff $r^*=r$, i.e.\ iff $a=2r$ (case $r<a$) or $a+11=2r$ (case $r>a$). For even $a$
only the first equation has an integer solution, $r_0=a/2\in[1,a-1]$; for odd $a$ only the second,
$r_0=(a+11)/2\in[a+1,10]$ (as $a\le5$). So $r\mapsto r^*$ has exactly one fixed point $r_0$ in $R$, $X_1$ has exactly
one fixed entry, and $R\setminus\{r_0\}$ splits into four two-element orbits $\{b,b^*\}$. Let $o=[(b,\sigma),(c,\tau)]$
with $b<c$. If $\iota(o)=o$, the residues satisfy $\{b^*,c^*\}=\{b,c\}$; $b^*=b$ and $c^*=c$ is impossible (only one
fixed residue), so $c=b^*$ and $b=c^*$. In that case
$\iota(o)=\mathsf{sortTok}([(c,\sigma\oplus f_b),(b,\tau\oplus f_c)])=[(b,\tau\oplus f_c),(c,\sigma\oplus f_b)]$,
which equals $o$ iff $\tau\oplus f_c=\sigma$ and $\sigma\oplus f_b=\tau$; as $f_c=f_{b^*}=f_b$, both conditions say
$\tau=\sigma\oplus f_b$. So each of the four orbits $\{b<c\}$ carries exactly two fixed options (one for each
$\sigma$), and $F=9$: one in $X_1$ and eight in $X_2$.

\emph{Counting.} The lists $O$, $X_1$, $X_2$ satisfy the hypotheses of Lemma~\ref{cb:lem:orbits} with this $\iota$
(Lemma~\ref{cb:lem:options}(c); $\iota$ maps $\mathcal O$ to itself, preserves the number of tokens and is an
involution), and then $k=\mathsf{kept}_t$. Hence the numbers of kept entries are:
for $\mathsf{H3}$, $10+(190-10)/2=100$ in $O$, $0+10/2=5$ in $X_1$ and $10+(180-10)/2=95$ in $X_2$;
for $\mathsf{H6}(a,s)$, $9+(153-9)/2=81$ in $O$, $1+(9-1)/2=5$ in $X_1$ and $8+(144-8)/2=76$ in $X_2$.

\emph{The six components.} (F1) is $\card O=190$, resp.\ $153$. (F2): $\mathsf{all}(O,\mathsf{isOption}(U,\cdot))=
\mathsf{true}$ by Lemma~\ref{cb:lem:options}(b); the list $\mathsf{map}(\mathsf{layerCode},O)$ is duplicate-free
because $O$ is and $\mathsf{layerCode}$ is injective on valid layers (Lemma~\ref{cb:lem:decode}(b)), so
$\mathsf{distinctCodes}(\cdot,0)=\mathsf{true}$ by Lemma~\ref{cb:lem:codes}. (F3): for $o\in O$, $\iota(o)\in\mathcal O$,
so $\mathsf{isOption}(U,\iota(o))=\mathsf{true}$, and $\iota(\iota(o))=o$, so $\mathsf{beqLayer}(\iota(\iota(o)),o)=
\mathsf{true}$ by Lemma~\ref{cb:lem:beq-refl}. (F4): by Lemma~\ref{cb:lem:beq-refl},
$\mathsf{beqLayer}(\iota(o),o)=\mathsf{true}$ iff $\iota(o)=o$, so the value is $F=10$, resp.\ $9$. (F5):
$\mathsf{keptOptions}(t)=\mathsf{filter}(\mathsf{kept}_t,O)$ has $100$, resp.\ $81$, entries. (F6): for $o\in O$, by
Lemma~\ref{cb:lem:orbits}(a),(b) at least one of $\mathsf{kept}_t(o)$, $\mathsf{kept}_t(\iota(o))$ is $\mathsf{true}$;
and if $\mathsf{beqLayer}(\iota(o),o)=\mathsf{false}$, then $\iota(o)\ne o$ and not both are $\mathsf{true}$. This
proves Facts~\ref{cb:fact:refl-h3} and~\ref{cb:fact:refl-h6}.

\emph{Top cases and strata.} Unfolding Definition~\ref{cb:def:topcases} with $\mathsf{range}(5)=[0,1,2,3,4]$ gives the
list of Fact~\ref{cb:fact:topcases}. In $\mathsf{keptOptions}(t)$ the entries with one token are the kept entries of
$X_1$ and those with two tokens the kept entries of $X_2$. In the order of $\mathsf{topCases}$ ($\mathsf{H3}$ first,
then the ten $\mathsf{H6}(a,s)$), the counts are therefore $5$ eleven times, and $95$ followed by $76$ ten times. This
is Fact~\ref{cb:fact:strata}.
\end{proof}

\begin{proposition}[Proof of roots\_match; written step]\label{cb:prop:roots-match}
$\card{\mathsf{roots}}=910$, and for every $\rho\in\mathsf{roots}$,
$\mathsf{decodeRoot}(\mathsf{rootCode}(\rho))=\rho$ and
$\mathsf{beqNode}(\mathsf{decodeRoot}(\mathsf{rootCode}(\rho)),\rho)=\mathsf{true}$. Hence
Fact~\ref{cb:fact:roots-match} holds.
\end{proposition}

\begin{proof}
By (L3), $\card{\mathsf{roots}}=\sum_t\card{\mathsf{keptOptions}(t)}$ over the 11 entries of $\mathsf{topCases}$, which
is $100+10\cdot81=910$ by Proposition~\ref{cb:prop:refl}. By (L2) and (L4), every $\rho\in\mathsf{roots}$ is
$\langle[t,o],[]\rangle$ with $t\in\mathsf{topCases}$ and $o\in\mathsf{keptOptions}(t)$, so
$o\in\mathcal O(\mathsf{map}(u\mapsto u.1,t))$ (Lemma~\ref{cb:lem:options}(b)). The layer $t$ is valid
($\mathsf{H3}=[(0,0)]$, or $\mathsf{H6}(a,s)=[(0,0),(a,s)]$ with $a\le5$ and $s\le1$, by the unfolding of
$\mathsf{topCases}$ in Proposition~\ref{cb:prop:refl}), and $o$ is valid by Lemma~\ref{cb:lem:options}(c).
Lemma~\ref{cb:lem:decode}(c) gives $\mathsf{decodeRoot}(\mathsf{rootCode}(\rho))=\rho$, and
Lemma~\ref{cb:lem:beq-refl} gives $\mathsf{beqNode}(\rho,\rho)=\mathsf{true}$.
\end{proof}

\begin{corollary}[Roots and certificates; written step]\label{cb:cor:roots}
(Computer-assisted: the proof uses Facts~\ref{cb:fact:roots-codes} and~\ref{cb:fact:roots-tsv}.) The following hold.
\begin{enumerate}[label=(\alph*)]
\item $\mathsf{map}(\mathsf{rootCode},\mathsf{roots})=\mathsf{map}(c\mapsto c.\mathsf{root},\mathsf{certs})$ as lists,
  and $\card{\mathsf{certs}}=910$.
\item The 910 entries of $\mathsf{roots}$ are pairwise distinct nodes, and
  $\{\mathsf{rootCode}(\rho):\rho\in\mathsf{roots}\}=\mathrm{set}(\mathsf{tsvRootCodes})$.
\end{enumerate}
\end{corollary}

\begin{proof}
(a) Lemma~\ref{cb:lem:beqlist} with $eq=\mathsf{Nat.beq}$ (whose value $\mathsf{true}$ implies equality) turns
Fact~\ref{cb:fact:roots-codes} into the equality of lists; then
$\card{\mathsf{certs}}=\card{\mathsf{map}(\mathsf{rootCode},\mathsf{roots})}=910$ by (L3) and
Fact~\ref{cb:fact:roots-match}.
(b) By Fact~\ref{cb:fact:roots-tsv} and Lemma~\ref{cb:lem:codes}, $\mathsf{map}(c\mapsto c.\mathsf{root},\mathsf{certs})$
is duplicate-free, hence so is $\mathsf{map}(\mathsf{rootCode},\mathsf{roots})$ by (a); two positions of
$\mathsf{roots}$ holding the same node would give equal codes, so $\mathsf{roots}$ is duplicate-free. The equality of
the $\mathsf{codeSet}$ values and Lemma~\ref{cb:lem:codes} give
$\mathrm{set}(\mathsf{map}(c\mapsto c.\mathsf{root},\mathsf{certs}))=\mathrm{set}(\mathsf{tsvRootCodes})$, and (a)
identifies the left side with the set of root codes of the roots.
\end{proof}

The roots are distinct also by a hand argument (distinct top cases and distinct options give distinct nodes), but their
identification with $\mathsf{tsvRootCodes}$ rests on the computer-assisted Facts~\ref{cb:fact:roots-codes}
and~\ref{cb:fact:roots-tsv}.

\subsection{The threefold cover theorem}\label{cb:sec:c3}

\begin{theorem}[Roots covered; \leanref{C3Cover.all_roots_covered}{C3LEAN/C3Cover.lean:2251}]\label{cb:thm:c3}
(Computer-assisted: the proof uses Facts~\ref{cb:fact:batches} and~\ref{cb:fact:roots-codes}.) For every
$\rho\in\mathsf{roots}$, $\mathsf{Covered}(S^\star)(\rho)$, where $S^\star(n)$ means that there is
$k\in\mathsf{certifiedKeys}$ with $\mathsf{decodeKey}(k)=n$.
\end{theorem}

The theorem is Lean-checked (axioms \lean{propext}, \lean{Quot.sound}) and has no hypotheses. Its proof uses the
computer-assisted Facts~\ref{cb:fact:batches} (through Lemma~\ref{cb:lem:all-checked}) and~\ref{cb:fact:roots-codes},
and Fact~\ref{cb:fact:roots-match}, which also has the written proof of Proposition~\ref{cb:prop:roots-match}.

\begin{proof}
We follow the Lean proof. Let $\rho\in\mathsf{roots}$.
\begin{enumerate}[label=(\arabic*)]
\item By Fact~\ref{cb:fact:roots-match} (or Proposition~\ref{cb:prop:roots-match}), we have
  $\mathsf{beqNode}(\mathsf{decodeRoot}(\mathsf{rootCode}(\rho)),\rho)=\mathsf{true}$, so
  $\mathsf{decodeRoot}(\mathsf{rootCode}(\rho))=\rho$ by Lemma~\ref{cb:lem:beqnode}.
\item By Corollary~\ref{cb:cor:roots}(a), $\mathsf{map}(\mathsf{rootCode},\mathsf{roots})=\mathsf{map}(c\mapsto
  c.\mathsf{root},\mathsf{certs})$. As $\mathsf{rootCode}(\rho)$ is an entry of the left side (L4), there is
  $c\in\mathsf{certs}$ with $c.\mathsf{root}=\mathsf{rootCode}(\rho)$.
\item By Lemma~\ref{cb:lem:all-checked}, $\mathsf{checkRoot}(\mathsf{decodeRoot}(c.\mathsf{root}),c)=\mathsf{true}$.
  By (1) and (2), $\mathsf{decodeRoot}(c.\mathsf{root})=\mathsf{decodeRoot}(\mathsf{rootCode}(\rho))=\rho$, so
  $\mathsf{checkRoot}(\rho,c)=\mathsf{true}$.
\item By Lemma~\ref{cb:lem:run}(b), $\mathsf{Covered}(S_c)(\rho)$, with $S_c(n)$: there is $k\in\mathsf{keys}(c)$ with
  $\mathsf{decodeKey}(k)=n$.
\item If $S_c(n)$ holds with witness $k\in\mathsf{keys}(c)$, then $k\in\mathsf{flatMap}(\mathsf{keys},\mathsf{certs})=
  \mathsf{certifiedKeys}$ by (L2), since $c\in\mathsf{certs}$; so $S^\star(n)$. Lemma~\ref{cb:lem:mono} turns
  $\mathsf{Covered}(S_c)(\rho)$ into $\mathsf{Covered}(S^\star)(\rho)$.\qedhere
\end{enumerate}
\end{proof}

\begin{remark}[What the theorem says, and how the project uses it]\label{cb:rem:use}
\emph{The leaf set.} Theorem~\ref{cb:thm:c3} is only as informative as the predicate $S^\star$, which is read off the
certificates themselves. A certificate that names its root as the only leaf satisfies the checker; this is the case for
554 of the 910 roots (our count from the certificate data; see the paragraph after Fact~\ref{cb:fact:batches}), and
for those roots the statement is immediate. The substantial
claim of the project, that every node satisfying $S^\star$ names a run record in which both checkers are recorded as
accepting the proof (the criterion of \dirn{C3U2/tools/compose.py}; \dirn{C3LEAN/README.md}, \S4) and whose CNF hash
matches the base formula plus the units of the node, is checked outside Lean by
the Python scripts \dirn{C3LEAN/tools/audit_keys.py} and \dirn{C3LEAN/review/indep_audit.py} (73{,}567 leaves,
73{,}707 records, no mismatch; \dirn{C3LEAN/review/logs/indep_audit.json}). The review demonstrates the gap with a
certificate that collapses an inner $\sigma$-split of root~0 into one leaf naming a node without such a record and
still passes (\lean{C3Cover.rev_uncertified_leaf_passes}, \dirn{C3LEAN/review/RevCtrl_pass.lean:307}). The verdicts
themselves are recorded fields, and the proofs were not kept (\dirn{C3LEAN/README.md}, \S4). The project later recorded that
acceptance by \texttt{lrat-check} alone is not a certificate and that \texttt{lrat-check} verdicts are supporting only
(lines 111--114 of \dirn{WRITEUP/note/symmetric_extremal.tex}; \dirn{LRATCHECK20261105}); of the two recorded
verdicts, only the \texttt{cake\_lpr} verdict counts as a certificate (\dirn{WRITEUP/SYMMETRIC.md}, \S3).

\emph{Intended use (not established here).} By the induction principle of Definition~\ref{cb:def:covered}, the theorem
gives $P(\rho)$ for every root $\rho$ and every $S^\star$-closed predicate $P$. In the project, $P(n)$ is ``no set of
the class (threefold symmetric, 33 points, general position, no convex $7$-gon, with its canonical labelling) is
contained in the node $n$''. That $P$ is $S^\star$-closed amounts to three claims: (i) $P(n)$ for every leaf, which is
the record-based fact above together with the soundness of the CNF encoding; (ii) the rule (layer), which is Lemma~4
of \dirn{C3-S/soundness.md} and, for nodes with a nonempty cube, rule~3 of \dirn{C3U2/sigma_split.md} (Lean's rule
(layer) applies to such nodes, whereas the project's composer \dirn{C3U2/tools/compose.py} also requires a layer-split
record there, \dirn{C3LEAN/README.md}, \S5, finding~3); (iii) the rule (sigma), which is rule~2 of
\dirn{C3U2/sigma_split.md}. The passage from ``$P$ holds at every root'' to ``the class is empty'' is the cover
theorem with the reduced cut in \dirn{C3H6M/soundness_h6m.md}, which uses Lemma~5 of \dirn{C3-S/soundness.md} and
Lemma~H6M; the $90$ options below $\mathsf{H3}$ and the $720$ options below the ten $\mathsf{H6}(a,s)$ that these
reductions drop ($190-100$ and $10\cdot(153-81)$) are not roots of the file, so Theorem~\ref{cb:thm:c3} says nothing
about them. Facts~\ref{cb:fact:refl-h3}
and~\ref{cb:fact:refl-h6} record that the computed second layers have the properties those reductions need:
$\mathsf{refl}$ is an involution of the options, and the kept options contain exactly one member of each pair and every
fixed option. None of (i)--(iii), and none of the written arguments cited here, is Lean-checked or proved in this
document. The project's own list of what the Lean check leaves open, with the status it records for each part, is
\S6 of \dirn{C3LEAN/README.md}; its trust base for the threefold theorem (Theorem~1 of
\dirn{WRITEUP/SYMMETRIC.md}) is \S3 of that file.
\end{remark}

\subsection{The central 34-point split}\label{cb:sec:c2}

Every Lean reference in a statement head of this subsection is to \dirn{C2FINAL/Central34Orbits.lean}, namespace
\lean{Central34}; the statements tagged ``written step'' have no Lean counterpart.

\begin{definition}[Inner classes; \leanref{Central34.inner}{C2FINAL/Central34Orbits.lean:21}]\label{cb:def:inner}
For $a\in\N$,
\[
\mathsf{inner}(a)=\mathsf{filter}\big(i\mapsto(i\ne0)\wedge(i\ne a),\ \mathsf{range}(17)\big).
\]
\end{definition}

\begin{definition}[Pairs; \leanref{Central34.choose2}{C2FINAL/Central34Orbits.lean:24}]\label{cb:def:choose2}
For lists of naturals,
\[
\mathsf{choose2}([])=[],\qquad
\mathsf{choose2}(x\mathbin{::}xs)=\mathsf{map}\big(y\mapsto[x,y],\ xs\big)\mathbin{+\!\!+}\mathsf{choose2}(xs).
\]
\end{definition}

\begin{definition}[Triples; \leanref{Central34.choose3}{C2FINAL/Central34Orbits.lean:29}]\label{cb:def:choose3}
For lists of naturals,
\[
\mathsf{choose3}([])=[],\qquad
\mathsf{choose3}(x\mathbin{::}xs)=\mathsf{map}\big(W\mapsto x\mathbin{::}W,\ \mathsf{choose2}(xs)\big)
\mathbin{+\!\!+}\mathsf{choose3}(xs).
\]
\end{definition}

\begin{definition}[Subcases; \leanref{Central34.subcases}{C2FINAL/Central34Orbits.lean:34}]\label{cb:def:subcases}
For $a\in\N$, $\mathsf{subcases}(a)=\mathsf{choose2}(\mathsf{inner}(a))\mathbin{+\!\!+}\mathsf{choose3}(\mathsf{inner}(a))$.
\end{definition}

\begin{definition}[Masks; \leanref{Central34.mask}{C2FINAL/Central34Orbits.lean:37}]\label{cb:def:mask}
For a list $V$ of naturals,
\[
\mathsf{mask}(V)=\mathsf{foldl}\big((acc,i)\mapsto acc\mathbin{\mathsf{or}}2^i,\ 0,\ V\big).
\]
\end{definition}

\begin{definition}[Central reflection; \leanref{Central34.refl}{C2FINAL/Central34Orbits.lean:40}]\label{cb:def:c2refl}
For $a\in\N$ and a list $V$ of naturals,
\[
\mathsf{refl}(a,V)=\mathsf{map}\big(i\mapsto((a+17)\mathbin{\dot-}i)\bmod17,\ V\big).
\]
(This is \lean{Central34.refl}, not to be confused with \lean{C3Cover.refl}.)
\end{definition}

\begin{definition}[Bit sets; \leanref{Central34.bitset}{C2FINAL/Central34Orbits.lean:43}]\label{cb:def:bitset}
For a list $L$ of lists of naturals,
\[
\mathsf{bitset}(L)=\mathsf{foldl}\big((acc,W)\mapsto acc\mathbin{\mathsf{or}}2^{\mathsf{mask}(W)},\ 0,\ L\big).
\]
\end{definition}

\begin{definition}[Uncovered subcases; \leanref{Central34.uncovered}{C2FINAL/Central34Orbits.lean:46}]\label{cb:def:uncovered}
For $a\in\N$ and a list $C$ of lists of naturals, with the local abbreviation $B=\mathsf{bitset}(C)$,
\[
\mathsf{uncovered}(a,C)=\mathsf{filter}\Big(V\mapsto\neg\big(\mathsf{testBit}(B,\mathsf{mask}(V))\vee
\mathsf{testBit}(B,\mathsf{mask}(\mathsf{refl}(a,V)))\big),\ \mathsf{subcases}(a)\Big).
\]
\end{definition}

\begin{definition}[Closed lists; \leanref{Central34.closedP1}{C2FINAL/Central34Orbits.lean:51}]\label{cb:def:closedP}
The lists $\mathsf{closedP}_1,\dots,\mathsf{closedP}_8$ (Lean \lean{closedP1}, \dots, \lean{closedP8}, at lines 51,
78, 105, 132, 159, 186, 213, 240; the data end at line 264) are explicit lists of lists of naturals with $287$ entries
each, except $288$ for $\mathsf{closedP}_6$.
\end{definition}

\begin{definition}[Closed subcases; \leanref{Central34.closed}{C2FINAL/Central34Orbits.lean:267}]\label{cb:def:closed}
For $a\in\N$: $\mathsf{closed}(a)=\mathsf{closedP}_a$ if $1\le a\le8$, and $\mathsf{closed}(a)=[]$ otherwise.
\end{definition}

By Lemma~\ref{cb:lem:bits}, $\beta(\mathsf{mask}(V))=\mathrm{set}(V)$, so $\mathsf{mask}(V)=\mathsf{mask}(W)$ iff
$\mathrm{set}(V)=\mathrm{set}(W)$; and $\beta(\mathsf{bitset}(L))=\{\mathsf{mask}(W):W\in L\}$, so
$\mathsf{testBit}(\mathsf{bitset}(L),\mathsf{mask}(V))=\mathsf{true}$ iff $\mathrm{set}(V)=\mathrm{set}(W)$ for some
$W\in L$. For $1\le a\le16$ put $I_a=\{1,\dots,16\}\setminus\{a\}$.

\begin{remark}[Intended reading]\label{cb:rem:c2-reading}
In the project (header of the file and \dirn{C2FINAL/README.md}, \S1), a centrally symmetric $34$-point set consists of
$17$ antipodal classes $0,\dots,16$; in the hull case $P_a$, $1\le a\le8$, the hull is the parallelogram on the classes
$0$ and $a$, and the second convex layer is spanned by $2$ or $3$ of the other $15$ classes, which gives the subcases;
the reflection $i\mapsto(a-i)\bmod17$ is reduction~B. The lists $\mathsf{closedP}_a$ were generated by
\dirn{C2FINAL/gen_lean.py} as the subcases with a run record in which both checkers, \texttt{cake\_lpr} and
\texttt{lrat-check}, are recorded as accepting the proof and whose CNF hash was
rebuilt and matched by \dirn{C2FINAL/hashmatch.py} (or a re-solve of \dirn{C2FINAL/README.md}, \S4). This reading is not
used in any proof.
\end{remark}

\begin{lemma}[The subcases; written step]\label{cb:lem:c2-subsets}
Let $l$ be a strictly increasing list of $k$ naturals. Then $\mathsf{choose2}(l)$ is duplicate-free with $\binom k2$
entries, which are exactly the lists $[x,y]$ with $x,y\in l$, $x<y$; and $\mathsf{choose3}(l)$ is duplicate-free with
$\binom k3$ entries, which are exactly the lists $[x,y,z]$ with $x,y,z\in l$, $x<y<z$. Consequently, for
$1\le a\le16$, $\mathsf{subcases}(a)$ is a duplicate-free list of $\binom{15}2+\binom{15}3=105+455=560$ strictly
increasing lists, and $V\mapsto\mathrm{set}(V)$ is a bijection from its entries onto the set of subsets of $I_a$ with
two or three elements.
\end{lemma}

\begin{proof}
Induction on $l$; for $l=[]$ both lists are empty. Let $l=x\mathbin{::}xs$, so $xs$ is strictly increasing with $k-1$
entries, all larger than $x$. The first part of $\mathsf{choose2}(l)$ consists of the $k-1$ distinct lists $[x,y]$,
$y\in xs$, all increasing and with head $x$; by induction $\mathsf{choose2}(xs)$ is duplicate-free with
$\binom{k-1}2$ increasing entries, with heads in $xs$, hence different from $x$. So $\mathsf{choose2}(l)$ is
duplicate-free with $k-1+\binom{k-1}2=\binom k2$ entries. Every increasing pair $[u,v]$ of elements of $l$ occurs: if
$u=x$, then $v\in xs$ and $[x,v]$ is in the first part; otherwise $u,v\in xs$ (as $v>u>x$) and the pair is in
$\mathsf{choose2}(xs)$ by induction. Likewise the first part of $\mathsf{choose3}(l)$ consists of the
$\binom{k-1}2$ distinct lists $x\mathbin{::}W$ with $W\in\mathsf{choose2}(xs)$, increasing with head $x$, and
$\mathsf{choose3}(xs)$ contributes $\binom{k-1}3$ increasing triples with heads in $xs$; the total is
$\binom{k-1}2+\binom{k-1}3=\binom k3$, without repetition, and every increasing triple of elements of $l$ occurs (the
same case distinction on its first element). For $1\le a\le16$, $\mathsf{inner}(a)$ is the increasing list of the
$15$ elements of $I_a$. The pairs and the triples are disjoint by length, and a strictly increasing list is determined
by its set of entries; this gives the bijection.
\end{proof}

\begin{theorem}[Orbits; \leanref{Central34.subcases_reflection}{C2FINAL/Central34Orbits.lean:273}]\label{cb:thm:c2-refl}
(Computer-assisted in Lean, by kernel evaluation; written proof below.) The statement is
\[
\mathsf{all}\big(\mathsf{range}(8),\ k\mapsto Q(k+1)\big)=\mathsf{true},
\]
where for $a\in\N$, with $S_a=\mathsf{bitset}(\mathsf{subcases}(a))$, $Q(a)$ is the Boolean conjunction of
\begin{enumerate}[label=(\roman*)]
\item $\card{\mathsf{subcases}(a)}=560$;
\item $\mathsf{all}\big(\mathsf{subcases}(a),\ V\mapsto\mathrm{Img}_a(V)\wedge\mathrm{Inv}_a(V)\big)$, where
  $\mathrm{Img}_a(V)=\mathsf{testBit}(S_a,\mathsf{mask}(\mathsf{refl}(a,V)))$ and
  $\mathrm{Inv}_a(V)=\big(\mathsf{mask}(\mathsf{refl}(a,\mathsf{refl}(a,V)))=\mathsf{mask}(V)\big)$;
\item $\card{\mathsf{filter}\big(V\mapsto\mathsf{mask}(V)\le\mathsf{mask}(\mathsf{refl}(a,V)),\
  \mathsf{subcases}(a)\big)}=287$;
\item $\card{\mathsf{filter}\big(V\mapsto\mathsf{mask}(V)=\mathsf{mask}(\mathsf{refl}(a,V)),\
  \mathsf{subcases}(a)\big)}=14$.
\end{enumerate}
In words: for every $a\in\{1,\dots,8\}$, the statements (i)--(iv) hold.
\end{theorem}

The theorem is Lean-checked according to the record (paragraph ``Axioms'' in \S\ref{cb:sec:overview}), by kernel
evaluation of the definitions. The following written proof does not use the
evaluation.

\begin{proof}
Since $\mathsf{range}(8)=[0,\dots,7]$, the statement says that $Q(a)=\mathsf{true}$ for $a=1,\dots,8$. Fix such an
$a$ and let $\rho_a(i)=((a+17)\mathbin{\dot-}i)\bmod17$ for $i\in\{0,\dots,16\}$. As $i<a+17$,
$\rho_a(i)=(a+17-i)\bmod17$, so $\rho_a(i)\equiv a-i\pmod{17}$ and $\rho_a(i)\in\{0,\dots,16\}$; hence
$\rho_a(\rho_a(i))\equiv i$, and $\rho_a(\rho_a(i))=i$. Also $\rho_a(0)=a$ and $\rho_a(a)=0$. So $\rho_a$ is an
involution of $\{0,\dots,16\}$ that maps $I_a$ onto itself. For $V\in\mathsf{subcases}(a)$ the entries of $V$ lie in
$I_a$, $\mathsf{refl}(a,V)=\mathsf{map}(\rho_a,V)$ and $\mathrm{set}(\mathsf{refl}(a,V))=\rho_a(\mathrm{set}(V))$.

By Lemma~\ref{cb:lem:c2-subsets} there is, for each $V\in\mathsf{subcases}(a)$, a unique entry $V'$ of
$\mathsf{subcases}(a)$ with $\mathrm{set}(V')=\rho_a(\mathrm{set}(V))$, since $\rho_a(\mathrm{set}(V))\subseteq I_a$
has $\card{\mathrm{set}(V)}\in\{2,3\}$ elements; and $(V')'=V$ because $\rho_a\circ\rho_a$ is the identity. Moreover
$\mathsf{mask}(\mathsf{refl}(a,V))=\mathsf{mask}(V')$.

(i) is Lemma~\ref{cb:lem:c2-subsets}. (ii): $\mathrm{Img}_a(V)=\mathsf{true}$, because the number
$\mathsf{mask}(\mathsf{refl}(a,V))=\mathsf{mask}(V')$ lies in $\beta(S_a)$, as $V'\in\mathsf{subcases}(a)$; and
$\mathrm{Inv}_a(V)=\mathsf{true}$, because
$\mathrm{set}(\mathsf{refl}(a,\mathsf{refl}(a,V)))=\rho_a(\rho_a(\mathrm{set}(V)))=\mathrm{set}(V)$, so the two masks
agree.

(iv): $\mathsf{mask}(V)=\mathsf{mask}(\mathsf{refl}(a,V))$ iff $\mathrm{set}(V)=\rho_a(\mathrm{set}(V))$, i.e.\ iff
$\mathrm{set}(V)$ is a union of orbits of $\rho_a$ on $I_a$. A point $i\in I_a$ is fixed iff $2i\equiv a\pmod{17}$; as
$2\cdot9\equiv1\pmod{17}$, this means $i\equiv9a\pmod{17}$, so there is exactly one fixed point $x_0=9a\bmod17$ in
$\{0,\dots,16\}$, and $x_0\in I_a$ because $9a\not\equiv0$ and $9a\not\equiv a\pmod{17}$ for $1\le a\le8$ ($x_0=9,1,10,
2,11,3,12,4$ for $a=1,\dots,8$). The other $14$ elements of $I_a$ form $7$ two-element orbits. A union of orbits with
two elements is one of these $7$ orbits; one with three elements is $\{x_0\}$ together with one of them. So exactly
$7+7=14$ entries $V$ satisfy the condition.

(iii): if $\mathsf{mask}(V)=\mathsf{mask}(\mathsf{refl}(a,V))$, the condition $\le$ holds; this accounts for $14$
entries. The other $546$ entries are paired by $V\leftrightarrow V'$ with $V'\ne V$, and $\mathsf{mask}(V)$,
$\mathsf{mask}(V')$ are distinct naturals; as $\mathsf{mask}(\mathsf{refl}(a,V))=\mathsf{mask}(V')$ and
$\mathsf{mask}(\mathsf{refl}(a,V'))=\mathsf{mask}(V)$, exactly one member of each pair satisfies the condition. So the
count is $14+546/2=287$.
\end{proof}

\begin{fact}[Lists; \leanref{Central34.closed_are_subcases}{C2FINAL/Central34Orbits.lean:284}]\label{cb:fact:c2-sub}
(Computer-assisted.) The statement is
\[
\mathsf{all}\Big(\mathsf{range}(8),\ k\mapsto\mathsf{all}\big(\mathsf{closed}(k+1),\ V\mapsto
\mathsf{testBit}(\mathsf{bitset}(\mathsf{subcases}(k+1)),\mathsf{mask}(V))\big)\Big)=\mathsf{true}.
\]
Equivalently (by Lemma~\ref{cb:lem:c2-subsets} and the remarks after Definition~\ref{cb:def:closed}): for every
$a\in\{1,\dots,8\}$ and every $V\in\mathsf{closedP}_a$, $\mathrm{set}(V)$ is a subset of $I_a$ with two or three
elements.
\end{fact}

Lean-checked according to the record (paragraph ``Axioms'' in \S\ref{cb:sec:overview}), by kernel evaluation over
the explicit lists: for each $a$, the kernel computes the
masks of the $560$ subcases and their or-fold, and tests the mask of each of the $287$ (for $a=6$: $288$) entries of
$\mathsf{closedP}_a$ against it. No hand proof is given.

\begin{theorem}[Closure; \leanref{Central34.all_cases_closed}{C2FINAL/Central34Orbits.lean:291}]\label{cb:thm:c2}
(Computer-assisted.) The statement is
\[
\mathsf{all}\big(\mathsf{range}(8),\ k\mapsto\mathsf{isEmpty}(\mathsf{uncovered}(k+1,\mathsf{closed}(k+1)))\big)=
\mathsf{true};
\]
that is, $\mathsf{uncovered}(a,\mathsf{closedP}_a)=[]$ for every $a\in\{1,\dots,8\}$.
\end{theorem}

The theorem is Lean-checked according to the record (paragraph ``Axioms'' in \S\ref{cb:sec:overview}), by kernel
evaluation; no hand proof is given, and none is possible
without the explicit lists. What is evaluated: for each $a\in\{1,\dots,8\}$ the kernel computes
$B_a=\mathsf{bitset}(\mathsf{closedP}_a)$, the or-fold of $2^{\mathsf{mask}(W)}$ over the $287$ (for $a=6$: $288$)
entries $W$, and then, for each of the $560$ entries $V$ of $\mathsf{subcases}(a)$, the disjunction of the bits of
$B_a$ at $\mathsf{mask}(V)$ and at $\mathsf{mask}(\mathsf{refl}(a,V))$ (the second bit is needed only when the first
is $0$); the disjunction is $\mathsf{true}$ for every $V$, so the filtered list is empty. The two formulations in the statement agree because $\mathsf{range}(8)=[0,\dots,7]$,
$\mathsf{closed}(a)=\mathsf{closedP}_a$ for $1\le a\le8$, and $\mathsf{isEmpty}(l)=\mathsf{true}$ iff $l=[]$. The
statement is not vacuous: $\mathsf{uncovered}(a,[])=\mathsf{subcases}(a)$ has $560$ entries.

\begin{corollary}[Set form; written step]\label{cb:cor:c2}
(Computer-assisted: a consequence of Theorem~\ref{cb:thm:c2}.) Let $a\in\{1,\dots,8\}$, and for $i\in\{0,\dots,16\}$
let $\rho_a(i)=((a+17)\mathbin{\dot-}i)\bmod17$, the residue of $a-i$ modulo $17$. For every subset $T\subseteq I_a$
with two or three elements there is $W\in\mathsf{closedP}_a$ with $\mathrm{set}(W)=T$ or
$\mathrm{set}(W)=\rho_a(T)$.
\end{corollary}

\begin{proof}
This is a consequence of Theorem~\ref{cb:thm:c2}. By Lemma~\ref{cb:lem:c2-subsets}, $T=\mathrm{set}(V)$ for some
$V\in\mathsf{subcases}(a)$. Since $V$ is not an entry of $\mathsf{uncovered}(a,\mathsf{closedP}_a)=[]$,
$\mathsf{testBit}(B_a,\mathsf{mask}(V))$ or $\mathsf{testBit}(B_a,\mathsf{mask}(\mathsf{refl}(a,V)))$ is
$\mathsf{true}$, where $B_a=\mathsf{bitset}(\mathsf{closedP}_a)$. By the remarks after Definition~\ref{cb:def:closed},
there is $W\in\mathsf{closedP}_a$ with $\mathrm{set}(W)=\mathrm{set}(V)=T$ or
$\mathrm{set}(W)=\mathrm{set}(\mathsf{refl}(a,V))$; and $\mathrm{set}(\mathsf{refl}(a,V))=\rho_a(\mathrm{set}(V))=
\rho_a(T)$, because $\mathsf{refl}(a,V)=\mathsf{map}(\rho_a,V)$ by Definition~\ref{cb:def:c2refl}, the entries of $V$
lying in $I_a\subseteq\{0,\dots,16\}$.
\end{proof}

\begin{remark}[Scope of the central theorem]\label{cb:rem:c2-scope}
Theorem~\ref{cb:thm:c2} and Corollary~\ref{cb:cor:c2} say that the explicit lists meet every orbit $\{T,\rho_a(T)\}$ of
second-layer subcases, for each of the eight hull cases. That the listed subcases are refuted is not part of the
statement: the lists were selected outside Lean from run records (Remark~\ref{cb:rem:c2-reading}), and according to
\dirn{C2FINAL/README.md} (Status and \S9) most parallelogram pieces rest on recorded checker verdicts whose proofs
were deleted after checking. As in Remark~\ref{cb:rem:use}, of the two recorded checker verdicts only the
\texttt{cake\_lpr} verdict counts as a certificate; the project's trust base for its central theorem (Theorem~4 of
\dirn{WRITEUP/SYMMETRIC.md}) is \S3 of that file. The second-layer covers of the cases $P_a$, the $120$ hexagonal hull cases, the hull
cover, the reductions A ($P_a\sim P_{17-a}$) and B (the reflection), and the soundness of the encoder (\dirn{C2ENC})
are written arguments or records outside Lean and outside this document. By inspection of the data, the list
$\mathsf{closedP}_6$ contains both members of exactly one orbit, $[1,7]$ and $[5,16]$ (as sets,
$\{5,16\}=\rho_6(\{1,7\})$); this accounts for its $288$ entries. The header comment
of the file mentions a list \lean{pilotUnhashed} that the file does not define (stale text of the generator); it has
no effect on the theorems.
\end{remark}

\subsection{Coverage}\label{cb:sec:coverage}

\paragraph{Covered.} From \dirn{C3LEAN/C3Cover.lean}, every definition of the scope is covered by a definition above
(the generated family \lean{certs0}, \dots, \lean{certs126} by a single one):
\lean{m}, \lean{Tok}, \lean{Layer}, \lean{Node}, \lean{Node.used}, \lean{pairOptions}, \lean{layerOptions},
\lean{insTok}, \lean{sortTok}, \lean{topCases}, \lean{okey}, \lean{lexLe}, \lean{mirror}, \lean{phi}, \lean{refl},
\lean{keptOptions}, \lean{roots} (Definitions~\ref{cb:def:m}--\ref{cb:def:roots}); \lean{Node.free}, \lean{bitsOf},
\lean{addBits}, \lean{Node.sigmaChildren}, \lean{conflicts}, \lean{Node.layerChildren}
(Definitions~\ref{cb:def:free}--\ref{cb:def:layerchildren}); \lean{Covered} (Definition~\ref{cb:def:covered});
\lean{decodeTok}, \lean{decodeLayer}, \lean{decodeLayers}, \lean{decodeCube}, \lean{decodeKey}, \lean{tokCode},
\lean{layerCode}, \lean{rootCode}, \lean{decodeRoot}, \lean{beqList}, \lean{beqTok}, \lean{beqNode}, \lean{beqLayer}
(Definitions~\ref{cb:def:decodetok}--\ref{cb:def:beqlayer}); \lean{Cert}, \lean{unpack}, \lean{Cert.keys}, \lean{run},
\lean{checkRoot} (Definitions~\ref{cb:def:cert}--\ref{cb:def:checkroot}); \lean{tsvRootCodes},
\lean{certs0}--\lean{certs126}, \lean{certs}, \lean{bounds}, \lean{certifiedKeys}
(Definitions~\ref{cb:def:tsv}--\ref{cb:def:certifiedkeys}; the data are described, not reproduced); \lean{isOption},
\lean{distinctCodes}, \lean{codeSet}, \lean{reflectionFacts} (Definitions~\ref{cb:def:isoption}--%
\ref{cb:def:reflectionfacts}). The theorems: \lean{Covered.mono} (Lemma~\ref{cb:lem:mono}); \lean{beqList_eq},
\lean{beqTok_eq}, \lean{beqNode_eq} (Lemmas~\ref{cb:lem:beqlist}--\ref{cb:lem:beqnode}); \lean{run_sound} and
\lean{checkRoot_sound} (Lemma~\ref{cb:lem:run}); \lean{all_append_of} (Lemma~\ref{cb:lem:append});
\lean{batch0}--\lean{batch126} (Fact~\ref{cb:fact:batches}); \lean{all_checked} (Lemma~\ref{cb:lem:all-checked});
\lean{reflection_facts_H3} (Fact~\ref{cb:fact:refl-h3}); the ten \lean{reflection_facts_H6_1_0}, \dots,
\lean{reflection_facts_H6_5_1} (Fact~\ref{cb:fact:refl-h6}); \lean{topCases_eq} (Fact~\ref{cb:fact:topcases});
\lean{strata} (Fact~\ref{cb:fact:strata}); these four with the written proof of Proposition~\ref{cb:prop:refl};
\lean{roots_match} (Fact~\ref{cb:fact:roots-match}, with Proposition~\ref{cb:prop:roots-match});
\lean{roots_codes_match} and \lean{roots_eq_tsv} (Facts~\ref{cb:fact:roots-codes} and~\ref{cb:fact:roots-tsv}, with
Corollary~\ref{cb:cor:roots}); \lean{all_roots_covered} (Theorem~\ref{cb:thm:c3}). From
\dirn{C2FINAL/Central34Orbits.lean}: \lean{inner}, \lean{choose2}, \lean{choose3}, \lean{subcases}, \lean{mask},
\lean{refl}, \lean{bitset}, \lean{uncovered}, \lean{closedP1}--\lean{closedP8} (by a single definition), \lean{closed}
(Definitions~\ref{cb:def:inner}--\ref{cb:def:closed}); \lean{subcases_reflection} (Theorem~\ref{cb:thm:c2-refl}),
\lean{closed_are_subcases} (Fact~\ref{cb:fact:c2-sub}) and \lean{all_cases_closed} (Theorem~\ref{cb:thm:c2}, with
Corollary~\ref{cb:cor:c2}).

\paragraph{Not covered, deliberately.}
\begin{itemize}
\item The seven negative controls at the end of \dirn{C3LEAN/C3Cover.lean}, which are tests, not results:
  \lean{control_missing_record} (line 2266; root~0 with its fourth key removed is rejected),
  \lean{control_missing_sigma_child} (2270; root~0 without the subtree of the last child of its first $\sigma$-split),
  \lean{control_swapped_keys} (2274; two adjacent keys swapped), \lean{control_missing_layer_child} (2278; root~100,
  code $41560$, without the key of its last layer child), \lean{control_wrong_root} (2282; the certificate of root~0
  does not pass for root~1), \lean{control_strict_reduction} (2286; a strict comparison in the reduction would leave
  $810$ roots; this agrees with Proposition~\ref{cb:prop:refl}, as $810=(100-10)+10\cdot(81-9)$) and
  \lean{control_phi_mutant} (2292; $\mathsf{phi}$ without its $[s=0]$ term would keep options below $\mathsf{H6}(1,0)$
  whose root codes are not all in $\mathsf{tsvRootCodes}$). The \texttt{\#print axioms} commands (lines 2301--2327 of
  \dirn{C3LEAN/C3Cover.lean}, lines 296--298 of \dirn{C2FINAL/Central34Orbits.lean}) are not statements.
\item \dirn{C3LEAN/tools/prelude.lean} (SHA-256 \texttt{cdcf96de7739}\dots) is the verbatim source of lines 29--329 of
  \dirn{C3LEAN/C3Cover.lean}; its declarations are covered through that file.
\item \dirn{C3LEAN/Control_record.lean} (SHA-256 \texttt{5f755e897c30}\dots) and \dirn{C3LEAN/Control_sibling.lean}
  (\texttt{fb68b330e6e3}\dots): must-fail files, each a verbatim copy of the prelude plus one theorem (line 305)
  asserting that a mutated certificate of root~0 passes; they fail to compile, as intended
  (\dirn{C3LEAN/logs/control_record.log}, \dirn{C3LEAN/logs/control_sibling.log}).
\item \dirn{C3LEAN/review/RevCtrl_mut.lean} (\texttt{49bf2cce1b73}\dots): a copy of the prelude plus nine reviewer
  controls stating that mutated certificates are rejected (\lean{rev_extra_decision} has the same certificate as
  \lean{rev_extra_key}). \dirn{C3LEAN/review/RevCtrl_pass.lean} (\texttt{a135e2b92e6a}\dots): a copy of the prelude plus
  two controls stating that two non-genuine certificates pass (\lean{rev_uncertified_leaf_passes}, used in
  Remark~\ref{cb:rem:use}; \lean{rev_noncanonical_key_passes}) and the evaluation \lean{rev_parent_decodes} (used in
  Remark~\ref{cb:rem:keys}). \dirn{C3LEAN/review/RevDiff.lean} (\texttt{5aa3f6172026}\dots): a copy of the prelude plus
  an \texttt{\#eval} harness, without theorems of its own. These review files are tests and audits, not results;
  their prelude copies duplicate covered declarations.
\item The Lean package \dirn{C3L}: \dirn{C3L/C3Encoding.lean} (only imports) with \dirn{C3L/C3Encoding/Gen.lean},
  \dirn{C3L/C3Encoding/Case.lean} and \dirn{C3L/C3Encoding/Emit.lean}. It is a Lean program that writes, for
  parameters $m$ and $k$, the base CNF of the project's threefold encoding and the node CNFs (the base CNF followed by
  the unit clauses of a node) as DIMACS files, and it states no theorem (\dirn{C3L/README.md}). Its $41$ declarations
  are $39$ definitions, a structure and an abbreviation. In \dirn{C3L/C3Encoding/Gen.lean}, namespace \lean{C3}: the
  enumerations \lean{pairs}, \lean{triples}, \lean{quads}, \lean{perms2}, \lean{others}; the helpers \lean{key3},
  \lean{paritySort}, \lean{rot}, \lean{eps}; the variables and literal tables \lean{orbits}, the structure
  \lean{Tables}, \lean{mkTables}, \lean{Tables.s}, \lean{Tables.sig}, \lean{Tables.Cl}, \lean{Tables.Tl},
  \lean{Tables.invv}, \lean{Tables.Bvv}, \lean{Tables.B}, \lean{Tables.ccw}, \lean{chainVars}, \lean{numVars}; the
  clauses \lean{addCl}, \lean{push}, \lean{xor0}, \lean{secSig}, \lean{secSb}, \lean{secPrefix}, \lean{secForced},
  \lean{secXorB}, \lean{secXorT}, \lean{secChain} and the base formula \lean{baseClauses}. In
  \dirn{C3L/C3Encoding/Case.lean}, namespace \lean{C3}: the abbreviation \lean{Layer}, \lean{parseLayers},
  \lean{layerOrder} and \lean{layerUnits} (the unit clauses of a node). In \dirn{C3L/C3Encoding/Emit.lean}, root
  namespace: \lean{clauseLine}, \lean{body}, \lean{writeCnf} and \lean{main}, the output routines of the executable
  \texttt{c3\_dimacs}. That the output agrees with the CNF files of the project's Python generator is recorded in
  \dirn{C3L/logs} (comparisons outside Lean); the soundness theorem planned in \dirn{C3L/PLAN.md} (named
  \texttt{c3\_base\_sound} there) is not in the repository. CNF encodings are outside this document (item~(ii) of the
  paragraph ``What is not established here'' in \S\ref{cb:sec:overview}).
\end{itemize}

\paragraph{Treated in another section.} The Lean-checked files \dirn{AUDIT3/A1_reduction/SizeArith.lean} and
\dirn{AUDIT3/A1_reduction/verify_A1-F1/AssembledSizes.lean} (the size arithmetic of the reduction of the symmetric
$k=7$ statement to its case classes) and \dirn{ES7SPLIT/review/equivalence/BaekSeven.lean} (an instance of Baek's
theorem for abstract configurations) are outside the scope of this section; they are treated in the last section,
on size arithmetic.

\paragraph{Audit scripts of results outside Stage~1.} The five Lean files of \dirn{AUDIT2}, an adversarial review of
the project's mirror-class result, state no theorem of their own. \dirn{AUDIT2/A1/A1Dump.lean} defines \lean{cube}
(line~5) and \lean{main} (line~6), an input--output program that writes the data \lean{V2L.caseTable} (with the
leaves of the trees \lean{V2L.treeOf}), \lean{V2L.upTable} and \lean{V1L.hullLeaves} to a file, for comparison with a
Python decoding of their source files. \dirn{AUDIT2/A1/A1Audit.lean}, \dirn{AUDIT2/A2/Ax.lean},
\dirn{AUDIT2/A2/Ax2.lean} and \dirn{AUDIT2/A6/scripts/A6Check.lean} contain, apart from \texttt{import}, \texttt{open}
and \texttt{set\_option} lines and comments, only \texttt{\#check}, \texttt{\#print}, \texttt{\#print axioms} and
\texttt{\#eval} commands. The theorems they audit, such as \lean{V2L.mirrorExclusionLine_final},
\lean{V1L.hullVertexLine_of_leaves} and \lean{Mir.mirrorExclusionLine_of_unsat}, belong to the Lean developments in
\dirn{H3}, \dirn{V1-L} and \dirn{V2-L}, which lie outside Stage~1; this document does not restate them.

\section{Size arithmetic of the symmetric \texorpdfstring{$k=7$}{k=7} reduction}\label{sa:sec}

This section gives written proofs for three small Lean files of the research records that no other section covers.
Two of them check the size arithmetic of the reduction of the symmetric $k=7$ statement to its case classes; the third
instantiates Baek's theorem at $7$-gons.

\paragraph{The three files and their role.}
\begin{enumerate}[label=(\arabic*)]
\item \dirn{AUDIT3/A1_reduction/SizeArith.lean} (namespace \lean{A1}, $20$ theorems). The directory \dirn{AUDIT3}
  records an adversarial audit, run on 1~November 2026 against commit \dirn{9cf4949a5}, of the statement ``every set
  of at least $33$ points in general position with a nontrivial symmetry, isometric or affine, contains a convex
  $7$-gon''; it records an audit and claims no theorem (\dirn{AUDIT3/README.md}). Its component audit~A1 (report
  \dirn{AUDIT3/reports/A1_reduction.md}) re-derived by hand the reduction of this statement to the case theorems of
  the records for half-turns, threefold and fivefold rotations and reflections, to which the case classes of
  Definition~\ref{sr:def:classes} correspond. The auditor wrote this file to check in Lean, for all sizes $n$, the
  arithmetic of that reduction and the arithmetic behind the upper bounds of the table of the largest symmetric
  $7$-gon-free sizes (trust-table rows 4, 8 and~9 of the report).
\item \dirn{AUDIT3/A1_reduction/verify_A1-F1/AssembledSizes.lean} ($7$ theorems and one anonymous example). Finding
  A1-F1 of that audit (rated minor) is that the assembled theorem was not written anywhere. The verifier of the
  finding wrote this file to check, for all $n$, the bookkeeping of the assembled statement: the half-turn, threefold
  and fivefold cases, the two mirror cases and the fourfold case through the half-turn, together with a control
  showing that the half-turn bound fails at $32$ points (\dirn{AUDIT3/reports/A1_reduction.md}, lines 103--106). The
  file also states, at its line~24, that an order of at least $7$ gives at least one orbit other than the centre
  (\lean{big_order_orbit}); the report's description of the file, at its line~106, does not mention this lemma.
\item \dirn{ES7SPLIT/review/equivalence/BaekSeven.lean} (one theorem). It instantiates the vendored formalization
  \lean{main} of Baek's theorem, cited in this document as Theorem~\ref{mb:thm:baek}, at $7$-gons: $17$ elements of
  an abstract configuration contain a $4$-cap or a $7$-gon. The records of \dirn{ES7SPLIT} use it as the Lean-checked
  abstract form of the one published theorem that Theorems~5 and~6 and Corollary~7 of \dirn{ES7SPLIT/README.md}
  (written proofs there) use as a black box. It plays no role in the reduction.
\end{enumerate}

\paragraph{What the files do not contain.}
The statements of the first two files are facts about natural numbers. The geometric modelling that connects them
with point sets is not part of them: the header comment of \dirn{SizeArith.lean} (lines 4--8) states it and calls it
``audited by hand in the report, not formalised here''. In that modelling, a rotation of order $s$ about $c$ gives
$n=s\cdot k+e$, where $k$ is the number of orbits other than $\{c\}$, $e=1$ if $c\in S$ and $e=0$ otherwise, so
$e\le1$, and $e=0$ for even $s$; a reflection gives $n=2k+b$, where $k$ is the number of mirror pairs and $b\le2$ the
number of points on the axis; and deleting whole orbits, the centre or axis points keeps the symmetry and general
position. We prove these facts as a written step (Lemma~\ref{sa:lem:model}), which is not a declaration of these
files. Related facts are Lean-checked elsewhere in this document, in other forms: Theorem~\ref{ec:thm:compat}(a),(b),
an abstract form of the count $n=s\cdot k+e$ for a map $z\mapsto\zeta z$ of exact order $s$ in a commutative group
with zero; the selection of whole orbits (Lemma~\ref{sr:lem:orbit}) and the counting of half-turn pairs, rotation
orbits and mirror pairs in the Lean proof of Theorem~\ref{sr:thm:main}; the layer results
(Theorems~\ref{ll:thm:layers} and~\ref{ll:thm:card-peel}), through verbatim copies of which, and of
Theorem~\ref{ll:thm:reduction}, the Lean proofs of Theorems~\ref{ec:thm:fourfold} and~\ref{ec:thm:fivefold} pass
from a symmetric set to a subset of exactly $32$, respectively $35$, points (the file of these proofs repeats
lines 1--1147 of \dirn{LAYERLEAN/LayerLemma.lean} instead of importing it); and Lemmas~\ref{ll:lem:minperiod}
and~\ref{lg:lem:centre}, which the proof of
Lemma~\ref{sa:lem:model} uses. No Lean theorem connects these facts with the statements of the files of this
section. Remark~\ref{sa:rem:use} explains how Lemma~\ref{sa:lem:model} combines with the arithmetic and with the case
classes. The Lean-checked reduction of the symmetric statement to the case classes is Theorem~\ref{sr:thm:main} for
Euclidean isometries, with Theorem~\ref{ar:thm:main} for affine symmetries. Both are conditional on the classes,
and their Lean proofs do their own arithmetic: a text search of the repository's Lean files finds the names of the
declarations of \dirn{SizeArith.lean} and \dirn{AssembledSizes.lean} only in these two files, so no other Lean file
uses them.

\paragraph{Sources and evidence.}
Paths are relative to the research-records directory \dirn{cloud/es7-2026-10-19/results}, except those beginning
with \dirn{ErdosSzekeres/}, which are relative to the repository root. The table lists the three Lean files and
their logs by file name; their full paths are given in the list that follows it. The files are those of commit
\dirn{7e32ee259}; their line counts and SHA-256 hashes (computed with \mbox{\texttt{shasum -a 256}}) are:
\begin{center}\small
\begin{tabular}{@{}lrl@{}}
\toprule
file & lines & SHA-256\\
\midrule
\dirn{SizeArith.lean} & $103$ & \texttt{170384d81db76bb5e8d019e886f50b2545dc7581599f077c6431e08edb532418}\\
\dirn{SizeArith.lean.log} & $20$ & \texttt{b4bca1cd8542a5e39d2a8560570b573ece40533d9850e8b77f5fc33c41da5770}\\
\dirn{AssembledSizes.lean} & $39$ & \texttt{c08bdf200fb59f8353c999460a1069ecd7e3414874c4a79a8030c58dc36f2feb}\\
\dirn{AssembledSizes.lean.log} & $11$ & \texttt{1db4174c7c04edd5a20fad31800eaaf076a0af5cb480c08dcc1771018ae377c9}\\
\dirn{BaekSeven.lean} & $12$ & \texttt{d6d6636c78dd1a2b271ef8cd1dd09a9c869d5db8f3319f702921319eaf36c2d0}\\
\dirn{BaekSeven.out} & $3$ & \texttt{53dced1cd76f1ff0b2afb9ba3ba854ffe0ac1a1dbe895855cb6abfa108f8c747}\\
\bottomrule
\end{tabular}
\end{center}
\begin{itemize}
\item \dirn{AUDIT3/A1_reduction/SizeArith.lean}: namespace \lean{A1}, no \texttt{import} (core Lean only). Lines 1--9
  are the header comment; lines 16--80 hold the $20$ theorems; lines 84--103 are one \texttt{\#print axioms} command
  for each of them, in the same order. Its log \dirn{AUDIT3/A1_reduction/SizeArith.lean.log} is the output of these
  commands and contains no error and no warning. It reports the axioms \lean{propext}, \lean{Classical.choice},
  \lean{Quot.sound} for \lean{A1.c5_reduce}, \lean{A1.c4_reduce}, \lean{A1.c6_reduce}, \lean{A1.c2_reduce},
  \lean{A1.c5_gap} and \lean{A1.table_mirror}, and \lean{propext}, \lean{Quot.sound} for the other $14$ theorems.
\item \dirn{AUDIT3/A1_reduction/verify_A1-F1/AssembledSizes.lean}: root namespace, no \texttt{import}. The seven
  theorems are declared at lines 6, 9, 12, 15, 18, 21 and~24 (the last one with its proof on lines 25--27), the
  anonymous \texttt{example} at line~30, and lines 33--39 are a \texttt{\#print axioms} command for each theorem.
  Its log \dirn{AUDIT3/A1_reduction/verify_A1-F1/AssembledSizes.lean.log} contains the linter warning ``unused
  variable \texttt{hp}'' at line~24, column~41, then \lean{propext}, \lean{Classical.choice},
  \lean{Quot.sound} for \lean{halfturn_pairs} and \lean{propext}, \lean{Quot.sound} for the other six theorems, and a
  timing line; it contains no error.
\item \dirn{ES7SPLIT/review/equivalence/BaekSeven.lean}: imports \lean{ErdosSzekeres.Vendor.ETV.Main.Main} (hence
  Mathlib). Line~6 is \lean{baek_seven}, line~11 the command \mbox{\texttt{\#check @main}}, line~12 the command
  \mbox{\texttt{\#print axioms baek\_seven}}. Its output \dirn{ES7SPLIT/review/equivalence/BaekSeven.out} gives in
  lines 1--2 the type of \lean{main} (Remark~\ref{sa:rem:baek-seven}) and in line~3 the axioms \lean{propext},
  \lean{Classical.choice}, \lean{Quot.sound} of \lean{baek_seven}. Since \texttt{\#print axioms} reports every axiom
  in the dependency closure of a declaration, this line also covers the vendored theorem \lean{main}
  (\dirn{ErdosSzekeres/Vendor/ETV/Main/Main.lean}, SHA-256 prefix \texttt{83dd0a214204}) as it is used.
\end{itemize}
The first two files were checked with Lean~4.24.0 alone (\dirn{AUDIT3/README.md}, section ``Reproduction'';
\dirn{AUDIT3/reports/A1_reduction.md}, lines 106 and 262), the third with \mbox{\texttt{lake env lean}} from the
repository root (\dirn{ES7SPLIT/review/README.md}, section ``Reproduction''), that is, with the repository's
Lean~4.24.0 and
Mathlib \dirn{f897ebcf}. A text search of the three files finds no \texttt{sorry}, \texttt{admit},
\texttt{native\_decide} or \texttt{axiom} declaration. The logs were produced by the project and were not re-run for
this document.

\paragraph{Rigor.}
Every statement below whose head names a Lean declaration is \emph{Lean-checked}, with the axioms shown in the logs
above (all within \lean{propext}, \lean{Classical.choice}, \lean{Quot.sound}). The proofs given here are
\emph{written proofs (not refereed)}. No statement of this section is computer-assisted. None of them is a finite
computation: the Lean-checked statements are universally quantified, over natural numbers or, in
Theorem~\ref{sa:thm:baek-seven}, over linearly ordered types, configurations and finite sets, and the control of
Remark~\ref{sa:rem:control} is the negation of such a statement. The Lean proofs use \texttt{decide} only to close
side facts inside these proofs, never to evaluate a whole statement: the five
closed instances $s=2,\dots,6$ of the conclusion of Lemma~\ref{sa:lem:split}, the value $\binom62+2=17$ in
Theorem~\ref{sa:thm:baek-seven} and $32\le32$ in the control; the written proofs check these facts by hand as well.
By the convention that accompanies the rigor labels at the beginning of this document, such a use of \texttt{decide}
does not make a statement computer-assisted.
Lemma~\ref{sa:lem:model} and Remark~\ref{sa:rem:use} are written steps and are not Lean-checked.
Theorem~\ref{sa:thm:baek-seven} rests on Baek's theorem, Theorem~\ref{mb:thm:baek}, which is cited and
not proved in this document.

\paragraph{Conventions.}
In Subsections~\ref{sa:ssec:lower}--\ref{sa:ssec:assembled} all variables range over $\N=\{0,1,2,\dots\}$ (Lean
\texttt{Nat}), and each statement is universally quantified over its variables in the order in which Lean lists
them. No statement subtracts natural numbers. Lean's product \mbox{\texttt{2 * t * k}} is $(2t)k$ (multiplication
associates to the left). The Lean proofs use the tactic \texttt{omega}, a decision procedure for linear arithmetic
over $\N$ and
$\Z$ whose proof term is checked by the kernel; the written proofs are direct arguments.

%----------------------------------------------------------------------------------------------------------------------
\subsection{Lower bounds for the reduction}\label{sa:ssec:lower}

The first eleven theorems of \dirn{SizeArith.lean} (its part ``Reduction'', lines 13--56) turn ``at least $33$
points'' into the numbers of orbits or pairs that the case classes need (Remark~\ref{sa:rem:use}).

\begin{proposition}[rotations of order $2$ to $6$;\\
  \leanref{A1.c3_reduce}{AUDIT3/A1_reduction/SizeArith.lean:16},\\
  \leanref{A1.c5_reduce}{AUDIT3/A1_reduction/SizeArith.lean:19},\\
  \leanref{A1.c4_reduce}{AUDIT3/A1_reduction/SizeArith.lean:22},\\
  \leanref{A1.c6_reduce}{AUDIT3/A1_reduction/SizeArith.lean:25},\\
  \leanref{A1.c2_reduce}{AUDIT3/A1_reduction/SizeArith.lean:28}]\label{sa:prop:rot}
All variables range over $\N$.
\begin{enumerate}[label=(\alph*)]
\item \textup{(}\lean{A1.c3_reduce}\textup{)} For all $k,e\in\N$: if $e\le1$ and $33\le3k+e$, then $11\le k$.
\item \textup{(}\lean{A1.c5_reduce}\textup{)} For all $k,e\in\N$: if $e\le1$ and $33\le5k+e$, then $35\le5k+e$ and
  $7\le k$.
\item \textup{(}\lean{A1.c4_reduce}\textup{)} For every $k\in\N$: if $33\le4k$, then $36\le4k$, $8\le k$ and
  $17\le2k$.
\item \textup{(}\lean{A1.c6_reduce}\textup{)} For every $k\in\N$: if $33\le6k$, then $11\le2k$ and $17\le3k$.
\item \textup{(}\lean{A1.c2_reduce}\textup{)} For every $k\in\N$: if $33\le2k$, then $34\le2k$ and $17\le k$.
\end{enumerate}
\end{proposition}

\begin{proof}
(a) From $33\le3k+e$ and $e\le1$ we get $3k\ge32$; since $3\cdot10=30<32$, $k\ge11$.
(b) In the same way $5k\ge32$, and $5\cdot6=30<32$ gives $k\ge7$; then $5k+e\ge5k\ge35$.
(c) Since $4\cdot8=32<33\le4k$, we have $k\ge9$; hence $4k\ge36$, $k\ge8$ and $2k\ge18\ge17$. (So the statement records
the weaker bound $8\le k$.)
(d) Since $6\cdot5=30<33\le6k$, we have $k\ge6$; hence $2k\ge12\ge11$ and $3k\ge18\ge17$.
(e) Since $2\cdot16=32<33\le2k$, we have $k\ge17$, and then $2k\ge34$.
Lean proves each part with \texttt{omega}.
\end{proof}

\begin{proposition}[reflections;\\
  \leanref{A1.mirror_reduce_b0}{AUDIT3/A1_reduction/SizeArith.lean:31},\\
  \leanref{A1.mirror_reduce_b12}{AUDIT3/A1_reduction/SizeArith.lean:34}]\label{sa:prop:mirror}
All variables range over $\N$.
\begin{enumerate}[label=(\alph*)]
\item \textup{(}\lean{A1.mirror_reduce_b0}\textup{)} For every $k\in\N$: if $33\le2k$, then $17\le k$.
\item \textup{(}\lean{A1.mirror_reduce_b12}\textup{)} For all $k,b\in\N$: if $1\le b$, $b\le2$ and $33\le2k+b$,
  then $16\le k$.
\end{enumerate}
\end{proposition}

\begin{proof}
(a) As in Proposition~\ref{sa:prop:rot}(e), $2\cdot16=32<33\le2k$ gives $k\ge17$.
(b) From $33\le2k+b$ and $b\le2$ we get $2k\ge31$; since $2\cdot15=30<31$, $k\ge16$. The hypothesis $1\le b$ is not
used; it records that the statement is meant for one or two points on the axis. Lean proves both parts with
\texttt{omega}.
\end{proof}

\begin{proposition}[orders divisible by $2$, $3$ or $5$;\\
  \leanref{A1.even_order_reduce}{AUDIT3/A1_reduction/SizeArith.lean:37},\\
  \leanref{A1.mult3_order_reduce}{AUDIT3/A1_reduction/SizeArith.lean:42},\\
  \leanref{A1.mult5_order_reduce}{AUDIT3/A1_reduction/SizeArith.lean:47}]\label{sa:prop:order}
All variables range over $\N$.
\begin{enumerate}[label=(\alph*)]
\item \textup{(}\lean{A1.even_order_reduce}\textup{)} For all $t,k\in\N$: if $33\le(2t)k$, then $17\le tk$.
\item \textup{(}\lean{A1.mult3_order_reduce}\textup{)} For all $t,k,e\in\N$: if $e\le1$ and $33\le(3t)k+e$, then
  $11\le tk$.
\item \textup{(}\lean{A1.mult5_order_reduce}\textup{)} For all $t,k,e\in\N$: if $e\le1$ and $33\le(5t)k+e$, then
  $35\le5(tk)+e$.
\end{enumerate}
The Lean terms are \mbox{\texttt{2 * t * k}}, \mbox{\texttt{3 * t * k + e}} and \mbox{\texttt{5 * t * k + e}} in the
hypotheses, \mbox{\texttt{t * k}} in the conclusions of \textup{(a)} and \textup{(b)}, and
\mbox{\texttt{5 * (t * k) + e}} in the conclusion of \textup{(c)}.
\end{proposition}

\begin{proof}
By associativity, $(2t)k=2(tk)$, $(3t)k=3(tk)$ and $(5t)k=5(tk)$. Put $m=tk$. Then (a) is
Proposition~\ref{sa:prop:rot}(e) for $m$ (its conclusion $17\le m$), (b) is Proposition~\ref{sa:prop:rot}(a) for $m$,
and (c) is the first conclusion of Proposition~\ref{sa:prop:rot}(b) for $m$. The Lean proofs are the same: each first
proves the associativity identity with \lean{Nat.mul_assoc} and then calls \texttt{omega}, which treats the product
$tk$ as a single unknown.
\end{proof}

\begin{lemma}[the order split;\\
  \leanref{A1.order_split}{AUDIT3/A1_reduction/SizeArith.lean:52}]\label{sa:lem:split}\raggedright
For every $s\in\N$ with $2\le s$: ${s\bmod2=0}$, or ${s\bmod3=0}$, or ${s\bmod5=0}$, or ${7\le s}$.
\end{lemma}

Here $s\bmod q$ is the remainder of $s$ on division by $q$ (Lean \mbox{\texttt{s \% q}}); the Lean conclusion is the
disjunction \mbox{\texttt{s \% 2 = 0}} $\lor$ \mbox{\texttt{s \% 3 = 0}} $\lor$ \mbox{\texttt{s \% 5 = 0}} $\lor$
\mbox{\texttt{7} $\le$ \texttt{s}}.

\begin{proof}
Suppose that $s$ is divisible by none of $2$, $3$, $5$. Then $s\notin\{2,4,6\}$ (these are divisible by $2$),
$s\ne3$ and $s\ne5$. Since $s\ge2$, this leaves $s\ge7$. Lean argues in the same way: it splits on $s<7$ or $s\ge7$
(\lean{Nat.lt_or_ge}); in the first case \texttt{omega} gives $s\in\{2,3,4,5,6\}$, and each of the five resulting
closed instances of the statement is closed by \texttt{decide}, that is, by kernel evaluation of the decision procedure
for the whole closed disjunction (for $s=5$, for example, this evaluates the remainders of $5$ on division by $2$, $3$
and~$5$); in the second case the last alternative holds. The five closed instances are side facts inside the proof
of a statement about all $s\ge2$, which is not itself a finite computation, and the written proof establishes each of
them by hand, through the disjuncts $2\mid2$, $3\mid3$, $2\mid4$, $5\mid5$ and $2\mid6$ used above; hence the lemma
is Lean-checked and is not computer-assisted (see the paragraph ``Rigor'' at the beginning of this section).
\end{proof}

%----------------------------------------------------------------------------------------------------------------------
\subsection{Upper bounds of the size table}\label{sa:ssec:table}

The second part of \dirn{SizeArith.lean} (``Upper bounds of the table, given the case theorems'', lines 58--80) turns
bounds on the number of orbits or pairs into bounds on the size. The case theorems that give such bounds for sets
without $7$ points in convex position are not part of the statements: each part of Proposition~\ref{sa:prop:table}
that rests on a case theorem assumes the bound that the theorem gives as a hypothesis, on $k$ in parts (a), (e)
and~(h), on $2k$ in parts (f) and~(g), on $5k+e$ in part~(c), and on $k$, separately for $b=0$ and for $1\le b$, in
part~(i) (Remark~\ref{sa:rem:use}, item~(5)). Parts (b) and~(d), \lean{A1.c3_no32} and \lean{A1.c5_gap}, rest on no
case theorem: their only hypothesis is $e\le1$.

\begin{proposition}[upper bounds;\\
  \leanref{A1.table_c3}{AUDIT3/A1_reduction/SizeArith.lean:61},\\
  \leanref{A1.c3_no32}{AUDIT3/A1_reduction/SizeArith.lean:62},\\
  \leanref{A1.table_c5}{AUDIT3/A1_reduction/SizeArith.lean:65},\\
  \leanref{A1.c5_gap}{AUDIT3/A1_reduction/SizeArith.lean:66},\\
  \leanref{A1.table_c4}{AUDIT3/A1_reduction/SizeArith.lean:69},\\
  \leanref{A1.table_c4_only_c2}{AUDIT3/A1_reduction/SizeArith.lean:70},\\
  \leanref{A1.table_c6}{AUDIT3/A1_reduction/SizeArith.lean:73},\\
  \leanref{A1.table_c2}{AUDIT3/A1_reduction/SizeArith.lean:76},\\
  \leanref{A1.table_mirror}{AUDIT3/A1_reduction/SizeArith.lean:79}]\label{sa:prop:table}
All variables range over $\N$.
\begin{enumerate}[label=(\alph*)]
\item \textup{(}\lean{A1.table_c3}\textup{)} For all $k,e\in\N$: if $e\le1$ and $k\le10$, then $3k+e\le31$.
\item \textup{(}\lean{A1.c3_no32}\textup{)} For all $k,e\in\N$: if $e\le1$, then $3k+e\ne32$.
\item \textup{(}\lean{A1.table_c5}\textup{)} For all $k,e\in\N$: if $e\le1$ and $5k+e<35$, then $5k+e\le31$.
\item \textup{(}\lean{A1.c5_gap}\textup{)} For all $k,e\in\N$: if $e\le1$, then $5k+e\ne32$, $5k+e\ne33$ and
  $5k+e\ne34$.
\item \textup{(}\lean{A1.table_c4}\textup{)} For every $k\in\N$: if $k\le7$, then $4k\le28$.
\item \textup{(}\lean{A1.table_c4_only_c2}\textup{)} For every $k\in\N$: if $2k\le16$, then $4k\le32$.
\item \textup{(}\lean{A1.table_c6}\textup{)} For every $k\in\N$: if $2k\le10$, then $6k\le30$.
\item \textup{(}\lean{A1.table_c2}\textup{)} For every $k\in\N$: if $k\le16$, then $2k\le32$.
\item \textup{(}\lean{A1.table_mirror}\textup{)} For all $k,b\in\N$: if $b\le2$, $(b=0\Rightarrow k\le16)$ and
  $(1\le b\Rightarrow k\le15)$, then $2k+b\le32$.
\end{enumerate}
\end{proposition}

\begin{proof}
(a) $3k+e\le3\cdot10+1=31$.
(b) If $3k+e=32$, then $3k=32-e\in\{31,32\}$; but $31=3\cdot10+1$ and $32=3\cdot10+2$ are not multiples of $3$.
(c) From $5k\le5k+e<35$ we get $k\le6$, so $5k+e\le30+1=31$.
(d) If $5k+e=j$ with $j\in\{32,33,34\}$, then $5k=j-e$ lies between $31$ and $34$, strictly between the consecutive
multiples $30$ and $35$ of $5$; this is impossible.
(e) $4k\le4\cdot7=28$.
(f) $4k=2\cdot(2k)\le2\cdot16=32$.
(g) $6k=3\cdot(2k)\le3\cdot10=30$.
(h) $2k\le2\cdot16=32$.
(i) If $b=0$, then $k\le16$ and $2k+b=2k\le32$. Otherwise $1\le b$, so $k\le15$, and with $b\le2$ we get
$2k+b\le30+2=32$.
Lean proves each part with \texttt{omega}.
\end{proof}

%----------------------------------------------------------------------------------------------------------------------
\subsection{Size bookkeeping for the assembled statement (audit finding A1-F1)}\label{sa:ssec:assembled}

The file \dirn{AssembledSizes.lean} restates six of the statements above, or parts of them, with the size $n$ as an
explicit variable, adds \lean{big_order_orbit} and a negative control. Its declarations are in the root namespace,
and the file does not import \dirn{SizeArith.lean}.

\begin{proposition}[assembled sizes;\\
  \leanref{halfturn_pairs}{AUDIT3/A1_reduction/verify_A1-F1/AssembledSizes.lean:6},\\
  \leanref{threefold_orbits}{AUDIT3/A1_reduction/verify_A1-F1/AssembledSizes.lean:9},\\
  \leanref{fivefold_size}{AUDIT3/A1_reduction/verify_A1-F1/AssembledSizes.lean:12},\\
  \leanref{mirror_b0}{AUDIT3/A1_reduction/verify_A1-F1/AssembledSizes.lean:15},\\
  \leanref{mirror_b12}{AUDIT3/A1_reduction/verify_A1-F1/AssembledSizes.lean:18},\\
  \leanref{fourfold_via_halfturn}{AUDIT3/A1_reduction/verify_A1-F1/AssembledSizes.lean:21}]\label{sa:prop:assembled}
All variables range over $\N$.
\begin{enumerate}[label=(\alph*)]
\item \textup{(}\lean{halfturn_pairs}\textup{)} For all $n,k\in\N$: if $n=2k$ and $33\le n$, then $17\le k$ and
  $34\le n$.
\item \textup{(}\lean{threefold_orbits}\textup{)} For all $n,k,e\in\N$: if $e\le1$, $n=3k+e$ and $33\le n$, then
  $11\le k$.
\item \textup{(}\lean{fivefold_size}\textup{)} For all $n,k,e\in\N$: if $e\le1$, $n=5k+e$ and $33\le n$, then
  $35\le n$.
\item \textup{(}\lean{mirror_b0}\textup{)} For all $n,k\in\N$: if $n=2k$ and $33\le n$, then $17\le k$.
\item \textup{(}\lean{mirror_b12}\textup{)} For all $n,k,b\in\N$: if $1\le b$, $b\le2$, $n=2k+b$ and $33\le n$, then
  $16\le k$.
\item \textup{(}\lean{fourfold_via_halfturn}\textup{)} For all $n,k\in\N$: if $n=4k$ and $33\le n$, then
  $17\le2k$.
\end{enumerate}
\end{proposition}

\begin{proof}
Substitute the given value of $n$ into $33\le n$. (a) Proposition~\ref{sa:prop:rot}(e) gives $17\le k$ and
$34\le2k=n$. (b) Proposition~\ref{sa:prop:rot}(a). (c) Proposition~\ref{sa:prop:rot}(b) gives $35\le5k+e=n$.
(d) Proposition~\ref{sa:prop:mirror}(a). (e) Proposition~\ref{sa:prop:mirror}(b); as there, $1\le b$ is not used.
(f) Proposition~\ref{sa:prop:rot}(c) gives $17\le2k$. Lean proves each part directly with \texttt{omega}.
\end{proof}

\begin{lemma}[an orbit for an order of at least $7$;\\
  \leanref{big_order_orbit}{AUDIT3/A1_reduction/verify_A1-F1/AssembledSizes.lean:24}]\label{sa:lem:big-order}\raggedright
For all $n,k,e,p\in\N$: if $7\le p$, $e\le1$, $n=pk+e$ and $33\le n$, then $1\le k$.
\end{lemma}

\begin{proof}
If $k=0$, then $n=p\cdot0+e=e\le1<33$, which contradicts $33\le n$. Hence $k\ge1$. We follow the Lean proof, which
splits on $k$ (the product $pk$ is not linear in the unknowns): for $k=0$, \texttt{simp} turns the hypothesis
$n=p\cdot0+e$ into $n=e$, and \texttt{omega} derives a contradiction from $n=e$, $e\le1$ and $33\le n$; for $k=k'+1$
the goal $1\le k'+1$ is closed by \texttt{omega}.
\end{proof}

\begin{remark}[the hypothesis $7\le p$]\label{sa:rem:big-order}
The proof does not use the hypothesis $7\le p$ (Lean binder \texttt{hp}): the conclusion holds for every $p\in\N$,
given $e\le1$, $n=pk+e$ and $33\le n$. The compile log records this as the warning ``unused variable \texttt{hp}''
at line~24, column~41. The name of the lemma and the comment above it (``prime order $p\ge7$'') describe the intended
use (Remark~\ref{sa:rem:use}, item~(1)); the statement does not require $p$ to be prime.
\end{remark}

\begin{remark}[the negative control, \dirn{AssembledSizes.lean:30}]\label{sa:rem:control}
Line~30 of \dirn{AssembledSizes.lean} is an anonymous \texttt{example}, not a named declaration. It states
\[
\neg\bigl(\forall n,k\in\N:\ n=2k\ \Rightarrow\ (32\le n\ \Rightarrow\ 17\le k)\bigr),
\]
that is, Proposition~\ref{sa:prop:assembled}(a) fails when $33$ is replaced by $32$. \emph{Proof.} The numbers $n=32$
and $k=16$ satisfy $32=2\cdot16$ and $32\le32$, but not $17\le16$. The Lean proof is the same:
\texttt{intro h; have := h 32 16 rfl (by decide); omega}, where \texttt{rfl} proves $32=2\cdot16$ by evaluation,
\texttt{decide} proves $32\le32$, and \texttt{omega} refutes $17\le16$. As the example has no name, no
\texttt{\#print axioms} line refers to it; the compile log contains no error, so Lean accepted it. The comment above
the example (line~29) reads ``control: the half-turn line is false below 33 (32 points = 16 pairs; the central
32-point example exists)''. The A1 report describes that example, \dirn{SYM2/sym7_32points.txt}, as half-turn
symmetric about $(0,0)$, not containing the centre, with no collinear triple, and with largest subset in convex
position of $6$ points (\dirn{AUDIT3/reports/A1_reduction.md}, lines 93 and~107); the set is not examined in this
section.
\end{remark}

%----------------------------------------------------------------------------------------------------------------------
\subsection{Baek's theorem at \texorpdfstring{$7$}{7}-gons}\label{sa:ssec:baek}

We recall the notions of Theorem~\ref{mb:thm:baek}. They are defined in Lean in
\dirn{ErdosSzekeres/Vendor/ETV/Config/Defs.lean}, lines 20--67, with the list predicates \lean{List.In} (line~22) and
\texttt{List.Chain3\textquotesingle} (line~46, through the inductive predicate \lean{List.Chain3} of line~42) of
\dirn{ErdosSzekeres/Vendor/ETV/Lib/List/Defs.lean} and the predicate \lean{List.IsChain}, which Mathlib takes from
its dependency Batteries. Let $(\gamma,<)$ be a linearly ordered set. A
\emph{configuration} $C$ on $\gamma$ is a ternary relation $\mathrm{Cup3}_C$ on $\gamma$ (in Lean a structure
\texttt{Config}~$\gamma$ with this relation as a field and a decidability witness, which plays no mathematical role);
$\mathrm{Cap3}_C(a,b,c)$ means $\neg\,\mathrm{Cup3}_C(a,b,c)$. A list $l=(l_0,\dots,l_{r-1})$ of elements of $\gamma$
is a \emph{$C$-cap} if $l_0<l_1<\dots<l_{r-1}$ and $\mathrm{Cap3}_C(l_i,l_{i+1},l_{i+2})$ for every $i$ with $i+2<r$,
and a \emph{$C$-cup} if the same holds with $\mathrm{Cup3}_C$ in place of $\mathrm{Cap3}_C$. For a finite set
$S\subseteq\gamma$, $\mathsf{HasNCap}_C(4,S)$ means that some $C$-cap of length $4$ has all its entries in $S$, and
$\mathsf{HasNGon}_C(7,S)$ means that there are lists $l$ and $l'$ with all entries in $S$ such that $\card l\ge2$, $l$
is a $C$-cap, $\card{l'}\ge2$, $l'$ is a $C$-cup, $l$ and $l'$ have the same first entry and the same last entry, and
$\card l+\card{l'}=7+2$. The relation $\mathrm{Cup3}_C$ is arbitrary, so these notions carry no geometry.

\begin{theorem}[Baek's theorem at $7$-gons;\\
  \leanref{baek_seven}{ES7SPLIT/review/equivalence/BaekSeven.lean:6}]\label{sa:thm:baek-seven}
Let $\gamma$ be a type \textup{(}Lean \mbox{\texttt{$\gamma$ : Type}}, the lowest universe\textup{)} with a linear
order $<$, let $C$ be a configuration on $\gamma$ \textup{(}\mbox{\texttt{C : Config}~$\gamma$}\textup{)}, and let
$S$ be
a finite subset of $\gamma$ \textup{(}\mbox{\texttt{S : Finset}~$\gamma$}\textup{)} with $17\le\card S$. Then
$\mathsf{HasNCap}_C(4,S)$ or $\mathsf{HasNGon}_C(7,S)$.
\end{theorem}

The theorem is Lean-checked (\dirn{BaekSeven.out}, line~3). Its written proof rests on the cited
Theorem~\ref{mb:thm:baek}.

\begin{proof}
We have $\binom{4+2}{2}+2=\binom62+2=\frac{6\cdot5}{2}+2=15+2=17\le\card S$. Theorem~\ref{mb:thm:baek} with $n=4$
and the same $\gamma$, $C$ and $S$ gives $\mathsf{HasNCap}_C(4,S)$ or $\mathsf{HasNGon}_C(4+3,S)$, and $4+3=7$. The
Lean proof
is the same. Its first line, \texttt{have e : Nat.choose (4 + 2) 2 + 2 = 17 := by decide}, evaluates the binomial
coefficient in the kernel (the value $\binom62=15$ computed above by hand). Its second line,
\texttt{exact main 4 C S (e} $\blacktriangleright$ \texttt{h)}, applies \lean{main} with $n=4$: the term
\texttt{e} $\blacktriangleright$ \texttt{h} (Lean's rewriting operator) turns the hypothesis \texttt{h}, which says
$17\le\card S$, into the hypothesis $\binom{4+2}{2}+2\le\card S$ of \lean{main} at $n=4$, and Lean accepts the
conclusion \mbox{\texttt{C.HasNGon (4 + 3) S}} as \mbox{\texttt{C.HasNGon 7 S}} because $4+3$ evaluates to $7$.
\end{proof}

\begin{remark}[comparison with Theorem~\ref{mb:thm:baek}; scope]\label{sa:rem:baek-seven}
(1) Lines 1--2 of \dirn{BaekSeven.out} are the output of \mbox{\texttt{\#check @main}}. They say that \lean{main}
takes, in this order, an implicit type $\gamma$ in an arbitrary universe (\mbox{\texttt{Type u\_1}}), an instance
\mbox{\texttt{LinearOrder}~$\gamma$}, $n\in\N$, \mbox{\texttt{C : Config}~$\gamma$} and
\mbox{\texttt{S : Finset}~$\gamma$}, and that it derives from the hypothesis
\mbox{\texttt{(n + 2).choose 2 + 2} $\le$ \texttt{S.card}}, that is $\binom{n+2}{2}+2\le\card S$, the conclusion
\mbox{\texttt{C.HasNCap 4 S}} $\lor$ \mbox{\texttt{C.HasNGon (n + 3) S}}. This is the statement of
Theorem~\ref{mb:thm:baek}, with the same notions; it is also the declaration at
\dirn{ErdosSzekeres/Vendor/ETV/Main/Main.lean:87} together with the variables declared at line~26 of that file. The
cap length $4$ is fixed; only the gon size $n+3$ and the threshold depend on $n$.

(2) Theorem~\ref{sa:thm:baek-seven} is the case $n=4$. The only other difference is the universe: \lean{baek_seven}
is stated for $\gamma$ in the lowest universe \texttt{Type}, \lean{main} for types in every universe. Read
mathematically, both statements are about an arbitrary linearly ordered set.

(3) The theorem concerns abstract configurations; it mentions no points of the plane, and
\dirn{BaekSeven.lean} contains no bridge from point sets to configurations. The records of \dirn{ES7SPLIT} say:
``The step from the abstract form to real point sets is not Lean-checked for $n = 7$; the repository's bridge covers
the octagon case only. Here that step rests on the published paper.'' (\dirn{ES7SPLIT/README.md}, lines 78--79, in
the rigor paragraph of its section~2; the status paragraph of \dirn{ES7SPLIT/proofs/split_theorem.md}, lines 13--15,
says the same.) These records were committed in \dirn{9793c06e0} (2~November 2026, 01:03 +0900) and predate
\lean{MirrorOdd.capacity16} (\dirn{MIRRORETV/lean/MirrorMain.lean:1529}, added in \dirn{83c80f38f} at 03:41 the same
day), a geometric form of the same instance of Theorem~\ref{mb:thm:baek}: a configuration of $n$ points in the sense
of Definition~\ref{mm:def:config} with no $4$-cap and no $7$ points in convex position has $n\le16$. That theorem is
Lean-checked: the Lean proof of Theorem~\ref{mb:thm:seven} uses it, and the axioms printed for that theorem cover it.
Whether it covers the use of Baek's theorem in the $z$-split arguments of \dirn{ES7SPLIT} is not examined here; the
theorems of \dirn{ES7SPLIT} remain written proofs. The source of Baek's theorem is J.~Baek, \emph{On the
Erd\H{o}s--Tuza--Valtr conjecture}, Eur.\ J.\ Combin.\ 124 (2025) 104085; arXiv:2206.04260.
\end{remark}

%----------------------------------------------------------------------------------------------------------------------
\subsection{How the arithmetic is used}\label{sa:ssec:use}

The following lemma is the modelling of the header comment of \dirn{SizeArith.lean}, proved here as a written step.
For $c\in\R^2$ and $\theta\in\R$, $\rho_{c,\theta}$ is the rotation of Definition~\ref{fd:def:rot}; in complex
notation it is $z\mapsto c+e^{i\theta}(z-c)$. For a map $f$ we write $f^{[i]}$ for its $i$-th iterate. The names
\lean{HalfTurn34}, \lean{Threefold33}, \lean{Fivefold35}, \lean{MirrorOneAxis33}, \lean{MirrorNoAxis34} and
\lean{Fourfold32} denote the case classes of Definition~\ref{sr:def:classes}, in the namespace
\lean{ErdosSzekeres.SymReduction}.

\begin{lemma}[orbit counts; written step]\label{sa:lem:model}
Let $S\subseteq\R^2$ be a finite set and $n=\card S$.
\begin{enumerate}[label=(\alph*)]
\item \emph{Rotations.} Let $c\in\R^2$, let $s\in\N$ with $0<s$, let $\rho=\rho_{c,2\pi/s}$, and assume
  $\mathrm{Invariant}\,\rho\,S$ \textup{(}Definition~\ref{fd:def:inv}\textup{)}. For $y\in\R^2$ put
  $O(y)=\{\rho^{[i]}(y):0\le i<s\}$. For every $y\in S\setminus\{c\}$ the set $O(y)$ is contained in
  $S\setminus\{c\}$, has exactly $s$ elements and is mapped into itself by $\rho$; any two of these sets are equal or
  disjoint; and $n=s\cdot k+e$, where $k$ is the number of distinct sets $O(y)$ with $y\in S\setminus\{c\}$, and $e=1$
  if $c\in S$, $e=0$ otherwise. If moreover $s$ is even, $S$ is in general position and $S\not\subseteq\{c\}$, then
  $e=0$.
\item \emph{Reflections.} Let $a,u\in\R^2$ with $u_0^2+u_1^2=1$, let $\sigma=\mathrm{reflLine}_{a,u}$, that is
  \[
  \sigma(p)=\bigl(a_0+2du_0-(p_0-a_0),\;a_1+2du_1-(p_1-a_1)\bigr),\qquad d=(p_0-a_0)u_0+(p_1-a_1)u_1,
  \]
  the map of the classes \lean{MirrorOneAxis33} and \lean{MirrorNoAxis34} of Definition~\ref{sr:def:classes}, and
  let $L=\{a+tu:t\in\R\}$. Assume $\mathrm{Invariant}\,\sigma\,S$ and that $S$ is in general position. Then
  $S\setminus L$ is the disjoint union of $k$ two-element sets of the form $\{p,\sigma(p)\}$, each mapped into itself
  by $\sigma$, and $n=2k+b$ with $b=\card{S\cap L}\le2$.
\item \emph{Deleting.} Let $T\subseteq S$. If $S$ is in general position, so is $T$; and
  $\mathrm{HasConvexNGon}\,7\,T$ implies $\mathrm{HasConvexNGon}\,7\,S$. In \textup{(a)}, a union of sets $O(y)$,
  together with $c$ or not when $c\in S$, is mapped into itself by $\rho$; in \textup{(b)}, a union of sets
  $\{p,\sigma(p)\}$ and of points of $S\cap L$ is mapped into itself by $\sigma$.
\end{enumerate}
\end{lemma}

\begin{proof}
(a) Put $\omega=e^{2\pi i/s}$, so that $\rho(z)=c+\omega(z-c)$. Then $\rho(c)=c$, and $\rho$ is injective because
$\omega\ne0$; so every iterate $\rho^{[i]}$ is injective. Hence $\rho$ maps $S\setminus\{c\}$ into itself: for
$y\in S\setminus\{c\}$ we have $\rho(y)\in S$, and $\rho(y)=c=\rho(c)$ would give $y=c$. By induction on $i$,
$\rho^{[i]}(y)\in S\setminus\{c\}$ for all $i$, so $O(y)\subseteq S\setminus\{c\}$. Let $y\in S\setminus\{c\}$. By
Lemma~\ref{ll:lem:minperiod} (with $m=s$), the minimal period of $y$ under $\rho$ is $s$: $\rho^{[s]}(y)=y$ and
$\rho^{[i]}(y)\ne y$ for $0<i<s$. The points $\rho^{[i]}(y)$, $0\le i<s$, are pairwise distinct: if
$\rho^{[i]}(y)=\rho^{[j]}(y)$ with $0\le i<j<s$, then $\rho^{[i]}\bigl(\rho^{[j-i]}(y)\bigr)=\rho^{[i]}(y)$, and
injectivity of $\rho^{[i]}$ gives $\rho^{[j-i]}(y)=y$ with $0<j-i<s$, a contradiction. So $\card{O(y)}=s$. Further,
$\rho(\rho^{[i]}(y))=\rho^{[i+1]}(y)\in O(y)$ for $i+1<s$, and $\rho(\rho^{[s-1]}(y))=\rho^{[s]}(y)=y\in O(y)$; so
$\rho$ maps $O(y)$ into itself. Consequently, for $x\in O(y)$ all points $\rho^{[i]}(x)$ lie in $O(y)$, so
$O(x)\subseteq O(y)$, and $O(x)=O(y)$ because both sets have $s$ elements. If $z\in O(y)\cap O(y')$, then
$O(y)=O(z)=O(y')$. Every $y\in S\setminus\{c\}$ lies in $O(y)$ (take $i=0$). Therefore the distinct sets $O(y)$ form a
partition of $S\setminus\{c\}$ into $k$ blocks of size $s$, and $n=\card{S\setminus\{c\}}+e=s\cdot k+e$. If $s$ is
even, $S$ is in general position and some $p\in S$ differs from $c$, then Lemma~\ref{lg:lem:centre} (with $m=s$)
gives $c\notin S$, that is, $e=0$.

(b) Write $\langle x,y\rangle=x_0y_0+x_1y_1$, so that $d=\langle p-a,u\rangle$ and $\sigma(p)-a=2du-(p-a)$.
\emph{$\sigma$ is an involution:} for $p'=\sigma(p)$ we get $\langle p'-a,u\rangle=2d\langle u,u\rangle-d=d$, hence
$\sigma(p')-a=2du-(p'-a)=2du-2du+(p-a)=p-a$, that is, $\sigma(\sigma(p))=p$; in particular $\sigma$ is injective.
\emph{The fixed points of $\sigma$ are the points of $L$:} $\sigma(p)=p$ means $2du=2(p-a)$, so $p=a+du\in L$;
conversely, for $p=a+tu$ we have $d=t\langle u,u\rangle=t$ and $\sigma(p)=a+2tu-tu=p$. (This is also shown at the
end of Definition~\ref{sr:def:classes}.) Now let $p\in S\setminus L$. Then $\sigma(p)\in S$ and $\sigma(p)\ne p$; and
$\sigma(p)\notin L$, since otherwise $\sigma(\sigma(p))=\sigma(p)$, that is, $p=\sigma(p)$. So $\{p,\sigma(p)\}$ is a
two-element subset of $S\setminus L$, mapped into itself by $\sigma$ because $\sigma(\sigma(p))=p$. If two such sets
$\{p,\sigma(p)\}$ and $\{q,\sigma(q)\}$ meet, then $q$ or $\sigma(q)$ lies in $\{p,\sigma(p)\}$, and applying
$\sigma$ shows that both $q$ and $\sigma(q)$ do; so the two sets are equal. Hence the distinct sets $\{p,\sigma(p)\}$
partition $S\setminus L$, $\card{S\setminus L}=2k$, and $n=2k+b$ with $b=\card{S\cap L}$. Finally $b\le2$: if
$x,y,z\in S\cap L$ were pairwise distinct, all differences of two of them would be real multiples of $u$, so the
linear span of these differences would have dimension at most $1$ and $\{x,y,z\}$ would be collinear, contrary to
general position (Definition~\ref{fd:def:gp}). This is the fact that at most two points of a set in general position
lie on a line.

(c) The first two claims are Lemma~\ref{fd:lem:mono}(b) and~(c). A union of sets each of which is mapped into itself
by a map is mapped into itself by that map; $\{c\}$ is mapped into itself by $\rho$ since $\rho(c)=c$, and each point
of $S\cap L$ is fixed by $\sigma$ by (b).
\end{proof}

\begin{remark}[How the arithmetic is used]\label{sa:rem:use}
This remark is a written step; it is not Lean-checked. It follows the reduction re-derived by the auditor in
\dirn{AUDIT3/reports/A1_reduction.md} (Summary and trust-table rows 2--9) and shows where each statement of this
section enters. Let $S$ be a finite set in general position with $n=\card S\ge33$.

(1) \emph{Rotations.} Let $c\in\R^2$, $s\ge2$, $\rho=\rho_{c,2\pi/s}$, and let $S$ be mapped into itself by $\rho$.
By Lemma~\ref{sa:lem:model}(a), $n=sk+e$ with $e\le1$, and $e=0$ for even $s$ (as $n\ge2$, $S\not\subseteq\{c\}$). By
Lemma~\ref{sa:lem:model}(c), a union $T$ of some of the sets $O(y)$ is in general position, mapped into itself by
$\rho$ and does not contain $c$.
\begin{itemize}
\item $s=2$: as $\cos\pi=-1$ and $\sin\pi=0$, $\rho=\rho_{c,\pi}$ is the half-turn $p\mapsto2c-p$ (Lean
  \mbox{\texttt{halfTurn c}}), and $n=2k$. Proposition~\ref{sa:prop:rot}(e) (or~\ref{sa:prop:assembled}(a)) gives
  $k\ge17$ and $n\ge34$. A union $T$ of $17$ pairs has $34$ points: these are the hypotheses of \lean{HalfTurn34}
  (Definition~\ref{sr:def:classes}) for $T$ and $c$.
\item $s=3$: Proposition~\ref{sa:prop:rot}(a) (or~\ref{sa:prop:assembled}(b)) gives $k\ge11$; $11$ orbits give
  $33$ points without $c$, invariant under $\rho_{c,2\pi/3}$: the hypotheses of \lean{Threefold33}.
\item $s=5$: Proposition~\ref{sa:prop:rot}(b) gives $n\ge35$ and $k\ge7$ (Proposition~\ref{sa:prop:assembled}(c)
  gives $n\ge35$ only); $7$ orbits give $35$ points without $c$: the hypotheses of \lean{Fivefold35}. (The comment of
  \lean{A1.c5_reduce}, ``$T_{C5}$ applies directly'', refers to the fivefold case theorem of the records, which admits
  at least $35$ points and the centre; the Lean class asks for exactly $35$ points without the centre, which $k\ge7$
  supplies.)
\item $s=4$: Proposition~\ref{sa:prop:rot}(c) gives $n=4k\ge36$, $k\ge8$ and $2k\ge17$
  (Proposition~\ref{sa:prop:assembled}(f) gives $2k\ge17$ only). Rotations about $c$ compose by adding angles (in
  complex notation, $c+e^{i\alpha}\bigl(c+e^{i\beta}(z-c)-c\bigr)=c+e^{i(\alpha+\beta)}(z-c)$), so
  $\rho^{[2]}=\rho_{c,\pi}$ is the
  half-turn; Lemma~\ref{sa:lem:model}(a) for it gives $n=2k'$, and $k'=2k\ge17$, so the case $s=2$ applies. (Eight of
  the $k$ orbits would give the $32$ points of the class \lean{Fourfold32}. Among the reduction theorems and case
  lemmas stated in this document, only Theorem~\ref{sr:thm:variant} and the case lemma
  \lean{SymReduction.case_fourfold} that its proof uses take this class as a hypothesis (so do some checks of that
  theorem in the review file \dirn{SYMLEAN/review/VariantChecks.lean}, among them its group form
  \lean{SymReduction.VariantChecks.group_form}); a same-text copy, \lean{LayerLemma.Fourfold32}, is the hypothesis of
  Theorem~\ref{ec:thm:fourfold} and of its layer form \lean{LayerLemma.fourfold_heptagon_of_erase}, which its proof
  uses;
  item~(5) uses the class for the size table.)
\item $s=6$: Proposition~\ref{sa:prop:rot}(d) gives $2k\ge11$ and $3k\ge17$. Here $\rho^{[2]}=\rho_{c,2\pi/3}$ has
  $2k$ orbits of size $3$ and $\rho^{[3]}=\rho_{c,\pi}$ has $3k$ pairs (Lemma~\ref{sa:lem:model}(a) for these
  rotations, with $n=6k$), so both the threefold and the half-turn case apply.
\item General $s$: by Lemma~\ref{sa:lem:split}, $2\mid s$, $3\mid s$, $5\mid s$, or $s\ge7$. If $s=mt$ with
  $m\in\{2,3,5\}$, then $\rho^{[t]}=\rho_{c,2\pi t/s}=\rho_{c,2\pi/m}$ maps $S$ into itself, and
  Lemma~\ref{sa:lem:model}(a) for it gives $n=mk'+e$ with the same $e$; from $mk'=n-e=mtk$ we get $k'=tk$.
  Proposition~\ref{sa:prop:order} is the arithmetic of this step: $tk\ge17$ pairs for even $s$ (where $e=0$),
  $tk\ge11$ orbits of size $3$ if $3\mid s$, and $35\le5(tk)+e$ if $5\mid s$, hence $5tk\ge34$ and $tk\ge7$. The cases
  $s=2,3,5$ then apply to $\rho^{[t]}$. If $s\ge7$, Lemma~\ref{sa:lem:big-order} (with $p=s$) gives $k\ge1$, so $S$
  contains a set $O(y)$ with $s\ge7$ points. They all have the same distance from $c$ as $y$, because
  $\abs{\omega(z-c)}=\abs{z-c}$, so $O(y)$ is in convex position by Lemma~\ref{sr:lem:concyclic}; any $7$ of its
  points are in convex position (Lemma~\ref{fd:lem:mono}(a)), so $\mathrm{HasConvexNGon}\,7\,S$
  (Definition~\ref{fd:def:ngon}). No case class is needed here.
\end{itemize}

(2) \emph{Reflections.} Let $S$ be mapped into itself by $\sigma=\mathrm{reflLine}_{a,u}$ with $u_0^2+u_1^2=1$. By
Lemma~\ref{sa:lem:model}(b), $n=2k+b$ with $b\le2$. If $b=0$, Proposition~\ref{sa:prop:mirror}(a)
(or~\ref{sa:prop:assembled}(d)) gives $k\ge17$, and $17$ pairs form a $34$-point set that is in general position and
mapped into itself by $\sigma$ (Lemma~\ref{sa:lem:model}(c)) and has no fixed point: the hypotheses of
\lean{MirrorNoAxis34}. If $b\in\{1,2\}$, Proposition~\ref{sa:prop:mirror}(b) (or~\ref{sa:prop:assembled}(e)) gives
$k\ge16$, and $16$ pairs together with one point $x\in S\cap L$ form a $33$-point set that is in general position and
mapped into itself by $\sigma$ (Lemma~\ref{sa:lem:model}(c)) and whose only fixed point is $x$: the hypotheses of
\lean{MirrorOneAxis33}. (The fixed points of $\sigma$ are the points of $L$, as shown in the proof of
Lemma~\ref{sa:lem:model}(b), and the pairs lie in $S\setminus L$.)

(3) \emph{All symmetries.} By Theorem~\ref{sr:thm:structure}(b), a Euclidean isometry $g\ne\mathrm{id}$ that maps $S$
into itself is either a rotation, in complex notation $z\mapsto c+a(z-c)$ with $a\ne1$ and $a^m=1$ for some $m>0$, or
a map $\mathrm{reflLine}_{A,u}$ with $u_0^2+u_1^2=1$. In the first case let $s\ge2$ be the order of $a$. Then $a$ is a
primitive $s$-th root of unity, so $e^{2\pi i/s}=a^j$ for some $j\in\N$, because a primitive $s$-th root of unity
generates all $s$-th roots of unity. As $g^{[j]}(z)=c+a^j(z-c)$ (induction on $j$), $\rho_{c,2\pi/s}=g^{[j]}$, which
maps $S$ into itself (here $c$ is read as a point); so item~(1) applies. In the second case item~(2) applies.
Affine symmetries are reduced to Euclidean isometries by Theorem~\ref{ar:thm:transfer}, which gives
Theorem~\ref{ar:thm:main}.

(4) \emph{Status.} Items (1)--(3) derive $\mathrm{HasConvexNGon}\,7\,S$ from the five case classes \lean{HalfTurn34},
\lean{Threefold33}, \lean{Fivefold35}, \lean{MirrorOneAxis33} and \lean{MirrorNoAxis34}, used as hypotheses. Except
for the orders $s\ge7$ in item~(1), which need no class, items (1) and~(2) exhibit a subset $T$ that satisfies the
hypotheses of one class, and the class then gives $\mathrm{HasConvexNGon}\,7\,T$, hence
$\mathrm{HasConvexNGon}\,7\,S$ by Lemma~\ref{sa:lem:model}(c); item~(3) reduces every Euclidean isometry other than
the identity to items (1) and~(2), and affine symmetries to Euclidean isometries. The Lean-checked form of this
implication is Theorem~\ref{sr:thm:main} for Euclidean isometries, with Theorem~\ref{ar:thm:main} for affine
symmetries. The Lean proof of Theorem~\ref{sr:thm:main} passes to a power of the rotation that is the half-turn or has
odd prime order, instead of using Lemma~\ref{sa:lem:split}, and it does its own arithmetic. The two mirror classes
are Lean-checked theorems (Theorem~\ref{mb:thm:classes}). The classes \lean{HalfTurn34}, \lean{Threefold33} and
\lean{Fivefold35} are not proved in Lean or in this document (Theorem~\ref{mb:thm:main3} is conditional on them);
the records that supply them are described in the remark that follows Definition~\ref{sr:def:classes}. Of the
three ingredients, the modelling (Lemma~\ref{sa:lem:model}), the case classes and the arithmetic, the files of this
section check only the arithmetic.

(5) \emph{The size table.} Proposition~\ref{sa:prop:table} is the arithmetic of the upper bounds in the size table of
\dirn{WRITEUP/SYMMETRIC.md}, \S2 (trust-table row~9 of the A1 report). For a set without $7$ points in convex position,
the case classes bound $k$ through item~(1) or~(2), and the parts of Proposition~\ref{sa:prop:table} turn these
bounds into bounds on $n$. Here $S$ need not have $33$ points: if the number of orbits or pairs to which a class is
applied below ($k$, or $2k$ for $\rho^{[2]}$ when $s=4$ or $s=6$) reached the number that the class requires, the
subset $T$ built in item~(1) or~(2) would satisfy the hypotheses of the class, which would then give $7$ points of
$S$ in convex position; for even $s$, $e=0$ as soon as $k\ge1$ (Lemma~\ref{sa:lem:model}(a)), and $k=0$ gives
$n=e\le1$. For threefold sets, \lean{Threefold33} gives $k\le10$, so $n\le31$ by part~(a), and
$n\ne32$ in any case by part~(b). For fivefold sets, \lean{Fivefold35} gives $k\le6$, so $n=5k+e<35$ and $n\le31$ by
part~(c); $32$, $33$ and $34$ are not fivefold sizes by part~(d). For fourfold sets, \lean{Fourfold32} gives $k\le7$,
so $n\le28$ by part~(e), while \lean{HalfTurn34} alone gives only $2k\le16$ and $n\le32$ by part~(f). For sixfold
sets, \lean{Threefold33} applied to the $2k$ orbits of size $3$ gives $2k\le10$, so $n\le30$ by part~(g). For
half-turns, \lean{HalfTurn34} gives $k\le16$, so $n\le32$ by part~(h). For reflections, the two mirror classes give
$k\le16$ if $b=0$ and $k\le15$ if $b\ge1$, so $n\le32$ by part~(i).
Except for the mirror row, these bounds are conditional on the classes named, which are hypotheses here: the rows
for half-turns, threefold and sixfold sets and the bound $32$ of part~(f) on \lean{HalfTurn34} or \lean{Threefold33},
the fourfold bound $28$ on \lean{Fourfold32}, and the fivefold bound $31$ on \lean{Fivefold35}. For the last two,
Lean-checked conditional forms exist that include the geometric modelling: Theorem~\ref{ec:thm:fourfold}, with the
same-text copy \lean{LayerLemma.Fourfold32} as hypothesis, gives $\mathrm{HasConvexNGon}\,7\,S$ for every set $S$ in
general position with at least $29$ points that is mapped into itself by $\rho_{c,2\pi/4}$, so such a set without $7$
points in convex position has at most $28$ points; Theorem~\ref{ec:thm:fivefold}, with the same-text copy
\lean{LayerLemma.Fivefold35}, does the same with $\rho_{c,2\pi/5}$, at least $32$ points and the bound $31$. For
reflections the bound $n\le32$ is Lean-checked with no class hypothesis (Theorem~\ref{mb:thm:refl32}(b) for
$\mathrm{reflLine}_{a,u}$, derived from Theorem~\ref{mb:thm:seven}(b)), and the two mirror classes are themselves
Lean-checked theorems (Theorem~\ref{mb:thm:classes}). The matching lower bounds rest on example sets of the records.
None of this is part of the files of this section.
\end{remark}

{\raggedright
\paragraph{Coverage.}
Covered, by label. In \dirn{AUDIT3/A1_reduction/SizeArith.lean}: \lean{A1.c3_reduce}, \lean{A1.c5_reduce},
\lean{A1.c4_reduce}, \lean{A1.c6_reduce}, \lean{A1.c2_reduce} (Proposition~\ref{sa:prop:rot});
\lean{A1.mirror_reduce_b0}, \lean{A1.mirror_reduce_b12} (Proposition~\ref{sa:prop:mirror});
\lean{A1.even_order_reduce}, \lean{A1.mult3_order_reduce}, \lean{A1.mult5_order_reduce}
(Proposition~\ref{sa:prop:order}); \lean{A1.order_split} (Lemma~\ref{sa:lem:split}); \lean{A1.table_c3},
\lean{A1.c3_no32}, \lean{A1.table_c5}, \lean{A1.c5_gap}, \lean{A1.table_c4}, \lean{A1.table_c4_only_c2},
\lean{A1.table_c6}, \lean{A1.table_c2}, \lean{A1.table_mirror} (Proposition~\ref{sa:prop:table}). In
\dirn{AUDIT3/A1_reduction/verify_A1-F1/AssembledSizes.lean}: \lean{halfturn_pairs}, \lean{threefold_orbits},
\lean{fivefold_size}, \lean{mirror_b0}, \lean{mirror_b12}, \lean{fourfold_via_halfturn}
(Proposition~\ref{sa:prop:assembled}); \lean{big_order_orbit} (Lemma~\ref{sa:lem:big-order}, with
Remark~\ref{sa:rem:big-order}); the anonymous example of line~30 (Remark~\ref{sa:rem:control}). In
\dirn{ES7SPLIT/review/equivalence/BaekSeven.lean}: \lean{baek_seven} (Theorem~\ref{sa:thm:baek-seven}); the vendored
theorem \lean{main} that it instantiates is the cited Theorem~\ref{mb:thm:baek}. That is all $28$ theorems of the
three files and the one example.

Not covered, deliberately: the \texttt{\#print axioms} commands (\dirn{SizeArith.lean}, lines 84--103;
\dirn{AssembledSizes.lean}, lines 33--39; \dirn{BaekSeven.lean}, line~12) and the \texttt{\#check} command
(\dirn{BaekSeven.lean}, line~11), which produce the logs named in the Sources paragraph and are not results; and the
comments of the files, among them the modelling in the header of \dirn{SizeArith.lean}, which these files do not
formalise and which is treated here as the written step Lemma~\ref{sa:lem:model}. Lemma~\ref{sa:lem:model} and
Remark~\ref{sa:rem:use} are not declarations of the three files. Out of scope: the scripts and other logs of
\dirn{AUDIT3/A1_reduction} and its subdirectories, and the other files of \dirn{ES7SPLIT/review/equivalence}; none of
them contains a Lean declaration.\par}


\appendix
\begingroup\esmapbreaks
% appendix_map.tex -- generated by tools/make_appendix.py from sections/*.map.tsv; do not edit by hand.
\section{Statement map}\label{app:map}

Each row lists a numbered statement of this document, the Lean declarations that it restates, each with
its file and line at commit \texttt{7e32ee259}, and a note. For a remark, the Lean column names the
declarations that the remark discusses; a remark does not restate a Lean statement. A path that begins with
\dirn{ErdosSzekeres/}, \dirn{docs/}, \dirn{experiments/} or \dirn{scripts/} is relative to the repository
root; every other path is relative to \dirn{cloud/es7-2026-10-19/results}. A statement introduced only for
the written proofs is marked ``written step'', one that restates Mathlib definitions or facts is marked
``Mathlib'', and a remark without a Lean declaration has a dash. Labelled remarks appear where a section
lists them. The notes are those of the section maps.

{\footnotesize
\begin{longtable}{@{}>{\raggedright\arraybackslash}p{0.10\linewidth}>{\raggedright\arraybackslash}p{0.49\linewidth}>{\raggedright\arraybackslash}p{0.35\linewidth}@{}}
\toprule
Statement & Lean declaration; file and line & Note\\
\midrule
\endfirsthead
\toprule
Statement & Lean declaration; file and line & Note\\
\midrule
\endhead
\bottomrule
\endlastfoot
\multicolumn{3}{@{}l}{\normalsize\textbf{Section~\ref{fd:sec}}}\\[2pt]
Def.~\ref{fd:def:point} & \hangindent=1em\hangafter=1 \lean{ErdosSzekeres.Point}\newline\dirn{ErdosSzekeres/Definitions.lean:9} & Finset vs Set conventions\\
Rem.~\ref{fd:rem:norm} & -- & sup norm of the Lean point type; the only topological argument is the compactness step in Lemma~\ref{fd:lem:separation} (Heine-Borel, extreme value theorem; Rudin cited)\\
Def.~\ref{fd:def:orient} & \hangindent=1em\hangafter=1 \lean{ErdosSzekeres.orient}\newline\dirn{ErdosSzekeres/Orientation.lean:9} & any commutative ring\\
Def.~\ref{fd:def:conv} & Mathlib & Mathlib definitions restated: Convex, \texttt{convexHull}, \texttt{AffineIndependent}, \texttt{AffineMap} (map with $f(w+p)=$ \texttt{f.\allowbreak{}linear} $w+fp$), \texttt{AffineEquiv}; no repository declaration\\
Def.~\ref{fd:def:gp} & \hangindent=1em\hangafter=1 \lean{ErdosSzekeres.GeneralPosition}\newline\dirn{ErdosSzekeres/Definitions.lean:12} & uses Mathlib Collinear (rank of \texttt{vectorSpan} at most 1)\\
Def.~\ref{fd:def:cp} & \hangindent=1em\hangafter=1 \lean{ErdosSzekeres.ConvexPosition}\newline\dirn{ErdosSzekeres/Definitions.lean:17} & Mathlib \texttt{ConvexIndependent} of the inclusion\\
Def.~\ref{fd:def:ngon} & \hangindent=1em\hangafter=1 \lean{ErdosSzekeres.HasConvexNGon}\newline\dirn{ErdosSzekeres/Definitions.lean:21} & also defines K-free; no emptiness condition\\
Rem.~\ref{fd:rem:ngon} & -- & \texttt{HasConvexNGon} 0 S always holds (no set is 0-free); K-free is upward closed; \texttt{HasConvexNGon} n S implies $n\le|S|$\\
Def.~\ref{fd:def:forces} & \hangindent=1em\hangafter=1 \lean{ErdosSzekeres.ForcesConvexNGon}\newline\dirn{ErdosSzekeres/Definitions.lean:25} & \\
Def.~\ref{fd:def:inv} & \hangindent=1em\hangafter=1 \lean{ErdosSzekeres.LayerLemma.Invariant}\newline\dirn{LAYERLEAN/LayerLemma.lean:105}\par \hangindent=1em\hangafter=1 \lean{ErdosSzekeres.SymReduction.Invariant}\newline\dirn{SYMLEAN/SymmetricReduction.lean:271}\par \hangindent=1em\hangafter=1 \lean{ErdosSzekeres.OrbitConsistency.Invariant}\newline\dirn{LEANEXT/orbit/OrbitConsistency.lean:47} & maps into; textually identical copies\\
Lem.~\ref{fd:lem:inv-image} & \hangindent=1em\hangafter=1 \lean{ErdosSzekeres.LayerLemma.iterate_mem}\newline\dirn{LAYERLEAN/LayerLemma.lean:52} & part (c) is \texttt{iterate\_\allowbreak{}mem}, stated as in Lean for an arbitrary set X (Lean: arbitrary type); (a),(b) written step for maps of the plane (the argument of (a) is the step at line 165 of the proof of \texttt{LayerLemma.\allowbreak{}map\_\allowbreak{}mem\_\allowbreak{}hullVerts}, \texttt{LAYERLEAN/\allowbreak{}LayerLemma.\allowbreak{}lean:162})\\
Def.~\ref{fd:def:rot} & \hangindent=1em\hangafter=1 \lean{ErdosSzekeres.LayerLemma.rotAbout}\newline\dirn{LAYERLEAN/LayerLemma.lean:298}\par \hangindent=1em\hangafter=1 \lean{ErdosSzekeres.SymReduction.rotAbout}\newline\dirn{SYMLEAN/SymmetricReduction.lean:400}\par \hangindent=1em\hangafter=1 \lean{ErdosSzekeres.OrbitConsistency.rotAbout}\newline\dirn{LEANEXT/orbit/OrbitConsistency.lean:93} & textually identical copies; counterclockwise\\
Rem.~\ref{fd:rem:angles} & -- & angle $2\pi/m$ with m natural (cast explicit; $m=0$ gives the identity); Sixfold66 and Sevenfold70 use real numerals 6 and 7\\
Lem.~\ref{fd:lem:rot} & \hangindent=1em\hangafter=1 \lean{ErdosSzekeres.LayerLemma.rotAbout_eq_crot}\newline\dirn{LAYERLEAN/LayerLemma.lean:302}\par \hangindent=1em\hangafter=1 \lean{ErdosSzekeres.LayerLemma.rotAff}\newline\dirn{LAYERLEAN/LayerLemma.lean:311}\par \hangindent=1em\hangafter=1 \lean{ErdosSzekeres.LayerLemma.rotAbout_iterate}\newline\dirn{LAYERLEAN/LayerLemma.lean:465}\par \hangindent=1em\hangafter=1 \lean{ErdosSzekeres.LayerLemma.rotAbout_injective}\newline\dirn{LAYERLEAN/LayerLemma.lean:319} & (a) \texttt{rotAbout\_\allowbreak{}eq\_\allowbreak{}crot}; (b) \texttt{rotAff} (affine packaging: linear part $R_{\theta},\rho(p+v)=\rho(p)+R_{\theta}v$); (c) \texttt{rotAbout\_\allowbreak{}iterate}; (d) \texttt{rotAbout\_\allowbreak{}injective} (injectivity only); written steps: the matrix form $\rho(p)=c+R_{\theta}\ (p-c)=R_{\theta}p+(c-R_{\theta}c)$ of Definition~\ref{fd:def:rot} and $\rho(c)=c$ in (b) (the latter proved in Lean inside \texttt{invariant\_\allowbreak{}erase\_\allowbreak{}centre}, \texttt{LAYERLEAN/\allowbreak{}LayerLemma.\allowbreak{}lean:325}), composition law and $\rho_{0}=\mathrm{id}$ in (c), bijectivity and inverse in (d), (e), (f) (counterparts named in Coverage)\\
Lem.~\ref{fd:lem:orient} & \hangindent=1em\hangafter=1 \lean{ErdosSzekeres.orient_swap}\newline\dirn{ErdosSzekeres/Orientation.lean:12}\par \hangindent=1em\hangafter=1 \lean{ErdosSzekeres.orient_cycle}\newline\dirn{ErdosSzekeres/Orientation.lean:15}\par \hangindent=1em\hangafter=1 \lean{ErdosSzekeres.orient_translate}\newline\dirn{ErdosSzekeres/Orientation.lean:18} & (a) \texttt{orient\_\allowbreak{}swap}; (b) \texttt{orient\_\allowbreak{}cycle}; (c) \texttt{orient\_\allowbreak{}translate}; (d),(e) written step (real case of (d) in Remark~\ref{fd:rem:orient-lean}); any commutative ring\\
Rem.~\ref{fd:rem:orient-lean} & \hangindent=1em\hangafter=1 \lean{ErdosSzekeres.OrbitConsistency.orient3_perm}\newline\dirn{LEANEXT/orbit/OrbitConsistency.lean:178}\par \hangindent=1em\hangafter=1 \lean{ErdosSzekeres.OrbitConsistency.orient3}\newline\dirn{LEANEXT/orbit/OrbitConsistency.lean:168} & \texttt{orient3\_\allowbreak{}perm} is the case $R=$ real numbers of Lemma~\ref{fd:lem:orient}(d) (Lemma~\ref{oc:lem:perm})\\
Lem.~\ref{fd:lem:affine-orient} & \hangindent=1em\hangafter=1 \lean{ErdosSzekeres.SymReduction.orient_affine}\newline\dirn{SYMLEAN/AffineReduction.lean:1241}\par \hangindent=1em\hangafter=1 \lean{ErdosSzekeres.OrbitConsistency.orient_rotAbout_inv}\newline\dirn{LEANEXT/orbit/OrbitConsistency.lean:117} & main formula \texttt{orient\_\allowbreak{}affine} (g a map with $g(p)=Lp+w$ for all p; Lean names w as b); verbatim copy at \texttt{SYMLEAN/\allowbreak{}mirror/\allowbreak{}ReductionThreeClasses.\allowbreak{}lean:1267}; (i) \texttt{orient\_\allowbreak{}rotAbout\_\allowbreak{}inv} (stated for the \texttt{OrbitConsistency} copy of \texttt{rotAbout}); (ii) written step\\
Rem.~\ref{fd:rem:affine-orient-lean} & \hangindent=1em\hangafter=1 \lean{ErdosSzekeres.SymReduction.orient_affine}\newline\dirn{SYMLEAN/mirror/ReductionThreeClasses.lean:1267}\par \hangindent=1em\hangafter=1 \lean{ErdosSzekeres.OrbitConsistency.orient_rotAbout_inv}\newline\dirn{LEANEXT/orbit/OrbitConsistency.lean:117}\par \hangindent=1em\hangafter=1 \lean{ErdosSzekeres.Suk.orient_reflectY}\newline\dirn{ErdosSzekeres/Suk/Reflection.lean:47} & verbatim copy of \texttt{orient\_\allowbreak{}affine}; (i) is \texttt{orient\_\allowbreak{}rotAbout\_\allowbreak{}inv}, also cited in the head of Lemma~\ref{fd:lem:affine-orient} (Lemma~\ref{oc:lem:rot-inv}); the case $o=0,\phi=0$ of (ii), the reflection $(p0,p1)\to(p0,-p1)$, is \texttt{orient\_\allowbreak{}reflectY}\\
Lem.~\ref{fd:lem:collinear-mem} & written step & written step; Mathlib \texttt{collinear\_\allowbreak{}iff\_\allowbreak{}of\_\allowbreak{}mem} reproved\\
Lem.~\ref{fd:lem:collinear} & \hangindent=1em\hangafter=1 \lean{ErdosSzekeres.orient_eq_zero_iff_collinear}\newline\dirn{ErdosSzekeres/OrientationGeometry.lean:36}\par \hangindent=1em\hangafter=1 \lean{ErdosSzekeres.generalPosition_iff_orient_ne_zero}\newline\dirn{ErdosSzekeres/OrientationGeometry.lean:66} & (a) \texttt{orient\_\allowbreak{}eq\_\allowbreak{}zero\_\allowbreak{}iff\_\allowbreak{}collinear}; (b) \texttt{generalPosition\_\allowbreak{}iff\_\allowbreak{}orient\_\allowbreak{}ne\_\allowbreak{}zero}; includes degenerate triples\\
Lem.~\ref{fd:lem:hull} & written step & written step (standard facts on convex hulls)\\
Lem.~\ref{fd:lem:cp-hull} & \hangindent=1em\hangafter=1 \lean{ErdosSzekeres.convexPosition_iff}\newline\dirn{ErdosSzekeres/Basic.lean:15} & Lean proof is Mathlib \texttt{convexIndependent\_\allowbreak{}set\_\allowbreak{}iff\_\allowbreak{}notMem\_\allowbreak{}convexHull\_\allowbreak{}diff}; direct proof given\\
Lem.~\ref{fd:lem:mono} & \hangindent=1em\hangafter=1 \lean{ErdosSzekeres.ConvexPosition.mono}\newline\dirn{ErdosSzekeres/Basic.lean:11}\par \hangindent=1em\hangafter=1 \lean{ErdosSzekeres.GeneralPosition.mono}\newline\dirn{ErdosSzekeres/Basic.lean:6}\par \hangindent=1em\hangafter=1 \lean{ErdosSzekeres.HasConvexNGon.mono}\newline\dirn{ErdosSzekeres/Basic.lean:20}\par \hangindent=1em\hangafter=1 \lean{MirrorOdd.hasConvexNGon_of_le}\newline\dirn{MIRRORETV/lean/MirrorMain.lean:2182} & (a) \texttt{ConvexPosition.\allowbreak{}mono}; (b) \texttt{GeneralPosition.\allowbreak{}mono}; (c) \texttt{HasConvexNGon.\allowbreak{}mono}; (d) \texttt{MirrorOdd.\allowbreak{}hasConvexNGon\_\allowbreak{}of\_\allowbreak{}le}, also verbatim at \texttt{SYMLEAN/\allowbreak{}mirror/\allowbreak{}ReductionThreeClasses.\allowbreak{}lean:3741}; axioms of (d) through its use in the proof of \texttt{MirrorOdd.\allowbreak{}mirror\_\allowbreak{}odd\_\allowbreak{}general} (\texttt{MIRRORETV/\allowbreak{}lean/\allowbreak{}MirrorMain.\allowbreak{}lean:2498})\\
Lem.~\ref{fd:lem:ngon-basic} & \hangindent=1em\hangafter=1 \lean{ErdosSzekeres.HasConvexNGon.card_le}\newline\dirn{ErdosSzekeres/Basic.lean:25}\par \hangindent=1em\hangafter=1 \lean{ErdosSzekeres.ConvexPosition.hasConvexNGon}\newline\dirn{ErdosSzekeres/Basic.lean:31} & (a) \texttt{HasConvexNGon.\allowbreak{}card\_\allowbreak{}le}; (b) \texttt{ConvexPosition.\allowbreak{}hasConvexNGon}; part (a) is not printed in the axiom audit, and its evidence is of kind (a) of the rigor labels (Section~\ref{fd:sec}): build record \texttt{docs/\allowbreak{}VALIDATION.\allowbreak{}md:73}, the axiom audit as Lean output of the same library, text search\\
Lem.~\ref{fd:lem:forces-basic} & \hangindent=1em\hangafter=1 \lean{ErdosSzekeres.ForcesConvexNGon.mono}\newline\dirn{ErdosSzekeres/Basic.lean:47}\par \hangindent=1em\hangafter=1 \lean{ErdosSzekeres.forcesConvexNGon_iff_exact}\newline\dirn{ErdosSzekeres/Basic.lean:37} & (a) \texttt{ForcesConvexNGon.\allowbreak{}mono}; (b) \texttt{forcesConvexNGon\_\allowbreak{}iff\_\allowbreak{}exact}\\
Lem.~\ref{fd:lem:small} & \hangindent=1em\hangafter=1 \lean{ErdosSzekeres.convexPosition_of_card_three}\newline\dirn{ErdosSzekeres/Basic.lean:58} & last assertion of (c) only; (a),(b),(c) otherwise written step\\
Lem.~\ref{fd:lem:affine-cp} & \hangindent=1em\hangafter=1 \lean{ErdosSzekeres.SymReduction.convexPosition_image_affine}\newline\dirn{SYMLEAN/AffineReduction.lean:1263}\par \hangindent=1em\hangafter=1 \lean{ErdosSzekeres.SymReduction.generalPosition_image_affine}\newline\dirn{SYMLEAN/AffineReduction.lean:1278}\par \hangindent=1em\hangafter=1 \lean{ErdosSzekeres.SymReduction.generalPosition_image_affineEquiv}\newline\dirn{SYMLEAN/AffineReduction.lean:1290}\par \hangindent=1em\hangafter=1 \lean{ErdosSzekeres.SymReduction.hasConvexNGon_image_affine}\newline\dirn{SYMLEAN/AffineReduction.lean:1298}\par \hangindent=1em\hangafter=1 \lean{ErdosSzekeres.SymReduction.injective_of_affine}\newline\dirn{SYMLEAN/AffineReduction.lean:1305} & Lean forms of the forward directions: (b') \texttt{convexPosition\_\allowbreak{}image\_\allowbreak{}affine}; (c') \texttt{generalPosition\_\allowbreak{}image\_\allowbreak{}affine} and \texttt{generalPosition\_\allowbreak{}image\_\allowbreak{}affineEquiv}; (d') \texttt{hasConvexNGon\_\allowbreak{}image\_\allowbreak{}affine}; (e') \texttt{injective\_\allowbreak{}of\_\allowbreak{}affine} (converse of injective $\Rightarrow\det A\ne0$); (a) is Mathlib \texttt{AffineMap.\allowbreak{}image\_\allowbreak{}convexHull}; reverse directions of (b)-(d) and injective $\Rightarrow\det A\ne0$ written step; $h(p)=Lp+w$ (Lean names w as b); the affine equivalence of (c') is written $\phi$\\
Rem.~\ref{fd:rem:affine-cp-lean} & \hangindent=1em\hangafter=1 \lean{ErdosSzekeres.HasConvexNGon.linearEquiv}\newline\dirn{ErdosSzekeres/Suk/Reflection.lean:6}\par \hangindent=1em\hangafter=1 \lean{ErdosSzekeres.Suk.reflectY}\newline\dirn{ErdosSzekeres/Suk/Reflection.lean:34}\par \hangindent=1em\hangafter=1 \lean{ErdosSzekeres.Suk.reflectY_x}\newline\dirn{ErdosSzekeres/Suk/Reflection.lean:42}\par \hangindent=1em\hangafter=1 \lean{ErdosSzekeres.Suk.reflectY_y}\newline\dirn{ErdosSzekeres/Suk/Reflection.lean:43}\par \hangindent=1em\hangafter=1 \lean{ErdosSzekeres.Suk.reflectY_reflectY}\newline\dirn{ErdosSzekeres/Suk/Reflection.lean:44}\par \hangindent=1em\hangafter=1 \lean{ErdosSzekeres.Suk.generalPosition_reflectY}\newline\dirn{ErdosSzekeres/Suk/Reflection.lean:52}\par \hangindent=1em\hangafter=1 \lean{ErdosSzekeres.Suk.hasConvexNGon_of_reflectY}\newline\dirn{ErdosSzekeres/Suk/Reflection.lean:64} & main-library special cases of Lemma~\ref{fd:lem:affine-cp} (linear bijections; the reflection \texttt{reflectY} in both directions); \texttt{reflectY\_\allowbreak{}x}, \texttt{reflectY\_\allowbreak{}y} and \texttt{reflectY\_\allowbreak{}reflectY} are not printed in the axiom audit: evidence of kind (a) of the rigor labels (Section~\ref{fd:sec}): build record \texttt{docs/\allowbreak{}VALIDATION.\allowbreak{}md:73}, the axiom audit as Lean output of the same library, text search\\
Lem.~\ref{fd:lem:caratheodory} & written step & written step; Lean uses Mathlib \texttt{Caratheodory.\allowbreak{}minCardFinsetOfMemConvexHull}\\
Lem.~\ref{fd:lem:four-subsets} & \hangindent=1em\hangafter=1 \lean{ErdosSzekeres.Suk.convexPosition_of_subsets_card_le_four}\newline\dirn{ErdosSzekeres/Suk/FourPoint.lean:10}\par \hangindent=1em\hangafter=1 \lean{ErdosSzekeres.Suk.convexPosition_of_four_subsets}\newline\dirn{ErdosSzekeres/Suk/FourPoint.lean:40} & (a) \texttt{convexPosition\_\allowbreak{}of\_\allowbreak{}subsets\_\allowbreak{}card\_\allowbreak{}le\_\allowbreak{}four}; (b) \texttt{convexPosition\_\allowbreak{}of\_\allowbreak{}four\_\allowbreak{}subsets}; no general position assumed\\
Lem.~\ref{fd:lem:halfplane} & \hangindent=1em\hangafter=1 \lean{ErdosSzekeres.Suk.orient_affine_combination}\newline\dirn{ErdosSzekeres/Suk/SupportingLines.lean:6}\par \hangindent=1em\hangafter=1 \lean{ErdosSzekeres.Suk.convex_orient_positive}\newline\dirn{ErdosSzekeres/Suk/SupportingLines.lean:11}\par \hangindent=1em\hangafter=1 \lean{ErdosSzekeres.Suk.convex_orient_nonnegative}\newline\dirn{ErdosSzekeres/Suk/SupportingLines.lean:21} & (a) \texttt{orient\_\allowbreak{}affine\_\allowbreak{}combination}; (b) \texttt{convex\_\allowbreak{}orient\_\allowbreak{}positive}; (c) \texttt{convex\_\allowbreak{}orient\_\allowbreak{}nonnegative}\\
Lem.~\ref{fd:lem:strict-supports} & \hangindent=1em\hangafter=1 \lean{ErdosSzekeres.Suk.convexPosition_of_strict_supports}\newline\dirn{ErdosSzekeres/Suk/SupportingLines.lean:28} & \\
Lem.~\ref{fd:lem:support-sum} & \hangindent=1em\hangafter=1 \lean{ErdosSzekeres.Suk.orient_support_sum}\newline\dirn{ErdosSzekeres/Suk/SupportingLines.lean:40} & \\
Lem.~\ref{fd:lem:four-positive} & \hangindent=1em\hangafter=1 \lean{ErdosSzekeres.Suk.convexPosition_four_of_positive}\newline\dirn{ErdosSzekeres/Suk/SupportingLines.lean:47} & twelve support inequalities as in Lean\\
Lem.~\ref{fd:lem:pos-triples} & \hangindent=1em\hangafter=1 \lean{ErdosSzekeres.Suk.convexPosition_of_positive_triples}\newline\dirn{ErdosSzekeres/Suk/OrientationClique.lean:9} & requires $4\le\operatorname{card}s$ and p injective on all of $\iota$; Lean uses decide only for side facts (order facts among 0,1,2,3 in Fin 4 and \texttt{Finset.\allowbreak{}univ} $=$ \{0,1,2,3\} in Fin 4) inside the proof of a statement that is not itself a finite computation: not computer-assisted\\
Rem.~\ref{fd:rem:pos-triples} & -- & finiteness and injectivity hypotheses; use of decide for side facts (not computer-assisted); variant for $\operatorname{card}s\le3$ (labelled as a variant, proved from Lemma~\ref{fd:lem:small})\\
Lem.~\ref{fd:lem:neg-triples} & \hangindent=1em\hangafter=1 \lean{ErdosSzekeres.Suk.convexPosition_of_negative_triples}\newline\dirn{ErdosSzekeres/Suk/OrientationClique.lean:33} & not printed in the axiom audit: evidence of kind (a) of the rigor labels (Section~\ref{fd:sec}): build record \texttt{docs/\allowbreak{}VALIDATION.\allowbreak{}md:73}, the axiom audit as Lean output of the same library, text search\\
Lem.~\ref{fd:lem:pos-cycle} & \hangindent=1em\hangafter=1 \lean{ErdosSzekeres.hasConvexNGon_of_positive_cycle}\newline\dirn{ErdosSzekeres/ConvexCycleCuts.lean:11} & all increasing triples of the embedding e\\
Rem.~\ref{fd:rem:pos-cycle} & \hangindent=1em\hangafter=1 \lean{ErdosSzekeres.ConvexCycleCut}\newline\dirn{ErdosSzekeres/ConvexCycleCuts.lean:8}\par \hangindent=1em\hangafter=1 \lean{ErdosSzekeres.convexCycleCut_of_avoids}\newline\dirn{ErdosSzekeres/ConvexCycleCuts.lean:31} & clause family of the SMT coordinate-repair searches (\texttt{docs/\allowbreak{}SAT.\allowbreak{}md}); \texttt{convexCycleCut\_\allowbreak{}of\_\allowbreak{}avoids} is the contrapositive of Lemma~\ref{fd:lem:pos-cycle}\\
Lem.~\ref{fd:lem:separation} & written step & written step; both directions proved (nearest point for the hard direction; Heine-Borel and extreme value theorem cited from Rudin); main-library variant of the easy direction: \texttt{Omission.\allowbreak{}BoundaryGeometry.\allowbreak{}exposed\_\allowbreak{}not\_\allowbreak{}in\_\allowbreak{}hull} (\texttt{ErdosSzekeres/\allowbreak{}Omission/\allowbreak{}BoundaryGeometry.\allowbreak{}lean:30}), named in Coverage\\
\addlinespace
\multicolumn{3}{@{}l}{\normalsize\textbf{Section~\ref{ll:sec}}}\\[2pt]
Def.~\ref{ll:def:orbit} & \hangindent=1em\hangafter=1 \lean{ErdosSzekeres.LayerLemma.orbitFS}\newline\dirn{LAYERLEAN/LayerLemma.lean:48} & \\
Lem.~\ref{ll:lem:orbit} & \hangindent=1em\hangafter=1 \lean{ErdosSzekeres.LayerLemma.iterate_mem}\newline\dirn{LAYERLEAN/LayerLemma.lean:52}\par \hangindent=1em\hangafter=1 \lean{ErdosSzekeres.LayerLemma.mem_orbitFS_self}\newline\dirn{LAYERLEAN/LayerLemma.lean:59}\par \hangindent=1em\hangafter=1 \lean{ErdosSzekeres.LayerLemma.orbitFS_subset}\newline\dirn{LAYERLEAN/LayerLemma.lean:62}\par \hangindent=1em\hangafter=1 \lean{ErdosSzekeres.LayerLemma.orbitFS_invariant}\newline\dirn{LAYERLEAN/LayerLemma.lean:68}\par \hangindent=1em\hangafter=1 \lean{ErdosSzekeres.LayerLemma.card_orbitFS_of_minimalPeriod}\newline\dirn{LAYERLEAN/LayerLemma.lean:80}\par \hangindent=1em\hangafter=1 \lean{ErdosSzekeres.LayerLemma.mem_orbitFS_of_not_disjoint}\newline\dirn{LAYERLEAN/LayerLemma.lean:91} & \\
Def.~\ref{ll:def:inv} & \hangindent=1em\hangafter=1 \lean{ErdosSzekeres.LayerLemma.Invariant}\newline\dirn{LAYERLEAN/LayerLemma.lean:105} & \\
Def.~\ref{ll:def:hv} & \hangindent=1em\hangafter=1 \lean{ErdosSzekeres.LayerLemma.hullVerts}\newline\dirn{LAYERLEAN/LayerLemma.lean:111} & \\
Lem.~\ref{ll:lem:hv-basic} & \hangindent=1em\hangafter=1 \lean{ErdosSzekeres.LayerLemma.mem_hullVerts}\newline\dirn{LAYERLEAN/LayerLemma.lean:114}\par \hangindent=1em\hangafter=1 \lean{ErdosSzekeres.LayerLemma.hullVerts_subset}\newline\dirn{LAYERLEAN/LayerLemma.lean:118} & \\
Lem.~\ref{ll:lem:hv} & \hangindent=1em\hangafter=1 \lean{ErdosSzekeres.LayerLemma.convexPosition_hullVerts}\newline\dirn{LAYERLEAN/LayerLemma.lean:121}\par \hangindent=1em\hangafter=1 \lean{ErdosSzekeres.LayerLemma.extreme_mem_hullVerts}\newline\dirn{LAYERLEAN/LayerLemma.lean:130}\par \hangindent=1em\hangafter=1 \lean{ErdosSzekeres.LayerLemma.subset_convexHull_hullVerts}\newline\dirn{LAYERLEAN/LayerLemma.lean:146}\par \hangindent=1em\hangafter=1 \lean{ErdosSzekeres.LayerLemma.hullVerts_nonempty}\newline\dirn{LAYERLEAN/LayerLemma.lean:139} & (b)-(d) proved by elementary arguments, different from the Lean proofs (Lean: (b) by two Mathlib facts on extreme points, (c) by the Krein-Milman theorem, (d) by the Krein-Milman lemma)\\
Lem.~\ref{ll:lem:map-hv} & \hangindent=1em\hangafter=1 \lean{ErdosSzekeres.LayerLemma.map_mem_hullVerts}\newline\dirn{LAYERLEAN/LayerLemma.lean:162} & \\
Thm.~\ref{ll:thm:core} & \hangindent=1em\hangafter=1 \lean{ErdosSzekeres.LayerLemma.hullVerts_eq_orbitFS}\newline\dirn{LAYERLEAN/LayerLemma.lean:189} & no general position\\
Lem.~\ref{ll:lem:inv-sdiff} & \hangindent=1em\hangafter=1 \lean{ErdosSzekeres.LayerLemma.invariant_sdiff_hullVerts}\newline\dirn{LAYERLEAN/LayerLemma.lean:217} & \\
Rem.~\ref{ll:rem:core} & -- & remark on hypotheses and degenerate cases of \texttt{hullVerts\_\allowbreak{}eq\_\allowbreak{}orbitFS}\\
Def.~\ref{ll:def:complex} & \hangindent=1em\hangafter=1 \lean{ErdosSzekeres.LayerLemma.toC}\newline\dirn{LAYERLEAN/LayerLemma.lean:231}\par \hangindent=1em\hangafter=1 \lean{ErdosSzekeres.LayerLemma.ofC}\newline\dirn{LAYERLEAN/LayerLemma.lean:234}\par \hangindent=1em\hangafter=1 \lean{ErdosSzekeres.LayerLemma.crot}\newline\dirn{LAYERLEAN/LayerLemma.lean:252} & \\
Lem.~\ref{ll:lem:complex} & \hangindent=1em\hangafter=1 \lean{ErdosSzekeres.LayerLemma.toC_re}\newline\dirn{LAYERLEAN/LayerLemma.lean:236}\par \hangindent=1em\hangafter=1 \lean{ErdosSzekeres.LayerLemma.toC_im}\newline\dirn{LAYERLEAN/LayerLemma.lean:237}\par \hangindent=1em\hangafter=1 \lean{ErdosSzekeres.LayerLemma.ofC_zero}\newline\dirn{LAYERLEAN/LayerLemma.lean:238}\par \hangindent=1em\hangafter=1 \lean{ErdosSzekeres.LayerLemma.ofC_one}\newline\dirn{LAYERLEAN/LayerLemma.lean:239}\par \hangindent=1em\hangafter=1 \lean{ErdosSzekeres.LayerLemma.toC_ofC}\newline\dirn{LAYERLEAN/LayerLemma.lean:241}\par \hangindent=1em\hangafter=1 \lean{ErdosSzekeres.LayerLemma.ofC_toC}\newline\dirn{LAYERLEAN/LayerLemma.lean:244}\par \hangindent=1em\hangafter=1 \lean{ErdosSzekeres.LayerLemma.toC_injective}\newline\dirn{LAYERLEAN/LayerLemma.lean:248}\par \hangindent=1em\hangafter=1 \lean{ErdosSzekeres.LayerLemma.toC_crot}\newline\dirn{LAYERLEAN/LayerLemma.lean:254}\par \hangindent=1em\hangafter=1 \lean{ErdosSzekeres.LayerLemma.crot_one}\newline\dirn{LAYERLEAN/LayerLemma.lean:257}\par \hangindent=1em\hangafter=1 \lean{ErdosSzekeres.LayerLemma.crot_comp}\newline\dirn{LAYERLEAN/LayerLemma.lean:262}\par \hangindent=1em\hangafter=1 \lean{ErdosSzekeres.LayerLemma.crot_iterate}\newline\dirn{LAYERLEAN/LayerLemma.lean:267}\par \hangindent=1em\hangafter=1 \lean{ErdosSzekeres.LayerLemma.crot_eq_self_iff}\newline\dirn{LAYERLEAN/LayerLemma.lean:273}\par \hangindent=1em\hangafter=1 \lean{ErdosSzekeres.LayerLemma.crot_injective}\newline\dirn{LAYERLEAN/LayerLemma.lean:287} & includes the seven @[simp] lemmas (\texttt{toC\_\allowbreak{}re}, \texttt{toC\_\allowbreak{}im}, \texttt{ofC\_\allowbreak{}zero}, \texttt{ofC\_\allowbreak{}one}, \texttt{toC\_\allowbreak{}ofC}, \texttt{ofC\_\allowbreak{}toC}, \texttt{toC\_\allowbreak{}crot})\\
Def.~\ref{ll:def:rot} & \hangindent=1em\hangafter=1 \lean{ErdosSzekeres.LayerLemma.rotAbout}\newline\dirn{LAYERLEAN/LayerLemma.lean:298} & \\
Lem.~\ref{ll:lem:rot} & \hangindent=1em\hangafter=1 \lean{ErdosSzekeres.LayerLemma.rotAbout_eq_crot}\newline\dirn{LAYERLEAN/LayerLemma.lean:302}\par \hangindent=1em\hangafter=1 \lean{ErdosSzekeres.LayerLemma.rotAff}\newline\dirn{LAYERLEAN/LayerLemma.lean:311}\par \hangindent=1em\hangafter=1 \lean{ErdosSzekeres.LayerLemma.rotAbout_injective}\newline\dirn{LAYERLEAN/LayerLemma.lean:319}\par \hangindent=1em\hangafter=1 \lean{ErdosSzekeres.LayerLemma.rotAbout_iterate}\newline\dirn{LAYERLEAN/LayerLemma.lean:465} & \texttt{rotAff} is a def (affine-map structure); \texttt{rotAbout\_\allowbreak{}iterate} is line 465\\
Lem.~\ref{ll:lem:erase} & \hangindent=1em\hangafter=1 \lean{ErdosSzekeres.LayerLemma.invariant_erase_centre}\newline\dirn{LAYERLEAN/LayerLemma.lean:325} & \\
Lem.~\ref{ll:lem:minperiod} & \hangindent=1em\hangafter=1 \lean{ErdosSzekeres.LayerLemma.minimalPeriod_rotAbout}\newline\dirn{LAYERLEAN/LayerLemma.lean:336} & \\
Thm.~\ref{ll:thm:core-rot} & \hangindent=1em\hangafter=1 \lean{ErdosSzekeres.LayerLemma.hullVerts_eq_orbitFS_rotAbout}\newline\dirn{LAYERLEAN/LayerLemma.lean:364} & \\
Def.~\ref{ll:def:peel} & \hangindent=1em\hangafter=1 \lean{ErdosSzekeres.LayerLemma.peel}\newline\dirn{LAYERLEAN/LayerLemma.lean:374} & \\
Lem.~\ref{ll:lem:peel-basic} & \hangindent=1em\hangafter=1 \lean{ErdosSzekeres.LayerLemma.peel_succ}\newline\dirn{LAYERLEAN/LayerLemma.lean:378}\par \hangindent=1em\hangafter=1 \lean{ErdosSzekeres.LayerLemma.peel_subset}\newline\dirn{LAYERLEAN/LayerLemma.lean:381} & \\
Lem.~\ref{ll:lem:antitone} & \hangindent=1em\hangafter=1 \lean{ErdosSzekeres.LayerLemma.peel_antitone}\newline\dirn{LAYERLEAN/LayerLemma.lean:460} & \\
Lem.~\ref{ll:lem:inv-peel} & \hangindent=1em\hangafter=1 \lean{ErdosSzekeres.LayerLemma.invariant_peel}\newline\dirn{LAYERLEAN/LayerLemma.lean:386} & \\
Thm.~\ref{ll:thm:layers} & \hangindent=1em\hangafter=1 \lean{ErdosSzekeres.LayerLemma.layers_are_orbits}\newline\dirn{LAYERLEAN/LayerLemma.lean:402} & \\
Lem.~\ref{ll:lem:exists-layer} & \hangindent=1em\hangafter=1 \lean{ErdosSzekeres.LayerLemma.exists_layer}\newline\dirn{LAYERLEAN/LayerLemma.lean:419} & \\
Thm.~\ref{ll:thm:card-peel} & \hangindent=1em\hangafter=1 \lean{ErdosSzekeres.LayerLemma.card_peel}\newline\dirn{LAYERLEAN/LayerLemma.lean:1037} & \\
Cor.~\ref{ll:cor:decomp} & written step & written step (not a Lean declaration): decomposition of a set of exactly $m\cdot t$ points into t nested orbits\\
Thm.~\ref{ll:thm:reduction} & \hangindent=1em\hangafter=1 \lean{ErdosSzekeres.LayerLemma.convexNGon_of_layers}\newline\dirn{LAYERLEAN/LayerLemma.lean:1066} & takes the premise H (every invariant centre-free set of exactly $m\cdot t$ points in general position has a convex K-gon) as an explicit hypothesis\\
Def.~\ref{ll:def:66-70} & \hangindent=1em\hangafter=1 \lean{ErdosSzekeres.LayerLemma.Sixfold66}\newline\dirn{LAYERLEAN/LayerLemma.lean:1102}\par \hangindent=1em\hangafter=1 \lean{ErdosSzekeres.LayerLemma.Sevenfold70}\newline\dirn{LAYERLEAN/LayerLemma.lean:1107} & propositions (Lean def \dots{} : Prop), not proved in Lean; real numerals 6 and 7 in the angle\\
Thm.~\ref{ll:thm:sixfold} & \hangindent=1em\hangafter=1 \lean{ErdosSzekeres.LayerLemma.sixfold_octagon}\newline\dirn{LAYERLEAN/LayerLemma.lean:1113} & takes hypotheses: Sixfold66; uses Lemma~\ref{lg:lem:centre} (argument also included)\\
Thm.~\ref{ll:thm:sevenfold} & \hangindent=1em\hangafter=1 \lean{ErdosSzekeres.LayerLemma.sevenfold_octagon}\newline\dirn{LAYERLEAN/LayerLemma.lean:1136} & takes hypotheses: Sevenfold70\\
Rem.~\ref{ll:rem:hyp} & -- & status of \texttt{Sixfold66/\allowbreak{}Sevenfold70}: not proved in Lean; supported by \texttt{C6UPPER/\allowbreak{}C7UPPER} \texttt{cake\_\allowbreak{}lpr} verdicts (computer-assisted: LRAT proofs checked by \texttt{cake\_\allowbreak{}lpr} outside Lean; the project's records call such a result a certificate, a label this document does not use; encoder B's formula alone, encoder A's with a symmetry-breaking unit clause) plus the written soundness proof in \texttt{WRITEUP/\allowbreak{}note/\allowbreak{}k8\_\allowbreak{}soundness.\allowbreak{}tex} (not refereed; gaps G1-G4); ENCBLEAN Lean route for $m=6$ assumes \texttt{EncB.\allowbreak{}CubesUnsat64} (trust items E1-E6)\\
Rem.~\ref{ll:rem:numerals} & -- & real-cast numerals 6, 7 versus (m : R)\\
\addlinespace
\multicolumn{3}{@{}l}{\normalsize\textbf{Section~\ref{lg:sec}}}\\[2pt]
Lem.~\ref{lg:lem:layer} & written step & written step; collects the opening steps of the Lean proofs of \texttt{strict\_\allowbreak{}nesting}, \texttt{lin\_\allowbreak{}pos\_\allowbreak{}of\_\allowbreak{}peel} and \texttt{dist2\_\allowbreak{}le\_\allowbreak{}of\_\allowbreak{}later}; uses Theorem~\ref{ll:thm:core-rot}, Lemma~\ref{ll:lem:inv-peel}, Lemma~\ref{ll:lem:hv}, Lemma~\ref{ll:lem:antitone}\\
Lem.~\ref{lg:lem:rotfacts} & written step & written step; elementary facts about \texttt{rotAbout} (composition, rotation by 0 and $2\pi$, distance to the centre); the Lean proofs obtain the corresponding steps through the complex form of \texttt{rotAbout}, by direct computation or through minimal periods\\
Lem.~\ref{lg:lem:rot-iter} & \hangindent=1em\hangafter=1 \lean{ErdosSzekeres.LayerLemma.rotAbout_iterate}\newline\dirn{LAYERLEAN/LayerLemma.lean:465} & Lean-checked; axioms in \texttt{LEANEXT/\allowbreak{}review/\allowbreak{}axioms.\allowbreak{}log}; proof here via the matrix composition rule; Lean uses the complex form\\
Cor.~\ref{lg:cor:period} & written step & written step; $\rho^{m}=\mathrm{id}$ for $0<m$ (Lean uses the minimal period instead)\\
Lem.~\ref{lg:lem:orient-rot} & \hangindent=1em\hangafter=1 \lean{ErdosSzekeres.LayerLemma.orient_rotAbout}\newline\dirn{LAYERLEAN/LayerLemma.lean:472} & Lean-checked; axioms in \texttt{LEANEXT/\allowbreak{}review/\allowbreak{}axioms.\allowbreak{}log}\\
Lem.~\ref{lg:lem:trig} & \hangindent=1em\hangafter=1 \lean{ErdosSzekeres.LayerLemma.trig_layer}\newline\dirn{LAYERLEAN/LayerLemma.lean:480} & Lean-checked; axioms in \texttt{LAYERLEAN/\allowbreak{}logs/\allowbreak{}axioms.\allowbreak{}log} and \texttt{LEANEXT/\allowbreak{}review/\allowbreak{}axioms.\allowbreak{}log}\\
Lem.~\ref{lg:lem:regular} & \hangindent=1em\hangafter=1 \lean{ErdosSzekeres.LayerLemma.orient_regular_nonneg}\newline\dirn{LAYERLEAN/LayerLemma.lean:488} & Lean-checked; axioms in \texttt{LAYERLEAN/\allowbreak{}logs/\allowbreak{}axioms.\allowbreak{}log} and \texttt{LEANEXT/\allowbreak{}review/\allowbreak{}axioms.\allowbreak{}log}; its proof uses \texttt{trig\_\allowbreak{}layer} (Lemma~\ref{lg:lem:trig}, stated just before)\\
Def.~\ref{lg:def:lin} & \hangindent=1em\hangafter=1 \lean{ErdosSzekeres.LayerLemma.lin}\newline\dirn{LAYERLEAN/LayerLemma.lean:625} & Lean definition\\
Lem.~\ref{lg:lem:lin-rot} & \hangindent=1em\hangafter=1 \lean{ErdosSzekeres.LayerLemma.lin_rotAbout}\newline\dirn{LAYERLEAN/LayerLemma.lean:794} & Lean-checked; axioms in \texttt{LEANEXT/\allowbreak{}review/\allowbreak{}axioms.\allowbreak{}log}\\
Lem.~\ref{lg:lem:lin-orbit} & \hangindent=1em\hangafter=1 \lean{ErdosSzekeres.LayerLemma.lin_orbit_pos}\newline\dirn{LAYERLEAN/LayerLemma.lean:801} & Lean-checked; axioms in \texttt{LEANEXT/\allowbreak{}review/\allowbreak{}axioms.\allowbreak{}log}\\
Lem.~\ref{lg:lem:next-prev} & \hangindent=1em\hangafter=1 \lean{ErdosSzekeres.LayerLemma.orient_next_prev_pos}\newline\dirn{LAYERLEAN/LayerLemma.lean:830} & Lean-checked; axioms in \texttt{LEANEXT/\allowbreak{}review/\allowbreak{}axioms.\allowbreak{}log}\\
Lem.~\ref{lg:lem:halfplane-convex} & \hangindent=1em\hangafter=1 \lean{ErdosSzekeres.LayerLemma.convex_orient_nonneg}\newline\dirn{LAYERLEAN/LayerLemma.lean:516} & Lean-checked; axioms in \texttt{LEANEXT/\allowbreak{}review/\allowbreak{}axioms.\allowbreak{}log}\\
Def.~\ref{lg:def:dist2} & \hangindent=1em\hangafter=1 \lean{ErdosSzekeres.LayerLemma.dist2}\newline\dirn{LAYERLEAN/LayerLemma.lean:999} & Lean definition\\
Lem.~\ref{lg:lem:disk} & \hangindent=1em\hangafter=1 \lean{ErdosSzekeres.LayerLemma.convex_disk}\newline\dirn{LAYERLEAN/LayerLemma.lean:1002} & Lean-checked; axioms in \texttt{LAYERLEAN/\allowbreak{}logs/\allowbreak{}axioms.\allowbreak{}log} and \texttt{LEANEXT/\allowbreak{}review/\allowbreak{}axioms.\allowbreak{}log}\\
Thm.~\ref{lg:thm:nesting} & \hangindent=1em\hangafter=1 \lean{ErdosSzekeres.LayerLemma.strict_nesting}\newline\dirn{LAYERLEAN/LayerLemma.lean:529} & Lean-checked; axioms in \texttt{LAYERLEAN/\allowbreak{}logs/\allowbreak{}axioms.\allowbreak{}log} and \texttt{LEANEXT/\allowbreak{}review/\allowbreak{}axioms.\allowbreak{}log}; $x\ne\rho x$ proved here via \texttt{lin\_\allowbreak{}orbit\_\allowbreak{}pos}; Lean uses the minimal period\\
Lem.~\ref{lg:lem:centre} & \hangindent=1em\hangafter=1 \lean{ErdosSzekeres.LayerLemma.centre_not_mem_of_even}\newline\dirn{LAYERLEAN/LayerLemma.lean:582} & Lean-checked; axioms in \texttt{LAYERLEAN/\allowbreak{}logs/\allowbreak{}axioms.\allowbreak{}log} and \texttt{LEANEXT/\allowbreak{}review/\allowbreak{}axioms.\allowbreak{}log}\\
Thm.~\ref{lg:thm:radius} & \hangindent=1em\hangafter=1 \lean{ErdosSzekeres.LayerLemma.dist2_le_of_later}\newline\dirn{LAYERLEAN/LayerLemma.lean:1012} & Lean-checked; axioms in \texttt{LAYERLEAN/\allowbreak{}logs/\allowbreak{}axioms.\allowbreak{}log} and \texttt{LEANEXT/\allowbreak{}review/\allowbreak{}axioms.\allowbreak{}log}; weak inequality only, $n\le n'$\\
Lem.~\ref{lg:lem:cramer} & written step & written step; the planar identity behind \texttt{orient\_\allowbreak{}trans\_\allowbreak{}halfplane}, \texttt{fan\_\allowbreak{}step} and the key identity of \texttt{lin\_\allowbreak{}pos\_\allowbreak{}of\_\allowbreak{}peel}; Lean checks only these three instances, each by ring\\
Lem.~\ref{lg:lem:trans} & \hangindent=1em\hangafter=1 \lean{ErdosSzekeres.LayerLemma.orient_trans_halfplane}\newline\dirn{LAYERLEAN/LayerLemma.lean:629} & Lean-checked; axioms in \texttt{LAYERLEAN/\allowbreak{}logs/\allowbreak{}axioms.\allowbreak{}log} and \texttt{LEANEXT/\allowbreak{}review/\allowbreak{}axioms.\allowbreak{}log}; key identity derived here from Lemma~\ref{lg:lem:cramer}; Lean uses ring\\
Lem.~\ref{lg:lem:fan-step} & \hangindent=1em\hangafter=1 \lean{ErdosSzekeres.LayerLemma.fan_step}\newline\dirn{LAYERLEAN/LayerLemma.lean:641} & Lean-checked; axioms in \texttt{LAYERLEAN/\allowbreak{}logs/\allowbreak{}axioms.\allowbreak{}log} and \texttt{LEANEXT/\allowbreak{}review/\allowbreak{}axioms.\allowbreak{}log}; key identity derived here from Lemma~\ref{lg:lem:cramer}; Lean uses ring\\
Lem.~\ref{lg:lem:orient-lin} & \hangindent=1em\hangafter=1 \lean{ErdosSzekeres.LayerLemma.orient_eq_lin}\newline\dirn{LAYERLEAN/LayerLemma.lean:653} & Lean-checked; axioms in \texttt{LEANEXT/\allowbreak{}review/\allowbreak{}axioms.\allowbreak{}log}\\
Lem.~\ref{lg:lem:fan-pos} & \hangindent=1em\hangafter=1 \lean{ErdosSzekeres.LayerLemma.fan_positive}\newline\dirn{LAYERLEAN/LayerLemma.lean:658} & Lean-checked; axioms in \texttt{LAYERLEAN/\allowbreak{}logs/\allowbreak{}axioms.\allowbreak{}log} and \texttt{LEANEXT/\allowbreak{}review/\allowbreak{}axioms.\allowbreak{}log}\\
Def.~\ref{lg:def:fanpt} & \hangindent=1em\hangafter=1 \lean{ErdosSzekeres.LayerLemma.fanPt}\newline\dirn{LAYERLEAN/LayerLemma.lean:694} & Lean definition\\
Thm.~\ref{lg:thm:fan} & \hangindent=1em\hangafter=1 \lean{ErdosSzekeres.LayerLemma.hasConvexNGon_of_fan}\newline\dirn{LAYERLEAN/LayerLemma.lean:700} & Lean-checked; axioms in \texttt{LAYERLEAN/\allowbreak{}logs/\allowbreak{}axioms.\allowbreak{}log} and \texttt{LEANEXT/\allowbreak{}review/\allowbreak{}axioms.\allowbreak{}log}; uses \texttt{ErdosSzekeres.\allowbreak{}Suk.\allowbreak{}convexPosition\_\allowbreak{}of\_\allowbreak{}positive\_\allowbreak{}triples} (Lemma~\ref{fd:lem:pos-triples}); no general position\\
Lem.~\ref{lg:lem:halfplane} & \hangindent=1em\hangafter=1 \lean{ErdosSzekeres.LayerLemma.lin_pos_of_peel}\newline\dirn{LAYERLEAN/LayerLemma.lean:864} & Lean-checked; axioms in \texttt{LAYERLEAN/\allowbreak{}logs/\allowbreak{}axioms.\allowbreak{}log} and \texttt{LEANEXT/\allowbreak{}review/\allowbreak{}axioms.\allowbreak{}log}; $\rho(p')=p$ proved here via Corollary~\ref{lg:cor:period} (\texttt{rotAbout\_\allowbreak{}iterate} and rotation by $2\pi$); Lean uses the minimal period; key identity derived here from Lemma~\ref{lg:lem:cramer}; Lean uses ring\\
Thm.~\ref{lg:thm:chain} & \hangindent=1em\hangafter=1 \lean{ErdosSzekeres.LayerLemma.chain_clauses_sound}\newline\dirn{LAYERLEAN/LayerLemma.lean:914} & Lean-checked; axioms in \texttt{LAYERLEAN/\allowbreak{}logs/\allowbreak{}axioms.\allowbreak{}log} and \texttt{LEANEXT/\allowbreak{}review/\allowbreak{}axioms.\allowbreak{}log}; Lean conclusion False; geometric statement only (that the encoders emit these clauses is not in Lean)\\
Def.~\ref{lg:def:polar} & \hangindent=1em\hangafter=1 \lean{ErdosSzekeres.LayerLemma.polar}\newline\dirn{LAYERLEAN/LayerLemma.lean:932} & Lean definition\\
Lem.~\ref{lg:lem:polar} & \hangindent=1em\hangafter=1 \lean{ErdosSzekeres.LayerLemma.orient_polar}\newline\dirn{LAYERLEAN/LayerLemma.lean:935} & Lean-checked; axioms in \texttt{LAYERLEAN/\allowbreak{}logs/\allowbreak{}axioms.\allowbreak{}log} and \texttt{LEANEXT/\allowbreak{}review/\allowbreak{}axioms.\allowbreak{}log}\\
Lem.~\ref{lg:lem:ruleA} & \hangindent=1em\hangafter=1 \lean{ErdosSzekeres.LayerLemma.orient_polar_centre_inside}\newline\dirn{LAYERLEAN/LayerLemma.lean:943} & Lean-checked; axioms in \texttt{LAYERLEAN/\allowbreak{}logs/\allowbreak{}axioms.\allowbreak{}log} and \texttt{LEANEXT/\allowbreak{}review/\allowbreak{}axioms.\allowbreak{}log}\\
Lem.~\ref{lg:lem:sinadd} & \hangindent=1em\hangafter=1 \lean{ErdosSzekeres.LayerLemma.sin_add_lt}\newline\dirn{LAYERLEAN/LayerLemma.lean:956} & Lean-checked; axioms in \texttt{LAYERLEAN/\allowbreak{}logs/\allowbreak{}axioms.\allowbreak{}log} and \texttt{LEANEXT/\allowbreak{}review/\allowbreak{}axioms.\allowbreak{}log}\\
Lem.~\ref{lg:lem:ruleB} & \hangindent=1em\hangafter=1 \lean{ErdosSzekeres.LayerLemma.orient_polar_halfplane}\newline\dirn{LAYERLEAN/LayerLemma.lean:974} & Lean-checked; axioms in \texttt{LAYERLEAN/\allowbreak{}logs/\allowbreak{}axioms.\allowbreak{}log} and \texttt{LEANEXT/\allowbreak{}review/\allowbreak{}axioms.\allowbreak{}log}\\
\addlinespace
\multicolumn{3}{@{}l}{\normalsize\textbf{Section~\ref{oc:sec}}}\\[2pt]
Def.~\ref{oc:def:basic} & \hangindent=1em\hangafter=1 \lean{ErdosSzekeres.OrbitConsistency.Invariant}\newline\dirn{LEANEXT/orbit/OrbitConsistency.lean:47}\par \hangindent=1em\hangafter=1 \lean{ErdosSzekeres.OrbitConsistency.rotAbout}\newline\dirn{LEANEXT/orbit/OrbitConsistency.lean:93}\par \hangindent=1em\hangafter=1 \lean{ErdosSzekeres.OrbitConsistency.polar}\newline\dirn{LEANEXT/orbit/OrbitConsistency.lean:112} & copies of \texttt{ErdosSzekeres.\allowbreak{}LayerLemma.\allowbreak{}Invariant} (\texttt{LAYERLEAN/\allowbreak{}LayerLemma.\allowbreak{}lean:105}), \texttt{rotAbout} (:298), polar (:932); identical formulas\\
Rem.~\ref{oc:rem:copies} & -- & identification of the duplicated definitions\\
Lem.~\ref{oc:lem:rot-add} & written step & written step: composition of rotations; $R_{0},R_{2\pi},R_{\pi}$\\
Lem.~\ref{oc:lem:iterate} & \hangindent=1em\hangafter=1 \lean{ErdosSzekeres.OrbitConsistency.rotAbout_iterate}\newline\dirn{LEANEXT/orbit/OrbitConsistency.lean:105}\par \hangindent=1em\hangafter=1 \lean{ErdosSzekeres.LayerLemma.rotAbout_iterate}\newline\dirn{LAYERLEAN/LayerLemma.lean:465} & copied declaration (outside the new content); proof by addition formulas instead of complex exponentials; axioms: \texttt{OrbitConsistency} copy covered transitively by \texttt{LEANEXT/\allowbreak{}orbit/\allowbreak{}logs/\allowbreak{}axioms.\allowbreak{}log} (used by \texttt{antipode\_\allowbreak{}mem\_\allowbreak{}of\_\allowbreak{}even}), \texttt{LayerLemma} copy printed in \texttt{LEANEXT/\allowbreak{}review/\allowbreak{}axioms.\allowbreak{}log:33}\\
Cor.~\ref{oc:cor:period} & written step & written step: $\rho^{m}=\mathrm{id}$; $\rho^{m/2}$ is the half-turn for even m\\
Lem.~\ref{oc:lem:iter-mem} & \hangindent=1em\hangafter=1 \lean{ErdosSzekeres.OrbitConsistency.iterate_mem}\newline\dirn{LEANEXT/orbit/OrbitConsistency.lean:39}\par \hangindent=1em\hangafter=1 \lean{ErdosSzekeres.LayerLemma.iterate_mem}\newline\dirn{LAYERLEAN/LayerLemma.lean:52} & copied declaration (outside the new content); axioms: \texttt{OrbitConsistency} copy covered transitively by \texttt{LEANEXT/\allowbreak{}orbit/\allowbreak{}logs/\allowbreak{}axioms.\allowbreak{}log} (used by \texttt{antipode\_\allowbreak{}mem\_\allowbreak{}of\_\allowbreak{}even}), \texttt{LayerLemma} copy printed in \texttt{LEANEXT/\allowbreak{}review/\allowbreak{}axioms.\allowbreak{}log:1} (propext, \texttt{Quot.\allowbreak{}sound} only)\\
Lem.~\ref{oc:lem:rot-inv} & \hangindent=1em\hangafter=1 \lean{ErdosSzekeres.OrbitConsistency.orient_rotAbout_inv}\newline\dirn{LEANEXT/orbit/OrbitConsistency.lean:117} & Lean-checked; axioms printed in \texttt{LEANEXT/\allowbreak{}orbit/\allowbreak{}logs/\allowbreak{}axioms.\allowbreak{}log}\\
Def.~\ref{oc:def:preserves} & \hangindent=1em\hangafter=1 \lean{ErdosSzekeres.OrbitConsistency.PreservesOrient}\newline\dirn{LEANEXT/orbit/OrbitConsistency.lean:124} & \\
Def.~\ref{oc:def:rotcomp} & \hangindent=1em\hangafter=1 \lean{ErdosSzekeres.OrbitConsistency.RotComp}\newline\dirn{LEANEXT/orbit/OrbitConsistency.lean:145} & inductive predicate; any $m\ (2\pi/0=0)$\\
Lem.~\ref{oc:lem:rotcomp-pres} & \hangindent=1em\hangafter=1 \lean{ErdosSzekeres.OrbitConsistency.preservesOrient_id}\newline\dirn{LEANEXT/orbit/OrbitConsistency.lean:127}\par \hangindent=1em\hangafter=1 \lean{ErdosSzekeres.OrbitConsistency.PreservesOrient.comp}\newline\dirn{LEANEXT/orbit/OrbitConsistency.lean:129}\par \hangindent=1em\hangafter=1 \lean{ErdosSzekeres.OrbitConsistency.PreservesOrient.iterate}\newline\dirn{LEANEXT/orbit/OrbitConsistency.lean:134}\par \hangindent=1em\hangafter=1 \lean{ErdosSzekeres.OrbitConsistency.preservesOrient_rotAbout}\newline\dirn{LEANEXT/orbit/OrbitConsistency.lean:141}\par \hangindent=1em\hangafter=1 \lean{ErdosSzekeres.OrbitConsistency.RotComp.preservesOrient}\newline\dirn{LEANEXT/orbit/OrbitConsistency.lean:150}\par \hangindent=1em\hangafter=1 \lean{ErdosSzekeres.OrbitConsistency.rotComp_iterate}\newline\dirn{LEANEXT/orbit/OrbitConsistency.lean:157} & Lean-checked; axioms printed in \texttt{LEANEXT/\allowbreak{}orbit/\allowbreak{}logs/\allowbreak{}axioms.\allowbreak{}log} (\texttt{RotComp.\allowbreak{}preservesOrient}, \texttt{rotComp\_\allowbreak{}iterate}; helpers transitively); \texttt{PreservesOrient.\allowbreak{}iterate} has no recorded axiom output (compilation of the whole file recorded in \texttt{LEANEXT/\allowbreak{}orbit/\allowbreak{}logs/\allowbreak{}build.\allowbreak{}log}, exit 0)\\
Def.~\ref{oc:def:orient3} & \hangindent=1em\hangafter=1 \lean{ErdosSzekeres.OrbitConsistency.orient3}\newline\dirn{LEANEXT/orbit/OrbitConsistency.lean:168} & \\
Fact~\ref{oc:fact:sign} & Mathlib & cited standard facts (Mathlib \texttt{Equiv.\allowbreak{}Perm.\allowbreak{}sign}, \texttt{Equiv.\allowbreak{}Perm.\allowbreak{}sign\_\allowbreak{}mul}, \texttt{Equiv.\allowbreak{}Perm.\allowbreak{}sign\_\allowbreak{}swap}, \texttt{Equiv.\allowbreak{}Perm.\allowbreak{}swap\_\allowbreak{}induction\_\allowbreak{}on}, \texttt{Int.\allowbreak{}units\_\allowbreak{}eq\_\allowbreak{}one\_\allowbreak{}or})\\
Lem.~\ref{oc:lem:swap} & \hangindent=1em\hangafter=1 \lean{ErdosSzekeres.OrbitConsistency.orient3_swap}\newline\dirn{LEANEXT/orbit/OrbitConsistency.lean:171} & Lean-checked; axioms printed in \texttt{LEANEXT/\allowbreak{}orbit/\allowbreak{}logs/\allowbreak{}axioms.\allowbreak{}log}; proof via orient symmetries instead of case expansion\\
Lem.~\ref{oc:lem:perm} & \hangindent=1em\hangafter=1 \lean{ErdosSzekeres.OrbitConsistency.orient3_perm}\newline\dirn{LEANEXT/orbit/OrbitConsistency.lean:178} & Lean-checked; axioms printed in \texttt{LEANEXT/\allowbreak{}orbit/\allowbreak{}logs/\allowbreak{}axioms.\allowbreak{}log}\\
Lem.~\ref{oc:lem:rot-perm} & \hangindent=1em\hangafter=1 \lean{ErdosSzekeres.OrbitConsistency.orient3_rot_perm}\newline\dirn{LEANEXT/orbit/OrbitConsistency.lean:194} & Lean-checked; axioms printed in \texttt{LEANEXT/\allowbreak{}orbit/\allowbreak{}logs/\allowbreak{}axioms.\allowbreak{}log}\\
Lem.~\ref{oc:lem:two-reps} & \hangindent=1em\hangafter=1 \lean{ErdosSzekeres.OrbitConsistency.sign_eq_of_two_reps}\newline\dirn{LEANEXT/orbit/OrbitConsistency.lean:204}\par \hangindent=1em\hangafter=1 \lean{ErdosSzekeres.OrbitConsistency.orient3_of_common_rep}\newline\dirn{LEANEXT/orbit/OrbitConsistency.lean:216} & Lean-checked; axioms printed in \texttt{LEANEXT/\allowbreak{}orbit/\allowbreak{}logs/\allowbreak{}axioms.\allowbreak{}log}\\
Def.~\ref{oc:def:litval} & \hangindent=1em\hangafter=1 \lean{ErdosSzekeres.OrbitConsistency.litVal}\newline\dirn{LEANEXT/orbit/OrbitConsistency.lean:223} & \\
Thm.~\ref{oc:thm:literal} & \hangindent=1em\hangafter=1 \lean{ErdosSzekeres.OrbitConsistency.literal_sound}\newline\dirn{LEANEXT/orbit/OrbitConsistency.lean:228} & Lean-checked; axioms printed in \texttt{LEANEXT/\allowbreak{}orbit/\allowbreak{}logs/\allowbreak{}axioms.\allowbreak{}log}\\
Thm.~\ref{oc:thm:encoder} & \hangindent=1em\hangafter=1 \lean{ErdosSzekeres.OrbitConsistency.encoder_literal_sound}\newline\dirn{LEANEXT/orbit/OrbitConsistency.lean:253} & Lean-checked; axioms printed in \texttt{LEANEXT/\allowbreak{}orbit/\allowbreak{}logs/\allowbreak{}axioms.\allowbreak{}log}; encoders compute parity by inversion counting, not modelled in Lean (Remark~\ref{oc:rem:parity})\\
Rem.~\ref{oc:rem:parity} & -- & inversion-count parity vs sign: six-case written check; code correspondence outside Lean; Python evidence only\\
Rem.~\ref{oc:rem:literal-use} & -- & how the encoder uses the literal rule\\
Lem.~\ref{oc:lem:antipode} & \hangindent=1em\hangafter=1 \lean{ErdosSzekeres.OrbitConsistency.antipode_mem_of_even}\newline\dirn{LEANEXT/orbit/OrbitConsistency.lean:273}\par \hangindent=1em\hangafter=1 \lean{ErdosSzekeres.LayerLemma.antipode_mem_of_even}\newline\dirn{LEANEXT/merged/Merged.lean:1160} & Lean-checked; axioms printed in \texttt{LEANEXT/\allowbreak{}orbit/\allowbreak{}logs/\allowbreak{}axioms.\allowbreak{}log} (\texttt{OrbitConsistency}); Merged copy covered transitively by \texttt{LEANEXT/\allowbreak{}merged/\allowbreak{}axioms.\allowbreak{}log}\\
Lem.~\ref{oc:lem:same-ray} & \hangindent=1em\hangafter=1 \lean{ErdosSzekeres.OrbitConsistency.same_ray_false}\newline\dirn{LEANEXT/orbit/OrbitConsistency.lean:295}\par \hangindent=1em\hangafter=1 \lean{ErdosSzekeres.LayerLemma.same_ray_false}\newline\dirn{LEANEXT/merged/Merged.lean:1181} & Lean-checked; axioms printed in \texttt{LEANEXT/\allowbreak{}orbit/\allowbreak{}logs/\allowbreak{}axioms.\allowbreak{}log} (\texttt{OrbitConsistency}); Merged copy covered transitively by \texttt{LEANEXT/\allowbreak{}merged/\allowbreak{}axioms.\allowbreak{}log}\\
Rem.~\ref{oc:rem:ray} & -- & $l>0$ cannot be weakened to $l\ne0$\\
Lem.~\ref{oc:lem:polar-rot} & \hangindent=1em\hangafter=1 \lean{ErdosSzekeres.OrbitConsistency.rotAbout_polar}\newline\dirn{LEANEXT/orbit/OrbitConsistency.lean:355}\par \hangindent=1em\hangafter=1 \lean{ErdosSzekeres.OrbitConsistency.polar_add_int}\newline\dirn{LEANEXT/orbit/OrbitConsistency.lean:383}\par \hangindent=1em\hangafter=1 \lean{ErdosSzekeres.LayerLemma.rotAbout_polar}\newline\dirn{LEANEXT/merged/Merged.lean:1241}\par \hangindent=1em\hangafter=1 \lean{ErdosSzekeres.LayerLemma.polar_add_int}\newline\dirn{LEANEXT/merged/Merged.lean:1269} & Lean-checked; axioms covered transitively by \texttt{LEANEXT/\allowbreak{}orbit/\allowbreak{}logs/\allowbreak{}axioms.\allowbreak{}log} and \texttt{LEANEXT/\allowbreak{}merged/\allowbreak{}axioms.\allowbreak{}log}\\
Lem.~\ref{oc:lem:phase-nat} & \hangindent=1em\hangafter=1 \lean{ErdosSzekeres.OrbitConsistency.eq_iterate_of_phase_eq}\newline\dirn{LEANEXT/orbit/OrbitConsistency.lean:363}\par \hangindent=1em\hangafter=1 \lean{ErdosSzekeres.LayerLemma.eq_iterate_of_phase_eq}\newline\dirn{LEANEXT/merged/Merged.lean:1249} & Lean-checked; axioms printed in \texttt{LEANEXT/\allowbreak{}orbit/\allowbreak{}logs/\allowbreak{}axioms.\allowbreak{}log} (\texttt{OrbitConsistency}); Merged copy covered transitively by \texttt{LEANEXT/\allowbreak{}merged/\allowbreak{}axioms.\allowbreak{}log}\\
Lem.~\ref{oc:lem:phase-int} & \hangindent=1em\hangafter=1 \lean{ErdosSzekeres.OrbitConsistency.eq_iterate_of_phase_eq_int}\newline\dirn{LEANEXT/orbit/OrbitConsistency.lean:389}\par \hangindent=1em\hangafter=1 \lean{ErdosSzekeres.LayerLemma.eq_iterate_of_phase_eq_int}\newline\dirn{LEANEXT/merged/Merged.lean:1275} & Lean-checked; axioms printed in \texttt{LEANEXT/\allowbreak{}orbit/\allowbreak{}logs/\allowbreak{}axioms.\allowbreak{}log} (\texttt{OrbitConsistency}); Merged copy covered transitively by \texttt{LEANEXT/\allowbreak{}merged/\allowbreak{}axioms.\allowbreak{}log}\\
Thm.~\ref{oc:thm:phase-orbit} & \hangindent=1em\hangafter=1 \lean{ErdosSzekeres.OrbitConsistency.phase_ne_of_not_mem_orbit}\newline\dirn{LEANEXT/orbit/OrbitConsistency.lean:420}\par \hangindent=1em\hangafter=1 \lean{ErdosSzekeres.LayerLemma.phase_ne_of_not_mem_orbit}\newline\dirn{LEANEXT/merged/Merged.lean:1304} & axioms printed in \texttt{LEANEXT/\allowbreak{}orbit/\allowbreak{}logs/\allowbreak{}axioms.\allowbreak{}log} and \texttt{LEANEXT/\allowbreak{}merged/\allowbreak{}axioms.\allowbreak{}log}\\
Rem.~\ref{oc:rem:phase-orbit} & -- & the doc-comment reading (nonzero phases) needs facts outside the statement\\
Lem.~\ref{oc:lem:disjoint} & \hangindent=1em\hangafter=1 \lean{ErdosSzekeres.LayerLemma.not_mem_hullVerts_peel_of_lt}\newline\dirn{LEANEXT/merged/Merged.lean:1316}\par \hangindent=1em\hangafter=1 \lean{ErdosSzekeres.LayerLemma.not_mem_hullVerts_peel_of_ne}\newline\dirn{LEANEXT/merged/Merged.lean:1322} & Lean-checked; axioms covered transitively by \texttt{LEANEXT/\allowbreak{}merged/\allowbreak{}axioms.\allowbreak{}log}\\
Lem.~\ref{oc:lem:closed} & \hangindent=1em\hangafter=1 \lean{ErdosSzekeres.LayerLemma.iterate_mem_hullVerts_peel}\newline\dirn{LEANEXT/merged/Merged.lean:1330} & Lean-checked; axioms covered transitively by \texttt{LEANEXT/\allowbreak{}merged/\allowbreak{}axioms.\allowbreak{}log}; proof uses $\rho^{m}=\mathrm{id}$ directly instead of \texttt{minimalPeriod\_\allowbreak{}rotAbout}\\
Thm.~\ref{oc:thm:layers-ray} & \hangindent=1em\hangafter=1 \lean{ErdosSzekeres.LayerLemma.not_same_ray_of_distinct_layers}\newline\dirn{LEANEXT/merged/Merged.lean:1350} & Lean-checked; axioms printed in \texttt{LEANEXT/\allowbreak{}merged/\allowbreak{}axioms.\allowbreak{}log}\\
Thm.~\ref{oc:thm:layers-phase} & \hangindent=1em\hangafter=1 \lean{ErdosSzekeres.LayerLemma.phase_ne_of_distinct_layers}\newline\dirn{LEANEXT/merged/Merged.lean:1367} & Lean-checked; axioms printed in \texttt{LEANEXT/\allowbreak{}merged/\allowbreak{}axioms.\allowbreak{}log}\\
Rem.~\ref{oc:rem:layers} & -- & $c\notin S$ derivable via Lemma~\ref{lg:lem:centre}; example showing $\pi/m$ cannot replace $2\pi/m$\\
Def.~\ref{oc:def:keys} & Mathlib & Mathlib definitions: \texttt{toIcoDiv/\allowbreak{}toIcoMod}, \texttt{Int.\allowbreak{}ediv/\allowbreak{}Int.\allowbreak{}emod}, lexicographic order (\texttt{Prod.\allowbreak{}Lex}), \texttt{StrictMonoOn}\\
Lem.~\ref{oc:lem:toico} & written step & written step (Mathlib \texttt{toIcoMod\_\allowbreak{}eq\_\allowbreak{}iff})\\
Lem.~\ref{oc:lem:cell} & \hangindent=1em\hangafter=1 \lean{KeyOrder.cell_lt_iff}\newline\dirn{LEANEXT/keys/KeyOrder.lean:23}\par \hangindent=1em\hangafter=1 \lean{KeyOrder.cell_eq_iff}\newline\dirn{LEANEXT/keys/KeyOrder.lean:45} & Lean-checked; axioms printed by \texttt{KeyOrder.\allowbreak{}lean:151-158}, output in \texttt{ENCBLEAN/\allowbreak{}logs/\allowbreak{}m4\_\allowbreak{}build.\allowbreak{}log:3-10}\\
Lem.~\ref{oc:lem:rank} & \hangindent=1em\hangafter=1 \lean{KeyOrder.cell_lt_iff_rank}\newline\dirn{LEANEXT/keys/KeyOrder.lean:62}\par \hangindent=1em\hangafter=1 \lean{KeyOrder.cell_lt_iff_rank_finset}\newline\dirn{LEANEXT/keys/KeyOrder.lean:71} & Lean-checked; axioms printed by \texttt{KeyOrder.\allowbreak{}lean:151-158}, output in \texttt{ENCBLEAN/\allowbreak{}logs/\allowbreak{}m4\_\allowbreak{}build.\allowbreak{}log:3-10}\\
Lem.~\ref{oc:lem:reduce} & \hangindent=1em\hangafter=1 \lean{KeyOrder.reduce_mod_two_pi}\newline\dirn{LEANEXT/keys/KeyOrder.lean:81} & Lean-checked; axioms printed by \texttt{KeyOrder.\allowbreak{}lean:151-158}, output in \texttt{ENCBLEAN/\allowbreak{}logs/\allowbreak{}m4\_\allowbreak{}build.\allowbreak{}log:3-10}\\
Lem.~\ref{oc:lem:antipode-cell} & \hangindent=1em\hangafter=1 \lean{KeyOrder.antipode_odd}\newline\dirn{LEANEXT/keys/KeyOrder.lean:100}\par \hangindent=1em\hangafter=1 \lean{KeyOrder.antipode_even}\newline\dirn{LEANEXT/keys/KeyOrder.lean:115} & Lean-checked; axioms printed by \texttt{KeyOrder.\allowbreak{}lean:151-158}, output in \texttt{ENCBLEAN/\allowbreak{}logs/\allowbreak{}m4\_\allowbreak{}build.\allowbreak{}log:3-10}; oddness of m not assumed in \texttt{antipode\_\allowbreak{}odd}\\
Thm.~\ref{oc:thm:key} & \hangindent=1em\hangafter=1 \lean{KeyOrder.reduced_lt_iff_key}\newline\dirn{LEANEXT/keys/KeyOrder.lean:139} & Lean-checked; axioms printed by \texttt{KeyOrder.\allowbreak{}lean:151-158}, output in \texttt{ENCBLEAN/\allowbreak{}logs/\allowbreak{}m4\_\allowbreak{}build.\allowbreak{}log:3-10}\\
Rem.~\ref{oc:rem:halfopen} & -- & half-open intervals are needed\\
Rem.~\ref{oc:rem:keys-outside} & -- & encoder decomposition, local ranks and key-to-circular-order step are outside Lean\\
\addlinespace
\multicolumn{3}{@{}l}{\normalsize\textbf{Section~\ref{pc:sec}}}\\[2pt]
Def.~\ref{pc:def:polar} & \hangindent=1em\hangafter=1 \lean{PhaseClassification.polar}\newline\dirn{LEANEXT/phase/PhaseClassification.lean:41}\par \hangindent=1em\hangafter=1 \lean{ErdosSzekeres.LayerLemma.dist2}\newline\dirn{LAYERLEAN/LayerLemma.lean:999} & \texttt{PhaseClassification.\allowbreak{}polar} is a separate constant with the formula of \texttt{LayerLemma.\allowbreak{}polar} (\texttt{LAYERLEAN/\allowbreak{}LayerLemma.\allowbreak{}lean:932}); dist2 repeated at \texttt{LEANEXT/\allowbreak{}linked/\allowbreak{}Linked.\allowbreak{}lean:999} and \texttt{LEANEXT/\allowbreak{}capstone/\allowbreak{}Capstone2.\allowbreak{}lean:1007}\\
Lem.~\ref{pc:lem:polar-facts} & \hangindent=1em\hangafter=1 \lean{PhaseClassification.orient_polar}\newline\dirn{LEANEXT/phase/PhaseClassification.lean:44}\par \hangindent=1em\hangafter=1 \lean{PhaseClassification.orient_polar_centre}\newline\dirn{LEANEXT/phase/PhaseClassification.lean:407}\par \hangindent=1em\hangafter=1 \lean{PhaseClassification.orient_polar_centre_inside}\newline\dirn{LEANEXT/phase/PhaseClassification.lean:51}\par \hangindent=1em\hangafter=1 \lean{PhaseClassification.sin_add_lt}\newline\dirn{LEANEXT/phase/PhaseClassification.lean:64}\par \hangindent=1em\hangafter=1 \lean{PhaseClassification.orient_polar_halfplane}\newline\dirn{LEANEXT/phase/PhaseClassification.lean:80} & (a),(c),(d),(e) are copies of \texttt{LAYERLEAN/\allowbreak{}LayerLemma.\allowbreak{}lean:935-996} (Lemma~\ref{lg:lem:polar}, Lemma~\ref{lg:lem:ruleA}, Lemma~\ref{lg:lem:sinadd}, Lemma~\ref{lg:lem:ruleB}); proofs repeated; axioms covered transitively by \texttt{LEANEXT/\allowbreak{}phase/\allowbreak{}logs/\allowbreak{}axioms.\allowbreak{}log}; head cites (a), parts (b)-(e) cite their own declarations\\
Lem.~\ref{pc:lem:periodic} & \hangindent=1em\hangafter=1 \lean{PhaseClassification.polar_add_int_two_pi}\newline\dirn{LEANEXT/phase/PhaseClassification.lean:106}\par \hangindent=1em\hangafter=1 \lean{PhaseClassification.polar_add_two_pi}\newline\dirn{LEANEXT/phase/PhaseClassification.lean:111}\par \hangindent=1em\hangafter=1 \lean{PhaseClassification.polar_toIcoMod}\newline\dirn{LEANEXT/phase/PhaseClassification.lean:485} & \texttt{polar\_\allowbreak{}toIcoMod} in \texttt{LEANEXT/\allowbreak{}phase/\allowbreak{}logs/\allowbreak{}axioms.\allowbreak{}log}; the other two used by logged theorems\\
Lem.~\ref{pc:lem:exists-polar} & \hangindent=1em\hangafter=1 \lean{PhaseClassification.exists_polar}\newline\dirn{LEANEXT/phase/PhaseClassification.lean:459} & in \texttt{LEANEXT/\allowbreak{}phase/\allowbreak{}logs/\allowbreak{}axioms.\allowbreak{}log}; written proof via arccos (Lean uses \texttt{Complex.\allowbreak{}arg})\\
Lem.~\ref{pc:lem:helpers} & \hangindent=1em\hangafter=1 \lean{ErdosSzekeres.Capstone.polar_eq}\newline\dirn{LEANEXT/capstone/Capstone.lean:2001}\par \hangindent=1em\hangafter=1 \lean{ErdosSzekeres.Capstone.polar_zero}\newline\dirn{LEANEXT/capstone/Capstone2.lean:2051}\par \hangindent=1em\hangafter=1 \lean{ErdosSzekeres.LayerLemma.dist2_polar}\newline\dirn{LEANEXT/linked/Linked.lean:1390}\par \hangindent=1em\hangafter=1 \lean{ErdosSzekeres.Capstone.exists_angle}\newline\dirn{LEANEXT/capstone/Capstone2.lean:2056}\par \hangindent=1em\hangafter=1 \lean{ErdosSzekeres.Capstone.sqrt_dist2_pos}\newline\dirn{LEANEXT/capstone/Capstone2.lean:2064}\par \hangindent=1em\hangafter=1 \lean{ErdosSzekeres.Capstone.ne_of_mem}\newline\dirn{LEANEXT/capstone/Capstone2.lean:2073} & no saved axioms output in LEANEXT; used (via the byte-identical Capstone2 copy) in proofs on which \texttt{EncB.\allowbreak{}Geo.\allowbreak{}labelling\_\allowbreak{}of\_\allowbreak{}counterexample} and \texttt{EncB.\allowbreak{}Geo.\allowbreak{}modelOfLabelling} depend, whose axioms are pinned at \texttt{ENCBLEAN/\allowbreak{}geo/\allowbreak{}ES8Geo/\allowbreak{}Audit.\allowbreak{}lean:45-49}; the copy is compiled against verbatim copies of the ten repository modules it needs (\texttt{ENCBLEAN/\allowbreak{}geo/\allowbreak{}VENDOR.\allowbreak{}sha256}, hashes equal to the repository files)\\
Def.~\ref{pc:def:arc} & \hangindent=1em\hangafter=1 \lean{PhaseClassification.inArc}\newline\dirn{LEANEXT/phase/PhaseClassification.lean:117}\par \hangindent=1em\hangafter=1 \lean{PhaseClassification.arcHasAntipode}\newline\dirn{LEANEXT/phase/PhaseClassification.lean:120}\par \hangindent=1em\hangafter=1 \lean{PhaseClassification.Alternating}\newline\dirn{LEANEXT/phase/PhaseClassification.lean:126}\par \hangindent=1em\hangafter=1 \lean{PhaseClassification.RunMidB}\newline\dirn{LEANEXT/phase/PhaseClassification.lean:132}\par \hangindent=1em\hangafter=1 \lean{PhaseClassification.RunMidD}\newline\dirn{LEANEXT/phase/PhaseClassification.lean:136}\par \hangindent=1em\hangafter=1 \lean{PhaseClassification.RunMidA}\newline\dirn{LEANEXT/phase/PhaseClassification.lean:140} & $H_{1},H_{2},H_{3}$ are notation of the write-up, not Lean names\\
Rem.~\ref{pc:rem:reading} & -- & explanatory; gap notation $g_{1},g_{2},g_{3}$\\
Lem.~\ref{pc:lem:inarc} & \hangindent=1em\hangafter=1 \lean{PhaseClassification.inArc_of}\newline\dirn{LEANEXT/phase/PhaseClassification.lean:143}\par \hangindent=1em\hangafter=1 \lean{PhaseClassification.not_inArc}\newline\dirn{LEANEXT/phase/PhaseClassification.lean:147} & \\
Lem.~\ref{pc:lem:shift} & written step & written step; makes explicit the have-steps hrun' in the Lean proofs of \texttt{orient\_\allowbreak{}pos\_\allowbreak{}of\_\allowbreak{}runMidD/\allowbreak{}A} (\texttt{LEANEXT/\allowbreak{}phase/\allowbreak{}PhaseClassification.\allowbreak{}lean} lines 323-335, 346-362); proves the equivalences\\
Lem.~\ref{pc:lem:empty} & \hangindent=1em\hangafter=1 \lean{PhaseClassification.arc1_empty}\newline\dirn{LEANEXT/phase/PhaseClassification.lean:161}\par \hangindent=1em\hangafter=1 \lean{PhaseClassification.arc2_empty}\newline\dirn{LEANEXT/phase/PhaseClassification.lean:169}\par \hangindent=1em\hangafter=1 \lean{PhaseClassification.arc2_empty'}\newline\dirn{LEANEXT/phase/PhaseClassification.lean:177}\par \hangindent=1em\hangafter=1 \lean{PhaseClassification.arc3_empty}\newline\dirn{LEANEXT/phase/PhaseClassification.lean:185}\par \hangindent=1em\hangafter=1 \lean{PhaseClassification.arc3_empty'}\newline\dirn{LEANEXT/phase/PhaseClassification.lean:193}\par \hangindent=1em\hangafter=1 \lean{PhaseClassification.arc1_empty'}\newline\dirn{LEANEXT/phase/PhaseClassification.lean:201} & same lift integers m as the Lean proofs\\
Lem.~\ref{pc:lem:arc12} & \hangindent=1em\hangafter=1 \lean{PhaseClassification.arc12_nonempty}\newline\dirn{LEANEXT/phase/PhaseClassification.lean:229} & written proof uses a coarser three-case split than the Lean proof\\
Thm.~\ref{pc:thm:alternating} & \hangindent=1em\hangafter=1 \lean{PhaseClassification.alternating_iff_gaps_lt_pi}\newline\dirn{LEANEXT/phase/PhaseClassification.lean:210} & no genericity hypothesis\\
Rem.~\ref{pc:rem:alt-boundary} & -- & written step: a gap equal to $\pi$ makes alternation fail\\
Prop.~\ref{pc:prop:exactly-one} & \hangindent=1em\hangafter=1 \lean{PhaseClassification.alternating_exactly_one}\newline\dirn{LEANEXT/phase/PhaseClassification.lean:386} & \\
Thm.~\ref{pc:thm:runB} & \hangindent=1em\hangafter=1 \lean{PhaseClassification.runMidB_iff}\newline\dirn{LEANEXT/phase/PhaseClassification.lean:251} & hypothesis $\alpha+2\pi-\delta\ne\pi$ needed for $\Rightarrow$\\
Rem.~\ref{pc:rem:runB-example} & -- & written step: example $\delta-\alpha=\pi$ shows the gap hypothesis is needed\\
Prop.~\ref{pc:prop:gaps} & \hangindent=1em\hangafter=1 \lean{PhaseClassification.gap_classification}\newline\dirn{LEANEXT/phase/PhaseClassification.lean:268} & Lean proves only the disjunction (docstring says exactly one)\\
Rem.~\ref{pc:rem:gaps} & -- & written steps: counterexample without genericity; disjuncts mutually exclusive (not Lean-checked)\\
Thm.~\ref{pc:thm:trichotomy} & \hangindent=1em\hangafter=1 \lean{PhaseClassification.phase_trichotomy}\newline\dirn{LEANEXT/phase/PhaseClassification.lean:285} & states the disjunction and Alternating $\to$ no run; mutual exclusion of runs not stated in Lean\\
Rem.~\ref{pc:rem:runs} & -- & written step: runs mutually exclusive under the ordering (not Lean-checked)\\
Thm.~\ref{pc:thm:rule-alt} & \hangindent=1em\hangafter=1 \lean{PhaseClassification.orient_pos_of_alternating}\newline\dirn{LEANEXT/phase/PhaseClassification.lean:301} & rule A; on the ENCBLEAN route\\
Thm.~\ref{pc:thm:rule-run} & \hangindent=1em\hangafter=1 \lean{PhaseClassification.orient_pos_of_runMidB}\newline\dirn{LEANEXT/phase/PhaseClassification.lean:310}\par \hangindent=1em\hangafter=1 \lean{PhaseClassification.orient_pos_of_runMidD}\newline\dirn{LEANEXT/phase/PhaseClassification.lean:318}\par \hangindent=1em\hangafter=1 \lean{PhaseClassification.orient_pos_of_runMidA}\newline\dirn{LEANEXT/phase/PhaseClassification.lean:342} & rule B; \texttt{runMidB} on the ENCBLEAN route\\
Thm.~\ref{pc:thm:phase} & \hangindent=1em\hangafter=1 \lean{PhaseClassification.orient_pos_of_phase}\newline\dirn{LEANEXT/phase/PhaseClassification.lean:372} & \\
Prop.~\ref{pc:prop:centre} & \hangindent=1em\hangafter=1 \lean{PhaseClassification.centre_inside_of_alternating}\newline\dirn{LEANEXT/phase/PhaseClassification.lean:414} & \\
Prop.~\ref{pc:prop:halfplane} & \hangindent=1em\hangafter=1 \lean{PhaseClassification.halfplane_of_runMidB}\newline\dirn{LEANEXT/phase/PhaseClassification.lean:431} & \\
Thm.~\ref{pc:thm:sign} & \hangindent=1em\hangafter=1 \lean{PhaseClassification.orient_sign_of_phase}\newline\dirn{LEANEXT/phase/PhaseClassification.lean:496} & case $\beta'=\delta'$ not covered; not combined with the layer theorems in Lean\\
Rem.~\ref{pc:rem:sign-scope} & -- & scope remark\\
Thm.~\ref{pc:thm:word} & \hangindent=1em\hangafter=1 \lean{PhaseClassification.word_dichotomy}\newline\dirn{LEANEXT/phase/PhaseClassification.lean:539} & computer-assisted in Lean; written proof by hand (8 words) given; Lean proof: decide (evaluated during elaboration and again by the kernel over the 64 functions Fin $6\to$ Bool; the conclusion only for the 8 satisfying the premise), no \texttt{native\_\allowbreak{}decide}\\
Rem.~\ref{pc:rem:word-scope} & -- & link between sorted directions and the Boolean word is not in LEANEXT (\texttt{LEANEXT/\allowbreak{}phase/\allowbreak{}REPORT.\allowbreak{}md}); Lean-checked encoder-B counterpart of the run link in ENCBLEAN (\texttt{EncB.\allowbreak{}wordSpec}, \texttt{EncB.\allowbreak{}Geo.\allowbreak{}Arc.\allowbreak{}arcFrame}, \texttt{EncB.\allowbreak{}Geo.\allowbreak{}labelling\_\allowbreak{}ruleB}) noted; their axioms: \#guard\_msgs pin at \texttt{ENCBLEAN/\allowbreak{}EncB/\allowbreak{}SoundWord.\allowbreak{}lean:228-229} (with \texttt{ENCBLEAN/\allowbreak{}logs/\allowbreak{}axioms.\allowbreak{}log:446}), \#print axioms output in \texttt{ENCBLEAN/\allowbreak{}logs/\allowbreak{}m8\_\allowbreak{}fix\_\allowbreak{}build.\allowbreak{}log:29,35}, and the \texttt{modelOfLabelling} pin at \texttt{ENCBLEAN/\allowbreak{}geo/\allowbreak{}ES8Geo/\allowbreak{}Audit.\allowbreak{}lean:45-46}\\
Thm.~\ref{pc:thm:layers} & \hangindent=1em\hangafter=1 \lean{ErdosSzekeres.LayerLemma.orient_pos_of_layers}\newline\dirn{LEANEXT/linked/Linked.lean:1398}\par \hangindent=1em\hangafter=1 \lean{ErdosSzekeres.Capstone.orient_pos_of_layers_runMidB}\newline\dirn{LEANEXT/capstone/Capstone.lean:1977}\par \hangindent=1em\hangafter=1 \lean{ErdosSzekeres.Capstone.orient_pos_of_layers_runMidD}\newline\dirn{LEANEXT/capstone/Capstone2.lean:2011}\par \hangindent=1em\hangafter=1 \lean{ErdosSzekeres.Capstone.orient_pos_of_layers_runMidA}\newline\dirn{LEANEXT/capstone/Capstone2.lean:2031}\par \hangindent=1em\hangafter=1 \lean{ErdosSzekeres.Capstone.orient_pos_of_layers_runMidB_intrinsic}\newline\dirn{LEANEXT/capstone/Capstone2.lean:2082}\par \hangindent=1em\hangafter=1 \lean{ErdosSzekeres.Capstone.orient_pos_of_layers_runMidD_intrinsic}\newline\dirn{LEANEXT/capstone/Capstone2.lean:2098}\par \hangindent=1em\hangafter=1 \lean{ErdosSzekeres.Capstone.orient_pos_of_layers_runMidA_intrinsic}\newline\dirn{LEANEXT/capstone/Capstone2.lean:2115} & uses \texttt{dist2\_\allowbreak{}le\_\allowbreak{}of\_\allowbreak{}later} (Theorem~\ref{lg:thm:radius}); no saved axiom output (reports state standard axioms); polar forms and the angle conditions ((a): $0<\beta-\alpha,0<\delta-\beta,\delta-\alpha<\pi$; (b)-(e): sorted angles, gap condition, pattern) are hypotheses; head cites (a), parts (b)-(e) cite the six Capstone declarations\\
Cor.~\ref{pc:cor:alt-polar} & \hangindent=1em\hangafter=1 \lean{ErdosSzekeres.Capstone.orient_pos_of_alternating_polar}\newline\dirn{LEANEXT/capstone/Capstone.lean:1992} & no saved axiom output (\texttt{LEANEXT/\allowbreak{}capstone/\allowbreak{}REPORT.\allowbreak{}md} states standard axioms)\\
Rem.~\ref{pc:rem:layers-scope} & -- & scope remark: what the layer theorems assume\\
\addlinespace
\multicolumn{3}{@{}l}{\normalsize\textbf{Section~\ref{ec:sec}}}\\[2pt]
Def.~\ref{ec:def:cases} & \hangindent=1em\hangafter=1 \lean{ErdosSzekeres.LayerLemma.Fivefold35}\newline\dirn{LEANEXT/es7corr/ES7Corollaries.lean:1157}\par \hangindent=1em\hangafter=1 \lean{ErdosSzekeres.LayerLemma.Fourfold32}\newline\dirn{LEANEXT/es7corr/ES7Corollaries.lean:1164} & the case classes: Prop definitions used only as hypotheses; never proved in Lean (supported outside Lean by SYM5 and C4CUT: computer-assisted, with written arguments)\\
Rem.~\ref{ec:rem:cases} & -- & status of the case classes \texttt{Fivefold35/\allowbreak{}Fourfold32}: SAT certificates (SYM5, C4CUT; only \texttt{cake\_\allowbreak{}lpr} verdicts count, lrat-check is secondary) plus written arguments outside Lean (encoder soundness, normal forms, centre perturbation; the records' reviews and refereeing were internal to the project); layer lemma and even-order ray fact have Lean-checked forms (Theorem~\ref{ll:thm:layers}; Lemma~\ref{oc:lem:same-ray}, Theorem~\ref{oc:thm:phase-orbit}) not linked to the formulas; distinct from the SYMLEAN constants\\
Thm.~\ref{ec:thm:fivefold-erase} & \hangindent=1em\hangafter=1 \lean{ErdosSzekeres.LayerLemma.fivefold_heptagon_of_erase}\newline\dirn{LEANEXT/es7corr/ES7Corollaries.lean:1170} & takes hypotheses: Fivefold35; proof via Theorem~\ref{ll:thm:reduction}\\
Thm.~\ref{ec:thm:fivefold} & \hangindent=1em\hangafter=1 \lean{ErdosSzekeres.LayerLemma.fivefold_heptagon}\newline\dirn{LEANEXT/es7corr/ES7Corollaries.lean:1183} & takes hypotheses: Fivefold35\\
Thm.~\ref{ec:thm:fivefold-nc} & \hangindent=1em\hangafter=1 \lean{ErdosSzekeres.LayerLemma.fivefold_heptagon_of_centre_not_mem}\newline\dirn{LEANEXT/es7corr/ES7Corollaries.lean:1193} & takes hypotheses: Fivefold35\\
Thm.~\ref{ec:thm:fourfold-erase} & \hangindent=1em\hangafter=1 \lean{ErdosSzekeres.LayerLemma.fourfold_heptagon_of_erase}\newline\dirn{LEANEXT/es7corr/ES7Corollaries.lean:1203} & takes hypotheses: Fourfold32; proof via Theorem~\ref{ll:thm:reduction}\\
Thm.~\ref{ec:thm:fourfold} & \hangindent=1em\hangafter=1 \lean{ErdosSzekeres.LayerLemma.fourfold_heptagon}\newline\dirn{LEANEXT/es7corr/ES7Corollaries.lean:1216} & takes hypotheses: Fourfold32; uses Lemma~\ref{lg:lem:centre}\\
Rem.~\ref{ec:rem:es7-sharp} & -- & conditional bounds 31 and 28; SYMMAX examples are exact computation outside Lean\\
Rem.~\ref{ec:rem:orbits} & -- & alternative route by orbit counting (not the Lean route)\\
Lem.~\ref{ec:lem:sphere} & \hangindent=1em\hangafter=1 \lean{RefuteC6.sphere_subset_extremePoints}\newline\dirn{ES8SYM/lean/C6check.lean:24} & over C; any real r\\
Lem.~\ref{ec:lem:ci-circle} & \hangindent=1em\hangafter=1 \lean{RefuteC6.convexIndependent_of_norm_eq}\newline\dirn{ES8SYM/lean/C6check.lean:43} & any index type\\
Def.~\ref{ec:def:zeta} & \hangindent=1em\hangafter=1 \lean{RefuteC6.zeta}\newline\dirn{ES8SYM/lean/C6check.lean:60} & $\exp(2\pi i/m)$, m cast to C\\
Lem.~\ref{ec:lem:zeta} & \hangindent=1em\hangafter=1 \lean{RefuteC6.zeta_prim}\newline\dirn{ES8SYM/lean/C6check.lean:62} & Mathlib \texttt{Complex.\allowbreak{}isPrimitiveRoot\_\allowbreak{}exp}\\
Fact~\ref{ec:fact:prim} & written step & written step: standard properties of primitive roots (Mathlib lemmas)\\
Lem.~\ref{ec:lem:orbit8} & \hangindent=1em\hangafter=1 \lean{RefuteC6.orbit_injective}\newline\dirn{ES8SYM/lean/C6check.lean:66}\par \hangindent=1em\hangafter=1 \lean{RefuteC6.orbit_convexIndependent}\newline\dirn{ES8SYM/lean/C6check.lean:75} & \\
Thm.~\ref{ec:thm:big-order} & \hangindent=1em\hangafter=1 \lean{RefuteC6.part_a}\newline\dirn{ES8SYM/lean/C6check.lean:85} & over C for an arbitrary invariant Set; no finiteness, no general position\\
Rem.~\ref{ec:rem:big-order} & -- & scope of \texttt{part\_\allowbreak{}a}; planar link not in Lean; a rotation of order m is reduced to rotAbout(c,2$\pi$/m) by a power\\
Cor.~\ref{ec:cor:planar} & written step & written step: planar form of \texttt{part\_\allowbreak{}a} (not in Lean)\\
Prop.~\ref{ec:prop:thr} & \hangindent=1em\hangafter=1 \lean{RefuteC6.thr_m4}\newline\dirn{ES8SYM/lean/C6check.lean:108}\par \hangindent=1em\hangafter=1 \lean{RefuteC6.thr_m5}\newline\dirn{ES8SYM/lean/C6check.lean:109}\par \hangindent=1em\hangafter=1 \lean{RefuteC6.thr_m6}\newline\dirn{ES8SYM/lean/C6check.lean:110}\par \hangindent=1em\hangafter=1 \lean{RefuteC6.thr_m7}\newline\dirn{ES8SYM/lean/C6check.lean:111}\par \hangindent=1em\hangafter=1 \lean{RefuteC6.m4_with_centre_counter}\newline\dirn{ES8SYM/lean/C6check.lean:115} & arithmetic (omega in Lean); \texttt{m4\_\allowbreak{}with\_\allowbreak{}centre\_\allowbreak{}counter} has no axiom line\\
Rem.~\ref{ec:rem:thr} & -- & geometric part of ES8SYM Lemma 2(b) is not in this file\\
Def.~\ref{ec:def:polar} & \hangindent=1em\hangafter=1 \lean{C7.cross}\newline\dirn{ES8SYM/lean/PolarOrient.lean:45}\par \hangindent=1em\hangafter=1 \lean{C7.pt}\newline\dirn{ES8SYM/lean/PolarOrient.lean:48}\par \hangindent=1em\hangafter=1 \lean{C7.coef}\newline\dirn{ES8SYM/lean/PolarOrient.lean:63}\par \hangindent=1em\hangafter=1 \lean{C7.phase}\newline\dirn{ES8SYM/lean/PolarOrient.lean:94}\par \hangindent=1em\hangafter=1 \lean{C7.ngonPt}\newline\dirn{ES8SYM/lean/PolarOrient.lean:97} & \\
Lem.~\ref{ec:lem:det-mul} & \hangindent=1em\hangafter=1 \lean{C7.det_mul_eq}\newline\dirn{ES8SYM/lean/PolarOrient.lean:24} & \\
Lem.~\ref{ec:lem:polar-u} & \hangindent=1em\hangafter=1 \lean{C7.orient_polar}\newline\dirn{ES8SYM/lean/PolarOrient.lean:37} & not the radius-form \texttt{orient\_\allowbreak{}polar} of LAYERLEAN\\
Lem.~\ref{ec:lem:polar-cross} & \hangindent=1em\hangafter=1 \lean{C7.orient_polar_cross}\newline\dirn{ES8SYM/lean/PolarOrient.lean:50} & \\
Prop.~\ref{ec:prop:layered} & \hangindent=1em\hangafter=1 \lean{C7.orient_layered}\newline\dirn{ES8SYM/lean/PolarOrient.lean:68} & \\
Cor.~\ref{ec:cor:scale} & \hangindent=1em\hangafter=1 \lean{C7.orient_layered_scale}\newline\dirn{ES8SYM/lean/PolarOrient.lean:80} & \\
Lem.~\ref{ec:lem:ngon-pt} & \hangindent=1em\hangafter=1 \lean{C7.ngonPt_eq_pt}\newline\dirn{ES8SYM/lean/PolarOrient.lean:100} & \\
Thm.~\ref{ec:thm:nested} & \hangindent=1em\hangafter=1 \lean{C7.nested_regular_orient}\newline\dirn{ES8SYM/lean/PolarOrient.lean:108} & assumes only $r_{o}>0$ (no nesting, no general position)\\
Lem.~\ref{ec:lem:centre-u} & \hangindent=1em\hangafter=1 \lean{C7.orient_centre}\newline\dirn{ES8SYM/lean/PolarOrient.lean:116} & \\
Rem.~\ref{ec:rem:polar} & -- & cross versus orient; distinct from Lemma~\ref{lg:lem:polar}\\
Lem.~\ref{ec:lem:pc-basic} & \hangindent=1em\hangafter=1 \lean{PCompat.ne_zero_of_order}\newline\dirn{PRIMEFOLD/compat/lean/PCompatCore.lean:36}\par \hangindent=1em\hangafter=1 \lean{PCompat.pow_inj_of_order}\newline\dirn{PRIMEFOLD/compat/lean/PCompatCore.lean:43}\par \hangindent=1em\hangafter=1 \lean{PCompat.pow_mul_mem}\newline\dirn{PRIMEFOLD/compat/lean/PCompatCore.lean:58} & commutative group with zero\\
Lem.~\ref{ec:lem:pc-exist} & \hangindent=1em\hangafter=1 \lean{PCompat.exists_card0}\newline\dirn{PRIMEFOLD/compat/lean/PCompatCore.lean:129}\par \hangindent=1em\hangafter=1 \lean{PCompat.exists_card1}\newline\dirn{PRIMEFOLD/compat/lean/PCompatCore.lean:160} & takes hypotheses hr0, hdisj on r (not discharged for C in Lean)\\
Def.~\ref{ec:def:compat} & \hangindent=1em\hangafter=1 \lean{PCompat.Compat}\newline\dirn{PRIMEFOLD/compat/lean/PCompatCore.lean:205} & \\
Prop.~\ref{ec:prop:parity} & \hangindent=1em\hangafter=1 \lean{PCompat.compat_even}\newline\dirn{PRIMEFOLD/compat/lean/PCompatCore.lean:208}\par \hangindent=1em\hangafter=1 \lean{PCompat.compat_odd}\newline\dirn{PRIMEFOLD/compat/lean/PCompatCore.lean:235} & \\
Thm.~\ref{ec:thm:compat} & \hangindent=1em\hangafter=1 \lean{PCompat.dvd_card}\newline\dirn{PRIMEFOLD/compat/lean/PCompatCore.lean:67}\par \hangindent=1em\hangafter=1 \lean{PCompat.card_mod}\newline\dirn{PRIMEFOLD/compat/lean/PCompatCore.lean:112}\par \hangindent=1em\hangafter=1 \lean{PCompat.exists_iff}\newline\dirn{PRIMEFOLD/compat/lean/PCompatCore.lean:174}\par \hangindent=1em\hangafter=1 \lean{PCompat.classification}\newline\dirn{PRIMEFOLD/compat/lean/PCompatCore.lean:253} & \texttt{exists\_\allowbreak{}iff} takes hypotheses hr0, hdisj on r\\
Prop.~\ref{ec:prop:ex341} & \hangindent=1em\hangafter=1 \lean{PCompat.ex341}\newline\dirn{PRIMEFOLD/compat/lean/PCompatCore.lean:277} & computer-assisted in Lean; written proof by hand (Lean: kernel decide; axioms propext only)\\
Rem.~\ref{ec:rem:compat} & -- & Fermat little theorem and general-position existence not formalized\\
Lem.~\ref{ec:lem:arith} & \hangindent=1em\hangafter=1 \lean{L0Check.arith}\newline\dirn{PRIMEFOLD/compat/lean/L0Check.lean:21} & \\
Lem.~\ref{ec:lem:orbit-dvd} & \hangindent=1em\hangafter=1 \lean{L0Check.orbit_dvd}\newline\dirn{PRIMEFOLD/compat/lean/L0Check.lean:60} & written proof derives it from Theorem~\ref{ec:thm:compat}(a); Lean re-proves it independently\\
Thm.~\ref{ec:thm:l0} & \hangindent=1em\hangafter=1 \lean{L0Check.L0}\newline\dirn{PRIMEFOLD/compat/lean/L0Check.lean:122}\par \hangindent=1em\hangafter=1 \lean{L0Check.L0_rot}\newline\dirn{PRIMEFOLD/compat/lean/L0Check.lean:144} & no general position assumed\\
Rem.~\ref{ec:rem:l0} & -- & parity excludes the centre for even m; Fermat not used\\
\addlinespace
\multicolumn{3}{@{}l}{\normalsize\textbf{Section~\ref{sr:sec}}}\\[2pt]
Def.~\ref{sr:def:iso} & \hangindent=1em\hangafter=1 \lean{ErdosSzekeres.SymReduction.eDist2}\newline\dirn{SYMLEAN/SymmetricReduction.lean:265}\par \hangindent=1em\hangafter=1 \lean{ErdosSzekeres.SymReduction.IsEucIsometry}\newline\dirn{SYMLEAN/SymmetricReduction.lean:268} & \\
Rem.~\ref{sr:rem:sup} & -- & remark: Mathlib sup metric on Point; why \texttt{eDist2} is used\\
Def.~\ref{sr:def:inv} & \hangindent=1em\hangafter=1 \lean{ErdosSzekeres.SymReduction.Invariant}\newline\dirn{SYMLEAN/SymmetricReduction.lean:271} & same notion as Definition~\ref{fd:def:inv} (maps into)\\
Def.~\ref{sr:def:complex} & \hangindent=1em\hangafter=1 \lean{ErdosSzekeres.SymReduction.toC}\newline\dirn{SYMLEAN/SymmetricReduction.lean:274}\par \hangindent=1em\hangafter=1 \lean{ErdosSzekeres.SymReduction.ofC}\newline\dirn{SYMLEAN/SymmetricReduction.lean:277} & \\
Lem.~\ref{sr:lem:complex} & \hangindent=1em\hangafter=1 \lean{ErdosSzekeres.SymReduction.toC_re}\newline\dirn{SYMLEAN/SymmetricReduction.lean:279}\par \hangindent=1em\hangafter=1 \lean{ErdosSzekeres.SymReduction.toC_im}\newline\dirn{SYMLEAN/SymmetricReduction.lean:280}\par \hangindent=1em\hangafter=1 \lean{ErdosSzekeres.SymReduction.ofC_zero}\newline\dirn{SYMLEAN/SymmetricReduction.lean:281}\par \hangindent=1em\hangafter=1 \lean{ErdosSzekeres.SymReduction.ofC_one}\newline\dirn{SYMLEAN/SymmetricReduction.lean:282}\par \hangindent=1em\hangafter=1 \lean{ErdosSzekeres.SymReduction.toC_ofC}\newline\dirn{SYMLEAN/SymmetricReduction.lean:284}\par \hangindent=1em\hangafter=1 \lean{ErdosSzekeres.SymReduction.ofC_toC}\newline\dirn{SYMLEAN/SymmetricReduction.lean:287}\par \hangindent=1em\hangafter=1 \lean{ErdosSzekeres.SymReduction.toC_injective}\newline\dirn{SYMLEAN/SymmetricReduction.lean:291}\par \hangindent=1em\hangafter=1 \lean{ErdosSzekeres.SymReduction.eDist2_eq_normSq}\newline\dirn{SYMLEAN/SymmetricReduction.lean:294}\par \hangindent=1em\hangafter=1 \lean{ErdosSzekeres.SymReduction.eDist2_pos_of_ne}\newline\dirn{SYMLEAN/SymmetricReduction.lean:297}\par \hangindent=1em\hangafter=1 \lean{ErdosSzekeres.SymReduction.point_ext}\newline\dirn{SYMLEAN/SymmetricReduction.lean:502} & \\
Def.~\ref{sr:def:maps} & \hangindent=1em\hangafter=1 \lean{ErdosSzekeres.SymReduction.crot}\newline\dirn{SYMLEAN/SymmetricReduction.lean:349}\par \hangindent=1em\hangafter=1 \lean{ErdosSzekeres.SymReduction.rotAbout}\newline\dirn{SYMLEAN/SymmetricReduction.lean:400}\par \hangindent=1em\hangafter=1 \lean{ErdosSzekeres.SymReduction.halfTurn}\newline\dirn{SYMLEAN/SymmetricReduction.lean:413}\par \hangindent=1em\hangafter=1 \lean{ErdosSzekeres.SymReduction.reflLine}\newline\dirn{SYMLEAN/SymmetricReduction.lean:423} & \texttt{rotAbout} is the formula of Definition~\ref{fd:def:rot}\\
Def.~\ref{sr:def:classes} & \hangindent=1em\hangafter=1 \lean{ErdosSzekeres.SymReduction.HalfTurn34}\newline\dirn{SYMLEAN/SymmetricReduction.lean:467}\par \hangindent=1em\hangafter=1 \lean{ErdosSzekeres.SymReduction.Threefold33}\newline\dirn{SYMLEAN/SymmetricReduction.lean:472}\par \hangindent=1em\hangafter=1 \lean{ErdosSzekeres.SymReduction.Fivefold35}\newline\dirn{SYMLEAN/SymmetricReduction.lean:477}\par \hangindent=1em\hangafter=1 \lean{ErdosSzekeres.SymReduction.MirrorOneAxis33}\newline\dirn{SYMLEAN/SymmetricReduction.lean:482}\par \hangindent=1em\hangafter=1 \lean{ErdosSzekeres.SymReduction.MirrorNoAxis34}\newline\dirn{SYMLEAN/SymmetricReduction.lean:488}\par \hangindent=1em\hangafter=1 \lean{ErdosSzekeres.SymReduction.Fourfold32}\newline\dirn{SYMLEAN/SymmetricReduction.lean:496} & Props used only as hypotheses; none proved in this section\\
Rem.~\ref{sr:rem:records} & -- & remark: project records supplying each class, with the rigor labels that the records claim (C4CUT states Fourfold32 without one)\\
Lem.~\ref{sr:lem:iterate} & \hangindent=1em\hangafter=1 \lean{ErdosSzekeres.SymReduction.iterate_mem}\newline\dirn{SYMLEAN/SymmetricReduction.lean:69}\par \hangindent=1em\hangafter=1 \lean{ErdosSzekeres.SymReduction.Invariant.iterate}\newline\dirn{SYMLEAN/SymmetricReduction.lean:309} & \\
Lem.~\ref{sr:lem:pigeon} & written step & written step (local statement pigeon in \texttt{isometry\_\allowbreak{}structure})\\
Def.~\ref{sr:def:minper} & written step & written step (Mathlib \texttt{Function.\allowbreak{}minimalPeriod})\\
Lem.~\ref{sr:lem:minper} & written step & written step (Mathlib minimal-period facts, proved)\\
Def.~\ref{sr:def:orbitFS} & \hangindent=1em\hangafter=1 \lean{ErdosSzekeres.SymReduction.orbitFS}\newline\dirn{SYMLEAN/SymmetricReduction.lean:65} & \\
Lem.~\ref{sr:lem:orbitFS} & \hangindent=1em\hangafter=1 \lean{ErdosSzekeres.SymReduction.card_orbitFS_of_minimalPeriod}\newline\dirn{SYMLEAN/SymmetricReduction.lean:111}\par \hangindent=1em\hangafter=1 \lean{ErdosSzekeres.SymReduction.mem_orbitFS_self}\newline\dirn{SYMLEAN/SymmetricReduction.lean:90}\par \hangindent=1em\hangafter=1 \lean{ErdosSzekeres.SymReduction.orbitFS_subset}\newline\dirn{SYMLEAN/SymmetricReduction.lean:93}\par \hangindent=1em\hangafter=1 \lean{ErdosSzekeres.SymReduction.orbitFS_invariant}\newline\dirn{SYMLEAN/SymmetricReduction.lean:99}\par \hangindent=1em\hangafter=1 \lean{ErdosSzekeres.SymReduction.card_orbitFS}\newline\dirn{SYMLEAN/SymmetricReduction.lean:76} & \\
Lem.~\ref{sr:lem:orbit-core} & written step & written step (the induction shared by \texttt{orbit\_\allowbreak{}select} and \texttt{orbit\_\allowbreak{}select\_\allowbreak{}minimalPeriod})\\
Lem.~\ref{sr:lem:orbit} & \hangindent=1em\hangafter=1 \lean{ErdosSzekeres.SymReduction.orbit_select_minimalPeriod}\newline\dirn{SYMLEAN/SymmetricReduction.lean:123}\par \hangindent=1em\hangafter=1 \lean{ErdosSzekeres.SymReduction.orbit_select}\newline\dirn{SYMLEAN/SymmetricReduction.lean:194} & \\
Lem.~\ref{sr:lem:crot} & \hangindent=1em\hangafter=1 \lean{ErdosSzekeres.SymReduction.toC_crot}\newline\dirn{SYMLEAN/SymmetricReduction.lean:351}\par \hangindent=1em\hangafter=1 \lean{ErdosSzekeres.SymReduction.crot_one}\newline\dirn{SYMLEAN/SymmetricReduction.lean:354}\par \hangindent=1em\hangafter=1 \lean{ErdosSzekeres.SymReduction.crot_comp}\newline\dirn{SYMLEAN/SymmetricReduction.lean:359}\par \hangindent=1em\hangafter=1 \lean{ErdosSzekeres.SymReduction.crot_iterate}\newline\dirn{SYMLEAN/SymmetricReduction.lean:364}\par \hangindent=1em\hangafter=1 \lean{ErdosSzekeres.SymReduction.crot_eq_self_iff}\newline\dirn{SYMLEAN/SymmetricReduction.lean:370}\par \hangindent=1em\hangafter=1 \lean{ErdosSzekeres.SymReduction.crot_injective}\newline\dirn{SYMLEAN/SymmetricReduction.lean:384}\par \hangindent=1em\hangafter=1 \lean{ErdosSzekeres.SymReduction.eDist2_crot_center}\newline\dirn{SYMLEAN/SymmetricReduction.lean:394} & \\
Lem.~\ref{sr:lem:rot-half} & \hangindent=1em\hangafter=1 \lean{ErdosSzekeres.SymReduction.rotAbout_eq_crot}\newline\dirn{SYMLEAN/SymmetricReduction.lean:404}\par \hangindent=1em\hangafter=1 \lean{ErdosSzekeres.SymReduction.halfTurn_eq_crot}\newline\dirn{SYMLEAN/SymmetricReduction.lean:415} & \\
Lem.~\ref{sr:lem:refl} & \hangindent=1em\hangafter=1 \lean{ErdosSzekeres.SymReduction.toC_reflLine}\newline\dirn{SYMLEAN/SymmetricReduction.lean:427}\par \hangindent=1em\hangafter=1 \lean{ErdosSzekeres.SymReduction.reflLine_involutive}\newline\dirn{SYMLEAN/SymmetricReduction.lean:435}\par \hangindent=1em\hangafter=1 \lean{ErdosSzekeres.SymReduction.cross_of_reflLine_fixed}\newline\dirn{SYMLEAN/SymmetricReduction.lean:447}\par \hangindent=1em\hangafter=1 \lean{ErdosSzekeres.SymReduction.orient_eq_zero_of_cross}\newline\dirn{SYMLEAN/SymmetricReduction.lean:455} & \\
Lem.~\ref{sr:lem:toC-iso} & \hangindent=1em\hangafter=1 \lean{ErdosSzekeres.SymReduction.IsEucIsometry.injective}\newline\dirn{SYMLEAN/SymmetricReduction.lean:301}\par \hangindent=1em\hangafter=1 \lean{ErdosSzekeres.SymReduction.isometry_toC_conj}\newline\dirn{SYMLEAN/SymmetricReduction.lean:313} & \\
Thm.~\ref{sr:thm:structure} & \hangindent=1em\hangafter=1 \lean{ErdosSzekeres.SymReduction.isometry_complex_cases}\newline\dirn{SYMLEAN/SymmetricReduction.lean:321}\par \hangindent=1em\hangafter=1 \lean{ErdosSzekeres.SymReduction.isometry_structure}\newline\dirn{SYMLEAN/SymmetricReduction.lean:869} & written proof of part (a) by an elementary argument instead of the Mathlib results used in Lean (a construction and a theorem)\\
Lem.~\ref{sr:lem:supports} & \hangindent=1em\hangafter=1 \lean{ErdosSzekeres.Suk.convexPosition_of_strict_supports}\newline\dirn{ErdosSzekeres/Suk/SupportingLines.lean:28} & main-library dependency of \texttt{convexPosition\_\allowbreak{}of\_\allowbreak{}concyclic}; same declaration as Lemma~\ref{fd:lem:strict-supports} (proof repeated)\\
Lem.~\ref{sr:lem:concyclic} & \hangindent=1em\hangafter=1 \lean{ErdosSzekeres.SymReduction.convexPosition_of_concyclic}\newline\dirn{SYMLEAN/SymmetricReduction.lean:509} & \\
Cor.~\ref{sr:cor:concyclic} & \hangindent=1em\hangafter=1 \lean{ErdosSzekeres.SymReduction.hasConvexNGon_of_concyclic}\newline\dirn{SYMLEAN/SymmetricReduction.lean:522} & \\
Lem.~\ref{sr:lem:center} & \hangindent=1em\hangafter=1 \lean{ErdosSzekeres.SymReduction.center_not_mem_of_halfTurn}\newline\dirn{SYMLEAN/SymmetricReduction.lean:528} & \\
Def.~\ref{sr:def:prim} & written step & written step (Mathlib \texttt{IsPrimitiveRoot}, \texttt{orderOf})\\
Lem.~\ref{sr:lem:prim} & written step & written step (Mathlib facts on roots of unity, proved)\\
Lem.~\ref{sr:lem:rot345} & \hangindent=1em\hangafter=1 \lean{ErdosSzekeres.SymReduction.isPrimitiveRoot_rot3}\newline\dirn{SYMLEAN/SymmetricReduction.lean:584}\par \hangindent=1em\hangafter=1 \lean{ErdosSzekeres.SymReduction.isPrimitiveRoot_rot5}\newline\dirn{SYMLEAN/SymmetricReduction.lean:591}\par \hangindent=1em\hangafter=1 \lean{ErdosSzekeres.SymReduction.isPrimitiveRoot_rot4}\newline\dirn{SYMLEAN/SymmetricReduction.lean:598} & \\
Lem.~\ref{sr:lem:root-cases} & \hangindent=1em\hangafter=1 \lean{ErdosSzekeres.SymReduction.root_cases}\newline\dirn{SYMLEAN/SymmetricReduction.lean:552} & \\
Lem.~\ref{sr:lem:prime-cases} & \hangindent=1em\hangafter=1 \lean{ErdosSzekeres.SymReduction.prime_cases}\newline\dirn{SYMLEAN/SymmetricReduction.lean:572} & \\
Lem.~\ref{sr:lem:root-cases2} & \hangindent=1em\hangafter=1 \lean{ErdosSzekeres.SymReduction.root_cases2}\newline\dirn{SYMLEAN/SymmetricReduction.lean:607} & \\
Lem.~\ref{sr:lem:halfturn-facts} & \hangindent=1em\hangafter=1 \lean{ErdosSzekeres.SymReduction.halfTurn_facts}\newline\dirn{SYMLEAN/SymmetricReduction.lean:647} & \\
Lem.~\ref{sr:lem:case-halfturn} & \hangindent=1em\hangafter=1 \lean{ErdosSzekeres.SymReduction.case_halfTurn}\newline\dirn{SYMLEAN/SymmetricReduction.lean:670} & conditional; takes hypotheses: \texttt{HalfTurn34}\\
Lem.~\ref{sr:lem:rot-erase} & \hangindent=1em\hangafter=1 \lean{ErdosSzekeres.SymReduction.rot_erase_facts}\newline\dirn{SYMLEAN/SymmetricReduction.lean:676} & \\
Lem.~\ref{sr:lem:rot-big} & \hangindent=1em\hangafter=1 \lean{ErdosSzekeres.SymReduction.case_rot_big}\newline\dirn{SYMLEAN/SymmetricReduction.lean:711} & no general position assumed\\
Lem.~\ref{sr:lem:odd-prime} & \hangindent=1em\hangafter=1 \lean{ErdosSzekeres.SymReduction.case_odd_prime}\newline\dirn{SYMLEAN/SymmetricReduction.lean:736} & conditional; takes hypotheses: Threefold33, Fivefold35\\
Lem.~\ref{sr:lem:fourfold} & \hangindent=1em\hangafter=1 \lean{ErdosSzekeres.SymReduction.case_fourfold}\newline\dirn{SYMLEAN/SymmetricReduction.lean:763} & conditional; takes hypotheses: Fourfold32\\
Lem.~\ref{sr:lem:mirror} & \hangindent=1em\hangafter=1 \lean{ErdosSzekeres.SymReduction.case_mirror}\newline\dirn{SYMLEAN/SymmetricReduction.lean:811} & conditional; takes hypotheses: \texttt{MirrorOneAxis33}, \texttt{MirrorNoAxis34}\\
Thm.~\ref{sr:thm:main} & \hangindent=1em\hangafter=1 \lean{ErdosSzekeres.SymReduction.heptagon_of_nontrivial_isometry}\newline\dirn{SYMLEAN/SymmetricReduction.lean:997} & conditional; takes hypotheses: \texttt{HalfTurn34}, Threefold33, Fivefold35, \texttt{MirrorOneAxis33}, \texttt{MirrorNoAxis34}; axioms log \texttt{SYMLEAN/\allowbreak{}logs/\allowbreak{}SymmetricReductionAxioms.\allowbreak{}out}\\
Rem.~\ref{sr:rem:main} & -- & remark: conditional status of the main theorem; the unconditional statement is a project claim (\texttt{WRITEUP/\allowbreak{}SYMMETRIC.\allowbreak{}md}, Theorem A), established here only as the implication\\
Thm.~\ref{sr:thm:variant} & \hangindent=1em\hangafter=1 \lean{ErdosSzekeres.SymReduction.halfTurn_of_symmetric_counterexample}\newline\dirn{SYMLEAN/SymmetricReduction.lean:1024} & conditional; takes hypotheses: Threefold33, Fourfold32, Fivefold35, \texttt{MirrorOneAxis33}, \texttt{MirrorNoAxis34}; axioms log \texttt{SYMLEAN/\allowbreak{}logs/\allowbreak{}SymmetricReductionAxioms.\allowbreak{}out}\\
Rem.~\ref{sr:rem:variant} & -- & remark: per-symmetry, isometries only; the group form \texttt{VariantChecks.\allowbreak{}group\_\allowbreak{}form} (review file \texttt{SYMLEAN/\allowbreak{}review/\allowbreak{}VariantChecks.\allowbreak{}lean:1168}, log \texttt{SYMLEAN/\allowbreak{}logs/\allowbreak{}VariantChecks.\allowbreak{}out}) with a written argument, not restated as a numbered statement\\
Def.~\ref{sr:def:exclusion-line} & \hangindent=1em\hangafter=1 \lean{ErdosSzekeres.SymReduction.MirrorExclusionLineCopy}\newline\dirn{SYMLEAN/SymmetricReduction.lean:1052} & text of \texttt{Mir.\allowbreak{}MirrorExclusionLine} (\texttt{H3/\allowbreak{}MirrorEncoding/\allowbreak{}Line.\allowbreak{}lean:94})\\
Lem.~\ref{sr:lem:mirror-line} & \hangindent=1em\hangafter=1 \lean{ErdosSzekeres.SymReduction.mirrorOneAxis33_of_exclusionLine}\newline\dirn{SYMLEAN/SymmetricReduction.lean:1056} & conditional; takes hypotheses: \texttt{MirrorExclusionLineCopy}\\
Rem.~\ref{sr:rem:mirror-line} & -- & remark: the H3 statement is Lean-checked only conditionally on CNF facts (\texttt{V2L.\allowbreak{}mirrorExclusionLine\_\allowbreak{}final}) and is identified with \texttt{MirrorExclusionLineCopy} by its text; the converse and the equivalence (review file \texttt{SYMLEAN/\allowbreak{}review/\allowbreak{}ReviewSanity.\allowbreak{}lean:1593} and 1629) with a written argument, not restated as numbered statements\\
Rem.~\ref{sr:rem:redundant} & -- & remark: hypotheses of the classes implied by the others (not used)\\
\addlinespace
\multicolumn{3}{@{}l}{\normalsize\textbf{Section~\ref{ar:sec}}}\\[2pt]
Def.~\ref{ar:def:nsq} & \hangindent=1em\hangafter=1 \lean{ErdosSzekeres.SymReduction.nsq}\newline\dirn{SYMLEAN/AffineReduction.lean:1088} & \\
Lem.~\ref{ar:lem:mulvec} & \hangindent=1em\hangafter=1 \lean{ErdosSzekeres.SymReduction.mulVec2_fst}\newline\dirn{SYMLEAN/AffineReduction.lean:1097}\par \hangindent=1em\hangafter=1 \lean{ErdosSzekeres.SymReduction.mulVec2_snd}\newline\dirn{SYMLEAN/AffineReduction.lean:1101} & \\
Lem.~\ref{ar:lem:nsq} & \hangindent=1em\hangafter=1 \lean{ErdosSzekeres.SymReduction.nsq_nonneg}\newline\dirn{SYMLEAN/AffineReduction.lean:1090}\par \hangindent=1em\hangafter=1 \lean{ErdosSzekeres.SymReduction.eDist2_eq_nsq}\newline\dirn{SYMLEAN/AffineReduction.lean:1094} & \\
Def.~\ref{ar:def:gram} & \hangindent=1em\hangafter=1 \lean{ErdosSzekeres.SymReduction.gA}\newline\dirn{SYMLEAN/AffineReduction.lean:1106}\par \hangindent=1em\hangafter=1 \lean{ErdosSzekeres.SymReduction.gB}\newline\dirn{SYMLEAN/AffineReduction.lean:1107}\par \hangindent=1em\hangafter=1 \lean{ErdosSzekeres.SymReduction.gC}\newline\dirn{SYMLEAN/AffineReduction.lean:1108} & \\
Lem.~\ref{ar:lem:nsq-mulvec} & \hangindent=1em\hangafter=1 \lean{ErdosSzekeres.SymReduction.nsq_mulVec}\newline\dirn{SYMLEAN/AffineReduction.lean:1110} & \\
Def.~\ref{ar:def:avgq} & \hangindent=1em\hangafter=1 \lean{ErdosSzekeres.SymReduction.avgQ}\newline\dirn{SYMLEAN/AffineReduction.lean:1116}\par \hangindent=1em\hangafter=1 \lean{ErdosSzekeres.SymReduction.qA}\newline\dirn{SYMLEAN/AffineReduction.lean:1120}\par \hangindent=1em\hangafter=1 \lean{ErdosSzekeres.SymReduction.qB}\newline\dirn{SYMLEAN/AffineReduction.lean:1121}\par \hangindent=1em\hangafter=1 \lean{ErdosSzekeres.SymReduction.qC}\newline\dirn{SYMLEAN/AffineReduction.lean:1122} & \\
Lem.~\ref{ar:lem:avgq-eq} & \hangindent=1em\hangafter=1 \lean{ErdosSzekeres.SymReduction.avgQ_eq}\newline\dirn{SYMLEAN/AffineReduction.lean:1124} & \\
Lem.~\ref{ar:lem:avgq-pos} & \hangindent=1em\hangafter=1 \lean{ErdosSzekeres.SymReduction.avgQ_ge}\newline\dirn{SYMLEAN/AffineReduction.lean:1133}\par \hangindent=1em\hangafter=1 \lean{ErdosSzekeres.SymReduction.qA_pos}\newline\dirn{SYMLEAN/AffineReduction.lean:1152}\par \hangindent=1em\hangafter=1 \lean{ErdosSzekeres.SymReduction.qdisc_pos}\newline\dirn{SYMLEAN/AffineReduction.lean:1158} & each part assumes $0<n$\\
Lem.~\ref{ar:lem:avgq-inv} & \hangindent=1em\hangafter=1 \lean{ErdosSzekeres.SymReduction.avgQ_mulVec}\newline\dirn{SYMLEAN/AffineReduction.lean:1139} & assumes $L^{n}=1$ only (no positivity of n)\\
Def.~\ref{ar:def:lin2} & \hangindent=1em\hangafter=1 \lean{ErdosSzekeres.SymReduction.lin2}\newline\dirn{SYMLEAN/AffineReduction.lean:1167} & bundled linear map\\
Lem.~\ref{ar:lem:lin2} & \hangindent=1em\hangafter=1 \lean{ErdosSzekeres.SymReduction.lin2_fst}\newline\dirn{SYMLEAN/AffineReduction.lean:1174}\par \hangindent=1em\hangafter=1 \lean{ErdosSzekeres.SymReduction.lin2_snd}\newline\dirn{SYMLEAN/AffineReduction.lean:1175}\par \hangindent=1em\hangafter=1 \lean{ErdosSzekeres.SymReduction.orient_lin2}\newline\dirn{SYMLEAN/AffineReduction.lean:1177} & \\
Lem.~\ref{ar:lem:chol} & \hangindent=1em\hangafter=1 \lean{ErdosSzekeres.SymReduction.nsq_lin2_chol}\newline\dirn{SYMLEAN/AffineReduction.lean:1184}\par \hangindent=1em\hangafter=1 \lean{ErdosSzekeres.SymReduction.lin2_chol_inv_left}\newline\dirn{SYMLEAN/AffineReduction.lean:1191}\par \hangindent=1em\hangafter=1 \lean{ErdosSzekeres.SymReduction.lin2_chol_inv_right}\newline\dirn{SYMLEAN/AffineReduction.lean:1195} & (a) assumes $s\ne0$; (b),(c) assume $s\ne0$ and $t\ne0$\\
Lem.~\ref{ar:lem:iter} & \hangindent=1em\hangafter=1 \lean{ErdosSzekeres.SymReduction.iterate_sub_of_affine}\newline\dirn{SYMLEAN/AffineReduction.lean:1199} & \\
Lem.~\ref{ar:lem:periodic} & \hangindent=1em\hangafter=1 \lean{ErdosSzekeres.SymReduction.exists_periodic_of_injective}\newline\dirn{SYMLEAN/AffineReduction.lean:1208} & \\
Lem.~\ref{ar:lem:fix} & \hangindent=1em\hangafter=1 \lean{ErdosSzekeres.SymReduction.matrix_eq_one_of_fix}\newline\dirn{SYMLEAN/AffineReduction.lean:1217} & \\
Lem.~\ref{ar:lem:orient-affine} & \hangindent=1em\hangafter=1 \lean{ErdosSzekeres.SymReduction.orient_affine}\newline\dirn{SYMLEAN/AffineReduction.lean:1241} & \\
Lem.~\ref{ar:lem:inj} & \hangindent=1em\hangafter=1 \lean{ErdosSzekeres.SymReduction.injective_of_affine}\newline\dirn{SYMLEAN/AffineReduction.lean:1305} & \\
Lem.~\ref{ar:lem:matrix-form} & \hangindent=1em\hangafter=1 \lean{ErdosSzekeres.SymReduction.affineMap_eq_mulVec}\newline\dirn{SYMLEAN/AffineReduction.lean:1284} & \\
Lem.~\ref{ar:lem:gp-orient} & \hangindent=1em\hangafter=1 \lean{ErdosSzekeres.SymReduction.generalPosition_image_of_orient}\newline\dirn{SYMLEAN/AffineReduction.lean:1249} & injectivity of f not assumed\\
Lem.~\ref{ar:lem:gp-affine} & \hangindent=1em\hangafter=1 \lean{ErdosSzekeres.SymReduction.generalPosition_image_affine}\newline\dirn{SYMLEAN/AffineReduction.lean:1278}\par \hangindent=1em\hangafter=1 \lean{ErdosSzekeres.SymReduction.generalPosition_image_affineEquiv}\newline\dirn{SYMLEAN/AffineReduction.lean:1290} & \\
Lem.~\ref{ar:lem:cp-affine} & \hangindent=1em\hangafter=1 \lean{ErdosSzekeres.SymReduction.convexPosition_image_affine}\newline\dirn{SYMLEAN/AffineReduction.lean:1263} & \\
Lem.~\ref{ar:lem:ngon-affine} & \hangindent=1em\hangafter=1 \lean{ErdosSzekeres.SymReduction.hasConvexNGon_image_affine}\newline\dirn{SYMLEAN/AffineReduction.lean:1298} & \\
Def.~\ref{ar:def:E} & written step & written step: the proposition E(N,n), the type of the hypothesis \texttt{hE} (no Lean name)\\
Prop.~\ref{ar:prop:averaging} & written step & written step; $L^{m}=1$ and parts (b),(c) and $g\ne\mathrm{id}\Rightarrow h\ne\mathrm{id}$ are steps inside the Lean proof of \texttt{hasConvexNGon\_\allowbreak{}of\_\allowbreak{}affine\_\allowbreak{}symmetry}; $g^{m}=\mathrm{id}$ and the rest of (d) (orders) are not in Lean\\
Thm.~\ref{ar:thm:transfer} & \hangindent=1em\hangafter=1 \lean{ErdosSzekeres.SymReduction.hasConvexNGon_of_affine_symmetry}\newline\dirn{SYMLEAN/AffineReduction.lean:1330} & takes hypothesis \texttt{hE} (the Euclidean statement E(N,n) for the same N and n, a parameter) and $3\le N$\\
Rem.~\ref{ar:rem:hE} & -- & role of \texttt{hE}; order of g not tracked; order-preserving form not formalized\\
Cor.~\ref{ar:cor:transfer-forms} & \hangindent=1em\hangafter=1 \lean{ErdosSzekeres.SymReduction.hasConvexNGon_of_affineMap_symmetry}\newline\dirn{SYMLEAN/AffineReduction.lean:1423}\par \hangindent=1em\hangafter=1 \lean{ErdosSzekeres.SymReduction.hasConvexNGon_of_affineEquiv_symmetry}\newline\dirn{SYMLEAN/AffineReduction.lean:1433}\par \hangindent=1em\hangafter=1 \lean{ErdosSzekeres.SymReduction.hasConvexNGon_of_affine_coords_symmetry}\newline\dirn{SYMLEAN/AffineReduction.lean:1444} & takes hypothesis \texttt{hE} (E(N,n)) and $3\le N$\\
Thm.~\ref{ar:thm:main} & \hangindent=1em\hangafter=1 \lean{ErdosSzekeres.SymReduction.heptagon_of_nontrivial_affineEquiv}\newline\dirn{SYMLEAN/AffineReduction.lean:1501}\par \hangindent=1em\hangafter=1 \lean{ErdosSzekeres.SymReduction.heptagon_of_nontrivial_affineEquiv_image}\newline\dirn{SYMLEAN/AffineReduction.lean:1510}\par \hangindent=1em\hangafter=1 \lean{ErdosSzekeres.SymReduction.heptagon_of_nontrivial_affineMap}\newline\dirn{SYMLEAN/AffineReduction.lean:1489}\par \hangindent=1em\hangafter=1 \lean{ErdosSzekeres.SymReduction.heptagon_of_nontrivial_affine}\newline\dirn{SYMLEAN/AffineReduction.lean:1477}\par \hangindent=1em\hangafter=1 \lean{ErdosSzekeres.SymReduction.heptagon_of_nontrivial_affine_coords}\newline\dirn{SYMLEAN/AffineReduction.lean:1521} & takes hypotheses: \texttt{HalfTurn34}, Threefold33, Fivefold35, \texttt{MirrorOneAxis33}, \texttt{MirrorNoAxis34}\\
Rem.~\ref{ar:rem:classes} & -- & what supplies the five class hypotheses and with which rigor label; the mirror theorems concern byte-identical copies in \texttt{SYMLEAN/\allowbreak{}mirror/\allowbreak{}ReductionThreeClasses.\allowbreak{}lean}; the unconditional statement is a project claim (\texttt{WRITEUP/\allowbreak{}SYMMETRIC.\allowbreak{}md}, Theorem A), established here only conditionally\\
Rem.~\ref{ar:rem:scope} & -- & what the statements do not say\\
\addlinespace
\multicolumn{3}{@{}l}{\normalsize\textbf{Section~\ref{mm:sec}}}\\[2pt]
Def.~\ref{mm:def:mir} & \hangindent=1em\hangafter=1 \lean{MirrorOdd.mir}\newline\dirn{MIRRORETV/lean/MirrorMain.lean:16}\par \hangindent=1em\hangafter=1 \lean{MirrorOdd.Rset}\newline\dirn{MIRRORETV/lean/MirrorMain.lean:1598}\par \hangindent=1em\hangafter=1 \lean{MirrorOdd.Aset}\newline\dirn{MIRRORETV/lean/MirrorMain.lean:1599} & \\
Lem.~\ref{mm:lem:orient-mir} & \hangindent=1em\hangafter=1 \lean{MirrorOdd.mir_zero}\newline\dirn{MIRRORETV/lean/MirrorMain.lean:18}\par \hangindent=1em\hangafter=1 \lean{MirrorOdd.mir_one}\newline\dirn{MIRRORETV/lean/MirrorMain.lean:19}\par \hangindent=1em\hangafter=1 \lean{MirrorOdd.mir_mir}\newline\dirn{MIRRORETV/lean/MirrorMain.lean:1543}\par \hangindent=1em\hangafter=1 \lean{MirrorOdd.orient_mir3}\newline\dirn{MIRRORETV/lean/MirrorMain.lean:2187}\par \hangindent=1em\hangafter=1 \lean{MirrorOdd.orient_mir2}\newline\dirn{MIRRORETV/lean/MirrorMain.lean:2190}\par \hangindent=1em\hangafter=1 \lean{MirrorOdd.orient_cycl}\newline\dirn{MIRRORETV/lean/MirrorMain.lean:2193}\par \hangindent=1em\hangafter=1 \lean{MirrorOdd.orient_swap12}\newline\dirn{MIRRORETV/lean/MirrorMain.lean:2196}\par \hangindent=1em\hangafter=1 \lean{MirrorOdd.orient_rev}\newline\dirn{MIRRORETV/lean/MirrorMain.lean:2199} & part (e) is a written step (identity used inside \texttt{cap\_\allowbreak{}below\_\allowbreak{}lines})\\
Def.~\ref{mm:def:config} & \hangindent=1em\hangafter=1 \lean{ErdosSzekeres.Omission.Configuration}\newline\dirn{ErdosSzekeres/Omission/Realization.lean:9}\par \hangindent=1em\hangafter=1 \lean{ErdosSzekeres.Omission.Configuration.HasCap}\newline\dirn{ErdosSzekeres/Omission/Realization.lean:64}\par \hangindent=1em\hangafter=1 \lean{MirrorOdd.xset}\newline\dirn{MIRRORETV/lean/MirrorMain.lean:1549}\par \hangindent=1em\hangafter=1 \lean{MirrorOdd.emb}\newline\dirn{MIRRORETV/lean/MirrorMain.lean:1553}\par \hangindent=1em\hangafter=1 \lean{MirrorOdd.sortPt}\newline\dirn{MIRRORETV/lean/MirrorMain.lean:1558}\par \hangindent=1em\hangafter=1 \lean{MirrorOdd.cfg}\newline\dirn{MIRRORETV/lean/MirrorMain.lean:1584} & \texttt{HasCap} is the all-triples cap\\
Lem.~\ref{mm:lem:sort} & \hangindent=1em\hangafter=1 \lean{MirrorOdd.xset_card}\newline\dirn{MIRRORETV/lean/MirrorMain.lean:1551}\par \hangindent=1em\hangafter=1 \lean{MirrorOdd.emb_mem}\newline\dirn{MIRRORETV/lean/MirrorMain.lean:1555}\par \hangindent=1em\hangafter=1 \lean{MirrorOdd.sortPt_mem}\newline\dirn{MIRRORETV/lean/MirrorMain.lean:1560}\par \hangindent=1em\hangafter=1 \lean{MirrorOdd.sortPt_x}\newline\dirn{MIRRORETV/lean/MirrorMain.lean:1563}\par \hangindent=1em\hangafter=1 \lean{MirrorOdd.sortPt_strictMono}\newline\dirn{MIRRORETV/lean/MirrorMain.lean:1566}\par \hangindent=1em\hangafter=1 \lean{MirrorOdd.sortPt_injective}\newline\dirn{MIRRORETV/lean/MirrorMain.lean:1572}\par \hangindent=1em\hangafter=1 \lean{MirrorOdd.sortPt_image}\newline\dirn{MIRRORETV/lean/MirrorMain.lean:1576} & part (h) is the \texttt{generalPosition} field of \texttt{MirrorOdd.\allowbreak{}cfg} (\texttt{MIRRORETV/\allowbreak{}lean/\allowbreak{}MirrorMain.\allowbreak{}lean:1587})\\
Lem.~\ref{mm:lem:ngon} & \hangindent=1em\hangafter=1 \lean{MirrorOdd.hasConvexNGon_of_le}\newline\dirn{MIRRORETV/lean/MirrorMain.lean:2182}\par \hangindent=1em\hangafter=1 \lean{MirrorOdd.hasConvexNGon_affineMap}\newline\dirn{MIRRORETV/lean/MirrorMain.lean:1859} & \\
Lem.~\ref{mm:lem:slope} & \hangindent=1em\hangafter=1 \lean{MirrorOdd.cap_below_lines}\newline\dirn{MIRRORETV/lean/MirrorMain.lean:2204} & part (b) is a written step (inline in \texttt{mirror\_\allowbreak{}odd\_\allowbreak{}general} at \texttt{MIRRORETV/\allowbreak{}lean/\allowbreak{}MirrorMain.\allowbreak{}lean:2542-2557}; local claim interior at \texttt{MIRRORETV/\allowbreak{}lean/\allowbreak{}MirrorMain.\allowbreak{}lean:2060} and \texttt{MIRRORETV/\allowbreak{}lean/\allowbreak{}MirrorMain.\allowbreak{}lean:1716})\\
Thm.~\ref{mm:thm:mirror-polygon} & \hangindent=1em\hangafter=1 \lean{MirrorOdd.cap_mirror_polygon}\newline\dirn{MIRRORETV/lean/MirrorMain.lean:2233} & proof follows Lean (every triple of an explicit cyclic order, taken in list order, positive)\\
Thm.~\ref{mm:thm:axis-polygon} & \hangindent=1em\hangafter=1 \lean{MirrorOdd.axis_left_polygon}\newline\dirn{MIRRORETV/lean/MirrorMain.lean:2309}\par \hangindent=1em\hangafter=1 \lean{MirrorOdd.axis_right_polygon}\newline\dirn{MIRRORETV/lean/MirrorMain.lean:2403} & proof follows Lean\\
Lem.~\ref{mm:lem:cyc} & \hangindent=1em\hangafter=1 \lean{MirrorOdd.cyc4}\newline\dirn{MIRRORETV/lean/MirrorMain.lean:995}\par \hangindent=1em\hangafter=1 \lean{MirrorOdd.cyc5}\newline\dirn{MIRRORETV/lean/MirrorMain.lean:1026}\par \hangindent=1em\hangafter=1 \lean{MirrorOdd.cyc6}\newline\dirn{MIRRORETV/lean/MirrorMain.lean:1068}\par \hangindent=1em\hangafter=1 \lean{MirrorOdd.cyc7}\newline\dirn{MIRRORETV/lean/MirrorMain.lean:22}\par \hangindent=1em\hangafter=1 \lean{MirrorOdd.cyc8}\newline\dirn{MIRRORETV/lean/MirrorMain.lean:103} & machine-generated Lean code; instance of Lemma~\ref{fd:lem:pos-cycle}\\
Prop.~\ref{mm:prop:small-polygons} & \hangindent=1em\hangafter=1 \lean{MirrorOdd.cap_mirror_quadrilateral}\newline\dirn{MIRRORETV/lean/MirrorMain.lean:1126}\par \hangindent=1em\hangafter=1 \lean{MirrorOdd.cap_mirror_pentagon}\newline\dirn{MIRRORETV/lean/MirrorMain.lean:1282}\par \hangindent=1em\hangafter=1 \lean{MirrorOdd.cap_mirror_hexagon}\newline\dirn{MIRRORETV/lean/MirrorMain.lean:1168}\par \hangindent=1em\hangafter=1 \lean{MirrorOdd.cap_mirror_heptagon}\newline\dirn{MIRRORETV/lean/MirrorMain.lean:212}\par \hangindent=1em\hangafter=1 \lean{MirrorOdd.cap_mirror_octagon}\newline\dirn{MIRRORETV/lean/MirrorMain.lean:392}\par \hangindent=1em\hangafter=1 \lean{MirrorOdd.axis_left_pentagon}\newline\dirn{MIRRORETV/lean/MirrorMain.lean:1358}\par \hangindent=1em\hangafter=1 \lean{MirrorOdd.axis_right_pentagon}\newline\dirn{MIRRORETV/lean/MirrorMain.lean:1433}\par \hangindent=1em\hangafter=1 \lean{MirrorOdd.axis_left_heptagon}\newline\dirn{MIRRORETV/lean/MirrorMain.lean:639}\par \hangindent=1em\hangafter=1 \lean{MirrorOdd.axis_right_heptagon}\newline\dirn{MIRRORETV/lean/MirrorMain.lean:816} & machine-generated Lean code (linarith certificates); written proof differs: derived from Theorem~\ref{mm:thm:mirror-polygon} and Theorem~\ref{mm:thm:axis-polygon}\\
Lem.~\ref{mm:lem:count} & \hangindent=1em\hangafter=1 \lean{MirrorOdd.card_split}\newline\dirn{MIRRORETV/lean/MirrorMain.lean:1601}\par \hangindent=1em\hangafter=1 \lean{MirrorOdd.card_A_le}\newline\dirn{MIRRORETV/lean/MirrorMain.lean:1636} & \\
Lem.~\ref{mm:lem:injx} & \hangindent=1em\hangafter=1 \lean{MirrorOdd.injOn_x}\newline\dirn{MIRRORETV/lean/MirrorMain.lean:1647} & \\
Lem.~\ref{mm:lem:upper-caps} & written step & written step; isolates the cap argument of the local claims \texttt{keyT} and \texttt{hR} in \texttt{mirror\_\allowbreak{}odd\_\allowbreak{}general}, \texttt{mirror\_\allowbreak{}even\_\allowbreak{}general}, \texttt{mirror\_\allowbreak{}five}, \texttt{mirror\_\allowbreak{}six}\\
Thm.~\ref{mm:thm:odd} & \hangindent=1em\hangafter=1 \lean{MirrorOdd.mirror_odd_general}\newline\dirn{MIRRORETV/lean/MirrorMain.lean:2498} & takes a parametric hypothesis: \texttt{hM} (an assumed bound M for configurations with no (j+1)-cap and no convex (2j+1)-gon, not an unproved claim; instances discharged by Lean-checked theorems, see Remark~\ref{mm:rem:hm}); Lean j (used in this section) $=$ \texttt{README/\allowbreak{}PROOFS} j minus 1\\
Thm.~\ref{mm:thm:even} & \hangindent=1em\hangafter=1 \lean{MirrorOdd.mirror_even_general}\newline\dirn{MIRRORETV/lean/MirrorMain.lean:2599} & takes a parametric hypothesis: \texttt{hM} (an assumed bound M for configurations with no j-cap and no convex 2j-gon, not an unproved claim; instances discharged by Lean-checked theorems, see Remark~\ref{mm:rem:hm})\\
Rem.~\ref{mm:rem:hm} & -- & remark (no statement): (\texttt{hM}) is a parametric hypothesis (an assumption on M, not an unproved claim); Lean-checked instances (3,16), (4,22), (4,113) via \texttt{MirrorOdd.\allowbreak{}capacity16}, \texttt{MirrorOdd.\allowbreak{}capacity22g}, \texttt{C12Check.\allowbreak{}capacity113} (Baek; cups-caps for (4,113)); small j; relation to N(n,a,b); index shift\\
Def.~\ref{mm:def:frame} & \hangindent=1em\hangafter=1 \lean{MirrorOdd.lineRefl}\newline\dirn{MIRRORETV/lean/MirrorMain.lean:1879}\par \hangindent=1em\hangafter=1 \lean{MirrorOdd.frame}\newline\dirn{MIRRORETV/lean/MirrorMain.lean:1883}\par \hangindent=1em\hangafter=1 \lean{MirrorOdd.unframeLin}\newline\dirn{MIRRORETV/lean/MirrorMain.lean:1886}\par \hangindent=1em\hangafter=1 \lean{MirrorOdd.unframe}\newline\dirn{MIRRORETV/lean/MirrorMain.lean:1892} & \\
Lem.~\ref{mm:lem:frame} & \hangindent=1em\hangafter=1 \lean{MirrorOdd.unframe_apply}\newline\dirn{MIRRORETV/lean/MirrorMain.lean:1897}\par \hangindent=1em\hangafter=1 \lean{MirrorOdd.unframe_frame}\newline\dirn{MIRRORETV/lean/MirrorMain.lean:1900}\par \hangindent=1em\hangafter=1 \lean{MirrorOdd.frame_unframe}\newline\dirn{MIRRORETV/lean/MirrorMain.lean:1907}\par \hangindent=1em\hangafter=1 \lean{MirrorOdd.frame_lineRefl}\newline\dirn{MIRRORETV/lean/MirrorMain.lean:1914}\par \hangindent=1em\hangafter=1 \lean{MirrorOdd.orient_frame}\newline\dirn{MIRRORETV/lean/MirrorMain.lean:1920}\par \hangindent=1em\hangafter=1 \lean{MirrorOdd.frame_injective}\newline\dirn{MIRRORETV/lean/MirrorMain.lean:1925}\par \hangindent=1em\hangafter=1 \lean{MirrorOdd.unframe_injective}\newline\dirn{MIRRORETV/lean/MirrorMain.lean:1930} & \\
Thm.~\ref{mm:thm:line} & \hangindent=1em\hangafter=1 \lean{MirrorOdd.transport}\newline\dirn{MIRRORETV/lean/MirrorMain.lean:1937}\par \hangindent=1em\hangafter=1 \lean{MirrorOdd.mirror_odd_general_line}\newline\dirn{MIRRORETV/lean/MirrorMain.lean:2641}\par \hangindent=1em\hangafter=1 \lean{MirrorOdd.mirror_even_general_line}\newline\dirn{MIRRORETV/lean/MirrorMain.lean:2650} & part (a) is transport; parts (b), (c) take the parametric hypothesis \texttt{hM} (as in Theorem~\ref{mm:thm:odd} / Theorem~\ref{mm:thm:even})\\
Lem.~\ref{mm:lem:cup} & \hangindent=1em\hangafter=1 \lean{MirrorOdd.hasCap3_of}\newline\dirn{MIRRORETV/lean/MirrorMain.lean:1991}\par \hangindent=1em\hangafter=1 \lean{MirrorOdd.capacity_cup}\newline\dirn{MIRRORETV/lean/MirrorMain.lean:2002} & \\
Thm.~\ref{mm:thm:small} & \hangindent=1em\hangafter=1 \lean{MirrorOdd.mirror_four}\newline\dirn{MIRRORETV/lean/MirrorMain.lean:2019}\par \hangindent=1em\hangafter=1 \lean{MirrorOdd.mirror_five}\newline\dirn{MIRRORETV/lean/MirrorMain.lean:2040}\par \hangindent=1em\hangafter=1 \lean{MirrorOdd.mirror_six}\newline\dirn{MIRRORETV/lean/MirrorMain.lean:2129}\par \hangindent=1em\hangafter=1 \lean{MirrorOdd.mirror_four_line}\newline\dirn{MIRRORETV/lean/MirrorMain.lean:2159}\par \hangindent=1em\hangafter=1 \lean{MirrorOdd.mirror_five_line}\newline\dirn{MIRRORETV/lean/MirrorMain.lean:2165}\par \hangindent=1em\hangafter=1 \lean{MirrorOdd.mirror_six_line}\newline\dirn{MIRRORETV/lean/MirrorMain.lean:2171} & unconditional (no hypotheses beyond the stated ones)\\
\addlinespace
\multicolumn{3}{@{}l}{\normalsize\textbf{Section~\ref{mb:sec}}}\\[2pt]
Def.~\ref{mb:def:config} & \hangindent=1em\hangafter=1 \lean{Config}\newline\dirn{ErdosSzekeres/Vendor/ETV/Config/Defs.lean:20}\par \hangindent=1em\hangafter=1 \lean{Config.Cap3}\newline\dirn{ErdosSzekeres/Vendor/ETV/Config/Defs.lean:30}\par \hangindent=1em\hangafter=1 \lean{Config.Cap}\newline\dirn{ErdosSzekeres/Vendor/ETV/Config/Defs.lean:38}\par \hangindent=1em\hangafter=1 \lean{Config.Cup}\newline\dirn{ErdosSzekeres/Vendor/ETV/Config/Defs.lean:41}\par \hangindent=1em\hangafter=1 \lean{Config.Gon}\newline\dirn{ErdosSzekeres/Vendor/ETV/Config/Defs.lean:45}\par \hangindent=1em\hangafter=1 \lean{Config.NGon}\newline\dirn{ErdosSzekeres/Vendor/ETV/Config/Defs.lean:57}\par \hangindent=1em\hangafter=1 \lean{Config.HasNCap}\newline\dirn{ErdosSzekeres/Vendor/ETV/Config/Defs.lean:60}\par \hangindent=1em\hangafter=1 \lean{Config.HasNCup}\newline\dirn{ErdosSzekeres/Vendor/ETV/Config/Defs.lean:63}\par \hangindent=1em\hangafter=1 \lean{Config.HasNGon}\newline\dirn{ErdosSzekeres/Vendor/ETV/Config/Defs.lean:66} & definitions of Baek's abstract configurations (vendored formalization)\\
Rem.~\ref{mb:rem:config} & -- & remark (no statement): Lean encoding of Definition~\ref{mb:def:config}\\
Lem.~\ref{mb:lem:config-mono} & \hangindent=1em\hangafter=1 \lean{hasNCap_supset}\newline\dirn{ErdosSzekeres/Vendor/ETV/Config/Lemmas.lean:293}\par \hangindent=1em\hangafter=1 \lean{hasNGon_supset}\newline\dirn{ErdosSzekeres/Vendor/ETV/Config/Lemmas.lean:305}\par \hangindent=1em\hangafter=1 \lean{ncup_is_ngon}\newline\dirn{ErdosSzekeres/Vendor/ETV/Config/Lemmas.lean:275} & vendored lemmas; proved here\\
Thm.~\ref{mb:thm:baek} & \hangindent=1em\hangafter=1 \lean{main}\newline\dirn{ErdosSzekeres/Vendor/ETV/Main/Main.lean:87} & cited (J. Baek, Eur. J. Combin. 124 (2025) 104085; \texttt{arXiv:2206.\allowbreak{}04260}); vendored Lean formalization; not reproved here\\
Rem.~\ref{mb:rem:baek} & -- & remark (no statement): source, status, index shift main $n\leftrightarrow$ Baek at $N=n+3$\\
Thm.~\ref{mb:thm:capcup} & \hangindent=1em\hangafter=1 \lean{Config.cap_cup}\newline\dirn{ErdosSzekeres/Vendor/ETV/Main/CapCup.lean:46} & cited (Erdos-Szekeres 1935 cups-caps theorem); vendored Lean formalization; not reproved here\\
Rem.~\ref{mb:rem:capcup} & -- & remark (no statement): source of the cups-caps theorem\\
Lem.~\ref{mb:lem:baek8} & \hangindent=1em\hangafter=1 \lean{ErdosSzekeres.Omission.Baek.four_cap_or_eight_gon}\newline\dirn{ErdosSzekeres/Omission/Baek.lean:12}\par \hangindent=1em\hangafter=1 \lean{ErdosSzekeres.Omission.Baek.capacity22}\newline\dirn{ErdosSzekeres/Omission/Baek.lean:16} & instance $n=5$ of the cited Theorem~\ref{mb:thm:baek}\\
Def.~\ref{mb:def:color} & \hangindent=1em\hangafter=1 \lean{ErdosSzekeres.Omission.Configuration.cup}\newline\dirn{ErdosSzekeres/Omission/Realization.lean:33}\par \hangindent=1em\hangafter=1 \lean{ErdosSzekeres.Omission.ActualRanks.colorConfig}\newline\dirn{ErdosSzekeres/Omission/ActualRanks.lean:28}\par \hangindent=1em\hangafter=1 \lean{ErdosSzekeres.Omission.ListGeometry.pointAt}\newline\dirn{ErdosSzekeres/Omission/ListGeometry.lean:13} & definitions (bridge): cup colouring, \texttt{colorConfig} \texttt{c.\allowbreak{}cup} true, \texttt{pointAt}\\
Lem.~\ref{mb:lem:signs} & \hangindent=1em\hangafter=1 \lean{ErdosSzekeres.Omission.Configuration.cup_true_iff}\newline\dirn{ErdosSzekeres/Omission/Realization.lean:39}\par \hangindent=1em\hangafter=1 \lean{ErdosSzekeres.Omission.ListGeometry.cup_sign}\newline\dirn{ErdosSzekeres/Omission/ListGeometry.lean:19}\par \hangindent=1em\hangafter=1 \lean{ErdosSzekeres.Omission.ListGeometry.cap_sign}\newline\dirn{ErdosSzekeres/Omission/ListGeometry.lean:24} & bridge lemma; \texttt{cup\_\allowbreak{}true\_\allowbreak{}iff} takes Fin indices, \texttt{cup\_\allowbreak{}sign} and \texttt{cap\_\allowbreak{}sign} natural-number indices with bounds; \texttt{cup\_\allowbreak{}false\_\allowbreak{}iff} is not on the proof paths and not covered\\
Lem.~\ref{mb:lem:four} & \hangindent=1em\hangafter=1 \lean{ErdosSzekeres.Omission.PlanarChains.four_point_identity_left}\newline\dirn{ErdosSzekeres/Omission/PlanarChains.lean:15}\par \hangindent=1em\hangafter=1 \lean{ErdosSzekeres.Omission.PlanarChains.four_point_identity_right}\newline\dirn{ErdosSzekeres/Omission/PlanarChains.lean:19}\par \hangindent=1em\hangafter=1 \lean{ErdosSzekeres.Omission.PlanarChains.positive_four}\newline\dirn{ErdosSzekeres/Omission/PlanarChains.lean:23}\par \hangindent=1em\hangafter=1 \lean{ErdosSzekeres.Omission.PlanarChains.negative_four}\newline\dirn{ErdosSzekeres/Omission/PlanarChains.lean:35} & bridge lemma; identities checked symbolically\\
Lem.~\ref{mb:lem:chain} & \hangindent=1em\hangafter=1 \lean{ErdosSzekeres.Omission.PlanarChains.chain3_all}\newline\dirn{ErdosSzekeres/Omission/PlanarChains.lean:58} & bridge lemma\\
Lem.~\ref{mb:lem:lists} & \hangindent=1em\hangafter=1 \lean{ErdosSzekeres.Omission.ListGeometry.cup_list_all}\newline\dirn{ErdosSzekeres/Omission/ListGeometry.lean:39}\par \hangindent=1em\hangafter=1 \lean{ErdosSzekeres.Omission.ListGeometry.cap_list_all}\newline\dirn{ErdosSzekeres/Omission/ListGeometry.lean:53}\par \hangindent=1em\hangafter=1 \lean{ErdosSzekeres.Omission.ListGeometry.hasCap_of_list}\newline\dirn{ErdosSzekeres/Omission/ListGeometry.lean:85} & bridge lemma\\
Lem.~\ref{mb:lem:boundary} & \hangindent=1em\hangafter=1 \lean{ErdosSzekeres.Omission.BoundaryGeometry.boundary_polygon}\newline\dirn{ErdosSzekeres/Omission/BoundaryGeometry.lean:403}\par \hangindent=1em\hangafter=1 \lean{ErdosSzekeres.Omission.BoundaryGeometry.Boundary.card}\newline\dirn{ErdosSzekeres/Omission/BoundaryGeometry.lean:359}\par \hangindent=1em\hangafter=1 \lean{ErdosSzekeres.Omission.BoundaryGeometry.boundary_convex_position}\newline\dirn{ErdosSzekeres/Omission/BoundaryGeometry.lean:317} & bridge lemma; stated after the change of variables $v=$ reflect(upperReflected), see Remark~\ref{mb:rem:boundary}\\
Rem.~\ref{mb:rem:boundary} & -- & remark (no statement): Lean encoding of Lemma~\ref{mb:lem:boundary}\\
Prop.~\ref{mb:prop:gon} & \hangindent=1em\hangafter=1 \lean{MirrorOdd.hasNGon_heptagon}\newline\dirn{MIRRORETV/lean/MirrorMain.lean:1512}\par \hangindent=1em\hangafter=1 \lean{ErdosSzekeres.Omission.BaekGeometry.hasNGon_octagon}\newline\dirn{ErdosSzekeres/Omission/BaekGeometry.lean:94}\par \hangindent=1em\hangafter=1 \lean{C12Check.hasNGon_nonagon}\newline\dirn{MIRRORETV/lean/MirrorNine_append.lean:99} & bridge; three Lean declarations ($m=7$, 8, 9) whose proofs are identical up to the value of m\\
Prop.~\ref{mb:prop:cap16} & \hangindent=1em\hangafter=1 \lean{MirrorOdd.capacity16}\newline\dirn{MIRRORETV/lean/MirrorMain.lean:1529} & uses the cited Theorem~\ref{mb:thm:baek} $(n=4)$; binomial value $C(6,2)=15$ by kernel decide\\
Lem.~\ref{mb:lem:gc22} & \hangindent=1em\hangafter=1 \lean{ErdosSzekeres.Omission.BaekGeometry.geometric_capacity22}\newline\dirn{ErdosSzekeres/Omission/BaekGeometry.lean:111} & uses the cited Theorem~\ref{mb:thm:baek} $(n=5)$\\
Prop.~\ref{mb:prop:cap22} & \hangindent=1em\hangafter=1 \lean{MirrorOdd.capacity22g}\newline\dirn{MIRRORETV/lean/MirrorMain.lean:1812} & uses the cited Theorem~\ref{mb:thm:baek} $(n=5)$ through Lemma~\ref{mb:lem:gc22}\\
Def.~\ref{mb:def:mir} & \hangindent=1em\hangafter=1 \lean{MirrorOdd.mir}\newline\dirn{MIRRORETV/lean/MirrorMain.lean:16}\par \hangindent=1em\hangafter=1 \lean{MirrorOdd.Rset}\newline\dirn{MIRRORETV/lean/MirrorMain.lean:1598}\par \hangindent=1em\hangafter=1 \lean{MirrorOdd.Aset}\newline\dirn{MIRRORETV/lean/MirrorMain.lean:1599} & definitions\\
Lem.~\ref{mb:lem:split} & \hangindent=1em\hangafter=1 \lean{MirrorOdd.card_split}\newline\dirn{MIRRORETV/lean/MirrorMain.lean:1601}\par \hangindent=1em\hangafter=1 \lean{MirrorOdd.card_A_le}\newline\dirn{MIRRORETV/lean/MirrorMain.lean:1636}\par \hangindent=1em\hangafter=1 \lean{MirrorOdd.injOn_x}\newline\dirn{MIRRORETV/lean/MirrorMain.lean:1647} & proved here (also used inside Theorem~\ref{mm:thm:odd} and Theorem~\ref{mm:thm:even})\\
Lem.~\ref{mb:lem:sort} & \hangindent=1em\hangafter=1 \lean{MirrorOdd.cfg}\newline\dirn{MIRRORETV/lean/MirrorMain.lean:1584}\par \hangindent=1em\hangafter=1 \lean{MirrorOdd.sortPt_image}\newline\dirn{MIRRORETV/lean/MirrorMain.lean:1576} & the sorting construction of Definition~\ref{mm:def:config}\\
Lem.~\ref{mb:lem:interior} & written step & written step (local claim interior in the Lean proof of \texttt{mirror\_\allowbreak{}seven\_\allowbreak{}of})\\
Lem.~\ref{mb:lem:polygons} & \hangindent=1em\hangafter=1 \lean{MirrorOdd.cap_mirror_heptagon}\newline\dirn{MIRRORETV/lean/MirrorMain.lean:212}\par \hangindent=1em\hangafter=1 \lean{MirrorOdd.cap_mirror_octagon}\newline\dirn{MIRRORETV/lean/MirrorMain.lean:392}\par \hangindent=1em\hangafter=1 \lean{MirrorOdd.axis_left_heptagon}\newline\dirn{MIRRORETV/lean/MirrorMain.lean:639}\par \hangindent=1em\hangafter=1 \lean{MirrorOdd.axis_right_heptagon}\newline\dirn{MIRRORETV/lean/MirrorMain.lean:816} & Lean proofs are generated code (linarith with hints found offline); written proof derives the statements from Theorem~\ref{mm:thm:mirror-polygon} and Theorem~\ref{mm:thm:axis-polygon} (different proof)\\
Rem.~\ref{mb:rem:generated} & -- & remark (no statement): how Lean proves Lemma~\ref{mb:lem:polygons}\\
Lem.~\ref{mb:lem:seven-of} & \hangindent=1em\hangafter=1 \lean{MirrorOdd.mirror_seven_of}\newline\dirn{MIRRORETV/lean/MirrorMain.lean:1676} & takes hypotheses: Hcap, HL, HR (the statements of Lemma~\ref{mb:lem:polygons} (a), (c), (d); parametric premises, discharged in Theorem~\ref{mb:thm:seven})\\
Def.~\ref{mb:def:frame} & \hangindent=1em\hangafter=1 \lean{MirrorOdd.lineRefl}\newline\dirn{MIRRORETV/lean/MirrorMain.lean:1879}\par \hangindent=1em\hangafter=1 \lean{MirrorOdd.frame}\newline\dirn{MIRRORETV/lean/MirrorMain.lean:1883}\par \hangindent=1em\hangafter=1 \lean{MirrorOdd.unframe}\newline\dirn{MIRRORETV/lean/MirrorMain.lean:1892} & definitions\\
Lem.~\ref{mb:lem:frame} & \hangindent=1em\hangafter=1 \lean{MirrorOdd.frame_lineRefl}\newline\dirn{MIRRORETV/lean/MirrorMain.lean:1914}\par \hangindent=1em\hangafter=1 \lean{MirrorOdd.orient_frame}\newline\dirn{MIRRORETV/lean/MirrorMain.lean:1920}\par \hangindent=1em\hangafter=1 \lean{MirrorOdd.unframe_frame}\newline\dirn{MIRRORETV/lean/MirrorMain.lean:1900}\par \hangindent=1em\hangafter=1 \lean{MirrorOdd.frame_unframe}\newline\dirn{MIRRORETV/lean/MirrorMain.lean:1907}\par \hangindent=1em\hangafter=1 \lean{MirrorOdd.frame_injective}\newline\dirn{MIRRORETV/lean/MirrorMain.lean:1925}\par \hangindent=1em\hangafter=1 \lean{MirrorOdd.unframe_injective}\newline\dirn{MIRRORETV/lean/MirrorMain.lean:1930} & polynomial identities checked symbolically\\
Lem.~\ref{mb:lem:affine} & \hangindent=1em\hangafter=1 \lean{MirrorOdd.hasConvexNGon_affineMap}\newline\dirn{MIRRORETV/lean/MirrorMain.lean:1859} & proved here\\
Lem.~\ref{mb:lem:transport} & \hangindent=1em\hangafter=1 \lean{MirrorOdd.transport}\newline\dirn{MIRRORETV/lean/MirrorMain.lean:1937} & proved here\\
Thm.~\ref{mb:thm:seven} & \hangindent=1em\hangafter=1 \lean{MirrorOdd.mirror_seven}\newline\dirn{MIRRORETV/lean/MirrorMain.lean:1803}\par \hangindent=1em\hangafter=1 \lean{MirrorOdd.mirror_seven_line}\newline\dirn{MIRRORETV/lean/MirrorMain.lean:1969} & rests on the cited Theorem~\ref{mb:thm:baek} and on Theorem~\ref{mm:thm:mirror-polygon}, Theorem~\ref{mm:thm:axis-polygon}\\
Rem.~\ref{mb:rem:seven} & -- & remark (no statement): \texttt{mirror\_\allowbreak{}seven'} as a second derivation\\
Thm.~\ref{mb:thm:eight} & \hangindent=1em\hangafter=1 \lean{MirrorOdd.mirror_eight}\newline\dirn{MIRRORETV/lean/MirrorMain.lean:1820}\par \hangindent=1em\hangafter=1 \lean{MirrorOdd.mirror_eight_line}\newline\dirn{MIRRORETV/lean/MirrorMain.lean:1977} & rests on the cited Theorem~\ref{mb:thm:baek} and on Theorem~\ref{mm:thm:mirror-polygon}\\
Rem.~\ref{mb:rem:eight} & -- & remark (no statement): \texttt{mirror\_\allowbreak{}eight'} as a second derivation\\
Thm.~\ref{mb:thm:nine} & \hangindent=1em\hangafter=1 \lean{C12Check.config_capacity113}\newline\dirn{MIRRORETV/lean/MirrorNine_append.lean:24}\par \hangindent=1em\hangafter=1 \lean{C12Check.mirror_nine}\newline\dirn{MIRRORETV/lean/MirrorNine_append.lean:131}\par \hangindent=1em\hangafter=1 \lean{C12Check.mirror_nine_line}\newline\dirn{MIRRORETV/lean/MirrorNine_append.lean:136} & with \texttt{config\_\allowbreak{}capacity113} as part (a); (a) uses the cited Theorem~\ref{mb:thm:baek} $(n=6)$ and Theorem~\ref{mb:thm:capcup} $(a=3,b=6),C(8,2)=28$ and $C(9,3)=84$ by kernel decide; (b),(c) via Theorem~\ref{mm:thm:odd} and Theorem~\ref{mm:thm:line} $(j=4,M=113)$\\
Prop.~\ref{mb:prop:cap113} & \hangindent=1em\hangafter=1 \lean{C12Check.capacity113}\newline\dirn{MIRRORETV/lean/MirrorNine_append.lean:116} & uses Theorem~\ref{mb:thm:nine}(a) (\texttt{config\_\allowbreak{}capacity113}); stated and proved between the proofs of Theorem~\ref{mb:thm:nine}(a) and (b),(c)\\
Def.~\ref{mb:def:reflline} & \hangindent=1em\hangafter=1 \lean{ErdosSzekeres.SymReduction.reflLine}\newline\dirn{SYMLEAN/mirror/ReductionThreeClasses.lean:449} & definition (same text as \texttt{SYMLEAN/\allowbreak{}SymmetricReduction.\allowbreak{}lean:423})\\
Thm.~\ref{mb:thm:refl32} & \hangindent=1em\hangafter=1 \lean{ErdosSzekeres.SymReduction.reflLine_eq_lineRefl}\newline\dirn{SYMLEAN/mirror/ReductionThreeClasses.lean:4253}\par \hangindent=1em\hangafter=1 \lean{ErdosSzekeres.SymReduction.card_le_32_of_reflLine}\newline\dirn{SYMLEAN/mirror/ReductionThreeClasses.lean:4264} & with \texttt{reflLine\_\allowbreak{}eq\_\allowbreak{}lineRefl}\\
Thm.~\ref{mb:thm:classes} & \hangindent=1em\hangafter=1 \lean{ErdosSzekeres.SymReduction.mirrorOneAxis33_proved}\newline\dirn{SYMLEAN/mirror/ReductionThreeClasses.lean:4274}\par \hangindent=1em\hangafter=1 \lean{ErdosSzekeres.SymReduction.mirrorNoAxis34_proved}\newline\dirn{SYMLEAN/mirror/ReductionThreeClasses.lean:4282} & the axis hypotheses of both classes are unused\\
Thm.~\ref{mb:thm:main3} & \hangindent=1em\hangafter=1 \lean{ErdosSzekeres.SymReduction.heptagon_of_nontrivial_isometry_3}\newline\dirn{SYMLEAN/mirror/ReductionThreeClasses.lean:4292}\par \hangindent=1em\hangafter=1 \lean{ErdosSzekeres.SymReduction.heptagon_of_nontrivial_affineEquiv_3}\newline\dirn{SYMLEAN/mirror/ReductionThreeClasses.lean:4303} & takes hypotheses: \texttt{HalfTurn34}, Threefold33, Fivefold35 (not proved in Lean; computer-assisted claims supplied by the records \texttt{C2FINAL/\allowbreak{}C2C34}, \texttt{C3FINAL/\allowbreak{}C3LEAN}, SYM5, resting also on written encoder-soundness and perturbation proofs and in part on verdicts recorded at solve time; see Remark~\ref{mb:rem:main3}); conditional\\
Rem.~\ref{mb:rem:main3} & -- & remark (no statement): status of the three hypotheses of Theorem~\ref{mb:thm:main3}, with the records and rigor labels of their evidence; the case-split bookkeeping is Lean-checked according to the record (Theorem~\ref{cb:thm:c2}) or Lean-checked (Theorem~\ref{cb:thm:c3}), in both cases by kernel evaluation, so computer-assisted\\
Rem.~\ref{mb:rem:axis} & -- & remark (no statement): \texttt{MirrorExclusionLineCopy} follows from Theorem~\ref{mb:thm:refl32}(b)\\
\addlinespace
\multicolumn{3}{@{}l}{\normalsize\textbf{Section~\ref{cb:sec}}}\\[2pt]
Lem.~\ref{cb:lem:bits} & written step & written step; binary digits of an or-fold\\
Def.~\ref{cb:def:m} & \hangindent=1em\hangafter=1 \lean{C3Cover.m}\newline\dirn{C3LEAN/C3Cover.lean:34} & definition\\
Def.~\ref{cb:def:tok} & \hangindent=1em\hangafter=1 \lean{C3Cover.Tok}\newline\dirn{C3LEAN/C3Cover.lean:38} & definition\\
Def.~\ref{cb:def:layer} & \hangindent=1em\hangafter=1 \lean{C3Cover.Layer}\newline\dirn{C3LEAN/C3Cover.lean:41} & definition\\
Def.~\ref{cb:def:node} & \hangindent=1em\hangafter=1 \lean{C3Cover.Node}\newline\dirn{C3LEAN/C3Cover.lean:45} & definition\\
Def.~\ref{cb:def:used} & \hangindent=1em\hangafter=1 \lean{C3Cover.Node.used}\newline\dirn{C3LEAN/C3Cover.lean:51} & definition\\
Rem.~\ref{cb:rem:reading} & -- & commentary; intended meaning, not used in proofs\\
Def.~\ref{cb:def:pairoptions} & \hangindent=1em\hangafter=1 \lean{C3Cover.pairOptions}\newline\dirn{C3LEAN/C3Cover.lean:54} & definition\\
Def.~\ref{cb:def:layeroptions} & \hangindent=1em\hangafter=1 \lean{C3Cover.layerOptions}\newline\dirn{C3LEAN/C3Cover.lean:61} & definition; rem is a local abbreviation inside \texttt{layerOptions}\\
Def.~\ref{cb:def:instok} & \hangindent=1em\hangafter=1 \lean{C3Cover.insTok}\newline\dirn{C3LEAN/C3Cover.lean:66} & definition\\
Def.~\ref{cb:def:sorttok} & \hangindent=1em\hangafter=1 \lean{C3Cover.sortTok}\newline\dirn{C3LEAN/C3Cover.lean:71} & definition\\
Def.~\ref{cb:def:topcases} & \hangindent=1em\hangafter=1 \lean{C3Cover.topCases}\newline\dirn{C3LEAN/C3Cover.lean:77} & definition\\
Def.~\ref{cb:def:okey} & \hangindent=1em\hangafter=1 \lean{C3Cover.okey}\newline\dirn{C3LEAN/C3Cover.lean:82} & definition\\
Def.~\ref{cb:def:lexle} & \hangindent=1em\hangafter=1 \lean{C3Cover.lexLe}\newline\dirn{C3LEAN/C3Cover.lean:85} & definition\\
Def.~\ref{cb:def:mirror} & \hangindent=1em\hangafter=1 \lean{C3Cover.mirror}\newline\dirn{C3LEAN/C3Cover.lean:91} & definition\\
Def.~\ref{cb:def:phi} & \hangindent=1em\hangafter=1 \lean{C3Cover.phi}\newline\dirn{C3LEAN/C3Cover.lean:95} & definition\\
Def.~\ref{cb:def:refl} & \hangindent=1em\hangafter=1 \lean{C3Cover.refl}\newline\dirn{C3LEAN/C3Cover.lean:101} & definition\\
Def.~\ref{cb:def:kept} & \hangindent=1em\hangafter=1 \lean{C3Cover.keptOptions}\newline\dirn{C3LEAN/C3Cover.lean:109} & definition\\
Def.~\ref{cb:def:roots} & \hangindent=1em\hangafter=1 \lean{C3Cover.roots}\newline\dirn{C3LEAN/C3Cover.lean:113} & definition\\
Def.~\ref{cb:def:free} & \hangindent=1em\hangafter=1 \lean{C3Cover.Node.free}\newline\dirn{C3LEAN/C3Cover.lean:116} & definition\\
Def.~\ref{cb:def:bitsof} & \hangindent=1em\hangafter=1 \lean{C3Cover.bitsOf}\newline\dirn{C3LEAN/C3Cover.lean:120} & definition\\
Def.~\ref{cb:def:addbits} & \hangindent=1em\hangafter=1 \lean{C3Cover.addBits}\newline\dirn{C3LEAN/C3Cover.lean:123} & definition\\
Def.~\ref{cb:def:sigmachildren} & \hangindent=1em\hangafter=1 \lean{C3Cover.Node.sigmaChildren}\newline\dirn{C3LEAN/C3Cover.lean:127} & definition; bits is a local abbreviation\\
Def.~\ref{cb:def:conflicts} & \hangindent=1em\hangafter=1 \lean{C3Cover.conflicts}\newline\dirn{C3LEAN/C3Cover.lean:132} & definition\\
Def.~\ref{cb:def:layerchildren} & \hangindent=1em\hangafter=1 \lean{C3Cover.Node.layerChildren}\newline\dirn{C3LEAN/C3Cover.lean:137} & definition\\
Rem.~\ref{cb:rem:children} & -- & commentary; reading of the children, not used in proofs\\
Def.~\ref{cb:def:covered} & \hangindent=1em\hangafter=1 \lean{C3Cover.Covered}\newline\dirn{C3LEAN/C3Cover.lean:146} & inductive predicate; least predicate closed under the three constructors\\
Lem.~\ref{cb:lem:mono} & \hangindent=1em\hangafter=1 \lean{C3Cover.Covered.mono}\newline\dirn{C3LEAN/C3Cover.lean:153} & Lean-checked (dependency of \texttt{all\_\allowbreak{}roots\_\allowbreak{}covered}); written proof by rule induction\\
Def.~\ref{cb:def:decodetok} & \hangindent=1em\hangafter=1 \lean{C3Cover.decodeTok}\newline\dirn{C3LEAN/C3Cover.lean:165} & definition\\
Def.~\ref{cb:def:decodelayer} & \hangindent=1em\hangafter=1 \lean{C3Cover.decodeLayer}\newline\dirn{C3LEAN/C3Cover.lean:167} & definition\\
Def.~\ref{cb:def:decodelayers} & \hangindent=1em\hangafter=1 \lean{C3Cover.decodeLayers}\newline\dirn{C3LEAN/C3Cover.lean:172} & definition\\
Def.~\ref{cb:def:decodecube} & \hangindent=1em\hangafter=1 \lean{C3Cover.decodeCube}\newline\dirn{C3LEAN/C3Cover.lean:180} & definition\\
Def.~\ref{cb:def:decodekey} & \hangindent=1em\hangafter=1 \lean{C3Cover.decodeKey}\newline\dirn{C3LEAN/C3Cover.lean:191} & definition\\
Rem.~\ref{cb:rem:keys} & -- & commentary; uses the control \texttt{rev\_\allowbreak{}parent\_\allowbreak{}decodes} (\texttt{C3LEAN/\allowbreak{}review/\allowbreak{}RevCtrl\_\allowbreak{}pass.\allowbreak{}lean:315}); \texttt{decodeKey} not injective\\
Def.~\ref{cb:def:tokcode} & \hangindent=1em\hangafter=1 \lean{C3Cover.tokCode}\newline\dirn{C3LEAN/C3Cover.lean:223} & definition\\
Def.~\ref{cb:def:layercode} & \hangindent=1em\hangafter=1 \lean{C3Cover.layerCode}\newline\dirn{C3LEAN/C3Cover.lean:226} & definition\\
Def.~\ref{cb:def:rootcode} & \hangindent=1em\hangafter=1 \lean{C3Cover.rootCode}\newline\dirn{C3LEAN/C3Cover.lean:232} & definition\\
Def.~\ref{cb:def:decoderoot} & \hangindent=1em\hangafter=1 \lean{C3Cover.decodeRoot}\newline\dirn{C3LEAN/C3Cover.lean:238} & definition\\
Def.~\ref{cb:def:beqlist} & \hangindent=1em\hangafter=1 \lean{C3Cover.beqList}\newline\dirn{C3LEAN/C3Cover.lean:194} & definition\\
Def.~\ref{cb:def:beqtok} & \hangindent=1em\hangafter=1 \lean{C3Cover.beqTok}\newline\dirn{C3LEAN/C3Cover.lean:199} & definition\\
Def.~\ref{cb:def:beqnode} & \hangindent=1em\hangafter=1 \lean{C3Cover.beqNode}\newline\dirn{C3LEAN/C3Cover.lean:202} & definition\\
Def.~\ref{cb:def:beqlayer} & \hangindent=1em\hangafter=1 \lean{C3Cover.beqLayer}\newline\dirn{C3LEAN/C3Cover.lean:335} & definition\\
Lem.~\ref{cb:lem:beqlist} & \hangindent=1em\hangafter=1 \lean{C3Cover.beqList_eq}\newline\dirn{C3LEAN/C3Cover.lean:204} & Lean-checked (dependency of printed theorems); written proof\\
Lem.~\ref{cb:lem:beqtok} & \hangindent=1em\hangafter=1 \lean{C3Cover.beqTok_eq}\newline\dirn{C3LEAN/C3Cover.lean:213} & Lean-checked (dependency of printed theorems); written proof\\
Lem.~\ref{cb:lem:beqnode} & \hangindent=1em\hangafter=1 \lean{C3Cover.beqNode_eq}\newline\dirn{C3LEAN/C3Cover.lean:218} & Lean-checked (dependency of printed theorems); written proof\\
Lem.~\ref{cb:lem:beq-refl} & written step & written step; reflexivity of the Boolean equalities\\
Def.~\ref{cb:def:cert} & \hangindent=1em\hangafter=1 \lean{C3Cover.Cert}\newline\dirn{C3LEAN/C3Cover.lean:245} & definition\\
Def.~\ref{cb:def:unpack} & \hangindent=1em\hangafter=1 \lean{C3Cover.unpack}\newline\dirn{C3LEAN/C3Cover.lean:251} & definition\\
Def.~\ref{cb:def:keys} & \hangindent=1em\hangafter=1 \lean{C3Cover.Cert.keys}\newline\dirn{C3LEAN/C3Cover.lean:259} & definition\\
Def.~\ref{cb:def:run} & \hangindent=1em\hangafter=1 \lean{C3Cover.run}\newline\dirn{C3LEAN/C3Cover.lean:262} & definition\\
Def.~\ref{cb:def:checkroot} & \hangindent=1em\hangafter=1 \lean{C3Cover.checkRoot}\newline\dirn{C3LEAN/C3Cover.lean:277} & definition\\
Rem.~\ref{cb:rem:cert} & -- & commentary; how a certificate is read; final test and fuel not needed for soundness\\
Lem.~\ref{cb:lem:run} & \hangindent=1em\hangafter=1 \lean{C3Cover.run_sound}\newline\dirn{C3LEAN/C3Cover.lean:279}\par \hangindent=1em\hangafter=1 \lean{C3Cover.checkRoot_sound}\newline\dirn{C3LEAN/C3Cover.lean:323} & Lean-checked, axioms propext and \texttt{Quot.\allowbreak{}sound}; written proof by induction on fuel\\
Lem.~\ref{cb:lem:append} & \hangindent=1em\hangafter=1 \lean{C3Cover.all_append_of}\newline\dirn{C3LEAN/C3Cover.lean:327} & Lean-checked (dependency of \texttt{all\_\allowbreak{}checked}); written proof\\
Def.~\ref{cb:def:tsv} & \hangindent=1em\hangafter=1 \lean{C3Cover.tsvRootCodes}\newline\dirn{C3LEAN/C3Cover.lean:413} & generated data (910 root codes); identification with the TSV file checked in Python only\\
Def.~\ref{cb:def:certsN} & \hangindent=1em\hangafter=1 \lean{C3Cover.certs0}\newline\dirn{C3LEAN/C3Cover.lean:423} & generated data; family certs0--certs126 at lines 423-1711 (certs126 at 1668)\\
Def.~\ref{cb:def:certs} & \hangindent=1em\hangafter=1 \lean{C3Cover.certs}\newline\dirn{C3LEAN/C3Cover.lean:1714} & generated data; right-nested append\\
Def.~\ref{cb:def:bounds} & \hangindent=1em\hangafter=1 \lean{C3Cover.bounds}\newline\dirn{C3LEAN/C3Cover.lean:1717} & generated data; unused by any theorem\\
Def.~\ref{cb:def:certifiedkeys} & \hangindent=1em\hangafter=1 \lean{C3Cover.certifiedKeys}\newline\dirn{C3LEAN/C3Cover.lean:1737} & generated data; also defines the leaf predicate S-star of the main theorem\\
Def.~\ref{cb:def:isoption} & \hangindent=1em\hangafter=1 \lean{C3Cover.isOption}\newline\dirn{C3LEAN/C3Cover.lean:340} & definition\\
Def.~\ref{cb:def:distinctcodes} & \hangindent=1em\hangafter=1 \lean{C3Cover.distinctCodes}\newline\dirn{C3LEAN/C3Cover.lean:346} & definition\\
Def.~\ref{cb:def:codeset} & \hangindent=1em\hangafter=1 \lean{C3Cover.codeSet}\newline\dirn{C3LEAN/C3Cover.lean:351} & definition\\
Def.~\ref{cb:def:reflectionfacts} & \hangindent=1em\hangafter=1 \lean{C3Cover.reflectionFacts}\newline\dirn{C3LEAN/C3Cover.lean:356} & definition\\
Lem.~\ref{cb:lem:codes} & written step & written step; characterization of \texttt{distinctCodes} and \texttt{codeSet}\\
Fact~\ref{cb:fact:batches} & \hangindent=1em\hangafter=1 \lean{C3Cover.batch0}\newline\dirn{C3LEAN/C3Cover.lean:1740} & computer-assisted (decide +kernel); family batch0--batch126 at lines 1740+4N; no hand proof\\
Lem.~\ref{cb:lem:all-checked} & \hangindent=1em\hangafter=1 \lean{C3Cover.all_checked}\newline\dirn{C3LEAN/C3Cover.lean:2247} & Lean-checked, axioms propext; computer-assisted (joins the 127 computer-assisted batch facts by \texttt{all\_\allowbreak{}append\_\allowbreak{}of})\\
Fact~\ref{cb:fact:refl-h3} & \hangindent=1em\hangafter=1 \lean{C3Cover.reflection_facts_H3}\newline\dirn{C3LEAN/C3Cover.lean:367} & Lean-checked, axioms propext; computer-assisted in Lean (decide +kernel); written proof in Proposition~\ref{cb:prop:refl}\\
Fact~\ref{cb:fact:refl-h6} & \hangindent=1em\hangafter=1 \lean{C3Cover.reflection_facts_H6_1_0}\newline\dirn{C3LEAN/C3Cover.lean:370} & Lean-checked, axioms propext; computer-assisted in Lean (decide +kernel); ten theorems \texttt{reflection\_\allowbreak{}facts\_\allowbreak{}H6\_\allowbreak{}1\_\allowbreak{}0}--\texttt{H6\_\allowbreak{}5\_\allowbreak{}1} at lines 370+3(2a+s-2); written proof in Proposition~\ref{cb:prop:refl}\\
Fact~\ref{cb:fact:topcases} & \hangindent=1em\hangafter=1 \lean{C3Cover.topCases_eq}\newline\dirn{C3LEAN/C3Cover.lean:401} & Lean-checked, no axioms; computer-assisted in Lean (decide +kernel); written proof in Proposition~\ref{cb:prop:refl}\\
Fact~\ref{cb:fact:strata} & \hangindent=1em\hangafter=1 \lean{C3Cover.strata}\newline\dirn{C3LEAN/C3Cover.lean:406} & Lean-checked, no axioms; computer-assisted in Lean (decide +kernel); written proof in Proposition~\ref{cb:prop:refl}\\
Fact~\ref{cb:fact:roots-match} & \hangindent=1em\hangafter=1 \lean{C3Cover.roots_match}\newline\dirn{C3LEAN/C3Cover.lean:1720} & Lean-checked, no axioms; computer-assisted in Lean (decide +kernel); written proof in Proposition~\ref{cb:prop:roots-match}\\
Fact~\ref{cb:fact:roots-codes} & \hangindent=1em\hangafter=1 \lean{C3Cover.roots_codes_match}\newline\dirn{C3LEAN/C3Cover.lean:1725} & computer-assisted (decide +kernel on generated certificate data), no axioms; no hand proof\\
Fact~\ref{cb:fact:roots-tsv} & \hangindent=1em\hangafter=1 \lean{C3Cover.roots_eq_tsv}\newline\dirn{C3LEAN/C3Cover.lean:1730} & computer-assisted (decide +kernel on generated data), axioms propext; no hand proof\\
Lem.~\ref{cb:lem:sort} & written step & written step\\
Lem.~\ref{cb:lem:options} & written step & written step; structure and count of \texttt{layerOptions}\\
Lem.~\ref{cb:lem:decode} & written step & written step; decoding lemma (\texttt{decodeRoot} after \texttt{rootCode} is the identity on valid depth-2 nodes)\\
Lem.~\ref{cb:lem:lex} & written step & written step; \texttt{lexLe} is a total order\\
Lem.~\ref{cb:lem:orbits} & written step & written step; one kept member per orbit of the reflection\\
Prop.~\ref{cb:prop:refl} & written step & written step; hand proof of reflection\_facts\_*, \texttt{topCases\_\allowbreak{}eq}, strata\\
Prop.~\ref{cb:prop:roots-match} & written step & written step; hand proof of \texttt{roots\_\allowbreak{}match}\\
Cor.~\ref{cb:cor:roots} & written step & written step; computer-assisted (uses \texttt{roots\_\allowbreak{}codes\_\allowbreak{}match} and \texttt{roots\_\allowbreak{}eq\_\allowbreak{}tsv})\\
Thm.~\ref{cb:thm:c3} & \hangindent=1em\hangafter=1 \lean{C3Cover.all_roots_covered}\newline\dirn{C3LEAN/C3Cover.lean:2251} & Lean-checked, axioms propext and \texttt{Quot.\allowbreak{}sound}; no hypotheses; computer-assisted (uses the computer-assisted batch facts and \texttt{roots\_\allowbreak{}codes\_\allowbreak{}match}); leaf set is the certificates' own key list (the run records of the leaves are checked outside Lean)\\
Rem.~\ref{cb:rem:use} & -- & commentary; scope and intended use (geometry, SAT verdicts, run records of the leaves not established)\\
Def.~\ref{cb:def:inner} & \hangindent=1em\hangafter=1 \lean{Central34.inner}\newline\dirn{C2FINAL/Central34Orbits.lean:21} & definition\\
Def.~\ref{cb:def:choose2} & \hangindent=1em\hangafter=1 \lean{Central34.choose2}\newline\dirn{C2FINAL/Central34Orbits.lean:24} & definition\\
Def.~\ref{cb:def:choose3} & \hangindent=1em\hangafter=1 \lean{Central34.choose3}\newline\dirn{C2FINAL/Central34Orbits.lean:29} & definition\\
Def.~\ref{cb:def:subcases} & \hangindent=1em\hangafter=1 \lean{Central34.subcases}\newline\dirn{C2FINAL/Central34Orbits.lean:34} & definition\\
Def.~\ref{cb:def:mask} & \hangindent=1em\hangafter=1 \lean{Central34.mask}\newline\dirn{C2FINAL/Central34Orbits.lean:37} & definition\\
Def.~\ref{cb:def:c2refl} & \hangindent=1em\hangafter=1 \lean{Central34.refl}\newline\dirn{C2FINAL/Central34Orbits.lean:40} & definition\\
Def.~\ref{cb:def:bitset} & \hangindent=1em\hangafter=1 \lean{Central34.bitset}\newline\dirn{C2FINAL/Central34Orbits.lean:43} & definition\\
Def.~\ref{cb:def:uncovered} & \hangindent=1em\hangafter=1 \lean{Central34.uncovered}\newline\dirn{C2FINAL/Central34Orbits.lean:46} & definition\\
Def.~\ref{cb:def:closedP} & \hangindent=1em\hangafter=1 \lean{Central34.closedP1}\newline\dirn{C2FINAL/Central34Orbits.lean:51} & generated data; family \texttt{closedP1}--\texttt{closedP8} at lines 51-264 (287 entries each, 288 for \texttt{closedP6})\\
Def.~\ref{cb:def:closed} & \hangindent=1em\hangafter=1 \lean{Central34.closed}\newline\dirn{C2FINAL/Central34Orbits.lean:267} & definition\\
Rem.~\ref{cb:rem:c2-reading} & -- & commentary; intended meaning, not used in proofs\\
Lem.~\ref{cb:lem:c2-subsets} & written step & written step; subcases enumerate the 2- and 3-subsets\\
Thm.~\ref{cb:thm:c2-refl} & \hangindent=1em\hangafter=1 \lean{Central34.subcases_reflection}\newline\dirn{C2FINAL/Central34Orbits.lean:273} & Lean-checked according to the record (no saved Lean output; \texttt{C2FINAL/\allowbreak{}README.\allowbreak{}md} section 6 records propext only); computer-assisted in Lean (decide +kernel); written proof given\\
Fact~\ref{cb:fact:c2-sub} & \hangindent=1em\hangafter=1 \lean{Central34.closed_are_subcases}\newline\dirn{C2FINAL/Central34Orbits.lean:284} & Lean-checked according to the record (no saved Lean output; \texttt{C2FINAL/\allowbreak{}README.\allowbreak{}md} section 6 records propext only); computer-assisted (decide +kernel on generated lists); no hand proof\\
Thm.~\ref{cb:thm:c2} & \hangindent=1em\hangafter=1 \lean{Central34.all_cases_closed}\newline\dirn{C2FINAL/Central34Orbits.lean:291} & Lean-checked according to the record (no saved Lean output; \texttt{C2FINAL/\allowbreak{}README.\allowbreak{}md} section 6 records propext only); computer-assisted (decide +kernel on generated lists \texttt{closedP1}--\texttt{closedP8}); no hypotheses\\
Cor.~\ref{cb:cor:c2} & written step & written step; computer-assisted (a consequence of Theorem~\ref{cb:thm:c2})\\
Rem.~\ref{cb:rem:c2-scope} & -- & commentary; scope (refutation of the listed subcases, covers, reductions A and B not established)\\
\addlinespace
\multicolumn{3}{@{}l}{\normalsize\textbf{Section~\ref{sa:sec}}}\\[2pt]
Prop.~\ref{sa:prop:rot} & \hangindent=1em\hangafter=1 \lean{A1.c3_reduce}\newline\dirn{AUDIT3/A1_reduction/SizeArith.lean:16}\par \hangindent=1em\hangafter=1 \lean{A1.c5_reduce}\newline\dirn{AUDIT3/A1_reduction/SizeArith.lean:19}\par \hangindent=1em\hangafter=1 \lean{A1.c4_reduce}\newline\dirn{AUDIT3/A1_reduction/SizeArith.lean:22}\par \hangindent=1em\hangafter=1 \lean{A1.c6_reduce}\newline\dirn{AUDIT3/A1_reduction/SizeArith.lean:25}\par \hangindent=1em\hangafter=1 \lean{A1.c2_reduce}\newline\dirn{AUDIT3/A1_reduction/SizeArith.lean:28} & linear arithmetic over N (omega in Lean); core Lean, no Mathlib\\
Prop.~\ref{sa:prop:mirror} & \hangindent=1em\hangafter=1 \lean{A1.mirror_reduce_b0}\newline\dirn{AUDIT3/A1_reduction/SizeArith.lean:31}\par \hangindent=1em\hangafter=1 \lean{A1.mirror_reduce_b12}\newline\dirn{AUDIT3/A1_reduction/SizeArith.lean:34} & linear arithmetic over N (omega in Lean); the hypothesis $1\le b$ of \texttt{mirror\_\allowbreak{}reduce\_\allowbreak{}b12} is not needed\\
Prop.~\ref{sa:prop:order} & \hangindent=1em\hangafter=1 \lean{A1.even_order_reduce}\newline\dirn{AUDIT3/A1_reduction/SizeArith.lean:37}\par \hangindent=1em\hangafter=1 \lean{A1.mult3_order_reduce}\newline\dirn{AUDIT3/A1_reduction/SizeArith.lean:42}\par \hangindent=1em\hangafter=1 \lean{A1.mult5_order_reduce}\newline\dirn{AUDIT3/A1_reduction/SizeArith.lean:47} & \texttt{Nat.\allowbreak{}mul\_\allowbreak{}assoc}, then omega with $t\cdot k$ as one unknown\\
Lem.~\ref{sa:lem:split} & \hangindent=1em\hangafter=1 \lean{A1.order_split}\newline\dirn{AUDIT3/A1_reduction/SizeArith.lean:52} & omega, then decide on the five closed cases $s=$ 2--6: side facts inside the proof of a statement over all $s\ge2$, which is not a finite computation, and checked by hand in the written proof; not computer-assisted\\
Prop.~\ref{sa:prop:table} & \hangindent=1em\hangafter=1 \lean{A1.table_c3}\newline\dirn{AUDIT3/A1_reduction/SizeArith.lean:61}\par \hangindent=1em\hangafter=1 \lean{A1.c3_no32}\newline\dirn{AUDIT3/A1_reduction/SizeArith.lean:62}\par \hangindent=1em\hangafter=1 \lean{A1.table_c5}\newline\dirn{AUDIT3/A1_reduction/SizeArith.lean:65}\par \hangindent=1em\hangafter=1 \lean{A1.c5_gap}\newline\dirn{AUDIT3/A1_reduction/SizeArith.lean:66}\par \hangindent=1em\hangafter=1 \lean{A1.table_c4}\newline\dirn{AUDIT3/A1_reduction/SizeArith.lean:69}\par \hangindent=1em\hangafter=1 \lean{A1.table_c4_only_c2}\newline\dirn{AUDIT3/A1_reduction/SizeArith.lean:70}\par \hangindent=1em\hangafter=1 \lean{A1.table_c6}\newline\dirn{AUDIT3/A1_reduction/SizeArith.lean:73}\par \hangindent=1em\hangafter=1 \lean{A1.table_c2}\newline\dirn{AUDIT3/A1_reduction/SizeArith.lean:76}\par \hangindent=1em\hangafter=1 \lean{A1.table_mirror}\newline\dirn{AUDIT3/A1_reduction/SizeArith.lean:79} & linear arithmetic over N (omega in Lean); the bounds that the case theorems give are hypotheses, not part of the statements: on k in \texttt{table\_\allowbreak{}c3}, \texttt{table\_\allowbreak{}c4}, \texttt{table\_\allowbreak{}c2}; on 2k in \texttt{table\_\allowbreak{}c4\_\allowbreak{}only\_\allowbreak{}c2}, \texttt{table\_\allowbreak{}c6}; on 5k+e in \texttt{table\_\allowbreak{}c5}; on k for $b=0$ and for $1\le b$ in \texttt{table\_\allowbreak{}mirror}; \texttt{c3\_\allowbreak{}no32} and \texttt{c5\_\allowbreak{}gap} assume only $e\le1$\\
Prop.~\ref{sa:prop:assembled} & \hangindent=1em\hangafter=1 \lean{halfturn_pairs}\newline\dirn{AUDIT3/A1_reduction/verify_A1-F1/AssembledSizes.lean:6}\par \hangindent=1em\hangafter=1 \lean{threefold_orbits}\newline\dirn{AUDIT3/A1_reduction/verify_A1-F1/AssembledSizes.lean:9}\par \hangindent=1em\hangafter=1 \lean{fivefold_size}\newline\dirn{AUDIT3/A1_reduction/verify_A1-F1/AssembledSizes.lean:12}\par \hangindent=1em\hangafter=1 \lean{mirror_b0}\newline\dirn{AUDIT3/A1_reduction/verify_A1-F1/AssembledSizes.lean:15}\par \hangindent=1em\hangafter=1 \lean{mirror_b12}\newline\dirn{AUDIT3/A1_reduction/verify_A1-F1/AssembledSizes.lean:18}\par \hangindent=1em\hangafter=1 \lean{fourfold_via_halfturn}\newline\dirn{AUDIT3/A1_reduction/verify_A1-F1/AssembledSizes.lean:21} & verifier file of finding A1-F1; restatements of six \texttt{SizeArith.\allowbreak{}lean} statements, or parts of them (\texttt{fivefold\_\allowbreak{}size}: first conclusion of \texttt{c5\_\allowbreak{}reduce}; \texttt{fourfold\_\allowbreak{}via\_\allowbreak{}halfturn}: third conclusion of \texttt{c4\_\allowbreak{}reduce}), with explicit n (omega in Lean)\\
Lem.~\ref{sa:lem:big-order} & \hangindent=1em\hangafter=1 \lean{big_order_orbit}\newline\dirn{AUDIT3/A1_reduction/verify_A1-F1/AssembledSizes.lean:24} & hypothesis hp $(7\le p)$ unused (linter warning in the log); case split on k, then omega\\
Rem.~\ref{sa:rem:big-order} & -- & the hypothesis $7\le p$ (binder hp) of \texttt{big\_\allowbreak{}order\_\allowbreak{}orbit} is unused (linter warning in \texttt{AssembledSizes.\allowbreak{}lean.\allowbreak{}log}); the statement does not require p to be prime\\
Rem.~\ref{sa:rem:control} & \hangindent=1em\hangafter=1 --\newline\dirn{AUDIT3/A1_reduction/verify_A1-F1/AssembledSizes.lean:30} & anonymous example (no declaration name): the half-turn bound fails at 32 points ($n=32$, $k=16$); accepted by Lean (no error in \texttt{AssembledSizes.\allowbreak{}lean.\allowbreak{}log}) but not covered by \#print axioms\\
Thm.~\ref{sa:thm:baek-seven} & \hangindent=1em\hangafter=1 \lean{baek_seven}\newline\dirn{ES7SPLIT/review/equivalence/BaekSeven.lean:6} & instance $n=4$ of the cited Theorem~\ref{mb:thm:baek}; decide only for the side fact $C(6,2)+2=17$, also checked by hand; the statement is not a finite computation, so not computer-assisted; abstract configurations only, no bridge to point sets\\
Rem.~\ref{sa:rem:baek-seven} & -- & compares \texttt{baek\_\allowbreak{}seven} with the \#check @main output (\texttt{BaekSeven.\allowbreak{}out} lines 1-2) and Theorem~\ref{mb:thm:baek} ($n=4$; universe); abstract configurations only; the ES7SPLIT remark that the step to point sets is not Lean-checked for $n=7$ predates \texttt{MirrorOdd.\allowbreak{}capacity16} (\texttt{MIRRORETV/\allowbreak{}lean/\allowbreak{}MirrorMain.\allowbreak{}lean:1529}), used in Theorem~\ref{mb:thm:seven}\\
Lem.~\ref{sa:lem:model} & written step & written step: orbit counts $n=s\cdot k+e\ (e\le1$, $e=0$ for even s) and $n=2k+b\ (b\le2)$, the modelling of the \texttt{SizeArith.\allowbreak{}lean} header (not formalised there); related Lean-checked forms elsewhere: Theorem~\ref{ec:thm:compat} (a), (b); Lemma~\ref{sr:lem:orbit}; Theorem~\ref{ll:thm:layers} and Theorem~\ref{ll:thm:card-peel}\\
Rem.~\ref{sa:rem:use} & -- & written step: how Lemma~\ref{sa:lem:model}, the arithmetic and the five case classes (hypotheses) combine; \texttt{HalfTurn34}, Threefold33, Fivefold35 not proved; mirror classes Lean-checked (Theorem~\ref{mb:thm:classes}); Lean-checked forms of the implication Theorem~\ref{sr:thm:main} and Theorem~\ref{ar:thm:main}; table bounds conditional except the mirror row (Theorem~\ref{mb:thm:refl32}); conditional forms Theorem~\ref{ec:thm:fourfold}, Theorem~\ref{ec:thm:fivefold}\\
\addlinespace
\end{longtable}}

\endgroup

% bibliography.tex -- the works cited in this document. The sections cite them inline (see Works cited); the
% front matter (Scope of Stage 1) cites every entry by key. Entries are copied from the project's paper P1
% (WRITEUP/P1/p1.tex) where it has them; Roc70 and Rud76 are not in P1.
\begin{thebibliography}{99}\small
\setlength{\itemsep}{1pt}\setlength{\parsep}{0pt}
\bibitem{Bae25} J.~Baek, On the Erd\H{o}s--Tuza--Valtr conjecture, \emph{European J. Combin.} 124 (2025), 104085;
arXiv:2206.04260. Theorem numbers refer to the arXiv version.
\bibitem{Lean4} L.~de Moura and S.~Ullrich, The Lean~4 theorem prover and programming language, in \emph{CADE-28},
LNCS 12699, Springer, 2021, 625--635.
\bibitem{DT} drat-trim, \url{https://github.com/marijnheule/drat-trim}: issue \#31 (22 July 2022), commits
\texttt{cbd2915} (19 October 2022, fix) and \texttt{5cd8ac3} (21 April 2024, regression), issue \#42 and pull request
\#43 (29 September 2026); master at \texttt{2e3b2dc}.
\bibitem{Dum25} B.~Dumitru, Notes on the 33-point Erd\H{o}s--Szekeres problem, arXiv:2512.24061 (2025).
\bibitem{ES35} P.~Erd\H{o}s and G.~Szekeres, A combinatorial problem in geometry, \emph{Compositio Math.} 2 (1935),
463--470.
\bibitem{Mathlib} The mathlib Community, The Lean mathematical library, in \emph{CPP 2020}, ACM, 2020, 367--381.
\bibitem{Roc70} R.~T.~Rockafellar, \emph{Convex Analysis}, Princeton University Press, Princeton, 1970.
\bibitem{Rud76} W.~Rudin, \emph{Principles of Mathematical Analysis}, 3rd ed., McGraw--Hill, New York, 1976.
\bibitem{Sch26} M.~Schroeder, Erd\H{o}s--Szekeres: upper bounds for convex heptagons and octagons, archive version~1,
19 September 2026, doi:10.5281/zenodo.22842891 (unreviewed).
\bibitem{SMQH25} B.~Subercaseaux, E.~Mackey, L.~Qian and M.~J.~H.~Heule, Automated symmetric constructions in discrete
geometry, arXiv:2506.00224v1 (2025); also in \emph{Intelligent Computer Mathematics (CICM 2025)}, LNCS 16136 (LNAI),
Springer, 2026 (published October 2025), 29--47,
doi:10.1007/978-3-032-07021-0\_3 (Crossref metadata; chapter text not accessed).
\bibitem{CakeLPR} Y.~K.~Tan, M.~J.~H.~Heule and M.~O.~Myreen, Verified propagation redundancy and compositional UNSAT
checking in CakeML, \emph{Int. J. Softw. Tools Technol. Transf.} 25 (2023), 167--184; conference version
\emph{cake\_lpr: verified propagation redundancy checking in CakeML}, TACAS 2021, LNCS 12652, 223--241.
\end{thebibliography}


\end{document}
